% Gerado automaticamente. Nao editar.
\providecommand{\fvMeshCW}[3]{
  \draw[->,loopcol,line width=1.1pt] (#1)
    arc[start angle=0,end angle=-305,radius=#2cm];
  \node[loopcol,fill=white,inner sep=1pt] at (#1){\footnotesize #3};}

\section{Eletrostática}

\subsection{Cargas elétricas}
\stepcounter{questao}
\resposta{(c)}
\resolucao{A conservação da carga afirma que, em um sistema isolado, cargas podem
ser transferidas de um corpo para outro, mas a soma algébrica permanece a mesma.
As alternativas (a) e (b) descrevem a \emph{lei de atração e repulsão}, e não a
conservação; (d) e (e) são falsas.}
\stepcounter{questao}
\resposta{(a)}
\resolucao{Na eletrização por atrito há apenas \emph{transferência} de elétrons de
um corpo para o outro. Pela conservação da carga, o que um perde o outro ganha:
as cargas finais têm o mesmo módulo e sinais opostos.}
\stepcounter{questao}
\resposta{(d)}
\resolucao{Mesma justificativa da questão anterior. Na série triboelétrica o vidro
cede elétrons à lã: o vidro fica positivo e a lã, negativa.}
\stepcounter{questao}
\resposta{(c)}
\resolucao{A lã (posição 3) está acima do PVC (posição 5), logo cede elétrons ao PVC.
O PVC \emph{ganha} elétrons (fica negativo) e a lã \emph{perde} elétrons (fica positiva).}
\stepcounter{questao}
\resposta{(a) sinais opostos; (b) negativo}
\resolucao{(a) A atração ocorre entre cargas de sinais contrários (ou entre corpo
eletrizado e corpo neutro; como o enunciado fala em \emph{cargas}, os sinais são
opostos).\\
(b) $Q_1$ é repelida por $Q_3>0$, logo $Q_1>0$. Como $Q_1$ e $Q_2$ se atraem,
$Q_2<0$.}
\stepcounter{questao}
\resposta{Apenas (a)}
\resolucao{(a) Correta. (b) Errada: o corpo neutro possui igual número de prótons
e elétrons, não ausência de cargas. (c) Errada: prótons estão presos ao núcleo;
nos sólidos só os elétrons se movem. (d) Errada: por indução, um corpo neutro é
atraído por um corpo eletrizado. (e) Errada: materiais iguais têm a mesma
tendência de ceder ou receber elétrons, e o atrito não produz eletrização
apreciável.}
\stepcounter{questao}
\resposta{(a)}
\resolucao{Gotículas de mesmo sinal se repelem, o que espalha o jato e evita
acúmulos; simultaneamente, todas são atraídas pela peça de sinal oposto,
inclusive contornando-a (efeito \emph{wrap-around}). O aproveitamento da tinta
passa de cerca de \SI{50}{\percent} para mais de \SI{90}{\percent}.}
\stepcounter{questao}
\resposta{(a) falta de elétrons; (b) $n = \pot{3.0}{14}$ elétrons}
\resolucao{(a) A carga final é positiva, logo o corpo \emph{perdeu} elétrons
(ficou com falta).\\
(b) Da quantização da carga, $Q = n\,e$:
\[
n=\frac{Q}{e}=\frac{\SI{48e-6}{\coulomb}}{\SI{1.6e-19}{\coulomb}}
 = \pot{3.0}{14}\ \text{elétrons}.
\]}
\stepcounter{questao}
\resposta{(e)}
\resolucao{A atração é compatível com duas situações: objeto negativo, ou objeto
\emph{neutro} polarizado por indução. Como não há informação adicional, a única
afirmação segura é que ele \emph{pode} estar neutro --- por isso (a) e (d) estão
erradas. A polarização ocorre tanto em condutores quanto em isolantes, o que
elimina (b) e (c).}
\stepcounter{questao}
\resposta{(d)}
\resolucao{A barra negativa \emph{repele} os elétrons livres do metal, que migram
para a extremidade \emph{mais distante}, B. Logo B fica negativa e A, positiva
--- o que elimina (a), (b) e (c). Como não houve aterramento nem contato, a carga
total de AB continua nula, eliminando (e). O fenômeno descrito é exatamente a
indução eletrostática.}
\stepcounter{questao}
\resposta{(d)}
\resolucao{O corpo positivo \emph{atrai} os elétrons livres: X fica negativa e Y,
positiva --- o que torna (a), (b) e (c) falsas. Como a extremidade negativa X está
\emph{mais próxima} da barra do que a positiva Y, a atração supera a repulsão e
existe força resultante de atração entre os corpos.}
\stepcounter{questao}
\resposta{(b)}
\resolucao{Com as esferas em contato, elas formam um único condutor: o bastão
positivo atrai os elétrons para $x$ e deixa $y$ com falta de elétrons.
\textbf{Separando as esferas com o bastão ainda presente}, essa separação de
cargas fica "aprisionada": $x$ permanece negativa e $y$, positiva.
(Se o bastão fosse afastado \emph{antes} da separação, ambas voltariam a ser
neutras.)}
\stepcounter{questao}
\resposta{(b)}
\resolucao{O bastão negativo repele os elétrons livres do disco, que descem para
as folhas. O disco fica \emph{positivo} e as folhas ficam \emph{negativas}; como
ambas têm o mesmo sinal, elas se repelem e se abrem. A carga total do
eletroscópio, porém, continua nula --- o que elimina (a) e (d).}
\stepcounter{questao}
\resposta{(c)}
\resolucao{A \emph{repulsão} só ocorre entre corpos eletrizados de mesmo sinal.
Como cada bola repele alguma outra, cada uma das três precisa estar eletrizada:
apenas a hipótese III explica o fenômeno. Um arranjo possível é $(+,+,-)$.}
\stepcounter{questao}
\resposta{(d)}
\resolucao{Se as três estivessem carregadas, haveria pelo menos um par de mesmo
sinal (princípio da casa dos pombos com dois sinais e três esferas), gerando
repulsão --- III fica descartada. Com apenas uma carregada (I), ela atrai as duas
neutras por indução e as neutras se atraem mutuamente por polarização mútua.
Com duas carregadas de sinais opostos (II), elas se atraem entre si e cada uma
atrai a neutra. Logo I e II são possíveis.}
\stepcounter{questao}
\resposta{(d)}
\resolucao{Cargas de mesmo sinal se repeliriam, o que elimina (a) e (c).
Em (b), as duas de mesmo sinal se repeliriam. Em (e), as duas esferas neutras
não interagiriam apreciavelmente entre si. Resta (d): as duas eletrizadas se
atraem por terem sinais opostos, e cada uma atrai a neutra por indução.}
\stepcounter{questao}
\resposta{(e)}
\resolucao{Duas condições devem ser satisfeitas simultaneamente:
\textbf{(i)} bastão e papel, eletrizados por atrito, têm sinais \emph{opostos};
\textbf{(ii)} bastão e pêndulo se repelem, logo têm sinais \emph{iguais}.
Somente a alternativa (e) satisfaz ambas (papel $+$, bastão $-$, pêndulo $-$).
Em (a) o par bastão/papel tem sinais iguais; em (b) e (d) o bastão e o pêndulo
têm sinais opostos; em (c) o par bastão/papel tem sinais iguais.}
\stepcounter{questao}
\resposta{(e)}
\resolucao{As esferas são idênticas, então a carga total se divide igualmente:
\[
Q'_A=Q'_B=\frac{Q_A+Q_B}{2}=\frac{(+5)+(-1)}{2}=\SI{+2}{\coulomb}.
\]
A esfera A passou de $+5$ para $+2$: ela \emph{ganhou} \SI{3}{\coulomb} em
elétrons. Como os elétrons são negativos, eles se deslocam \emph{para} o corpo de
maior potencial, ou seja, de B para A.}
\stepcounter{questao}
\resposta{(c)}
\resolucao{Carga final de cada esfera:
$Q' = \dfrac{3+1}{2} = \SI{+2}{\micro\coulomb}$.
A esfera K passou de $+3$ para $+2\,\si{\micro\coulomb}$: recebeu
$\SI{-1}{\micro\coulomb}$ (isto é, elétrons). Portanto os elétrons vão de L para K
transportando $\SI{-1}{\micro\coulomb}$. O que impulsiona o fluxo é a diferença de
\emph{potencial}, e não o sinal das cargas --- por isso (a) é falsa.}
\stepcounter{questao}
\resposta{(e)}
\resolucao{$Q' = \dfrac{5Q + (-3Q)}{2} = \dfrac{2Q}{2} = Q$ em cada esfera.}
\stepcounter{questao}
\resposta{(e)}
\resolucao{Após o contato, cada esfera fica com
$Q' = \dfrac{\pot{8.0}{-8}+0}{2} = \pot{4.0}{-8}\,\si{\coulomb}$.
A carga transferida tem módulo $\Delta Q = \pot{4.0}{-8}\,\si{\coulomb}$, logo
\[
n=\frac{\Delta Q}{e}=\frac{\pot{4.0}{-8}}{\pot{1.6}{-19}}=\pot{2.5}{11}
\ \text{elétrons}.
\]
Note que a esfera eletrizada é positiva: os elétrons vão da esfera neutra para ela.}
\stepcounter{questao}
\resposta{Não: a carga não é múltipla inteira de $e$.}
\resolucao{Pela quantização, toda carga deve satisfazer $Q=n\,e$ com $n$ inteiro.
Aqui,
\[
n=\frac{\pot{2.4}{-19}}{\pot{1.6}{-19}}=1{,}5,
\]
que não é inteiro. Não existe carga livre igual a $1{,}5\,e$, logo a medida está
incorreta (erro experimental ou de cálculo).}
\stepcounter{questao}
\resposta{(d)}
\resolucao{\textbf{1\textsuperscript{o} contato (A--B):}
$Q_A = Q_B = \dfrac{Q+0}{2} = \dfrac{Q}{2}$.\\
\textbf{2\textsuperscript{o} contato (A--C):}
$Q_A = Q_C = \dfrac{Q/2+0}{2} = \dfrac{Q}{4}$.\\
Situação final: $A=\dfrac{Q}{4}$, $B=\dfrac{Q}{2}$, $C=\dfrac{Q}{4}$.
O erro mais comum é esquecer que B \emph{conserva} a carga obtida no primeiro
contato.}
\stepcounter{questao}
\resposta{(a)}
\resolucao{A cada contato com uma esfera \emph{neutra} e idêntica, a carga de D
é dividida por 2:
\[
Q \xrightarrow{\text{A}} \frac{Q}{2}
  \xrightarrow{\text{B}} \frac{Q}{4}
  \xrightarrow{\text{C}} \frac{Q}{8}.
\]
No último contato, C fica com o mesmo valor de D: ambas com $\dfrac{Q}{8}$.
De modo geral, após $n$ contatos com esferas neutras idênticas,
$Q_D = \dfrac{Q}{2^{\,n}}$.}
\stepcounter{questao}
\resposta{(e)}
\resolucao{\textbf{Contato X--Y:} $Q_X = Q_Y = \dfrac{0+Q}{2} = \dfrac{Q}{2}$.\\
\textbf{Contato X--Z:} agora X tem $+\dfrac{Q}{2}$ e Z tem $-Q$:
\[
Q_X=\frac{\dfrac{Q}{2}+(-Q)}{2}=\frac{-\dfrac{Q}{2}}{2}=-\frac{Q}{4}.
\]
Observe a inversão de sinal: X termina \emph{negativa}. A alternativa (a) é a
armadilha típica de quem soma as cargas iniciais de Y e Z.}
\stepcounter{questao}
\resposta{(b)}
\resolucao{\textbf{Contato C--A:}
$Q_C = Q_A = \dfrac{0+16}{2} = \SI{8}{\micro\coulomb}$.\\
\textbf{Contato C--B:}
$Q_C = Q_B = \dfrac{8+4}{2} = \SI{6}{\micro\coulomb}$.\\
Situação final: $A=\SI{8}{\micro\coulomb}$, $B=Q_C=\SI{6}{\micro\coulomb}$.
Confira a conservação: $16+4+0 = 8+6+6 = \SI{20}{\micro\coulomb}$.}
\stepcounter{questao}
\resposta{(c)}
\resolucao{\textbf{Contato A--B:}
$q_A = q_B = \dfrac{0+6}{2} = \SI{3}{\micro\coulomb}$.\\
\textbf{Contato A--C:}
$q_A = q_C = \dfrac{3+7}{2} = \SI{5}{\micro\coulomb}$.\\
Final: $A=\SI{5}{\micro\coulomb}$, $B=\SI{3}{\micro\coulomb}$,
$C=\SI{5}{\micro\coulomb}$.
Verificação: $0+6+7 = 5+3+5 = \SI{13}{\micro\coulomb}$.}
\stepcounter{questao}
\resposta{$Q_1 = \dfrac{5Q}{3}$ (raio $R$) e $Q_2 = \dfrac{10Q}{3}$ (raio $2R$)}
\resolucao{Sejam $Q_1$ e $Q_2$ as cargas finais. Aplicando as duas condições:
\[
\textbf{(i) conservação: } Q_1+Q_2 = 3Q+2Q = 5Q,
\qquad
\textbf{(ii) proporcionalidade: } \frac{Q_1}{R}=\frac{Q_2}{2R}
\ \Rightarrow\ Q_2 = 2Q_1 .
\]
Substituindo (ii) em (i): $Q_1 + 2Q_1 = 5Q \Rightarrow Q_1 = \dfrac{5Q}{3}$ e
$Q_2 = \dfrac{10Q}{3}$.
Fisicamente, a proporcionalidade decorre de as esferas ficarem com o \emph{mesmo
potencial} após o contato: $V = \dfrac{\ke Q}{R}$ igual para ambas.}
\stepcounter{questao}
\resposta{(a) $Q_n = Q_0/2^{\,n}$; \quad (b) $n = 7$}
\resolucao{(a) Cada contato com uma esfera idêntica e neutra reparte a carga em
duas partes iguais, logo $Q_1 = Q_0/2$, $Q_2 = Q_0/4$ e, por indução,
\[
\boxed{\,Q_n = \frac{Q_0}{2^{\,n}}\,}.
\]
(b) Queremos $\dfrac{1}{2^{\,n}} < 0{,}01 \Rightarrow 2^{\,n} > 100$. Como
$2^6 = 64$ e $2^7 = 128$, são necessários $n = 7$ contatos.
Repare que a carga nunca se anula: ela apenas tende a zero, o que ilustra por
que o aterramento (contato com um condutor "infinito") é a forma prática de
descarregar um corpo.}
\stepcounter{questao}
\resposta{(a) \SI{10}{\micro\coulomb}; (b) $\pot{6.25}{13}$;
(c) $V=\SI{360}{\kilo\volt}$, $E=\SI{1.44}{\mega\volt\per\meter}$ --- sem descarga}
\resolucao{(a) $Q = I\,t = \pot{1}{-6}\cdot10 = \SI{10}{\micro\coulomb}$.\\
(b) $n = \dfrac{Q}{e} = \dfrac{\pot{1}{-5}}{\pot{1.6}{-19}} = \pot{6.25}{13}$
elétrons --- \emph{removidos} da cúpula (que fica positiva).\\
(c) Para uma esfera condutora,
\[
V=\frac{\ke Q}{R}=\frac{9\cdot10^9\cdot\pot{1}{-5}}{0{,}25}=\SI{360}{\kilo\volt},
\qquad
E_{\text{sup}}=\frac{\ke Q}{R^{2}}=\frac{V}{R}=\frac{\pot{3.6}{5}}{0{,}25}
=\SI{1.44}{\mega\volt\per\meter}.
\]
Como $\SI{1.44}{} < \SI{3}{\mega\volt\per\meter}$, \textbf{não há ruptura} --- a
carga continua se acumulando.\\
\textbf{Limite do gerador:} a descarga ocorreria quando
$E_{\text{sup}} = \SI{3}{\mega\volt\per\meter}$, ou seja
$V_{\max}=E\,R = 3\cdot10^6\cdot0{,}25=\SI{750}{\kilo\volt}$.
\textbf{Consequência de projeto:} como $V_{\max} \propto R$, cúpulas
\emph{maiores} sustentam tensões maiores --- por isso os geradores de demonstração
têm esferas polidas e volumosas, sem pontas (onde $E$ dispararia).}
\stepcounter{questao}
\resposta{(a) $\approx\pot{9.5}{21}$; (b) $\approx\SI{1517}{\coulomb}$;
(c) razão $\sim10^9$}
\resolucao{(a) O número de mols é $n=\dfrac{1}{63{,}5}$ mol, logo
\[
N=\frac{1}{63{,}5}\cdot\pot{6.02}{23}\approx\pot{9.5}{21}\ \text{elétrons livres}.
\]
(b) $|Q| = N\,e = \pot{9.5}{21}\cdot\pot{1.6}{-19}\approx\SI{1517}{\coulomb}$.\\
(c) A razão entre essa carga e a obtida por atrito é
\[
\frac{1517}{\pot{1}{-6}}\approx\pot{1.5}{9},
\]
ou seja, cerca de \textbf{um bilhão e meio de vezes} maior.

\textbf{Interpretação --- por que isso importa.} Eletrizar um corpo por atrito
remove ou acrescenta uma fração absolutamente ínfima dos elétrons disponíveis
(da ordem de $10^{-9}$). O material permanece, para todos os efeitos,
\emph{eletricamente neutro em sua estrutura}; a "carga" que percebemos é apenas
um desequilíbrio superficial minúsculo.\\
Isso explica dois fatos:
\begin{itemize}[leftmargin=1.6em,itemsep=1pt,topsep=2pt]
\item por que um condutor pode fornecer corrente indefinidamente sem "acabar" os
      elétrons;
\item por que forças elétricas macroscópicas são tão intensas --- se
      conseguíssemos separar \emph{toda} a carga de $\SI{1}{\gram}$ de cobre e
      colocá-la a $\SI{1}{\meter}$ de distância de outra igual, a força seria da
      ordem de $10^{16}$\,N, milhões de vezes o peso de um transatlântico.
\end{itemize}}
\stepcounter{questao}
\resposta{(a) $Q_A=\dfrac{Q}{12}$, $Q_B=\dfrac{2Q}{3}$, $Q_C=\dfrac{Q}{4}$;
(b) soma $=Q$; (c) $Q_A^{(n)}=Q\displaystyle\prod_{j=2}^{n}\frac{1}{j+1}$}
\resolucao{\textbf{Contato A--B.} A carga total $Q$ se reparte proporcionalmente
aos raios ($R$ e $2R$), ou seja, na razão $1:2$:
\[
Q_A=\frac{R}{R+2R}\,Q=\frac{Q}{3},
\qquad
Q_B=\frac{2R}{3R}\,Q=\frac{2Q}{3}.
\]
\textbf{Contato A--C.} Agora $A$ traz $\dfrac{Q}{3}$ e encontra $C$ (raio $3R$,
neutra). A repartição é na razão $R:3R = 1:3$:
\[
Q_A=\frac{1}{4}\cdot\frac{Q}{3}=\frac{Q}{12},
\qquad
Q_C=\frac{3}{4}\cdot\frac{Q}{3}=\frac{Q}{4}.
\]
\textbf{(b) Conservação:}
$\dfrac{Q}{12}+\dfrac{2Q}{3}+\dfrac{Q}{4}
=\dfrac{Q+8Q+3Q}{12}=\dfrac{12Q}{12}=Q$. \checkmark

\textbf{(c) Generalização.} Ao tocar uma esfera de raio $jR$, a esfera $A$ retém a
fração $\dfrac{R}{R+jR}=\dfrac{1}{1+j}$ da carga que possuía. Logo, após a
sequência $2R, 3R, \dots, nR$:
\[
Q_A^{(n)} = Q\prod_{j=2}^{n}\frac{1}{1+j}
= Q\cdot\frac{1}{3}\cdot\frac{1}{4}\cdots\frac{1}{n+1}
= \frac{2\,Q}{(n+1)!}.
\]
Confira para $n=3$: $\dfrac{2Q}{4!}=\dfrac{2Q}{24}=\dfrac{Q}{12}$. \checkmark\\
\textbf{Por que isso importa:} a carga de $A$ decai \emph{fatorialmente} --- muito
mais rápido que o decaimento $2^{-n}$ das esferas idênticas. É o princípio do
\emph{gerador de Van de Graaff}: transferir carga para um condutor de raio muito
maior esvazia quase completamente o menor.}
\stepcounter{questao}
\resposta{$Q=+2e=\pot{3.2}{-19}\,\si{\coulomb}$ nos dois instantes}
\resolucao{(a) \textbf{Antes:} $Q=5(+e)+3(-e)=+2e=\pot{3.2}{-19}\,\si{\coulomb}$.\\
\textbf{Depois:} dois pares próton--elétron formam átomos neutros, restando
$3$ prótons livres e $1$ elétron livre:
\[
Q=3(+e)+1(-e)=+2e=\pot{3.2}{-19}\,\si{\coulomb}.
\]
(b) \textbf{Nada mudou} --- e não poderia mudar. A reorganização apenas
\emph{redistribuiu} as cargas dentro do sistema; nenhuma foi criada nem
destruída, e o sistema é isolado.\\
\textbf{Ponto conceitual:} um corpo \emph{neutro} não é um corpo \emph{sem
carga} --- ele tem tantas cargas positivas quanto negativas. Agrupar cargas em
átomos neutros esconde-as, mas não as elimina do balanço.}
\stepcounter{questao}
\resposta{$n=\pot{3.1}{10}$ elétrons; $\Delta m=\pot{2.8}{-20}\,\si{\kilogram}$}
\resolucao{(a)
\[
n=\frac{|Q|}{e}=\frac{\pot{5}{-9}}{\pot{1.6}{-19}}=\pot{3.125}{10}.
\]
Trinta bilhões de elétrons --- número enorme em termos absolutos.\\
(b)
\[
\Delta m=n\,m_e=\pot{3.125}{10}\times\pot{9.11}{-31}
=\pot{2.85}{-20}\,\si{\kilogram}.
\]
Comparando com os $\SI{20}{\gram}$ do bastão:
\[
\frac{\Delta m}{m}=\frac{\pot{2.85}{-20}}{\pot{2}{-2}}=\pot{1.4}{-18}.
\]
\textbf{Conclusão:} a eletrização por atrito \emph{não} altera a massa de forma
mensurável --- o desvio está $15$ ordens de grandeza abaixo da melhor balança
analítica. É por isso que a massa não serve como indicador de eletrização, e
por que o fenômeno passou séculos sendo descrito apenas por seus efeitos de
força.}
\stepcounter{questao}
\resposta{$Q_A=Q/2^{n}$; $Q_{B_k}=Q/2^{k}$; a soma reconstitui $Q$}
\resolucao{(a) Cada contato entre esferas \emph{idênticas} reparte igualmente a
carga presente. Como cada $B_k$ chega neutra, $A$ perde metade do que tem:
\[
Q\to\frac{Q}{2}\to\frac{Q}{4}\to\cdots
\;\Longrightarrow\;
Q_A=\frac{Q}{2^{n}}.
\]
(b) A esfera $B_k$ leva metade do que $A$ possuía imediatamente antes, ou seja
$\dfrac{1}{2}\cdot\dfrac{Q}{2^{k-1}}=\dfrac{Q}{2^{k}}$.\\
(c) \textbf{Conservação:}
\[
Q_A+\sum_{k=1}^{n}Q_{B_k}
=\frac{Q}{2^{n}}+Q\sum_{k=1}^{n}\frac{1}{2^{k}}
=\frac{Q}{2^{n}}+Q\left(1-\frac{1}{2^{n}}\right)=Q.\ \checkmark
\]
A soma geométrica $\sum 1/2^{k}$ tende a $1$, e o termo residual de $A$ fecha
exatamente a diferença.\\
\textbf{Leitura:} nenhuma carga se perde --- ela apenas se dilui. Com
$n=10$, $A$ retém menos de $0{,}1\%$ de $Q$, mas os $99{,}9\%$ restantes estão
distribuídos pelas dez esferas retiradas.\\
\textbf{Uso prático:} é assim que se \emph{descarrega} progressivamente um corpo
sem aterrá-lo, e também como se obtêm frações controladas de uma carga de
referência.}
\stepcounter{questao}
\resposta{$Q_{\max}=\SI{3.33}{\micro\coulomb}$; $V=\SI{300}{\kilo\volt}$}
\resolucao{(a) O campo na superfície de uma esfera condutora vale
$E=\ke Q/R^{2}$. Impondo $E=E_{rup}$:
\[
Q_{\max}=\frac{E_{rup}R^{2}}{\ke}=\frac{\pot{3}{6}\times0{,}01}{\pot{9}{9}}
=\pot{3.33}{-6}\,\si{\coulomb}=\SI{3.33}{\micro\coulomb}.
\]
(b)
\[
V=\frac{\ke Q_{\max}}{R}=E_{rup}\,R=\pot{3}{6}(0{,}10)=\SI{300}{\kilo\volt}.
\]
\textbf{Relação notável:} $V_{\max}=E_{rup}R$ --- o potencial máximo é
\emph{proporcional ao raio}. Dobrar o raio dobra a tensão suportável, e é por
isso que geradores de Van de Graaff têm cúpulas grandes e lisas.\\
(c) \textbf{Acima do limite}, o ar ioniza e passa a conduzir: a carga escapa por
descarga corona (o brilho azulado e o chiado característicos) ou por faísca.
A esfera \emph{não consegue} reter mais carga --- ela se autolimita.\\
\textbf{Consequência de projeto:} para armazenar mais carga é preciso aumentar
$R$, usar meio de maior rigidez (óleo, $\mathrm{SF_6}$, vácuo) ou eliminar
irregularidades da superfície, onde o campo local é muito maior --- assunto da
questão seguinte.}
\stepcounter{questao}
\resposta{$Q_R/Q_r=10$; $\sigma_r/\sigma_R=10$;
$E_R=\SI{300}{\kilo\volt\per\meter}$ e $E_r=\SI{3}{\mega\volt\per\meter}$}
\resolucao{(a) Ligadas por fio, as duas ficam no \emph{mesmo potencial}:
\[
V=\frac{\ke Q}{R}\Rightarrow Q\propto R
\;\Longrightarrow\;
\frac{Q_R}{Q_r}=\frac{R}{r}=10.
\]
A densidade superficial é $\sigma=Q/(4\pi R^{2})\propto R/R^{2}=1/R$:
\[
\frac{\sigma_r}{\sigma_R}=\frac{R}{r}=10.
\]
\textbf{A esfera maior tem mais carga, mas a menor tem carga mais
concentrada.}\\
(b) $E=V/R$ na superfície de cada uma:
\[
E_R=\frac{\pot{30}{3}}{0{,}10}=\SI{300}{\kilo\volt\per\meter};
\qquad
E_r=\frac{\pot{30}{3}}{0{,}01}=\SI{3}{\mega\volt\per\meter}.
\]
(c) \textbf{O campo na esfera pequena atingiu exatamente a rigidez do ar.} A
ionização começa \emph{ali}, embora as duas estejam ao mesmo potencial ---
regiões de menor raio de curvatura concentram campo. É o \textbf{poder das
pontas}.\\
\textbf{Aplicações:}
\begin{itemize}
\item \textbf{Para-raios:} a ponta afiada ioniza o ar e estabelece o caminho
preferencial da descarga, protegendo a estrutura.
\item \textbf{Impressão e pintura eletrostática:} eletrodos pontiagudos ionizam
o ar para carregar tinta ou papel.
\end{itemize}
\textbf{E a contrapartida:} em equipamentos de alta tensão evita-se qualquer
ponta ou aresta viva --- daí os terminais arredondados e as blindagens
toroidais.}
\stepcounter{questao}
\resposta{$\approx\pot{2}{23}$ átomos; fração de $\pot{1.6}{-13}$}
\resolucao{(a)
\[
N=\frac{m}{M}N_A=\frac{20}{60}\times\pot{6.02}{23}\approx\pot{2.0}{23}.
\]
(b) Já se sabe que $\pot{3.1}{10}$ elétrons foram transferidos:
\[
f=\frac{\pot{3.1}{10}}{\pot{2.0}{23}}=\pot{1.6}{-13},
\]
ou cerca de \textbf{um átomo em seis trilhões}.\\
(c) \textbf{Duas leituras que se complementam.}\\
\emph{Do ponto de vista microscópico}, a eletrização por atrito é um fenômeno
de superfície absolutamente marginal --- envolve uma fração desprezível da
matéria, e apenas a camada atômica mais externa.\\
\emph{Do ponto de vista macroscópico}, esses mesmos
$\pot{3.1}{10}$ elétrons produzem forças facilmente visíveis, capazes de
levantar pedaços de papel contra o peso.\\
\textbf{O contraste explica a intensidade da força elétrica:} ela é tão maior
que a gravitacional que um desequilíbrio de uma parte em $10^{13}$ já produz
efeitos observáveis. E explica também por que a matéria comum é neutra com
altíssima precisão --- qualquer desequilíbrio maior seria violentamente
desfeito.}
\stepcounter{questao}
\resposta{$\SI{90}{\micro\newton}$ e $\SI{9}{\milli\newton}$;
$q\approx\SI{52}{\nano\coulomb}$}
\resolucao{(a)
\[
F=\frac{\ke q^{2}}{d^{2}}:\quad
q=\SI{5}{\nano\coulomb}\Rightarrow\SI{9.0e-5}{\newton};
\qquad
q=\SI{50}{\nano\coulomb}\Rightarrow\SI{9.0e-3}{\newton}.
\]
O peso vale $P=mg=\pot{1}{-3}(9{,}8)=\SI{9.8}{\milli\newton}$.\\
(b) Impondo $F=P$:
\[
q=d\sqrt{\frac{P}{\ke}}=0{,}05\sqrt{\frac{\pot{9.8}{-3}}{\pot{9}{9}}}
=0{,}05\times\pot{1.04}{-6}=\SI{52}{\nano\coulomb}.
\]
(c) \textbf{Perfeitamente viável.} Cinquenta nanocoulombs correspondem a
$\pot{3.3}{11}$ elétrons --- ordem de grandeza obtida rotineiramente por
atrito, como mostrou a questão anterior.\\
\textbf{É por isso que a eletrostática é observável a olho nu:} um simples
canudo atritado sustenta pedaços de papel contra a gravidade de todo o planeta.
\\
\textbf{Verificação do limite:} essas esferas suportariam essa carga? Para
$R=\SI{5}{\milli\meter}$, o campo na superfície seria
$\ke q/R^{2}=\SI{1.9}{\mega\volt\per\meter}$ --- abaixo da rigidez do ar, mas
já na mesma ordem. Cargas muito maiores escapariam por descarga.}
\stepcounter{questao}
\resposta{Conservada em ambos: $0=+e-e+0$ e $-e+e=0$}
\resolucao{(a) \textbf{Decaimento beta.} Antes: nêutron, carga $0$. Depois:
próton ($+e$), elétron ($-e$) e antineutrino ($0$):
\[
0=(+e)+(-e)+0=0.\ \checkmark
\]
Note que \emph{o número de partículas mudou} e que partículas foram
\emph{criadas} --- o elétron não estava dentro do nêutron. Ainda assim a carga
total se conserva.\\
(b) \textbf{Aniquilação.} Antes: $-e+(+e)=0$. Depois: dois fótons, ambos de
carga nula. \checkmark\\
Aqui a \emph{matéria} desaparece por completo, convertida em radiação --- e a
carga continua conservada, porque já era nula.\\
(c) \textbf{Alcance do princípio.} A conservação da carga é uma das leis mais
robustas da física:
\begin{itemize}
\item Vale em processos \textbf{químicos}, \textbf{nucleares} e de
\textbf{física de partículas}, sem exceção conhecida.
\item Vale mesmo quando partículas são criadas ou destruídas --- o que a
distingue da conservação do número de partículas, que \emph{não} é uma lei.
\item Experimentalmente, o decaimento do elétron (que violaria a conservação) foi
buscado e não observado; o limite atual para sua vida média é superior a
$10^{26}$ anos.
\end{itemize}
\textbf{Origem teórica:} a conservação da carga decorre de uma simetria
fundamental do eletromagnetismo, do mesmo modo que a conservação da energia
decorre da invariância no tempo. Ela não é um fato empírico isolado, mas
consequência da estrutura da teoria.}
\stepcounter{questao}
\resposta{O eletrômetro indica exatamente a carga do corpo, independentemente da
posição}
\resolucao{(a) \textbf{Indução total.} O corpo de carga $+Q$ induz $-Q$ na
superfície \emph{interna} do recipiente --- exatamente $-Q$, porque todas as
linhas de campo que saem do corpo terminam nela. Por conservação, sobra
$+Q$ na superfície \emph{externa}, e é essa carga que o eletrômetro mede.\\
\textbf{Resultado:} a leitura é $+Q$, o valor exato da carga introduzida.\\
(b) \textbf{Independência da posição.} O campo elétrico dentro do metal é nulo
em equilíbrio. Aplicando a lei de Gauss a uma superfície inteiramente contida no
metal, o fluxo é zero --- logo a carga interna total (corpo mais superfície
interna) é nula, \emph{qualquer que seja a posição do corpo}. A carga externa
fica sempre $+Q$, distribuída de acordo apenas com a geometria \emph{externa}
do recipiente.\\
\textbf{Consequência prática:} não é preciso centralizar o corpo nem controlar
sua orientação --- o que torna o copo de Faraday o método padrão de medição
absoluta de carga.\\
(c) \textbf{Se o corpo tocar a parede:} a carga $+Q$ neutraliza a
$-Q$ induzida, o corpo sai neutro e o recipiente fica com $+Q$ na superfície
externa. \textbf{A leitura do eletrômetro não muda} --- continua $+Q$.\\
Isso permite \emph{transferir} integralmente a carga de um corpo para o
recipiente, e é a base do método clássico de acumulação sucessiva de carga em
geradores eletrostáticos.}
\stepcounter{questao}
\resposta{Blinda o interior contra campos externos; o aterramento é necessário
para blindar o exterior contra fontes internas}
\resolucao{(a) \textbf{Campo externo não penetra.} As cargas livres do condutor
se redistribuem até anular o campo em seu interior --- é a condição de
equilíbrio eletrostático. Qualquer campo externo induz uma distribuição
superficial cujo campo cancela exatamente o campo aplicado na região interna.\\
Isso vale \emph{mesmo sem aterramento}: a gaiola de Faraday isolada blinda o
interior. É o efeito que protege passageiros de um carro atingido por raio.\\
(b) \textbf{Sentido inverso: só com aterramento.} Uma carga $+Q$ colocada
\emph{dentro} da gaiola induz $-Q$ na face interna e $+Q$ na externa --- e essa
carga externa produz campo do lado de fora. \textbf{A blindagem falha.}\\
Com a malha \textbf{aterrada}, a carga $+Q$ da face externa escoa para a terra,
e o campo externo desaparece. Só então a blindagem funciona nos dois sentidos.\\
(c) \textbf{Resumo da diferença.}
\begin{center}
\begin{tabular}{lcc}
\hline
& não aterrada & aterrada\\
\hline
Protege o interior de campos externos & sim & sim\\
Impede que fontes internas afetem o exterior & não & sim\\
\hline
\end{tabular}
\end{center}
\textbf{Limitações práticas:} a malha só é eficaz para aberturas muito menores
que o comprimento de onda do sinal a blindar --- em eletrostática (campo
constante), qualquer malha fina basta; em radiofrequência, a exigência é
severa. E a impedância da conexão de terra importa: um fio longo de aterramento
tem indutância e degrada a blindagem em altas frequências.}

\gabarito[1.1 Cargas elétricas]

\subsection{Lei de Coulomb}
\stepcounter{questao}
\resposta{$F = \SI{33.75}{\newton}$, atrativa}
\resolucao{Com $d = \SI{4}{\centi\meter} = \SI{0.04}{\meter}$:
\[
F=\ke\frac{|q_1||q_2|}{d^2}
=9\cdot10^9\cdot\frac{(\pot{2}{-6})(\pot{3}{-6})}{(0{,}04)^2}
=\SI{33.75}{\newton}.
\]
Como as cargas têm sinais opostos, a força é \textbf{atrativa}.}
\stepcounter{questao}
\resposta{$F = \SI{562.5}{\newton}$, repulsiva}
\resolucao{$F = 9\cdot10^9\cdot\dfrac{(\pot{5}{-6})^2}{(0{,}02)^2}
=\SI{562.5}{\newton}$, \textbf{repulsiva} (sinais iguais).}
\stepcounter{questao}
\resposta{(d)}
\resolucao{Como $F \propto \dfrac{1}{d^2}$, reduzir $d$ à metade \emph{multiplica}
a força por 4:
\[
F' = 4F = 4\cdot\pot{42}{-5} = \pot{168}{-5}\,\si{\newton}.
\]}
\stepcounter{questao}
\resposta{(e)}
\resolucao{Situação inicial: $F = \ke\dfrac{q\cdot(q/2)}{d^2} = \dfrac{\ke q^2}{2d^2}$.\\
Situação nova: $F' = \ke\dfrac{q\cdot 2q}{(d/2)^2} = \dfrac{2\ke q^2}{d^2/4}
= \dfrac{8\ke q^2}{d^2}$.\\
Dividindo: $\dfrac{F'}{F} = \dfrac{8\ke q^2/d^2}{\ke q^2/(2d^2)} = 16$, logo
$F' = 16F$. O produto das cargas quadruplicou e a distância caiu à metade
(fator 4 na força): $4\times4 = 16$.}
\stepcounter{questao}
\resposta{(d)}
\resolucao{Reduzir cada carga à metade divide o produto $q_1 q_2$ por 4, o que
sozinho reduziria $F$ a $1/4$. Para $F$ permanecer constante, o fator $1/d^2$
precisa \emph{quadruplicar}, ou seja, $d^2$ deve cair a $1/4$, o que significa $d$
reduzido à \textbf{metade}:
\[
F=\ke\frac{(q_1/2)(q_2/2)}{d'^2}=\ke\frac{q_1q_2}{4d'^2}
\stackrel{!}{=}\ke\frac{q_1q_2}{d^2}
\ \Rightarrow\ d'^2=\frac{d^2}{4}\ \Rightarrow\ d'=\frac{d}{2}.
\]}
\stepcounter{questao}
\resposta{$F_R \approx \SI{6.24}{\newton}$}
\resolucao{Cada par de cargas exerce uma força de módulo
\[
F=\ke\frac{q^2}{d^2}=9\cdot10^9\cdot\frac{(\pot{1}{-6})^2}{(0{,}05)^2}
=\SI{3.6}{\newton}.
\]
Sobre a carga escolhida atuam \emph{duas} forças iguais de \SI{3.6}{\newton}, com
$\SI{60}{\degree}$ entre elas (geometria do equilátero). A resultante é
\[
F_R = \sqrt{F^2+F^2+2F^2\cos 60^\circ}=F\sqrt{3}
=3{,}6\sqrt{3}\approx\SI{6.24}{\newton},
\]
dirigida para fora do triângulo, ao longo da bissetriz do vértice.}
\stepcounter{questao}
\resposta{$F_R \approx \SI{1.72}{\newton}$}
\resolucao{Sobre a carga de um vértice atuam três forças repulsivas:
duas dos vértices \emph{adjacentes} (ao longo dos lados, perpendiculares entre si)
e uma do vértice \emph{oposto} (ao longo da diagonal).\\
\textbf{Força de um lado} ($d = \SI{10}{\centi\meter}$):
$F_L = \ke\dfrac{q^2}{d^2} = 9\cdot10^9\dfrac{(\pot{1}{-6})^2}{(0{,}1)^2}
= \SI{0.9}{\newton}$.\\
\textbf{Força da diagonal} ($d' = d\sqrt2$):
$F_D = \ke\dfrac{q^2}{(d\sqrt2)^2} = \dfrac{F_L}{2} = \SI{0.45}{\newton}$.\\
As duas forças de lado, perpendiculares, têm resultante
$F_L\sqrt2 \approx \SI{1.27}{\newton}$, na direção da diagonal --- \emph{mesma}
direção de $F_D$. Somando:
\[
F_R = F_L\sqrt2 + F_D \approx 1{,}27 + 0{,}45 = \SI{1.72}{\newton}.
\]}
\stepcounter{questao}
\resposta{$q = \SI{0.2}{\micro\coulomb}$}
\resolucao{No equilíbrio, a força elétrica (para cima) equilibra o peso (para
baixo): $qE = mg$, logo
\[
q=\frac{mg}{E}=\frac{\pot{2}{-3}\cdot10}{\pot{1}{5}}
=\pot{2}{-7}\,\si{\coulomb}=\SI{0.2}{\micro\coulomb}.
\]
Para a força ser para cima (mesmo sentido de $E$), a carga deve ser
\emph{positiva}. É o princípio do experimento de Millikan, que mediu a carga do
elétron equilibrando gotas de óleo.}
\stepcounter{questao}
\resposta{$x = \SI{10}{\centi\meter}$ a partir de A}
\resolucao{No equilíbrio, as forças que $Q_1$ e $Q_2$ exercem sobre $Q_3$ têm
mesmo módulo e sentidos opostos (o que só é possível \emph{entre} as cargas, já
que ambas são positivas). Igualando os módulos, com $d = \SI{30}{\centi\meter}$:
\[
\ke\frac{Q_1 Q_3}{x^2}=\ke\frac{4Q_1\,Q_3}{(d-x)^2}
\ \Rightarrow\
\frac{1}{x^2}=\frac{4}{(d-x)^2}
\ \Rightarrow\
\frac{d-x}{x}=2.
\]
Logo $d - x = 2x \Rightarrow x = \dfrac{d}{3} = \SI{10}{\centi\meter}$. O ponto de
equilíbrio fica \emph{mais próximo} da carga menor ($Q_1$), como esperado.
Note que o valor de $Q_3$ e seu sinal não influenciam a posição.}
\stepcounter{questao}
\resposta{(a)}
\resolucao{A força devida ao par $+Q$/$-Q$ já está sobre a diagonal
$-Q\,$--$\,+Q$. Para que a resultante \emph{total} continue sobre essa mesma
diagonal, a contribuição do par $q_1$/$q_2$ (na outra diagonal) não pode ter
componente perpendicular a ela: as forças de $q_1$ e $q_2$ sobre a carga central
devem se \emph{cancelar} nessa direção transversal. Isso ocorre quando $q_1$ e
$q_2$ têm o mesmo módulo e o mesmo sinal (por simetria em relação à diagonal
principal), produzindo forças opostas sobre a carga de prova que se anulam.}
\stepcounter{questao}
\resposta{(a) $x = \SI{6}{\centi\meter}$; (b) $x=\SI{30}{\centi\meter}$ (fora do
segmento); (c) estável para $q_3<0$, instável para $q_3>0$}
\resolucao{(a) Igualando os módulos das forças (a carga $q_3$ cancela-se):
\[
\frac{\ke\,9}{x^{2}}=\frac{\ke\,4}{(10-x)^{2}}
\;\Longrightarrow\;
\frac{3}{x}=\frac{2}{10-x}
\;\Longrightarrow\;
30-3x=2x
\;\Longrightarrow\;
x=\SI{6}{\centi\meter}.
\]
(b) Ao elevar ao quadrado, a equação admite também $\dfrac{3}{x}=-\dfrac{2}{10-x}$,
que fornece $x = \SI{30}{\centi\meter}$. Nesse ponto, \emph{fora} do segmento,
as duas forças apontam no \textbf{mesmo sentido} (ambas as cargas são positivas e
empurram $q_3$ para longe) --- elas se somam e nunca se cancelam. A raiz é
espúria, introduzida pela elevação ao quadrado.

(c) \textbf{Análise de estabilidade.} Suponha $q_3$ deslocada ligeiramente para a
direita (aproximando-se de $q_2$):
\begin{itemize}[leftmargin=1.6em,itemsep=1pt,topsep=2pt]
\item Se $q_3 > 0$: a repulsão de $q_2$ (agora mais próxima) supera a de $q_1$, e
      $q_3$ é empurrada \emph{ainda mais} para longe do ponto de equilíbrio ---
      \textbf{instável}.
\item Se $q_3 < 0$: a atração por $q_2$ cresce, puxando $q_3$ \emph{de volta}?
      Não --- a atração por $q_2$ a puxa para longe de $x=6$. Nesse eixo,
      portanto, também é instável.
\end{itemize}
\textbf{Análise completa:} para $q_3<0$, o equilíbrio é estável na direção
\emph{transversal} ao eixo (qualquer deslocamento lateral gera força restauradora),
mas instável ao longo do eixo. Para $q_3>0$, ocorre o oposto.\\
\textbf{Conclusão profunda --- teorema de Earnshaw:} \emph{nenhuma} configuração
de cargas estáticas produz equilíbrio estável em \textbf{todas} as direções. Sempre
há ao menos uma direção instável. É por isso que não se consegue levitar um objeto
apenas com cargas fixas --- são necessários campos variáveis (armadilhas de Paul)
ou diamagnetismo.}
\stepcounter{questao}
\resposta{(a) demonstração; (b) $\approx\pot{4.2}{42}$; (c) ver comentário}
\resolucao{(a) Ambas as forças seguem a lei do inverso do quadrado:
\[
F_e=\ke\frac{e^{2}}{d^{2}},
\qquad
F_g=G\frac{m_e^{2}}{d^{2}}
\;\Longrightarrow\;
\frac{F_e}{F_g}=\frac{\ke\,e^{2}}{G\,m_e^{2}}.
\]
O fator $d^{2}$ \textbf{cancela-se}: a razão é uma constante universal.\\
(b) Numericamente:
\[
\frac{F_e}{F_g}=\frac{(9\cdot10^{9})(\pot{1.6}{-19})^{2}}
{(\pot{6.67}{-11})(\pot{9.11}{-31})^{2}}\approx\pot{4.2}{42}.
\]
(c) \textbf{Duas conclusões profundas:}
\begin{itemize}[leftmargin=1.6em,itemsep=2pt,topsep=3pt]
\item \textbf{Na escala atômica, a gravidade é irrelevante.} A força elétrica é
      $10^{42}$ vezes maior --- por isso a estrutura de átomos, moléculas,
      cristais e de toda a química é governada \emph{exclusivamente} por forças
      eletromagnéticas. Nenhum modelo atômico precisa incluir gravidade.
\item \textbf{Na escala astronômica, ocorre o oposto.} Corpos celestes são
      eletricamente \emph{neutros} (as cargas positivas e negativas se cancelam
      quase perfeitamente), então a força elétrica resultante é praticamente
      nula. Já a massa só se \emph{acumula} --- não existe "massa negativa". Por
      isso a gravidade, apesar de absurdamente mais fraca, domina o cosmos.
\end{itemize}
\emph{A natureza usa a força fraca para construir galáxias e a força forte para
construir átomos --- precisamente porque uma se cancela e a outra não.}}
\stepcounter{questao}
\resposta{(a) $d\approx\SI{6.9}{\centi\meter}$; (b) $F\approx\SI{0.88}{\milli\newton}$;
(c) $q\approx\SI{22}{\nano\coulomb}$}
\resolucao{(a) Cada esfera se afasta $L\sin\theta$ da vertical, logo
\[
d=2L\sin\theta=2(0{,}20)(0{,}174)\approx\SI{0.069}{\meter}=\SI{6.9}{\centi\meter}.
\]
(b) No equilíbrio de cada esfera (tração, peso e força elétrica),
a decomposição fornece $\tan\theta = \dfrac{F_e}{mg}$:
\[
F_e = mg\tan\theta=(\pot{0.5}{-3})(10)(0{,}176)
\approx\pot{8.8}{-4}\,\si{\newton}=\SI{0.88}{\milli\newton}.
\]
(c) Da lei de Coulomb, com as duas cargas iguais:
\[
F_e=\ke\frac{q^{2}}{d^{2}}
\;\Longrightarrow\;
q=d\sqrt{\frac{F_e}{\ke}}
=0{,}069\sqrt{\frac{\pot{8.8}{-4}}{9\cdot10^{9}}}
\approx\pot{2.2}{-8}\,\si{\coulomb}=\SI{22}{\nano\coulomb}.
\]
\textbf{Ordem de grandeza:} apenas $\sim\pot{1.4}{11}$ elétrons em excesso ---
uma quantidade ínfima --- já produz um afastamento visível. \emph{Isso mostra
quão intensa é a força elétrica em escala macroscópica.}}
\stepcounter{questao}
\resposta{$F_R \approx \SI{0.82}{\newton}$, dirigida ao centro do quadrado}
\resolucao{\textbf{Forças envolvidas.} Sobre a carga positiva escolhida atuam:
\begin{itemize}[leftmargin=1.6em,itemsep=1pt,topsep=2pt]
\item duas forças dos vizinhos \emph{adjacentes} (que são \textbf{negativos}):
      são de \textbf{atração}, ao longo dos lados;
\item uma força da carga da \emph{diagonal} (que é \textbf{positiva}):
      \textbf{repulsão}, ao longo da diagonal.
\end{itemize}
\textbf{Módulos:}
\[
F_L=\ke\frac{q^{2}}{a^{2}}=9\cdot10^{9}\frac{(\pot{1}{-6})^{2}}{(0{,}1)^{2}}
=\SI{0.9}{\newton},
\qquad
F_D=\ke\frac{q^{2}}{(a\sqrt2)^{2}}=\frac{F_L}{2}=\SI{0.45}{\newton}.
\]
\textbf{Composição.} As duas atrações, perpendiculares entre si, têm resultante
$F_L\sqrt2 \approx \SI{1.27}{\newton}$ apontando \emph{para o centro} do quadrado
(ao longo da diagonal). A repulsão da diagonal aponta \emph{para fora}, na mesma
reta. Logo elas se \textbf{subtraem}:
\[
F_R=F_L\sqrt2-F_D = 1{,}27-0{,}45\approx\SI{0.82}{\newton},
\]
dirigida \textbf{para o centro} do quadrado.\\
\textbf{Contraste com o caso de cargas iguais} (visto anteriormente): lá as três
forças eram repulsivas e se \emph{somavam} ($F_R=F_L\sqrt2+F_D=\SI{1.72}{\newton}$),
apontando para fora. Aqui, alternar os sinais inverte o sentido e reduz o módulo
para menos da metade. \emph{A geometria é a mesma; a configuração de sinais muda
tudo.}}
\stepcounter{questao}
\resposta{(a) a $d/3$ de $q_1$; (b) $q_3=-\dfrac{4q}{9}$; (c) instável}
\resolucao{(a) Para $q_3$ estar em equilíbrio, ela deve ficar onde os campos de
$q_1$ e $q_2$ se cancelam (calculado anteriormente):
\[
\frac{q}{x^{2}}=\frac{4q}{(d-x)^{2}}
\;\Rightarrow\;
\frac{d-x}{x}=2
\;\Rightarrow\;
x=\frac{d}{3}.
\]
(b) Mas isso garante apenas o equilíbrio de $q_3$. Para que \emph{$q_1$} também
esteja equilibrada, a repulsão de $q_2$ sobre ela deve ser cancelada pela atração
de $q_3$:
\[
\underbrace{\ke\frac{q\cdot 4q}{d^{2}}}_{\text{repulsão de }q_2}
=\underbrace{\ke\frac{q\,|q_3|}{(d/3)^{2}}}_{\text{atração de }q_3}
\;\Longrightarrow\;
\frac{4q}{d^{2}}=\frac{9|q_3|}{d^{2}}
\;\Longrightarrow\;
|q_3|=\frac{4q}{9}.
\]
O sinal deve ser \textbf{negativo} (para atrair), logo
$q_3 = -\dfrac{4q}{9}$.\\
\textbf{Verificação em $q_2$:} a repulsão de $q_1$ sobre $q_2$ é
$\ke\dfrac{4q^2}{d^2}$; a atração de $q_3$ (à distância $2d/3$) é
$\ke\dfrac{4q\cdot(4q/9)}{(2d/3)^{2}}=\ke\dfrac{(16q^2/9)}{(4d^2/9)}
=\ke\dfrac{4q^{2}}{d^{2}}$. \checkmark\ São iguais --- as três estão em
equilíbrio.\\
(c) \textbf{Instável}, pelo teorema de Earnshaw: qualquer deslocamento infinitesimal
de qualquer das cargas destrói o equilíbrio, e as forças resultantes \emph{afastam}
ainda mais o sistema da configuração original. Não existe arranjo eletrostático
puro que seja estável.}
\stepcounter{questao}
\resposta{(a) $q \approx \SI{5.8}{\micro\coulomb}$; (b) $T\approx\SI{0.115}{\newton}$;
(c) não --- o ângulo cresce menos que proporcionalmente}
\resolucao{\textbf{(a)} No equilíbrio, decompondo a tração:
\[
T\sin\theta = F_e = qE,
\qquad
T\cos\theta = P = mg
\quad\Longrightarrow\quad
\tan\theta=\frac{qE}{mg}.
\]
Isolando $q$:
\[
q=\frac{mg\tan\theta}{E}
=\frac{(\pot{10}{-3})(10)(0{,}577)}{\pot{1}{4}}
\approx\pot{5.8}{-6}\,\si{\coulomb}=\SI{5.8}{\micro\coulomb}.
\]
\textbf{(b)} $T = \dfrac{mg}{\cos\theta} = \dfrac{0{,}1}{0{,}866}
\approx\SI{0.115}{\newton}$.\\
\textbf{(c)} \textbf{Não.} A relação é $\tan\theta \propto E$, e \emph{não}
$\theta \propto E$. Dobrando o campo:
\[
\tan\theta' = 2\tan30^\circ = 1{,}155
\;\Rightarrow\;
\theta' = \arctan(1{,}155)\approx 49{,}1^\circ,
\]
que é \emph{menos} que $2\times30^\circ = 60^\circ$. A função tangente cresce mais
rápido que o ângulo, então o ângulo responde de forma \emph{sublinear}. No limite
$E\to\infty$, $\theta\to90^\circ$ mas nunca o atinge --- o fio jamais fica
horizontal.}
\stepcounter{questao}
\resposta{(a) nula, por simetria; (b) $F=\dfrac{7\ke Q^{2}}{4a^{2}}$ dirigida para
$+x$; (c) ver comentário}
\resolucao{(a) A carga central $-2Q$ está \emph{equidistante} das duas cargas
$+Q$, que são iguais. As duas atrações têm mesmo módulo e sentidos opostos: a
resultante é \textbf{nula} por simetria.\\
(b) Sobre a carga em $x=0$ atuam:
\begin{itemize}[leftmargin=1.6em,itemsep=1pt,topsep=2pt]
\item atração de $-2Q$ (distância $a$):
      $F_1=\ke\dfrac{Q\cdot2Q}{a^{2}}=\dfrac{2\ke Q^{2}}{a^{2}}$, para $+x$;
\item repulsão de $+Q$ (distância $2a$):
      $F_2=\ke\dfrac{Q^{2}}{4a^{2}}$, para $-x$.
\end{itemize}
Resultante:
\[
F=F_1-F_2=\frac{2\ke Q^{2}}{a^{2}}-\frac{\ke Q^{2}}{4a^{2}}
=\frac{\ke Q^{2}}{a^{2}}\left(2-\frac14\right)
=\frac{7\ke Q^{2}}{4a^{2}},
\]
dirigida para $+x$ (para dentro do conjunto).\\
(c) \textbf{Por que a carga total nula não implica campo nulo.} A carga total é
$Q-2Q+Q=0$, então \emph{a longas distâncias} o campo decai muito mais rápido que
o de uma carga isolada. Mas ele não é nulo: as cargas ocupam \emph{posições
diferentes}, e essa separação espacial cria momentos de ordem superior.\\
\textbf{Hierarquia dos decaimentos:}
\begin{center}\small\sffamily
\begin{tabular}{@{}l c c@{}}
\toprule
\textbf{Configuração} & \textbf{Carga total} & \textbf{Campo distante}\\
\midrule
Monopolo (carga isolada) & $\neq0$ & $\propto 1/r^{2}$\\
Dipolo ($+q,-q$)         & $0$     & $\propto 1/r^{3}$\\
Quadrupolo (este caso)   & $0$     & $\propto 1/r^{4}$\\
\bottomrule
\end{tabular}
\end{center}
\emph{Cada cancelamento sucessivo acelera o decaimento --- é o princípio da
expansão multipolar, base da análise de antenas e de moléculas polares.}}
\stepcounter{questao}
\resposta{(a) $F=\dfrac{\ke\,q(Q-q)}{d^{2}}$; (b) $q=Q/2$;
(c) $F_{\max}=\dfrac{\ke Q^{2}}{4d^{2}}$}
\resolucao{(a) Pela lei de Coulomb,
\[
F(q)=\ke\,\frac{q\,(Q-q)}{d^{2}}=\frac{\ke}{d^{2}}\left(Qq-q^{2}\right).
\]
(b) É uma \textbf{parábola} com concavidade para baixo em função de $q$. O máximo
ocorre no vértice. Derivando:
\[
\frac{\mathrm{d}F}{\mathrm{d}q}=\frac{\ke}{d^{2}}\left(Q-2q\right)=0
\;\Longrightarrow\;
\boxed{\,q=\frac{Q}{2}\,}
\]
(Sem cálculo: o vértice de $Qq-q^2$ está em $q = -\dfrac{b}{2a}=\dfrac{Q}{2}$.)\\
(c) Substituindo:
\[
F_{\max}=\frac{\ke}{d^{2}}\cdot\frac{Q}{2}\cdot\frac{Q}{2}
=\frac{\ke\,Q^{2}}{4d^{2}}.
\]
\textbf{Interpretação.} A força é máxima quando a carga é dividida
\emph{igualmente}. Isso decorre da desigualdade entre médias: para soma fixa
($q+(Q-q)=Q$), o produto é máximo quando as parcelas são iguais.\\
\textbf{Casos-limite:} se $q\to0$ ou $q\to Q$, uma das partes fica sem carga e
$F\to0$ --- coerente. \emph{Este é um problema de otimização disfarçado de
eletrostática: reconhecer a estrutura matemática vale mais que decorar a
fórmula.}}
\stepcounter{questao}
\resposta{$F=\SI{6}{\milli\newton}$, repulsiva}
\resolucao{Convertendo antes de substituir --- este é o erro mais comum:
$q_1=\pot{2}{-8}\,\si{\coulomb}$, $q_2=\pot{3}{-8}\,\si{\coulomb}$,
$r=\pot{3}{-2}\,\si{\meter}$.
\[
F=\ke\frac{q_1q_2}{r^{2}}
 =\frac{\pot{9}{9}\cdot\pot{2}{-8}\cdot\pot{3}{-8}}{(\pot{3}{-2})^{2}}
 =\frac{\pot{5.4}{-6}}{\pot{9}{-4}}=\pot{6}{-3}\,\si{\newton}.
\]
Ambas positivas $\Rightarrow$ repulsão.}
\stepcounter{questao}
\resposta{(a) $q=\SI{2}{\micro\coulomb}$; (b) $n=\pot{1.25}{13}$ elétrons}
\resolucao{(a) De $F=\ke q^{2}/r^{2}$,
\[
q=\sqrt{\frac{Fr^{2}}{\ke}}=\sqrt{\frac{0{,}90\cdot(0{,}20)^{2}}{\pot{9}{9}}}
 =\sqrt{\pot{4}{-12}}=\pot{2}{-6}\,\si{\coulomb}.
\]
(b) Repulsão com cargas iguais e sinal positivo indica \emph{falta} de
elétrons: $n=q/e=\pot{2}{-6}/\pot{1.6}{-19}=\pot{1.25}{13}$.}
\stepcounter{questao}
\resposta{$r=\SI{0.30}{\meter}$}
\resolucao{Isolando $r$ na lei de Coulomb:
\[
r=\sqrt{\frac{\ke q_1q_2}{F}}
 =\sqrt{\frac{\pot{9}{9}\cdot(\pot{5}{-6})^{2}}{2{,}5}}
 =\sqrt{\frac{0{,}225}{2{,}5}}=\sqrt{0{,}09}=\SI{0.30}{\meter}.
\]}
\stepcounter{questao}
\resposta{(a)}
\resolucao{$F\propto 1/r^{2}$. Dobrar $r$ divide a força por $2^{2}=4$:
$F'=F/4$. Note que a natureza (atrativa) não muda --- ela depende apenas dos
sinais, nunca da distância.}
\stepcounter{questao}
\resposta{$F'=\SI{0.20}{\newton}$}
\resolucao{Em um meio dielétrico, $F_{\text{meio}}=F_{\text{vácuo}}/\varepsilon_r$,
pois $\ke$ é substituída por $\ke/\varepsilon_r$:
\[
F'=\frac{0{,}80}{4}=\SI{0.20}{\newton}.
\]
O dielétrico se polariza e blinda parcialmente as cargas --- por isso a força
\emph{sempre} diminui ($\varepsilon_r\ge 1$).}
\stepcounter{questao}
\resposta{(b)}
\resolucao{A carga \emph{líquida} da esfera continua nula, mas o bastão
positivo atrai elétrons livres para a face próxima e repele carga positiva
para a face distante (\textbf{indução eletrostática}). A face próxima, negativa,
está \emph{mais perto} do bastão que a face distante, positiva. Como
$F\propto 1/r^{2}$, a atração vence: resultante \textbf{atrativa}. É o mesmo
princípio pelo qual um canudo atritado atrai pedaços de papel neutros.}
\stepcounter{questao}
\resposta{$F=\SI{9}{\kilo\newton}$ (o peso de cerca de \SI{900}{\kilogram})}
\resolucao{
\[
F=\frac{\pot{9}{9}\cdot 1\cdot 1}{(10^{3})^{2}}=\pot{9}{3}\,\si{\newton}.
\]
Nove mil newtons --- o peso de quase uma tonelada --- a \emph{um quilômetro}
de distância. Isso mostra por que o coulomb é uma unidade enorme para a
eletrostática: cargas realmente separáveis por atrito ficam na faixa de
$\si{\nano\coulomb}$ a $\si{\micro\coulomb}$.}
\stepcounter{questao}
\resposta{$F=\SI{2.35}{\newton}$, no sentido $+x$}
\resolucao{Trate cada par separadamente e some \emph{com sinal}.
\[
F_{12}=\frac{\pot{9}{9}\cdot\pot{3}{-6}\cdot\pot{2}{-6}}{(0{,}20)^{2}}
=\SI{1.35}{\newton}\quad(\text{repulsão}\Rightarrow +x)
\]
\[
F_{32}=\frac{\pot{9}{9}\cdot\pot{2}{-6}\cdot\pot{5}{-6}}{(0{,}30)^{2}}
=\SI{1.00}{\newton}\quad(\text{atração por } q_3 \Rightarrow +x)
\]
As duas apontam no mesmo sentido: $F=1{,}35+1{,}00=\SI{2.35}{\newton}$ em $+x$.
\textbf{Cuidado:} não some as cargas antes; a lei de Coulomb vale par a par
(princípio da superposição).}
\stepcounter{questao}
\resposta{(a) $\SI{1.6}{\newton}$, atrativa; (b) $\SI{0.9}{\newton}$, repulsiva}
\resolucao{(a) $F=\dfrac{\pot{9}{9}\cdot\pot{8}{-6}\cdot\pot{2}{-6}}{(0{,}30)^{2}}
=\dfrac{0{,}144}{0{,}09}=\SI{1.6}{\newton}$, atrativa (sinais opostos).\\
(b) Esferas \emph{idênticas} em contato dividem a carga total igualmente:
\[
q'=\frac{q_1+q_2}{2}=\frac{+8-2}{2}=\SI{+3}{\micro\coulomb}\ \text{cada}.
\]
\[
F'=\frac{\pot{9}{9}\cdot(\pot{3}{-6})^{2}}{0{,}09}=\SI{0.9}{\newton},
\ \text{agora repulsiva}.
\]
A força \emph{e} a natureza mudaram: o contato conserva a carga total, não a
individual.}
\stepcounter{questao}
\resposta{(a) $\SI{0.24}{\newton}$; (b) $r'=r/\sqrt{2{,}5}\approx 0{,}63\,r$}
\resolucao{(a) $F'=F/\varepsilon_r=0{,}60/2{,}5=\SI{0.24}{\newton}$.\\
(b) Queremos $\dfrac{\ke q_1q_2}{\varepsilon_r r'^{2}}=\dfrac{\ke q_1q_2}{r^{2}}$,
ou seja $\varepsilon_r r'^{2}=r^{2}$:
\[
r'=\frac{r}{\sqrt{\varepsilon_r}}=\frac{r}{\sqrt{2{,}5}}=0{,}632\,r.
\]
Aproximar as cargas em cerca de $37\%$ compensa exatamente o efeito do meio.}
\stepcounter{questao}
\resposta{(d)}
\resolucao{Os efeitos são multiplicativos e independentes:
\[
F'=\ke\frac{(3q_1)q_2}{(d/2)^{2}}=3\cdot 4\cdot \ke\frac{q_1q_2}{d^{2}}=12F.
\]
Fator $3$ da carga $\times$ fator $4$ da distância $=12$.}
\stepcounter{questao}
\resposta{$F_e=\pot{9.2}{-12}\,\si{\newton}$; $P=\pot{9.8}{-11}\,\si{\newton}$;
$F_e/P\approx 0{,}094$}
\resolucao{Carga da gota: $q=ne=\pot{2.0}{5}\cdot\pot{1.6}{-19}
=\pot{3.2}{-14}\,\si{\coulomb}$.
\[
F_e=\frac{\pot{9}{9}\cdot(\pot{3.2}{-14})^{2}}{(\pot{1.0}{-3})^{2}}
=\pot{9.2}{-12}\,\si{\newton}.
\]
Peso: $P=mg=\pot{1.0}{-11}\cdot 9{,}8=\pot{9.8}{-11}\,\si{\newton}$.
Razão $\approx 9{,}4\%$: a repulsão entre gotas é pequena diante do peso, mas
não desprezível --- é ela que degrada a nitidez quando as gotas são emitidas
muito próximas.}
\stepcounter{questao}
\resposta{$x=\SI{20}{\centi\meter}$}
\resolucao{Como as duas são positivas, o ponto de força nula está
\emph{entre} elas. Seja $x$ a distância a $q_1$:
\[
\frac{\ke q_1 q_3}{x^{2}}=\frac{\ke q_2 q_3}{(0{,}50-x)^{2}}
\;\Longrightarrow\;
\frac{4}{x^{2}}=\frac{9}{(0{,}50-x)^{2}}.
\]
Extraindo a raiz (ambos os lados positivos):
$\dfrac{2}{x}=\dfrac{3}{0{,}50-x}\Rightarrow 2(0{,}50-x)=3x
\Rightarrow x=\SI{0.20}{\meter}$.
O ponto fica \emph{mais próximo da carga menor}, como esperado. Note que $q_3$
se cancela: a posição não depende dela.}
\stepcounter{questao}
\resposta{(c)}
\resolucao{A lei de Coulomb é simétrica no produto $q_1q_2$: ambas valem
$\ke\,q\cdot 100q/r^{2}$. Formam um par ação--reação (3\textsuperscript{a} lei
de Newton): mesmo módulo, mesma direção, sentidos opostos. O que \emph{difere}
é a \textbf{aceleração}, se as massas forem diferentes --- confundir força com
aceleração é o erro clássico aqui.}
\stepcounter{questao}
\resposta{$F=\pot{8.2}{-8}\,\si{\newton}$;
$a=\pot{9.0}{22}\,\si{\meter\per\second\squared}$}
\resolucao{
\[
F=\frac{\pot{9}{9}\cdot(\pot{1.6}{-19})^{2}}{(\pot{5.3}{-11})^{2}}
=\frac{\pot{2.304}{-28}}{\pot{2.81}{-21}}=\pot{8.2}{-8}\,\si{\newton}.
\]
\[
a=\frac{F}{m_e}=\frac{\pot{8.2}{-8}}{\pot{9.11}{-31}}
=\pot{9.0}{22}\,\si{\meter\per\second\squared}.
\]
Uma força minúscula produz aceleração de $10^{22}\,\si{\meter\per\second\squared}$
porque a massa do elétron é ainda mais minúscula.}
\stepcounter{questao}
\resposta{$F'=F/3$}
\resolucao{Antes: $F=\ke\dfrac{Q\cdot 3Q}{d^{2}}=\dfrac{3\ke Q^{2}}{d^{2}}$.\\
Após o contato, cada esfera fica com $\dfrac{Q+3Q}{2}=2Q$. A $2d$:
\[
F'=\ke\frac{(2Q)(2Q)}{(2d)^{2}}=\ke\frac{4Q^{2}}{4d^{2}}=\frac{\ke Q^{2}}{d^{2}}
=\frac{F}{3}.
\]
O produto das cargas subiu (de $3Q^2$ para $4Q^2$), mas o fator $4$ da
distância dominou.}
\stepcounter{questao}
\resposta{$F=\SI{90}{\newton}$ contra $P=\SI{9.8}{\newton}$; razão $\approx 9{,}2$}
\resolucao{
\[
F=\frac{\pot{9}{9}\cdot(\pot{1}{-6})^{2}}{(\pot{1}{-2})^{2}}
=\frac{\pot{9}{-3}}{\pot{1}{-4}}=\SI{90}{\newton};
\qquad P=1{,}0\cdot 9{,}8=\SI{9.8}{\newton}.
\]
Um milionésimo de coulomb, a um centímetro, produz nove vezes o peso de um
quilograma. Essa desproporção entre a força elétrica e a gravitacional é a razão
pela qual é a matéria macroscópica é praticamente neutra: qualquer desequilíbrio
apreciável de carga seria violentamente desfeito.}
\stepcounter{questao}
\resposta{Um único ponto: $x=\SI{60}{\centi\meter}$}
\resolucao{Com \textbf{sinais opostos}, o ponto não pode estar entre as cargas
(ali as duas forças apontam para o mesmo lado). Deve estar \emph{fora} do
segmento e, mais especificamente, do lado da carga de \emph{menor} módulo ---
só assim a maior proximidade compensa o menor valor. Testando $x>0{,}30$:
\[
\frac{\ke\cdot 4\mu}{x^{2}}=\frac{\ke\cdot 1\mu}{(x-0{,}30)^{2}}
\;\Longrightarrow\;
\frac{2}{x}=\frac{1}{x-0{,}30}
\;\Longrightarrow\;
2x-0{,}60=x \Rightarrow x=\SI{0.60}{\meter}.
\]
\textbf{Verificação:}
$\ke\cdot4\mu/(0{,}60)^{2}=\ke\cdot1\mu/(0{,}30)^{2}$ --- confere. A raiz
$x=0{,}20$ da equação quadrática completa é \emph{espúria}: corresponde ao
ponto interno, onde as forças se somam em vez de se cancelar.}
\stepcounter{questao}
\resposta{(a) nula; (b) $\SI{1.8}{\newton}$, apontando \emph{para} o vértice vazio}
\resolucao{(a) Em um hexágono regular, o centro dista $a$ de todos os vértices
e as cargas ocupam posições diametralmente opostas duas a duas. Cada par
contribui com forças de mesmo módulo e sentidos opostos: resultante nula, por
simetria.\\
(b) Use o argumento da superposição \emph{ao contrário}: 
$\vec{F}_{6}= \vec{F}_{5}+\vec{F}_{\text{removida}}=\vec{0}$, logo
$\vec{F}_{5}=-\vec{F}_{\text{removida}}$.
\[
|\vec{F}_{5}|=\frac{\ke Qq}{a^{2}}
=\frac{\pot{9}{9}\cdot\pot{2}{-6}\cdot\pot{1}{-6}}{(0{,}10)^{2}}
=\SI{1.8}{\newton}.
\]
A carga removida \emph{repelia} $q$ para longe do vértice; sem ela, a resultante
aponta \textbf{para} o vértice vazio. Este truque de simetria evita somar cinco
vetores.}
\stepcounter{questao}
\resposta{$F=\SI{0.75}{\newton}$, a $\SI{216.9}{\degree}$ (isto é, $\SI{36.9}{\degree}$ abaixo
do semieixo $-x$)}
\resolucao{As duas forças sobre $q_A$ são perpendiculares --- o que torna a
composição imediata.
\[
F_{AB}=\frac{\pot{9}{9}\cdot\pot{2}{-6}\cdot\pot{3}{-6}}{(0{,}30)^{2}}
=\SI{0.60}{\newton}\ \ (\text{repulsão}\Rightarrow -x)
\]
\[
F_{AC}=\frac{\pot{9}{9}\cdot\pot{2}{-6}\cdot\pot{4}{-6}}{(0{,}40)^{2}}
=\SI{0.45}{\newton}\ \ (\text{repulsão}\Rightarrow -y)
\]
\[
F=\sqrt{0{,}60^{2}+0{,}45^{2}}=\sqrt{0{,}5625}=\SI{0.75}{\newton},
\qquad
\tan\theta=\frac{0{,}45}{0{,}60}=0{,}75 \Rightarrow \theta=\SI{36.87}{\degree}.
\]
Um clássico triângulo $3$--$4$--$5$ escondido. \textbf{Atenção:} a hipotenusa
$BC=\SI{50}{\centi\meter}$ é irrelevante --- ela governa a força \emph{entre}
$q_B$ e $q_C$, não a força sobre $q_A$.}
\stepcounter{questao}
\resposta{$q_1=\SI{+0.6}{\micro\coulomb}$ e $q_2=\SI{-0.2}{\micro\coulomb}$
(ou a permuta)}
\resolucao{\textbf{Etapa 1 --- produto.} A atração indica sinais opostos:
\[
|q_1q_2|=\frac{F_1r^{2}}{\ke}=\frac{0{,}108\cdot 0{,}01}{\pot{9}{9}}
=\pot{1.2}{-13}\ \Rightarrow\ q_1q_2=-\pot{1.2}{-13}\,\si{\coulomb\squared}.
\]
\textbf{Etapa 2 --- soma.} Após o contato cada esfera fica com
$q'=(q_1+q_2)/2$:
\[
q'=\sqrt{\frac{F_2r^{2}}{\ke}}=\sqrt{\frac{0{,}036\cdot0{,}01}{\pot{9}{9}}}
=\pot{2}{-7}\,\si{\coulomb}
\ \Rightarrow\ q_1+q_2=\pot{4}{-7}\,\si{\coulomb}.
\]
\textbf{Etapa 3 --- soma e produto.} As cargas são raízes de
$x^{2}-Sx+P=0$ com $S=\pot{4}{-7}$ e $P=-\pot{1.2}{-13}$:
\[
x=\frac{\pot{4}{-7}\pm\sqrt{\pot{1.6}{-13}+\pot{4.8}{-13}}}{2}
 =\frac{\pot{4}{-7}\pm\pot{8}{-7}}{2}
 \Rightarrow x=\pot{6}{-7}\ \text{ou}\ -\pot{2}{-7}.
\]
A repulsão final confirma que a soma é não nula. Se as cargas fossem simétricas
($q$ e $-q$), o contato as neutralizaria e a força final seria zero.}
\stepcounter{questao}
\resposta{$d\approx\SI{30.3}{\centi\meter}$}
\resolucao{Sobre o plano liso, o equilíbrio na direção da superfície envolve
apenas a componente do peso e a repulsão elétrica (a normal é perpendicular):
\[
\frac{\ke Qq}{d^{2}}=mg\sin\theta.
\]
\[
mg\sin\SI{30}{\degree}=0{,}020\cdot 9{,}8\cdot 0{,}5=\SI{0.098}{\newton};
\qquad
d=\sqrt{\frac{\pot{9}{9}\cdot(\pot{1}{-6})^{2}}{0{,}098}}
=\sqrt{0{,}0918}=\SI{0.303}{\meter}.
\]
\textbf{Estabilidade:} se o bloco descer, $d$ diminui e a repulsão cresce como
$1/d^{2}$, empurrando-o de volta; se subir, a repulsão cai e o peso o traz de
volta. O equilíbrio é \emph{estável}.}
\stepcounter{questao}
\resposta{(a) as componentes verticais se cancelam; (b) $F=\SI{0.82}{\newton}$,
no sentido $-x$ (de $+q$ para $-q$)}
\resolucao{(a) Cada carga dista
$r=\sqrt{0{,}05^{2}+0{,}12^{2}}=\SI{0.13}{\meter}$ de $Q$ (triângulo
$5$--$12$--$13$). A carga $+q$ \emph{repele} $Q$ e $-q$ a \emph{atrai}: as duas
forças têm o mesmo módulo e as componentes em $y$ apontam em sentidos opostos,
cancelando-se. Restam as componentes em $x$, que se \textbf{somam}.\\
(b) Cada força vale $F_1=\ke qQ/r^{2}$ e sua componente horizontal é
$F_1\cdot(a/r)$, com $a=\SI{0.05}{\meter}$:
\[
F=2\,\frac{\ke qQ}{r^{2}}\cdot\frac{a}{r}=\frac{2\ke qQ\,a}{r^{3}}
=\frac{2\cdot\pot{9}{9}\cdot\pot{1}{-6}\cdot\pot{2}{-6}\cdot 0{,}05}
       {(0{,}13)^{3}}=\SI{0.82}{\newton}.
\]
Note a dependência $1/r^{3}$: longe do dipolo a força cai \emph{mais rápido}
que a de uma carga isolada, porque as duas cargas quase se cancelam.}
\stepcounter{questao}
\resposta{(a) $\log F=\log C+n\log r$; (b) $n\approx-2{,}00$; (c)
$\ke\approx\pot{9.1}{9}\,\si{\newton\meter\squared\per\coulomb\squared}$
(erro $\approx 1\%$)}
\resolucao{(a) Tomando logaritmos, uma lei de potência vira uma reta:
$\log F = \log C + n\log r$. O coeficiente angular do gráfico
$\log F \times \log r$ é o expoente procurado.\\
(b) Ajustando os cinco pontos por mínimos quadrados obtém-se
$n=-1{,}997$ e $C=\pot{9.09}{-5}$. O expoente reproduz $-2$ com desvio de
$0{,}15\%$ --- dentro do que se espera de medidas de força dessa ordem. Um teste
rápido sem regressão: dobrar $r$ de $0{,}020$ para $0{,}040$ deveria dividir $F$
por $4$; de fato $0{,}223/0{,}0558=4{,}00$.\\
(c) Como $C=\ke q^{2}$ e $q^{2}=\pot{1}{-14}\,\si{\coulomb\squared}$:
\[
\ke=\frac{C}{q^{2}}=\frac{\pot{9.09}{-5}}{\pot{1}{-14}}
=\pot{9.09}{9}\,\si{\newton\meter\squared\per\coulomb\squared},
\]
contra $\pot{8.99}{9}$ tabelado --- erro relativo de $1{,}1\%$.\\
\textbf{Observação de método:} a linearização logarítmica é preferível a ajustar
$F\times 1/r^{2}$ diretamente, porque \emph{não pressupõe} o expoente --- ela o
mede. Esse é exatamente o argumento que Coulomb precisou construir em 1785.}
\stepcounter{questao}
\resposta{(b) $x=R/\sqrt{2}\approx\SI{7.07}{\centi\meter}$;
(c) $F_{\max}=\dfrac{2\ke Qq}{3\sqrt{3}\,R^{2}}=\SI{6.93}{\newton}$}
\resolucao{(a) Divida o anel em elementos $\mathrm{d} Q$. Cada um dista
$s=\sqrt{x^{2}+R^{2}}$ de $q$ e contribui com $\ke q\,\mathrm{d} Q/s^{2}$. Por
simetria, as componentes perpendiculares ao eixo se cancelam aos pares; sobra a
projeção axial, fator $\cos\alpha=x/s$:
\[
F=\int \frac{\ke q\,\mathrm{d} Q}{s^{2}}\cdot\frac{x}{s}
 =\frac{\ke q x}{(x^{2}+R^{2})^{3/2}}\int \mathrm{d} Q
 =\frac{\ke Q q\,x}{(x^{2}+R^{2})^{3/2}}.
\]
(b) Derivando e igualando a zero:
\[
\frac{\mathrm{d} F}{\mathrm{d} x}\propto (x^{2}+R^{2})^{3/2}-x\cdot 3x(x^{2}+R^{2})^{1/2}=0
\;\Rightarrow\; x^{2}+R^{2}-3x^{2}=0 \;\Rightarrow\; x=\frac{R}{\sqrt{2}}.
\]
(c) Substituindo $x=R/\sqrt2$:
\[
F_{\max}=\frac{2\,\ke Qq}{3\sqrt{3}\,R^{2}}
=\frac{2\cdot\pot{9}{9}\cdot\pot{1}{-5}\cdot\pot{2}{-6}}{3\sqrt3\cdot 0{,}01}
=\frac{0{,}36}{0{,}0520}=\SI{6.93}{\newton}.
\]
\textbf{Casos-limite:} em $x=0$ a força é nula (simetria perfeita); para
$x\gg R$, $F\to\ke Qq/x^{2}$ --- o anel se comporta como carga puntiforme.
Entre esses dois zeros deve existir um máximo, e ele está em $R/\sqrt2$.}
\stepcounter{questao}
\resposta{(b) $F\to\ke Qq/d^{2}$; (c) $F=\SI{1.2}{\newton}$}
\resolucao{(a) Densidade linear $\lambda=Q/L$. Um elemento $\mathrm{d} x$ a uma
distância $x$ da carga carrega $\mathrm{d} Q=\lambda\,\mathrm{d} x$ e todas as contribuições
são \emph{colineares} --- por isso a integral é escalar:
\[
F=\int_{d}^{d+L}\frac{\ke q\lambda\,\mathrm{d} x}{x^{2}}
 =\ke q\lambda\left[-\frac{1}{x}\right]_{d}^{d+L}
 =\ke q\lambda\left(\frac{1}{d}-\frac{1}{d+L}\right)
 =\frac{\ke q\lambda L}{d(d+L)}=\frac{\ke Qq}{d(d+L)}.
\]
(b) Se $L\ll d$, então $d+L\approx d$ e $F\to\ke Qq/d^{2}$: recuperamos a lei de
Coulomb puntiforme, como deve ser.\\
(c) $F=\dfrac{\pot{9}{9}\cdot\pot{4}{-6}\cdot\pot{1}{-6}}{0{,}10\cdot 0{,}30}
=\dfrac{0{,}036}{0{,}03}=\SI{1.2}{\newton}$.\\
\textbf{Compare:} se toda a carga estivesse concentrada no centro do bastão
($\SI{20}{\centi\meter}$), teríamos $\SI{0.9}{\newton}$. A força real é maior
porque as porções mais próximas pesam com $1/x^{2}$ --- concentrar carga no
centro geométrico \emph{subestima} a força.}
\stepcounter{questao}
\resposta{(a) por simetria; (b) estável, com $F\approx-4\ke Qq\,x/d^{3}$;
(c) instável, com $F\approx+2\ke Qq\,y/d^{3}$ --- não há equilíbrio estável em
todas as direções}
\resolucao{(a) As duas forças têm o mesmo módulo e sentidos opostos: resultante
nula, independentemente do sinal e do valor de $q$.\\
(b) Deslocando $q>0$ para $x$ pequeno:
\[
F_x=\ke Qq\left[\frac{1}{(d+x)^{2}}-\frac{1}{(d-x)^{2}}\right].
\]
Com $(d\pm x)^{-2}\approx d^{-2}(1\mp 2x/d)$:
\[
F_x\approx\frac{\ke Qq}{d^{2}}\left(-\frac{4x}{d}\right)=-\frac{4\ke Qq}{d^{3}}\,x.
\]
Sinal negativo $\Rightarrow$ força \textbf{restauradora}: estável nessa direção
(e o movimento é harmônico simples).\\
(c) Deslocando para $y$ pequeno, cada carga dista $\sqrt{d^{2}+y^{2}}\approx d$ e
repele $q$ com componente transversal $\ke Qq\,y/d^{3}$; as duas se somam:
\[
F_y\approx +\frac{2\ke Qq}{d^{3}}\,y,
\]
sinal positivo $\Rightarrow$ força que \textbf{afasta}: instável.\\
\textbf{Conclusão.} O equilíbrio é do tipo \emph{ponto de sela}: estável em uma
direção, instável na outra. Este é um caso particular do \textbf{teorema de
Earnshaw}: nenhuma configuração de cargas fixas pode manter uma carga puntiforme
em equilíbrio estável apenas por forças eletrostáticas. É por isso que modelos
puramente eletrostáticos do átomo não se sustentam.}
\stepcounter{questao}
\resposta{$F=\dfrac{\pi^{2}}{12}\cdot\dfrac{\ke Qq}{a^{2}}
\approx 0{,}822\,\dfrac{\ke Qq}{a^{2}}$, atrativa (sentido $+x$)}
\resolucao{(a) A $n$-ésima carga está em $x=na$ e tem sinal $(-1)^{n}Q$. Uma
carga negativa \emph{atrai} $q$ (puxa no sentido $+x$); uma positiva
\emph{repele} (sentido $-x$). Tomando $+x$ como positivo:
\[
F=\frac{\ke Qq}{a^{2}}\left(\frac{1}{1^{2}}-\frac{1}{2^{2}}+\frac{1}{3^{2}}
-\frac{1}{4^{2}}+\cdots\right)
=\frac{\ke Qq}{a^{2}}\sum_{n=1}^{\infty}\frac{(-1)^{n+1}}{n^{2}}.
\]
(b) Usando o resultado dado:
\[
F=\frac{\pi^{2}}{12}\cdot\frac{\ke Qq}{a^{2}}\approx 0{,}822\,\frac{\ke Qq}{a^{2}}.
\]
\textbf{Interpretação.} Apesar de haver infinitas cargas, a força é
\emph{finita} e vale menos que a da primeira carga sozinha: a alternância de
sinais produz cancelamento progressivo, e o decaimento $1/n^{2}$ garante
convergência absoluta. Se todas as cargas fossem negativas, a soma seria
$\pi^2/6\approx1{,}645$ --- exatamente o dobro.}
\stepcounter{questao}
\resposta{(b) $\omega=\sqrt{\ke Qq/(mR^{3})}=\SI{300}{\radian\per\second}$;
$T\approx\SI{20.9}{\milli\second}$}
\resolucao{(a) A força axial sobre a carga negativa é atrativa, dirigida ao
centro, de módulo $F=\ke Qq\,x/(x^{2}+R^{2})^{3/2}$. Para $x\ll R$, o
denominador tende a $R^{3}$:
\[
F_x\approx-\frac{\ke Qq}{R^{3}}\,x.
\]
Força proporcional ao deslocamento e de sinal oposto --- exatamente a forma
$F=-kx$ da lei de Hooke, com constante efetiva $k_{\text{ef}}=\ke Qq/R^{3}$.
Logo o MHS.\\
(b)
\[
\omega=\sqrt{\frac{k_{\text{ef}}}{m}}=\sqrt{\frac{\ke Qq}{mR^{3}}}
=\sqrt{\frac{\pot{9}{9}\cdot\pot{5}{-6}\cdot\pot{2}{-6}}
{\pot{1}{-3}\cdot(0{,}10)^{3}}}
=\sqrt{\frac{0{,}09}{\pot{1}{-6}}}=\SI{300}{\radian\per\second},
\]
\[
T=\frac{2\pi}{\omega}=\SI{20.9}{\milli\second}.
\]
\textbf{Observe:} o período não depende de $x_0$ --- é a assinatura do MHS. Fora
da aproximação $x\ll R$, a força deixa de ser linear e o movimento continua
periódico, mas com período dependente da amplitude.}
\stepcounter{questao}
\resposta{(c) $q\approx\SI{19.8}{\nano\coulomb}$ e
$\mathrm{d} q/\mathrm{d} t\approx\SI{-0.49}{\nano\coulomb\per\second}$}
\resolucao{(a) Para cada esfera: tensão, peso e repulsão. Do triângulo de
forças, $\tan\theta=F/(mg)$. Com $x\ll L$,
$\tan\theta\approx\sin\theta=(x/2)/L$. Igualando:
\[
\frac{x}{2L}=\frac{\ke q^{2}}{x^{2}mg}
\;\Longrightarrow\;
x^{3}=\frac{2\ke q^{2}L}{mg}.
\]
(b) Derivando implicitamente em relação ao tempo:
\[
3x^{2}\frac{\mathrm{d} x}{\mathrm{d} t}=\frac{4\ke L}{mg}\,q\,\frac{\mathrm{d} q}{\mathrm{d} t}.
\]
De (a), $\dfrac{2\ke L}{mg}=\dfrac{x^{3}}{q^{2}}$; substituindo:
\[
3x^{2}\frac{\mathrm{d} x}{\mathrm{d} t}=\frac{2x^{3}}{q}\frac{\mathrm{d} q}{\mathrm{d} t}
\;\Longrightarrow\;
\frac{\mathrm{d} x}{\mathrm{d} t}=\frac{2x}{3q}\frac{\mathrm{d} q}{\mathrm{d} t}.
\]
(c) Primeiro $q$, por (a):
\[
q=\sqrt{\frac{x^{3}mg}{2\ke L}}
=\sqrt{\frac{(0{,}060)^{3}\cdot 0{,}0020\cdot 9{,}8}{2\cdot\pot{9}{9}\cdot0{,}60}}
=\pot{1.98}{-8}\,\si{\coulomb}.
\]
Com $\mathrm{d} x/\mathrm{d} t=-\pot{1.0}{-3}\,\si{\meter\per\second}$ (aproximação):
\[
\frac{\mathrm{d} q}{\mathrm{d} t}=\frac{3q}{2x}\frac{\mathrm{d} x}{\mathrm{d} t}
=\frac{3\cdot\pot{1.98}{-8}}{2\cdot 0{,}060}\cdot(-\pot{1}{-3})
=-\pot{4.9}{-10}\,\si{\coulomb\per\second}.
\]
\textbf{Leitura física:} medir uma \emph{velocidade} converte-se em medir uma
\emph{corrente de fuga} de meio nanoampère --- pequena demais para muitos
amperímetros, mas visível a olho nu pelo movimento das esferas.}
\stepcounter{questao}
\resposta{$F=\SI{6.9}{\newton}$, ao longo do eixo de simetria}
\resolucao{(a) Para cada elemento do fio existe outro simétrico em relação ao
eixo vertical. Suas contribuições têm componentes perpendiculares ao eixo iguais
e opostas --- que se cancelam --- e componentes ao longo do eixo que se somam.\\
(b) Densidade linear $\lambda=Q/(\pi R)$. O elemento em $\theta$ tem
$\mathrm{d} Q=\lambda R\,\mathrm{d}\theta$, dista $R$ do centro e contribui com
$\ke q\,\mathrm{d} Q/R^{2}$; sua projeção no eixo vale $\sin\theta$:
\[
F=\int_{0}^{\pi}\frac{\ke q\lambda R\,\mathrm{d}\theta}{R^{2}}\sin\theta
 =\frac{\ke q\lambda}{R}\underbrace{\int_{0}^{\pi}\sin\theta\,\mathrm{d}\theta}_{=2}
 =\frac{2\ke q\lambda}{R}=\frac{2\ke Qq}{\pi R^{2}}.
\]
\[
F=\frac{2\cdot\pot{9}{9}\cdot\pot{3}{-6}\cdot\pot{1}{-6}}
{\pi\,(0{,}050)^{2}}=\frac{0{,}054}{\pot{7.85}{-3}}=\SI{6.9}{\newton}.
\]
\textbf{Compare:} se a mesma carga $Q$ estivesse concentrada em um ponto a $R$,
teríamos $\ke Qq/R^{2}=\SI{10.8}{\newton}$. O fator $2/\pi\approx 0{,}64$ mede
exatamente a perda por cancelamento vetorial ao longo do arco. Um anel
\emph{completo} levaria esse fator a zero.}
\stepcounter{questao}
\resposta{(a) $q=-\dfrac{Q(2\sqrt{2}+1)}{4}$; (b) $q\approx\SI{-1.91}{\micro\coulomb}$;
(c) não --- o equilíbrio é instável}
\resolucao{(a) Por simetria, basta impor o equilíbrio de \emph{um} vértice; a
resultante das outras três cargas aponta ao longo da diagonal, para fora.
\begin{itemize}
\item Duas vizinhas, a distância $a$: cada uma contribui $\ke Q^{2}/a^{2}$, e
suas projeções sobre a diagonal valem $\cos\SI{45}{\degree}=\sqrt2/2$ cada, somando
$\sqrt{2}\,\ke Q^{2}/a^{2}$.
\item A oposta, a distância $a\sqrt{2}$: contribui
$\ke Q^{2}/(2a^{2})$, já ao longo da diagonal.
\end{itemize}
\[
F_{\text{vértices}}=\frac{\ke Q^{2}}{a^{2}}\left(\sqrt{2}+\frac{1}{2}\right).
\]
A carga central dista $a/\sqrt{2}$ do vértice, logo atua com
$F_{c}=\ke|q|Q/(a^{2}/2)=2\ke|q|Q/a^{2}$ e deve ser \textbf{atrativa} ---
portanto $q<0$. Igualando:
\[
2|q|=Q\left(\sqrt2+\frac12\right)
\;\Longrightarrow\;
|q|=\frac{Q\,(2\sqrt{2}+1)}{4}\approx 0{,}957\,Q.
\]
(b) $q\approx-0{,}957\cdot 2=\SI{-1.91}{\micro\coulomb}$.\\
(c) \textbf{Não.} A condição imposta zera a resultante sobre os vértices, mas
não sobre a carga central --- que já está em equilíbrio por simetria, embora
instável. Mais fundamentalmente, o teorema de Earnshaw proíbe equilíbrio
eletrostático estável: qualquer perturbação de $q$ ao longo de uma diagonal a
faz cair sobre um vértice. O equilíbrio existe, mas não sobrevive a uma
perturbação infinitesimal.}

\gabarito[1.2 Lei de Coulomb]

\subsection{Campo elétrico}
\stepcounter{questao}
\resposta{$E = \pot{1.2}{9}\,\si{\newton\per\coulomb}$}
\resolucao{Com $d = \SI{10}{\milli\meter} = \SI{0.01}{\meter}$:
\[
E=\ke\frac{|q|}{d^2}
=8\cdot10^9\cdot\frac{\pot{15}{-6}}{(0{,}01)^2}=\pot{1.2}{9}\,\si{\newton\per\coulomb}.
\]
O campo aponta \emph{para} a carga, por ela ser negativa.}
\stepcounter{questao}
\resposta{(c)}
\resolucao{Como $E \propto \dfrac{1}{d^2}$, dobrar a distância divide o campo por
$2^2 = 4$: o novo valor é $E/4$.}
\stepcounter{questao}
\resposta{(a) $E=\pot{1.0}{7}\,\si{\newton\per\coulomb}$, vertical para baixo;
(b) $F = \pot{3}{4}\,\si{\newton}$}
\resolucao{(a) $E = \dfrac{F}{q} = \dfrac{\pot{1}{-2}}{\pot{1}{-9}}
= \pot{1.0}{7}\,\si{\newton\per\coulomb}$. Como a carga de prova é positiva, o
campo tem o \emph{mesmo} sentido da força: vertical, para baixo.\\
(b) $F = qE = \pot{3}{-3}\cdot\pot{1}{7} = \pot{3}{4}\,\si{\newton}$,
também vertical para baixo (carga positiva).}
\stepcounter{questao}
\resposta{$E = \pot{9.0}{5}\,\si{\newton\per\coulomb}$, radial, afastando-se de $Q$}
\resolucao{$E = \ke\dfrac{Q}{d^2}
= 9\cdot10^9\cdot\dfrac{\pot{4}{-6}}{(0{,}2)^2} = \pot{9.0}{5}\,\si{\newton\per\coulomb}$.
A direção é a da reta que une $Q$ a M, e o sentido é \emph{para fora} de $Q$
(carga positiva).}
\stepcounter{questao}
\resposta{(b)}
\resolucao{O campo se anula \emph{entre} as cargas, onde os módulos se igualam:
\[
\ke\frac{Q}{d_1^2}=\ke\frac{4Q}{d_2^2}
\ \Rightarrow\
\frac{d_2^2}{d_1^2}=4
\ \Rightarrow\
d_2 = 2d_1.
\]
O ponto de campo nulo fica \emph{mais próximo} da carga menor.}
\stepcounter{questao}
\resposta{(b)}
\resolucao{Em M, cada carga produz um campo de mesmo módulo (mesma distância,
mesmo $|Q|$).\\
\textbf{Caso I ($+Q$, $-Q$):} os dois campos apontam no \emph{mesmo} sentido (para
longe da positiva e em direção à negativa), então se \textbf{somam}: campo
resultante \emph{não-nulo}, $E = 2\cdot\dfrac{\ke Q}{d^2}$.\\
\textbf{Caso II ($+Q$, $+Q$):} os campos apontam em sentidos \emph{opostos} e se
\textbf{cancelam}: campo resultante \emph{nulo}.\\
Cuidado: para o \emph{potencial} a conclusão se inverte (nulo no caso I, não-nulo
no caso II), porque potencial é escalar e campo é vetorial.}
\stepcounter{questao}
\resposta{A $\SI{4}{\meter}$ de A (e $\SI{5}{\meter}$ de B)}
\resolucao{Como as cargas têm o mesmo sinal, o campo se anula entre elas. Seja $x$
a distância a partir de A; a distância a B é $9 - x$. Igualando os módulos:
\[
\ke\frac{q_1}{x^2}=\ke\frac{q_2}{(9-x)^2}
\ \Rightarrow\
\frac{20}{x^2}=\frac{64}{(9-x)^2}.
\]
Extraindo a raiz (grandezas positivas):
$\dfrac{9-x}{x}=\sqrt{\dfrac{64}{20}}=\dfrac{8}{\sqrt{20}}=\dfrac{8}{2\sqrt5}
=\dfrac{4}{\sqrt5}$. Resolvendo,
$9 - x = \dfrac{4}{\sqrt5}x \Rightarrow x = \dfrac{9}{1+4/\sqrt5}\approx\SI{4}{\meter}$.
O ponto neutro fica a cerca de \SI{4}{\meter} de A e \SI{5}{\meter} de B --- mais
perto da carga menor, como sempre ocorre.}
\stepcounter{questao}
\resposta{$E = \SI{20}{\newton\per\coulomb}$}
\resolucao{Como $E \propto \dfrac{1}{d^2}$, dobrar a distância (de \SI{3}{} para
\SI{6}{\meter}) divide o campo por 4:
\[
E(6)=\frac{E(3)}{4}=\frac{80}{4}=\SI{20}{\newton\per\coulomb}.
\]
Formalmente, $E(3)\cdot3^2 = E(6)\cdot6^2$, ou seja,
$80\cdot9 = E(6)\cdot36 \Rightarrow E(6) = \SI{20}{\newton\per\coulomb}$.}
\stepcounter{questao}
\resposta{A $\SI{6}{\centi\meter}$ de $q_1$}
\resolucao{Cargas de mesmo sinal: o campo se anula \emph{entre} elas. Seja $x$ a
distância a $q_1$; a distância a $q_2$ é $(10 - x)$. Igualando os módulos dos
campos:
\[
\frac{\ke q_1}{x^2}=\frac{\ke q_2}{(10-x)^2}
\ \Rightarrow\
\frac{9}{x^2}=\frac{4}{(10-x)^2}.
\]
Extraindo a raiz: $\dfrac{3}{x}=\dfrac{2}{10-x} \Rightarrow 3(10-x)=2x
\Rightarrow 30 = 5x \Rightarrow x = \SI{6}{\centi\meter}$.
O ponto neutro fica mais perto da carga \emph{menor} ($q_2$), a
\SI{4}{\centi\meter} dela.}
\stepcounter{questao}
\resposta{(a) a $2d$ da carga $+4Q$ (isto é, a $d$ além de $-Q$);
(b) e (c) ver comentário}
\resolucao{\textbf{(a)} Seja $x$ a distância medida a partir de $+4Q$, com o ponto
além de $-Q$ (logo a distância a $-Q$ é $x-d$). Igualando os módulos:
\[
\frac{\ke\,4Q}{x^{2}}=\frac{\ke\,Q}{(x-d)^{2}}
\;\Longrightarrow\;
4(x-d)^2 = x^2
\;\Longrightarrow\;
2(x-d)=\pm x.
\]
A raiz fisicamente válida é $2x-2d = x \Rightarrow \boxed{x = 2d}$, ou seja, a
$2d$ da carga maior e a $d$ além da carga menor.
(A outra raiz, $x = \tfrac{2d}{3}$, cai \emph{entre} as cargas, onde os campos
apontam no \emph{mesmo} sentido e portanto não podem se cancelar --- deve ser
descartada.)

\textbf{(b)} Entre cargas de sinais \emph{opostos}, os dois vetores campo apontam
no mesmo sentido (saindo da positiva, entrando na negativa) e se \emph{somam};
nunca se anulam ali. Fora do segmento, eles têm sentidos contrários, e o
cancelamento só é possível onde o campo da carga \emph{menor} --- que está mais
próxima --- consegue igualar o da maior. Por isso o ponto neutro fica sempre do
lado da carga de menor módulo.

\textbf{(c) Não.} Potencial é grandeza \emph{escalar}:
\[
V = \ke\frac{4Q}{2d}+\ke\frac{(-Q)}{d}
= \frac{2\ke Q}{d}-\frac{\ke Q}{d}
= \frac{\ke Q}{d}\neq 0.
\]
\textbf{Conclusão conceitual:} campo nulo e potencial nulo são condições
\emph{independentes}. Neste caso, $\vec E = 0$ mas $V \neq 0$. No ponto médio de um
dipolo simétrico ocorre o oposto: $V = 0$ mas $\vec E \neq 0$.}
\stepcounter{questao}
\resposta{(a) $a\approx\SI{1.76e15}{\meter\per\second\squared}$;
(b) $t=\SI{2.5}{\nano\second}$, $y\approx\SI{5.5}{\milli\meter}$;
(c) $\theta\approx\SI{12.4}{\degree}$}
\resolucao{Este é um \textbf{lançamento de projétil} eletrostático: movimento
uniforme na horizontal e uniformemente acelerado na vertical.\\
(a) $a=\dfrac{eE}{m_e}=\dfrac{(\pot{1.6}{-19})(\pot{1}{4})}{\pot{9.11}{-31}}
\approx\pot{1.76}{15}\,\si{\meter\per\second\squared}$
--- cerca de $10^{14}$ vezes a gravidade, o que justifica desprezar o peso.\\
(b) Na horizontal (MU):
$t=\dfrac{L}{v_0}=\dfrac{0{,}05}{\pot{2}{7}}=\pot{2.5}{-9}\,\si{\second}
=\SI{2.5}{\nano\second}$.\\
Na vertical (MUV, partindo do repouso):
\[
y=\tfrac12 a t^{2}=\tfrac12(\pot{1.76}{15})(\pot{2.5}{-9})^{2}
\approx\pot{5.5}{-3}\,\si{\meter}=\SI{5.5}{\milli\meter}.
\]
(c) A velocidade vertical na saída é $v_y = a\,t \approx
\pot{4.4}{6}\,\si{\meter\per\second}$, logo
\[
\tan\theta=\frac{v_y}{v_0}=\frac{\pot{4.4}{6}}{\pot{2}{7}}=0{,}22
\;\Longrightarrow\;
\theta\approx\SI{12.4}{\degree}.
\]
\textbf{Aplicação:} é exatamente esse o mecanismo de deflexão dos antigos tubos de
raios catódicos (TV e osciloscópios) e das impressoras a jato de tinta, que
defletem gotículas carregadas para formar caracteres.}
\stepcounter{questao}
\resposta{(a) $\vec E=0$; (b) $\vec E=0$; (c) $E\approx\pot{5.1}{6}\,\si{\newton\per\coulomb}$}
\resolucao{Cada carga dista $\dfrac{a\sqrt2}{2}=\SI{0.0707}{\meter}$ do centro,
produzindo individualmente
\[
E_1=\frac{\ke q}{(a\sqrt2/2)^{2}}=\frac{9\cdot10^{9}\cdot\pot{1}{-6}}{0{,}005}
=\pot{1.8}{6}\,\si{\newton\per\coulomb}.
\]
\textbf{(a) Todas positivas.} Os quatro vetores apontam radialmente para fora,
aos pares diametralmente opostos: cancelam-se dois a dois. $\vec E = 0$.\\
\textbf{(b) Alternadas ($+,-,+,-$).} As cargas diametralmente opostas têm o
\emph{mesmo} sinal, e seus campos continuam se cancelando aos pares.
$\vec E = 0$ novamente.\\
\textbf{(c) Adjacentes ($+,+,-,-$).} Agora cada par diametralmente oposto tem
sinais \emph{contrários}: os campos se \textbf{somam} em vez de cancelar. Os dois
pares contribuem na mesma direção (perpendicular à linha que separa as positivas
das negativas):
\[
E=2\,(E_1\sqrt2)=2\sqrt2\,E_1\approx 2{,}83\cdot\pot{1.8}{6}
\approx\pot{5.1}{6}\,\si{\newton\per\coulomb}.
\]
\textbf{Lição central.} Em (a) e (b), configurações \emph{muito} diferentes
produzem o mesmo resultado nulo. Em (c), com as mesmas quatro cargas, apenas
reposicionadas, o campo é máximo. \emph{O campo resultante depende da geometria
tanto quanto dos valores das cargas} --- e a simetria é a ferramenta que revela
isso sem nenhum cálculo pesado.}
\stepcounter{questao}
\resposta{(a) $x=-d$ (a uma distância $d$ à esquerda de $-Q$);
(b) e (c) ver comentário; $V=\dfrac{\ke Q}{d}$}
\resolucao{(a) Como as cargas têm sinais \emph{opostos}, entre elas os campos
apontam no mesmo sentido e se somam --- o cancelamento só pode ocorrer fora do
segmento. Seja $x<0$ a posição procurada; a distância a $-Q$ é $|x|$ e a $+4Q$ é
$|x|+d$:
\[
\frac{\ke Q}{|x|^{2}}=\frac{\ke\,4Q}{(|x|+d)^{2}}
\;\Longrightarrow\;
\frac{|x|+d}{|x|}=2
\;\Longrightarrow\;
|x|=d.
\]
O ponto está em $x=-d$.\\
(b) Do lado da carga \emph{maior}, seu campo sempre domina (está mais perto
\emph{e} é mais intensa) --- nunca poderia ser equilibrado. Do lado da carga
menor, a proximidade compensa a menor intensidade, e existe exatamente um ponto
de equilíbrio.\\
(c) O potencial é escalar:
\[
V=\ke\frac{(-Q)}{d}+\ke\frac{4Q}{2d}
=-\frac{\ke Q}{d}+\frac{2\ke Q}{d}
=\frac{\ke Q}{d}\neq0.
\]
\textbf{Mais um exemplo de que $\vec E=0$ não implica $V=0$} --- e vice-versa. O
campo mede a \emph{taxa de variação} do potencial, não o seu valor.}
\stepcounter{questao}
\resposta{(b) $E(0)=0$ e $E\to\ke Q/x^{2}$; (c) demonstração}
\resolucao{(a) \textbf{Argumento de simetria.} Para cada elemento de carga do
anel existe outro diametralmente oposto. As componentes de campo
\emph{perpendiculares} ao eixo desses dois elementos são iguais e opostas,
cancelando-se. Restam apenas as componentes \emph{paralelas} ao eixo, que se
somam. Logo $\vec E$ é axial.

(b) \textbf{No centro} ($x=0$): substituindo,
\[
E(0)=\frac{\ke Q\cdot 0}{R^{3}}=0.
\]
Faz sentido: no centro, as contribuições de elementos opostos se cancelam
\emph{integralmente}.\\
\textbf{Longe} ($x\gg R$): então $\left(x^{2}+R^{2}\right)^{3/2}\approx x^{3}$ e
\[
E\approx\frac{\ke Q x}{x^{3}}=\frac{\ke Q}{x^{2}},
\]
o campo de uma carga puntiforme --- como esperado, pois de longe o anel "vira"
um ponto.

(c) \textbf{Máximo.} Derivando $E(x)$ e igualando a zero:
\[
\frac{\mathrm{d}E}{\mathrm{d}x}
=\ke Q\,\frac{\left(x^{2}+R^{2}\right)^{3/2}-x\cdot\tfrac32\left(x^{2}+R^{2}\right)^{1/2}(2x)}
{\left(x^{2}+R^{2}\right)^{3}}=0.
\]
O numerador, dividido por $\left(x^{2}+R^{2}\right)^{1/2}$, dá
\[
\left(x^{2}+R^{2}\right)-3x^{2}=0
\;\Longrightarrow\;
R^{2}=2x^{2}
\;\Longrightarrow\;
\boxed{\,x=\frac{R}{\sqrt2}\,}
\]
\textbf{Interpretação física do máximo.} Perto do centro, as contribuições
axiais são pequenas (os vetores são quase perpendiculares ao eixo). Muito longe,
a distância cresce e o campo decai. Entre os dois regimes existe um ponto ótimo,
a $x=R/\sqrt2\approx0{,}71R$.\\
\textbf{Onde isso é usado:} essa geometria é a base das \emph{bobinas de
Helmholtz} (versão magnética do mesmo cálculo) e dos anéis focalizadores em
aceleradores de partículas e microscópios eletrônicos.}
\stepcounter{questao}
\resposta{(b) demonstração; (c) $E\to \ke Q/a^{2}$ (carga puntiforme)}
\resolucao{(a) Tome um elemento de comprimento $\mathrm{d}x$ à distância $x$ de P
(com $x$ variando de $a$ até $a+L$). Sua carga é
$\mathrm{d}q=\lambda\,\mathrm{d}x$ e contribui com
\[
\mathrm{d}E=\frac{\ke\,\mathrm{d}q}{x^{2}}=\frac{\ke\lambda\,\mathrm{d}x}{x^{2}}.
\]
Como todos os elementos estão alinhados com P, as contribuições têm a
\emph{mesma direção} e somam-se escalarmente:
\[
E=\int_{a}^{a+L}\frac{\ke\lambda}{x^{2}}\,\mathrm{d}x.
\]
(b) Integrando:
\[
E=\ke\lambda\left[-\frac{1}{x}\right]_{a}^{a+L}
=\ke\lambda\left(\frac{1}{a}-\frac{1}{a+L}\right)
=\ke\lambda\,\frac{L}{a(a+L)}.
\]
Como $\lambda L = Q$:
\[
\boxed{\;E=\frac{\ke\,Q}{a\,(a+L)}\;}
\]
(c) \textbf{Limite $a\gg L$:} então $a+L\approx a$ e
\[
E\approx\frac{\ke Q}{a^{2}},
\]
que é exatamente o campo de uma \textbf{carga puntiforme} $Q$.\\
\textbf{Interpretação:} de longe, qualquer distribuição finita de carga
"parece" puntiforme. Esse teste de consistência --- verificar se a expressão
obtida reproduz o caso conhecido em algum limite --- é uma das melhores formas de
validar um resultado. \emph{Se o limite não fechasse, haveria erro na
integração.}\\
\textbf{Limite oposto ($a\ll L$):} $E\approx\dfrac{\ke Q}{aL}=\dfrac{\ke\lambda}{a}$
--- campo de um fio "infinito" visto de perto, decaindo com $1/a$ em vez de
$1/a^{2}$.}
\stepcounter{questao}
\resposta{(a) demonstração; (b) e (c) ver comentário}
\resolucao{(a) Sobre a mediatriz, cada carga dista
$d=\sqrt{r^{2}+\ell^{2}}$ do ponto e produz campo de módulo
$E_1=\dfrac{\ke q}{r^{2}+\ell^{2}}$. As componentes \emph{paralelas} à mediatriz
se cancelam (simetria); restam as componentes ao longo do eixo do dipolo, cada uma
valendo $E_1\cos\alpha$, com $\cos\alpha=\dfrac{\ell}{\sqrt{r^{2}+\ell^{2}}}$:
\[
E=2E_1\cos\alpha
=\frac{2\ke q\,\ell}{\left(r^{2}+\ell^{2}\right)^{3/2}}.
\]
Para $r\gg\ell$, $\left(r^{2}+\ell^{2}\right)^{3/2}\approx r^{3}$, e como
$p=2q\ell$:
\[
\boxed{\;E\approx\frac{\ke\,p}{r^{3}}\;}
\]
(b) Porque a carga \emph{total} do dipolo é nula. De longe, os campos das duas
cargas quase se cancelam --- sobra apenas o efeito da pequena separação entre
elas, que é de ordem superior. Daí o decaimento $1/r^{3}$ em vez de $1/r^{2}$.\\
(c) Em campo \textbf{uniforme}:
\begin{itemize}[leftmargin=1.6em,itemsep=1pt,topsep=2pt]
\item \textbf{Força resultante nula}: as forças $+qE$ e $-qE$ sobre as duas cargas
      têm mesmo módulo e sentidos opostos.
\item \textbf{Torque não nulo}: como elas atuam em \emph{pontos diferentes},
      formam um binário de momento
      \[
      \tau = p\,E\sin\theta,
      \]
      que tende a \emph{alinhar} o dipolo com o campo.
\end{itemize}
\textbf{Aplicações:} é esse torque que alinha moléculas polares (como a água) em
um campo --- princípio do forno de micro-ondas, onde o campo oscilante faz as
moléculas girarem e o atrito molecular gera calor. Em campo \emph{não} uniforme,
surge também força resultante --- razão pela qual um objeto neutro é atraído por
um corpo eletrizado.}
\stepcounter{questao}
\resposta{$E=\SI{300}{\kilo\newton\per\coulomb}$}
\resolucao{
\[
E=\frac{\ke Q}{r^{2}}=\frac{\pot{9}{9}\times\pot{3}{-6}}{(0{,}30)^{2}}
=\frac{\pot{2.7}{4}}{0{,}09}=\pot{3}{5}\,\si{\newton\per\coulomb}.
\]
O campo aponta \textbf{para longe} da carga, por ela ser positiva.\\
\textbf{Atenção à potência:} o campo cai com $1/r^{2}$, o potencial com
$1/r$. Dobrar a distância divide $E$ por quatro e $V$ por dois.}
\stepcounter{questao}
\resposta{$E=\SI{180}{\kilo\newton\per\coulomb}$, apontando \emph{para} a carga}
\resolucao{
\[
E=\frac{\pot{9}{9}\times\pot{5}{-6}}{0{,}25}
=\pot{1.8}{5}\,\si{\newton\per\coulomb}.
\]
\textbf{Sentido:} cargas negativas \emph{atraem} cargas de prova positivas, logo
o campo aponta \textbf{para} a carga fonte.\\
\textbf{Regra mnemônica:} o campo \emph{sai} do positivo e \emph{entra} no
negativo --- é a mesma regra que orienta as linhas de campo.}
\stepcounter{questao}
\resposta{$F=\SI{10}{\milli\newton}$, no sentido do campo}
\resolucao{
\[
F=qE=\pot{2}{-6}\times5000=\pot{1}{-2}\,\si{\newton}=\SI{10}{\milli\newton}.
\]
Como a carga é positiva, a força tem o \emph{mesmo} sentido de $\vec E$. Se
fosse negativa, teria sentido oposto --- e apenas o sentido mudaria, não o
módulo.}
\stepcounter{questao}
\resposta{$Q=\SI{20}{\nano\coulomb}$}
\resolucao{
\[
Q=\frac{Er^{2}}{\ke}=\frac{2000\times0{,}09}{\pot{9}{9}}
=\pot{2}{-8}\,\si{\coulomb}=\SI{20}{\nano\coulomb}.
\]
Uma carga de nanocoulombs --- ordem de grandeza típica da eletrização por
atrito --- já produz milhares de newtons por coulomb a trinta centímetros.}
\stepcounter{questao}
\resposta{$r=\SI{0.6}{\meter}$}
\resolucao{
\[
r=\sqrt{\frac{\ke Q}{E}}=\sqrt{\frac{\pot{9}{9}\times\pot{8}{-6}}{\pot{2}{5}}}
=\sqrt{0{,}36}=\SI{0.6}{\meter}.
\]
\textbf{Superfície de campo constante:} todos os pontos a
$\SI{0.6}{\meter}$ da carga têm o mesmo \emph{módulo} de campo --- mas
\emph{direções} diferentes, pois o campo é radial. Diferentemente das
equipotenciais, essas superfícies não têm significado especial.}
\stepcounter{questao}
\resposta{(c)}
\resolucao{Por definição, $\vec E=\vec F/q_0$. Dobrar $q_0$ dobra a força
\emph{e} o denominador --- a razão permanece.\\
\textbf{Ponto conceitual:} o campo é uma propriedade \textbf{do ponto do
espaço}, criada pelas cargas-fonte. A carga de prova apenas o \emph{revela};
ela não o cria nem o modifica.\\
\textbf{Ressalva real:} se $q_0$ for grande demais, ela perturba as
cargas-fonte (por indução, em condutores) e a medida deixa de valer. Daí a
definição rigorosa usar o limite $q_0\to0$.}
\stepcounter{questao}
\resposta{$\si{\volt}=\si{\joule\per\coulomb}$ e
$\si{\joule}=\si{\newton\meter}$}
\resolucao{
\[
\frac{\si{\volt}}{\si{\meter}}
=\frac{\si{\joule}}{\si{\coulomb}\cdot\si{\meter}}
=\frac{\si{\newton}\cdot\si{\meter}}{\si{\coulomb}\cdot\si{\meter}}
=\frac{\si{\newton}}{\si{\coulomb}}.\ \checkmark
\]
\textbf{As duas leituras do campo.} $\si{\newton\per\coulomb}$ enfatiza a
\emph{força} por unidade de carga; $\si{\volt\per\meter}$ enfatiza a
\emph{taxa de queda do potencial} por unidade de comprimento. São descrições
equivalentes de $\vec E=-\nabla V$.\\
Na prática, usa-se $\si{\volt\per\meter}$ quando há tensões envolvidas (placas,
cabos) e $\si{\newton\per\coulomb}$ em problemas de dinâmica de partículas.}
\stepcounter{questao}
\resposta{(b)}
\resolucao{A força vale $\vec F=q\vec E$ com $q<0$, logo $\vec F$ é
\textbf{antiparalela} a $\vec E$.\\
\textbf{Consequência energética:} a carga negativa move-se espontaneamente no
sentido de $V$ \emph{crescente} --- o oposto de uma carga positiva. Em ambos os
casos, a energia potencial $U=qV$ \emph{diminui}, como deve ocorrer em qualquer
movimento espontâneo.}
\stepcounter{questao}
\resposta{$F=\SI{2}{\milli\newton}$, oposta a $\vec E$;
$a=\SI{0.5}{\meter\per\second\squared}$}
\resolucao{(a) $F=|q|E=\pot{1}{-6}(2000)=\pot{2}{-3}\,\si{\newton}$, com sentido
\emph{oposto} ao campo.\\
(b)
\[
a=\frac{F}{m}=\frac{\pot{2}{-3}}{\pot{4}{-3}}=\SI{0.5}{\meter\per\second\squared}.
\]
\textbf{Comparação com a gravidade:} $a$ é apenas $5\%$ de $g$. Para objetos
macroscópicos, a força elétrica só compete com o peso se a carga for
relativamente grande ou a massa muito pequena --- daí os experimentos
eletrostáticos clássicos usarem pedaços de papel e gotas de óleo.}
\stepcounter{questao}
\resposta{$a=\pot{9.6}{11}\,\si{\meter\per\second\squared}$}
\resolucao{
\[
a=\frac{eE}{m_p}=\frac{(\pot{1.6}{-19})(\pot{1}{4})}{\pot{1.67}{-27}}
=\pot{9.58}{11}\,\si{\meter\per\second\squared}.
\]
\textbf{Cerca de $10^{11}$ vezes $g$} --- a gravidade é completamente
desprezível na dinâmica de partículas carregadas.\\
\textbf{Para o elétron}, a aceleração seria $1836$ vezes maior
($\pot{1.76}{15}\,\si{\meter\per\second\squared}$), pela razão das massas. É por
isso que, em tubos e aceleradores, são os elétrons que respondem primeiro.}
\stepcounter{questao}
\resposta{$E=\SI{100}{\kilo\newton\per\coulomb}$, no sentido $-x$}
\resolucao{
\[
E_1=\frac{\pot{9}{9}(\pot{4}{-6})}{(0{,}60)^{2}}
=\pot{1}{5}\,\si{\newton\per\coulomb}\ (\text{sentido }+x);
\]
\[
E_2=\frac{\pot{9}{9}(\pot{2}{-6})}{(0{,}30)^{2}}
=\pot{2}{5}\,\si{\newton\per\coulomb}\ (\text{sentido }-x).
\]
\[
E=2\times10^{5}-1\times10^{5}=\pot{1}{5}\,\si{\newton\per\coulomb}\ (-x).
\]
\textbf{Observe:} a carga \emph{menor} produziu o campo \emph{maior}, por estar
duas vezes mais próxima --- e $1/r^{2}$ vence o fator $2$ de carga.
Compare os fatores: carga $\times\frac12$, distância$^{2}$
$\times\frac14$.}
\stepcounter{questao}
\resposta{$E=\SI{115.2}{\kilo\newton\per\coulomb}$, ao longo de $+y$}
\resolucao{Cada carga dista $r=\sqrt{0{,}30^{2}+0{,}40^{2}}=\SI{0.5}{\meter}$
(triângulo $3$--$4$--$5$):
\[
E_1=E_2=\frac{\pot{9}{9}(\pot{2}{-6})}{0{,}25}
=\pot{7.2}{4}\,\si{\newton\per\coulomb}.
\]
\textbf{Por simetria}, as componentes horizontais se cancelam e as verticais se
somam, cada uma multiplicada por $\cos\theta=0{,}40/0{,}50=0{,}8$:
\[
E=2(\pot{7.2}{4})(0{,}8)=\pot{1.152}{5}\,\si{\newton\per\coulomb}.
\]
\textbf{Uso da simetria:} reconhecê-la elimina metade do trabalho. Sempre que a
configuração for simétrica em relação ao ponto de interesse, identifique quais
componentes se anulam \emph{antes} de calcular.}
\stepcounter{questao}
\resposta{$E=\SI{300}{\kilo\volt\per\meter}$}
\resolucao{
\[
E=\frac{U}{d}=\frac{600}{\pot{2}{-3}}=\pot{3}{5}\,\si{\volt\per\meter}.
\]
\textbf{Campo uniforme:} entre placas extensas e próximas, $E$ é praticamente o
mesmo em todos os pontos --- e independe da posição. É a configuração de
referência para acelerar e defletir partículas.\\
\textbf{Verificação de segurança:} $\SI{300}{\kilo\volt\per\meter}$ é um décimo
da rigidez do ar ($\SI{3}{\mega\volt\per\meter}$) --- há folga confortável.
Reduzir a separação para $\SI{0.2}{\milli\meter}$ com a mesma tensão provocaria
ruptura.}
\stepcounter{questao}
\resposta{$E=\SI{226}{\kilo\newton\per\coulomb}$}
\resolucao{Para um condutor em equilíbrio,
\[
E=\frac{\sigma}{\varepsilon_0}=\frac{\pot{2}{-6}}{\pot{8.85}{-12}}
=\pot{2.26}{5}\,\si{\newton\per\coulomb}.
\]
\textbf{Fato notável:} o campo \textbf{não depende da distância} à placa (desde
que ela seja extensa e o ponto próximo). Diferentemente da carga puntiforme, não
há queda com $1/r^{2}$.\\
\textbf{Cuidado com a fórmula:} para um \emph{plano isolante} com carga nas duas
faces, o campo vale $\sigma/2\varepsilon_0$ --- metade. A diferença é que no
condutor toda a carga fica em uma face, e o campo interno é nulo.}
\stepcounter{questao}
\resposta{$E_1=E_2=\SI{900}{\kilo\newton\per\coulomb}$ --- iguais}
\resolucao{
\[
E_1=\frac{\pot{9}{9}(\pot{9}{-6})}{0{,}09}=\pot{9}{5};
\qquad
E_2=\frac{\pot{9}{9}(\pot{4}{-6})}{0{,}04}=\pot{9}{5}\,\si{\newton\per\coulomb}.
\]
\textbf{Coincidência não é acaso:} as razões $Q/r^{2}$ são iguais, pois
$9/9=4/4$. Sempre que $Q\propto r^{2}$, o campo se mantém.\\
\textbf{Consequência:} se essas duas cargas estivessem em lados opostos de um
ponto, o campo resultante ali seria \emph{nulo} --- e essa é exatamente a
condição do "ponto de campo nulo" entre cargas de mesmo sinal.}
\stepcounter{questao}
\resposta{$E=\SI{5}{\kilo\newton\per\coulomb}$, a $\ang{53.1}$ do eixo $x$}
\resolucao{Sendo perpendiculares, compõem-se por Pitágoras:
\[
E=\sqrt{3^{2}+4^{2}}=\SI{5}{\kilo\newton\per\coulomb};
\qquad
\theta=\arctan\frac{4}{3}=\ang{53.1}.
\]
\textbf{Contraste com o potencial:} se esses mesmos pontos tivessem potenciais
de $3$ e $\SI{4}{\kilo\volt}$, o total seria simplesmente
$\SI{7}{\kilo\volt}$ --- soma algébrica, sem geometria.\\
É essa diferença que torna o potencial a ferramenta preferida para cálculos, e o
campo a ferramenta preferida para entender forças e movimento.}
\stepcounter{questao}
\resposta{(b)}
\resolucao{A \textbf{densidade de linhas} é proporcional ao módulo do campo ---
essa é a convenção que dá utilidade quantitativa ao desenho.\\
\textbf{Três regras das linhas de campo:}
\begin{enumerate}
\item Nascem em cargas positivas e morrem em negativas (ou no infinito).
\item São tangentes a $\vec E$ em cada ponto.
\item \textbf{Nunca se cruzam} --- em um cruzamento, o campo teria duas
direções, o que é impossível.
\end{enumerate}
\textbf{Cuidado com (a):} potencial e campo são grandezas distintas. Linhas
próximas indicam campo intenso, o que corresponde a equipotenciais próximas ---
mas não diz nada sobre o \emph{valor} de $V$.}
\stepcounter{questao}
\resposta{$E=\SI{2}{\kilo\newton\per\coulomb}$}
\resolucao{
\[
E_{meio}=\frac{E_{\text{vácuo}}}{\varepsilon_r}=\frac{8}{4}
=\SI{2}{\kilo\newton\per\coulomb}.
\]
\textbf{Mecanismo:} o dielétrico se \emph{polariza} --- suas moléculas se
orientam e criam um campo interno oposto, que cancela parcialmente o campo
original.\\
\textbf{Consequência tecnológica:} o dielétrico reduz o campo para a mesma
carga, o que permite armazenar \emph{mais} carga para a mesma tensão --- é
exatamente por isso que $C=\varepsilon_r C_0$ em capacitores.}
\stepcounter{questao}
\resposta{$E=\SI{19.6}{\kilo\newton\per\coulomb}$}
\resolucao{Equilíbrio entre peso e força elétrica:
\[
qE=mg
\;\Longrightarrow\;
E=\frac{mg}{q}=\frac{(\pot{2}{-6})(9{,}8)}{\pot{1}{-9}}
=\pot{1.96}{4}\,\si{\newton\per\coulomb}.
\]
\textbf{Orientação:} o campo deve produzir força \emph{para cima}. Se a gota for
positiva, $\vec E$ aponta para cima; se negativa, para baixo.\\
\textbf{Este é o princípio do experimento de Millikan}, que determinou a carga
elementar. Medindo $E$ e $m$ de gotas suspensas, ele encontrou que todas as
cargas eram múltiplas inteiras de $\pot{1.6}{-19}\,\si{\coulomb}$ --- a
primeira evidência direta da quantização.}
\stepcounter{questao}
\resposta{(a) $E=\SI{159}{\kilo\newton\per\coulomb}$ ao longo de $+y$;
(b) $\vec E=\vec0$}
\resolucao{(a) Cada carga dista
$r=\sqrt{2}\,(0{,}20)=\SI{0.283}{\meter}$:
\[
E_1=E_2=\frac{\pot{9}{9}(\pot{1}{-6})}{0{,}08}
=\pot{1.125}{5}\,\si{\newton\per\coulomb}.
\]
As componentes em $x$ cancelam; as em $y$ somam-se, cada uma multiplicada por
$\cos\ang{45}=0{,}707$:
\[
E=2(\pot{1.125}{5})(0{,}707)=\pot{1.59}{5}\,\si{\newton\per\coulomb}.
\]
(b) \textbf{No ponto médio}, as duas contribuições têm módulos iguais e sentidos
\emph{opostos} (cada carga empurra para longe de si):
\[
\vec E=\vec 0.
\]
\textbf{Mas o potencial ali não é nulo:}
$V=2\ke Q/0{,}20=\SI{90}{\kilo\volt}$. É o caso complementar ao do dipolo, em
que $V=0$ e $\vec E\ne\vec0$ --- e mostra novamente que as duas condições são
independentes.\\
\textbf{Estabilidade:} uma carga positiva no ponto médio está em equilíbrio, mas
\emph{instável} na direção perpendicular ao eixo --- conforme o teorema de
Earnshaw.}
\stepcounter{questao}
\resposta{$E=\dfrac{2\ke qa}{(y^{2}+a^{2})^{3/2}}\to\dfrac{2\ke qa}{y^{3}}$}
\resolucao{(a) Cada carga dista $r=\sqrt{y^{2}+a^{2}}$ e produz
$\ke q/r^{2}$. Por simetria, as componentes ao longo de $y$ se cancelam (uma
aponta para cima, a outra para baixo) e restam as componentes \emph{paralelas ao
eixo do dipolo}, cada uma multiplicada por $a/r$:
\[
E=2\cdot\frac{\ke q}{y^{2}+a^{2}}\cdot\frac{a}{\sqrt{y^{2}+a^{2}}}
=\frac{2\ke qa}{(y^{2}+a^{2})^{3/2}}.
\]
(b) Para $y\gg a$, $(y^{2}+a^{2})^{3/2}\to y^{3}$:
\[
E\approx\frac{2\ke qa}{y^{3}}=\frac{\ke p}{y^{3}},
\]
com $p=2qa$ o momento dipolar.\\
(c) \textbf{O campo do dipolo cai com $1/y^{3}$}, muito mais rápido que o
$1/r^{2}$ de uma carga isolada.\\
\textbf{Razão física:} de longe, as duas cargas quase se cancelam --- o sistema
é \emph{globalmente neutro}. O que sobra é o efeito de sua \emph{separação},
uma correção de ordem superior.\\
\textbf{Padrão geral:} carga isolada $\sim1/r^{2}$; dipolo
$\sim1/r^{3}$; quadrupolo $\sim1/r^{4}$. Quanto mais neutra a distribuição, mais
rápido o campo se extingue --- e é por isso que a matéria comum, neutra, não
exerce força elétrica a distância.}
\stepcounter{questao}
\resposta{Entre: $E=\SI{113}{\kilo\newton\per\coulomb}$; fora: $E=0$}
\resolucao{(a) Cada plano contribui com $\sigma/2\varepsilon_0$, e
\textbf{entre} as placas as duas contribuições têm o \emph{mesmo} sentido (do
positivo para o negativo):
\[
E=2\cdot\frac{\sigma}{2\varepsilon_0}=\frac{\sigma}{\varepsilon_0}
=\frac{\pot{1}{-6}}{\pot{8.85}{-12}}=\pot{1.13}{5}\,\si{\newton\per\coulomb}.
\]
(b) \textbf{Fora}, as duas contribuições têm sentidos \emph{opostos} e módulos
iguais:
\[
E=\frac{\sigma}{2\varepsilon_0}-\frac{\sigma}{2\varepsilon_0}=0.
\]
(c) \textbf{Justificativa:} o campo de um plano infinito é uniforme e aponta
para longe dele (se positivo) ou para ele (se negativo). Na região externa,
qualquer ponto está do mesmo lado das duas placas --- e os campos se opõem.
Entre elas, o ponto está "depois" da positiva e "antes" da negativa, e os campos
se somam.\\
\textbf{Consequência: blindagem natural.} O par de placas confina completamente
seu campo na região interna --- não há campo escapando para fora. É por isso que
um capacitor de placas paralelas não interfere no circuito vizinho, e por que a
energia armazenada está inteiramente entre as armaduras.}
\stepcounter{questao}
\resposta{$0$; $\SI{4.5}{\mega\volt\per\meter}$;
$\SI{500}{\kilo\volt\per\meter}$ --- \textbf{insustentável}}
\resolucao{(a) \textbf{Interior:} $E=0$, condição de equilíbrio de qualquer
condutor.\\
\textbf{Superfície:} toda a carga se comporta como concentrada no centro, para
pontos externos:
\[
E=\frac{\ke Q}{R^{2}}=\frac{\pot{9}{9}(\pot{5}{-6})}{0{,}01}
=\pot{4.5}{6}\,\si{\volt\per\meter}.
\]
\textbf{A $\SI{30}{\centi\meter}$:}
$E=\dfrac{\pot{4.5}{4}}{0{,}09}=\pot{5}{5}\,\si{\volt\per\meter}$.\\
(b) \textbf{Insustentável.} O campo na superfície,
$\SI{4.5}{\mega\volt\per\meter}$, excede a rigidez dielétrica do ar
($\SI{3}{\mega\volt\per\meter}$) em $50\%$. O ar ioniza e a carga escapa por
descarga corona.\\
\textbf{Carga máxima real:}
$Q_{\max}=E_{rup}R^{2}/\ke=\SI{3.33}{\micro\coulomb}$. A esfera simplesmente
\emph{não consegue} reter os $\SI{5}{\micro\coulomb}$ enunciados.\\
\textbf{Lição de método:} sempre verifique se o resultado é fisicamente
possível. Um problema pode ser aritmeticamente consistente e descrever uma
situação inexistente.}
\stepcounter{questao}
\resposta{$a=\pot{1.76}{14}\,\si{\meter\per\second\squared}$;
$t=\SI{15.1}{\nano\second}$; $v=\pot{2.65}{6}\,\si{\meter\per\second}$;
$\SI{20}{\electronvolt}$}
\resolucao{(a)
\[
a=\frac{eE}{m_e}=\frac{(\pot{1.6}{-19})(1000)}{\pot{9.11}{-31}}
=\pot{1.756}{14}\,\si{\meter\per\second\squared};
\]
\[
t=\sqrt{\frac{2d}{a}}=\sqrt{\frac{2(0{,}02)}{\pot{1.756}{14}}}
=\pot{1.51}{-8}\,\si{\second}=\SI{15.1}{\nano\second}.
\]
(b) $v=at=(\pot{1.756}{14})(\pot{1.51}{-8})
=\pot{2.65}{6}\,\si{\meter\per\second}$.\\
\textbf{Energia:} o mais rápido é usar a tensão percorrida,
$U=Ed=1000(0{,}02)=\SI{20}{\volt}$:
\[
W=eU=\SI{20}{\electronvolt}=\pot{3.2}{-18}\,\si{\joule}.
\]
\emph{(Conferindo: $\frac12m_ev^{2}=\frac12(\pot{9.11}{-31})(\pot{2.65}{6})^{2}
=\pot{3.2}{-18}\,\si{\joule}$ \checkmark.)}\\
\textbf{Verificação relativística:} $v/c=0{,}9\%$ --- plenamente no regime
clássico.}
\stepcounter{questao}
\resposta{$1/r^{2}$, $1/r$ e constante}
\resolucao{(a)
\begin{center}
\begin{tabular}{lcc}
\hline
Distribuição & Campo & Queda\\
\hline
Carga puntiforme & $\dfrac{\ke Q}{r^{2}}$ & $1/r^{2}$\\[6pt]
Fio infinito ($\lambda$) & $\dfrac{\lambda}{2\pi\varepsilon_0 r}$ & $1/r$\\[6pt]
Plano infinito ($\sigma$) & $\dfrac{\sigma}{2\varepsilon_0}$ & constante\\
\hline
\end{tabular}
\end{center}
(b) \textbf{Origem geométrica.} Pela lei de Gauss, o fluxo total é fixo pela
carga envolvida, e o campo é o fluxo dividido pela área da superfície:
\begin{itemize}
\item \textbf{Ponto:} a superfície é uma \emph{esfera}, de área
$4\pi r^{2}$ --- o campo cai com $1/r^{2}$.
\item \textbf{Fio:} a superfície é um \emph{cilindro}, de área lateral
$2\pi r\ell$ --- cresce apenas com $r$, e o campo cai com $1/r$.
\item \textbf{Plano:} a superfície é um \emph{prisma} cuja área \emph{não
depende} de $r$ --- e o campo não cai.
\end{itemize}
\textbf{Regra unificadora:} o campo cai como $1/r^{n}$, onde $n$ é o número de
dimensões em que a fonte é \emph{limitada}. Ponto: limitado em três
($n=2$). Fio: em duas ($n=1$). Plano: em uma ($n=0$).\\
\textbf{Consequência prática:} perto de uma placa carregada extensa, afastar-se
não ajuda --- o campo permanece. É o oposto da intuição formada com cargas
puntiformes.}
\stepcounter{questao}
\resposta{Não necessariamente --- só vale se a metade removida contribuísse com
exatamente metade do campo}
\resolucao{(a) \textbf{Não, em geral.} O campo é a soma \emph{vetorial} das
contribuições de todos os elementos da haste, e cada elemento contribui com
módulo e direção diferentes, dependendo de sua distância e posição angular em
relação ao ponto.\\
Se o ponto estiver \emph{sobre a mediatriz} da haste, as duas metades
contribuem igualmente e o campo de fato cai à metade --- mas a \emph{direção}
resultante também muda, pois as componentes que antes se cancelavam agora não
se cancelam mais.\\
Se o ponto estiver \emph{sobre o prolongamento} da haste, a metade mais próxima
contribui muito mais que a distante, e remover a distante reduz o campo bem
menos que $50\%$.\\
(b) \textbf{A queda é exatamente à metade quando:}
\begin{itemize}
\item as duas metades contribuem com vetores \emph{idênticos} --- o que exige
simetria completa em relação ao ponto; \emph{e}
\item essa condição só se realiza em pontos muito distantes, onde a haste inteira
se comporta como carga puntiforme e a divisão é simplesmente $Q\to Q/2$.
\end{itemize}
\textbf{Lição:} a superposição vale \emph{sempre}, mas ela é vetorial. Reduzir
a fonte pela metade não reduz o campo pela metade, exceto em casos de simetria
particular.}
\stepcounter{questao}
\resposta{Duas cargas iguais $+Q$; no ponto médio $\vec E=\vec0$ e
$V=2\ke Q/a$}
\resolucao{(a) \textbf{Arranjo:} duas cargas iguais $Q=+\SI{2}{\micro\coulomb}$
em $(\pm\SI{0.20}{\meter},\,0)$. No ponto médio:
\[
\vec E=\vec 0
\quad\text{(contribuições opostas)};
\]
\[
V=2\cdot\frac{\ke Q}{0{,}20}
=2\cdot\frac{\pot{9}{9}(\pot{2}{-6})}{0{,}20}
=\pot{1.8}{5}\,\si{\volt}=\SI{180}{\kilo\volt}.
\]
Campo nulo, potencial de cento e oitenta mil volts.\\
(b) \textbf{Por que é possível.} As duas grandezas se relacionam por
$\vec E=-\nabla V$: o campo é a \emph{taxa de variação} de $V$, não seu valor.
Uma função pode ter valor grande e derivada nula --- é exatamente o que ocorre
em um ponto \emph{estacionário}.\\
No ponto médio, $V$ tem um mínimo ao longo do eixo que une as cargas: o valor é
alto, mas a inclinação é zero.\\
(c) \textbf{Sim, o inverso também ocorre.} Substituindo uma das cargas por
$-Q$, o ponto médio passa a ter $V=0$ (contribuições opostas se cancelam
algebricamente) e $\vec E\ne\vec0$ (os vetores apontam para o mesmo lado e se
somam).\\
\textbf{Síntese:} $V=0$ e $\vec E=\vec0$ são condições \textbf{logicamente
independentes}. Nenhuma implica a outra, em nenhum dos dois sentidos.}
\stepcounter{questao}
\resposta{Saem de $+q$ e entram em $-q$; não se cruzam porque $\vec E$ tem
direção única em cada ponto}
\resolucao{(a) \textbf{Traçado.} As linhas saem radialmente de $+q$, curvam-se e
entram em $-q$. A linha que segue o eixo do dipolo é reta; as demais formam
arcos progressivamente mais abertos. Algumas partem para o infinito e retornam.\\
A densidade de linhas é máxima entre as cargas --- onde o campo é mais intenso
--- e decresce com $1/r^{3}$ ao se afastar, conforme demonstrado antes.\\
(b) \textbf{Demonstração.} Suponha, por absurdo, que duas linhas se cruzem em
um ponto $P$. Por definição, a linha de campo é \emph{tangente} a $\vec E$ em
cada ponto. Se duas linhas distintas passam por $P$, elas têm tangentes
distintas ali --- e portanto $\vec E(P)$ teria \textbf{duas direções}
simultaneamente.\\
Mas $\vec E$ é uma função vetorial \emph{bem definida}: a cada ponto do espaço
corresponde \emph{um} vetor. Contradição. Logo, as linhas não se cruzam.\\
\textbf{Corolário:} o mesmo argumento vale para linhas de corrente, linhas de
campo magnético e qualquer outro campo vetorial.\\
(c) \textbf{Pontos de campo nulo são a exceção.} Onde $\vec E=\vec0$, não há
direção definida --- e ali as linhas \emph{podem} se encontrar, sem contradição.
São os \textbf{pontos de sela} do potencial.\\
Em um dipolo ($+q$ e $-q$) não existe tal ponto no espaço finito. Mas em duas
cargas \emph{iguais}, o ponto médio é exatamente um ponto de campo nulo --- e
ali as linhas se aproximam de todos os lados sem se tocarem, definindo uma
separatriz.}
\stepcounter{questao}
\resposta{$E=\SI{11.4}{\volt\per\meter}$; $U=\SI{0.114}{\volt}$;
$t=\SI{100}{\nano\second}$; $v=\pot{2}{5}\,\si{\meter\per\second}$}
\resolucao{(a)
\[
E=\frac{m_ea}{e}=\frac{(\pot{9.11}{-31})(\pot{2}{12})}{\pot{1.6}{-19}}
=\SI{11.4}{\volt\per\meter}.
\]
(b) $U=Ed=11{,}4(0{,}01)=\SI{0.114}{\volt}$.\\
(c)
\[
t=\sqrt{\frac{2d}{a}}=\sqrt{\frac{0{,}02}{\pot{2}{12}}}
=\pot{1}{-7}\,\si{\second}=\SI{100}{\nano\second};
\]
\[
v=at=(\pot{2}{12})(\pot{1}{-7})=\pot{2}{5}\,\si{\meter\per\second}.
\]
\textbf{O resultado é surpreendente e instrutivo.} Uma aceleração de
$\pot{2}{12}\,\si{\meter\per\second\squared}$ --- duzentos bilhões de vezes
$g$ --- exige apenas $\SI{0.114}{\volt}$, tensão menor que a de uma pilha
descarregada.\\
\textbf{Razão:} a razão carga/massa do elétron é enorme
($\pot{1.76}{11}\,\si{\coulomb\per\kilogram}$). Campos elétricos modestos
produzem acelerações colossais em partículas leves.\\
\textbf{Ressalva prática:} com tensão tão baixa, o projeto seria dominado por
efeitos parasitas --- diferenças de função-trabalho entre os metais das placas
(décimos de volt), cargas superficiais e ruído térmico. Um projeto real usaria
tensões de dezenas de volts e aceleração proporcionalmente maior.}
\stepcounter{questao}
\resposta{$-q$ na parede da cavidade, $+q$ na superfície externa; campo nulo no
metal e não nulo fora}
\resolucao{(a) \textbf{Parede da cavidade:} aplicando a lei de Gauss a uma
superfície inteiramente dentro do metal (onde $\vec E=\vec0$), o fluxo é nulo e
a carga envolvida também. Como a superfície envolve $+q$ e a parede interna:
\[
q+q_{interna}=0
\;\Longrightarrow\;
q_{interna}=-q.
\]
\textbf{Superfície externa:} o condutor era neutro, então a carga total deve
permanecer nula:
\[
q_{externa}=+q.
\]
(b) \textbf{Dentro do metal:} $\vec E=\vec0$, por equilíbrio.\\
\textbf{Fora:} a carga $+q$ da superfície externa produz campo --- e ele é
\textbf{esfericamente simétrico} se a superfície externa for esférica,
\emph{independentemente} de onde a carga esteja dentro da cavidade.\\
\textbf{Consequência notável:} de fora, é impossível saber a posição da carga
interna, nem a forma da cavidade. O condutor "apaga" essa informação.\\
(c) \textbf{Se o condutor for aterrado:} a carga $+q$ da superfície externa
escoa para a terra. Resta apenas $-q$ na parede da cavidade, e
\[
\vec E=\vec 0\quad\text{em todo o exterior}.
\]
\textbf{A blindagem torna-se completa nos dois sentidos.} É exatamente esse o
princípio da gaiola de Faraday aterrada --- e a diferença entre blindar contra
campos externos (que dispensa aterramento) e impedir que fontes internas
irradiem (que o exige).}
\stepcounter{questao}
\resposta{$E=\dfrac{\ke Qr}{R^{3}}$ dentro e $\dfrac{\ke Q}{r^{2}}$ fora;
máximo na superfície}
\resolucao{(a) \textbf{Fora ($r>R$):} pela lei de Gauss com superfície esférica,
toda a carga está envolvida:
\[
E=\frac{\ke Q}{r^{2}},
\]
idêntico ao de uma carga puntiforme no centro.\\
\textbf{Dentro ($r<R$):} apenas a fração de carga interna à superfície gaussiana
conta. Como a densidade é uniforme, essa fração é $(r/R)^{3}$:
\[
E=\frac{\ke\,Q(r/R)^{3}}{r^{2}}=\frac{\ke Q\,r}{R^{3}},
\]
\textbf{proporcional a $r$} --- cresce linearmente do centro à superfície.\\
(b) \textbf{Máximo na superfície} ($r=R$), onde as duas expressões coincidem em
$E=\ke Q/R^{2}$. O campo é contínuo ali, embora com derivadas diferentes de
cada lado.\\
\emph{No centro, $E=0$ por simetria.}\\
(c) \textbf{Comparação com a esfera condutora.}
\begin{center}
\begin{tabular}{lcc}
\hline
Região & Isolante & Condutora\\
\hline
$r<R$ & $\ke Qr/R^{3}$ & $0$\\
$r=R$ & $\ke Q/R^{2}$ & $\ke Q/R^{2}$\\
$r>R$ & $\ke Q/r^{2}$ & $\ke Q/r^{2}$\\
\hline
\end{tabular}
\end{center}
\textbf{São indistinguíveis do lado de fora} --- a lei de Gauss só "enxerga" a
carga total envolvida, não sua distribuição. A diferença está inteiramente no
interior: no condutor a carga migra para a superfície e o campo interno é nulo;
no isolante ela permanece onde foi colocada.}
\stepcounter{questao}
\resposta{Integrar $\dfrac{\mathrm{d}\vec r}{\mathrm{d}s}
=\dfrac{\vec E}{|\vec E|}$ a partir de pontos distribuídos em torno das cargas}
\resolucao{(a) \textbf{Algoritmo.} A linha de campo é, por definição, tangente a
$\vec E$. Parametrizando pelo comprimento de arco $s$:
\[
\frac{\mathrm{d}\vec r}{\mathrm{d}s}=\frac{\vec E(\vec r)}{|\vec E(\vec r)|}.
\]
Numericamente, com passo $h$:
\begin{enumerate}
\item Escolha um ponto inicial $\vec r_0$.
\item Calcule $\vec E(\vec r_n)$ e normalize.
\item Avance: $\vec r_{n+1}=\vec r_n+h\,\hat E$ (Euler) ou, melhor, por
Runge--Kutta de quarta ordem.
\item Repita até sair do domínio ou atingir uma carga negativa.
\end{enumerate}
\textbf{Normalizar é essencial:} sem isso, o passo seria proporcional a
$|\vec E|$ e a linha "voaria" perto das cargas e travaria longe delas.\\
(b) \textbf{Pontos de partida.} Distribua-os em pequenos círculos ao redor de
cada carga \emph{positiva}, com espaçamento angular uniforme. Para que a
densidade de linhas represente corretamente a intensidade, o número de linhas
por carga deve ser \emph{proporcional a $|q|$} --- oito linhas para $+q$,
dezesseis para $+2q$.\\
(c) \textbf{Dificuldades e soluções.}
\begin{itemize}
\item \textbf{Singularidades:} $|\vec E|\to\infty$ nas cargas. Solução: iniciar
a uma distância pequena mas finita, e interromper a integração ao entrar em um
raio mínimo em torno de uma carga negativa.
\item \textbf{Pontos de campo nulo:} ali $\hat E$ é indefinido e o algoritmo
trava ou oscila. Solução: detectar $|\vec E|<\epsilon$ e encerrar a linha.
\item \textbf{Acúmulo de erro:} o método de Euler desvia sistematicamente em
regiões de curvatura alta. Solução: Runge--Kutta e passo adaptativo, menor onde
o campo varia rápido.
\item \textbf{Linhas que escapam:} defina um domínio e encerre ao ultrapassá-lo.
\end{itemize}
\textbf{Verificação do resultado:} as linhas traçadas devem ser perpendiculares
às equipotenciais calculadas independentemente --- teste que detecta quase todo
erro de implementação.}
\stepcounter{questao}
\resposta{$U_{\max}=\SI{2}{\kilo\volt}$;
$E=\SI{2e7}{\volt\per\meter}$;
$\sigma=\SI{531}{\micro\coulomb\per\meter\squared}$}
\resolucao{(a)
\[
U_{\max}=E_{rup}\,d=20\times0{,}1=\SI{2}{\kilo\volt};
\qquad
E=\pot{2}{7}\,\si{\volt\per\meter}.
\]
(b) No dielétrico, $E=\dfrac{\sigma}{\varepsilon_r\varepsilon_0}$, logo
\[
\sigma=\varepsilon_r\varepsilon_0E
=3(\pot{8.85}{-12})(\pot{2}{7})
=\pot{5.31}{-4}\,\si{\coulomb\per\meter\squared}.
\]
(c) \textbf{O compromisso.} A capacitância por unidade de área é
$C/A=\varepsilon_r\varepsilon_0/d$: reduzir $d$ \emph{aumenta} $C$. Mas
$U_{\max}=E_{rup}d$: reduzir $d$ \emph{diminui} a tensão suportável.\\
O produto revela o que realmente importa:
\[
\frac{C\,U_{\max}}{A}=\varepsilon_r\varepsilon_0E_{rup},
\]
que \textbf{não depende de $d$}. E a energia por unidade de volume,
\[
u=\frac12\varepsilon_r\varepsilon_0E_{rup}^{2},
\]
também não. \textbf{A geometria não decide nada} --- só o material.\\
\textbf{Consequência:} não existe truque de projeto que aumente a energia
armazenada por volume. Todo avanço em capacitores vem de materiais com
$\varepsilon_r$ ou $E_{rup}$ maiores --- e esses dois parâmetros costumam ser
antagônicos, o que torna o problema difícil.}
\stepcounter{questao}
\resposta{Centro: $\vec E=\vec0$ por simetria; longe: $E\to\ke Q/x^{2}$;
com meio anel, campo não nulo no centro}
\resolucao{(a) \textbf{No centro,} para cada elemento de carga existe outro
diametralmente oposto, à mesma distância $R$. Suas contribuições têm módulos
iguais e sentidos opostos, cancelando-se aos pares:
\[
\vec E_{centro}=\vec 0.
\]
\textbf{Note:} o \emph{potencial} no centro não é nulo --- vale
$V=\ke Q/R$, pois potenciais somam-se sem cancelamento entre contribuições de
mesmo sinal. Mais um caso de $\vec E=\vec0$ com $V\ne0$.\\
(b) \textbf{Muito longe} ($x\gg R$), todos os elementos estão praticamente à
mesma distância $x$ e praticamente na mesma direção. O anel se comporta como uma
carga puntiforme:
\[
E\to\frac{\ke Q}{x^{2}}.
\]
\textbf{Este é o teste que toda expressão deve passar:} qualquer distribuição
finita, vista de longe, deve reproduzir a lei de Coulomb.\\
(c) \textbf{Com meio anel}, a simetria de cancelamento se perde. Cada elemento
não tem mais seu par oposto, e o campo no centro \textbf{deixa de ser nulo}.\\
Pelo argumento da superposição ao contrário: como o anel inteiro dá campo nulo,
o campo do semianel restante é o \emph{negativo} do campo do semianel removido
--- ou seja, os dois têm o mesmo módulo. Integrando (como no caso do fio
semicircular), obtém-se
\[
E=\frac{2\ke Q'}{\pi R^{2}},
\]
com $Q'=Q/2$ a carga de cada metade, dirigido ao longo do eixo de simetria.}

\gabarito[1.3 Campo elétrico]

\subsection{Potencial elétrico e superfícies equipotenciais}
\stepcounter{questao}
\resposta{(b)}
\resolucao{Por definição, o potencial é constante sobre uma equipotencial (isso
elimina c). Como não há variação de potencial ao longo dela, o trabalho para mover
uma carga sobre a superfície é nulo (elimina d). O campo, que aponta na direção de
maior variação do potencial, é sempre \emph{perpendicular} às equipotenciais ---
nunca tangente, o que elimina (a) e (e).}
\stepcounter{questao}
\resposta{$U_{AB} = \SI{500}{\volt}$}
\resolucao{Em campo uniforme, $U = E\,d = \pot{1}{4}\cdot0{,}05 = \SI{500}{\volt}$
(com A no potencial maior, a favor do campo).}
\stepcounter{questao}
\resposta{(a) $E = \SI{1000}{\volt\per\meter}$; (b) $W = \SI{2e-4}{\joule}$}
\resolucao{(a) Entre superfícies adjacentes, $U = \SI{20}{\volt}$ em
$d = \SI{2}{\centi\meter}$:
\[
E=\frac{U}{d}=\frac{20}{0{,}02}=\SI{1000}{\volt\per\meter}.
\]
(b) $U_{AC} = V_A - V_C = 60 - 20 = \SI{40}{\volt}$:
\[
W = q\,U_{AC} = \pot{5}{-6}\cdot40 = \pot{2}{-4}\,\si{\joule}=\SI{2e-4}{\joule}.
\]
O trabalho é positivo porque a carga positiva se move no sentido do campo
(de maior para menor potencial).}
\stepcounter{questao}
\resposta{$U_{AB} = \SI{20}{\volt}$; $U_{BC} = \SI{20}{\volt}$}
\resolucao{A ddp é sempre a diferença dos potenciais:
\[
U_{AB}=V_A-V_B=30-10=\SI{20}{\volt},
\qquad
U_{BC}=V_B-V_C=10-(-10)=\SI{20}{\volt}.
\]
Embora $U_{AB}=U_{BC}$, note que $V_C$ é \emph{negativo}: perto da carga negativa
o potencial fica abaixo de zero. O ponto médio B, equidistante das duas cargas de
módulos iguais, tem potencial intermediário.}
\stepcounter{questao}
\resposta{(b)}
\resolucao{A carga positiva se desloca de A (maior potencial) para C (menor
potencial), ou seja, \emph{a favor} do campo elétrico. O trabalho da força
elétrica é $W = q(V_A - V_C) > 0$. O trabalho \emph{não} depende do caminho (força
conservativa), o que elimina (d).}
\stepcounter{questao}
\resposta{(b)}
\resolucao{Linhas de força saindo indicam carga \textbf{positiva}. O potencial de
uma carga puntiforme, $V = \ke Q/d$, \emph{diminui} à medida que a distância
cresce: $V_A > V_C$. As equipotenciais, mais afastadas, têm menor potencial. O
campo não é uniforme --- é radial e enfraquece com $1/d^2$.}
\stepcounter{questao}
\resposta{(b)}
\resolucao{As linhas de força \emph{nascem} nas cargas positivas e \emph{morrem}
nas negativas. Como saem de A, esta é \textbf{positiva}; como entram em B, esta é
\textbf{negativa}. A configuração é a de um dipolo elétrico.}
\stepcounter{questao}
\resposta{(c)}
\resolucao{Para uma carga puntiforme, $E = \dfrac{\ke Q}{d^2}$ e
$V = \dfrac{\ke Q}{d}$. Dividindo, $\dfrac{V}{E} = d$, logo
\[
V = E\cdot d = 900\cdot0{,}20 = \SI{180}{\volt}.
\]
Essa relação $V = E\,d$ vale para carga puntiforme; não confundir com o campo
uniforme, onde também $U = E\,d$, mas por outro motivo.}
\stepcounter{questao}
\resposta{(c)}
\resolucao{Em campo uniforme, $U = E\,d = 120\cdot10 = \SI{1200}{\volt}$. Como o
campo aponta para baixo (de M para N), M está a maior potencial:
$V_M - V_N = \SI{1200}{\volt}$.}
\stepcounter{questao}
\resposta{$V_M = 0$}
\resolucao{O potencial é uma grandeza \emph{escalar}: soma-se com sinal. No ponto
médio, cada carga está a $d = \SI{10}{\centi\meter}$:
\[
V_M = \ke\frac{Q_1}{d}+\ke\frac{Q_2}{d}
= \ke\frac{(+2)+(-2)}{d}\,\si{\micro\coulomb}=0.
\]
O potencial é nulo, mas o \emph{campo} em M \emph{não} é zero --- os vetores campo
das duas cargas apontam no mesmo sentido e se somam. Este é o ponto central: campo
nulo e potencial nulo são condições independentes.}
\stepcounter{questao}
\resposta{(a) $V = \SI{1.5e5}{\volt}$; (b) $W = \SI{0.3}{\joule}$}
\resolucao{(a) $V = \ke\dfrac{Q}{d}
= 9\cdot10^9\cdot\dfrac{\pot{5}{-6}}{0{,}3}=\pot{1.5}{5}\,\si{\volt}$.\\
(b) O trabalho da força externa iguala a energia potencial adquirida (partindo do
infinito, onde $V=0$):
\[
W = q\,(V - V_\infty) = \pot{2}{-6}\cdot\pot{1.5}{5} = \SI{0.3}{\joule}.
\]
Como ambas as cargas são positivas, elas se repelem, e o trabalho externo é
\emph{positivo} (precisamos empurrar $q$ contra a repulsão). Se as cargas tivessem
sinais opostos, o trabalho externo seria negativo.}
\stepcounter{questao}
\resposta{$E_c = \SI{3.2e-17}{\joule}$; $v \approx \SI{8.4e6}{\meter\per\second}$}
\resolucao{Toda a energia potencial elétrica se converte em energia cinética:
\[
E_c = |q|\,U = \pot{1.6}{-19}\cdot200 = \pot{3.2}{-17}\,\si{\joule}.
\]
Da definição de energia cinética, $E_c = \tfrac12 m v^2$:
\[
v=\sqrt{\frac{2E_c}{m}}=\sqrt{\frac{2\cdot\pot{3.2}{-17}}{\pot{9.1}{-31}}}
\approx\SI{8.4e6}{\meter\per\second}.
\]
É o princípio de funcionamento dos antigos tubos de raios catódicos e dos
aceleradores de partículas: usar uma ddp para acelerar cargas.}
\stepcounter{questao}
\resposta{(a)}
\resolucao{O potencial varia com $1/d$: dobrando a distância (de \SI{2}{} para
\SI{4}{\meter}), ele cai à metade, $V(4) = \SI{20}{\volt}$.\\
Para o campo, usamos $E = \dfrac{V}{d}$ (válido para carga puntiforme):
$E(2) = \dfrac{40}{2} = \SI{20}{\volt\per\meter}$.}
\stepcounter{questao}
\resposta{(b)}
\resolucao{Para carga positiva, $V = \ke Q/d$ \emph{decresce} com a distância. Como
A é a superfície mais interna (menor $d$) e C a mais externa (maior $d$):
$V_A > V_B > V_C$. O sinal de $Q$ determina a tendência; o \emph{valor} de $Q$ não
é necessário para ordenar os potenciais.}
\stepcounter{questao}
\resposta{(a) $E=\SI{5.4e6}{\newton\per\coulomb}$, $V=\SI{540}{\kilo\volt}$;
(b) $E=\SI{6e5}{\newton\per\coulomb}$, $V=\SI{180}{\kilo\volt}$; (c) ver comentário}
\resolucao{(a) Fora da esfera (inclusive na superfície), ela se comporta como uma
carga puntiforme concentrada no centro:
\[
E(R)=\frac{\ke Q}{R^{2}}=\frac{9\cdot10^9\cdot\pot{6}{-6}}{(0{,}1)^2}
=\pot{5.4}{6}\,\si{\newton\per\coulomb},
\]
\[
V(R)=\frac{\ke Q}{R}=\frac{9\cdot10^9\cdot\pot{6}{-6}}{0{,}1}
=\pot{5.4}{5}\,\si{\volt}=\SI{540}{\kilo\volt}.
\]
(b) A $\SI{30}{\centi\meter}$ (três vezes o raio): $E$ cai com $1/d^2$ (fator 9) e
$V$ com $1/d$ (fator 3):
\[
E=\frac{5{,}4\cdot10^6}{9}=\pot{6}{5}\,\si{\newton\per\coulomb},
\qquad
V=\frac{540}{3}=\SI{180}{\kilo\volt}.
\]
(c) \textbf{A distinção conceitual central.} Em um condutor em equilíbrio
eletrostático, toda a carga se distribui na \emph{superfície} e o campo interno é
nulo --- se não fosse, as cargas livres se moveriam e não haveria equilíbrio.\\
Mas o potencial \emph{não} é nulo: ele é a \emph{integral} do campo desde o
infinito até o ponto. Ao entrar na esfera, o campo passa a ser zero, então o
potencial \textbf{para de variar} --- ele permanece constante e igual ao valor da
superfície:
\[
V_{\text{interno}} = V(R) = \SI{540}{\kilo\volt}.
\]
Isso é coerente com $E = -\dfrac{\mathrm{d}V}{\mathrm{d}d}$: campo nulo significa
potencial \emph{constante}, não necessariamente nulo. Todo o volume do condutor é
uma única região equipotencial.\\
\textbf{Aplicação:} é o princípio da \textbf{gaiola de Faraday} --- dentro de um
condutor fechado o campo é nulo, protegendo o conteúdo, ainda que o conjunto esteja
a um potencial elevadíssimo. É por isso que se está seguro dentro de um carro
atingido por um raio.}
\stepcounter{questao}
\resposta{(a) $0$, \SI{0.09}{\joule}, \SI{0.18}{\joule};
(b) $U = \SI{0.27}{\joule}$; (c) \SI{0.27}{\joule}}
\resolucao{\textbf{(a) Trabalho carga a carga.}
\begin{itemize}[leftmargin=1.6em,itemsep=1pt,topsep=2pt]
\item \textbf{1\textsuperscript{a} carga:} não há campo algum ainda, logo
      $W_1 = 0$.
\item \textbf{2\textsuperscript{a} carga:} vem contra o campo da primeira, a uma
      distância final $L$:
      \[
      W_2=\ke\frac{q^2}{L}=9\cdot10^9\frac{(\pot{1}{-6})^2}{0{,}1}=\SI{0.09}{\joule}.
      \]
\item \textbf{3\textsuperscript{a} carga:} sofre a ação das \emph{duas} anteriores,
      ambas à distância $L$:
      \[
      W_3=2\,\ke\frac{q^2}{L}=\SI{0.18}{\joule}.
      \]
\end{itemize}
\textbf{(b) Energia total da configuração.} É a soma dos trabalhos --- ou,
equivalentemente, a soma sobre os \emph{três pares} distintos:
\[
U = W_1+W_2+W_3 = 0+0{,}09+0{,}18 = \SI{0.27}{\joule}
= 3\,\ke\frac{q^2}{L}. \checkmark
\]
Note que a resposta \emph{não depende} da ordem em que as cargas foram trazidas
--- consequência de a força elétrica ser conservativa.

\textbf{(c)} Ao serem liberadas, toda a energia potencial armazenada se converte
em cinética (conservação de energia, sem dissipação):
\[
E_{c,\text{total}} = U = \SI{0.27}{\joule}.
\]
\textbf{Cuidado com o erro clássico:} a energia total \emph{não} é
$3\times W_3$ nem a soma de "cada carga vezes o potencial das outras" sem dividir
por 2 --- isso contaria cada par duas vezes. A fórmula geral é
$U=\tfrac12\sum_i q_i V_i$, onde $V_i$ é o potencial criado \emph{pelas demais}
na posição de $i$.}
\stepcounter{questao}
\resposta{(a) $Q_1=Q/4$, $Q_2=3Q/4$; (b) $E_1/E_2=3$; (c) ver comentário}
\resolucao{(a) Ligadas por um condutor, as esferas atingem o \emph{mesmo
potencial}:
\[
\frac{\ke Q_1}{R}=\frac{\ke Q_2}{3R}
\;\Longrightarrow\;
Q_2 = 3Q_1.
\]
Com $Q_1+Q_2=Q$: $Q_1=\dfrac{Q}{4}$ e $Q_2=\dfrac{3Q}{4}$.\\
(b) Campos nas superfícies:
\[
E_1=\frac{\ke Q_1}{R^{2}}=\frac{\ke Q}{4R^{2}},
\qquad
E_2=\frac{\ke Q_2}{(3R)^{2}}=\frac{3\ke Q}{4\cdot9R^{2}}=\frac{\ke Q}{12R^{2}}.
\]
\[
\frac{E_1}{E_2}=3.
\]
A esfera \textbf{menor}, embora tenha \emph{menos} carga, apresenta campo
\textbf{três vezes mais intenso} em sua superfície.\\
\textbf{Relação geral:} como $V=\ke Q/R$ e $E=\ke Q/R^{2}$, tem-se
$E = V/R$. Com $V$ comum às duas, \emph{o campo é inversamente proporcional ao
raio}.\\
(c) \textbf{Poder das pontas.} Uma "ponta" é uma região de raio de curvatura
muito pequeno. Sendo $E \propto 1/R$, o campo ali fica extremamente intenso ---
suficiente para ionizar o ar e provocar descarga, mesmo com o objeto a um
potencial moderado.\\
\textbf{Consequências práticas:}
\begin{itemize}[leftmargin=1.6em,itemsep=1pt,topsep=2pt]
\item \textbf{Para-raios} são pontiagudos \emph{de propósito}: concentram o campo
      e induzem a descarga em local controlado.
\item Equipamentos de alta tensão têm cantos \emph{arredondados} e cúpulas
      polidas, justamente para \emph{evitar} campos intensos e o efeito corona.
\end{itemize}}
\stepcounter{questao}
\resposta{(a) $U\approx\SI{0.49}{\joule}$; (b) $W\approx\SI{0.51}{\joule}$;
(c) $\approx\SI{0.49}{\joule}$}
\resolucao{(a) Há $\binom{4}{2}=6$ pares: \textbf{4 lados} (distância $a$) e
\textbf{2 diagonais} (distância $a\sqrt2$):
\[
U=4\cdot\frac{\ke q^{2}}{a}+2\cdot\frac{\ke q^{2}}{a\sqrt2}
=\frac{\ke q^{2}}{a}\left(4+\frac{2}{\sqrt2}\right)
=\frac{\ke q^{2}}{a}\left(4+\sqrt2\right).
\]
Com $\dfrac{\ke q^{2}}{a}=\dfrac{9\cdot10^{9}(\pot{1}{-6})^{2}}{0{,}1}
=\SI{0.09}{\joule}$:
\[
U=0{,}09\,(4+1{,}41)\approx\SI{0.49}{\joule}.
\]
(b) A quinta carga, no centro, dista $\dfrac{a\sqrt2}{2}=\SI{0.0707}{\meter}$ de
\emph{cada} uma das quatro. O potencial no centro é
\[
V_c=4\cdot\frac{\ke q}{a\sqrt2/2}
=4\cdot\frac{9\cdot10^{9}\cdot\pot{1}{-6}}{0{,}0707}
\approx\pot{5.09}{5}\,\si{\volt},
\]
e o trabalho externo é
\[
W=q\,V_c=\pot{1}{-6}\cdot\pot{5.09}{5}\approx\SI{0.51}{\joule}.
\]
(c) Sem a quinta carga, ao liberar as quatro, toda a energia potencial da
configuração converte-se em cinética:
\[
E_{c}=U\approx\SI{0.49}{\joule}.
\]
\textbf{Atenção ao método.} Em (a) somam-se os \emph{pares} (cada par uma vez
só); em (b) usa-se $W=qV$ porque as quatro cargas já estavam posicionadas e
imóveis. Confundir os dois procedimentos --- contando pares duas vezes --- é o
erro mais comum neste tipo de problema.}
\stepcounter{questao}
\resposta{(a) $-Q$ na interna, $+Q$ na externa; (b) e (c) ver comentário}
\resolucao{(a) Dentro do \emph{material} do condutor ($a<r<b$) o campo deve ser
\textbf{nulo} em equilíbrio. Aplicando a lei de Gauss a uma superfície esférica
nessa região, a carga interna total deve ser zero:
\[
Q + q_{\text{int}} = 0 \;\Longrightarrow\; q_{\text{int}}=-Q.
\]
Como a casca é neutra no total, a superfície externa fica com
$q_{\text{ext}}=+Q$.\\
(b) \textbf{Comportamento do campo:}
\begin{center}\small\sffamily
\begin{tabular}{@{}l l l@{}}
\toprule
\textbf{Região} & \textbf{Carga envolvida} & \textbf{Campo}\\
\midrule
$r<a$      & $+Q$        & $E=\ke Q/r^{2}$ (radial, para fora)\\
$a<r<b$    & $Q-Q=0$     & $E=0$ (dentro do condutor)\\
$r>b$      & $Q-Q+Q=+Q$  & $E=\ke Q/r^{2}$\\
\bottomrule
\end{tabular}
\end{center}
Fora da casca, o campo é \emph{idêntico} ao que existiria sem ela --- a casca
neutra não blinda o exterior.\\
(c) \textbf{Aterrando a casca}, a superfície externa perde sua carga $+Q$ para a
terra (os elétrons sobem para neutralizá-la). Resultado:
\[
q_{\text{int}}=-Q,\qquad q_{\text{ext}}=0
\;\Longrightarrow\;
E = 0 \ \text{ para } r>b.
\]
\textbf{Agora sim há blindagem completa}: nenhum campo escapa para fora.\\
\textbf{Aplicação --- blindagem eletrostática.} É por isso que cabos coaxiais e
gabinetes de equipamentos sensíveis têm sua malha metálica \textbf{aterrada}:
uma blindagem flutuante (não aterrada) não impede o campo de vazar. O aterramento
é o que transforma o condutor em uma verdadeira gaiola de Faraday.}
\stepcounter{questao}
\resposta{$V=\SI{60}{\kilo\volt}$}
\resolucao{
\[
V=\frac{\ke Q}{r}=\frac{\pot{9}{9}\times\pot{2}{-6}}{0{,}30}
=\frac{\pot{1.8}{4}}{0{,}30}=\pot{6}{4}\,\si{\volt}=\SI{60}{\kilo\volt}.
\]
\textbf{Atenção à potência de $r$:} o potencial cai com $1/r$, enquanto o campo
cai com $1/r^{2}$. Confundir as duas é o erro mais comum do tópico --- e o
resultado sai com a dimensão errada.}
\stepcounter{questao}
\resposta{$V=\SI{-60}{\kilo\volt}$}
\resolucao{
\[
V=\frac{\pot{9}{9}\times(-\pot{4}{-6})}{0{,}60}=-\pot{6}{4}\,\si{\volt}
=\SI{-60}{\kilo\volt}.
\]
\textbf{O sinal é do potencial, não do módulo.} Diferentemente do campo --- cuja
\emph{orientação} é dada por um vetor --- o potencial é um escalar com sinal
próprio. Cargas negativas produzem potencial negativo em toda parte.}
\stepcounter{questao}
\resposta{$r=\SI{0.6}{\meter}$}
\resolucao{
\[
r=\frac{\ke Q}{V}=\frac{\pot{9}{9}\times\pot{3}{-6}}{\pot{4.5}{4}}
=\frac{\pot{2.7}{4}}{\pot{4.5}{4}}=\SI{0.6}{\meter}.
\]
\textbf{Superfície equipotencial:} \emph{todos} os pontos a
$\SI{0.6}{\meter}$ da carga estão a $\SI{45}{\kilo\volt}$. Em torno de uma
carga puntiforme, as equipotenciais são esferas concêntricas.}
\stepcounter{questao}
\resposta{$Q=-\SI{1}{\micro\coulomb}$}
\resolucao{
\[
Q=\frac{Vr}{\ke}=\frac{(-\pot{1.8}{4})(0{,}5)}{\pot{9}{9}}
=-\pot{1}{-6}\,\si{\coulomb}=-\SI{1}{\micro\coulomb}.
\]
O sinal negativo do potencial revelou diretamente o sinal da carga --- vantagem
prática de trabalhar com uma grandeza escalar.}
\stepcounter{questao}
\resposta{(b)}
\resolucao{O potencial é \textbf{escalar}: basta somar as contribuições
$\ke q_k/r_k$ com seus sinais, sem ângulos nem componentes.\\
\textbf{Vantagem prática enorme:} calcular $V$ de dez cargas exige dez divisões
e uma soma; calcular $\vec E$ exigiria dez vetores decompostos em componentes.\\
A alternativa (d) descreve a propriedade oposta --- justamente por o campo
eletrostático ser \emph{conservativo}, o potencial de um ponto independe do
caminho, e é isso que permite defini-lo.}
\stepcounter{questao}
\resposta{$U=\SI{0.6}{\milli\joule}$}
\resolucao{
\[
U=qV=\pot{3}{-6}\times200=\pot{6}{-4}\,\si{\joule}=\SI{0.6}{\milli\joule}.
\]
\textbf{Distinção essencial:} $V$ é uma propriedade \emph{do ponto} --- existe
mesmo sem carga alguma ali. Já $U$ é uma propriedade \emph{do sistema}
carga-campo, e só faz sentido quando a carga está presente.\\
A relação $U=qV$ é o análogo elétrico de $E_p=mgh$: potencial faz o papel de
$gh$, e carga faz o papel de massa.}
\stepcounter{questao}
\resposta{$W=\SI{60}{\micro\joule}$; $\Delta U=\SI{-60}{\micro\joule}$}
\resolucao{(a)
\[
W=q\left(V_A-V_B\right)=\pot{2}{-6}(50-20)=\pot{6}{-5}\,\si{\joule}
=\SI{60}{\micro\joule}.
\]
(b) $\Delta U=-W=\SI{-60}{\micro\joule}$.\\
\textbf{Coerência dos sinais:} a carga é positiva e foi de maior para menor
potencial --- movimento espontâneo. A força elétrica realiza trabalho
\emph{positivo} e a energia potencial \emph{diminui}, convertendo-se em energia
cinética.\\
\textbf{Regra:} cargas positivas "descem" o potencial espontaneamente; cargas
negativas "sobem".}
\stepcounter{questao}
\resposta{$V=0$}
\resolucao{Somando algebricamente:
\[
V=\frac{\pot{9}{9}(\pot{3}{-6})}{0{,}30}
+\frac{\pot{9}{9}(-\pot{2}{-6})}{0{,}20}
=\pot{9}{4}-\pot{9}{4}=\SI{0}{\volt}.
\]
\textbf{Cuidado:} $V=0$ \textbf{não} significa campo nulo nem ausência de força.
As duas cargas produzem campos que \emph{não} se cancelam --- elas têm sinais
opostos, e seus campos apontam para o mesmo lado ao longo da linha que as une.\\
Uma carga colocada em $P$ teria energia potencial nula, mas sofreria força.
Essa distinção é aprofundada no bloco seguinte.}
\stepcounter{questao}
\resposta{$U=\SI{20}{\volt}$}
\resolucao{
\[
U=E\,d=500\times0{,}04=\SI{20}{\volt},
\]
com o ponto \emph{a montante} do campo em potencial maior --- o campo aponta
sempre no sentido de $V$ decrescente.\\
\textbf{Se o deslocamento fosse perpendicular} ao campo, a diferença seria
\emph{nula}: os dois pontos estariam na mesma equipotencial. Só a componente do
deslocamento \emph{ao longo} do campo contribui.\\
\textbf{Unidades:} $\si{\volt\per\meter}$ e $\si{\newton\per\coulomb}$ são a
mesma coisa --- verifique multiplicando por metro e por coulomb.}
\stepcounter{questao}
\resposta{$\SI{500}{\electronvolt}=\pot{8}{-17}\,\si{\joule}$;
$v=\pot{1.33}{7}\,\si{\meter\per\second}$}
\resolucao{(a)
\[
W=eU=\pot{1.6}{-19}(500)=\pot{8}{-17}\,\si{\joule}.
\]
Em elétron-volts, a leitura é imediata: $\SI{500}{\electronvolt}$ --- \emph{por
definição}, $\SI{1}{\electronvolt}$ é a energia que a carga elementar adquire
sob $\SI{1}{\volt}$.\\
(b) Toda a energia vira cinética:
\[
v=\sqrt{\frac{2W}{m}}=\sqrt{\frac{2(\pot{8}{-17})}{\pot{9.11}{-31}}}
=\sqrt{\pot{1.756}{14}}=\pot{1.33}{7}\,\si{\meter\per\second}.
\]
\textbf{Verificação de consistência:} isso é $4{,}4\%$ da velocidade da luz ---
ainda no regime não relativístico, onde $\frac12mv^{2}$ é válido. Acima de cerca
de $\SI{50}{\kilo\electronvolt}$, a correção relativística passa a ser
necessária.}
\stepcounter{questao}
\resposta{$x=\SI{0.4}{\meter}$}
\resolucao{Impondo $V=0$ com $x$ medido a partir de $q_1$:
\[
\frac{\ke(4\mu)}{x}+\frac{\ke(-1\mu)}{0{,}5-x}=0
\;\Longrightarrow\;
\frac{4}{x}=\frac{1}{0{,}5-x}.
\]
\[
4(0{,}5-x)=x
\;\Longrightarrow\;
2=5x
\;\Longrightarrow\;
x=\SI{0.4}{\meter}.
\]
\textbf{Contraste importante com o campo:} o ponto de \emph{campo} nulo dessas
mesmas cargas fica \emph{fora} do segmento (a $\SI{1}{\meter}$ de $q_1$),
porque o campo é vetorial e só pode se cancelar onde os dois vetores forem
opostos. Já o potencial é escalar e se anula sempre que as contribuições tiverem
módulos iguais e sinais opostos --- o que ocorre dentro do segmento.\\
\textbf{Existe também um segundo ponto de $V=0$} fora do segmento, do lado da
carga menor, em $x=\SI{0.667}{\meter}$. A superfície completa de $V=0$ é uma
esfera.}
\stepcounter{questao}
\resposta{$V=0$; $E=\SI{900}{\kilo\volt\per\meter}$}
\resolucao{(a) As duas contribuições têm módulos iguais e sinais opostos:
\[
V=\frac{\ke Q}{0{,}20}-\frac{\ke Q}{0{,}20}=0.
\]
(b) Os \emph{campos}, ao contrário, apontam ambos da carga positiva para a
negativa --- eles se \textbf{somam}:
\[
E=2\times\frac{\ke Q}{(0{,}20)^{2}}
=2\times\frac{\pot{9}{9}(\pot{2}{-6})}{0{,}04}
=\pot{9}{5}\,\si{\volt\per\meter}.
\]
(c) \textbf{Não há contradição.} O potencial não determina o campo --- quem o
determina é o \emph{gradiente} do potencial:
\[
\vec E=-\nabla V.
\]
Uma função pode valer zero em um ponto e ter derivada grande ali. É exatamente o
caso: caminhando ao longo do eixo, $V$ passa de positivo a negativo cruzando o
zero com inclinação acentuada.\\
\textbf{Analogia:} a altitude do nível do mar é zero, mas isso nada diz sobre a
inclinação do terreno naquele ponto. É a inclinação --- e não a altitude --- que
faz uma bola rolar.\\
\textbf{Consequência:} uma carga solta no ponto médio \emph{não} fica parada,
apesar de $V=0$ e $U=0$.}
\stepcounter{questao}
\resposta{$E=\SI{50}{\volt\per\meter}$ no sentido $+x$;
$W=\SI{200}{\micro\joule}$}
\resolucao{(a)
\[
E_x=-\frac{\mathrm{d}V}{\mathrm{d}x}=-(-50)=\SI{+50}{\volt\per\meter}.
\]
Campo uniforme, apontando no sentido \emph{crescente} de $x$ --- ou seja, no
sentido em que $V$ \emph{decresce}, como sempre.\\
(b) $V(1)=\SI{150}{\volt}$ e $V(3)=\SI{50}{\volt}$:
\[
W=q\left(V_1-V_3\right)=\pot{2}{-6}(150-50)=\SI{200}{\micro\joule}.
\]
\emph{(Conferindo por força: $W=qE\,d=\pot{2}{-6}(50)(2)
=\SI{200}{\micro\joule}$ \checkmark.)}\\
(c) Como $V$ depende \emph{apenas} de $x$, as equipotenciais são
\textbf{planos perpendiculares ao eixo} $x$. O plano $x=\SI{1}{\meter}$ é toda
a superfície com $V=\SI{150}{\volt}$.\\
\textbf{Regra geral:} equipotenciais são sempre \textbf{perpendiculares} às
linhas de campo. Se não fossem, haveria componente de campo ao longo da
superfície, e mover uma carga sobre ela exigiria trabalho --- contradizendo a
definição de equipotencial.}
\stepcounter{questao}
\resposta{$W_{ext}=U=\SI{0.225}{\joule}$}
\resolucao{(a) e (b) Como a carga parte do repouso no infinito e chega em
repouso, todo o trabalho externo vira energia potencial:
\[
U=\frac{\ke Qq}{r}=\frac{\pot{9}{9}(\pot{5}{-6})(\pot{1}{-6})}{0{,}20}
=\frac{\pot{4.5}{-2}}{0{,}20}=\SI{0.225}{\joule}.
\]
\emph{(Equivalentemente, $U=qV$ com $V=\ke Q/r=\SI{225}{\kilo\volt}$.)}\\
(c) \textbf{Sinal positivo:} as cargas se repelem, então aproximá-las exige
trabalho \emph{contra} a força elétrica. O agente externo gasta energia, que
fica armazenada no sistema.\\
\textbf{Se as cargas tivessem sinais opostos}, $U$ seria negativa: elas se
atraem, e seria preciso \emph{segurá-las} durante a aproximação --- o agente
externo receberia energia.\\
\textbf{Verificação de coerência:} soltando a carga, ela é repelida e ganha
$\SI{0.225}{\joule}$ de energia cinética ao voltar ao infinito. A energia
potencial é integralmente recuperável --- marca de um campo conservativo.}
\stepcounter{questao}
\resposta{$V$ nunca se anula (exceto no infinito); $\vec E$ se anula no ponto
médio}
\resolucao{(a) Ambas as cargas são positivas, logo cada uma contribui com
potencial \textbf{positivo} em todo ponto finito. A soma de dois positivos nunca
é zero:
\[
V=\frac{\ke Q}{r_1}+\frac{\ke Q}{r_2}>0\quad\text{sempre}.
\]
$V\to0$ apenas no limite $r\to\infty$.\\
(b) O campo é vetorial e \emph{pode} se cancelar. No ponto médio do segmento,
os dois vetores têm módulos iguais e sentidos opostos:
\[
\vec E=\vec 0\ \text{no ponto médio}.
\]
(c) \textbf{A diferença é a natureza das grandezas.}
\begin{itemize}
\item $V$ é escalar e, para cargas de mesmo sinal, as contribuições não podem se
cancelar --- só se somam.
\item $\vec E$ é vetorial e se cancela sempre que os vetores forem opostos, o que
não impõe restrição alguma sobre os sinais das cargas.
\end{itemize}
\textbf{Contraste com o caso de cargas opostas} ($+Q$ e $-Q$): ali ocorre
exatamente o inverso --- $V=0$ no ponto médio, mas $\vec E\neq\vec 0$.\\
\textbf{Síntese:} $V=0$ e $\vec E=\vec 0$ são condições \emph{independentes}.
Nenhuma implica a outra, em nenhum sentido.}
\stepcounter{questao}
\resposta{$W_{AB}=q(V_A-V_B)$, com
$V=\ke Q/r$}
\resolucao{(a) Para uma carga fixa $Q$ e uma carga de prova $q$, a força é
radial: $\vec F=\dfrac{\ke Qq}{r^{2}}\hat r$. Em um deslocamento
$\mathrm{d}\vec\ell$, apenas a componente radial contribui, pois
$\vec F\cdot\mathrm{d}\vec\ell=F\,\mathrm{d}r$:
\[
W_{AB}=\int_{r_A}^{r_B}\frac{\ke Qq}{r^{2}}\,\mathrm{d}r
=\ke Qq\left[-\frac{1}{r}\right]_{r_A}^{r_B}
=\ke Qq\left(\frac{1}{r_A}-\frac{1}{r_B}\right).
\]
\textbf{O resultado depende apenas de $r_A$ e $r_B$} --- as componentes
transversais do caminho não contribuem, por serem perpendiculares à força.\\
(b) \textbf{Generalização:} para $n$ cargas, a força total é a soma vetorial das
forças individuais, e o trabalho total é a soma dos trabalhos. Como cada parcela
independe do caminho, a soma também independe.\\
Definindo $V=\sum\ke Q_k/r_k$, obtém-se
\[
\boxed{W_{AB}=q\left(V_A-V_B\right)}
\]
(c) \textbf{Por que isso permite definir $V$.} Se o trabalho dependesse do
caminho, não haveria como atribuir um \emph{número} a cada ponto --- o "potencial"
seria ambíguo. A independência do caminho garante que
$V(P)=W_{\infty\to P}/q$ é bem definido.\\
\textbf{Formulação equivalente:} $\oint\vec E\cdot\mathrm{d}\vec\ell=0$ em
qualquer curva fechada. O campo eletrostático é \textbf{conservativo}, e essa é
exatamente a lei das malhas de Kirchhoff em sua forma de campo.\\
\textbf{Ressalva importante:} isso vale para campos \emph{eletrostáticos}.
Campos elétricos induzidos por variação de fluxo magnético \emph{não} são
conservativos, e para eles o potencial escalar deixa de existir --- assunto da
lei de Faraday.}
\stepcounter{questao}
\resposta{Condição: $|q_1|/r_1=|q_2|/r_2$; com $+4\,\si{\micro\coulomb}$ e
$-1\,\si{\micro\coulomb}$ a $\SI{0.5}{\meter}$, o ponto é
$x=\SI{0.4}{\meter}$}
\resolucao{(a) \textbf{Condição para $V=0$:}
\[
\frac{\ke q_1}{r_1}+\frac{\ke q_2}{r_2}=0
\;\Longrightarrow\;
\frac{|q_1|}{r_1}=\frac{|q_2|}{r_2},
\]
com $q_1$ e $q_2$ de sinais \emph{opostos}.\\
\textbf{Condição para $\vec E=\vec0$ no mesmo ponto:}
$\dfrac{|q_1|}{r_1^{2}}=\dfrac{|q_2|}{r_2^{2}}$.\\
As duas só coincidem se $r_1=r_2$ --- e então $|q_1|=|q_2|$, que é o caso
excluído pelo enunciado. \textbf{Com módulos diferentes, é impossível anular
ambos no mesmo ponto.}\\
(b) Tome $q_1=+\SI{4}{\micro\coulomb}$ em $x=0$ e
$q_2=-\SI{1}{\micro\coulomb}$ em $x=\SI{0.5}{\meter}$. De
$\dfrac{4}{x}=\dfrac{1}{0{,}5-x}$ vem $x=\SI{0.4}{\meter}$.\\
\emph{(Confira: $r_1=0{,}4$ e $r_2=0{,}1$, com
$4/0{,}4=1/0{,}1=10$ \checkmark.)}\\
(c) \textbf{Campo no ponto.} Ambos apontam no sentido $+x$ (afastando-se de
$q_1$, aproximando-se de $q_2$) e portanto se somam:
\[
E=\frac{\pot{9}{9}(\pot{4}{-6})}{0{,}16}
+\frac{\pot{9}{9}(\pot{1}{-6})}{0{,}01}
=\pot{2.25}{5}+\pot{9}{5}=\pot{1.125}{6}\,\si{\volt\per\meter}.
\]
Mais de um megavolt por metro --- longe de nulo.\\
\textbf{Interpretação:} uma carga colocada ali teria energia potencial nula, mas
sofreria força intensa. \textbf{$V$ e $\vec E$ carregam informações
diferentes:} $V$ é o valor, $\vec E$ é a taxa de variação.}
\stepcounter{questao}
\resposta{Positiva: \textbf{instável}; negativa: estável em $x$ --- mas nenhuma
é estável em três dimensões}
\resolucao{(a) A energia potencial de uma carga positiva é $U=qV$, com
$q>0$ --- ela tem a \emph{mesma} forma de $V$ e, portanto, também um mínimo em
$x_0$. Mínimo de energia significa \textbf{equilíbrio estável}... \emph{ao longo
de $x$}.\\
\emph{(Atenção ao raciocínio: é o mínimo de $U$, e não o de $V$, que define
estabilidade. Para carga positiva os dois coincidem; para negativa, invertem-se.)}\\
(b) Para carga negativa, $U=qV$ com $q<0$ \textbf{inverte} a curva: onde $V$
tem mínimo, $U$ tem \emph{máximo}. Equilíbrio \textbf{instável} em $x$.\\
(c) \textbf{Teorema de Earnshaw.} Em uma região sem cargas, o potencial
eletrostático satisfaz a equação de Laplace, $\nabla^{2}V=0$, ou seja
\[
\frac{\partial^{2}V}{\partial x^{2}}
+\frac{\partial^{2}V}{\partial y^{2}}
+\frac{\partial^{2}V}{\partial z^{2}}=0.
\]
Se $V$ tivesse um mínimo \emph{tridimensional} em $x_0$, as três segundas
derivadas seriam positivas e a soma não poderia ser zero. \textbf{Contradição.}\\
Logo, um mínimo em uma direção obriga a haver \emph{máximo} em outra --- todo
ponto crítico é um \textbf{ponto de sela}.\\
\textbf{Conclusão:} nenhuma carga pode ser mantida em equilíbrio estável apenas
por forças eletrostáticas, seja ela positiva ou negativa. A análise
unidimensional do item (a) é enganosa justamente por ignorar as outras duas
direções.\\
\textbf{Como se contorna na prática:} campos variáveis no tempo (armadilha de
Paul), campos magnéticos, ou realimentação ativa --- todas soluções que escapam
da hipótese eletrostática.}
\stepcounter{questao}
\resposta{Campo interno nulo $\Rightarrow$ $V$ constante em todo o condutor; o
campo externo é \textbf{perpendicular} à superfície}
\resolucao{(a) \textbf{Condição de equilíbrio:} não há movimento de cargas
livres, logo o campo dentro do metal é \textbf{nulo}. Se houvesse campo, as
cargas se moveriam e o equilíbrio não existiria.\\
Tomando dois pontos $A$ e $B$ quaisquer do condutor e um caminho
\emph{interno} entre eles:
\[
V_A-V_B=\int_A^B\vec E\cdot\mathrm{d}\vec\ell=\int_A^B\vec 0\cdot\mathrm{d}\vec\ell=0.
\]
Portanto $V_A=V_B$ --- \textbf{todo o condutor é equipotencial}.\\
(b) O argumento não distinguiu superfície de interior: ele vale para
\emph{qualquer} par de pontos ligados por caminho dentro do metal. Superfície e
interior estão no mesmo potencial, e o condutor inteiro é um volume
equipotencial.\\
(c) \textbf{Orientação do campo externo.} Se houvesse componente
\emph{tangencial} à superfície, ela empurraria as cargas livres ao longo dela --
contradizendo o equilíbrio. Logo o campo externo é \textbf{perpendicular} à
superfície em todo ponto.\\
\textbf{Consequências, todas verificáveis:}
\begin{itemize}
\item A carga em excesso reside \emph{integralmente} na superfície externa.
\item A densidade superficial é maior onde a curvatura é maior (poder das
pontas).
\item Uma cavidade vazia dentro do condutor tem campo nulo --- é a blindagem
eletrostática.
\item Ligar dois condutores por fio iguala seus potenciais, e a carga se
redistribui proporcionalmente às capacitâncias.
\end{itemize}}
\stepcounter{questao}
\resposta{$\vec E=(20;\,-10y)\,\si{\volt\per\meter}$; em $(2;1)$,
$\vec E=(20;\,-10)$ com $|E|=\SI{22.4}{\volt\per\meter}$}
\resolucao{(a) $\vec E=-\nabla V$, componente a componente:
\[
E_x=-\frac{\partial V}{\partial x}=-(-20)=\SI{+20}{\volt\per\meter};
\qquad
E_y=-\frac{\partial V}{\partial y}=-(10y)=-10y.
\]
Note que $E_x$ é \emph{uniforme}, enquanto $E_y$ cresce linearmente com $y$.\\
(b) Em $(2;\,1)$:
\[
\vec E=(20;\,-10)\,\si{\volt\per\meter};
\qquad
|E|=\sqrt{400+100}=\SI{22.4}{\volt\per\meter},
\]
formando $\ang{-26.6}$ com o eixo $x$.\\
\textbf{Observe:} o valor de $V$ no ponto ($V=100-40+5=\SI{65}{\volt}$) não
entrou no cálculo do campo. Só as \emph{derivadas} importam --- somar uma
constante a $V$ não altera $\vec E$.\\
(c) \textbf{Equipotenciais:} $100-20x+5y^{2}=\text{const}$, ou seja
\[
x=\frac{5y^{2}+100-C}{20},
\]
uma família de \textbf{parábolas} com eixo horizontal.\\
\textbf{Verificação de coerência:} o vetor $\vec E$ calculado deve ser
perpendicular à parábola que passa por $(2;1)$. A tangente à curva ali tem
inclinação $\mathrm{d}y/\mathrm{d}x=2$, e o campo tem inclinação
$-10/20=-1/2$ --- produto $-1$, portanto perpendiculares \checkmark.}
\stepcounter{questao}
\resposta{Perto: esferas em torno de cada carga; longe: esferas em torno do par;
o ponto médio é ponto de sela do potencial}
\resolucao{(a) \textbf{Equipotenciais.}
\begin{itemize}
\item \emph{Muito perto de uma carga:} a contribuição dela domina, e as
superfícies são praticamente \textbf{esferas} centradas nessa carga.
\item \emph{Muito longe do par:} o sistema parece uma única carga $2Q$, e as
superfícies voltam a ser esferas, agora centradas no \textbf{centro do par}.
\item \emph{Na região intermediária:} as superfícies deixam de ser esferas e
assumem forma de "amendoim", chegando a se fundir em uma única superfície que
envolve as duas cargas.
\end{itemize}
(b) \textbf{Linhas de campo.} Saem radialmente de cada carga, curvam-se
afastando-se uma da outra e seguem para o infinito. \textbf{No ponto médio o
campo é nulo} --- as linhas não passam por ali. Não há linha ligando as duas
cargas, pois ambas são fontes.\\
(c) \textbf{Regras gerais.}
\begin{enumerate}
\item As linhas de campo são \textbf{perpendiculares} às equipotenciais em todo
ponto.
\item O campo aponta no sentido de $V$ \textbf{decrescente}.
\item Onde as equipotenciais estão \textbf{mais próximas}, o campo é mais
\textbf{intenso} --- pois $|E|=|\mathrm{d}V/\mathrm{d}n|$.
\item Equipotenciais nunca se cruzam (um ponto teria dois potenciais); linhas de
campo também não (haveria duas direções de força).
\item Onde $\vec E=\vec0$, as duas regras de perpendicularidade perdem sentido:
é um \textbf{ponto de sela} do potencial, e ali as equipotenciais podem se
cruzar.
\end{enumerate}
\textbf{O ponto médio é justamente esse caso excepcional:} $V$ tem mínimo ao
longo da linha que une as cargas e máximo na direção perpendicular --- exatamente
a estrutura de sela exigida pelo teorema de Earnshaw.}

\gabarito[1.4 Potencial elétrico e superfícies equipotenciais]

\section{Conceitos Iniciais: Leis de Ohm, Potência, Fontes e Kirchhoff}

\subsection{Corrente elétrica e carga}
\stepcounter{questao}
\resposta{(c)}
\resolucao{Convertendo o tempo para segundos, $t = 2\cdot 60 = \SI{120}{\second}$:
\[
Q = I\,t = 3\cdot 120 = \SI{360}{\coulomb}.
\]
O erro típico (alternativa e) é usar $t = \SI{2}{\minute}$ diretamente, sem
converter para segundos.}
\stepcounter{questao}
\resposta{(c)}
\resolucao{$I = \dfrac{Q}{t} = \dfrac{120}{60} = \SI{2}{\ampere}$.}
\stepcounter{questao}
\resposta{(e)}
\resolucao{$Q = I\,t = 2\cdot(5\cdot60) = 2\cdot 300 = \SI{600}{\coulomb}$.}
\stepcounter{questao}
\resposta{(a) \SI{32}{\milli\coulomb}; (b) $n = \pot{2.0}{17}$ elétrons}
\resolucao{(a) Em $t = \SI{1}{\second}$:
$Q = I\,t = \pot{32}{-3}\cdot 1 = \SI{32}{\milli\coulomb}$.\\
(b) $n = \dfrac{Q}{e} = \dfrac{\pot{32}{-3}}{\pot{1.6}{-19}} = \pot{2.0}{17}$
elétrons por segundo.}
\stepcounter{questao}
\resposta{$Q = \SI{20}{\coulomb}$}
\resolucao{A carga é a \emph{área} sob o gráfico $i \times t$. Como a corrente é
constante e igual a $\SI{4}{\ampere}$ durante $\SI{5}{\second}$, a área é um
retângulo:
\[
Q = i\,t = 4\cdot 5 = \SI{20}{\coulomb}.
\]
Essa interpretação por área vale mesmo quando a corrente varia --- basta calcular
a área da figura correspondente.}
\stepcounter{questao}
\resposta{$Q = \SI{16}{\coulomb}$}
\resolucao{A carga é a \emph{área} sob o gráfico $i \times t$. A figura é composta
de um retângulo (de $0$ a \SI{2}{\second}) e um triângulo (de $2$ a \SI{6}{\second}):
\[
Q = \underbrace{4\cdot2}_{\text{retângulo}}
  + \underbrace{\tfrac12\cdot4\cdot4}_{\text{triângulo}}
  = 8 + 8 = \SI{16}{\coulomb}.
\]
Mesmo com corrente variável, a área continua valendo como carga total.}
\stepcounter{questao}
\resposta{(d)}
\resolucao{Convertendo o tempo: $\SI{10}{\hour} = 10\cdot3600 = \SI{36000}{\second}$.
\[
Q = I\,t = 5\cdot36000 = \SI{180000}{\coulomb}.
\]
O erro comum é esquecer de converter horas em segundos.}
\stepcounter{questao}
\resposta{(a) $Q = \SI{32}{\coulomb}$; (b) $i_m = \SI{4}{\ampere}$;
(c) $t = \SI{4}{\second}$}
\resolucao{(a) A carga é a área sob o gráfico --- aqui, um trapézio de bases
\SI{6}{} e \SI{2}{\ampere} e altura \SI{8}{\second}:
\[
Q=\frac{(6+2)}{2}\cdot 8 = \SI{32}{\coulomb}.
\]
(b) $i_m = \dfrac{Q}{\Delta t} = \dfrac{32}{8} = \SI{4}{\ampere}$, que coincide
com a média aritmética $\dfrac{6+2}{2}$ --- válido \emph{porque} a variação é
linear.\\
(c) Como a reta decresce uniformemente, ela cruza o valor médio exatamente no
\emph{meio} do intervalo: $t = \SI{4}{\second}$.
Conferindo pela equação da reta $i(t) = 6 - 0{,}5t$:
$i(4) = 6-2 = \SI{4}{\ampere}$. \checkmark\\
\textbf{Cuidado:} se a corrente variasse de forma \emph{não} linear (por exemplo,
exponencialmente), a corrente média \textbf{não} seria a média dos extremos ---
seria preciso calcular a área real sob a curva.}
\stepcounter{questao}
\resposta{(a) $v_d \approx \SI{0.74}{\milli\meter\per\second}$;
(b) $\approx\SI{1360}{\second}$ ($\approx\SI{23}{\minute}$); (c) ver comentário}
\resolucao{\textbf{(a)} A corrente relaciona-se com a velocidade de deriva por
$I = n\,e\,A\,v_d$. Isolando:
\[
v_d=\frac{I}{n\,e\,A}
=\frac{10}{(\pot{8.5}{28})(\pot{1.6}{-19})(\pot{1}{-6})}
\approx\pot{7.4}{-4}\,\si{\meter\per\second}=\SI{0.74}{\milli\meter\per\second}.
\]
\textbf{(b)} $t = \dfrac{1}{\pot{7.4}{-4}} \approx \SI{1360}{\second}
\approx \SI{23}{\minute}$ --- um elétron leva aproximadamente 23 minutos para percorrer um metro.

\textbf{(c) A resolução do aparente paradoxo.} A lâmpada acende
instantaneamente porque \textbf{não é preciso que um elétron viaje da chave até a
lâmpada}. O que se propaga rapidamente é o \emph{campo elétrico} ao longo do
condutor, a uma velocidade próxima à da luz ($\sim\SI{2e8}{\meter\per\second}$
em cabos reais). Esse campo põe em movimento, quase simultaneamente,
\emph{todos} os elétrons livres ao longo de todo o fio --- inclusive os que já
estavam dentro do filamento da lâmpada.\\
\textbf{Analogia:} é como um tubo já cheio de água. Ao empurrar na entrada, sai
água na saída imediatamente, embora cada molécula individual se desloque
lentamente. A informação (pressão) viaja rápido; a matéria, devagar.\\
\textbf{Ordens de grandeza envolvidas:} deriva $\sim\SI{e-4}{\meter\per\second}$;
agitação térmica dos elétrons $\sim\SI{e6}{\meter\per\second}$; sinal elétrico
$\sim\SI{e8}{\meter\per\second}$ --- doze ordens de grandeza separando a mais
lenta da mais rápida.}
\stepcounter{questao}
\resposta{$Q=\SI{60}{\coulomb}$}
\resolucao{
\[
Q=It=1{,}5\times40=\SI{60}{\coulomb}.
\]
\textbf{Fundamento:} a corrente é a taxa de passagem de carga,
$i=\mathrm{d}q/\mathrm{d}t$. Sendo constante, a integral vira produto simples.
Um ampère é, por definição, um coulomb por segundo.}
\stepcounter{questao}
\resposta{$I=\SI{2}{\ampere}$}
\resolucao{
\[
I=\frac{Q}{t}=\frac{90}{45}=\SI{2}{\ampere}.
\]
\textbf{Atenção ao termo "média":} o resultado é a corrente \emph{média} no
intervalo. Ela pode ter variado --- inclusive muito --- sem que isso altere a
carga total. Duas formas de onda completamente diferentes podem ter a mesma
média.}
\stepcounter{questao}
\resposta{$t=\SI{240}{\second}=\SI{4}{\minute}$}
\resolucao{
\[
t=\frac{Q}{I}=\frac{60}{0{,}250}=\SI{240}{\second}.
\]
\textbf{Erro comum:} usar $I=\SI{250}{\ampere}$ e obter
$\SI{0.24}{\second}$. Converta o prefixo \emph{antes} de dividir --- o
resultado errado é mil vezes menor e ainda assim parece plausível.}
\stepcounter{questao}
\resposta{$n=\pot{1}{20}$ elétrons}
\resolucao{
\[
n=\frac{Q}{e}=\frac{16}{\pot{1.6}{-19}}=\pot{1}{20}.
\]
Cem quintilhões de elétrons --- e essa carga passa em apenas
$\SI{8}{\second}$ com corrente de $\SI{2}{\ampere}$.\\
\textbf{Ordem de grandeza:} o coulomb é uma unidade \emph{enorme} em termos de
partículas. É por isso que a quantização da carga é imperceptível em circuitos:
um erro de um elétron em $10^{20}$ não se mede.}
\stepcounter{questao}
\resposta{(b)}
\resolucao{O sentido \textbf{convencional} é o do movimento de cargas
\emph{positivas} --- escolha feita antes de se conhecer o elétron. Em metais,
os portadores são elétrons (negativos), que se movem no sentido
\textbf{oposto}.\\
\textbf{Por que a convenção não atrapalha:} uma carga $-q$ movendo-se para a
esquerda é eletricamente equivalente a $+q$ movendo-se para a direita. Todas as
leis de circuitos funcionam identicamente.\\
\textbf{Onde a distinção importa:} em efeitos que dependem do \emph{portador
real}, como o efeito Hall --- cujo sinal revela se os portadores são elétrons ou
lacunas, e é assim que se distingue um semicondutor tipo N de um tipo P.}
\stepcounter{questao}
\resposta{$I=\SI{2}{\ampere}$; $n=\pot{1.5}{21}$}
\resolucao{
\[
I=\frac{240}{120}=\SI{2}{\ampere};
\qquad
n=\frac{240}{\pot{1.6}{-19}}=\pot{1.5}{21}\ \text{elétrons}.
\]
\textbf{Taxa por segundo:} $\pot{1.25}{19}$ elétrons atravessam a seção a cada
segundo. Se fosse possível contá-los um por vez, a $\SI{1}{\per\second}$,
levaria cerca de $400$ bilhões de anos --- trinta vezes a idade do universo.}
\stepcounter{questao}
\resposta{$Q=\SI{7200}{\coulomb}$; $t=\SI{8}{\hour}$}
\resolucao{(a)
\[
Q=2000\,\si{\milli\ampere\hour}=2\,\si{\ampere}\times3600\,\si{\second}
=\SI{7200}{\coulomb}.
\]
(b)
\[
t=\frac{2000\,\si{\milli\ampere\hour}}{250\,\si{\milli\ampere}}=\SI{8}{\hour}.
\]
\textbf{Por que se usa $\si{\ampere\hour}$:} a conta acima sai em uma linha,
enquanto com coulombs seria preciso converter duas vezes. A unidade é prática,
mas é preciso lembrar que ela mede \emph{carga}, não energia.\\
\textbf{Ressalva:} a capacidade nominal pressupõe uma corrente de descarga
específica. Descargas muito rápidas entregam menos carga total --- efeito
Peukert.}
\stepcounter{questao}
\resposta{$Q=\SI{100}{\milli\coulomb}$; $I_{\text{méd}}=\SI{1}{\ampere}$}
\resolucao{(a) $Q=5\times0{,}020=\SI{0.1}{\coulomb}$.\\
(b) O período é $\SI{100}{\milli\second}$:
\[
I_{\text{méd}}=\frac{Q}{T}=\frac{0{,}1}{0{,}1}=\SI{1}{\ampere},
\]
ou, equivalentemente, $I_{pico}\times D=5(0{,}20)$.\\
\textbf{Advertência que vale repetir:} a corrente média de
$\SI{1}{\ampere}$ descreve corretamente a \emph{carga} transportada, mas
\emph{não} o aquecimento. Para efeito Joule, o que vale é o valor eficaz,
$I_{rms}=5\sqrt{0{,}2}=\SI{2.24}{\ampere}$ --- que dissipa cinco vezes mais que
$\SI{1}{\ampere}$ contínuo.}
\stepcounter{questao}
\resposta{$J=\SI{10}{\ampere\per\milli\meter\squared}=\pot{1}{7}\,
\si{\ampere\per\meter\squared}$}
\resolucao{
\[
J=\frac{I}{A}=\frac{15}{1{,}5}=\SI{10}{\ampere\per\milli\meter\squared}
=\frac{15}{\pot{1.5}{-6}}=\pot{1}{7}\,\si{\ampere\per\meter\squared}.
\]
\textbf{Por que $J$ importa mais que $I$:} o aquecimento e a durabilidade de um
condutor dependem da \emph{concentração} de corrente, não do total. Valores
usuais em instalações prediais ficam entre $4$ e
$\SI{6}{\ampere\per\milli\meter\squared}$ --- este condutor está bem acima, e
aqueceria demais.\\
\textbf{Comparação:} trilhas de circuito impresso admitem
$\SI{20}{\ampere\per\milli\meter\squared}$ ou mais, porque dissipam calor
diretamente para o ar e para a placa.}
\stepcounter{questao}
\resposta{$Q=\SI{10}{\milli\coulomb}$; $U=\SI{100}{\volt}$}
\resolucao{(a) $Q=It=\pot{2}{-3}\times5=\SI{10}{\milli\coulomb}$.\\
(b) $U=\dfrac{Q}{C}=\dfrac{\pot{1}{-2}}{\pot{1}{-4}}=\SI{100}{\volt}$.\\
\textbf{Observação:} com corrente \emph{constante}, a tensão cresce
\textbf{linearmente} no tempo --- uma rampa de
$\SI{20}{\volt\per\second}$. É exatamente o oposto do carregamento por fonte de
tensão através de resistor, que produz crescimento exponencial.\\
\textbf{Aplicação:} é assim que se geram rampas precisas em geradores de sinal e
em conversores analógico-digitais de dupla rampa.}
\stepcounter{questao}
\resposta{$I=\SI{5}{\ampere}$ para a direita}
\resolucao{Cargas positivas indo para a direita e cargas negativas indo para a
esquerda produzem corrente convencional \textbf{no mesmo sentido} --- elas se
\emph{somam}:
\[
I=3+2=\SI{5}{\ampere}\ \text{para a direita}.
\]
\textbf{Erro comum:} subtrair, obtendo $\SI{1}{\ampere}$. O sinal negativo do
íon \emph{e} o sentido oposto do movimento se cancelam, resultando em
contribuição positiva.\\
\textbf{Regra:} corrente convencional $=$ (cargas positivas no sentido adotado)
$+$ (cargas negativas no sentido contrário). Em eletrólitos e em plasmas, os
dois tipos de portador contribuem simultaneamente.}
\stepcounter{questao}
\resposta{$t=\SI{3}{\hour}$ teórico; na prática, $20$ a $40\%$ mais}
\resolucao{(a)
\[
t=\frac{1200}{400}=\SI{3}{\hour}.
\]
(b) \textbf{Três razões para o tempo real ser maior.}
\begin{itemize}
\item \textbf{Rendimento de carga:} parte da carga injetada não fica armazenada
--- ela se perde em reações secundárias e em aquecimento. O rendimento
coulômbico típico fica entre $85$ e $95\%$.
\item \textbf{Fase de tensão constante:} carregadores modernos aplicam corrente
constante até certa tensão e depois \emph{reduzem} a corrente
progressivamente. Essa fase final é lenta e pode responder por metade do tempo
total.
\item \textbf{Limitação térmica:} se a bateria aquecer, o carregador reduz a
corrente para protegê-la.
\end{itemize}
\textbf{Estimativa prática:} multiplique o tempo teórico por $1{,}3$ a
$1{,}4$.}
\stepcounter{questao}
\resposta{$t\approx\SI{0.78}{\second}$}
\resolucao{
\[
t=\frac{I^{2}t}{I^{2}}=\frac{50}{64}=\SI{0.78}{\second}.
\]
\textbf{Por que a especificação usa $I^{2}t$:} a energia dissipada no elemento
fusível é $RI^{2}t$, e a fusão ocorre quando essa energia atinge o valor
necessário para fundir o metal. Como $R$ é fixo, o parâmetro característico é
$I^{2}t$.\\
\textbf{Consequência:} dobrar a corrente reduz o tempo de atuação por
\textbf{quatro}. Com $\SI{16}{\ampere}$, o mesmo fusível abriria em
$\SI{0.20}{\second}$; com $\SI{80}{\ampere}$, em
$\SI{7.8}{\milli\second}$.\\
\textbf{Note a diferença conceitual:} aqui a grandeza relevante é
$I^{2}t$ (energia), e não $It$ (carga). São perguntas distintas sobre a mesma
corrente.}
\stepcounter{questao}
\resposta{$Q=\SI{6}{\coulomb}$; $I_{\text{méd}}=\SI{3}{\ampere}$}
\resolucao{(a)
\[
Q=\int_0^{2}(1+2t)\,\mathrm{d}t=\Big[t+t^{2}\Big]_0^{2}=2+4=\SI{6}{\coulomb}.
\]
(b) $I_{\text{méd}}=Q/T=6/2=\SI{3}{\ampere}$.\\
\textbf{Média dos extremos:} $i(0)=1$ e $i(2)=\SI{5}{\ampere}$, cuja média é
$\SI{3}{\ampere}$ --- \emph{coincide}.\\
\textbf{Mas cuidado:} a coincidência ocorre porque a função é \textbf{linear}.
Para $i(t)=t^{2}$ no mesmo intervalo, a média verdadeira seria
$\frac{1}{2}\int_0^{2}t^{2}\mathrm{d}t=\SI{1.33}{\ampere}$, enquanto a média
dos extremos daria $\SI{2}{\ampere}$ --- erro de $50\%$.\\
\textbf{Regra:} a média dos extremos só vale para variação linear. No caso
geral, é preciso integrar.}
\stepcounter{questao}
\resposta{$Q=\SI{15}{\milli\coulomb}$; $I_{\text{méd}}=\SI{3}{\ampere}$}
\resolucao{(a) A carga é a \textbf{área} do triângulo:
\[
Q=\frac{1}{2}\,I_{pico}\,T=\frac{1}{2}(6)(\pot{5}{-3})
=\pot{1.5}{-2}\,\si{\coulomb}=\SI{15}{\milli\coulomb}.
\]
\textbf{Observe:} a repartição $2+3$ do tempo \emph{não} importa para a área ---
qualquer triângulo de base $\SI{5}{\milli\second}$ e altura
$\SI{6}{\ampere}$ tem a mesma área. A assimetria afetaria o valor eficaz, não a
carga.\\
(b) $I_{\text{méd}}=Q/T=0{,}015/0{,}005=\SI{3}{\ampere}$, exatamente metade do pico
--- característica de qualquer forma triangular.}
\stepcounter{questao}
\resposta{$Q=\SI{18}{\coulomb}$; $I_{\text{méd}}=\SI{3}{\ampere}$}
\resolucao{(a) Some as três áreas:
\[
Q=\underbrace{\tfrac12(1)(4)}_{2}+\underbrace{(3)(4)}_{12}
+\underbrace{\tfrac12(2)(4)}_{4}=\SI{18}{\coulomb}.
\]
(b) O intervalo total é $\SI{6}{\second}$:
$I_{\text{méd}}=18/6=\SI{3}{\ampere}$.\\
\textbf{Atalho para trapézios:} a área vale
$I_{pico}\times\dfrac{B+b}{2}$, com $B=\SI{6}{\second}$ (base maior) e
$b=\SI{3}{\second}$ (base menor):
$4\times4{,}5=\SI{18}{\coulomb}$ \checkmark.}
\stepcounter{questao}
\resposta{Líquida: $\SI{0}{\coulomb}$; total: $\SI{24}{\coulomb}$}
\resolucao{(a)
\[
Q_{\text{líq}}=4(3)+(-6)(2)=12-12=\SI{0}{\coulomb}.
\]
(b) Em módulo, passaram $12+12=\SI{24}{\coulomb}$.\\
(c) \textbf{As duas respostas são corretas para perguntas diferentes.}
\begin{itemize}
\item A \textbf{carga líquida} nula significa que nada se acumula: um capacitor
nesse circuito voltaria à tensão inicial a cada ciclo, e uma bateria não se
carregaria nem descarregaria.
\item A \textbf{carga transportada} de $\SI{24}{\coulomb}$ é o que
efetivamente atravessou a seção. Em eletrólise, ela é que produziria depósito
--- se o processo fosse irreversível.
\item O \textbf{aquecimento} depende de outra grandeza ainda:
$I_{rms}=\sqrt{\frac{4^{2}(3)+6^{2}(2)}{5}}=\SI{4.9}{\ampere}$.
\end{itemize}
\textbf{Moral:} "quanta corrente passou" é uma pergunta ambígua. Especifique se
o interesse é acúmulo, transporte ou dissipação.}
\stepcounter{questao}
\resposta{$Q=\SI{0.1}{\coulomb}$; $63{,}2\%$}
\resolucao{(a)
\[
Q=\int_0^{\infty}I_0e^{-t/\tau}\,\mathrm{d}t
=I_0\tau=2(0{,}050)=\SI{0.1}{\coulomb}.
\]
\textbf{Interpretação geométrica:} a carga equivale à corrente inicial mantida
por exatamente $\tau$ --- o retângulo $I_0\times\tau$ tem a mesma área da
exponencial completa.\\
(b)
\[
\frac{Q(\tau)}{Q}=1-e^{-1}=63{,}2\%.
\]
\textbf{Compare com a energia:} na mesma descarga, a energia decai com
constante $\tau/2$ (pois $p\propto i^{2}$), e $86{,}5\%$ dela é dissipada em um
$\tau$. \textbf{Carga e energia têm ritmos diferentes} --- a energia se esgota
mais depressa.}
\stepcounter{questao}
\resposta{$Q=\SI{3600}{\coulomb}$; $m=\SI{1.18}{\gram}$}
\resolucao{(a) $Q=2\times1800=\SI{3600}{\coulomb}$.\\
(b) Pela lei de Faraday da eletrólise:
\[
m=\frac{QM}{zF}=\frac{3600(63{,}5)}{2(96500)}
=\frac{228600}{193000}=\SI{1.18}{\gram}.
\]
\textbf{Significado de $z=2$:} cada íon $\mathrm{Cu^{2+}}$ precisa de
\emph{dois} elétrons para se neutralizar. São necessários, portanto,
$2F$ coulombs por mol depositado.\\
\textbf{Conexão conceitual:} esta é a medida mais direta de que a corrente
transporta carga em quantidades proporcionais à matéria depositada --- foi
justamente por experimentos assim que Faraday inferiu a existência de uma
unidade elementar de carga, décadas antes da descoberta do elétron.}
\stepcounter{questao}
\resposta{$J=\SI{8}{\ampere\per\milli\meter\squared}$ --- excessiva; adotar
$\SI{4}{\milli\meter\squared}$ ou mais}
\resolucao{(a) $J=20/2{,}5=\SI{8}{\ampere\per\milli\meter\squared}$.\\
(b) O valor está $60\%$ acima do limite usual. A seção mínima seria
\[
A=\frac{20}{5}=\SI{4}{\milli\meter\squared},
\]
e a bitola comercial adequada é $\SI{4}{\milli\meter\squared}$ (ou
$\SI{6}{\milli\meter\squared}$, com folga).\\
\textbf{O que acontece se mantivermos $\SI{2.5}{\milli\meter\squared}$:} a
potência dissipada por metro é $\rho J^{2}A=\rho I^{2}/A$, que com a seção
menor é $1{,}6$ vezes maior. A temperatura do condutor sobe, a isolação
envelhece mais depressa e a resistência aumenta --- realimentando o
aquecimento.\\
\textbf{Ressalva:} o limite de $\SI{5}{\ampere\per\milli\meter\squared}$ é
apenas uma regra de bolso. O valor normativo depende do tipo de isolação, do
método de instalação (eletroduto, bandeja, ao ar livre), do agrupamento de
circuitos e da temperatura ambiente.}
\stepcounter{questao}
\resposta{$Q=\SI{47}{\milli\coulomb}$; $t\approx\SI{39}{\minute}$}
\resolucao{(a) $Q=CU=\pot{470}{-6}(100)=\SI{47}{\milli\coulomb}$.\\
(b) Com $I$ constante:
\[
t=\frac{Q}{I}=\frac{\pot{4.7}{-2}}{\pot{2}{-5}}=\SI{2350}{\second}
\approx\SI{39}{\minute}.
\]
(c) \textbf{A hipótese é otimista.} A corrente de fuga é aproximadamente
proporcional à tensão ($I=U/R_{fuga}$), de modo que ela \emph{diminui} durante
a descarga --- e o processo é exponencial, não linear.\\
Com $R_{fuga}=100/\pot{2}{-5}=\SI{5}{\mega\ohm}$, tem-se
$\tau=R_{fuga}C=\SI{2350}{\second}$ --- o \emph{mesmo} número, mas com
significado diferente: após $\SI{39}{\minute}$ resta $37\%$ da tensão, não
zero. A descarga prática ($<1\%$) leva $5\tau\approx\SI{3.3}{\hour}$.\\
\textbf{Implicação de segurança:} um banco de capacitores pode reter tensão
letal por horas. Resistores de descarga permanentes não são opcionais.}
\stepcounter{questao}
\resposta{Mesma carga ($\SI{20}{\coulomb}$); $I_{rms}=2$ e
$\SI{4.47}{\ampere}$; aquecimento $5$ vezes maior em (B)}
\resolucao{(a) Ambas transportam $Q=I_{\text{méd}}t=2(10)=\SI{20}{\coulomb}$ ---
por construção.\\
(b) \textbf{(A)} $I_{rms}=\SI{2}{\ampere}$ (constante).\\
\textbf{(B)} $I_{rms}=\sqrt{0{,}2(10)^{2}}=\sqrt{20}=\SI{4.47}{\ampere}$.\\
(c)
\[
P_A=1(2)^{2}=\SI{4}{\watt};
\qquad
P_B=1(4{,}47)^{2}=\SI{20}{\watt}.
\]
\textbf{Cinco vezes mais}, exatamente o inverso do ciclo de trabalho.\\
\textbf{Síntese fundamental:} \emph{mesma carga, aquecimento completamente
diferente}. A carga depende da média; o aquecimento, do valor eficaz. Duas
grandezas independentes de uma mesma forma de onda.\\
\textbf{Consequência prática:} circuitos pulsados exigem condutores
dimensionados pelo valor \emph{eficaz}, ainda que a carga transportada seja
modesta.}
\stepcounter{questao}
\resposta{$\SI{5}{\milli\coulomb}$ por pulso; $I_{\text{méd}}=\SI{1}{\ampere}$;
$D=2\%$; $I_{rms}=\SI{7.07}{\ampere}$}
\resolucao{(a) $Q=50\times\pot{100}{-6}=\SI{5}{\milli\coulomb}$.
Com $200$ pulsos por segundo,
\[
I_{\text{méd}}=200\times\pot{5}{-3}=\SI{1}{\ampere}.
\]
(b) $D=\dfrac{\pot{100}{-6}}{1/200}=\dfrac{\pot{100}{-6}}{\pot{5}{-3}}
=0{,}02=2\%$.\\
(c) $I_{rms}=I_{pico}\sqrt{D}=50\sqrt{0{,}02}=\SI{7.07}{\ampere}$.\\
\textbf{Três números, três finalidades.}
\begin{itemize}
\item $I_{\text{méd}}=\SI{1}{\ampere}$ dimensiona a \textbf{fonte} e a autonomia da
bateria.
\item $I_{rms}=\SI{7.07}{\ampere}$ dimensiona o \textbf{aquecimento} de
condutores e componentes.
\item $I_{pico}=\SI{50}{\ampere}$ dimensiona a \textbf{corrente máxima}
suportável pelos semicondutores e a queda instantânea na fonte.
\end{itemize}
Usar o número errado leva a projetos ou superdimensionados ou destruídos.}
\stepcounter{questao}
\resposta{$\SI{0.233}{\ampere\hour}$ por ciclo; $I_{\text{méd}}=\SI{0.56}{\ampere}$;
autonomia $\approx\SI{4.3}{\hour}$}
\resolucao{(a) Convertendo os tempos para horas:
\[
Q_{ciclo}=0{,}4\left(\tfrac{20}{60}\right)+1{,}2\left(\tfrac{5}{60}\right)
=0{,}1333+0{,}1000=\SI{0.2333}{\ampere\hour}.
\]
O ciclo dura $\SI{25}{\minute}$, logo
\[
I_{\text{méd}}=\frac{0{,}2333}{25/60}=\SI{0.56}{\ampere}.
\]
(b)
\[
n=\frac{2{,}4}{0{,}2333}=10{,}3\ \text{ciclos}
\;\Longrightarrow\;
t=10{,}3(25)=\SI{257}{\minute}=\SI{4.3}{\hour}.
\]
\emph{(Conferindo: $2{,}4/0{,}56=\SI{4.3}{\hour}$ \checkmark.)}\\
\textbf{Observe:} o patamar de $\SI{1.2}{\ampere}$ dura apenas um quinto do
tempo do outro, mas consome $43\%$ da carga.\\
(c) \textbf{Por que a autonomia real é menor.}
\begin{itemize}
\item \textbf{Efeito Peukert:} a capacidade útil diminui com correntes mais
altas. Os picos de $\SI{1.2}{\ampere}$ extraem menos carga útil do que a
proporção sugere.
\item \textbf{Tensão de corte:} o equipamento para de funcionar antes de a
bateria se esgotar, ao atingir sua tensão mínima --- e isso ocorre \emph{mais
cedo} durante um pico de corrente, por causa da queda em $r$.
\item \textbf{Resistência interna crescente} ao longo da descarga, o que agrava
o item anterior.
\item \textbf{Temperatura:} capacidade menor no frio.
\end{itemize}
\textbf{Estimativa realista:} $\SI{4.3}{\hour}$ é um teto teórico; espere de
$3{,}5$ a $\SI{4}{\hour}$.}
\stepcounter{questao}
\resposta{Erro de $+100\%$ --- o dobro do valor real}
\resolucao{(a) \textbf{Valor real:}
$Q=\frac12(6)(0{,}005)=\SI{15}{\milli\coulomb}$.\\
\textbf{Estimativa:} $6(0{,}005)=\SI{30}{\milli\coulomb}$.\\
\[
\text{erro}=\frac{30-15}{15}=+100\%.
\]
(b) \textbf{Generalização.} Definindo o \emph{fator de forma}
$k=I_{\text{méd}}/I_{pico}$, a estimativa por pico erra sempre por um fator $1/k$:
\begin{center}
\begin{tabular}{lcc}
\hline
Forma de onda & $k$ & erro de $I_{pico}T$\\
\hline
Retangular (contínua) & $1{,}00$ & $0\%$\\
Trapezoidal (exemplo anterior) & $0{,}75$ & $+33\%$\\
Triangular & $0{,}50$ & $+100\%$\\
Exponencial ($T=5\tau$) & $0{,}20$ & $+400\%$\\
Pulsada ($D=2\%$) & $0{,}02$ & $+4900\%$\\
\hline
\end{tabular}
\end{center}
\textbf{O erro é sempre para mais} --- a estimativa por pico é
\emph{conservadora}.\\
(c) \textbf{Regra prática.}
\begin{itemize}
\item Se a intenção é \textbf{garantir margem} (dimensionar uma bateria, por
exemplo), a estimativa por pico é segura, mas pode superdimensionar
absurdamente.
\item Se a intenção é \textbf{prever} a carga, integre --- ou ao menos use o
fator de forma da onda em questão.
\item \textbf{Nunca} use $I_{pico}T$ para estimar autonomia: o
superdimensionamento chega a ordens de grandeza em sinais pulsados.
\end{itemize}}
\stepcounter{questao}
\resposta{Condição: $I_1t_1+I_2t_2=10$ com $t_1+t_2=4$; a solução mais
uniforme minimiza o aquecimento}
\resolucao{(a) \textbf{Condição:}
\[
I_1t_1+I_2t_2=\SI{10}{\coulomb};
\qquad
t_1+t_2=\SI{4}{\second}.
\]
Duas equações e quatro incógnitas --- o problema tem \textbf{dois graus de
liberdade}, e admite infinitas soluções. É preciso um critério adicional.\\
(b) \textbf{Solução A:} $\SI{4}{\ampere}$ por $\SI{1}{\second}$ e
$\SI{2}{\ampere}$ por $\SI{3}{\second}$.
\[
Q=4+6=\SI{10}{\coulomb}\ \checkmark;
\qquad
I_{rms}=\sqrt{\frac{16(1)+4(3)}{4}}=\sqrt{7}=\SI{2.65}{\ampere}.
\]
\textbf{Solução B:} $\SI{8}{\ampere}$ por $\SI{1}{\second}$ e
$\SI{0.667}{\ampere}$ por $\SI{3}{\second}$.
\[
Q=8+2=\SI{10}{\coulomb}\ \checkmark;
\qquad
I_{rms}=\sqrt{\frac{64(1)+0{,}444(3)}{4}}=\sqrt{16{,}3}=\SI{4.04}{\ampere}.
\]
(c) \textbf{A solução A é preferível.} Ela tem valor eficaz $34\%$ menor e,
portanto, aquecimento $2{,}3$ vezes menor --- entregando exatamente a mesma
carga.\\
\textbf{Princípio geral:} para uma dada carga em um dado tempo, o valor eficaz é
\textbf{mínimo} quando a corrente é \emph{constante}
($I=10/4=\SI{2.5}{\ampere}$, $I_{rms}=\SI{2.5}{\ampere}$). Qualquer
variação aumenta $I_{rms}$ acima da média --- é a mesma desigualdade
$\langle i^{2}\rangle\ge\langle i\rangle^{2}$ vista antes.\\
\textbf{Conclusão de projeto:} se a aplicação permitir, use corrente constante.
Se exigir patamares, mantenha-os o mais próximos possível entre si.}
\stepcounter{questao}
\resposta{$\int_0^{T}i\,\mathrm{d}t=0$; os efeitos térmicos e mecânicos
permanecem}
\resolucao{(a) \textbf{Condição:}
\[
\int_0^{T}i(t)\,\mathrm{d}t=0,
\]
ou seja, as áreas positiva e negativa devem ser iguais em módulo. Em termos de
patamares, $I_{+}t_{+}=|I_{-}|t_{-}$.\\
(b) \textbf{Não implica ausência de efeitos.} O que se anula é apenas o
\emph{acúmulo} de carga. Permanecem:
\begin{itemize}
\item \textbf{Aquecimento:} $P=RI_{rms}^{2}$, com
$I_{rms}=\sqrt{\frac1T\int i^{2}\mathrm{d}t}>0$ sempre.
\item \textbf{Forças magnéticas} sobre condutores, que dependem de $i$
instantâneo.
\item \textbf{Desgaste eletroquímico}, se as reações direta e inversa não forem
perfeitamente simétricas.
\end{itemize}
(c) \textbf{Duas aplicações da condição imposta deliberadamente.}
\begin{itemize}
\item \textbf{Estimulação elétrica médica} (marca-passos, eletroestimuladores):
pulsos bifásicos com carga líquida rigorosamente nula evitam acúmulo de
subprodutos eletroquímicos no tecido e corrosão dos eletrodos. É requisito
normativo, não preferência.
\item \textbf{Acionamento de transformadores:} carga líquida não nula
introduziria componente contínua, que satura o núcleo magnético e provoca
corrente de magnetização excessiva. Conversores em ponte são projetados para
garantir simetria.
\end{itemize}
\textbf{E há um terceiro caso instrutivo:} em galvanoplastia usa-se
\emph{deliberadamente} carga líquida não nula --- ali o acúmulo é o objetivo.}
\stepcounter{questao}
\resposta{Média para carga, eficaz para aquecimento, instantâneo para pico}
\resolucao{(a) \textbf{Os instrumentos medem grandezas diferentes de uma mesma
forma de onda.} Um amperímetro de bobina móvel responde à média (por inércia
mecânica); um osciloscópio mostra o valor instantâneo; um multímetro "true RMS"
calcula o valor eficaz.\\
Para a corrente pulsada de $\SI{50}{\ampere}$ com $D=2\%$, os três leriam
$\SI{1}{\ampere}$, $\SI{50}{\ampere}$ e $\SI{7.07}{\ampere}$ --- todos
corretos, todos respondendo perguntas distintas.\\
(b)
\begin{center}
\begin{tabular}{ll}
\hline
Objetivo & Grandeza\\
\hline
Carga transferida, autonomia de bateria & \textbf{média}\\
Aquecimento, dimensionamento de condutor & \textbf{eficaz}\\
Corrente máxima em semicondutor, saturação & \textbf{pico}\\
\hline
\end{tabular}
\end{center}
(c) \textbf{Projeto de medição completa.} Use um \emph{resistor shunt} de baixo
valor e observe a tensão sobre ele com osciloscópio digital:
\begin{enumerate}
\item O \textbf{traço instantâneo} fornece o pico diretamente.
\item A função \textbf{média} do osciloscópio integra a forma de onda e entrega
$I_{\text{méd}}$ --- e portanto a carga, multiplicando pelo tempo.
\item A função \textbf{RMS} entrega $I_{rms}$.
\end{enumerate}
\textbf{Cuidados:} a largura de banda do osciloscópio e do shunt deve ser muito
maior que a taxa de subida dos pulsos, sob pena de arredondá-los e subestimar o
pico. E a indutância parasita do shunt introduz erro proporcional a
$\mathrm{d}i/\mathrm{d}t$ --- use shunts coaxiais em pulsos rápidos.}
\stepcounter{questao}
\resposta{O erro relativo decresce como $1/\sqrt{N}$; integrar piora se houver
\emph{deriva} (ruído de média não nula)}
\resolucao{(a) \textbf{Por que integrar ajuda.} A carga é
$Q=\int i\,\mathrm{d}t$. Se o ruído tem média zero, suas contribuições
positivas e negativas tendem a se cancelar na integral, enquanto o sinal
verdadeiro se acumula. A relação sinal-ruído \emph{melhora} com o tempo.\\
(b) \textbf{Decrescimento do erro.} Para $N$ amostras independentes de ruído de
desvio $\sigma$, o desvio da \emph{soma} é $\sigma\sqrt{N}$, enquanto a soma do
sinal cresce com $N$. Logo o erro relativo cai como
\[
\frac{\sigma\sqrt{N}}{SN}=\frac{\sigma}{S\sqrt{N}}\propto\frac{1}{\sqrt{T}}.
\]
\textbf{Quantificando:} integrar por $\SI{100}{\second}$ em vez de
$\SI{1}{\second}$ reduz o erro relativo por um fator $10$. Para ganhar outro
fator $10$, seriam necessários $\SI{10000}{\second}$ --- rendimento
decrescente.\\
(c) \textbf{Quando integrar piora.} Se o ruído \textbf{não} tiver média zero ---
isto é, se houver \emph{offset} ou \emph{deriva} no instrumento --- o erro
\textbf{também se acumula}, e cresce \emph{linearmente} com $T$:
\[
\text{erro}=\varepsilon_{offset}\times T.
\]
Enquanto o ruído aleatório cai com $\sqrt{T}$, o erro sistemático cresce com
$T$. Existe, portanto, um \textbf{tempo ótimo} de integração, além do qual a
deriva domina.\\
\textbf{Prática de laboratório:} meça a leitura com \emph{entrada em curto}
antes do ensaio, para quantificar o \emph{offset}, e subtraia-o. Sem essa
correção, integrações longas são pior que inúteis --- elas produzem números
precisos e errados.}
\stepcounter{questao}
\resposta{Com $C\approx\SI{1}{\pico\farad}$, basta $\SI{1}{\pico\coulomb}$ para
$\SI{1}{\volt}$ --- um desequilíbrio de $\SI{1}{\micro\ampere}$ por
$\SI{1}{\micro\second}$}
\resolucao{(a) Um nó real (uma solda, um trecho de trilha) tem capacitância
parasita para o ambiente da ordem de $C\approx\SI{1}{\pico\farad}$. Se uma
carga $\Delta Q$ se acumulasse ali, o potencial subiria:
\[
\Delta V=\frac{\Delta Q}{C}.
\]
Para $\Delta V=\SI{1}{\volt}$:
\[
\Delta Q=\pot{1}{-12}\,\si{\coulomb}
=\frac{\pot{1}{-12}}{\pot{1.6}{-19}}=\pot{6.25}{6}\ \text{elétrons}.
\]
Apenas seis milhões de elétrons --- número minúsculo diante dos
$10^{19}$ que atravessam a seção a cada segundo com $\SI{2}{\ampere}$.\\
(b) O desequilíbrio necessário é irrisório:
\[
\Delta I\,\Delta t=\SI{1}{\pico\coulomb}
\;\Longrightarrow\;
\SI{1}{\micro\ampere}\ \text{durante}\ \SI{1}{\micro\second},
\]
ou $\SI{1}{\nano\ampere}$ durante $\SI{1}{\milli\second}$. Em um circuito de
$\SI{2}{\ampere}$, isso corresponde a um desequilíbrio relativo de
$\pot{5}{-7}$ --- meio milionésimo.\\
(c) \textbf{Conclusão: a LKC não é uma aproximação, é uma consequência
autorregulada.} Qualquer tentativa de acumular carga em um nó eleva
imediatamente seu potencial, e esse potencial \emph{expulsa} o excesso --- as
correntes se reajustam em nanossegundos.\\
\textbf{Onde a lei realmente falha:} em frequências altíssimas, quando o tempo
de propagação pelo circuito deixa de ser desprezível, correntes de deslocamento
por capacitâncias parasitas passam a importar e a LKC "de nó concentrado" perde
sentido. É preciso então tratar o circuito por linhas de transmissão ou por
campos.\\
\textbf{Critério prático:} a LKC vale enquanto as dimensões do circuito forem
muito menores que o comprimento de onda do sinal --- a chamada aproximação de
\emph{parâmetros concentrados}.}
\stepcounter{questao}
\resposta{$v_d=\SI{0.29}{\milli\meter\per\second}$; $\SI{57}{\minute}$ para
$\SI{1}{\meter}$; o sinal leva $\SI{5}{\nano\second}$}
\resolucao{(a)
\[
v_d=\frac{I}{nAe}
=\frac{10}{(\pot{8.5}{28})(\pot{2.5}{-6})(\pot{1.6}{-19})}
=\frac{10}{\pot{3.4}{4}}=\pot{2.94}{-4}\,\si{\meter\per\second}.
\]
Menos de um terço de milímetro por segundo.\\
(b)
\[
t=\frac{1}{\pot{2.94}{-4}}=\SI{3400}{\second}\approx\SI{57}{\minute}.
\]
(c) \textbf{O sinal viaja a cerca de $\pot{2}{8}\,\si{\meter\per\second}$}
(dois terços de $c$, em cabo típico), percorrendo o metro em
$\SI{5}{\nano\second}$ --- \textbf{$\pot{6.8}{11}$ vezes mais rápido} que os
elétrons.\\
\textbf{Resolução do paradoxo aparente.} O que se propaga rapidamente não são os
elétrons, mas o \textbf{campo eletromagnético} que os põe em movimento. Ao
fechar o interruptor, o campo se estabelece ao longo de todo o condutor quase
instantaneamente, e \emph{todos} os elétrons começam a se mover praticamente ao
mesmo tempo --- cada um localmente, e devagar.\\
\textbf{Analogia:} um cano cheio de água. Ao abrir a torneira, a água sai
imediatamente na outra ponta, embora nenhuma molécula tenha percorrido o cano.
O que se propaga é a \emph{pressão}, não a matéria.\\
\textbf{Consequência prática:} a lâmpada acende no instante em que se aciona o
interruptor, mas os elétrons que estavam junto dele levariam quase uma hora para
chegar até ela.}

\gabarito[2.1 Corrente elétrica e carga]

\subsection{Lei de Ohm e resistividade}
\stepcounter{questao}
\resposta{(b)}
\resolucao{De $R = \rho\dfrac{L}{A}$: a resistência é \emph{diretamente}
proporcional ao comprimento $L$ e \emph{inversamente} proporcional à área $A$.
O material entra pela resistividade $\rho$, que por sua vez varia com a
temperatura --- o que torna (c) e (d) incorretas.}
\stepcounter{questao}
\resposta{(a) $k=\SI{2}{\ohm\per\meter}$; (b) \SI{10}{\ohm}; (c) 2ª lei de Ohm}
\resolucao{(a) A reta passa por $(3, 6)$, logo
$k = \dfrac{R}{L} = \dfrac{6}{3} = \SI{2}{\ohm\per\meter}$.\\
(b) $R = kL = 2\cdot 5 = \SI{10}{\ohm}$.\\
(c) A proporcionalidade direta entre $R$ e $L$ (com $A$ e material fixos) confirma
a segunda lei de Ohm, $R = \rho\dfrac{L}{A}$, em que $k = \dfrac{\rho}{A}$.}
\stepcounter{questao}
\resposta{$L \approx \SI{4.55}{\meter}$}
\resolucao{Isolando $L$ na segunda lei de Ohm (com
$A = \SI{0.5}{\milli\meter\squared} = \pot{0.5}{-6}\,\si{\meter\squared}$):
\[
L=\frac{R\,A}{\rho}=\frac{10\cdot\pot{0.5}{-6}}{\pot{1.1}{-6}}
=\frac{\pot{5}{-6}}{\pot{1.1}{-6}}\approx\SI{4.55}{\meter}.
\]}
\stepcounter{questao}
\resposta{(c)}
\resolucao{A área da seção é proporcional ao quadrado do diâmetro, $A \propto D^2$.
Ao dobrar $D$, a área quadruplica ($A' = 4A$); ao dobrar $L$, o comprimento
duplica ($L' = 2L$):
\[
R'=\rho\frac{L'}{A'}=\rho\frac{2L}{4A}=\frac{1}{2}\,\rho\frac{L}{A}=\frac{R}{2}.
\]
A armadilha (alternativa b) é considerar só o efeito do comprimento e esquecer que
a área cresce com o \emph{quadrado} do diâmetro.}
\stepcounter{questao}
\resposta{(b)}
\resolucao{Para um condutor ôhmico, $U = R\,i$: a resistência é a \emph{inclinação}
da reta. Tomando o ponto $(4, 16)$:
\[
R=\frac{U}{i}=\frac{16}{4}=\SI{4}{\ohm}.
\]
A reta passar pela origem confirma que o resistor obedece à lei de Ohm.}
\stepcounter{questao}
\resposta{$A_1/A_2 = 1$ (áreas iguais)}
\resolucao{Como o material é o mesmo, $\rho$ é comum. De $R = \rho L/A$:
\[
\frac{R_1}{R_2}=\frac{L_1/A_1}{L_2/A_2}=\frac{L_1}{L_2}\cdot\frac{A_2}{A_1}.
\]
Com $\dfrac{R_1}{R_2}=\dfrac{6}{3}=2$ e $\dfrac{L_1}{L_2}=2$:
\[
2 = 2\cdot\frac{A_2}{A_1}\ \Rightarrow\ \frac{A_2}{A_1}=1
\ \Rightarrow\ A_1 = A_2.
\]
Ou seja, dobrar o comprimento já explica a resistência dobrada; as seções são
iguais.}
\stepcounter{questao}
\resposta{(c)}
\resolucao{Com $L$ e $A$ iguais, $R = \rho\dfrac{L}{A}$ é proporcional a $\rho$:
\[
\frac{R_{\text{Al}}}{R_{\text{Cu}}}=\frac{\rho_{\text{Al}}}{\rho_{\text{Cu}}}
=\frac{2{,}8}{1{,}7}\approx1{,}6.
\]
Por isso, para a mesma resistência, um fio de alumínio precisa ser mais grosso
que um de cobre --- o que se compensa por ele ser mais leve e barato, motivo do
seu uso em linhas de transmissão.}
\stepcounter{questao}
\resposta{$R = \SI{34}{\milli\ohm}$}
\resolucao{Com $A = \SI{1}{\milli\meter\squared} = \pot{1}{-6}\,\si{\meter\squared}$:
\[
R=\rho\frac{L}{A}=\pot{1.7}{-8}\cdot\frac{2}{\pot{1}{-6}}
=\pot{3.4}{-2}\,\si{\ohm}=\SI{34}{\milli\ohm}.
\]
Resistência bem baixa, como esperado de um bom condutor em curto comprimento.}
\stepcounter{questao}
\resposta{$R'=4R$}
\resolucao{Pela segunda lei de Ohm, $R=\rho L/A$. Logo,
\[
R'=\rho\frac{2L}{A/2}=4\rho\frac{L}{A}=4R.
\]}
\stepcounter{questao}
\resposta{(a) $I_{\text{cobre}} = 10\,I_{\text{aço}}$;
(b) $P_{\text{cobre}} = 10\,P_{\text{aço}}$}
\resolucao{Como $L$ e $A$ são iguais, $R = \rho L/A \Rightarrow
R_{\text{aço}} = 10\,R_{\text{cobre}}$.\\
(a) Sob a mesma tensão, $I = U/R$, logo a corrente é inversamente proporcional a
$R$:
\[
\frac{I_{\text{cobre}}}{I_{\text{aço}}}
=\frac{R_{\text{aço}}}{R_{\text{cobre}}}=10
\ \Rightarrow\ I_{\text{cobre}}=10\,I_{\text{aço}}.
\]
(b) $P = U^2/R$; com $U$ igual, a potência também é inversamente proporcional a
$R$:
\[
\frac{P_{\text{cobre}}}{P_{\text{aço}}}=\frac{R_{\text{aço}}}{R_{\text{cobre}}}=10.
\]
O cobre, menos resistivo, conduz mais corrente e dissipa mais potência sob a mesma
tensão --- por isso é preferido em condutores.}
\stepcounter{questao}
\resposta{(a) $U_{\text{Al}}/U_{\text{Cu}}\approx1{,}65$;
(b) $\SI{4.53}{\volt}$ e $\SI{7.47}{\volt}$; (c) no alumínio}
\resolucao{(a) Com $L$ e $A$ iguais, $R \propto \rho$; e como estão em série
(mesma corrente), $U = RI \propto \rho$:
\[
\frac{U_{\text{Al}}}{U_{\text{Cu}}}=\frac{\rho_{\text{Al}}}{\rho_{\text{Cu}}}
=\frac{2{,}8}{1{,}7}\approx1{,}65.
\]
(b) Pelo divisor de tensão, a repartição segue as resistividades:
\[
U_{\text{Cu}}=12\cdot\frac{1{,}7}{4{,}5}\approx\SI{4.53}{\volt},
\qquad
U_{\text{Al}}=12\cdot\frac{2{,}8}{4{,}5}\approx\SI{7.47}{\volt}.
\]
(c) A potência é $P = UI$ com $I$ comum, logo \emph{maior tensão implica maior
potência}: o \textbf{alumínio} aquece mais. Fisicamente, sendo mais resistivo,
ele opõe mais resistência à mesma corrente.\\
\textbf{Implicação prática:} emendas entre metais diferentes são pontos críticos
em instalações --- além do aquecimento desigual, os coeficientes de dilatação
distintos afrouxam o contato com o tempo, aumentando a resistência local e
realimentando o aquecimento. É a causa clássica de incêndios em conexões
Cu--Al, razão pela qual normas exigem conectores bimetálicos específicos.}
\stepcounter{questao}
\resposta{(a) $T = \SI{80}{\celsius}$; (b) a potência cai $\approx\SI{19}{\percent}$;
(c) ver comentário}
\resolucao{\textbf{(a)} Isolando a temperatura:
\[
12{,}4 = 10\big[1+\pot{4}{-3}(T-20)\big]
\;\Rightarrow\;
1{,}24 = 1+\pot{4}{-3}(T-20)
\;\Rightarrow\;
T-20 = \frac{0{,}24}{\pot{4}{-3}} = 60,
\]
logo $T = \SI{80}{\celsius}$.

\textbf{(b)} Com tensão fixa, $P = \dfrac{U^2}{R}$ é inversamente proporcional a
$R$:
\[
\frac{P_{\text{quente}}}{P_{\text{frio}}}=\frac{R_{\text{frio}}}{R_{\text{quente}}}
=\frac{10}{12{,}4}=0{,}806,
\]
ou seja, a potência dissipada \emph{cai} cerca de \SI{19.4}{\percent}.

\textbf{(c) Método da variação de resistência.} Como $\alpha$ do cobre é bem
conhecido e estável, mede-se a resistência do enrolamento a frio (com o motor
parado e em temperatura ambiente conhecida) e novamente logo após o desligamento.
Invertendo a fórmula:
\[
T = T_0 + \frac{1}{\alpha}\left(\frac{R}{R_0}-1\right),
\]
obtém-se a temperatura \emph{média} de todo o enrolamento --- justamente o que
importa para a vida útil do isolamento --- sem instalar nenhum sensor interno.
Esse é o procedimento normalizado em ensaios de elevação de temperatura de
motores e transformadores.\\
\textbf{Observação crítica:} o método fornece a média, não o \emph{ponto quente}.
Como a degradação do isolante segue uma lei exponencial com a temperatura, o ponto
mais quente pode falhar antes do que a média sugeriria --- por isso as normas
aplicam margens de segurança sobre esse valor.}
\stepcounter{questao}
\resposta{(a) $\SI{1.0e-10}{\ohm\metre\per\celsius}$; (b) $\SI{3.0e-8}{\ohm\metre}$}
\resolucao{A inclinação é
\[
\frac{\Delta\rho}{\Delta T}=\frac{(2.5-1.5)\times10^{-8}}{100}
=\SI{1.0e-10}{\ohm\metre\per\celsius}.
\]
Extrapolando até $\SI{150}{\celsius}$:
$\rho=1.5\times10^{-8}+150\times10^{-10}=\SI{3.0e-8}{\ohm\metre}$.}
\stepcounter{questao}
\resposta{$U=\SI{7.05}{\volt}$}
\resolucao{
\[
U=RI=470\times\pot{15}{-3}=\SI{7.05}{\volt}.
\]
\textbf{Cuidado com o prefixo:} usar $I=\SI{15}{\ampere}$ daria
$\SI{7050}{\volt}$ --- resultado absurdo, mas que passa despercebido se a
conversão não for feita explicitamente.}
\stepcounter{questao}
\resposta{$R=\SI{300}{\ohm}$}
\resolucao{
\[
R=\frac{U}{I}=\frac{9}{0{,}030}=\SI{300}{\ohm}.
\]
\textbf{Atenção à hipótese:} a razão $U/I$ sempre pode ser calculada, mas só se
chama \emph{resistência} no sentido pleno se o componente for ôhmico --- isto é,
se essa razão for a mesma para qualquer tensão aplicada. Em componentes não
lineares ela varia, e recebe o nome de \emph{resistência estática}.}
\stepcounter{questao}
\resposta{$I=\SI{5.45}{\milli\ampere}$}
\resolucao{
\[
I=\frac{U}{R}=\frac{12}{2200}=\pot{5.45}{-3}\,\si{\ampere}
=\SI{5.45}{\milli\ampere}.
\]
\textbf{Atalho útil:} volts divididos por quilohms dão \emph{miliampères}
diretamente --- $12/2{,}2=5{,}45$. Essa combinação de prefixos aparece o tempo
todo em eletrônica e poupa conversões.}
\stepcounter{questao}
\resposta{$R=\SI{0.34}{\ohm}$}
\resolucao{Convertendo a área para $\si{\meter\squared}$:
$A=\SI{2.5}{\milli\meter\squared}=\pot{2.5}{-6}\,\si{\meter\squared}$.
\[
R=\frac{\rho L}{A}=\frac{\pot{1.7}{-8}\times50}{\pot{2.5}{-6}}
=\SI{0.34}{\ohm}.
\]
\textbf{Erro clássico:} esquecer que
$\SI{1}{\milli\meter\squared}=\pot{1}{-6}\,\si{\meter\squared}$, e não
$\pot{1}{-3}$. A área é o \emph{quadrado} de um comprimento, logo o fator de
conversão também é quadrático.}
\stepcounter{questao}
\resposta{$G=\SI{0.04}{\ensuremath{\mathrm{S}}}$}
\resolucao{
\[
G=\frac{1}{R}=\frac{1}{25}=\SI{0.04}{\ensuremath{\mathrm{S}}}.
\]
A condutância mede a \emph{facilidade} de conduzir. Ela é a grandeza natural
quando os elementos estão em paralelo (somam-se diretamente) e na análise nodal.
Nas relações de material, define-se analogamente a \textbf{condutividade}
$\sigma=1/\rho$.}
\stepcounter{questao}
\resposta{(a)}
\resolucao{Como $R=\rho L/A$, tem-se $R\propto L$: dobrar $L$ dobra $R$.\\
\textbf{Interpretação física:} o dobro do comprimento equivale a dois fios
idênticos em \emph{série} --- e resistências em série se somam. A analogia
hidráulica ajuda: um cano duas vezes mais longo oferece o dobro de oposição ao
escoamento.}
\stepcounter{questao}
\resposta{(b)}
\resolucao{$R\propto1/A$: dobrar $A$ divide $R$ por dois.\\
\textbf{Interpretação:} o dobro da área equivale a dois fios idênticos em
\emph{paralelo}, e o paralelo de dois iguais vale metade.\\
\textbf{Cuidado com o diâmetro:} dobrar o \emph{diâmetro} quadruplica a área
($A=\pi d^{2}/4$) e divide $R$ por \textbf{quatro} --- confundir diâmetro com
área é o deslize mais comum aqui.}
\stepcounter{questao}
\resposta{$R_1/R_2=4$}
\resolucao{A área é proporcional ao quadrado do diâmetro:
\[
\frac{A_2}{A_1}=\left(\frac{2d}{d}\right)^{2}=4
\;\Longrightarrow\;
\frac{R_1}{R_2}=\frac{A_2}{A_1}=4.
\]
O fio mais grosso tem um quarto da resistência. É por isso que a bitola dos
condutores é especificada em $\si{\milli\meter\squared}$ (área) e não em
milímetros de diâmetro --- a área é o que entra linearmente na conta.}
\stepcounter{questao}
\resposta{$L=\SI{4}{\meter}$}
\resolucao{Isolando $L$:
\[
L=\frac{RA}{\rho}=\frac{22\times\pot{0.2}{-6}}{\pot{1.1}{-6}}
=\frac{\pot{4.4}{-6}}{\pot{1.1}{-6}}=\SI{4}{\meter}.
\]
\textbf{Por que níquel-cromo:} sua resistividade é cerca de $65$ vezes a do
cobre, o que permite resistências úteis com comprimentos manejáveis. Com cobre,
os mesmos $\SI{22}{\ohm}$ exigiriam $\SI{259}{\meter}$ de fio --- inviável.}
\stepcounter{questao}
\resposta{$A=\SI{11.2}{\milli\meter\squared}$; $d=\SI{3.78}{\milli\meter}$}
\resolucao{
\[
A=\frac{\rho L}{R}=\frac{\pot{2.8}{-8}\times200}{0{,}5}
=\pot{1.12}{-5}\,\si{\meter\squared}=\SI{11.2}{\milli\meter\squared}.
\]
\[
d=\sqrt{\frac{4A}{\pi}}=\sqrt{\frac{4\times\pot{1.12}{-5}}{\pi}}
=\SI{3.78}{\milli\meter}.
\]
\textbf{Valor comercial:} a bitola imediatamente superior é
$\SI{16}{\milli\meter\squared}$, que daria $R=\SI{0.35}{\ohm}$ --- com folga.
Adotar $\SI{10}{\milli\meter\squared}$ resultaria em $\SI{0.56}{\ohm}$,
acima do especificado.}
\stepcounter{questao}
\resposta{$\rho=\pot{4.2}{-8}\,\si{\ohm\meter}$ --- compatível com latão ou
zinco}
\resolucao{
\[
\rho=\frac{RA}{L}=\frac{2{,}5\times\pot{0.5}{-6}}{30}
=\pot{4.17}{-8}\,\si{\ohm\meter}.
\]
\textbf{Comparação:} cobre vale $\pot{1.7}{-8}$, alumínio
$\pot{2.8}{-8}$, zinco $\pot{5.9}{-8}$ e latão em torno de
$\pot{7}{-8}\,\si{\ohm\meter}$. O valor obtido está entre alumínio e zinco.\\
\textbf{Ressalva honesta:} a resistividade sozinha \emph{não} identifica o
material com segurança --- ligas de composições distintas podem coincidir. Ela
serve para descartar hipóteses (não é cobre puro, certamente) e como um entre
vários indícios.}
\stepcounter{questao}
\resposta{$R'=4R$}
\resolucao{\textbf{Volume constante:} $V=AL=A'L'$. Com $L'=2L$:
\[
A'=\frac{AL}{2L}=\frac{A}{2}.
\]
\textbf{Nova resistência:}
\[
R'=\frac{\rho L'}{A'}=\frac{\rho(2L)}{A/2}=4\,\frac{\rho L}{A}=4R.
\]
\textbf{Dois efeitos que se multiplicam:} o comprimento dobra (fator $2$) e a
área cai à metade (outro fator $2$).\\
\textbf{Não confunda} com o caso em que $L$ e $A$ são alterados
\emph{independentemente} --- ali os fatores podem se cancelar. No estiramento,
eles se reforçam, porque o volume amarra as duas dimensões.}
\stepcounter{questao}
\resposta{$R=\SI{61.7}{\ohm}$}
\resolucao{
\[
R=R_0\left[1+\alpha(T-T_0)\right]=50\left[1+0{,}0039(60)\right]
=50(1{,}234)=\SI{61.7}{\ohm}.
\]
Aumento de $23{,}4\%$ --- nada desprezível.\\
\textbf{Consequência para medição:} um ohmímetro que aplique corrente suficiente
para aquecer o componente lerá um valor \emph{maior} que o real a temperatura
ambiente. É por isso que medições de precisão especificam a corrente de teste e
a temperatura.}
\stepcounter{questao}
\resposta{Não --- a razão $U/I$ cresce de $100$ para $\SI{105}{\ohm}$}
\resolucao{Calculando $R=U/I$ em cada ponto:
\[
\frac{2}{0{,}020}=100;
\qquad
\frac{4}{0{,}040}=100;
\qquad
\frac{6}{0{,}057}=\SI{105}{\ohm}.
\]
Os dois primeiros pontos são consistentes; o terceiro destoa em $5\%$.\\
\textbf{Diagnóstico provável: autoaquecimento.} Em $\SI{6}{\volt}$ a potência é
$\SI{0.34}{\watt}$, contra $\SI{0.04}{\watt}$ no primeiro ponto --- oito vezes
mais. Se o material tiver coeficiente positivo, a resistência sobe com a
temperatura.\\
\textbf{Como confirmar:} refaça o terceiro ponto \emph{rapidamente}, antes que o
componente aqueça. Se a leitura voltar a $\SI{100}{\ohm}$, o material é ôhmico
--- e o desvio era térmico, não uma propriedade elétrica intrínseca.}
\stepcounter{questao}
\resposta{$A_{Al}/A_{Cu}=1{,}65$}
\resolucao{Para $R$ e $L$ iguais, $A\propto\rho$:
\[
\frac{A_{Al}}{A_{Cu}}=\frac{\rho_{Al}}{\rho_{Cu}}=\frac{2{,}8}{1{,}7}=1{,}65.
\]
O alumínio exige $65\%$ mais seção.\\
\textbf{Por que ainda assim se usa alumínio:} sua densidade é
$\SI{2.7}{\gram\per\centi\meter\cubed}$ contra $\SI{8.9}{}$ do cobre. Para a
mesma resistência, o cabo de alumínio pesa
$1{,}65\times(2{,}7/8{,}9)=0{,}50$ --- \textbf{metade} do peso, e custa bem
menos. É por isso que linhas de transmissão aéreas são de alumínio, onde o peso
é crítico, enquanto instalações prediais são de cobre, onde o espaço nos
eletrodutos é que limita.}
\stepcounter{questao}
\resposta{$R=\SI{1.13}{\ohm}$; $\Delta U=\SI{11.3}{\volt}$;
$P=\SI{113}{\watt}$}
\resolucao{
\[
R=\frac{\pot{1.7}{-8}\times100}{\pot{1.5}{-6}}=\SI{1.13}{\ohm};
\qquad
\Delta U=RI=\SI{11.3}{\volt};
\qquad
P=RI^{2}=\SI{113}{\watt}.
\]
\textbf{Avaliação:} em uma rede de $\SI{220}{\volt}$, essa queda representa
$5{,}2\%$ --- acima do limite usual de $4\%$ para circuitos terminais. E
$\SI{113}{\watt}$ dissipados ao longo do cabo é aquecimento real, que exige
verificação de capacidade de condução.\\
\textbf{Conclusão:} a bitola está subdimensionada para esse comprimento. O
critério de queda de tensão, e não o de corrente admissível, é que domina em
circuitos longos.}
\stepcounter{questao}
\resposta{(a) $L=\SI{3.93}{\meter}$; (b) $I=\SI{0.707}{\ampere}$}
\resolucao{(a) Seção:
$A=\dfrac{\pi d^{2}}{4}=\dfrac{\pi(\pot{5}{-4})^{2}}{4}
=\pot{1.963}{-7}\,\si{\meter\squared}$.
\[
L=\frac{RA}{\rho}=\frac{10\times\pot{1.963}{-7}}{\pot{5}{-7}}
=\SI{3.93}{\meter}.
\]
(b) $I=\sqrt{P/R}=\sqrt{5/10}=\SI{0.707}{\ampere}$, com
$U=\SI{7.07}{\volt}$.\\
(c) \textbf{Construção.} Quase quatro metros de fio precisam ser enrolados
sobre um corpo cerâmico. Três cuidados:
\begin{itemize}
\item \textbf{Espaçamento entre espiras}, para evitar curtos e permitir
circulação de ar.
\item \textbf{Área de dissipação suficiente} --- $\SI{5}{\watt}$ exigem corpo de
alguns centímetros, e a temperatura de superfície pode passar de
$\SI{100}{\celsius}$.
\item \textbf{Indutância parasita}: um enrolamento simples de $\SI{3.93}{\meter}$
comporta-se como bobina. Em corrente contínua isso é irrelevante, mas em sinais
rápidos é preciso enrolamento \emph{bifilar}, em que a ida e a volta se
cancelam.
\end{itemize}}
\stepcounter{questao}
\resposta{de $\SI{94.2}{\ohm}$ a $\SI{107.1}{\ohm}$ ($-5{,}8\%$ a
$+7{,}1\%$)}
\resolucao{\textbf{Deriva térmica:} $400\,\mathrm{ppm}=\pot{4}{-4}$ por
$\si{\celsius}$.
\[
\SI{70}{\celsius}:\ \Delta T=+50 \Rightarrow +2{,}0\%;
\qquad
\SI{0}{\celsius}:\ \Delta T=-20 \Rightarrow -0{,}8\%.
\]
\textbf{Combinação de pior caso} (os dois efeitos no mesmo sentido):
\[
R_{\max}=100(1{,}05)(1{,}020)=\SI{107.1}{\ohm};
\qquad
R_{\min}=100(0{,}95)(0{,}992)=\SI{94.2}{\ohm}.
\]
\textbf{Observações de projeto.}
\begin{itemize}
\item A tolerância de fabricação domina --- ela responde por $5$ dos
$7{,}1$ pontos.
\item A faixa é \emph{assimétrica}, porque a temperatura de operação se afasta
mais para cima que para baixo do valor de referência.
\item Se $7\%$ for demais, apertar a tolerância para $1\%$ traria a faixa a
$-1{,}8\%$/$+3{,}0\%$; trocar o coeficiente para
$25\,\mathrm{ppm}$ traria a $-5{,}0\%$/$+5{,}1\%$. \textbf{Ataque primeiro a
maior parcela.}
\end{itemize}}
\stepcounter{questao}
\resposta{$R_{est}=\SI{6}{\ohm}$; $r_{din}=\SI{25}{\ohm}$}
\resolucao{(a) $R_{est}=\dfrac{U}{I}=\dfrac{12{,}0}{2{,}00}=\SI{6}{\ohm}$.\\
(b) $r_{din}=\dfrac{\Delta U}{\Delta I}=\dfrac{0{,}5}{0{,}02}=\SI{25}{\ohm}$.\\
(c) \textbf{A dinâmica é mais de quatro vezes maior.} A razão é o
\textbf{autoaquecimento}: aumentar a tensão eleva a potência, o filamento
esquenta e sua resistência sobe --- de modo que a corrente cresce muito menos
que proporcionalmente.\\
\textbf{Cada grandeza responde a uma pergunta diferente:}
\begin{itemize}
\item $R_{est}$ governa o \emph{balanço de potência} no ponto de operação:
$P=U^{2}/R_{est}=\SI{24}{\watt}$ \checkmark.
\item $r_{din}$ governa a \emph{resposta a pequenas variações} --- é ela que
entra em qualquer modelo incremental.
\end{itemize}
\textbf{A lâmpada não é ôhmica em regime}, mas é aproximadamente ôhmica a
temperatura constante. É o acoplamento térmico, e não a física da condução, que
quebra a linearidade.}
\stepcounter{questao}
\resposta{$R\le\SI{0.33}{\ohm}$; $A\ge\SI{4.12}{\milli\meter\squared}$;
adotar $\SI{6}{\milli\meter\squared}$}
\resolucao{(a) Queda admissível: $0{,}03\times220=\SI{6.6}{\volt}$.
\[
R_{\max}=\frac{6{,}6}{20}=\SI{0.33}{\ohm}.
\]
(b) \textbf{Atenção ao comprimento:} a corrente percorre \emph{ida e volta},
logo $L=2\times40=\SI{80}{\meter}$.
\[
A\ge\frac{\rho L}{R_{\max}}=\frac{\pot{1.7}{-8}\times80}{0{,}33}
=\pot{4.12}{-6}\,\si{\meter\squared}=\SI{4.12}{\milli\meter\squared}.
\]
(c) A bitola comercial imediatamente superior é
$\SI{6}{\milli\meter\squared}$, que dá $R=\SI{0.227}{\ohm}$ e queda de
$\SI{4.5}{\volt}$ ($2{,}1\%$) --- dentro da especificação com folga.\\
\textbf{Verificação obrigatória do outro critério:} $\SI{6}{\milli\meter\squared}$
suporta cerca de $\SI{41}{\ampere}$ em eletroduto, bem acima dos
$\SI{20}{\ampere}$ da carga. \textbf{Aqui a queda de tensão é o critério
dominante} --- e é o que costuma ocorrer em circuitos longos. Em circuitos
curtos, a corrente admissível é que decide.}
\stepcounter{questao}
\resposta{$R'=n^{2}R$; $9R$ e $25R$; $d'/d=1/\sqrt{n}$}
\resolucao{(a) Volume constante: $AL=A'L'$ com $L'=nL$, logo
$A'=A/n$.
\[
R'=\frac{\rho(nL)}{A/n}=n^{2}\frac{\rho L}{A}=n^{2}R.
\]
(b) $n=3\Rightarrow9R$; $n=5\Rightarrow25R$.\\
(c) Como $A\propto d^{2}$ e $A'=A/n$:
\[
\frac{d'}{d}=\frac{1}{\sqrt{n}}
\Rightarrow
n=3:\ 0{,}577;
\qquad
n=5:\ 0{,}447.
\]
\textbf{Aplicação prática --- o extensômetro.} Um fio colado a uma peça que se
deforma sofre exatamente esse efeito. Para pequenas deformações
$\varepsilon=\Delta L/L$, a expansão de $n^{2}$ dá
\[
\frac{\Delta R}{R}\approx2\varepsilon,
\]
ou seja, o fio funciona como sensor de deformação com \emph{fator de
sensibilidade} próximo de $2$. É o princípio das células de carga e das balanças
eletrônicas.}
\stepcounter{questao}
\resposta{$R=\dfrac{\rho L\ln2}{A_0}=\SI{11.8}{\milli\ohm}$}
\resolucao{(a) Somando as resistências de fatias infinitesimais em série:
\[
R=\int_0^{L}\frac{\rho\,\mathrm{d}x}{A(x)}
=\frac{\rho}{A_0}\int_0^{L}\frac{\mathrm{d}x}{1+x/L}
=\frac{\rho L}{A_0}\Big[\ln(1+x/L)\Big]_0^{L}
=\frac{\rho L\ln2}{A_0}.
\]
\[
R=\frac{\pot{1.7}{-8}\times1\times0{,}693}{\pot{1}{-6}}
=\SI{11.8}{\milli\ohm}.
\]
(b)
\begin{center}
\begin{tabular}{lc}
\hline
Modelo & $R$\\
\hline
uniforme com $A_0$ (menor seção) & $\SI{17.0}{\milli\ohm}$\\
\textbf{seção variável (exato)} & $\SI{11.8}{\milli\ohm}$\\
uniforme com $\bar A=1{,}5A_0$ & $\SI{11.3}{\milli\ohm}$\\
\hline
\end{tabular}
\end{center}
\textbf{O valor exato é maior que o da seção média} --- porque a resistência
depende de $1/A$, e a média de $1/A$ é maior que $1/\bar A$ (desigualdade das
médias). O erro de usar a seção média é de $4{,}3\%$, para menos.\\
\textbf{Regra:} em seções variáveis, integrar $\rho\,\mathrm{d}x/A(x)$. Usar a
seção média subestima; usar a seção mínima superestima. As duas aproximações
servem como limites.}
\stepcounter{questao}
\resposta{$U_0=\SI{100}{\milli\volt}$;
$\SI{0.385}{\milli\volt\per\celsius}$}
\resolucao{(a) $U_0=IR_0=\pot{1}{-3}(100)=\SI{100}{\milli\volt}$.
\[
\frac{\mathrm{d}U}{\mathrm{d}T}=I\,R_0\,\alpha
=\pot{1}{-3}(100)(0{,}00385)=\SI{0.385}{\milli\volt\per\celsius}.
\]
(b) \textbf{Vantagem decisiva: linearidade exata.} Com $I$ constante,
$U=IR$ é \emph{proporcional} a $R$, e como $R$ é linear em $T$, a tensão também
é. Não há distorção alguma.\\
Com \emph{tensão} constante em um divisor, a saída seria
$U\,R/(R+R_{fixo})$ --- uma função \textbf{não linear} de $R$, exigindo
linearização posterior.\\
(c) \textbf{Autoaquecimento:} $P=RI^{2}=100(\pot{1}{-6})=\SI{0.1}{\milli\watt}$.
Com resistência térmica típica de $\SI{100}{\celsius\per\watt}$ ao ar, a
elevação é de $\SI{0.01}{\celsius}$ --- desprezível.\\
\textbf{Mas cuidado:} elevar a corrente para $\SI{10}{\milli\ampere}$ (para
ganhar sinal) multiplicaria a potência por $100$, levando o erro a
$\SI{1}{\celsius}$. \textbf{A corrente de excitação é um compromisso entre
relação sinal-ruído e erro de autoaquecimento} --- e $\SI{1}{\milli\ampere}$ é
o valor de consenso para Pt100.}
\stepcounter{questao}
\resposta{Leitura de $\SI{0.5}{\ohm}$ --- erro de $400\%$}
\resolucao{(a) O instrumento mede tudo o que está em série:
\[
R_{lido}=0{,}1+2(0{,}2)=\SI{0.5}{\ohm}.
\]
(b) Erro de $\SI{0.4}{\ohm}$, ou $400\%$ --- a leitura é \emph{cinco vezes} o
valor real. Para resistências pequenas, o método de dois fios é simplesmente
inutilizável.\\
(c) \textbf{Método de quatro fios (Kelvin).} Usam-se dois pares independentes:
\begin{itemize}
\item Um par \textbf{injeta} a corrente conhecida $I$. A queda nesses cabos
existe, mas é irrelevante --- ela apenas exige um pouco mais de tensão da fonte.
\item O outro par \textbf{mede} a tensão diretamente sobre o resistor, com um
voltímetro de altíssima impedância. Como praticamente não circula corrente
nesses cabos, sua resistência \emph{não produz queda}.
\end{itemize}
Assim $R=U_{medido}/I$ devolve exatamente $\SI{0.1}{\ohm}$.\\
\textbf{Quando usar:} sempre que $R$ for comparável ou menor que a resistência
dos cabos --- tipicamente abaixo de alguns ohms. É o método obrigatório para
shunts, enrolamentos de motor e contatos.}
\stepcounter{questao}
\resposta{$\SI{5}{\volt}$/$\SI{50}{\milli\ampere}$ e
$\SI{22.4}{\volt}$/$\SI{224}{\milli\ampere}$; ambos são ôhmicos}
\resolucao{(a)
\[
U_{\max}=\sqrt{PR}:\quad \sqrt{0{,}25(100)}=\SI{5}{\volt};
\qquad \sqrt{5(100)}=\SI{22.4}{\volt};
\]
\[
I_{\max}=\sqrt{P/R}:\quad \SI{50}{\milli\ampere}
\quad\text{e}\quad \SI{224}{\milli\ampere}.
\]
(b) \textbf{Sim, ambos obedecem igualmente.} A lei de Ohm relaciona $U$ e $I$
através de $R$ --- e $R$ é o mesmo. A potência nominal \emph{não} aparece na
lei.\\
(c) \textbf{A potência nominal limita a faixa de validade}, não a relação em
si. Ela informa até onde o componente pode operar \emph{antes de se
autodestruir ou sair da linearidade por aquecimento}.\\
\textbf{Analogia útil:} a lei de Ohm é como a relação entre força e deformação
em uma mola; a potência nominal é o limite de escoamento do material. Duas
molas de mesma constante e materiais diferentes obedecem à mesma lei --- até que
uma delas se rompa.\\
\textbf{Diferença física:} o resistor de $\SI{5}{\watt}$ é maior, com mais área
de troca térmica. Mesma resistência, mesma lei, capacidade de dissipação
$20$ vezes maior.}
\stepcounter{questao}
\resposta{$\Delta T=\SI{100}{\celsius}$; $R=\SI{140}{\ohm}$}
\resolucao{(a)
\[
\Delta T=R_{\text{térm}}P=100(1)=\SI{100}{\celsius};
\qquad
R=100\left[1+0{,}004(100)\right]=\SI{140}{\ohm}.
\]
Aumento de $40\%$ apenas por autoaquecimento --- muito acima da tolerância
típica de fabricação.\\
(b) \textbf{Sob tensão constante:} $P=U^{2}/R$ \emph{diminui} quando $R$ sobe.
Menos potência, menos aquecimento, $R$ estabiliza. \textbf{Realimentação
negativa --- estável.}\\
\textbf{Sob corrente constante:} $P=RI^{2}$ \emph{aumenta} quando $R$ sobe.
Mais potência, mais aquecimento, $R$ cresce ainda mais.
\textbf{Realimentação positiva.}\\
\textbf{Condição de estabilidade} no caso de corrente constante: o incremento de
potência por grau, $\alpha R I^{2}$, deve ser menor que a capacidade de
escoar calor, $1/R_{\text{térm}}$:
\[
\alpha R I^{2}R_{\text{térm}}<1.
\]
Aqui, $0{,}004(100)(0{,}01)(100)=0{,}4<1$ \checkmark --- estável, mas com apenas
$2{,}5$ vezes de margem. Dobrar a corrente levaria o produto a $1{,}6$ e
provocaria fuga térmica.\\
\textbf{Materiais com $\alpha<0$} (termistores NTC, carbono) invertem os dois
casos --- e são instáveis justamente sob \emph{tensão} constante.}
\stepcounter{questao}
\resposta{Até $\SI{5}{\volt}$; usar $20\%$ da potência nominal como teto
prático}
\resolucao{(a) $U_{\max}=\sqrt{PR}=\SI{5}{\volt}$ no limite absoluto. Para
evitar autoaquecimento --- que é justamente o que falsearia o ensaio --- adote
$20\%$ da potência nominal:
\[
P_{ensaio}=\SI{50}{\milli\watt}
\;\Longrightarrow\;
U_{\max}=\sqrt{0{,}05(100)}=\SI{2.24}{\volt}.
\]
(b) \textbf{Pontos:} cinco valores igualmente espaçados entre $0{,}5$ e
$\SI{2.2}{\volt}$, medindo $U$ e $I$ em cada um.\\
\textbf{Critério:} calcular $R=U/I$ em todos os pontos. O componente é ôhmico se
a dispersão dos valores ficar dentro da incerteza combinada dos instrumentos
(tipicamente $1$ a $2\%$). Alternativamente, ajustar uma reta a $U\times I$ e
verificar se o intercepto é nulo dentro da incerteza --- intercepto não nulo
indica tensão de limiar, como em diodos.\\
(c) \textbf{Cuidados.}
\begin{itemize}
\item \textbf{Medir rápido} e desligar entre pontos, para o componente não
aquecer cumulativamente.
\item \textbf{Ordem aleatória} dos pontos, e não crescente --- assim um eventual
aquecimento progressivo aparece como dispersão, e não como tendência que
poderia ser confundida com não linearidade real.
\item \textbf{Método de quatro fios} se $R$ for pequeno.
\item \textbf{Repetir o primeiro ponto ao final:} se o valor mudou, houve
aquecimento e o ensaio precisa ser refeito com potência menor.
\end{itemize}}
\stepcounter{questao}
\resposta{$\Delta R/R\approx2\varepsilon$; fator $\approx2$}
\resolucao{(a) Com volume constante, $L'=L(1+\varepsilon)$ e
$A'=A/(1+\varepsilon)$:
\[
R'=\frac{\rho L(1+\varepsilon)}{A/(1+\varepsilon)}
=R(1+\varepsilon)^{2}\approx R(1+2\varepsilon),
\]
para $\varepsilon\ll1$. Logo
\[
\frac{\Delta R}{R}\approx2\varepsilon.
\]
(b) O \textbf{fator de sensibilidade} (\emph{gauge factor}) é
$k=\dfrac{\Delta R/R}{\varepsilon}\approx2$.\\
\textbf{Ordem de grandeza prática:} uma deformação de
$\SI{1000}{\micro\meter\per\meter}$ ($\varepsilon=\pot{1}{-3}$) produz
$\Delta R/R=0{,}2\%$ --- em um extensômetro de $\SI{350}{\ohm}$, apenas
$\SI{0.7}{\ohm}$. Daí a necessidade de ponte de Wheatstone e amplificação de
instrumentação.\\
(c) \textbf{Efeito piezorresistivo.} Em metais, $\rho$ varia pouco com a
deformação e $k\approx2$ decorre quase só da geometria. Em
\textbf{semicondutores}, a deformação altera fortemente a mobilidade dos
portadores, e $k$ chega a $100$ ou mais --- ganho enorme de sensibilidade.\\
\textbf{Contrapartida:} extensômetros semicondutores são muito mais sensíveis à
temperatura e menos lineares. \textbf{A escolha entre metal e semicondutor é um
compromisso entre sensibilidade e estabilidade} --- e por isso células de carga
de precisão ainda usam extensômetros metálicos.}
\stepcounter{questao}
\resposta{$R=\dfrac{\rho L}{A_0}\cdot\dfrac{\ln(1+k)}{k}$; para $k=1$,
$\ln2$}
\resolucao{(a) Fatias de espessura $\mathrm{d}x$ estão em série:
\[
R=\int_0^{L}\frac{\rho\,\mathrm{d}x}{A_0(1+x/L)}
=\frac{\rho L}{A_0}\ln 2.
\]
(b) Com $A(x)=A_0(1+kx/L)$:
\[
R=\frac{\rho}{A_0}\int_0^{L}\frac{\mathrm{d}x}{1+kx/L}
=\frac{\rho L}{A_0k}\Big[\ln(1+kx/L)\Big]_0^{L}
=\frac{\rho L}{A_0}\cdot\frac{\ln(1+k)}{k}.
\]
(c) \textbf{Limite $k\to0$} (seção uniforme). Usando
$\ln(1+k)\approx k-k^{2}/2$:
\[
\frac{\ln(1+k)}{k}\to1
\;\Longrightarrow\;
R\to\frac{\rho L}{A_0}.\ \checkmark
\]
Recupera-se a fórmula elementar --- teste obrigatório de qualquer generalização.\\
\textbf{Limite $k\to\infty$} (seção crescendo muito rápido):
\[
\frac{\ln(1+k)}{k}\to0
\;\Longrightarrow\;
R\to0.
\]
A barra fica tão larga na maior parte de seu comprimento que sua resistência se
concentra apenas na região inicial estreita --- e o total tende a zero, embora
\emph{logaritmicamente devagar}. Mesmo com $k=1000$, o fator ainda vale
$0{,}0069$, e não zero.\\
\textbf{Leitura geral:} a resistência é dominada pela região de \emph{menor}
seção. É por isso que um estrangulamento local em um condutor concentra
praticamente toda a queda de tensão --- e todo o aquecimento.}
\stepcounter{questao}
\resposta{$I=\SI{0.222}{\ampere}$, $U=\SI{22.2}{\volt}$;
$P=\SI{4.94}{\watt}$, $\eta=92{,}6\%$}
\resolucao{(a) A fonte impõe $U=24-8I$ e o resistor impõe $U=100I$. Igualando:
\[
100I=24-8I
\;\Longrightarrow\;
108I=24
\;\Longrightarrow\;
I=\SI{0.222}{\ampere};
\qquad
U=\SI{22.2}{\volt}.
\]
(b) $P_L=UI=22{,}2(0{,}222)=\SI{4.94}{\watt}$.\\
Rendimento: $\eta=\dfrac{R_L}{R_L+R_{Th}}=\dfrac{100}{108}=92{,}6\%$.\\
(c) \textbf{Com $\SI{8}{\ohm}$:} $8I=24-8I\Rightarrow I=\SI{1.5}{\ampere}$ e
$U=\SI{12}{\volt}$.
\[
P_L=12(1{,}5)=\SI{18}{\watt};
\qquad
\eta=\frac{8}{16}=50\%.
\]
\textbf{Comparação reveladora.} O resistor de $\SI{8}{\ohm}$ está \emph{casado}
com a fonte e recebe a potência máxima ($\SI{18}{\watt}$, quase quatro vezes
mais), mas desperdiça metade da energia internamente. O de
$\SI{100}{\ohm}$ recebe pouco, porém com rendimento de $92{,}6\%$.\\
\textbf{Ponto de método:} o problema é a interseção de duas características
$U\times I$ --- uma reta de fonte e uma reta de carga. Se a carga fosse não
linear (um diodo, por exemplo), o procedimento seria idêntico: basta trocar a
segunda reta pela curva do dispositivo.}
\stepcounter{questao}
\resposta{Condição: $\alpha R I^{2}R_{\text{térm}}<1$; aqui $0{,}4$ --- estável;
$I_c=\SI{158}{\milli\ampere}$}
\resolucao{(a) Em equilíbrio, $\Delta T=R_{\text{térm}}P=R_{\text{térm}}RI^{2}$ e
$R=R_0(1+\alpha\Delta T)$. Substituindo,
\[
R=R_0\left(1+\alpha R_{\text{térm}}RI^{2}\right)
\;\Longrightarrow\;
R\left(1-\alpha R_{\text{térm}}R_0I^{2}\right)=R_0.
\]
Existe solução finita e estável se, e somente se, o parênteses for positivo:
\[
\boxed{\alpha\,R_0\,I^{2}\,R_{\text{térm}}<1}
\]
Fisicamente: o acréscimo de potência por grau de aquecimento deve ser
\emph{menor} que a capacidade de escoar calor. Caso contrário, cada grau
adicional gera mais calor do que consegue dissipar --- \textbf{fuga térmica}.\\
(b)
\[
0{,}004\times100\times(0{,}1)^{2}\times100=0{,}4<1\ \checkmark
\]
Estável, com margem de $2{,}5$ vezes. A resistência de equilíbrio é
$R=100/(1-0{,}4)=\SI{167}{\ohm}$, e a elevação, $\SI{167}{\celsius}$ ---
elevada, mas estável.\\
(c) Corrente crítica:
\[
I_c=\sqrt{\frac{1}{\alpha R_0R_{\text{térm}}}}=\sqrt{\frac{1}{0{,}004(100)(100)}}
=\sqrt{0{,}025}=\SI{158}{\milli\ampere}.
\]
\textbf{Margem estreita:} apenas $58\%$ acima da corrente de operação. E note
que a estabilidade \emph{matemática} não garante a \emph{sobrevivência} --- a
$\SI{167}{\celsius}$ o componente já pode estar além de sua temperatura máxima.\\
\textbf{Sob tensão constante} o problema não existe: $P=U^{2}/R$ decresce com
$R$, e a realimentação é sempre negativa.}
\stepcounter{questao}
\resposta{$\SI{3.87}{\milli\meter\squared}$ e
$\SI{15.5}{\milli\meter\squared}$; a corrente domina no cabo curto, a queda no
longo}
\resolucao{(a) Queda admissível: $0{,}04(220)=\SI{8.8}{\volt}$, logo
$R_{\max}=8{,}8/30=\SI{0.293}{\ohm}$.
\[
\SI{20}{\meter}:\ A\ge\frac{\pot{1.7}{-8}(40)}{0{,}293}
=\SI{2.32}{\milli\meter\squared};
\]
\[
\SI{80}{\meter}:\ A\ge\frac{\pot{1.7}{-8}(160)}{0{,}293}
=\SI{9.28}{\milli\meter\squared}.
\]
(lembrando o fator $2$ de ida e volta).\\
(b) \textbf{Cabo de $\SI{20}{\meter}$:} a queda exige
$\SI{2.32}{\milli\meter\squared}$, mas a corrente de $\SI{30}{\ampere}$ exige
pelo menos $\SI{6}{\milli\meter\squared}$. \textbf{O critério de corrente
domina.}\\
\textbf{Cabo de $\SI{80}{\meter}$:} a queda exige
$\SI{9.28}{\milli\meter\squared}$, acima do que a corrente pediria.
\textbf{O critério de queda domina}, e adota-se
$\SI{10}{\milli\meter\squared}$.\\
(c) \textbf{Regra geral.} Os dois critérios escalam de formas diferentes:
\begin{itemize}
\item \textbf{Corrente admissível:} depende só de $I$ --- \emph{independe do
comprimento}. É um critério de segurança térmica do próprio cabo.
\item \textbf{Queda de tensão:} $A\propto IL/\Delta U$ --- cresce
\emph{linearmente} com o comprimento. É um critério de qualidade da energia
entregue.
\end{itemize}
Existe um \textbf{comprimento de cruzamento} em que os dois coincidem; abaixo
dele manda a corrente, acima manda a queda. Aqui, com
$\SI{6}{\milli\meter\squared}$, o cruzamento ocorre por volta de
$\SI{52}{\meter}$.\\
\textbf{Procedimento correto:} calcular os \emph{dois} e adotar a maior seção.
Verificar apenas um deles é a origem mais comum de instalações que "funcionam"
mas apresentam tensão baixa na ponta.}
\stepcounter{questao}
\resposta{Corrente constante: $U=IR$, exatamente linear; divisor:
$U=U_0R/(R+R_f)$, não linear}
\resolucao{(a) \textbf{Corrente constante:} $U=IR$ --- proporcional a $R$.\\
\textbf{Divisor:} $U=U_0\dfrac{R}{R+R_f}$ --- função racional de $R$.\\
(b) \textbf{Linearidade.} A primeira é \emph{exatamente} linear, sempre. A
segunda só é aproximadamente linear para pequenas variações; expandindo em torno
de $R=R_f$:
\[
U\approx\frac{U_0}{2}\left[1+\frac{\Delta R}{2R_f}
-\frac{1}{4}\left(\frac{\Delta R}{R_f}\right)^{2}+\cdots\right].
\]
O termo quadrático produz erro de linearidade. Para
$\Delta R/R_f=10\%$, ele vale $0{,}25\%$ da leitura --- pequeno, mas suficiente
para limitar a exatidão em medições finas.\\
(c) \textbf{Sensibilidades.}
\[
\text{corrente constante:}\quad \frac{\mathrm{d}U}{\mathrm{d}R}=I;
\qquad
\text{divisor:}\quad \frac{\mathrm{d}U}{\mathrm{d}R}
=\frac{U_0R_f}{(R+R_f)^{2}}\bigg|_{R=R_f}=\frac{U_0}{4R_f}.
\]
Ambas podem ser ajustadas ao mesmo valor escolhendo $I=U_0/(4R_f)$ --- a
sensibilidade não distingue os métodos no ponto de operação.\\
\textbf{O que decide é o resto.}
\begin{itemize}
\item \textbf{Corrente constante:} linearidade exata, mas exige fonte de
corrente estável e o sinal parte de zero apenas se $R\to0$.
\item \textbf{Divisor:} circuito trivial e barato, e a saída pode ser centrada
em meia alimentação --- conveniente para conversores A/D. Em contrapartida,
carrega o erro de linearidade e a deriva de $R_f$ e de $U_0$.
\end{itemize}
\textbf{Solução de compromisso muito usada:} a \emph{ponte} de Wheatstone, que
cancela a componente comum e entrega diretamente o desequilíbrio --- combinando
boa linearidade com circuito passivo.}
\stepcounter{questao}
\resposta{A hipótese oculta é dimensões constantes; o erro é da ordem de
$\SI{0.05}{\percent}$ --- desprezível}
\resolucao{(a) De $R=\rho L/A$, se $L$ e $A$ forem constantes, então
$R\propto\rho$ e a lei se transfere diretamente. \textbf{A hipótese oculta é
justamente essa}: que a geometria não muda com a temperatura --- o que é falso,
pois todo metal se dilata.\\
(b) Com coeficiente de dilatação linear $\lambda$:
$L'=L(1+\lambda\Delta T)$ e $A'=A(1+\lambda\Delta T)^{2}$, logo
\[
\frac{L'}{A'}=\frac{L}{A}\cdot\frac{1}{1+\lambda\Delta T}.
\]
Assim
\[
R=\frac{\rho_0(1+\alpha\Delta T)}{1+\lambda\Delta T}\cdot\frac{L}{A}
\approx R_0\left[1+(\alpha-\lambda)\Delta T\right].
\]
\textbf{O efeito da dilatação é subtrair $\lambda$ de $\alpha$} --- a peça fica
mais longa (aumentando $R$) mas também mais grossa (reduzindo $R$), e o segundo
efeito predomina por ser quadrático na área.\\
(c) Para o cobre, $\alpha\approx\SI{3.9e-3}{\per\celsius}$ e
$\lambda\approx\SI{1.7e-5}{\per\celsius}$:
\[
\frac{\lambda}{\alpha}=\frac{1{,}7\times10^{-5}}{3{,}9\times10^{-3}}
=0{,}0044=0{,}44\%.
\]
Com $\Delta T=\SI{100}{\celsius}$, $R$ aumenta $39\%$ pela resistividade e
diminui $0{,}17\%$ pela dilatação --- efeito líquido de $38{,}8\%$.\\
\textbf{Conclusão:} ignorar a dilatação introduz erro de cerca de
$0{,}4\%$ \emph{do efeito térmico}, ou $0{,}17\%$ do valor de $R$. Desprezível
em quase toda aplicação --- mas não em padrões de resistência de laboratório,
onde se buscam partes por milhão.}
\stepcounter{questao}
\resposta{Erro $=2R_{cabo}/R_x$; é preciso $R_x\ge\SI{40}{\ohm}$}
\resolucao{(a) O instrumento lê $R_x+2R_{cabo}$, logo
\[
\varepsilon=\frac{2R_{cabo}}{R_x}.
\]
(b) Impondo $\varepsilon\le0{,}01$ com $R_{cabo}=\SI{0.2}{\ohm}$:
\[
R_x\ge\frac{2(0{,}2)}{0{,}01}=\SI{40}{\ohm}.
\]
\textbf{Regra prática:} o método de dois fios só é aceitável para resistências
de dezenas de ohms ou mais. Abaixo disso, o erro cresce rapidamente --- em
$\SI{1}{\ohm}$ ele já é de $40\%$.\\
(c) \textbf{O que o método de quatro fios \emph{não} elimina.}
\begin{itemize}
\item \textbf{Tensões termoelétricas.} Junções entre metais diferentes nos
contatos geram forças eletromotrizes de dezenas de microvolts --- comparáveis ao
sinal quando se mede miliohms. Combate-se invertendo o sentido da corrente e
tomando a média das duas leituras.
\item \textbf{Autoaquecimento} pela corrente de excitação, que altera o próprio
$R_x$.
\item \textbf{Localização dos contatos de tensão.} Eles devem estar
\emph{entre} os de corrente, e a distribuição de corrente na peça influencia o
que se mede. Em amostras curtas, a definição do que é "o resistor" torna-se
ambígua.
\item \textbf{Ruído e deriva} do amplificador, agravados porque o sinal é de
poucos milivolts.
\end{itemize}
\textbf{Síntese:} o método de quatro fios elimina o erro \emph{sistemático e
dominante} --- a resistência dos cabos. Os erros remanescentes são de outra
natureza e exigem outras técnicas.}
\stepcounter{questao}
\resposta{$R_{\text{aço}}=\SI{7.5}{\ohm}$, $R_{Al}=\SI{0.28}{\ohm}$,
$R_{eq}=\SI{0.270}{\ohm}$; o aço conduz $3{,}6\%$}
\resolucao{(a) Os dois materiais estão em \textbf{paralelo} (mesmo comprimento,
mesmos terminais):
\[
R_{\text{aço}}=\frac{\pot{1.5}{-7}(1000)}{\pot{20}{-6}}=\SI{7.5}{\ohm};
\qquad
R_{Al}=\frac{\pot{2.8}{-8}(1000)}{\pot{100}{-6}}=\SI{0.28}{\ohm};
\]
\[
R_{eq}=\frac{7{,}5(0{,}28)}{7{,}78}=\SI{0.270}{\ohm}.
\]
(b) Pelo divisor de corrente em condutâncias:
\[
\frac{I_{Al}}{I_T}=\frac{1/0{,}28}{1/0{,}28+1/7{,}5}=96{,}4\%;
\qquad
\frac{I_{\text{aço}}}{I_T}=3{,}6\%.
\]
(c) \textbf{A alma de aço praticamente não conduz --- e não é essa a sua
função.} Ela existe por razões \textbf{mecânicas}:
\begin{itemize}
\item O alumínio tem resistência à tração baixa e escoa sob carga permanente. O
aço sustenta o peso próprio do cabo entre torres, que podem distar centenas de
metros.
\item Reduz a flecha (o "barrigamento") do vão e, com ela, a altura necessária
das torres.
\item Suporta os esforços de vento e de acúmulo de gelo.
\end{itemize}
\textbf{Observe a divisão de papéis:} o alumínio conduz $96\%$ da corrente e o
aço sustenta praticamente toda a carga mecânica. É o princípio do cabo
\textbf{ACSR} (\emph{aluminium conductor steel-reinforced}), padrão em linhas de
transmissão aéreas --- um caso em que a análise elétrica revela que o componente
"elétrico" está ali por outro motivo.}
\stepcounter{questao}
\resposta{Autoaquecimento, não linearidade intrínseca, efeitos de alta
frequência e ruptura do meio}
\resolucao{(a) \textbf{Quatro mecanismos.}
\begin{enumerate}
\item \textbf{Autoaquecimento.} A resistência depende da temperatura, que
depende da potência dissipada. Exemplo: lâmpada incandescente, cuja resistência
a frio é cerca de dez vezes menor que a quente.
\item \textbf{Não linearidade intrínseca.} Junções semicondutoras seguem
relação exponencial, não linear. Exemplo: diodo, cuja "resistência" varia
ordens de grandeza em uma década de tensão.
\item \textbf{Efeitos de frequência.} O efeito pelicular concentra a corrente na
superfície do condutor, aumentando a resistência efetiva; somam-se indutância e
capacitância parasitas. Exemplo: um fio que tem $\SI{1}{\milli\ohm}$ em CC pode
apresentar dezenas de miliohms em megahertz.
\item \textbf{Ruptura do meio.} Acima da rigidez dielétrica, um isolante
torna-se condutor abruptamente. Exemplo: arco elétrico, cuja característica
$U\times I$ tem inclinação \emph{negativa}.
\end{enumerate}
\emph{(Acrescente-se ainda a saturação de velocidade dos portadores em campos
muito intensos e a supercondutividade, em que $R$ colapsa a zero.)}\\
(b) \textbf{A lei de Ohm é uma lei empírica e aproximada}, não fundamental. Ela
descreve o comportamento de \emph{certos} materiais (metais, em faixa limitada
de temperatura e frequência) e emerge de um modelo estatístico de colisões dos
portadores --- não de um princípio de conservação.\\
Compare com as leis de Kirchhoff, que decorrem da conservação de carga e de
energia e valem \emph{sempre}, para qualquer componente, linear ou não.\\
(c) \textbf{Quando a lei falha, procede-se assim:}
\begin{itemize}
\item Substituir $R$ pela \textbf{característica $U\times I$} completa do
dispositivo, obtida por medição ou por modelo.
\item Resolver o circuito por \textbf{interseção de características} (reta de
carga), como se faz com diodos.
\item Para pequenos sinais, linearizar em torno do ponto de operação usando a
\textbf{resistência dinâmica} --- e aí toda a maquinaria linear volta a valer.
\end{itemize}
\textbf{As leis de Kirchhoff continuam válidas em todos esses casos} --- é
sempre por elas que se começa.}

\gabarito[2.2 Lei de Ohm e resistividade]

\subsection{Potência elétrica e energia}
\stepcounter{questao}
\resposta{(a) \SI{2}{\ampere}; (b) \SI{6}{\ohm}}
\resolucao{(a) De $P = U\,I$: $I = \dfrac{P}{U} = \dfrac{24}{12} = \SI{2}{\ampere}$.\\
(b) $R = \dfrac{U}{I} = \dfrac{12}{2} = \SI{6}{\ohm}$
(ou $R = \dfrac{U^2}{P} = \dfrac{144}{24} = \SI{6}{\ohm}$).}
\stepcounter{questao}
\resposta{(a) \SI{2}{\ampere}; (b) \SI{110}{\ohm}}
\resolucao{(a) $I = \dfrac{P}{U} = \dfrac{440}{220} = \SI{2}{\ampere}$.\\
(b) $R = \dfrac{U^2}{P} = \dfrac{220^2}{440} = \SI{110}{\ohm}$.}
\stepcounter{questao}
\resposta{$\eta=95\%$; perdas de $\SI{5}{\watt}$}
\resolucao{$\eta=P_{\text{útil}}/P_{\mathrm{entrada}}=95/100=0{,}95$.
A diferença, $100-95=\SI{5}{\watt}$, é dissipada em fios, contatos e outros elementos.}
\stepcounter{questao}
\resposta{(d)}
\resolucao{A resistência da lâmpada é fixa, determinada por seus valores nominais:
$R = \dfrac{U_{\text{nom}}^2}{P_{\text{nom}}} = \dfrac{220^2}{60} \approx
\SI{807}{\ohm}$. Sob a nova tensão, $P' = \dfrac{U'^2}{R}$. Como a tensão caiu à
metade e $P \propto U^2$:
\[
P' = \frac{(110)^2}{R} = \frac{60}{4} = \SI{15}{\watt},
\]
ou seja, a quarta parte. Meia tensão dissipa um quarto da potência.}
\stepcounter{questao}
\resposta{(b)}
\resolucao{Mesma ideia da questão anterior: metade da tensão dissipa um quarto da
potência:
\[
P' = \frac{4400}{4} = \SI{1100}{\watt}.
\]
Por isso um chuveiro projetado para \SI{220}{\volt} "esquenta pouco" em rede de
\SI{110}{\volt}.}
\stepcounter{questao}
\resposta{(a) \SI{2}{\ampere}; (b) $P_1=\SI{40}{\watt}$, $P_2=\SI{80}{\watt}$}
\resolucao{(a) $I = \dfrac{60}{10+20} = \SI{2}{\ampere}$.\\
(b) Com a mesma corrente, $P = R\,I^2$:
\[
P_1 = 10\cdot 2^2 = \SI{40}{\watt},\qquad
P_2 = 20\cdot 2^2 = \SI{80}{\watt}.
\]
\textbf{Verificação:} $P_{\text{total}} = U\,I = 60\cdot2 = \SI{120}{\watt}
= 40+80$. Na série, o maior resistor dissipa a maior potência.}
\stepcounter{questao}
\resposta{(a) \SI{25}{\ampere}; (b) \SI{8.8}{\ohm}; (c) não}
\resolucao{(a) $I = \dfrac{P}{U} = \dfrac{5500}{220} = \SI{25}{\ampere}$.\\
(b) $R = \dfrac{U^2}{P} = \dfrac{220^2}{5500} = \SI{8.8}{\ohm}$.\\
(c) A corrente exigida (\SI{25}{\ampere}) \emph{excede} os \SI{20}{\ampere} do
disjuntor, que desarmaria. Seria necessário um disjuntor de pelo menos
\SI{25}{\ampere} (na prática, \SI{32}{\ampere}) e fiação compatível --- por isso
chuveiros exigem circuito exclusivo.}
\stepcounter{questao}
\resposta{$\SI{45}{\kilo\watt\hour}$; R\$~18,00}
\resolucao{A tarifa é $80/200=\text{R\$ }0{,}40/\si{\kilo\watt\hour}$.
O tempo mensal é $0{,}5\times30=\SI{15}{\hour}$, então
$E=3\times15=\SI{45}{\kilo\watt\hour}$ e o custo é $45\times0{,}40=\text{R\$ }18{,}00$.}
\stepcounter{questao}
\resposta{$I = \SI{13}{\ampere}$; $E = \SI{3.12}{\kilo\watt\hour}$}
\resolucao{Potência total: $P = 1440 + 2\cdot60 = \SI{1560}{\watt}$.\\
Corrente no fusível (todos em paralelo, sob \SI{120}{\volt}):
\[
I = \frac{P}{U} = \frac{1560}{120} = \SI{13}{\ampere}.
\]
Energia em $\SI{2}{\hour}$:
\[
E = P\,t = 1{,}560\,\si{\kilo\watt}\cdot2\,\si{\hour} = \SI{3.12}{\kilo\watt\hour}.
\]
O kWh é a unidade da conta de luz: potência (em kW) vezes tempo (em h).}
\stepcounter{questao}
\resposta{$E \approx \SI{127.5}{\kilo\watt\hour}$; custo $\approx$ R\$\,102,00}
\resolucao{Energia diária de cada aparelho ($E = P\,t$, em kWh):
\[
\text{lâmpada: } 0{,}1\cdot5 = \SI{0.5}{\kilo\watt\hour};\quad
\text{chuveiro: } 4{,}4\cdot0{,}5 = \SI{2.2}{\kilo\watt\hour};\quad
\text{geladeira: } 0{,}2\cdot8 = \SI{1.6}{\kilo\watt\hour}.
\]
Total diário: $0{,}5 + 2{,}2 + 1{,}6 = \SI{4.3}{\kilo\watt\hour}$. Em 30 dias:
\[
E = 4{,}3\cdot30 = \SI{129}{\kilo\watt\hour}
\quad\Rightarrow\quad
\text{custo} = 129\cdot0{,}80 \approx \text{R\$}\,103{,}00.
\]
(Valores próximos de \SI{127.5}{\kilo\watt\hour} conforme arredondamentos.) É
assim que se estima a conta de luz --- e fica claro que o chuveiro, apesar de
usado poucos minutos, é um dos maiores vilões pelo seu alto \emph{potência}.}
\stepcounter{questao}
\resposta{(a) $\SI{806.7}{\ohm}$ e $\SI{484}{\ohm}$;
(b) $\SI{23.4}{\watt}$ e $\SI{14.1}{\watt}$; (c) a de \SI{60}{\watt}}
\resolucao{\textbf{(a)} Das especificações nominais, $R = \dfrac{U^2}{P}$:
\[
R_{60}=\frac{220^2}{60}\approx\SI{806.7}{\ohm},
\qquad
R_{100}=\frac{220^2}{100}=\SI{484}{\ohm}.
\]
Note: \emph{a lâmpada de menor potência tem MAIOR resistência.}

\textbf{(b)} Em série, a corrente é comum:
\[
I=\frac{220}{806{,}7+484}=\frac{220}{1290{,}7}\approx\SI{0.170}{\ampere}.
\]
Com $P = R I^2$:
\[
P_{60}=806{,}7\,(0{,}170)^2\approx\SI{23.4}{\watt},
\qquad
P_{100}=484\,(0{,}170)^2\approx\SI{14.1}{\watt}.
\]
\textbf{Verificação:} $23{,}4+14{,}1 = \SI{37.5}{\watt} = 220\cdot0{,}170$.
\checkmark

\textbf{(c) A lâmpada de \SI{60}{\watt} brilha MAIS} --- resultado que contraria
a intuição de quem associa "100 W" a "mais brilhante".\\
\textbf{Por quê:} as especificações valem \emph{apenas} na tensão nominal, com
cada lâmpada recebendo os \SI{220}{\volt} completos. Em série, elas
\textbf{compartilham} a mesma corrente, e quem dita a potência passa a ser a
resistência: $P = RI^2$. A de \SI{60}{\watt}, tendo maior $R$, fica com maior
parcela da tensão e dissipa mais.\\
\textbf{Regra geral:}
\begin{center}\small\sffamily
\begin{tabular}{@{}l l l@{}}
\toprule
\textbf{Ligação} & \textbf{Grandeza comum} & \textbf{Brilha mais}\\
\midrule
Série    & corrente $I$ & a de \emph{maior} $R$ (menor potência nominal)\\
Paralelo & tensão $U$   & a de \emph{menor} $R$ (maior potência nominal)\\
\bottomrule
\end{tabular}
\end{center}
Observe ainda que \emph{ambas} ficam muito abaixo do nominal (\SI{37.5}{\watt} no
total contra \SI{160}{\watt} se ligadas corretamente em paralelo).}
\stepcounter{questao}
\resposta{(a) \SI{5500}{\watt} e \SI{11000}{\watt}; (b) \SI{25}{\ampere} e
\SI{50}{\ampere} --- disjuntor insuficiente; (c) $\approx\SI{72}{\celsius}$}
\resolucao{\textbf{(a)} Com $P = \dfrac{U^2}{R}$:
\[
P_{\text{verão}}=\frac{220^2}{8{,}8}=\SI{5500}{\watt},
\qquad
P_{\text{inverno}}=\frac{220^2}{4{,}4}=\SI{11000}{\watt}.
\]
\emph{Metade da resistência dobra a potência} --- e não a reduz, como a intuição
poderia sugerir.

\textbf{(b)} $I = \dfrac{P}{U}$:
$I_{\text{verão}} = \SI{25}{\ampere}$ e $I_{\text{inverno}} = \SI{50}{\ampere}$.
O disjuntor de \SI{40}{\ampere} \textbf{não é adequado} --- ele desarmaria na
posição inverno. Seria necessário um de \SI{63}{\ampere} com fiação
de $\SI{10}{\milli\meter\squared}$.

\textbf{(c) Balanço térmico.} A vazão mássica é
\[
\dot m = \SI{3}{\liter\per\minute} = \frac{3}{60}=\SI{0.05}{\kilogram\per\second}.
\]
Toda a potência elétrica vira calor na água:
\[
P = \dot m\,c\,\Delta T
\;\Rightarrow\;
\Delta T=\frac{P}{\dot m\,c}=\frac{11000}{0{,}05\cdot4180}\approx\SI{52.6}{\celsius}.
\]
Logo a saída fica a $20+52{,}6 \approx \SI{72}{\celsius}$ --- alta demais para
banho, o que na prática é ajustado aumentando a vazão.\\
\textbf{Síntese de engenharia:} este exercício conecta \emph{três} domínios ---
circuito elétrico (lei de Ohm), proteção (dimensionamento de disjuntor) e
termodinâmica (balanço de energia). É exatamente esse tipo de integração que
caracteriza o trabalho profissional real.}
\stepcounter{questao}
\resposta{$P=\SI{50}{\watt}$}
\resolucao{
\[
P=RI^{2}=8(2{,}5)^{2}=8(6{,}25)=\SI{50}{\watt}.
\]
\textbf{Escolha da fórmula:} conhecendo $R$ e $I$, use $RI^{2}$ diretamente.
Calcular $U=RI=\SI{20}{\volt}$ e depois $P=UI$ dá o mesmo resultado, mas
acrescenta um passo e uma oportunidade de erro.}
\stepcounter{questao}
\resposta{$P=\SI{1100}{\watt}$}
\resolucao{
\[
P=\frac{U^{2}}{R}=\frac{220^{2}}{44}=\frac{48400}{44}=\SI{1100}{\watt}.
\]
Note a dependência \emph{quadrática} na tensão: uma queda de $10\%$ na rede
($\SI{198}{\volt}$) reduziria a potência para
$\SI{891}{\watt}$ --- $19\%$ a menos. É por isso que chuveiros "esquentam
menos" em horários de pico.}
\stepcounter{questao}
\resposta{$\SI{900}{\kilo\joule}$}
\resolucao{
\[
E=Pt=250\,\si{\watt}\times3600\,\si{\second}=\SI{900000}{\joule}
=\SI{900}{\kilo\joule}.
\]
\textbf{Fator de conversão útil:}
$\SI{1}{\kilo\watt\hour}=\SI{3.6}{\mega\joule}$. O joule é pequeno demais para
contas de energia elétrica --- daí o uso do quilowatt-hora, que é uma unidade
de energia, não de potência. Confundir os dois é o erro mais comum do tópico.}
\stepcounter{questao}
\resposta{(c)}
\resolucao{As alternativas (a) e (b) exigiriam calcular $I$ antes. A (d) está
dimensionalmente errada: $\si{\volt}/\si{\ohm}$ é ampère, não watt.\\
\textbf{Verificação dimensional rápida:} $U^{2}/R$ dá
$\si{\volt\squared}/\si{\ohm}=\si{\volt}\cdot\si{\ampere}=\si{\watt}$
\checkmark. Sempre que houver dúvida entre fórmulas parecidas, confira a
dimensão --- ela elimina a maioria das alternativas erradas.}
\stepcounter{questao}
\resposta{$E=\SI{18}{\kilo\joule}$}
\resolucao{Convertendo o tempo para segundos:
\[
E=Pt=60\times(5\times60)=60(300)=\SI{18000}{\joule}=\SI{18}{\kilo\joule}.
\]
\textbf{Cuidado com as unidades:} $P$ em watts exige $t$ em \emph{segundos}
para obter joules. Se o tempo ficar em minutos, o resultado sai
$60\times$ menor --- erro silencioso, porque o número parece plausível.}
\stepcounter{questao}
\resposta{$I=\SI{25}{\ampere}$; $R=\SI{8.8}{\ohm}$}
\resolucao{
\[
I=\frac{P}{U}=\frac{5500}{220}=\SI{25}{\ampere};
\qquad
R=\frac{U^{2}}{P}=\frac{48400}{5500}=\SI{8.8}{\ohm}.
\]
\textbf{Conferência:} $R=U/I=220/25=\SI{8.8}{\ohm}$ \checkmark.\\
\textbf{Implicação de instalação:} $\SI{25}{\ampere}$ é uma corrente elevada.
Ela exige circuito exclusivo, condutor de seção adequada e disjuntor de
$\SI{32}{\ampere}$ --- e é por isso que chuveiro nunca compartilha circuito com
tomadas.}
\stepcounter{questao}
\resposta{$R=\SI{40.3}{\ohm}$; $I=\SI{5.45}{\ampere}$; disjuntor de
$\SI{10}{\ampere}$}
\resolucao{
\[
R=\frac{U^{2}}{P}=\frac{48400}{1200}=\SI{40.3}{\ohm};
\qquad
I=\frac{1200}{220}=\SI{5.45}{\ampere}.
\]
\textbf{Disjuntor:} o valor comercial imediatamente acima é
$\SI{10}{\ampere}$, que oferece folga de $83\%$ --- adequado, pois aquecedores
resistivos não têm corrente de partida elevada.\\
\textbf{Nota sobre o material:} $\SI{40.3}{\ohm}$ deve ser obtido com fio de
liga resistiva (níquel-cromo) de comprimento e seção calculados. E como a
resistência desses materiais varia pouco com a temperatura, a potência
permanece próxima do nominal em regime.}
\stepcounter{questao}
\resposta{$\SI{268.8}{\ohm}$ e $\SI{806.7}{\ohm}$}
\resolucao{
\[
R_{127}=\frac{127^{2}}{60}=\SI{268.8}{\ohm};
\qquad
R_{220}=\frac{220^{2}}{60}=\SI{806.7}{\ohm}.
\]
A lâmpada de tensão maior tem resistência \textbf{três vezes} maior --- razão
$(220/127)^{2}=3{,}0$.\\
\textbf{Consequência construtiva:} para a mesma potência, o filamento de
$\SI{220}{\volt}$ é mais longo e mais fino. Ele é, portanto, mais frágil
mecanicamente --- fato que se observa na prática, com lâmpadas de tensão maior
tendendo a durar menos sob vibração.}
\stepcounter{questao}
\resposta{$\SI{27}{\kilo\watt\hour}$; $\mathrm{R\$}\,25{,}65$}
\resolucao{Energia diária, com os tempos em horas:
\[
E=1{,}5\left(\tfrac13\right)+0{,}6\left(\tfrac23\right)
=0{,}50+0{,}40=\SI{0.90}{\kilo\watt\hour}.
\]
\[
E_{\text{mês}}=0{,}90(30)=\SI{27}{\kilo\watt\hour};
\qquad
\text{custo}=27(0{,}95)=\mathrm{R\$}\,25{,}65.
\]
\textbf{Observação:} o modo de $\SI{600}{\watt}$, embora tenha menos da metade
da potência, responde por $44\%$ do consumo --- porque funciona o dobro do
tempo. Em energia, \textbf{potência sozinha não decide}; o produto pelo tempo é
que conta.}
\stepcounter{questao}
\resposta{$P_{\text{média}}=\SI{5}{\watt}$}
\resolucao{
\[
P_{\text{média}}=P_{pico}\times D=50(0{,}10)=\SI{5}{\watt}.
\]
\textbf{Mas cuidado:} isso \emph{não} significa que um resistor de
$\SI{5}{\watt}$ sirva. Se a duração do pulso for comparável ou maior que a
constante térmica do componente, ele atinge temperaturas próximas às de
$\SI{50}{\watt}$ contínuos e pode falhar.\\
A questão de quando a potência média basta é retomada quantitativamente no bloco
de aprofundamento.}
\stepcounter{questao}
\resposta{$\eta=76{,}5\%$; $P_{entrada}=\SI{130.7}{\watt}$}
\resolucao{Rendimentos em cascata \textbf{multiplicam-se}:
\[
\eta=0{,}90\times0{,}85=0{,}765=76{,}5\%.
\]
\[
P_{entrada}=\frac{100}{0{,}765}=\SI{130.7}{\watt}.
\]
\textbf{Erro comum:} \emph{somar} ou fazer a média dos rendimentos, o que daria
$87{,}5\%$ --- valor otimista demais. Cada estágio desperdiça uma fração do que
recebe, e as perdas se compõem multiplicativamente.\\
\textbf{Perdas:} $\SI{13.1}{\watt}$ na fonte e $\SI{17.6}{\watt}$ no
conversor, totalizando $\SI{30.7}{\watt}$ dissipados.}
\stepcounter{questao}
\resposta{$P_{\text{média}}=\SI{200}{\watt}$; $E=\SI{66.7}{\watt\hour}$}
\resolucao{A potência média é a \textbf{média ponderada pelo tempo}:
\[
P_{\text{média}}=\frac{500(5)+100(15)}{20}=\frac{2500+1500}{20}=\SI{200}{\watt}.
\]
\[
E_{ciclo}=P_{\text{média}}\times t_{ciclo}=200\times\frac{20}{60}
=\SI{66.7}{\watt\hour}.
\]
\textbf{Atenção:} a média \emph{aritmética} das potências ($300\,\si{\watt}$)
estaria errada --- ela ignoraria que o patamar baixo dura três vezes mais.
Média de potência é sempre ponderada pelo tempo.}
\stepcounter{questao}
\resposta{(a) $\SI{1}{\ampere}$ e $\SI{10}{\watt}$ (errado); (b)
$\SI{40}{\watt}$; (c) a potência depende de $\langle i^{2}\rangle$}
\resolucao{(a) $I_{\text{média}}=0{,}25(4)+0{,}75(0)=\SI{1}{\ampere}$, e
$RI_{\text{média}}^{2}=10(1)=\SI{10}{\watt}$.\\
(b) A potência instantânea vale $\SI{160}{\watt}$ durante $25\%$ do tempo e
zero no resto:
\[
P_{\text{média}}=0{,}25(160)+0{,}75(0)=\SI{40}{\watt}.
\]
\textbf{Quatro vezes} o valor obtido em (a).\\
(c) A potência é \emph{quadrática} na corrente, e a média de um quadrado não é
o quadrado da média:
\[
P_{\text{média}}=R\left\langle i^{2}\right\rangle
\ne R\left\langle i\right\rangle^{2}.
\]
Define-se então o \textbf{valor eficaz}
$I_{rms}=\sqrt{\langle i^{2}\rangle}=\sqrt{0{,}25(16)}=\SI{2}{\ampere}$, com o
qual $P=RI_{rms}^{2}=\SI{40}{\watt}$ \checkmark.\\
\textbf{Definição operacional do valor eficaz:} é a corrente contínua que
produziria o \emph{mesmo aquecimento}. Aqui, $\SI{2}{\ampere}$ contínuos
aqueceriam tanto quanto o regime pulsado --- e é por isso que multímetros de
qualidade medem "true RMS".}
\stepcounter{questao}
\resposta{$P_{\text{média}}=\SI{5}{\watt}$; a decisão depende da constante térmica}
\resolucao{(a) $P_{\text{média}}=\SI{5}{\watt}$, metade do nominal --- aparentemente
seguro.\\
(b) \textbf{O critério é a comparação entre a duração do pulso e a constante
térmica $\tau_t$ do componente} (tipicamente de segundos a dezenas de segundos
para resistores de fio; menos para os de filme).
\begin{itemize}
\item Se $t_{pulso}\ll\tau_t$, a massa térmica \emph{integra} os pulsos: a
temperatura acompanha a potência \textbf{média}, e o critério de
$\SI{5}{\watt}$ vale.
\item Se $t_{pulso}\gg\tau_t$, o resistor atinge o equilíbrio térmico
\emph{durante} cada pulso: a temperatura corresponde a $\SI{50}{\watt}$
contínuos, e o componente queima.
\end{itemize}
(c) \textbf{Pulsos de $\SI{1}{\milli\second}$:} muito menores que qualquer
$\tau_t$ realista --- o resistor "enxerga" $\SI{5}{\watt}$. \textbf{Seguro.}\\
\textbf{Pulsos de $\SI{10}{\second}$:} comparáveis ou maiores que $\tau_t$ --- o
resistor esquenta como se recebesse $\SI{50}{\watt}$ por dez segundos.
\textbf{Falha provável.}\\
\textbf{Consequência prática:} fabricantes publicam curvas de \emph{potência de
pulso} justamente porque o valor nominal contínuo não responde a essa pergunta.
Nunca dimensione regime pulsado apenas pela média.}
\stepcounter{questao}
\resposta{(a) $78{,}2\%$; (b) $84{,}6\%$ contra $82{,}5\%$ --- melhorar o pior
rende mais}
\resolucao{(a) $\eta=0{,}92(0{,}85)=0{,}782=78{,}2\%$.\\
(b) \textbf{Melhorando o pior:} $0{,}92(0{,}92)=84{,}6\%$ --- ganho de
$6{,}4$ pontos.\\
\textbf{Melhorando o melhor:} $0{,}97(0{,}85)=82{,}5\%$ --- ganho de
$4{,}3$ pontos.\\
E note que o segundo caso exigiu uma melhoria maior em pontos absolutos
($5$ contra $7$) para render \emph{menos}.\\
(c) \textbf{Regra geral.} Como $\eta=\prod\eta_k$, o ganho relativo global é a
soma dos ganhos relativos:
\[
\frac{\Delta\eta}{\eta}=\sum_k\frac{\Delta\eta_k}{\eta_k}.
\]
Um mesmo acréscimo \emph{absoluto} rende mais no estágio de menor $\eta_k$,
pois ali representa maior acréscimo relativo. \textbf{Ataque sempre o elo mais
fraco} --- e é também nele que a margem de melhoria costuma ser maior e mais
barata.}
\stepcounter{questao}
\resposta{$\SI{35}{\kilo\watt\hour}$ e $\mathrm{R\$}\,33{,}29$ por aparelho;
$\mathrm{R\$}\,333$ para dez}
\resolucao{(a) Um ano tem $\SI{8760}{\hour}$:
\[
E=4\times8760=\SI{35040}{\watt\hour}=\SI{35.04}{\kilo\watt\hour};
\]
\[
\text{custo}=35{,}04(0{,}95)=\mathrm{R\$}\,33{,}29.
\]
(b) Dez aparelhos: $\SI{350}{\kilo\watt\hour}$ e cerca de
$\mathrm{R\$}\,333$ por ano --- comparável ao consumo anual de uma geladeira
eficiente.\\
(c) \textbf{Decisões possíveis, com seus custos.}
\begin{itemize}
\item \textbf{Filtro com chave} para desligar grupos de aparelhos: custo de
dezenas de reais, retorno em poucos meses. Melhor relação custo-benefício.
\item \textbf{Redesenho do produto:} normas modernas exigem espera abaixo de
$\SI{0.5}{\watt}$, o que reduziria o custo a menos de
$\mathrm{R\$}\,4{,}20$ por ano.
\item \textbf{Não fazer nada:} defensável para um aparelho isolado, insustentável
em escala --- o consumo agregado de espera responde por vários por cento do
consumo residencial nacional.
\end{itemize}
\textbf{O ponto de engenharia:} $\SI{4}{\watt}$ parecem irrelevantes porque a
potência é pequena. Mas energia é potência \emph{vezes tempo}, e
$\SI{8760}{\hour}$ é um multiplicador enorme.}
\stepcounter{questao}
\resposta{(a) $\SI{333}{\watt}$ ($33\%$); (b) potência \emph{triplicada} ---
destruição}
\resolucao{(a) $R=\dfrac{220^{2}}{1000}=\SI{48.4}{\ohm}$, e em
$\SI{127}{\volt}$:
\[
P=\frac{127^{2}}{48{,}4}=\SI{333}{\watt}
=\left(\frac{127}{220}\right)^{2}\times1000=33{,}3\%.
\]
\textbf{Consequência:} o aparelho funciona, mas com um terço da potência --- um
chuveiro esquenta pouco, um aquecedor não aquece. Situação incômoda, porém
\emph{segura}.\\
(b) $P'=\left(\dfrac{220}{127}\right)^{2}P=3{,}0\,P$ --- a potência
\textbf{triplica}, e a corrente quase dobra.\\
\textbf{Resultado:} superaquecimento imediato, fumaça e, muitas vezes, incêndio.
O disjuntor pode nem atuar, pois a corrente ($\SI{13.6}{\ampere}$ para um
aparelho de $\SI{1000}{\watt}$/$\SI{127}{\volt}$) fica abaixo do valor de
disparo --- a proteção existe para o \emph{circuito}, não para o aparelho.\\
\textbf{Assimetria fundamental:} subtensão é inconveniente; sobretensão é
destrutiva. Por isso a etiqueta de tensão merece conferência antes de ligar
qualquer aparelho resistivo.}
\stepcounter{questao}
\resposta{(a) $\SI{100}{\watt}$; (b) $\SI{1}{\joule}$; (c)
$t=\SI{6.93}{\milli\second}$}
\resolucao{(a) $P_0=RI_0^{2}=4(25)=\SI{100}{\watt}$.\\
(b) A potência decai como $e^{-2t/\tau}$:
\[
E=\int_0^{\infty}RI_0^{2}e^{-2t/\tau}\,\mathrm{d}t
=RI_0^{2}\cdot\frac{\tau}{2}
=100\times\frac{0{,}020}{2}=\SI{1}{\joule}.
\]
\textbf{Regra útil:} a energia equivale à potência inicial mantida por
$\tau/2$. A "constante de tempo da potência" é metade da constante de tempo da
corrente, porque a potência vai com o quadrado.\\
(c) Metade da energia:
\[
\int_0^{t}e^{-2s/\tau}\mathrm{d}s=\frac12\int_0^{\infty}
\;\Longrightarrow\;
1-e^{-2t/\tau}=\frac12
\;\Longrightarrow\;
t=\frac{\tau\ln2}{2}=\SI{6.93}{\milli\second}.
\]
\textbf{Interpretação:} metade de toda a energia é liberada em apenas
$0{,}35\,\tau$ --- o transitório é fortemente concentrado no início. É o pico
inicial, e não a duração, que dimensiona o componente.}
\stepcounter{questao}
\resposta{(a) $\mathrm{R\$}\,104{,}02$ contra $\mathrm{R\$}\,15{,}60$;
(b) cerca de $2 meses$}
\resolucao{(a) Horas anuais: $5\times365=\SI{1825}{\hour}$.
\[
E_{inc}=0{,}060(1825)=\SI{109.5}{\kilo\watt\hour}
\Rightarrow \mathrm{R\$}\,104{,}02;
\]
\[
E_{LED}=0{,}009(1825)=\SI{16.4}{\kilo\watt\hour}
\Rightarrow \mathrm{R\$}\,15{,}60.
\]
Economia anual: $\mathrm{R\$}\,88{,}42$.\\
(b)
\[
t=\frac{15{,}00}{88{,}42}=0{,}17\ \text{ano}\approx2 meses.
\]
\textbf{Retorno excepcional} --- poucos investimentos em engenharia se pagam em
dois meses.\\
\textbf{E há mais:} a lâmpada de LED dura tipicamente $\SI{25000}{\hour}$
contra $\SI{1000}{\hour}$ da incandescente, evitando cerca de $24$ trocas ao
longo da vida. Somando o custo das lâmpadas substituídas e da mão de obra, a
vantagem cresce ainda mais.\\
\textbf{Nota física:} a incandescente converte cerca de $95\%$ da energia em
calor e apenas $5\%$ em luz --- ela é, essencialmente, um aquecedor que
ilumina de lado.}
\stepcounter{questao}
\resposta{(a) $\SI{45.5}{\ampere}$, $\SI{414}{\watt}$; (b) $4{,}1\%$;
(c) $\SI{103}{\watt}$, $1{,}0\%$}
\resolucao{(a) $I=\dfrac{10000}{220}=\SI{45.5}{\ampere}$ e
\[
P_{cabo}=RI^{2}=0{,}2(45{,}5)^{2}=\SI{414}{\watt}.
\]
(b) $414/10000=4{,}1\%$ da potência útil vira calor no cabo.\\
(c) Em $\SI{440}{\volt}$: $I=\SI{22.7}{\ampere}$ e
$P_{cabo}=0{,}2(22{,}7)^{2}=\SI{103}{\watt}$, ou $1{,}0\%$.\\
\textbf{Dobrar a tensão reduziu a perda a um quarto.} A razão é direta: para a
mesma potência, $I\propto1/U$, e a perda vai com $I^{2}$ --- logo
$P_{perda}\propto1/U^{2}$.\\
\textbf{É esta a razão de existir a alta tensão na transmissão.} Linhas de
$\SI{500}{\kilo\volt}$ transportam gigawatts com perdas de poucos por cento; nas
mesmas linhas em baixa tensão, a energia não chegaria ao destino. O preço é o
isolamento e os transformadores nas duas pontas.}
\stepcounter{questao}
\resposta{$I_{rms}=\SI{2}{\ampere}$ contra $I_{\text{média}}=\SI{1}{\ampere}$;
$P=R I_{rms}^{2}$}
\resolucao{(a) Por definição, a potência média em um período $T$ é
\[
P_{\text{média}}=\frac{1}{T}\int_0^{T}R\,i^{2}(t)\,\mathrm{d}t
=R\left[\frac{1}{T}\int_0^{T}i^{2}\,\mathrm{d}t\right]
=R\,I_{rms}^{2},
\]
onde a definição de valor eficaz é justamente
$I_{rms}=\sqrt{\frac1T\int_0^{T}i^{2}\mathrm{d}t}$ --- a
\emph{raiz} da \emph{média} do \emph{quadrado}.\\
(b) Pela desigualdade de Cauchy--Schwarz (ou pelo fato de a variância ser não
negativa):
\[
\left\langle i^{2}\right\rangle-\left\langle i\right\rangle^{2}
=\operatorname{Var}(i)\ge0
\;\Longrightarrow\;
I_{rms}^{2}\ge I_{\text{média}}^{2}.
\]
A igualdade vale se e somente se a variância é nula, isto é, se $i(t)$ é
\textbf{constante}.\\
\textbf{Leitura:} qualquer flutuação da corrente aumenta o aquecimento em
relação ao que a média sugeriria. Corrente ondulada aquece mais que corrente
contínua de mesmo valor médio --- resultado com consequências diretas em
retificadores e conversores.\\
(c) $\langle i\rangle=0{,}25(4)=\SI{1}{\ampere}$;
$\langle i^{2}\rangle=0{,}25(16)=\SI{4}{\ampere\squared}$, logo
$I_{rms}=\SI{2}{\ampere}$.\\
Com $R=\SI{10}{\ohm}$: $P=10(4)=\SI{40}{\watt}$, e não os
$\SI{10}{\watt}$ que a corrente média sugeriria. Fator $4$ --- exatamente
$(I_{rms}/I_{\text{média}})^{2}$.}
\stepcounter{questao}
\resposta{$E=\dfrac{RI_0^{2}\tau}{2}$; equivale à potência inicial mantida por
$\tau/2$}
\resolucao{(a)
\[
E=\int_0^{\infty}R\,I_0^{2}e^{-2t/\tau}\,\mathrm{d}t
=RI_0^{2}\left[-\frac{\tau}{2}e^{-2t/\tau}\right]_0^{\infty}
=\frac{RI_0^{2}\tau}{2}.
\]
\textbf{Observe o fator $\tau/2$, e não $\tau$:} a potência decai com constante
$\tau/2$ porque é proporcional a $i^{2}$, e o quadrado de uma exponencial tem
metade da constante de tempo.\\
(b) Distribuindo $E$ por $5\tau$:
\[
P_{\text{média}}=\frac{RI_0^{2}\tau/2}{5\tau}=\frac{RI_0^{2}}{10}=\frac{P_0}{10}.
\]
Apenas um décimo da potência inicial --- confirmando que o transitório é curto
diante da janela considerada.\\
(c) \textbf{Pulso retangular} de amplitude $I_0$ e duração $\tau$:
$E_{ret}=RI_0^{2}\tau$, exatamente o \textbf{dobro}.\\
\textbf{Consequência de dimensionamento:} substituir o decaimento real por um
pulso retangular equivalente de duração $\tau$ é uma aproximação
\emph{conservadora} --- ela superestima a energia por um fator $2$. Útil como
estimativa rápida e segura, desde que se saiba que há folga embutida.}
\stepcounter{questao}
\resposta{$U_{\max}=\SI{5}{\volt}$, $I_{\max}=\SI{50}{\milli\ampere}$;
limitador de $\SI{380}{\ohm}$}
\resolucao{(a)
\[
U_{\max}=\sqrt{PR}=\sqrt{0{,}25(100)}=\SI{5}{\volt};
\qquad
I_{\max}=\sqrt{\frac{P}{R}}=\sqrt{0{,}0025}=\SI{50}{\milli\ampere}.
\]
(b) Com $\SI{24}{\volt}$ e a exigência $I\le\SI{50}{\milli\ampere}$:
\[
R_{total}\ge\frac{24}{0{,}050}=\SI{480}{\ohm}
\;\Longrightarrow\;
R_{lim}\ge480-100=\SI{380}{\ohm}.
\]
O valor E24 imediatamente superior é $\SI{390}{\ohm}$, que dá
$I=24/490=\SI{49}{\milli\ampere}$ \checkmark.\\
(c) $P_{lim}=390(0{,}049)^{2}=\SI{0.94}{\watt}$ --- especifique
$\SI{2}{\watt}$.\\
\textbf{Observação incômoda:} o limitador dissipa \emph{quatro vezes} mais que o
resistor protegido, e o conjunto consome $\SI{1.18}{\watt}$ para entregar
$\SI{0.24}{\watt}$ ao componente de interesse.\\
\textbf{Quando isso se justifica:} apenas se o resistor de
$\SI{100}{\ohm}$ tiver função específica (sensor, elemento de calibração) que o
torne insubstituível. Caso contrário, redimensionar o próprio resistor para uma
potência maior é mais simples e mais eficiente.}
\stepcounter{questao}
\resposta{(a) $1250$ e $\SI{1087}{\kilo\watt\hour}$, diferença de
$\SI{163}{\kilo\watt\hour}$; (b) cerca de $\SI{5.2}{year}$}
\resolucao{(a)
\[
E_A=\frac{1000}{0{,}80}=\SI{1250}{\kilo\watt\hour};
\qquad
E_B=\frac{1000}{0{,}92}=\SI{1087}{\kilo\watt\hour};
\]
diferença de $\SI{163}{\kilo\watt\hour}$ por ano.\\
(b) Economia anual: $163(0{,}95)=\mathrm{R\$}\,154{,}89$.
\[
t=\frac{800}{154{,}89}=\SI{5.2}{year}.
\]
(c) \textbf{O que a conta simples ignora.}
\begin{itemize}
\item \textbf{Reajuste da tarifa:} se a energia subir acima da inflação, o
retorno acelera. Historicamente é o que ocorre.
\item \textbf{Custo do capital:} $\mathrm{R\$}\,800$ aplicados renderiam algo
--- o retorno real é mais longo que o nominal.
\item \textbf{Vida útil:} se o equipamento durar $\SI{10}{year}$, B gera lucro
nos últimos cinco. Se durar $\SI{4}{year}$, a troca não se paga.
\item \textbf{Custo do calor rejeitado:} A dissipa $\SI{250}{\kilo\watt\hour}$
contra $\SI{87}{\kilo\watt\hour}$ de B. Em ambiente climatizado, esse calor
extra precisa ser \emph{removido} --- o que pode dobrar a economia efetiva.
\item \textbf{Regime de uso:} rendimentos variam com a carga, e o valor de
catálogo costuma ser o de plena carga.
\end{itemize}
\textbf{Conclusão:} $\SI{5.2}{year}$ é aceitável para equipamento industrial de
vida longa e uso contínuo; é ruim para uso intermitente ou equipamento de
substituição rápida. \textbf{A resposta depende do horizonte, não apenas do
número.}}
\stepcounter{questao}
\resposta{(a) $\SI{60}{\hour}$ por mês, ou $\SI{2}{\hour}$ por dia;
(b) $\SI{750}{\watt}$}
\resolucao{(a) O orçamento permite $\SI{90}{\kilo\watt\hour}$ por mês:
\[
t=\frac{90}{1{,}5}=\SI{60}{\hour\per month}
=\SI{2}{\hour\per day}.
\]
(b) Para $4\times30=\SI{120}{\hour}$ mensais dentro dos mesmos
$\SI{90}{\kilo\watt\hour}$:
\[
P=\frac{90}{120}=\SI{0.75}{\kilo\watt}=\SI{750}{\watt}.
\]
(c) \textbf{Depende do objetivo térmico.} As duas soluções entregam a
\emph{mesma energia} --- e portanto o mesmo calor total no mês.
\begin{itemize}
\item \textbf{$\SI{1.5}{\kilo\watt}$ por $\SI{2}{\hour}$:} aquece rápido e
atinge temperatura mais alta, mas o ambiente esfria nas outras
$\SI{22}{\hour}$.
\item \textbf{$\SI{750}{\watt}$ por $\SI{4}{\hour}$:} aquecimento mais suave e
prolongado, com temperatura média mais estável.
\end{itemize}
\textbf{Para conforto térmico, a segunda costuma ser superior} --- o ambiente
tem inércia térmica, e potência menor por mais tempo reduz as oscilações.\\
\textbf{E há uma terceira solução, melhor que ambas:} melhorar o
\emph{isolamento} do ambiente. Isso reduz a energia necessária, e não apenas a
redistribui --- atacando a causa em vez do sintoma.}
\stepcounter{questao}
\resposta{No ótimo, o custo anualizado do cabo iguala o custo anual da energia
perdida}
\resolucao{(a) Sejam $S$ a seção, $k_1$ o custo anualizado do cabo por unidade
de seção e $R=\rho\ell/S$. O custo total anual é
\[
C(S)=\underbrace{k_1S}_{\text{investimento}}
+\underbrace{\frac{\rho\ell}{S}I^{2}\,t\,c_E}_{\text{energia perdida}}
=k_1S+\frac{k_2}{S},
\]
com $k_2=\rho\ell I^{2}t\,c_E$.\\
(b) Derivando e anulando:
\[
\frac{\mathrm{d}C}{\mathrm{d}S}=k_1-\frac{k_2}{S^{2}}=0
\;\Longrightarrow\;
S^{*}=\sqrt{\frac{k_2}{k_1}}.
\]
A segunda derivada, $2k_2/S^{3}>0$, confirma o mínimo.\\
(c) \textbf{Substituindo $S^{*}$ nas duas parcelas:}
\[
k_1S^{*}=\sqrt{k_1k_2}
\qquad\text{e}\qquad
\frac{k_2}{S^{*}}=\sqrt{k_1k_2}.
\]
As duas são \textbf{iguais} no ótimo --- resultado conhecido como \emph{lei de
Kelvin}: \textbf{a seção econômica é aquela em que o custo anualizado do
condutor iguala o custo anual das perdas}.\\
\textbf{Leituras práticas.}
\begin{itemize}
\item Cabos usados \emph{muitas horas} por ano ($t$ grande) justificam seções
maiores --- $S^{*}\propto\sqrt{t}$.
\item Correntes maiores também: $S^{*}\propto I$.
\item Energia mais cara empurra a seção para cima: $S^{*}\propto\sqrt{c_E}$.
\end{itemize}
\textbf{Importante:} a seção econômica é frequentemente \emph{maior} que a
mínima exigida pela norma (que garante apenas segurança térmica e queda de
tensão). Dimensionar pelo mínimo normativo é seguro, mas costuma ser
economicamente subótimo em instalações de uso intenso.}
\stepcounter{questao}
\resposta{(a) $\SI{0.0658}{\kilogram\per\second}\approx
\SI{3.9}{\liter\per\minute}$; (b) $\mathrm{R\$}\,0{,}70$}
\resolucao{(a) A potência necessária para aquecer uma vazão mássica $\dot m$
por $\Delta T=\SI{20}{\kelvin}$ é $P=\dot m\,c\,\Delta T$:
\[
\dot m=\frac{P}{c\,\Delta T}=\frac{5500}{4180(20)}
=\SI{0.0658}{\kilogram\per\second}\approx\SI{3.9}{\liter\per\minute}.
\]
Vazão modesta --- e é exatamente por isso que abrir demais o registro faz a água
esfriar: a potência é fixa, e mais vazão significa menor $\Delta T$.\\
(b)
\[
E=5{,}5\,\si{\kilo\watt}\times\frac{8}{60}\,\si{\hour}
=\SI{0.733}{\kilo\watt\hour}
\;\Longrightarrow\;
\text{custo}=0{,}733(0{,}95)=\mathrm{R\$}\,0{,}70.
\]
Um banho diário custa cerca de $\mathrm{R\$}\,21$ por mês --- e para uma
família de quatro pessoas, mais de $\mathrm{R\$}\,80$.\\
(c) \textbf{Por que o rendimento é quase $100\%$.} A resistência está
\emph{imersa} na água, e \textbf{toda} a energia dissipada por efeito Joule vira
calor exatamente onde ele é desejado. Não há conversão intermediária, nem calor
"desperdiçado" --- o que seria perda em outro contexto é aqui o produto.\\
\textbf{Contraste instrutivo:} uma lâmpada incandescente tem os mesmos
$\approx100\%$ de conversão em calor, mas rendimento \emph{luminoso} de $5\%$.
O rendimento só se define em relação ao \emph{objetivo}: a mesma física é
excelente para aquecer e péssima para iluminar.}
\stepcounter{questao}
\resposta{Acréscimo de $\SI{140}{\kilo\watt\hour}$, equivalente a
$\SI{194}{\watt}$ contínuos}
\resolucao{(a) O mês tem cerca de $\SI{720}{\hour}$:
\[
P=\frac{140000}{720}=\SI{194}{\watt}\ \text{contínuos}.
\]
(b) \textbf{Hipóteses compatíveis com $\approx\SI{200}{\watt}$ permanentes.}
\begin{itemize}
\item \textbf{Fuga para o terra} por isolamento danificado --- especialmente em
resistência de chuveiro ou aquecedor de água. Uma fuga de
$\SI{0.88}{\ampere}$ em $\SI{220}{\volt}$ dá exatamente
$\SI{194}{\watt}$.
\item \textbf{Geladeira com vedação ruim ou gás insuficiente}, operando em
regime quase contínuo em vez de ciclado.
\item \textbf{Bomba ou motor} funcionando sem parar por falha de boia ou
pressostato.
\item \textbf{Ligação clandestina} a jusante do medidor.
\end{itemize}
(c) \textbf{Procedimento de investigação, em ordem.}
\begin{enumerate}
\item \textbf{Desligar o disjuntor geral} e verificar se o medidor continua
girando. Se sim, o problema está \emph{antes} do quadro --- ligação irregular ou
medidor defeituoso.
\item \textbf{Religar circuito por circuito}, observando o medidor a cada
etapa, para isolar o ramo responsável.
\item \textbf{Medir a corrente} de cada circuito suspeito com alicate
amperímetro, com todos os aparelhos desligados. Corrente não nula indica fuga.
\item \textbf{Testar o dispositivo DR}, se houver --- fuga acima de
$\SI{30}{\milli\ampere}$ deveria ter provocado desarme. A ausência de desarme
com fuga confirmada indica DR defeituoso, e é problema de segurança, não apenas
de conta.
\end{enumerate}
\textbf{Prioridade:} uma fuga de $\SI{0.88}{\ampere}$ é risco de choque e de
incêndio. O custo na conta é o sintoma menos grave.}

\gabarito[2.3 Potência elétrica e energia]

\subsection{Fontes reais e Leis de Kirchhoff}
\stepcounter{questao}
\resposta{(a) \SI{2}{\volt}; (b) \SI{0.4}{\ampere}}
\resolucao{(a) Em oposição, as fem se subtraem:
$U = 6 - 4 = \SI{2}{\volt}$.\\
(b) $I = \dfrac{U}{R} = \dfrac{2}{5} = \SI{0.4}{\ampere}$.
Se estivessem em série \emph{concordante}, teríamos $6+4 = \SI{10}{\volt}$ e
$I = \SI{2}{\ampere}$.}
\stepcounter{questao}
\resposta{$I_3 = \SI{8}{\ampere}$}
\resolucao{A lei dos nós afirma que a soma das correntes que entram é igual à das
que saem:
\[
I_3 = I_1 + I_2 = 5 + 3 = \SI{8}{\ampere}.
\]
Ela expressa a conservação da carga: o nó não acumula carga.}
\stepcounter{questao}
\resposta{$I = \SI{2}{\ampere}$; $U = \SI{11}{\volt}$}
\resolucao{A resistência interna está em série com a carga:
\[
I=\frac{\EMF}{R+r}=\frac{12}{5{,}5+0{,}5}=\SI{2}{\ampere}.
\]
A tensão nos terminais é a fem menos a queda interna:
\[
U=\EMF - r\,I = 12 - 0{,}5\cdot2 = \SI{11}{\volt}.
\]
A diferença de \SI{1}{\volt} "some" dentro da bateria --- é a queda em $r$. Quanto
maior a corrente, menor a tensão entregue: por isso baterias velhas (com $r$
elevado) fornecem menos tensão sob carga.}
\stepcounter{questao}
\resposta{$\EMF = \SI{14}{\volt}$; em vazio, $U = \SI{14}{\volt}$}
\resolucao{Da equação da fonte real, $U = \EMF - r\,I$:
\[
\EMF = U + r\,I = 12 + 0{,}4\cdot5 = \SI{14}{\volt}.
\]
Em vazio ($I = 0$), não há queda interna: $U = \EMF = \SI{14}{\volt}$. A tensão de
circuito aberto de uma fonte é sempre igual à sua fem.}
\stepcounter{questao}
\resposta{$I_5 = \SI{2}{\ampere}$, saindo do nó}
\resolucao{Pela lei dos nós, a soma das que entram é igual à soma das que saem.
Entram $6 + 4 = \SI{10}{\ampere}$; saem $3 + 5 = \SI{8}{\ampere}$ pelos ramos
conhecidos. Para equilibrar, $I_5$ deve \emph{sair} com:
\[
I_5 = 10 - 8 = \SI{2}{\ampere}.
\]
Como o sentido adotado na figura já era de saída, o valor positivo confirma essa
suposição.}
\stepcounter{questao}
\resposta{$I_1 = \SI{3}{\ampere}$, $I_2 = \SI{1}{\ampere}$, $I_3 = \SI{2}{\ampere}$}
\resolucao{\textbf{Lei dos nós} (no topo): $I_1 = I_2 + I_3$, logo $I_3 = I_1 - I_2$.\\
\textbf{Malha esquerda} ($E_1$, $R_1$, $R_2$):
\[
E_1 = R_1 I_1 + R_2 I_3 \;\Rightarrow\; 18 = 2I_1 + 6(I_1 - I_2).
\]
\textbf{Malha direita} ($E_2$, $R_3$, $R_2$):
\[
E_2 = R_3 I_2 + R_2 I_3 \;\Rightarrow\; 6 = 6I_2 + 6(I_1 - I_2)\cdot(-1)?
\]
Organizando o sistema (com $I_3 = I_1 - I_2$):
\[
\begin{cases}
8I_1 - 6I_2 = 18\\[2pt]
-6I_1 + 12I_2 = -6
\end{cases}
\;\Rightarrow\;
I_1 = \SI{3}{\ampere},\ I_2 = \SI{1}{\ampere}.
\]
Logo $I_3 = 3 - 1 = \SI{2}{\ampere}$.\\
\textbf{Verificação} pela malha esquerda: $2\cdot3 + 6\cdot2 = 6+12 =
\SI{18}{\volt} = E_1$. Resultado consistente.}
\stepcounter{questao}
\resposta{$I_1 = \SI{2.4}{\ampere}$, $I_2 = \SI{-0.6}{\ampere}$, $I_3 = \SI{1.8}{\ampere}$}
\resolucao{\textbf{Nó superior:} $I_1 + I_2 = I_3$.\\
\textbf{Malha esquerda:} $E_1 = R_1 I_1 + R_2 I_3
\Rightarrow 12 = 2I_1 + 4I_3$.\\
\textbf{Malha direita:} $E_2 = R_3 I_2 + R_2 I_3
\Rightarrow 6 = 2I_2 + 4I_3$.\\
Resolvendo o sistema (com $I_3 = I_1 + I_2$):
\[
I_1 = \SI{2.4}{\ampere},\quad I_2 = \SI{-0.6}{\ampere},\quad
I_3 = \SI{1.8}{\ampere}.
\]
O sinal \emph{negativo} de $I_2$ significa que a corrente real nesse ramo é
\emph{oposta} ao sentido adotado: a fonte $E_2$, mais fraca, na verdade
\textbf{recebe} corrente (está sendo carregada) em vez de fornecer.\\
\textbf{Verificação} (malha esquerda): $2\cdot2{,}4 + 4\cdot1{,}8 = 4{,}8+7{,}2
= \SI{12}{\volt} = E_1$. Resultado consistente. Este é um bom lembrete: adotar um sentido
errado não invalida a solução --- o sinal negativo se encarrega de corrigi-lo.}
\stepcounter{questao}
\resposta{(a) $V_{ac}=\SI{6}{\volt}$; (b) $V_{ab}=\SI{14}{\volt}$,
$V_{bc}=\SI{-8}{\volt}$; (c) ver comentário}
\resolucao{\textbf{Estratégia.} Como os terminais estão abertos, \emph{não circula
corrente} em nenhum ramo. Sem corrente, não há queda em resistência alguma: os
potenciais são determinados exclusivamente pelas fem, somadas com o \emph{sinal}
correto. Basta adotar uma referência e "caminhar" pelo circuito.

\medskip
\textbf{Passo 1 --- adotar a referência.} Seja o nó inferior do barramento
(ao qual se ligam os ramos de $b$ e $c$) o nó de referência: $V_{N_2}=0$.

\textbf{Passo 2 --- subir pela fonte de \SI{5}{\volt} vertical.} Ela tem o
terminal $+$ \emph{em cima} e $-$ embaixo. Indo de $N_2$ para $N_1$, entramos
pelo $-$ e saímos pelo $+$: \emph{ganhamos} \SI{5}{\volt}.
\[
V_{N_1}=V_{N_2}+5=\SI{5}{\volt}.
\]

\textbf{Passo 3 --- percorrer cada ramo até seu terminal.}
\begin{itemize}[leftmargin=1.7em,itemsep=2pt,topsep=3pt]
\item \textbf{Ramo de $a$} (fonte \SI{4}{\volt}, com $-$ à esquerda e $+$ à
      direita): indo de $N_1$ para $a$, entramos pelo $-$ e saímos pelo $+$,
      ganhando \SI{4}{\volt}:
      \[V_a = V_{N_1}+4 = 5+4=\SI{9}{\volt}.\]
\item \textbf{Ramo de $b$} (fonte \SI{5}{\volt}, com $+$ à esquerda e $-$ à
      direita): indo de $N_2$ para $b$, entramos pelo $+$ e saímos pelo $-$,
      \emph{perdendo} \SI{5}{\volt}:
      \[V_b = V_{N_2}-5 = \SI{-5}{\volt}.\]
\item \textbf{Ramo de $c$} (fonte \SI{3}{\volt}, com $-$ à esquerda e $+$ à
      direita): ganhamos \SI{3}{\volt}:
      \[V_c = V_{N_2}+3 = \SI{3}{\volt}.\]
\end{itemize}

\textbf{Passo 4 --- calcular as diferenças.}
\[
\boxed{\;V_{ac}=V_a-V_c = 9-3 = \SI{6}{\volt}\;}
\]
\[
V_{ab}=V_a-V_b = 9-(-5)=\SI{14}{\volt},
\qquad
V_{bc}=V_b-V_c = -5-3=\SI{-8}{\volt}.
\]
\textbf{Consistência:} $V_{ab}+V_{bc} = 14+(-8) = \SI{6}{\volt} = V_{ac}$
\checkmark --- como deve ser, pois a diferença de potencial é uma função de
estado (independe do caminho).

\medskip
\textbf{(c) Por que os fios não importam?} A queda em um resistor é $U=RI$. Com
$I=0$ em todos os ramos (terminais abertos), \emph{qualquer} resistência produz
queda nula. Os potenciais dos terminais dependem só das fem e de suas
polaridades.\\
\textbf{Atenção ao erro mais comum:} somar as fontes ignorando a polaridade
(obtendo, por exemplo, $4+5+3=12$). O sinal de cada fonte é determinado pelo
\emph{sentido em que a percorremos}: entrando pelo $-$ e saindo pelo $+$, soma-se;
o contrário, subtrai-se. \emph{Sempre marque a polaridade antes de escrever a
primeira equação.}}
\stepcounter{questao}
\resposta{(a) $V = \SI{12}{\volt}$; (b) $I_1 = 0$, $I_2 = \SI{6}{\ampere}$;
(c) ver comentário}
\resolucao{\textbf{(a)} Tomando o nó superior com potencial $V$ e o inferior como
terra, a lei dos nós exige que a soma das correntes que \emph{chegam} ao nó seja
igual à que sai pela carga:
\[
\frac{\EMF_1-V}{r_1}+\frac{\EMF_2-V}{r_2}=\frac{V}{R_L}.
\]
Substituindo:
\[
\frac{12-V}{0{,}2}+\frac{12{,}6-V}{0{,}1}=\frac{V}{2}.
\]
Multiplicando por 2: $10(12-V)+20(12{,}6-V)=V$, ou seja,
\[
120-10V+252-20V=V
\;\Longrightarrow\;
372 = 31V
\;\Longrightarrow\;
V = \SI{12}{\volt}.
\]

\textbf{(b)} As correntes de cada fonte:
\[
I_1=\frac{12-12}{0{,}2}=\SI{0}{\ampere},
\qquad
I_2=\frac{12{,}6-12}{0{,}1}=\SI{6}{\ampere},
\qquad
I_L=\frac{12}{2}=\SI{6}{\ampere}.\ \checkmark
\]
\textbf{(c) Interpretação e risco.} A bateria de \emph{maior} fem assume
praticamente toda a carga, enquanto a de menor fem fica ociosa. Se a diferença de
fem fosse maior --- por exemplo, uma bateria descarregada em paralelo com uma
carregada --- a corrente na bateria mais fraca ficaria \textbf{negativa}: ela
seria \emph{carregada} pela outra, com corrente limitada apenas pelas
resistências internas, tipicamente muito pequenas.\\
\textbf{Consequência prática:} correntes de circulação elevadíssimas, aquecimento
e risco de dano ou explosão. Por isso o paralelismo de baterias exige unidades
\emph{idênticas em fem, tecnologia e estado de carga}, ou o uso de diodos de
isolamento e controladores de balanceamento --- prática obrigatória em bancos de
baterias e sistemas fotovoltaicos.}
\stepcounter{questao}
\resposta{$V=\SI{5}{\volt}$}
\resolucao{
\[
V=\EMF-rI=6-1(1)=\SI{5}{\volt}.
\]
A diferença de $\SI{1}{\volt}$ é a queda \emph{interna}, que não chega à carga.
Ela cresce com a corrente --- é por isso que a tensão de uma pilha cai quando se
liga um aparelho que exige mais.}
\stepcounter{questao}
\resposta{$r=\SI{0.3}{\ohm}$}
\resolucao{
\[
r=\frac{\EMF-V}{I}=\frac{9-8{,}4}{2}=\SI{0.3}{\ohm}.
\]
\textbf{Método geral:} a queda interna dividida pela corrente. Uma única medição
sob carga, somada ao valor de $\EMF$, determina $r$ --- e $\EMF$ pode ser lida
em vazio com um voltímetro de alta impedância.}
\stepcounter{questao}
\resposta{$I=\SI{4}{\ampere}$}
\resolucao{Pela lei dos nós,
\[
\sum I_{entra}=\sum I_{sai}
\;\Longrightarrow\;
2{,}5+1{,}5=\SI{4}{\ampere}.
\]
\textbf{Fundamento:} a LKC expressa a \textbf{conservação da carga}. Um nó não
acumula carga --- tudo o que entra, sai. Isso vale independentemente dos
elementos ligados a ele, e mesmo que sejam não lineares.}
\stepcounter{questao}
\resposta{$I=\SI{1}{\ampere}$}
\resolucao{Pela LKT, a soma das quedas iguala a f.e.m.:
\[
12=4I+8I=12I
\;\Longrightarrow\;
I=\SI{1}{\ampere}.
\]
\textbf{Fundamento:} a LKT expressa a \textbf{conservação da energia}. Uma carga
que percorre a malha e retorna ao ponto de partida deve ter ganho e perdido a
mesma energia --- o potencial é uma função de estado.}
\stepcounter{questao}
\resposta{(a)}
\resolucao{A fonte \textbf{ideal} tem $r=0$: sua tensão terminal é
$\EMF$ independentemente da corrente, inclusive em curto-circuito --- o que
exigiria corrente infinita, revelando que se trata de uma idealização.\\
A fonte \textbf{real} tem $r>0$, e sua tensão cai como $\EMF-rI$.\\
\textbf{Quando a idealização é aceitável:} sempre que $r\ll R_{carga}$. Uma
fonte de bancada com $r=\SI{10}{\milli\ohm}$ alimentando
$\SI{100}{\ohm}$ é, para todos os efeitos, ideal.}
\stepcounter{questao}
\resposta{$I_{cc}=\SI{6}{\ampere}$}
\resolucao{
\[
I_{cc}=\frac{\EMF}{r}=\frac{1{,}5}{0{,}25}=\SI{6}{\ampere}.
\]
\textbf{Ressalva importante:} o valor pressupõe $r$ constante --- hipótese
otimista. Em uma pilha real, $r$ cresce com a corrente e com o aquecimento, e a
corrente efetiva é menor. Ainda assim, $\SI{6}{\ampere}$ em uma pilha pequena
significa $\SI{9}{\watt}$ dissipados \emph{dentro} dela: aquecimento rápido,
vazamento de eletrólito e risco de ruptura.}
\stepcounter{questao}
\resposta{$I=\SI{2}{\ampere}$, $V=\SI{10.8}{\volt}$;
$P_R=\SI{21.6}{\watt}$, $P_r=\SI{2.4}{\watt}$, $\eta=90\%$}
\resolucao{(a)
\[
I=\frac{\EMF}{r+R}=\frac{12}{6}=\SI{2}{\ampere};
\qquad
V=12-0{,}6(2)=\SI{10.8}{\volt}.
\]
(b)
\[
P_{total}=\EMF\,I=\SI{24}{\watt};
\quad
P_R=5{,}4(4)=\SI{21.6}{\watt};
\quad
P_r=0{,}6(4)=\SI{2.4}{\watt};
\]
\[
\eta=\frac{R}{R+r}=\frac{5{,}4}{6}=90\%.
\]
\textbf{Fórmula útil:} o rendimento de uma fonte real é sempre
$R/(R+r)$ --- ele depende só da razão entre as resistências, e não da tensão
nem da corrente.}
\stepcounter{questao}
\resposta{$I=\SI{5}{\ampere}$; $V=\SI{14}{\volt}$}
\resolucao{(a) As duas f.e.m. estão em \textbf{oposição}, e a diferença é que
impulsiona a corrente:
\[
I=\frac{14-12}{0{,}4}=\SI{5}{\ampere}.
\]
(b) A bateria está \textbf{recebendo} corrente, então a queda interna
\emph{soma-se} à f.e.m.:
\[
V=\EMF+rI=12+0{,}4(5)=\SI{14}{\volt}.
\]
\textbf{Inversão do sinal:} quando a fonte fornece, $V=\EMF-rI<\EMF$; quando
recebe, $V=\EMF+rI>\EMF$. É por isso que a tensão de uma bateria de automóvel
sobe de $\SI{12.6}{\volt}$ para cerca de $\SI{14}{\volt}$ com o motor ligado ---
sinal de que o alternador está carregando.}
\stepcounter{questao}
\resposta{$I_5=\SI{3}{\ampere}$, saindo}
\resolucao{
\[
3+4+2=6+I_5
\;\Longrightarrow\;
I_5=9-6=\SI{3}{\ampere}.
\]
O sinal positivo confirma o sentido admitido (saindo).\\
\textbf{Convenção alternativa:} adotar todas as correntes como \emph{entrando} e
impor soma nula. As que realmente saem aparecem com sinal negativo. As duas
convenções são equivalentes --- escolha uma e mantenha-a.}
\stepcounter{questao}
\resposta{$I=\SI{3}{\ampere}$; a de $\SI{6}{\volt}$ absorve}
\resolucao{
\[
I=\frac{24-6}{6}=\SI{3}{\ampere},
\]
no sentido imposto pela fonte maior.\\
\textbf{Balanço:} a de $\SI{24}{\volt}$ fornece $24(3)=\SI{72}{\watt}$; a de
$\SI{6}{\volt}$ \emph{absorve} $6(3)=\SI{18}{\watt}$; os resistores dissipam
$6(9)=\SI{54}{\watt}$. Soma: $18+54=\SI{72}{\watt}$ \checkmark.\\
\textbf{Critério para identificar:} a fonte \emph{absorve} quando a corrente
entra pelo seu terminal positivo. Se ela for recarregável, está sendo
carregada; se não for (uma pilha comum), a energia vira calor e o dispositivo
pode ser danificado.}
\stepcounter{questao}
\resposta{Série: $\SI{3}{\volt}$ e $\SI{0.4}{\ohm}$; paralelo:
$\SI{1.5}{\volt}$ e $\SI{0.1}{\ohm}$}
\resolucao{\textbf{Série:} as f.e.m. se somam e as resistências internas também:
$\EMF=\SI{3}{\volt}$, $r=\SI{0.4}{\ohm}$.\\
\textbf{Paralelo} (pilhas \emph{idênticas}): a f.e.m. permanece
$\SI{1.5}{\volt}$ e as resistências internas ficam em paralelo:
$r=\SI{0.1}{\ohm}$.\\
\textbf{Correntes de curto:} série $3/0{,}4=\SI{7.5}{\ampere}$; paralelo
$1{,}5/0{,}1=\SI{15}{\ampere}$.\\
\textbf{Advertência:} a associação em paralelo só é segura com pilhas
\emph{idênticas e no mesmo estado de carga}. Com f.e.m. diferentes, circula
corrente entre elas mesmo sem carga externa.}
\stepcounter{questao}
\resposta{$4$ de LKC e $4$ de LKT, totalizando $8$}
\resolucao{\textbf{LKC:} $n-1=4$ equações. A do quinto nó é combinação linear
das anteriores --- redundante.\\
\textbf{LKT:} $b-n+1=8-5+1=4$ equações, uma por malha independente.\\
\textbf{Total:} $4+4=8=b$ --- exatamente o número de correntes de ramo a
determinar. O sistema é quadrado e bem posto.\\
\textbf{Não é coincidência:} $(n-1)+(b-n+1)=b$ sempre. As leis de Kirchhoff
fornecem \emph{precisamente} as equações necessárias, nem uma a mais nem a
menos.}
\stepcounter{questao}
\resposta{$V=\SI{10.67}{\volt}$; $\SI{0.89}{\ampere}$ e
$\SI{1.78}{\ampere}$}
\resolucao{
\[
R_{par}=\frac{12(6)}{18}=\SI{4}{\ohm};
\qquad
I=\frac{12}{4{,}5}=\SI{2.67}{\ampere};
\qquad
V=12-0{,}5(2{,}67)=\SI{10.67}{\volt}.
\]
\emph{(Conferindo: $V=R_{par}I=4(2{,}67)=\SI{10.67}{\volt}$ \checkmark.)}
\[
I_{12}=\frac{10{,}67}{12}=\SI{0.89}{\ampere};
\qquad
I_{6}=\frac{10{,}67}{6}=\SI{1.78}{\ampere}.
\]
\textbf{Acoplamento pelas cargas:} se a carga de $\SI{6}{\ohm}$ for desligada,
a corrente total cai para $12/12{,}5=\SI{0.96}{\ampere}$ e a tensão sobe para
$\SI{11.52}{\volt}$ --- o que \emph{aumenta} a corrente da carga remanescente.
As cargas se influenciam através de $r$.}
\stepcounter{questao}
\resposta{$\SI{1.36}{\volt}$ contra $\SI{0.93}{\volt}$ --- queda de $32\%$}
\resolucao{\textbf{Nova:} $I=1{,}5/3{,}3=\SI{0.455}{\ampere}$ e
$V=3(0{,}455)=\SI{1.36}{\volt}$.\\
\textbf{Gasta:} $I=1{,}3/4{,}2=\SI{0.310}{\ampere}$ e
$V=3(0{,}310)=\SI{0.93}{\volt}$.\\
\textbf{Observe o que domina:} a f.e.m. caiu apenas $13\%$, mas a resistência
interna \textbf{quadruplicou} --- e é ela a responsável pela maior parte da
perda de desempenho.\\
\textbf{Consequência para diagnóstico:} medir a tensão de uma pilha \emph{em
vazio} é enganoso, pois ela mostra quase $\EMF$ mesmo já esgotada. O teste
correto é sob carga, ou a medição direta de $r$.}
\stepcounter{questao}
\resposta{Fornecido $\SI{40}{\watt}$; dissipado $4+\SI{36}{\watt}$}
\resolucao{
\[
I=\frac{20}{10}=\SI{2}{\ampere};
\qquad
P_{fonte}=\EMF\,I=\SI{40}{\watt};
\]
\[
P_r=1(4)=\SI{4}{\watt};
\qquad
P_R=9(4)=\SI{36}{\watt};
\qquad 4+36=40\ \checkmark
\]
\textbf{Cuidado com a distinção:} a potência \emph{gerada} usa $\EMF$, não a
tensão terminal. Usando $V=\SI{18}{\volt}$ obter-se-ia $\SI{36}{\watt}$ ---
que é a potência \emph{entregue aos terminais}, e não a produzida internamente.
A diferença de $\SI{4}{\watt}$ é exatamente o que $r$ consome.}
\stepcounter{questao}
\resposta{$V=\SI{8}{\volt}$; $I_1=\SI{4}{\ampere}$,
$I_2=\SI{-2}{\ampere}$, $I_R=\SI{2}{\ampere}$ --- a fonte 2 é
\textbf{carregada}}
\resolucao{(a) LKC no nó comum:
\[
\frac{V-12}{1}+\frac{V-6}{1}+\frac{V}{4}=0
\;\Longrightarrow\;
4V-48+4V-24+V=0
\;\Longrightarrow\;
V=\SI{8}{\volt}.
\]
\[
I_1=\frac{12-8}{1}=\SI{4}{\ampere};
\qquad
I_2=\frac{6-8}{1}=\SI{-2}{\ampere};
\qquad
I_R=\frac{8}{4}=\SI{2}{\ampere}.
\]
\textbf{LKC:} $4+(-2)=2$ \checkmark.\\
(b) \textbf{A fonte 2 tem corrente negativa} --- ela não fornece, \emph{recebe}
$\SI{2}{\ampere}$. A fonte 1, mais forte, alimenta a carga \emph{e} carrega a
fonte 2.\\
\textbf{Balanço completo:}
\begin{center}
\begin{tabular}{lr}
\hline
Fonte 1 fornece & $\SI{48}{\watt}$\\
$r_1$ dissipa & $\SI{16}{\watt}$\\
$R$ dissipa & $\SI{16}{\watt}$\\
Fonte 2 absorve (f.e.m.) & $\SI{12}{\watt}$\\
$r_2$ dissipa & $\SI{4}{\watt}$\\
\hline
\end{tabular}
\end{center}
Total consumido: $16+16+12+4=\SI{48}{\watt}$ \checkmark.\\
\textbf{Se a fonte 2 for recarregável}, os $\SI{12}{\watt}$ são armazenados. Se
não for, viram calor --- e o dispositivo aquece perigosamente.}
\stepcounter{questao}
\resposta{Série: $\SI{6}{\volt}$/$\SI{0.8}{\ohm}$; paralelo:
$\SI{1.5}{\volt}$/$\SI{0.05}{\ohm}$; $P_{\max}=\SI{11.25}{\watt}$ nos
\textbf{dois}}
\resolucao{(a)
\begin{center}
\begin{tabular}{lccc}
\hline
Arranjo & $\EMF$ & $r$ & $I_{cc}$\\
\hline
Série & $\SI{6}{\volt}$ & $\SI{0.8}{\ohm}$ & $\SI{7.5}{\ampere}$\\
Paralelo & $\SI{1.5}{\volt}$ & $\SI{0.05}{\ohm}$ & $\SI{30}{\ampere}$\\
\hline
\end{tabular}
\end{center}
(b)
\[
\text{Série: } \frac{6^{2}}{4(0{,}8)}=\SI{11.25}{\watt};
\qquad
\text{Paralelo: } \frac{1{,}5^{2}}{4(0{,}05)}=\SI{11.25}{\watt}.
\]
\textbf{Idênticas} --- e não por acaso. Ambos os arranjos usam as mesmas quatro
pilhas, e a energia total disponível é a mesma. A topologia decide \emph{como}
ela é entregue, não \emph{quanto}.\\
(c) \textbf{Critério de escolha.}
\begin{itemize}
\item \textbf{Série} para cargas de \emph{alta} resistência, que precisam de
tensão. A carga ótima é $\SI{0.8}{\ohm}$.
\item \textbf{Paralelo} para cargas de \emph{baixa} resistência, que precisam de
corrente. A carga ótima é $\SI{0.05}{\ohm}$.
\item \textbf{Série-paralelo} (dois ramos de duas) daria $\SI{3}{\volt}$ e
$\SI{0.2}{\ohm}$ --- e, de novo, $\SI{11.25}{\watt}$.
\end{itemize}
\textbf{Regra:} escolha o arranjo cuja resistência interna se aproxime da carga
--- ou, mais realisticamente, cuja \emph{tensão} atenda ao que a carga exige.}
\stepcounter{questao}
\resposta{A soma das correntes que atravessam qualquer superfície fechada é
nula}
\resolucao{(a) \textbf{LKC generalizada:} para qualquer superfície fechada que
corte o circuito, a soma algébrica das correntes que a atravessam é zero.\\
(b) \textbf{Justificativa:} a carga total dentro da superfície não pode variar
(em regime permanente). Como corrente é fluxo de carga, o que entra deve
igualar o que sai --- exatamente o mesmo argumento que vale para um nó, agora
aplicado a uma região que \emph{contém} vários nós.\\
Formalmente, escrevendo a LKC para cada nó interno e somando, as correntes dos
ramos \emph{internos} aparecem duas vezes com sinais opostos e se cancelam;
sobram apenas as que cruzam a fronteira.\\
(c) \textbf{Exemplo de uso.} Considere um bloco de circuito ligado ao restante
por três fios apenas. Se dois deles conduzem $\SI{5}{\ampere}$ e
$\SI{3}{\ampere}$ para dentro, o terceiro conduz necessariamente
$\SI{8}{\ampere}$ para fora --- \textbf{sem que seja preciso conhecer nada do
interior do bloco}, que pode conter dezenas de componentes.\\
\textbf{Aplicação prática:} é o que permite verificar a corrente de um circuito
inteiro medindo apenas os condutores de entrada, e é também o princípio do
\emph{dispositivo diferencial residual}, que compara fase e neutro --- qualquer
diferença indica corrente escapando por um terceiro caminho.}
\stepcounter{questao}
\resposta{$\EMF=\SI{12}{\volt}$, $r=\SI{0.5}{\ohm}$;
$I_{cc}=\SI{24}{\ampere}$}
\resolucao{(a) O modelo $V=\EMF-rI$ é uma reta. De dois pontos quaisquer:
\[
r=-\frac{\Delta V}{\Delta I}=-\frac{11{,}00-11{,}75}{2{,}00-0{,}50}
=\frac{0{,}75}{1{,}50}=\SI{0.5}{\ohm};
\]
\[
\EMF=11{,}75+0{,}5(0{,}50)=\SI{12}{\volt}.
\]
(b) \textbf{Verificação com o terceiro ponto:}
$12-0{,}5(1{,}00)=\SI{11.50}{\volt}$ \checkmark. Os três pontos são
\emph{exatamente} colineares --- o modelo linear descreve bem a fonte nessa
faixa.\\
\textbf{Nota de método:} com apenas dois pontos, sempre passa uma reta e a
linearidade não pode ser testada. O terceiro ponto é o que transforma o ajuste
em \emph{verificação}.\\
(c) $I_{cc}=\EMF/r=12/0{,}5=\SI{24}{\ampere}$.\\
\textbf{Cautela com a extrapolação:} o ensaio cobriu de $0{,}5$ a
$\SI{2}{\ampere}$, e o valor de curto está \emph{doze vezes} além do maior
ponto medido. Fontes reais aquecem, $r$ cresce e proteções atuam --- os
$\SI{24}{\ampere}$ são um \textbf{limite superior}, útil para dimensionar
proteção, não uma previsão.}
\stepcounter{questao}
\resposta{$5$ LKC $+$ $4$ LKT $=9$; as fontes de corrente fixam duas correntes
de ramo, restando $7$ incógnitas}
\resolucao{(a) LKC: $n-1=5$. LKT: $b-n+1=9-6+1=4$. Total: $9$ --- igual ao
número de correntes de ramo.\\
(b) \textbf{Cada fonte de corrente ideal fixa diretamente a corrente de seu
ramo.} Duas delas eliminam duas incógnitas, restando $7$. Em compensação, a
tensão sobre cada fonte de corrente é desconhecida, o que impede escrever a LKT
de malhas que a atravessem --- é preciso escolher malhas que as contornem, ou
tratar essas tensões como novas incógnitas.\\
\textbf{Resultado líquido:} sistema $7\times7$ em vez de $9\times9$.\\
(c) \textbf{Análise nodal:} $n-1=5$ equações, com as fontes de corrente entrando
\emph{diretamente} no vetor independente, sem nenhuma manipulação. Sistema
$5\times5$.\\
\textbf{Conclusão:} a análise nodal é claramente superior aqui. Ela é, na
verdade, uma forma organizada de aplicar a LKC com a LKT já embutida na escolha
dos potenciais --- e por isso produz o sistema menor sempre que há fontes de
corrente.}
\stepcounter{questao}
\resposta{$I=\SI{3.80}{\ampere}$ contra $\SI{5}{\ampere}$ previstos}
\resolucao{(a) A LKT dá
\[
6=(0{,}2+0{,}1I)I+1\cdot I
\;\Longrightarrow\;
0{,}1I^{2}+1{,}2I-6=0.
\]
\[
I=\frac{-1{,}2+\sqrt{1{,}44+2{,}4}}{0{,}2}
=\frac{-1{,}2+1{,}960}{0{,}2}=\SI{3.80}{\ampere}.
\]
\emph{(Verificação: $r=0{,}2+0{,}380=\SI{0.58}{\ohm}$ e
$6/(0{,}58+1)=\SI{3.80}{\ampere}$ \checkmark.)}\\
(b) \textbf{Modelo linear:} $I=6/1{,}2=\SI{5}{\ampere}$ --- \textbf{$32\%$
acima} do valor real.\\
\textbf{Duas lições.}\\
\emph{(i)} As \textbf{leis de Kirchhoff continuam válidas} --- elas não exigem
linearidade. O que muda é que a equação resultante deixa de ser linear e passa a
exigir solução de segundo grau.\\
\emph{(ii)} O modelo de $r$ constante é uma \emph{aproximação de pequenos
sinais}. Ele descreve bem a fonte perto do ponto onde $r$ foi medido, e falha
progressivamente longe dele. Baterias reais comportam-se assim, e é por isso
que fabricantes especificam $r$ para uma corrente de ensaio definida.}
\stepcounter{questao}
\resposta{$I\le\SI{1.2}{\ampere}$; $R\ge\SI{9.5}{\ohm}$;
$P=\SI{13.68}{\watt}$}
\resolucao{(a) De $V=\EMF-rI\ge11{,}4$:
\[
0{,}5I\le0{,}6
\;\Longrightarrow\;
I\le\SI{1.2}{\ampere}.
\]
(b) $R=V/I=11{,}4/1{,}2=\SI{9.5}{\ohm}$ --- cargas menores exigiriam mais
corrente e derrubariam a tensão.\\
(c) $P=VI=11{,}4(1{,}2)=\SI{13.68}{\watt}$, com rendimento
$9{,}5/10=95\%$.\\
\textbf{Contraste com a máxima transferência:} a potência máxima ocorreria em
$R=\SI{0.5}{\ohm}$, valendo $\SI{72}{\watt}$ --- mas com tensão terminal de
apenas $\SI{6}{\volt}$, muito abaixo do requisito.\\
\textbf{A restrição de tensão é quase sempre a que manda} em eletrônica, porque
circuitos alimentados simplesmente não funcionam abaixo de sua tensão mínima. A
saída de engenharia é reduzir $r$, não aceitar tensão menor.}
\stepcounter{questao}
\resposta{$I=\SI{1.54}{\ampere}$ (fonte dependente em oposição)}
\resolucao{(a) Percorrendo a malha, com a fonte dependente \emph{opondo-se}:
\[
20=2I+8I+3I=13I
\;\Longrightarrow\;
I=\frac{20}{13}=\SI{1.54}{\ampere}.
\]
\emph{(Se a fonte dependente auxiliasse, viria $20=10I-3I=7I$ e
$I=\SI{2.86}{\ampere}$ --- o sentido da fonte dependente altera
completamente o resultado, e por isso a convenção de polaridade precisa estar
explícita no esquema.)}\\
(b) \textbf{As leis de Kirchhoff decorrem de conservação de carga e de
energia}, e nenhuma dessas conservações depende de os elementos serem lineares
ou independentes. A LKT continua sendo "a soma das quedas em uma malha é nula";
a LKC continua sendo "o que entra em um nó, sai".\\
\textbf{O que muda:} a fonte dependente introduz uma \textbf{equação de
restrição} adicional, ligando seu valor a uma variável do circuito. Essa
equação é resolvida \emph{junto} com as de Kirchhoff.\\
\textbf{Consequência estrutural:} o termo migra para o lado das incógnitas, o
que pode tornar a matriz assimétrica --- como se observa na análise nodal e na
de malhas. A reciprocidade se perde; as leis, não.}
\stepcounter{questao}
\resposta{$I=\SI{2.4}{\ampere}$; $V_{carga}=\SI{22.8}{\volt}$;
$V_{fonte}=\SI{23.52}{\volt}$}
\resolucao{(a)
\[
I=\frac{24}{0{,}2+0{,}3+9{,}5}=\frac{24}{10}=\SI{2.4}{\ampere};
\]
\[
V_{carga}=9{,}5(2{,}4)=\SI{22.8}{\volt};
\qquad
V_{fonte}=24-0{,}2(2{,}4)=\SI{23.52}{\volt}.
\]
A diferença de $\SI{0.72}{\volt}$ entre os dois é exatamente a queda no cabo.\\
(b) \textbf{Distribuição das perdas:}
\begin{center}
\begin{tabular}{lrr}
\hline
Elemento & Potência & Fração\\
\hline
Carga ($\SI{9.5}{\ohm}$) & $\SI{54.72}{\watt}$ & $95{,}0\%$\\
Cabo ($\SI{0.3}{\ohm}$) & $\SI{1.73}{\watt}$ & $3{,}0\%$\\
Interna ($\SI{0.2}{\ohm}$) & $\SI{1.15}{\watt}$ & $2{,}0\%$\\
\hline
Total & $\SI{57.6}{\watt}$ & $100\%$\\
\hline
\end{tabular}
\end{center}
\textbf{Observação de diagnóstico:} medindo apenas nos terminais da fonte
($\SI{23.52}{\volt}$), a instalação pareceria melhor do que é. A tensão que
importa é a que chega à \emph{carga} --- e ela está $5\%$ abaixo da f.e.m.\\
\textbf{Regra de medição:} sempre meça no ponto de uso, não na fonte. As duas
leituras diferem exatamente pela queda no caminho.}
\stepcounter{questao}
\resposta{$\EMF=V_{oc}$ e $r=V_{oc}/I_{cc}$; o curto é destrutivo em fontes de
baixa impedância}
\resolucao{(a) \textbf{Em vazio} ($I=0$): $V_{oc}=\EMF$ --- a leitura direta,
desde que o voltímetro tenha impedância muito maior que $r$.\\
\textbf{Em curto} ($V=0$): $\EMF=rI_{cc}$, logo
\[
r=\frac{V_{oc}}{I_{cc}}.
\]
(b) \textbf{Riscos do curto.}
\begin{itemize}
\item Uma pilha AA de $\SI{1.5}{\volt}$ com $r=\SI{0.25}{\ohm}$ entrega
$\SI{6}{\ampere}$ e dissipa $\SI{9}{\watt}$ internamente --- aquecimento
rápido e vazamento.
\item Uma bateria de automóvel ($r\approx\SI{5}{\milli\ohm}$) chega a
$\SI{2500}{\ampere}$, com centelhamento, fusão de terminais e risco de
explosão por ignição do hidrogênio liberado.
\item Baterias de lítio podem entrar em fuga térmica e incendiar.
\item Além disso, $r$ \emph{aumenta} durante o curto, de modo que a leitura
subestima a corrente inicial e superestima $r$.
\end{itemize}
(c) \textbf{Alternativa segura: dois pontos de carga.} Meça $(I_1,V_1)$ e
$(I_2,V_2)$ com resistências conhecidas e moderadas, e ajuste a reta:
\[
r=-\frac{V_2-V_1}{I_2-I_1};
\qquad
\EMF=V_1+rI_1.
\]
\textbf{Comparação.} O método de dois pontos é seguro, não perturba a fonte, e
com três ou mais pontos ainda \emph{testa} a linearidade do modelo --- coisa que
o ensaio de curto não faz. O ensaio em curto só se justifica em fontes
projetadas para suportá-lo, com limitação eletrônica de corrente.}
\stepcounter{questao}
\resposta{$\sum_k v_ki_k=0$ --- consequência direta de LKC e LKT}
\resolucao{(a) \textbf{Convenção do receptor:} para cada ramo $k$, adote a
corrente $i_k$ entrando pelo terminal de maior potencial. A potência
$p_k=v_ki_k$ é então \emph{positiva} se o elemento absorve e \emph{negativa} se
fornece.\\
(b) \textbf{Demonstração.} Escreva cada tensão de ramo como diferença de
potenciais nodais: $v_k=V_a-V_b$, onde $a$ e $b$ são os nós do ramo $k$. Então
\[
\sum_k p_k=\sum_k (V_a-V_b)\,i_k.
\]
Reagrupando por \emph{nó} em vez de por ramo, cada potencial $V_n$ aparece
multiplicado pela soma algébrica das correntes que chegam ao nó $n$:
\[
\sum_k p_k=\sum_n V_n\left(\sum_{k\in n}\pm i_k\right)=\sum_n V_n\cdot 0=0,
\]
pois o parêntese é exatamente a \textbf{LKC} do nó $n$.\\
\emph{(A escrita $v_k=V_a-V_b$ já embute a \textbf{LKT}: ela pressupõe que o
potencial é bem definido em cada nó.)}\\
(c) \textbf{Independência da natureza dos elementos.} A demonstração usou
\emph{apenas} as duas leis de Kirchhoff e a definição $p=vi$. Em nenhum momento
foi preciso saber a relação entre $v_k$ e $i_k$ de cada elemento.\\
Portanto o resultado vale para resistores, fontes, diodos, capacitores,
indutores, elementos ativos --- lineares ou não, com ou sem memória. É o
\textbf{teorema de Tellegen}, e sua generalidade é notável: ele é uma
propriedade da \emph{topologia} do circuito, não dos componentes.\\
\textbf{Uso prático:} verificar o balanço de potência é o teste de consistência
mais completo de qualquer solução --- se ele não fecha, há erro em algum lugar.}
\stepcounter{questao}
\resposta{Usar supernó: uma LKC envolvendo $A$ e $B$, mais a restrição
$V_A-V_B=\EMF$}
\resolucao{(a) A corrente que atravessa uma fonte de tensão \textbf{ideal} é
determinada pelo circuito, não pela fonte --- não existe relação
$i=f(v)$ para ela. Ao escrever a LKC em $A$, essa corrente desconhecida aparece
como uma incógnita a mais, e o mesmo ocorre em $B$. Duas equações, três
incógnitas.\\
(b) \textbf{Formulação por supernó.} Englobe $A$ e $B$ em uma única superfície
fechada. Pela LKC generalizada, a corrente interna da fonte \emph{se cancela}
--- ela entra em um nó e sai do outro:
\[
\sum_{\text{ramos que cruzam a fronteira}} i=0.
\]
Essa é \emph{uma} equação. A segunda vem da própria fonte:
\[
V_A-V_B=\EMF.
\]
Duas equações, duas incógnitas ($V_A$ e $V_B$). \textbf{Sistema fechado.}\\
(c) \textbf{Corrente na fonte:} obtida \emph{depois}, aplicando a LKC a
\emph{um} dos dois nós isoladamente. Com $V_A$ e $V_B$ já conhecidos, todas as
demais correntes daquele nó são calculáveis, e a da fonte é o que falta para
fechar o balanço.\\
\textbf{Observação:} se um dos terminais da fonte for o próprio nó de
referência, o problema desaparece --- o outro potencial fica \emph{conhecido} e
deixa de ser incógnita. \textbf{Escolher o terra em um terminal de fonte é a
simplificação mais barata que existe.}}
\stepcounter{questao}
\resposta{Critério: resíduo compatível com $u_c\approx\sqrt{\sum u_k^{2}}$;
$\SI{0.35}{\ampere}$ indica erro real}
\resolucao{(a) \textbf{Procedimento.} Meça as três correntes
\emph{simultaneamente}, com três amperímetros --- ou, se houver apenas um,
faça as três medições sem alterar o circuito entre elas. Registre sentidos e
adote uma convenção única (por exemplo, positivo entrando).\\
\textbf{Cuidado essencial:} inserir o amperímetro \emph{altera} o circuito, pois
acrescenta sua resistência interna ao ramo. Com três inserções sucessivas, o
circuito muda três vezes e a verificação perde sentido. Use amperímetro de baixa
resistência ou, melhor, alicate amperímetro, que não abre o ramo.\\
(b) \textbf{Propagação.} Suponha
$I_1=(5{,}02\pm0{,}05)\,\si{\ampere}$,
$I_2=(3{,}01\pm0{,}03)\,\si{\ampere}$ entrando, e $I_3$ saindo. A previsão é
$I_1+I_2=\SI{8.03}{\ampere}$, com incerteza combinada
\[
u_c=\sqrt{0{,}05^{2}+0{,}03^{2}}=\SI{0.058}{\ampere}.
\]
Incluindo a incerteza da própria medição de $I_3$ (digamos
$\SI{0.08}{\ampere}$), o resíduo esperado tem incerteza
$\sqrt{0{,}058^{2}+0{,}08^{2}}=\SI{0.099}{\ampere}$.\\
\textbf{Critério de aceitação:} $|I_1+I_2-I_3|\le2u\approx
\SI{0.20}{\ampere}$ (dois desvios, cerca de $95\%$ de confiança).\\
(c) \textbf{Resíduo de $\SI{0.35}{\ampere}$} é $3{,}5$ vezes a incerteza
combinada --- \textbf{não é compatível com erro aleatório}. Causas prováveis, em
ordem:
\begin{itemize}
\item \textbf{Erro de sinal ou de sentido} em uma das leituras --- o mais comum.
\item \textbf{Um quarto ramo não contabilizado}, inclusive uma fuga para o terra
ou pela bancada.
\item \textbf{Instrumento descalibrado} ou operando fora de escala.
\item \textbf{Circuito alterado entre medições} (aquecimento, contato
intermitente).
\end{itemize}
\textbf{Conclusão importante:} a LKC \emph{não} pode ser violada --- ela é
conservação de carga. Um resíduo significativo é sempre erro de medição ou de
modelo, nunca da lei. O ensaio, portanto, verifica o \emph{procedimento}, e não
a física.}
\stepcounter{questao}
\resposta{$I=\SI{0.76}{\ampere}$; a célula gasta apresenta
$\SI{-0.24}{\volt}$ --- está sendo \textbf{invertida}}
\resolucao{(a)
\[
\EMF_{total}=3(1{,}5)+0{,}9=\SI{5.4}{\volt};
\qquad
r_{total}=3(0{,}2)+1{,}5=\SI{2.1}{\ohm};
\]
\[
I=\frac{5{,}4}{2{,}1+5}=\frac{5{,}4}{7{,}1}=\SI{0.76}{\ampere}.
\]
\textbf{Tensão de cada célula boa:} $1{,}5-0{,}2(0{,}76)=\SI{1.35}{\volt}$.\\
\textbf{Tensão da célula gasta:}
$0{,}9-1{,}5(0{,}76)=0{,}9-1{,}14=\SI{-0.24}{\volt}$.\\
(b) \textbf{A tensão da célula gasta é negativa} --- ou seja, ela está sendo
percorrida por corrente no sentido de \emph{carga}, forçada pelas três células
boas. Ela deixou de ser fonte e passou a ser carga, dissipando
$1{,}5(0{,}76)^{2}=\SI{0.87}{\watt}$ internamente.\\
\textbf{Em células não recarregáveis}, isso provoca a chamada \textbf{inversão
de polaridade}: reações internas indesejadas, geração de gás, aquecimento e
vazamento de eletrólito --- exatamente o que se observa quando uma pilha
"vaza" em um controle remoto esquecido ligado.\\
(c) \textbf{Implicações práticas.}
\begin{itemize}
\item \textbf{Nunca misture pilhas novas e usadas}, nem de marcas ou tecnologias
diferentes, em série. A mais fraca será invertida pelas demais.
\item \textbf{Substitua sempre o conjunto completo.}
\item Em baterias de lítio em série, o problema é grave a ponto de exigir
\textbf{circuito de balanceamento} (BMS), que monitora cada célula
individualmente e impede que qualquer uma seja levada à inversão.
\end{itemize}
\textbf{Diagnóstico:} a célula gasta domina o conjunto por sua \emph{resistência
interna}, não por sua f.e.m. --- $\SI{1.5}{\ohm}$ contra
$\SI{0.6}{\ohm}$ das outras três somadas.}
\stepcounter{questao}
\resposta{Duas regiões: fonte de tensão até $\SI{2}{\ampere}$, fonte de corrente
depois}
\resolucao{(a) \textbf{Característica em dois trechos.}
\begin{itemize}
\item \textbf{Região de tensão constante} ($I<\SI{2}{\ampere}$):
$V=20-0{,}1I$ --- reta quase horizontal.
\item \textbf{Região de corrente constante} ($V<\SI{19.8}{\volt}$):
$I=\SI{2}{\ampere}$ --- reta vertical.
\end{itemize}
O "joelho" fica em $(\SI{2}{\ampere};\ \SI{19.8}{\volt})$, com
$R_{joelho}=\SI{9.9}{\ohm}$.\\
(b)
\begin{center}
\begin{tabular}{lccl}
\hline
$R_L$ & $I$ & $V$ & Região\\
\hline
$\SI{50}{\ohm}$ & $\SI{0.399}{\ampere}$ & $\SI{19.96}{\volt}$ & tensão\\
$\SI{10}{\ohm}$ & $\SI{1.98}{\ampere}$ & $\SI{19.80}{\volt}$ & tensão (limite)\\
$\SI{2}{\ohm}$ & $\SI{2.00}{\ampere}$ & $\SI{4.00}{\volt}$ & corrente\\
\hline
\end{tabular}
\end{center}
Note o terceiro caso: o modelo linear preveria
$20/2{,}1=\SI{9.5}{\ampere}$ --- quase cinco vezes o real.\\
(c) \textbf{Por que o modelo falha.} A fonte real é um modelo \emph{linear} de
dois parâmetros, e sua característica é uma \emph{única} reta. A limitação
eletrônica introduz um segundo trecho, tornando a característica
\textbf{quebrada} --- e nenhuma reta descreve as duas regiões.\\
\textbf{Como proceder corretamente:}
\begin{itemize}
\item Determinar em qual região a carga opera, comparando $R_L$ com
$R_{joelho}=\SI{9.9}{\ohm}$.
\item $R_L>R_{joelho}$: use o modelo de \emph{fonte de tensão} $(\EMF,r)$.
\item $R_L<R_{joelho}$: use o modelo de \emph{fonte de corrente} ideal de
$\SI{2}{\ampere}$, e a tensão sai de $V=R_LI$.
\end{itemize}
\textbf{As leis de Kirchhoff seguem válidas nas duas regiões} --- o que muda é
apenas a relação constitutiva da fonte. É o mesmo procedimento usado com
diodos: identificar o trecho e aplicar o modelo correspondente.}
\stepcounter{questao}
\resposta{$I_1=\SI{3.6}{\ampere}$, $I_2=\SI{-1.2}{\ampere}$,
$I_R=\SI{2.4}{\ampere}$}
\resolucao{(a) Chamando $V$ a tensão do nó comum:
\[
\frac{V-18}{1}+\frac{V-12}{2}+\frac{V}{6}=0.
\]
Multiplicando por $6$: $6V-108+3V-36+V=0\Rightarrow10V=144
\Rightarrow V=\SI{14.4}{\volt}$.
\[
I_1=\frac{18-14{,}4}{1}=\SI{3.6}{\ampere};
\quad
I_2=\frac{12-14{,}4}{2}=\SI{-1.2}{\ampere};
\quad
I_R=\frac{14{,}4}{6}=\SI{2.4}{\ampere}.
\]
\textbf{LKC:} $3{,}6-1{,}2=2{,}4$ \checkmark.\\
(b) \textbf{Balanço.}
\begin{center}
\begin{tabular}{lr}
\hline
$\EMF_1$ fornece $18(3{,}6)$ & $\SI{64.8}{\watt}$\\
$r_1$ dissipa $1(3{,}6)^{2}$ & $\SI{12.96}{\watt}$\\
$R$ dissipa $6(2{,}4)^{2}$ & $\SI{34.56}{\watt}$\\
$\EMF_2$ absorve $12(1{,}2)$ & $\SI{14.4}{\watt}$\\
$r_2$ dissipa $2(1{,}2)^{2}$ & $\SI{2.88}{\watt}$\\
\hline
Total consumido & $\SI{64.8}{\watt}$\\
\hline
\end{tabular}
\end{center}
Fecha exatamente \checkmark.\\
(c) \textbf{Papéis.} A fonte 1 é a única \emph{geradora}: ela alimenta a carga
$R$ \emph{e} carrega a fonte 2. A fonte 2 comporta-se como \textbf{receptor},
porque sua f.e.m. ($\SI{12}{\volt}$) é menor que a tensão do nó
($\SI{14.4}{\volt}$) --- e corrente sempre entra por onde o potencial externo é
maior.\\
\textbf{Critério geral:} compare a f.e.m. de cada fonte com a tensão do nó a
que está ligada. Maior $\Rightarrow$ fornece; menor $\Rightarrow$ recebe. O
sinal da corrente calculada confirma automaticamente.}

\gabarito[2.4 Fontes reais e Leis de Kirchhoff]

\section{Associação de Resistores e Transformações}

\subsection{Circuitos em série}
\stepcounter{questao}
\resposta{(b)}
\resolucao{Só existe um caminho para a corrente, logo ela é a mesma em todos os
elementos --- (a) e (d) descrevem o circuito paralelo. A resistência equivalente
é a \emph{soma}, portanto maior que a maior delas, o que elimina (c).
A potência $P=R\,I^2$ depende de $R$: só é igual se os resistores forem iguais,
o que elimina (e).}
\stepcounter{questao}
\resposta{(a) \SI{20}{\ohm}; (b) \SI{2}{\ampere}; (c) \SI{8}{\volt}, \SI{12}{\volt} e \SI{20}{\volt}}
\resolucao{(a) $\Req = 4+6+10 = \SI{20}{\ohm}$.\\
(b) $I = \dfrac{U}{\Req} = \dfrac{40}{20} = \SI{2}{\ampere}$.\\
(c) $U_1 = 4\cdot 2 = \SI{8}{\volt}$, $U_2 = 6\cdot 2 = \SI{12}{\volt}$,
$U_3 = 10\cdot 2 = \SI{20}{\volt}$.\\
\textbf{Verificação:} $8+12+20 = \SI{40}{\volt}$ --- confere com a tensão da fonte.
Repare que a maior tensão cai sobre o maior resistor.}
\stepcounter{questao}
\resposta{(a) \SI{5}{\ohm}, \SI{2}{\ampere}; (b) \SI{10}{\ohm}, \SI{1.5}{\ampere};
(c) \SI{10}{\ohm}, \SI{2.5}{\ampere}; (d) \SI{10}{\ohm}, \SI{3}{\ampere}}
\resolucao{Em todos os casos, $\Req$ é a soma das resistências e
$I = E/\Req$:\\
(a) $\Req = 3+2 = \SI{5}{\ohm}$, $I = 10/5 = \SI{2}{\ampere}$.\\
(b) $\Req = 4+6 = \SI{10}{\ohm}$, $I = 15/10 = \SI{1.5}{\ampere}$.\\
(c) $\Req = 4{,}5+3+2{,}5 = \SI{10}{\ohm}$, $I = 25/10 = \SI{2.5}{\ampere}$.\\
(d) $\Req = 2{,}5+4+1+2{,}5 = \SI{10}{\ohm}$, $I = 30/10 = \SI{3}{\ampere}$.}
\stepcounter{questao}
\resposta{(d)}
\resolucao{Para $n$ resistores iguais em série, $\Req = n\,R$, logo
$n = \dfrac{105}{15} = 7$.}
\stepcounter{questao}
\resposta{$I = \SI{4}{\ampere}$; $U_{R_1} = \SI{10}{\volt}$}
\resolucao{Percorrendo a malha, $E_1$ impulsiona a corrente no sentido adotado e
$E_2$ se opõe. Pela lei das malhas:
\[
E_1 - E_2 = I\,(R_1+R_2+R_3)
\quad\Rightarrow\quad
I = \frac{30-10}{2{,}5+1{,}5+1} = \frac{20}{5} = \SI{4}{\ampere}.
\]
$U_{R_1} = R_1 I = 2{,}5\cdot 4 = \SI{10}{\volt}$.\\
\textbf{Observação:} se $E_2$ fosse maior que $E_1$, a corrente resultaria
negativa --- o que significa apenas que ela circula no sentido oposto ao adotado.
Nesse caso $E_2$ estaria \emph{sendo carregada} por $E_1$.}
\stepcounter{questao}
\resposta{(a) \SI{2}{\ampere}; (b) \SI{11}{\volt}; (c) \SI{2}{\watt}}
\resolucao{(a) A resistência interna está \emph{em série} com a carga:
\[
I=\frac{\EMF}{R+r}=\frac{12}{5{,}5+0{,}5}=\frac{12}{6}=\SI{2}{\ampere}.
\]
(b) $U = \EMF - r\,I = 12 - 0{,}5\cdot 2 = \SI{11}{\volt}$
(ou, equivalentemente, $U = R\,I = 5{,}5\cdot 2 = \SI{11}{\volt}$).\\
(c) $P_r = r\,I^2 = 0{,}5\cdot 2^2 = \SI{2}{\watt}$ --- energia perdida dentro da
própria bateria, que é o que a faz aquecer em curtos-circuitos.}
\stepcounter{questao}
\resposta{(a) $I=\SI{0.5}{\ampere}$ e $U=\SI{20}{\volt}$ por lâmpada;
(b) todas apagam}
\resolucao{(a) $\Req = 6\cdot 40 = \SI{240}{\ohm}$;
$I = \dfrac{120}{240} = \SI{0.5}{\ampere}$;
$U = 40\cdot 0{,}5 = \SI{20}{\volt}$ em cada lâmpada
(confira: $6\cdot 20 = \SI{120}{\volt}$).\\
(b) O filamento aberto interrompe o \emph{único} caminho da corrente:
$I = 0$ e todas as lâmpadas apagam. É exatamente esse o inconveniente dos antigos
cordões de luzes de Natal ligados em série.}
\stepcounter{questao}
\resposta{(b)}
\resolucao{Se $U_2 = \SI{20}{\volt}$, sobram $120-20 = \SI{100}{\volt}$ para
$R_1$ e $R_3$, que somam \SI{50}{\ohm}. Como a corrente é a mesma em toda a série:
\[
I=\frac{100}{50}=\SI{2}{\ampere}
\qquad\Rightarrow\qquad
R_2=\frac{U_2}{I}=\frac{20}{2}=\SI{10}{\ohm}.
\]}
\stepcounter{questao}
\resposta{(a) \SI{2.4}{\ampere}; (b) \SI{0.96}{\volt}; (c) \SI{23.04}{\volt};
(d) \SI{96}{\percent}}
\resolucao{Os dois condutores (ida e volta) ficam \emph{em série} com a carga:
\[
\Req = 0{,}2+9{,}6+0{,}2=\SI{10}{\ohm}
\quad\Rightarrow\quad
I=\frac{24}{10}=\SI{2.4}{\ampere}.
\]
(b) $U_{\text{cabo}} = 0{,}4\cdot 2{,}4 = \SI{0.96}{\volt}$.\\
(c) $U_{\text{carga}} = 24-0{,}96 = \SI{23.04}{\volt}$ (ou $9{,}6\cdot2{,}4$).\\
(d) Como a corrente é a mesma em tudo,
$\eta = \dfrac{R_c}{\Req} = \dfrac{9{,}6}{10} = 0{,}96 = \SI{96}{\percent}$.\\
\textbf{Por que isso importa:} a perda cresce com $I^2$. Se a mesma potência fosse
entregue com o dobro da tensão (e metade da corrente), a perda no cabo cairia à
quarta parte --- é o argumento a favor da transmissão em alta tensão.}
\stepcounter{questao}
\resposta{(a) 5 resistores (saída sobre 2 deles); (b) \SI{0.1}{\ampere}}
\resolucao{(a) Com $n$ resistores iguais em série, a tensão sobre $k$ deles é
\[
U_k = E\,\frac{k}{n} \quad\Rightarrow\quad 20 = 50\cdot\frac{k}{n}
\quad\Rightarrow\quad \frac{k}{n}=\frac{2}{5}.
\]
A menor solução inteira é $k=2$ e $n=5$: cinco resistores em série, com a saída
tomada sobre dois deles.\\
(b) $\Req = 5\cdot 100 = \SI{500}{\ohm}$ e
$I = \dfrac{50}{500} = \SI{0.1}{\ampere}$.\\
\textbf{Cuidado prático:} essa tensão de \SI{20}{\volt} só se mantém \emph{em
vazio}. Ao ligar uma carga em paralelo com os dois resistores, a tensão de saída
cai --- por isso divisores resistivos não servem como fonte de alimentação.}
\stepcounter{questao}
\resposta{(a) $R_{\text{linha}} \le \SI{0.417}{\ohm}$;
(b) $A \ge \SI{163}{\milli\meter\squared}$; (c) $\eta = \SI{96}{\percent}$,
$P_{\text{perdida}}\approx\SI{221}{\watt}$}
\resolucao{\textbf{(a)} A linha e a carga estão em série, formando um divisor de
tensão. Para a queda na linha ser no máximo \SI{4}{\percent}:
\[
\frac{R_{\text{linha}}}{R_{\text{linha}}+R_L}\le 0{,}04
\;\Longrightarrow\;
R_{\text{linha}} \le \frac{0{,}04}{0{,}96}\,R_L = \frac{10}{24}
\approx\SI{0.417}{\ohm}.
\]
\textbf{(b)} O comprimento \emph{total} de condutor é o dobro da distância
(ida e volta): $L = \SI{4000}{\meter}$. De $R=\rho L/A$:
\[
A \ge \frac{\rho L}{R_{\text{linha}}}
=\frac{(\pot{1.7}{-8})(4000)}{0{,}417}
\approx\pot{1.63}{-4}\,\si{\meter\squared}=\SI{163}{\milli\meter\squared}.
\]
Comercialmente, adotaria-se $\SI{185}{\milli\meter\squared}$ (valor normalizado
imediatamente superior).

\textbf{(c)} Com $R_{\text{linha}} = \SI{0.417}{\ohm}$:
\[
I=\frac{240}{10{,}417}\approx\SI{23}{\ampere},
\qquad
\eta=\frac{R_L}{R_L+R_{\text{linha}}}=\frac{10}{10{,}417}\approx\SI{96}{\percent},
\]
\[
P_{\text{perdida}}=R_{\text{linha}}I^{2}\approx 0{,}417\cdot23^{2}
\approx\SI{221}{\watt}.
\]
\textbf{Reflexão de projeto:} note o custo do critério. Exigir apenas
\SI{4}{\percent} de queda impôs um cabo de $\SI{163}{\milli\meter\squared}$ ---
muito caro. Se a mesma potência fosse transmitida em $\SI{2400}{\volt}$
(dez vezes mais), a corrente cairia a $1/10$ e a perda ($\propto I^2$) a
$1/100$, permitindo um cabo cem vezes menor. \textbf{É por isso que se transmite
energia em alta tensão} --- a economia está no cobre, não na energia.}
\stepcounter{questao}
\resposta{$R_{eq}=\SI{3.87}{\kilo\ohm}$}
\resolucao{Converta tudo para a mesma unidade \emph{antes} de somar:
\[
R_{eq}=1200+470+2200=\SI{3870}{\ohm}=\SI{3.87}{\kilo\ohm}.
\]
Em série a soma é direta --- não há inversos envolvidos. Misturar
$\si{\kilo\ohm}$ com $\si{\ohm}$ sem converter é o erro mais frequente aqui.}
\stepcounter{questao}
\resposta{$R_3=\SI{20}{\ohm}$}
\resolucao{
\[
R_3=R_{eq}-R_1-R_2=75-22-33=\SI{20}{\ohm}.
\]
A soma série é reversível: conhecidos o total e todas as parcelas menos uma,
a que falta sai por subtração --- diferente do paralelo, em que seria preciso
subtrair \emph{condutâncias}.}
\stepcounter{questao}
\resposta{$I=\SI{2}{\ampere}$}
\resolucao{
\[
R_{eq}=4+8=\SI{12}{\ohm}
\;\Longrightarrow\;
I=\frac{V}{R_{eq}}=\frac{24}{12}=\SI{2}{\ampere}.
\]
Essa corrente é a \emph{mesma} nos dois resistores --- não existe "corrente de
$R_1$" diferente da "corrente de $R_2$" em uma série.}
\stepcounter{questao}
\resposta{$U=\SI{34}{\volt}$}
\resolucao{Lei de Ohm aplicada ao elemento:
\[
U=R\,I=68\cdot0{,}5=\SI{34}{\volt}.
\]
Como $I$ é comum, a queda de cada resistor é \emph{diretamente proporcional} ao
seu valor: o maior resistor sempre fica com a maior parcela da tensão.}
\stepcounter{questao}
\resposta{(b)}
\resolucao{$R_{eq}$ só pode crescer ao se acrescentar um resistor em série
($R_{eq}=\sum R_k$, com todas as parcelas positivas). Com $V$ fixo,
$I=V/R_{eq}$ necessariamente diminui. Consequência prática: em uma cadeia de
lâmpadas, acrescentar mais uma faz \emph{todas} brilharem menos.}
\stepcounter{questao}
\resposta{(d)}
\resolucao{A soma $R_{eq}=15+45=\SI{60}{\ohm}$ é comutativa, a corrente
$I=V/R_{eq}$ depende só de $R_{eq}$, e cada queda $U_k=R_kI$ depende só do
próprio resistor e de $I$. Nada muda.\\
O que muda é o \emph{potencial} de cada ponto intermediário em relação ao terra
--- relevante apenas se algum ponto do circuito estiver aterrado ou se houver
medição referida a um nó específico.}
\stepcounter{questao}
\resposta{$I=\SI{2}{\ampere}$; $U_1=\SI{10}{\volt}$, $U_2=\SI{20}{\volt}$,
$U_3=\SI{30}{\volt}$}
\resolucao{
\[
R_{eq}=5+10+15=\SI{30}{\ohm}
\;\Longrightarrow\;
I=\frac{60}{30}=\SI{2}{\ampere}.
\]
\[
U_1=5\cdot2=10,\quad U_2=10\cdot2=20,\quad U_3=15\cdot2=\SI{30}{\volt}.
\]
\textbf{Verificação (LKT):} $10+20+30=\SI{60}{\volt}$ --- fecha com a fonte.
Note a proporcionalidade: as quedas estão entre si como $5:10:15=1:2:3$.}
\stepcounter{questao}
\resposta{$\SI{20}{\watt}$, $\SI{40}{\watt}$, $\SI{60}{\watt}$; total
$\SI{120}{\watt}$}
\resolucao{Com $I=\SI{2}{\ampere}$ comum, use $P=RI^{2}$:
\[
P_1=5\cdot4=20,\quad P_2=10\cdot4=40,\quad P_3=15\cdot4=\SI{60}{\watt}.
\]
Fonte: $P=VI=60\cdot2=\SI{120}{\watt}=20+40+60$ \checkmark.\\
\textbf{Regra da série:} como $I$ é comum, $P\propto R$. O maior resistor
dissipa mais --- exatamente o oposto do que ocorre no paralelo, onde $V$ é
comum e $P\propto 1/R$.}
\stepcounter{questao}
\resposta{$V=\SI{36}{\volt}$}
\resolucao{A medida sobre um único elemento revela a corrente da cadeia
inteira:
\[
I=\frac{18}{90}=\SI{0.2}{\ampere}.
\]
\[
R_{eq}=30+60+90=\SI{180}{\ohm}
\;\Longrightarrow\;
V=R_{eq}I=180\cdot0{,}2=\SI{36}{\volt}.
\]
Este é o procedimento padrão de bancada: medir a queda em um resistor conhecido
transforma-o em um \emph{sensor de corrente} (resistor \emph{shunt}), sem
abrir o circuito para inserir o amperímetro.}
\stepcounter{questao}
\resposta{A de $\SI{960}{\ohm}$; $P=\SI{9.6}{\watt}$ contra $\SI{2.4}{\watt}$}
\resolucao{
\[
I=\frac{120}{240+960}=\frac{120}{1200}=\SI{0.1}{\ampere}
\]
\[
P_{240}=240\cdot(0{,}1)^{2}=\SI{2.4}{\watt},\qquad
P_{960}=960\cdot(0{,}1)^{2}=\SI{9.6}{\watt}.
\]
Como a corrente é comum, $P\propto R$: a lâmpada de \emph{maior} resistência
brilha mais. Isso contraria a intuição de quem associa "mais potência" a "menos
resistência" --- válido apenas quando a \emph{tensão} é comum, isto é, no
paralelo. Uma lâmpada de $\SI{100}{\watt}$ e outra de $\SI{25}{\watt}$ ligadas
em série: a de $\SI{25}{\watt}$ é que brilha mais.}
\stepcounter{questao}
\resposta{$R_s=\SI{100}{\ohm}$; $P_s=\SI{2.25}{\watt}$}
\resolucao{
\[
R_{eq}=\frac{V}{I}=\frac{24}{0{,}150}=\SI{160}{\ohm}
\;\Longrightarrow\;
R_s=160-60=\SI{100}{\ohm}.
\]
\[
P_s=R_sI^{2}=100\cdot(0{,}15)^{2}=\SI{2.25}{\watt}.
\]
Especifique pelo menos $\SI{5}{\watt}$ (fator $2\times$ de folga). Observe que o
resistor limitador dissipa $\SI{2.25}{\watt}$ dos $\SI{3.6}{\watt}$ totais ---
$62\%$ da energia vira calor. Limitar corrente com resistor é simples, mas
nunca eficiente.}
\stepcounter{questao}
\resposta{Corrente e potência totais iguais ($\SI{0.5}{\ampere}$,
$\SI{25}{\watt}$); resistor mais solicitado: $\SI{12.5}{\watt}$ em (A) e
$\SI{22.5}{\watt}$ em (B)}
\resolucao{Em ambas: $I=50/100=\SI{0.5}{\ampere}$ e
$P_{tot}=50\cdot0{,}5=\SI{25}{\watt}$.\\
(A) $P=50\cdot0{,}25=\SI{12.5}{\watt}$ em cada.\\
(B) $P_{10}=\SI{2.5}{\watt}$ e $P_{90}=\SI{22.5}{\watt}$.\\
\textbf{Conclusão de projeto:} $R_{eq}$ igual \emph{não} significa
especificação igual. A montagem (B) exige um resistor de pelo menos
$\SI{25}{\watt}$, enquanto (A) se resolve com dois de $\SI{15}{\watt}$.
Distribuir a resistência em parcelas iguais distribui também o aquecimento.}
\stepcounter{questao}
\resposta{$I=\SI{250}{\milli\ampere}$; o fusível deve ficar acima da corrente
normal e abaixo da corrente de falha}
\resolucao{
\[
I=\frac{12}{48}=\SI{0.25}{\ampere}=\SI{250}{\milli\ampere}.
\]
Um fusível de $\SI{500}{\milli\ampere}$ dá folga de $2\times$: não atua em
transitórios normais de energização, mas interrompe bem antes das correntes de
curto. \textbf{Estimativa da falha:} se metade da cadeia entrar em curto,
$R=\SI{24}{\ohm}$ e $I=\SI{500}{\milli\ampere}$ --- exatamente no limiar; um
curto franco ($R\to0$) levaria a corrente a valores muito maiores, e o fusível
atua rapidamente. Um fusível \emph{muito} próximo da corrente nominal produz
desarmes espúrios; muito acima, não protege.}
\stepcounter{questao}
\resposta{de $\SI{0.2}{\ampere}$ a $\SI{0.6}{\ampere}$}
\resolucao{
\[
I_{\max}=\frac{30}{50+0}=\SI{0.6}{\ampere},
\qquad
I_{\min}=\frac{30}{50+100}=\SI{0.2}{\ampere}.
\]
A faixa é $3{:}1$ --- \textbf{não} $\infty{:}1$, porque a carga fixa de
$\SI{50}{\ohm}$ impõe um teto à corrente por mais que o reostato seja reduzido.
Para ampliar a faixa seria preciso um reostato de valor maior em relação à
carga. Note também que a relação $I(R)$ é hiperbólica, não linear: o controle é
muito mais "sensível" perto de $R=0$.}
\stepcounter{questao}
\resposta{(b)}
\resolucao{A queda é $U_k=V\dfrac{R_k}{R_{eq}}$. Dobrar $R_k$ dobra o
numerador, mas \emph{também} aumenta o denominador --- o crescimento é menos
que proporcional.\\
\textbf{Exemplo:} $R_1=R_2=\SI{10}{\ohm}$ sob $\SI{12}{\volt}$ dá
$U_1=\SI{6}{\volt}$. Com $R_1=\SI{20}{\ohm}$: $U_1=12\cdot20/30=\SI{8}{\volt}$
--- aumento de $33\%$, não de $100\%$. No limite $R_1\to\infty$, $U_1\to V$: a
queda satura no valor da fonte, nunca o ultrapassa.}
\stepcounter{questao}
\resposta{$E=\SI{54}{\kilo\watt\hour}$; custo $=\mathrm{R\$}\,43{,}20$}
\resolucao{
\[
E=P\,t=0{,}300\,\si{\kilo\watt}\times6\times30
=\SI{54}{\kilo\watt\hour}.
\]
\[
\text{Custo}=54\times0{,}80=\mathrm{R\$}\,43{,}20.
\]
A conta não depende de \emph{como} a resistência está distribuída na cadeia ---
só da potência total. A distribuição importa para o dimensionamento térmico de
cada componente, não para a fatura.}
\stepcounter{questao}
\resposta{Sim ($22+35+41=98\approx100$, dentro de $2\%$);
$\SI{440}{\ohm}$, $\SI{700}{\ohm}$ e $\SI{820}{\ohm}$}
\resolucao{\textbf{Consistência (LKT):} a soma das quedas deve igualar a tensão
da fonte. Aqui $22+35+41=\SI{98}{\volt}$, desvio de $\SI{2}{\volt}$ ($2\%$) ---
compatível com a exatidão típica de um multímetro de bancada. Um desvio de,
digamos, $\SI{20}{\volt}$ indicaria erro de medida ou um quarto elemento não
contabilizado (resistência de contato, por exemplo).\\
\textbf{Resistores:} com $I=\SI{0.050}{\ampere}$ comum a todos,
\[
R_1=\frac{22}{0{,}05}=440,\quad
R_2=\frac{35}{0{,}05}=700,\quad
R_3=\frac{41}{0{,}05}=\SI{820}{\ohm}.
\]}
\stepcounter{questao}
\resposta{(a) $\SI{10}{\watt}$ e $\SI{20}{\watt}$; (b) $\SI{25}{\watt}$ e
$\SI{50}{\watt}$}
\resolucao{(a) $I=30/30=\SI{1}{\ampere}$; $P_1=10\cdot1=\SI{10}{\watt}$,
$P_2=20\cdot1=\SI{20}{\watt}$.\\
(b) Com fator $2$: $\SI{20}{\watt}$ e $\SI{40}{\watt}$; os valores comerciais
imediatamente superiores são $\SI{25}{\watt}$ e $\SI{50}{\watt}$.\\
(c) Especificar \emph{ambos} como $\SI{25}{\watt}$ deixaria $R_2$ operando a
$80\%$ do nominal --- sem margem para elevação de temperatura ambiente ou
ventilação deficiente. Especificar \emph{ambos} como $\SI{50}{\watt}$ dobraria
volume e custo de $R_1$ sem ganho. \textbf{Regra:} em série, $P\propto R$;
dimensione resistor a resistor.}
\stepcounter{questao}
\resposta{(a) $\SI{62}{\ohm}$; (b) $\SI{1.24}{\ampere}$ e
$\SI{1.107}{\ampere}$; (c) queda de $10{,}7\%$}
\resolucao{(a)
$R=R_0[1+\alpha(T-T_0)]=50[1+0{,}004(60)]=50\cdot1{,}24=\SI{62}{\ohm}$.\\
(b) A $\SI{20}{\celsius}$: $R_{eq}=\SI{100}{\ohm}$ e
$I=124/100=\SI{1.24}{\ampere}$.\\
A $\SI{80}{\celsius}$: $R_{eq}=62+50=\SI{112}{\ohm}$ e
$I=124/112=\SI{1.107}{\ampere}$.\\
(c) Redução de $10{,}7\%$. Note que o resistor de fio variou $24\%$, mas a
corrente variou apenas $10{,}7\%$ --- o resistor de filme, estável, "dilui" o
efeito. Essa é a base do item de projeto explorado adiante: combinar
coeficientes para estabilizar a cadeia inteira.}
\stepcounter{questao}
\resposta{(a) $\SI{300}{\ohm}$, entre $\SI{285}{\ohm}$ e $\SI{315}{\ohm}$;
(b) $\pm5\%$; (c) $\pm2{,}33\%$}
\resolucao{(a) Em série as incertezas \emph{absolutas} se somam:
$\Delta R=5+10=\SI{15}{\ohm}$, logo $300\pm15$, isto é, de
$\SI{285}{\ohm}$ a $\SI{315}{\ohm}$.\\
(b) $15/300=5\%$ --- \emph{idêntica} às parcelas.\\
(c) Agora $\Delta R=5+2=\SI{7}{\ohm}$, e $7/300=2{,}33\%$.\\
\textbf{Padrão geral:} a tolerância relativa da série é a \emph{média
ponderada} das tolerâncias individuais, com pesos $R_k/R_{eq}$:
\[
t_{eq}=\frac{\sum R_k t_k}{\sum R_k}.
\]
Em (b), $t_1=t_2$ e a média devolve o mesmo valor. Em (c), o resistor
\emph{maior} domina o resultado --- por isso, ao apertar tolerâncias, comece
sempre pela maior parcela da cadeia.}
\stepcounter{questao}
\resposta{$I$: $\SI{1}{\ampere}\to\SI{1.5}{\ampere}$;
$P_1$: $10\to\SI{22.5}{\watt}$; $P_3$: $30\to\SI{67.5}{\watt}$}
\resolucao{\textbf{Antes:} $R_{eq}=\SI{60}{\ohm}$, $I=\SI{1}{\ampere}$,
$P=(10;20;30)\,\si{\watt}$, total $\SI{60}{\watt}$.\\
\textbf{Depois:} $R_{eq}=10+30=\SI{40}{\ohm}$, $I=60/40=\SI{1.5}{\ampere}$.
\[
P_1=10(1{,}5)^{2}=\SI{22.5}{\watt},\qquad P_3=30(1{,}5)^{2}=\SI{67.5}{\watt},
\]
total $\SI{90}{\watt}$ --- $50\%$ a mais consumido da fonte.\\
\textbf{Ponto crítico:} $P_1$ e $P_3$ subiram $125\%$ (fator $1{,}5^{2}=2{,}25$).
Um curto em \emph{um} componente pode destruir os \emph{outros}, que passam a
operar muito acima do nominal. Essa é a razão de as falhas em cadeia serem
frequentemente sequenciais, e não simultâneas.}
\stepcounter{questao}
\resposta{$I=0$; $\SI{0}{\volt}$ em $R_1$ e $R_3$; $\SI{60}{\volt}$ em $R_2$}
\resolucao{Com o circuito interrompido, $I=0$ em toda a série. Pela lei de Ohm,
$U=RI=0$ nos resistores íntegros --- mesmo estando eles perfeitamente
conectados à fonte.\\
Pela LKT, a soma das quedas deve continuar valendo $\SI{60}{\volt}$; como as
outras são nulas, \textbf{toda} a tensão da fonte aparece sobre a
descontinuidade: o voltímetro sobre $R_2$ lê $\SI{60}{\volt}$.\\
\textbf{Uso prático:} é assim que se localiza o componente aberto em uma cadeia
--- percorre-se a série com o voltímetro e o defeito é o único ponto onde
aparece tensão. Note que o voltímetro precisa ser de alta impedância; um
instrumento de baixa impedância fecharia o circuito e falsearia a leitura.}
\stepcounter{questao}
\resposta{(a) $\SI{200}{\ohm}$; (b) $\SI{220}{\ohm}\Rightarrow
\SI{13.6}{\milli\ampere}$; (c) $\SI{41}{\milli\watt}$}
\resolucao{(a) O resistor absorve a diferença entre a fonte e a queda do LED:
\[
R=\frac{V_{fonte}-V_{LED}}{I}=\frac{5{,}0-2{,}0}{0{,}015}=\SI{200}{\ohm}.
\]
(b) A E12 oferece $\SI{180}{\ohm}$ e $\SI{220}{\ohm}$. Escolha o
\textbf{maior} --- errar para menos aumenta a corrente e reduz a vida do LED:
\[
I=\frac{3{,}0}{220}=\SI{13.6}{\milli\ampere}.
\]
(c) $P=R I^{2}=220(0{,}0136)^{2}=\SI{41}{\milli\watt}$; um resistor de
$\frac{1}{8}\,\si{\watt}$ já basta com folga generosa.\\
\textbf{Observação:} o LED \emph{não} obedece à lei de Ohm --- sua queda é
aproximadamente constante. Por isso ele entra na conta como uma "fonte de
tensão" a ser subtraída, e não como uma resistência a ser somada.}
\stepcounter{questao}
\resposta{(a) $\SI{1.2}{\ampere}$, $\SI{11.4}{\volt}$; (b) $5{,}26\%$;
(c) $5\%$}
\resolucao{(a) $I=\dfrac{12}{9{,}5+0{,}5}=\SI{1.2}{\ampere}$;
$V_{term}=12-0{,}5(1{,}2)=\SI{11.4}{\volt}$.\\
(b) Tomando a tensão em vazio como referência de partida:
\[
\text{Reg}=\frac{V_{vazio}-V_{carga}}{V_{carga}}
=\frac{12-11{,}4}{11{,}4}=5{,}26\%.
\]
(c) $P_r=rI^{2}=0{,}5(1{,}44)=\SI{0.72}{\watt}$ contra
$P_{tot}=12(1{,}2)=\SI{14.4}{\watt}$: exatamente $5\%$, que é a razão
$r/(r+R_L)$.\\
\textbf{Compromisso:} reduzir $r$ melhora regulação \emph{e} rendimento
simultaneamente --- por isso $r$ baixo é sempre desejável em fontes. É diferente
da máxima transferência de potência, que exigiria $R_L=r$ e um rendimento de
apenas $50\%$.}
\stepcounter{questao}
\resposta{Total $\SI{120}{\watt}$ em ambas; pior componente:
$\SI{30}{\watt}$ em (A) e $\SI{100}{\watt}$ em (B) --- (A) é preferível}
\resolucao{$I=120/120=\SI{1}{\ampere}$ nas duas.\\
(A) $P=30\cdot1=\SI{30}{\watt}$ em cada um dos quatro.\\
(B) $P_{100}=\SI{100}{\watt}$; $P_{10}=\SI{10}{\watt}$ (dois).\\
Total: $\SI{120}{\watt}$ nos dois casos --- a fonte não distingue as montagens.\\
\textbf{Decisão:} (A) exige quatro resistores de $\SI{50}{\watt}$ nominais;
(B) exige um de $\SI{200}{\watt}$, componente caro, volumoso e com ponto quente
concentrado. Sempre que possível, \textbf{divida a dissipação} em parcelas
iguais: a área de troca térmica cresce e a temperatura máxima cai.}
\stepcounter{questao}
\resposta{$\SI{0.71}{\ampere}$ na partida e $\SI{4.0}{\ampere}$ em regime}
\resolucao{\textbf{Partida (frio):}
$I=\dfrac{100}{120+20}=\SI{0.71}{\ampere}$.\\
\textbf{Regime (aquecido):}
$I=\dfrac{100}{5+20}=\SI{4.0}{\ampere}$.\\
\textbf{Função:} sem o NTC, a corrente de energização seria imediatamente
$100/20=\SI{5}{\ampere}$. O termistor a limita a $\SI{0.71}{\ampere}$ e vai
"saindo do caminho" à medida que aquece --- um \textbf{limitador de corrente de
partida}. É o que protege capacitores de entrada e contatos de relé em fontes
chaveadas.\\
\textbf{Limitação:} em regime o NTC ainda dissipa $5(4)^{2}=\SI{80}{\watt}$, e
precisa de tempo para esfriar antes de um novo religamento --- por isso ele não
protege contra reenergizações rápidas e sucessivas.}
\stepcounter{questao}
\resposta{$R_1=\SI{30}{\ohm}$, $R_2=\SI{60}{\ohm}$, $R_3=\SI{90}{\ohm}$;
quedas $\SI{15}{\volt}$, $\SI{30}{\volt}$, $\SI{45}{\volt}$}
\resolucao{
\[
R_{eq}=\frac{V}{I}=\frac{90}{0{,}5}=\SI{180}{\ohm}.
\]
Como $R_{eq}=R_1+2R_1+3R_1=6R_1$:
\[
R_1=\frac{180}{6}=\SI{30}{\ohm},\quad R_2=\SI{60}{\ohm},\quad R_3=\SI{90}{\ohm}.
\]
Quedas: $U_k=R_kI$, ou seja $15$, $30$ e $\SI{45}{\volt}$; soma
$=\SI{90}{\volt}$ \checkmark.\\
Como as quedas guardam a mesma proporção $1{:}2{:}3$ dos resistores, poderiam
ter sido obtidas direto: $90$ dividido em $6$ partes de $\SI{15}{\volt}$.}
\stepcounter{questao}
\resposta{$n=13$; $I=\SI{39.3}{\milli\ampere}$}
\resolucao{Exige-se $R_{eq}\ge V/I_{\max}=24/0{,}040=\SI{600}{\ohm}$:
\[
47n\ge600 \;\Longrightarrow\; n\ge12{,}77 \;\Longrightarrow\; n=13.
\]
\[
R_{eq}=13\cdot47=\SI{611}{\ohm}
\;\Longrightarrow\;
I=\frac{24}{611}=\SI{39.3}{\milli\ampere}\ \le\ \SI{40}{\milli\ampere}\ \checkmark
\]
\textbf{Cuidado com o arredondamento:} $n=12$ daria
$R_{eq}=\SI{564}{\ohm}$ e $I=\SI{42.6}{\milli\ampere}$ --- \emph{acima} do
limite. Em restrições do tipo "no máximo", arredonde sempre \textbf{para cima},
mesmo que a fração seja pequena.}
\stepcounter{questao}
\resposta{Aumentar $R_k$ reduz $P_{tot}$; a potência do outro sempre diminui
--- é a do próprio $R_k$ que pode subir}
\resolucao{Com $I=V/(R_1+R_2)$:
\[
P=VI=\frac{V^{2}}{R_1+R_2}.
\]
Como $R_1+R_2$ cresce, $P$ diminui --- \emph{sempre}.\\
Para o outro resistor, $P_2=R_2I^{2}=\dfrac{V^{2}R_2}{(R_1+R_2)^{2}}$: com
$R_2$ fixo e $R_1$ crescendo, $P_2$ só pode \textbf{diminuir}.\\
Já a potência do \emph{próprio} $R_1$,
$P_1=\dfrac{V^{2}R_1}{(R_1+R_2)^{2}}$, cresce enquanto $R_1<R_2$, atinge um
máximo em $R_1=R_2$ e depois decresce. \textbf{Conclusão:} o enunciado
intuitivo "mais resistência, mais aquecimento" é falso em série --- o
aquecimento de um resistor é máximo quando ele iguala a soma dos demais, e cai
para os dois lados. Este resultado é retomado quantitativamente no bloco de
desafio.}
\stepcounter{questao}
\resposta{(a) $\SI{300}{\ohm}$, $\SI{400}{\ohm}$, $\SI{500}{\ohm}$;
$\SI{0.12}{\watt}$, $\SI{0.16}{\watt}$, $\SI{0.20}{\watt}$;
(b) $300/390/510\ \si{\ohm}$ $\Rightarrow$ $5{,}98/7{,}77/10{,}16\ \si{\volt}$;
(c) sim, erro máximo de $2{,}9\%$}
\resolucao{(a) Como a corrente é comum, $R_k=U_k/I$:
\[
R_1=\frac{6}{0{,}02}=300,\quad R_2=\frac{8}{0{,}02}=400,\quad
R_3=\frac{10}{0{,}02}=\SI{500}{\ohm}.
\]
Potências: $P_k=U_kI$ $\Rightarrow$ $0{,}12$; $0{,}16$;
$\SI{0.20}{\watt}$. Um resistor de $\frac{1}{2}\,\si{\watt}$ em cada posição
oferece folga confortável.\\
(b) A E24 tem $300$ e $510$, mas não $400$ --- o vizinho é $\SI{390}{\ohm}$.
Com $R_{eq}=300+390+510=\SI{1200}{\ohm}$, a corrente real é
$24/1200=\SI{20.0}{\milli\ampere}$ (coincidência favorável), e as quedas ficam
\[
U_1=\SI{6.00}{\volt},\quad U_2=\SI{7.80}{\volt},\quad U_3=\SI{10.20}{\volt}.
\]
(c) Desvios: $0\%$, $-2{,}5\%$ e $+2{,}0\%$ --- todos dentro de $\pm5\%$.
\textbf{Ressalva essencial:} este projeto só vale \emph{em vazio}. Qualquer
carga real conectada em paralelo com uma das quedas altera a corrente da cadeia
e derruba todas as tensões. Um divisor resistivo puro serve para referências de
altíssima impedância, nunca para alimentar cargas.}
\stepcounter{questao}
\resposta{(a) $t$; (b) $t/\sqrt{N}$; (c) $5\%$ contra $2{,}24\%$}
\resolucao{(a) Pior caso: todos desviam para o mesmo lado.
$\Delta R_{eq}=N(tR)$ e $R_{eq}=NR$, logo
\[
\frac{\Delta R_{eq}}{R_{eq}}=\frac{NtR}{NR}=t.
\]
A tolerância relativa \textbf{não melhora nem piora} --- resultado que
surpreende, mas decorre de a série somar tanto os valores quanto os desvios.\\
(b) Desvios independentes combinam-se em quadratura:
\[
u(R_{eq})=\sqrt{N}\,(tR)
\;\Longrightarrow\;
\frac{u(R_{eq})}{R_{eq}}=\frac{\sqrt{N}\,tR}{NR}=\frac{t}{\sqrt{N}}.
\]
(c) Para $N=5$, $t=5\%$: pior caso $5\%$; estatístico
$5/\sqrt5=2{,}24\%$.\\
\textbf{Qual usar:} o \emph{pior caso} para projetos de segurança e para lotes
pequenos, em que os desvios podem ser correlacionados (resistores do mesmo
rolo, mesma batelada, mesma temperatura). O \emph{estatístico} para produção em
série com componentes de origens independentes, onde o pior caso seria
excessivamente conservador. Note que a hipótese de independência é quase sempre
otimista na prática.}
\stepcounter{questao}
\resposta{(a) $R_1\alpha_1+R_2\alpha_2=0$; (b) $R_2=\SI{400}{\ohm}$,
$R_{eq}=\SI{500}{\ohm}$; (c) só é exata em primeira ordem e numa temperatura}
\resolucao{(a) Para pequenas variações,
$R_k(T)=R_k[1+\alpha_k\Delta T]$, logo
\[
R_{eq}(T)=(R_1+R_2)+(R_1\alpha_1+R_2\alpha_2)\Delta T.
\]
Anular o termo em $\Delta T$ exige
\[
R_1\alpha_1+R_2\alpha_2=0.
\]
(b) $R_2=-\dfrac{R_1\alpha_1}{\alpha_2}
=\dfrac{100\cdot4000}{1000}=\SI{400}{\ohm}$, e
$R_{eq}=100+400=\SI{500}{\ohm}$.\\
\textbf{Verificação a $\Delta T=\SI{50}{\celsius}$:}
$R_1\to100(1{,}20)=\SI{120}{\ohm}$ e $R_2\to400(0{,}95)=\SI{380}{\ohm}$; soma
$=\SI{500}{\ohm}$ \checkmark.\\
(c) Três limitações. \textbf{(i)} A compensação é de \emph{primeira ordem};
$\alpha$ real varia com $T$, e o cancelamento se degrada longe da temperatura
de projeto. \textbf{(ii)} Os dois resistores precisam estar \emph{à mesma
temperatura} --- exige acoplamento térmico, o que é difícil quando as
dissipações diferem (aqui, $P_2=4P_1$). \textbf{(iii)} O valor de $R_2$ fica
\emph{amarrado} pela razão $\alpha_1/\alpha_2$: não se pode escolher
livremente $R_{eq}$ \emph{e} exigir compensação. Materiais dedicados
(manganina, constantan, com $\alpha$ naturalmente próximo de zero) resolvem o
problema de forma mais direta.}
\stepcounter{questao}
\resposta{(a) $\SI{12}{\volt}$ com $\SI{1.8}{\volt}$ $\Rightarrow$
$R\ge\SI{510}{\ohm}$; (b) $\SI{13.3}{\milli\ampere}$;
(c) $\SI{0.204}{\watt}\Rightarrow\frac{1}{2}\,\si{\watt}$}
\resolucao{(a) A corrente é máxima quando a fonte está no \textbf{máximo} e a
queda do LED no \textbf{mínimo} --- é preciso combinar os extremos que atuam no
mesmo sentido, e não os extremos "nominais":
\[
R\ge\frac{12-1{,}8}{0{,}020}=\SI{510}{\ohm}.
\]
O valor $\SI{510}{\ohm}$ existe na série E24.\\
(b) O outro extremo (fonte mínima, queda máxima):
\[
I_{\min}=\frac{9-2{,}2}{510}=\SI{13.3}{\milli\ampere}.
\]
A corrente varia $13{,}3$--$\SI{20}{\milli\ampere}$, razão de $1{,}5{:}1$. Como
o fluxo luminoso é aproximadamente proporcional à corrente, o brilho varia
cerca de $33\%$ --- perceptível, porém aceitável para sinalização. Para brilho
constante seria preciso uma fonte de \emph{corrente}, não um resistor.\\
(c) $P=RI_{\max}^{2}=510(0{,}020)^{2}=\SI{0.204}{\watt}$; especifique
$\frac{1}{2}\,\si{\watt}$.\\
\textbf{Método:} em projeto com faixas, nunca use valores nominais --- enumere
os extremos e identifique qual combinação é crítica para \emph{cada} requisito.}
\stepcounter{questao}
\resposta{(a) $-\pot{1.33}{-4}\,\si{\ampere\per\ohm}$;
(b) $S=-R_1/(R_1+R_2)=-1/3$; (c) $S\to0$ e $S\to-1$}
\resolucao{(a)
\[
\frac{\partial I}{\partial R_1}=-\frac{V}{(R_1+R_2)^{2}}
=-\frac{12}{300^{2}}=-\pot{1.33}{-4}\,\si{\ampere\per\ohm}.
\]
Um aumento de $\SI{1}{\ohm}$ em $R_1$ reduz a corrente em
$\SI{0.133}{\milli\ampere}$ (de $\SI{40}{\milli\ampere}$).\\
(b)
\[
S=\frac{\partial I}{\partial R_1}\cdot\frac{R_1}{I}
=-\frac{V}{(R_1+R_2)^{2}}\cdot\frac{R_1(R_1+R_2)}{V}
=-\frac{R_1}{R_1+R_2}=-\frac13.
\]
Ou seja: $1\%$ de erro em $R_1$ produz apenas $0{,}33\%$ de erro em $I$.\\
(c) Se $R_1\ll R_2$, $S\to0$: a corrente é governada por $R_2$ e $R_1$ torna-se
irrelevante --- é o caso do resistor \emph{shunt} de medição, deliberadamente
pequeno para não perturbar o circuito. Se $R_1\gg R_2$, $S\to-1$: $R_1$ passa a
determinar sozinho a corrente, e toda a sua imprecisão é transferida
integralmente. \textbf{Regra de projeto:} coloque a tolerância apertada no
elemento cuja sensibilidade relativa se aproxima de $1$; nos demais, use
componentes baratos.}
\stepcounter{questao}
\resposta{(a) $\SI{1}{\ampere}$, $\SI{10}{\volt}$, $\SI{10}{\watt}$;
(b) $\SI{1.11}{\ampere}$, $\SI{11.1}{\volt}$, $\SI{12.3}{\watt}$;
(c) o curto é mais perigoso}
\resolucao{(a) $R_{eq}=\SI{100}{\ohm}$, $I=\SI{1}{\ampere}$,
$U=\SI{10}{\volt}$, $P=\SI{10}{\watt}$ cada; total $\SI{100}{\watt}$.\\
(b) $R_{eq}=\SI{90}{\ohm}$, $I=100/90=\SI{1.111}{\ampere}$ ($+11\%$),
$U=\SI{11.1}{\volt}$, $P=10(1{,}111)^{2}=\SI{12.35}{\watt}$ ($+23{,}5\%$);
total $\SI{111}{\watt}$.\\
(c) O \textbf{aberto} zera a corrente: o circuito para de funcionar, mas nada
sofre --- é uma falha \emph{segura} (\emph{fail-safe}), fácil de localizar com
voltímetro.\\
O \textbf{curto} é insidioso: o circuito continua operando, aparentemente
normal, mas cada resistor remanescente passa a dissipar $23{,}5\%$ a mais. Se
já operavam próximos do nominal, entram em sobrecarga --- e cada novo curto
agrava o quadro de forma \emph{acelerada}. Com dois curtos,
$P=\SI{15.6}{\watt}$ ($+56\%$); com três, $\SI{20.4}{\watt}$ ($+104\%$).
\textbf{Conclusão:} falhas que \emph{reduzem} a resistência de um circuito em
série tendem a ser progressivas e devem ser detectadas por monitoração de
corrente, não por observação funcional.}
\stepcounter{questao}
\resposta{(a) $r=\SI{0.5}{\ohm}$, $\EMF=\SI{12}{\volt}$;
(c) $I_{cc}=\SI{24}{\ampere}$, valor apenas indicativo}
\resolucao{(a) O modelo é $V=\EMF-rI$ --- uma reta. Dois pontos a determinam:
\[
r=-\frac{\Delta V}{\Delta I}=-\frac{10{,}5-11{,}5}{3-1}=\SI{0.5}{\ohm},
\]
\[
\EMF=V+rI=11{,}5+0{,}5(1)=\SI{12}{\volt}.
\]
(b) A medida em vazio exigiria um voltímetro de impedância \emph{infinita}; com
impedância finita, circula uma pequena corrente e a leitura já não é $\EMF$.
Além disso, baterias reais apresentam recuperação de tensão em repouso, o que
falseia a leitura em vazio. O ajuste por dois (ou mais) pontos de carga
\textbf{extrapola} para $I=0$ sem precisar chegar lá, e ainda entrega $r$ de
brinde. Com três ou mais pontos, uma regressão linear também revela se o modelo
é adequado.\\
(c) $I_{cc}=\EMF/r=12/0{,}5=\SI{24}{\ampere}$. \textbf{Ressalva:} o valor
pressupõe $r$ constante, hipótese que falha em correntes altas --- $r$ cresce
com a corrente e com a descarga, e há aquecimento. O número serve como
\emph{ordem de grandeza} para dimensionar proteção, jamais como previsão
precisa.}
\stepcounter{questao}
\resposta{(a) cinco em série; (b) $\SI{0.4}{\watt}$ cada, margem de $25\%$;
(c) $\SI{44.7}{\volt}$ e $\SI{44.7}{\milli\ampere}$}
\resolucao{(a) $5\times\SI{200}{\ohm}=\SI{1000}{\ohm}$ \checkmark.\\
(b) Em série a corrente é comum e os resistores são iguais, logo a potência se
reparte igualmente:
\[
P_k=\frac{2}{5}=\SI{0.4}{\watt}\ <\ \SI{0.5}{\watt}\ \checkmark
\]
Margem: $(0{,}5-0{,}4)/0{,}4=25\%$. É uma folga modesta --- em ambiente quente
ou sem ventilação, considere sete unidades de $\SI{150}{\ohm}$
($\SI{1050}{\ohm}$) ou resistores de $\SI{1}{\watt}$.\\
(c)
\[
V=\sqrt{PR}=\sqrt{2\cdot1000}=\SI{44.7}{\volt},
\qquad
I=\sqrt{P/R}=\sqrt{0{,}002}=\SI{44.7}{\milli\ampere}.
\]
\textbf{Por que a série resolve:} ela multiplica \emph{simultaneamente} o valor
e a capacidade de dissipação por $N$. Uma associação em paralelo de cinco
unidades também daria $\SI{2.5}{\watt}$, mas $\SI{40}{\ohm}$ --- resolveria a
potência e destruiria o valor.}
\stepcounter{questao}
\resposta{(a) $1$, $2$, $3$, $4$ e $\SI{5}{\volt}$; (b) pela quebra do degrau
uniforme de $\SI{1}{\volt}$; (c) desacopla os sensores entre si}
\resolucao{(a) Com $I=\SI{1}{\milli\ampere}$, cada sensor cai
$\SI{1}{\volt}$. Acumulando a partir do terra: $1$, $2$, $3$, $4$,
$\SI{5}{\volt}$.\\
(b) Se o \emph{terceiro} sensor for a $\SI{1.2}{\kilo\ohm}$, as tensões passam
a $1$, $2$, $3{,}2$, $4{,}2$, $\SI{5.2}{\volt}$. A \textbf{diferença entre nós
consecutivos} vale $1;\,1;\,1{,}2;\,1;\,1$ --- o degrau anômalo aponta
diretamente o sensor defeituoso, e sua magnitude dá o novo valor:
$\Delta R=\Delta U/I=0{,}2/0{,}001=\SI{200}{\ohm}$.\\
(c) Com \textbf{corrente constante}, a queda de cada sensor é $U_k=IR_k$ ---
depende \emph{só} do próprio sensor. A variação de um não afeta a leitura dos
demais, e a relação é rigorosamente linear.\\
Com fonte de \textbf{tensão}, a corrente seria $V/\sum R_k$: a mudança de um
sensor alteraria a corrente e, portanto, \emph{todas} as leituras
simultaneamente --- o sistema ficaria acoplado, com resposta não linear e sem
identificação direta do sensor. É exatamente por isso que sistemas de RTD e
extensômetros são excitados por corrente.}
\stepcounter{questao}
\resposta{(a) $R=R_f$; (b) $P_{\max}=V^{2}/(4R_f)$, rendimento $50\%$;
(c) a potência total é máxima em $R\to0$}
\resolucao{(a)
\[
P_R=RI^{2}=\frac{V^{2}R}{(R+R_f)^{2}}.
\]
\[
\frac{\mathrm{d}P_R}{\mathrm{d}R}
=V^{2}\frac{(R+R_f)^{2}-R\cdot2(R+R_f)}{(R+R_f)^{4}}
=V^{2}\frac{R_f-R}{(R+R_f)^{3}}=0
\;\Longrightarrow\;
R=R_f.
\]
A derivada é positiva para $R<R_f$ e negativa para $R>R_f$: máximo confirmado.\\
(b) $P_{\max}=\dfrac{V^{2}R_f}{(2R_f)^{2}}=\dfrac{V^{2}}{4R_f}$. A potência
total nesse ponto é $\dfrac{V^{2}}{2R_f}$, de modo que $R$ recebe exatamente
\textbf{metade} --- rendimento de $50\%$.\\
(c) A potência \emph{total} $P=\dfrac{V^{2}}{R+R_f}$ é monotonicamente
decrescente: ela é máxima quando $R\to0$, e aí $P\to V^{2}/R_f$ --- o dobro do
valor anterior, porém integralmente entregue ao bloco fixo.\\
\textbf{Distinção que importa:} "máxima potência \emph{no elemento}" e "máxima
potência \emph{do circuito}" são objetivos diferentes e incompatíveis. O
primeiro é o critério de casamento (relevante quando o bloco fixo é a
resistência interna de uma fonte); o segundo é o critério de eficiência. Este
mesmo antagonismo reaparece no teorema da máxima transferência de potência.}

\gabarito[3.1 Circuitos em série]

\subsection{Circuitos em paralelo}
\stepcounter{questao}
\resposta{(c)}
\resolucao{Todos os ramos estão ligados ao mesmo par de nós, logo compartilham a
mesma tensão. A corrente, essa sim, se divide --- e de forma \emph{inversamente}
proporcional às resistências. A resistência equivalente do paralelo é sempre
\textbf{menor} que a menor das resistências associadas, o que elimina (b) e (e).}
\stepcounter{questao}
\resposta{(a) \SI{1.2}{\ohm}, $I_t=\SI{25}{\ampere}$;
(b) \SI{4}{\ohm}, $I_t=\SI{6}{\ampere}$;
(c) \SI{2}{\ohm}, $I_t=\SI{7.5}{\ampere}$;
(d) \SI{20}{\ohm}, $I_t=\SI{0.45}{\ampere}$}
\resolucao{Em todos os casos a tensão da fonte aparece \emph{inteira} sobre cada
ramo, então cada corrente sai direto da lei de Ohm.\\
\textbf{(a)} $I_1=30/4=\SI{7.5}{\ampere}$, $I_2=30/3=\SI{10}{\ampere}$,
$I_3=30/6=\SI{5}{\ampere}$, $I_4=30/12=\SI{2.5}{\ampere}$;
$I_t=\SI{25}{\ampere}$ e $\Req=30/25=\SI{1.2}{\ohm}$.\\
\textbf{(b)} Três resistores iguais: $\Req = 12/3 = \SI{4}{\ohm}$;
cada ramo com $24/12=\SI{2}{\ampere}$ e $I_t = \SI{6}{\ampere}$.\\
\textbf{(c)} $\dfrac{1}{\Req}=\dfrac{1}{12}+\dfrac{1}{12}+\dfrac{1}{3}
=\dfrac{1+1+4}{12}=\dfrac{6}{12}$, logo $\Req=\SI{2}{\ohm}$.
$I_1=I_2=\SI{1.25}{\ampere}$, $I_3=\SI{5}{\ampere}$, $I_t=\SI{7.5}{\ampere}$.\\
\textbf{(d)} $\Req=\dfrac{30\cdot 60}{30+60}=\SI{20}{\ohm}$;
$I_1=\SI{0.3}{\ampere}$, $I_2=\SI{0.15}{\ampere}$, $I_t=\SI{0.45}{\ampere}$.\\
Note em (a) que $\Req$ ficou menor que o menor resistor (\SI{3}{\ohm}) --- sempre
é assim no paralelo.}
\stepcounter{questao}
\resposta{(d)}
\resolucao{Para $n$ resistores iguais em paralelo, $\Req = \dfrac{R}{n}$, logo
$n = \dfrac{60}{10} = 6$.}
\stepcounter{questao}
\resposta{$R = \SI{30}{\ohm}$}
\resolucao{Da definição de paralelo,
\[
\frac{1}{12}=\frac{1}{20}+\frac{1}{R}
\quad\Rightarrow\quad
\frac{1}{R}=\frac{1}{12}-\frac{1}{20}=\frac{5-3}{60}=\frac{2}{60}
\quad\Rightarrow\quad R=\SI{30}{\ohm}.
\]
\textbf{Atalho útil:} com dois resistores,
$R = \dfrac{R_1\,\Req}{R_1-\Req} = \dfrac{20\cdot 12}{20-12} = \SI{30}{\ohm}$.
Note que $\Req$ precisa necessariamente ser menor que \SI{20}{\ohm}; se o problema
pedisse $\Req = \SI{25}{\ohm}$, não haveria solução com resistor positivo.}
\stepcounter{questao}
\resposta{(b)}
\resolucao{A corrente se divide de forma \emph{inversamente} proporcional às
resistências: o menor resistor recebe a maior corrente. Pelo divisor de corrente,
\[
I_1 = I_t\,\frac{R_2}{R_1+R_2}=12\cdot\frac{12}{16}=\SI{9}{\ampere},
\qquad
I_2 = 12-9 = \SI{3}{\ampere}.
\]
\textbf{Verificação alternativa:} $\Req = \dfrac{4\cdot12}{16}=\SI{3}{\ohm}$ e
$U = 3\cdot 12 = \SI{36}{\volt}$; daí $I_1 = 36/4 = \SI{9}{\ampere}$ e
$I_2 = 36/12 = \SI{3}{\ampere}$. A armadilha da alternativa (a) é aplicar a
proporção direta, como se fosse divisor de tensão.}
\stepcounter{questao}
\resposta{(a) \SI{240}{\ohm} cada, $\Req=\SI{80}{\ohm}$;
(b) \SI{1.5}{\ampere}; (c) nada muda}
\resolucao{(a) De $P = \dfrac{U^2}{R}$:
$R = \dfrac{120^2}{60} = \SI{240}{\ohm}$ por lâmpada.
Três iguais em paralelo: $\Req = \dfrac{240}{3} = \SI{80}{\ohm}$.\\
(b) $I_t = \dfrac{120}{80} = \SI{1.5}{\ampere}$
(ou $3\times\SI{0.5}{\ampere}$ por lâmpada).\\
(c) As demais continuam com os mesmos \SI{120}{\volt} nos seus terminais, logo o
brilho \emph{não muda}: cada ramo é independente. É por isso que as instalações
residenciais são feitas em paralelo, e não em série.}
\stepcounter{questao}
\resposta{(d)}
\resolucao{O fio ideal curto-circuita $R_2$: os dois terminais de $R_2$ ficam no
\emph{mesmo potencial}, logo $U_{R_2} = 0$ e nenhuma corrente passa por ele.
Formalmente, $R_2 \parallel 0 = 0$. Resta apenas $R_1$:
\[
I=\frac{20}{10}=\SI{2}{\ampere}.
\]
\textbf{Erro comum:} calcular $10 + 20 = \SI{30}{\ohm}$ e responder
\SI{0.67}{\ampere}, ignorando o curto.}
\stepcounter{questao}
\resposta{(a) 4 resistores; (b) $\approx\SI{14.1}{\volt}$;
(c) \SI{8}{\watt}}
\resolucao{(a) $\Req = R/n \Rightarrow n = 100/25 = 4$ resistores.\\
(b) Todos ficam sob a \emph{mesma} tensão, então basta respeitar o limite de um
deles:
\[
P_{\max}=\frac{U^2}{R}\ \Rightarrow\ U_{\max}=\sqrt{P_{\max}R}=\sqrt{2\cdot100}
=\sqrt{200}\approx\SI{14.1}{\volt}.
\]
(c) Como os 4 resistores trabalham no limite simultaneamente,
$P_{\text{total}} = 4\cdot 2 = \SI{8}{\watt}$ --- que também pode ser conferido
por $P = U^2/\Req = 200/25 = \SI{8}{\watt}$.\\
\textbf{Ideia geral:} associar resistores iguais em paralelo (ou em série)
multiplica por $n$ a capacidade de dissipação do conjunto.}
\stepcounter{questao}
\resposta{$R_{AB} = \SI{45}{\ohm}$}
\resolucao{O fio ideal liga os dois terminais do resistor de \SI{30}{\ohm}
central, curto-circuitando-o. Ele deixa de conduzir e pode ser desconsiderado no circuito equivalente. Sobram
apenas os resistores de \SI{15}{\ohm} e \SI{30}{\ohm} em série:
\[
R_{AB}=15+30=\SI{45}{\ohm}.
\]
\textbf{Método seguro:} antes de calcular, redesenhe o circuito nomeando os nós.
Dois pontos ligados por fio ideal são \emph{o mesmo nó} --- qualquer componente
entre eles está curto-circuitado.}
\stepcounter{questao}
\resposta{(a) e (b) demonstrações; (c) $\SI{10}{\ohm}$ e $\SI{15}{\ohm}$}
\resolucao{\textbf{(a)} Seja $R_{\min}$ a menor das resistências. Da definição,
\[
\frac{1}{\Req}=\frac{1}{R_1}+\dots+\frac{1}{R_n}
> \frac{1}{R_{\min}},
\]
pois a soma inclui o termo $\dfrac{1}{R_{\min}}$ \emph{mais} outras parcelas
positivas. Invertendo a desigualdade (ambos os lados positivos):
$\Req < R_{\min}$. \checkmark\\
Interpretação: acrescentar um caminho \emph{sempre} facilita a passagem de
corrente.

\textbf{(b)} Se $R_k \to 0$, então $\dfrac{1}{R_k}\to\infty$, e a soma das
condutâncias diverge, levando $\Req \to 0$. Fisicamente, esse ramo é um
\textbf{curto-circuito}: ele desvia toda a corrente e anula a tensão sobre o
conjunto, tornando irrelevantes os demais resistores.

\textbf{(c)} Temos o sistema
\[
\begin{cases}
R_1+R_2 = 25\\[2pt]
\dfrac{R_1R_2}{R_1+R_2}=6 \;\Rightarrow\; R_1R_2 = 6\cdot25 = 150.
\end{cases}
\]
Soma e produto conhecidos: $R_1$ e $R_2$ são as raízes de
$x^2-25x+150=0$, isto é,
\[
x=\frac{25\pm\sqrt{625-600}}{2}=\frac{25\pm5}{2}
\;\Longrightarrow\;
R_1=\SI{15}{\ohm},\quad R_2=\SI{10}{\ohm}.
\]
\textbf{Verificação:} $15+10=25$ \checkmark\ e
$\dfrac{15\cdot10}{25}=\SI{6}{\ohm}$ \checkmark\\
\textbf{Observação:} o problema só tem solução real se
$(R_s)^2 \ge 4R_sR_p$, ou seja, $R_s \ge 4R_p$. Aqui, $25 \ge 24$ --- a condição é satisfeita por pequena margem. Se o enunciado pedisse $\Req=\SI{7}{\ohm}$ em paralelo com a mesma soma
$\SI{25}{\ohm}$, não existiriam resistores reais que satisfizessem.}
\stepcounter{questao}
\resposta{$R_{eq}=\SI{20}{\ohm}$}
\resolucao{Para dois resistores, a regra do produto pela soma é a mais rápida:
\[
R_{eq}=\frac{R_1R_2}{R_1+R_2}=\frac{30\cdot60}{90}=\SI{20}{\ohm}.
\]
\textbf{Conferência instantânea:} o resultado ($\SI{20}{\ohm}$) é menor que o
menor dos dois ($\SI{30}{\ohm}$) --- se não for, há erro. Cuidado: a regra
produto/soma vale \emph{apenas} para dois resistores.}
\stepcounter{questao}
\resposta{$G=\SI{0.5}{\ensuremath{\mathrm{S}}}$; $R_{eq}=\SI{2}{\ohm}$}
\resolucao{
\[
G=\frac14+\frac18+\frac18=0{,}25+0{,}125+0{,}125=\SI{0.5}{\ensuremath{\mathrm{S}}}
\;\Longrightarrow\;
R_{eq}=\frac{1}{0{,}5}=\SI{2}{\ohm}.
\]
Com três ou mais ramos, some as condutâncias e inverta \emph{uma única vez} no
final. Inverter parcialmente a cada passo é a origem mais comum de erro.}
\stepcounter{questao}
\resposta{$R_{eq}=\SI{30}{\ohm}$}
\resolucao{Para $n$ resistores \emph{iguais},
\[
R_{eq}=\frac{R}{n}=\frac{90}{3}=\SI{30}{\ohm}.
\]
O atalho vale só quando todos são iguais; com valores diferentes é obrigatório
somar condutâncias.}
\stepcounter{questao}
\resposta{(a)}
\resolucao{Cada novo ramo acrescenta uma condutância positiva:
$G_{eq}=\sum G_k$ só pode crescer, logo $R_{eq}$ só pode diminuir. Com $V$
fixo, $I=V/R_{eq}$ aumenta.\\
Note o contraste com a série, em que acrescentar elemento \emph{reduz} a
corrente. E note também que os ramos já existentes continuam com \emph{a mesma}
corrente --- quem cresce é o total.}
\stepcounter{questao}
\resposta{$I=\SI{0.5}{\ampere}$}
\resolucao{
\[
I=\frac{V}{R}=\frac{24}{48}=\SI{0.5}{\ampere}.
\]
No circuito RLC em paralelo, a tensão é \emph{comum} a todos os ramos: a corrente de cada um
depende só do próprio resistor. Não é preciso conhecer os demais ramos nem
$R_{eq}$.}
\stepcounter{questao}
\resposta{$I_t=\SI{10}{\ampere}$; $R_{eq}=\SI{4}{\ohm}$}
\resolucao{Pela LKC, as correntes de ramo se somam:
$I_t=2+3+5=\SI{10}{\ampere}$.
\[
R_{eq}=\frac{V}{I_t}=\frac{40}{10}=\SI{4}{\ohm}.
\]
Observe que não foi preciso conhecer os resistores individualmente --- embora
eles saiam de imediato: $20$, $13{,}3$ e $\SI{8}{\ohm}$.}
\stepcounter{questao}
\resposta{(b)}
\resolucao{Com $V$ comum, $I_k=V/R_k$: a corrente é \emph{inversamente}
proporcional à resistência. O ramo de menor $R$ é o caminho de menor oposição e
leva a maior corrente.\\
A alternativa (c) é sedutora, mas falsa: em circuitos de parâmetros
concentrados, a posição física dos ramos é irrelevante --- só importa a
topologia.}
\stepcounter{questao}
\resposta{$\SI{6}{\ampere}$, $\SI{3}{\ampere}$, $\SI{3}{\ampere}$;
$I_t=\SI{12}{\ampere}$; $R_{eq}=\SI{5}{\ohm}$}
\resolucao{Tensão comum, então cada ramo é independente:
\[
I_1=\frac{60}{10}=6,\quad I_2=I_3=\frac{60}{20}=\SI{3}{\ampere}.
\]
\[
I_t=6+3+3=\SI{12}{\ampere}
\;\Longrightarrow\;
R_{eq}=\frac{60}{12}=\SI{5}{\ohm}.
\]
\textbf{Conferência:} $G=0{,}1+0{,}05+0{,}05=\SI{0.2}{\ensuremath{\mathrm{S}}}$ e
$1/0{,}2=\SI{5}{\ohm}$ \checkmark.}
\stepcounter{questao}
\resposta{$\SI{360}{\watt}$, $\SI{180}{\watt}$, $\SI{180}{\watt}$; total
$\SI{720}{\watt}$}
\resolucao{Com $V$ comum, use $P=V^{2}/R$:
\[
P_1=\frac{3600}{10}=360,\qquad P_2=P_3=\frac{3600}{20}=\SI{180}{\watt}.
\]
Fonte: $P=VI_t=60\cdot12=\SI{720}{\watt}=360+180+180$ \checkmark.\\
\textbf{Regra do paralelo:} $P\propto1/R$ --- o \emph{menor} resistor dissipa
mais. Compare com a série, em que $P\propto R$. Trocar essas duas regras é o
erro conceitual mais frequente em associações.}
\stepcounter{questao}
\resposta{$I_1=\SI{9}{\ampere}$, $I_2=\SI{6}{\ampere}$}
\resolucao{Para \emph{dois} ramos, a corrente de um deles é proporcional à
resistência do \textbf{outro}:
\[
I_1=I_t\frac{R_2}{R_1+R_2}=15\cdot\frac{12}{20}=\SI{9}{\ampere},
\qquad
I_2=15\cdot\frac{8}{20}=\SI{6}{\ampere}.
\]
\textbf{Verificação:} $9+6=\SI{15}{\ampere}$ \checkmark, e a maior corrente ficou
no \emph{menor} resistor \checkmark. O "cruzamento" dos índices é o ponto que
mais confunde --- ele existe porque $I_k\propto G_k$, e escrever a fórmula em
condutâncias ($I_1=I_tG_1/(G_1+G_2)$) elimina a inversão.}
\stepcounter{questao}
\resposta{$R_3=\SI{50}{\ohm}$}
\resolucao{Trabalhe em condutâncias, onde a subtração é direta:
\[
G_{eq}=\frac{I_t}{V}=\frac{5}{100}=\SI{0.05}{\ensuremath{\mathrm{S}}},
\qquad
G_1+G_2=0{,}02+0{,}01=\SI{0.03}{\ensuremath{\mathrm{S}}}.
\]
\[
G_3=0{,}05-0{,}03=\SI{0.02}{\ensuremath{\mathrm{S}}}
\;\Longrightarrow\;
R_3=\SI{50}{\ohm}.
\]
Tentar subtrair \emph{resistências} aqui levaria a um resultado sem sentido:
no paralelo, o que é aditivo é $G$, não $R$.}
\stepcounter{questao}
\resposta{Em paralelo, a de $\SI{240}{\ohm}$ ($\SI{60}{\watt}$ contra
$\SI{15}{\watt}$) --- inverte-se em relação à série}
\resolucao{
\[
P_{240}=\frac{120^{2}}{240}=\SI{60}{\watt},
\qquad
P_{960}=\frac{120^{2}}{960}=\SI{15}{\watt}.
\]
Em paralelo, $V$ é comum e $P\propto1/R$: brilha mais a de \emph{menor}
resistência. Em série (mesma dupla), $I$ é comum e $P\propto R$: brilharia mais
a de $\SI{960}{\ohm}$.\\
\textbf{Conclusão que vale guardar:} não existe "a lâmpada mais potente" em
termos absolutos --- a potência depende da \emph{ligação}. É por isso que
lâmpadas residenciais, sempre em paralelo, obedecem à intuição usual
($\SI{100}{\watt}$ brilha mais que $\SI{25}{\watt}$), mas guirlandas em série
se comportam ao contrário.}
\stepcounter{questao}
\resposta{$R_{sh}=\SI{10}{\ohm}$}
\resolucao{O shunt fica em paralelo com o galvanômetro e desvia o excedente:
$I_{sh}=11-1=\SI{10}{\milli\ampere}$. Como a tensão é comum,
\[
R_{sh}I_{sh}=R_gI_g
\;\Longrightarrow\;
R_{sh}=\frac{R_gI_g}{I_{sh}}=\frac{100\cdot1}{10}=\SI{10}{\ohm}.
\]
A regra geral é $R_{sh}=\dfrac{R_g}{n-1}$, com $n=I_{total}/I_g$ o fator de
ampliação --- aqui $n=11$ e $R_{sh}=100/10$ \checkmark.}
\stepcounter{questao}
\resposta{$I_t=\SI{27.5}{\ampere}$; disjuntor de $\SI{32}{\ampere}$}
\resolucao{Cada carga drena $I=P/V$:
\[
I_1=\frac{4400}{220}=20,\quad
I_2=\frac{1100}{220}=5,\quad
I_3=\frac{550}{220}=\SI{2.5}{\ampere}.
\]
Total: $\SI{27.5}{\ampere}$ (as correntes se somam porque as cargas estão em
paralelo).\\
O disjuntor comercial imediatamente superior é o de $\SI{32}{\ampere}$ ---
$\SI{25}{\ampere}$ seria insuficiente. \textbf{Atenção:} escolhido o disjuntor,
a seção do condutor deve suportar \emph{a corrente do disjuntor}, não a da
carga; caso contrário, o cabo vira o elo fraco da proteção.}
\stepcounter{questao}
\resposta{$\SI{0.4}{\ohm}$ --- metade}
\resolucao{$R_{eq}=0{,}8/2=\SI{0.4}{\ohm}$.\\
Dobrar o número de condutores idênticos equivale a dobrar a seção: a
resistência cai à metade, e com ela a queda de tensão e as perdas
$RI^{2}$ na linha.\\
\textbf{Condição prática:} os condutores devem ter mesmo comprimento, seção e
material. Se um for mais curto, ele terá menor resistência e assumirá corrente
desproporcional --- podendo sobreaquecer enquanto o outro fica ocioso. Por isso
condutores em paralelo são sempre especificados com o mesmo lance.}
\stepcounter{questao}
\resposta{(a)}
\resolucao{$G=0{,}01+10^{-6}\approx\SI{0.010001}{\ensuremath{\mathrm{S}}}$, logo
$R_{eq}\approx\SI{99.99}{\ohm}$ --- diferença de $0{,}01\%$.\\
\textbf{Princípio geral:} no paralelo, quem manda é o \emph{menor} resistor;
ramos muito maiores são praticamente circuitos abertos. É exatamente por isso
que um voltímetro de alta impedância "não carrega" o circuito sob medição.}
\stepcounter{questao}
\resposta{$\SI{2200}{\watt}$ com uma; $\SI{4400}{\watt}$ com as duas}
\resolucao{Uma resistência:
$P=\dfrac{220^{2}}{22}=\SI{2200}{\watt}$.\\
Duas em paralelo: $R_{eq}=\SI{11}{\ohm}$ e
$P=\dfrac{220^{2}}{11}=\SI{4400}{\watt}$.\\
A potência \emph{dobra} porque a resistência caiu à metade sob a mesma tensão.
É assim que funcionam as chaves inverno/verão. \textbf{Consequência para a
instalação:} a corrente também dobra ($10\to\SI{20}{\ampere}$), e o
dimensionamento de disjuntor e condutor deve ser feito para a posição de maior
potência.}
\stepcounter{questao}
\resposta{Consistente; $\SI{5}{\ohm}$, $\SI{20}{\ohm}$, $\SI{20}{\ohm}$;
$P=\SI{120}{\watt}$}
\resolucao{\textbf{LKC:} $4+1+1=\SI{6}{\ampere}$ \checkmark.\\
\textbf{Resistências:} $R_k=V/I_k$ $\Rightarrow$ $20/4=5$ e
$20/1=\SI{20}{\ohm}$ (duas vezes).\\
\textbf{Potência:} $P=VI_t=20\cdot6=\SI{120}{\watt}$; por ramo,
$80+20+20=\SI{120}{\watt}$ \checkmark.\\
Dois testes independentes --- LKC e balanço de potência --- fecham. Se apenas o
primeiro fechasse, o erro estaria na leitura da tensão.}
\stepcounter{questao}
\resposta{$R_{eq}=\SI{3.33}{\ohm}$, $V=\SI{20}{\volt}$;
$\SI{4}{\ampere}$, $\SI{1}{\ampere}$, $\SI{1}{\ampere}$}
\resolucao{(a) $G=0{,}2+0{,}05+0{,}05=\SI{0.3}{\ensuremath{\mathrm{S}}}$,
$R_{eq}=1/0{,}3=\SI{3.33}{\ohm}$ e $V=I_tR_{eq}=6\cdot3{,}33=\SI{20}{\volt}$.\\
(b) Com três ou mais ramos, a forma correta do divisor usa condutâncias:
\[
I_k=I_t\frac{G_k}{G_{eq}}
\;\Rightarrow\;
I_1=6\cdot\frac{0{,}2}{0{,}3}=4,\quad
I_2=I_3=6\cdot\frac{0{,}05}{0{,}3}=\SI{1}{\ampere}.
\]
\textbf{Alerta:} a fórmula $I_1=I_tR_2/(R_1+R_2)$ \emph{não} se generaliza para
três ramos --- ela é um caso particular que só funciona com dois.\\
(c) \emph{LKC:} $4+1+1=6$ \checkmark. \emph{Lei de Ohm por ramo:}
$20/5=4$ e $20/20=1$ \checkmark.}
\stepcounter{questao}
\resposta{Ramos restantes inalterados ($\SI{3}{\ampere}$ cada);
$I_t$: $12\to\SI{6}{\ampere}$; $R_{eq}$: $5\to\SI{10}{\ohm}$}
\resolucao{Com fonte ideal, $V=\SI{60}{\volt}$ permanece fixo. Como cada ramo
depende \emph{só} da tensão comum e do próprio resistor, os dois ramos íntegros
continuam com $60/20=\SI{3}{\ampere}$ --- \textbf{nada muda para eles}.\\
O que muda é o agregado: $I_t=3+3=\SI{6}{\ampere}$ (queda de $50\%$) e
$R_{eq}=\SI{10}{\ohm}$.\\
\textbf{Por isso a ligação residencial é em paralelo:} a queima de uma lâmpada
não afeta as demais. Em série, uma falha derruba o circuito inteiro.\\
\emph{Ressalva importante:} isso depende de a fonte ser ideal. Com resistência
interna, $V$ sobe quando um ramo abre --- caso tratado no bloco de desafio.}
\stepcounter{questao}
\resposta{$R_{eq}\to0$; $I_t\to\infty$ (limitada só pela fonte e pela fiação);
os demais ramos ficam sem tensão}
\resolucao{Um ramo em curto tem $R=0$, portanto $G\to\infty$. Como
$G_{eq}=\sum G_k$, a condutância total explode e $R_{eq}\to0$:
\[
R_{eq}=\frac{1}{\infty+G_2+G_3}\to0.
\]
Com fonte ideal, $I_t=V/R_{eq}\to\infty$. Na prática, a corrente é limitada
pela resistência interna da fonte e pela impedância da fiação --- mas ainda
assim atinge dezenas ou centenas de ampères.\\
\textbf{Efeito sobre os demais:} o curto derruba a tensão do barramento a
praticamente zero. Os outros ramos, embora íntegros, \emph{param de funcionar}
--- eles são vítimas, não causa.\\
\textbf{Proteção:} é este o cenário que o disjuntor combate. Diferentemente do
ramo aberto (falha benigna e localizada), o curto em um ramo compromete o
sistema inteiro e libera energia suficiente para incendiar a instalação. Daí a
assimetria de tratamento: monitoram-se \emph{sobrecorrentes}, não
subcorrentes.}
\stepcounter{questao}
\resposta{Quatro em paralelo; $P_{\max}=\SI{4}{\watt}$;
$V_{\max}=\SI{14.8}{\volt}$}
\resolucao{\textbf{Valor:} $220/n=55\Rightarrow n=4$.\\
\textbf{Potência:} sendo iguais e sob a mesma tensão, os quatro dividem a carga
térmica igualmente, então $P_{tot}=4\times1=\SI{4}{\watt}$ \checkmark.\\
\textbf{Tensão máxima:} limita-se pelo componente individual:
\[
V_{\max}=\sqrt{P_1R_1}=\sqrt{1\cdot220}=\SI{14.8}{\volt}.
\]
\textbf{Conferência:} $P_{tot}=V^{2}/R_{eq}=220/55=\SI{4}{\watt}$ \checkmark.\\
\textbf{Observação de projeto:} o paralelo multiplica a capacidade térmica
\emph{e} divide o valor por $n$; a série multiplica ambos. Se o objetivo exige
valor \emph{maior} com mais potência, use série; valor \emph{menor}, paralelo.
Para manter o valor e aumentar a potência, é preciso uma matriz
série--paralelo de $n\times n$ unidades.}
\stepcounter{questao}
\resposta{$G=\SI{0.25}{\ensuremath{\mathrm{S}}}$; $R_{eq}=\SI{4}{\ohm}$}
\resolucao{
\[
G=0{,}100+0{,}050+0{,}050+0{,}025+0{,}025=\SI{0.25}{\ensuremath{\mathrm{S}}}
\;\Longrightarrow\;
R_{eq}=\SI{4}{\ohm}.
\]
\textbf{Três vantagens da condutância.} \emph{(i)} A operação é uma soma
simples, com um único inverso no fim --- contra quatro aplicações sucessivas da
regra produto/soma. \emph{(ii)} A distribuição de corrente sai de graça:
$I_k/I_t=G_k/G_{eq}$, aqui $40\%$, $20\%$, $20\%$, $10\%$ e $10\%$.
\emph{(iii)} Para incerteza, a propagação em $G$ é \emph{linear}, ao passo que
em $R$ envolve derivadas de uma função racional.\\
É também a razão pela qual é a análise nodal --- construída sobre a LKC --- adota
condutâncias como grandeza natural.}
\stepcounter{questao}
\resposta{(a) $R_{sh}=\SI{1}{\ohm}$; (b) $\SI{0.99}{\ohm}$}
\resolucao{(a) O shunt conduz $100-1=\SI{99}{\milli\ampere}$:
\[
R_{sh}=\frac{R_gI_g}{I_{sh}}=\frac{99\cdot1}{99}=\SI{1}{\ohm}.
\]
(b) $R_A=R_g\parallel R_{sh}=\dfrac{99\cdot1}{100}=\SI{0.99}{\ohm}$.\\
(c) O amperímetro é inserido \textbf{em série} no ramo sob medição, somando-se
à resistência do circuito. Se $R_A$ não for desprezível diante do circuito, a
própria medição altera a corrente que se pretende medir --- o chamado
\emph{erro de inserção}.\\
\textbf{Exemplo:} medindo um ramo de $\SI{100}{\ohm}$, o erro é
$0{,}99/100\approx1\%$; medindo um de $\SI{1}{\ohm}$, seria de $50\%$.\\
Note a dualidade: o \emph{amperímetro} deve ter resistência quase nula e vai em
série; o \emph{voltímetro}, resistência quase infinita e vai em paralelo.}
\stepcounter{questao}
\resposta{(a) $\SI{20}{\milli\ampere}$ e $\SI{2.02}{\ampere}$;
(b) $\SI{99.01}{\ohm}$, erro de $1\%$; (c) risco de choque}
\resolucao{(a) $I_{carga}=200/100=\SI{2}{\ampere}$;
$I_{fuga}=200/10000=\SI{20}{\milli\ampere}$; total $\SI{2.02}{\ampere}$.\\
(b) $R_{eq}=\dfrac{100\cdot10000}{10100}=\SI{99.01}{\ohm}$ --- desvio de
$\approx1\%$ em relação a $\SI{100}{\ohm}$.\\
(c) Do ponto de vista \emph{elétrico}, $1\%$ é irrelevante: nenhum instrumento
de painel acusaria a anomalia. Do ponto de vista da \textbf{segurança}, é
gravíssimo: $\SI{20}{\milli\ampere}$ atravessando um corpo humano já provoca
tetanização muscular, e a partir de $\SI{30}{\milli\ampere}$ o risco de
fibrilação é real.\\
\textbf{Por isso existe o DR (dispositivo diferencial residual):} ele não mede
a corrente total, e sim o \emph{desequilíbrio} entre fase e neutro --- que é
exatamente a corrente de fuga. Um DR de $\SI{30}{\milli\ampere}$ detecta o que
os $\SI{2}{\ampere}$ da carga mascaram completamente.}
\stepcounter{questao}
\resposta{$R_{eq}=\SI{80}{\ohm}$, entre $\SI{76}{\ohm}$ e $\SI{84}{\ohm}$
($\pm5\%$)}
\resolucao{Nominal: $R_{eq}=\dfrac{100\cdot400}{500}=\SI{80}{\ohm}$.\\
Como $R_{eq}$ cresce monotonicamente com $R_1$ e com $R_2$, os extremos ocorrem
com \emph{ambos} nos extremos:
\[
R_{eq}^{\min}=\frac{95\cdot380}{475}=\SI{76}{\ohm},
\qquad
R_{eq}^{\max}=\frac{105\cdot420}{525}=\SI{84}{\ohm}.
\]
Ou seja, $80\pm4$, isto é, $\pm5\%$ --- exatamente a tolerância das parcelas.\\
\textbf{Resultado notável:} o paralelo \emph{também} preserva a tolerância
relativa, assim como a série. Isso não é coincidência, e a razão é demonstrada
no bloco de desafio.}
\stepcounter{questao}
\resposta{$\SI{50}{\ohm}$ e $\SI{55.36}{\ohm}$; variação de $+10{,}7\%$}
\resolucao{A $\SI{20}{\celsius}$: $R_{eq}=100/2=\SI{50}{\ohm}$.\\
A $\SI{80}{\celsius}$: $R_1=100[1+0{,}004(60)]=\SI{124}{\ohm}$ e
\[
R_{eq}=\frac{124\cdot100}{224}=\SI{55.36}{\ohm}.
\]
Variação: $(55{,}36-50)/50=+10{,}7\%$.\\
\textbf{Comparação instrutiva:} $R_1$ variou $24\%$, mas $R_{eq}$ variou
$10{,}7\%$ --- o ramo estável "amortece" a deriva. Curiosamente, o valor é o
mesmo obtido para a \emph{corrente} da associação série equivalente: em ambos
os casos o elemento estável divide o efeito aproximadamente pela metade, porque
as duas resistências são comparáveis.}
\stepcounter{questao}
\resposta{(a) $\SI{9.901}{\ohm}$; (b) $\SI{9.910}{\ohm}$; (c) variação de
apenas $0{,}09\%$}
\resolucao{(a) $R_{eq}=\dfrac{1000\cdot10}{1010}=\SI{9.901}{\ohm}$.\\
(b) Com $R_1=\SI{1100}{\ohm}$:
$R_{eq}=\dfrac{1100\cdot10}{1110}=\SI{9.910}{\ohm}$.\\
(c) Uma variação de $10\%$ em $R_1$ produziu $0{,}09\%$ em $R_{eq}$ ---
atenuação de mais de cem vezes. O ramo dominante ($\SI{10}{\ohm}$) fixa
praticamente sozinho o resultado.\\
\textbf{Uso prático:} é assim que se faz um ajuste \emph{fino}. Colocando um
resistor grande em paralelo com um pequeno, obtém-se um deslocamento pequeno e
controlado --- técnica usada em calibração de shunts e de divisores de
precisão. A quantificação exata dessa sensibilidade aparece no bloco de
desafio.}
\stepcounter{questao}
\resposta{(a) $\SI{5.5}{\volt}$; $\SI{214.5}{\volt}$;
(b) $\SI{151.25}{\watt}$; (c) $2{,}5\%$ --- atende}
\resolucao{(a) $\Delta V=R_fI=0{,}2\cdot27{,}5=\SI{5.5}{\volt}$, logo as cargas
recebem $220-5{,}5=\SI{214.5}{\volt}$.\\
(b) $P_f=R_fI^{2}=0{,}2(27{,}5)^{2}=\SI{151.25}{\watt}$ --- dissipados em puro
aquecimento do cabo.\\
(c) $5{,}5/220=2{,}5\%$, abaixo do limite típico de $4\%$ para circuitos
terminais. \checkmark\\
\textbf{Observação sutil:} a fiação está em \emph{série} com o conjunto, mesmo
que as cargas estejam em paralelo entre si. É por isso que ligar mais uma carga
--- que não afeta as demais em uma rede ideal --- na prática \emph{reduz} a
tensão de todas: a corrente total sobe e a queda na fiação também. O efeito é
visível quando as luzes piscam ao ligar um chuveiro.}
\stepcounter{questao}
\resposta{$R=\SI{300}{\ohm}$}
\resolucao{
\[
\frac{1}{75}=\frac{1}{100}+\frac{1}{R}
\;\Longrightarrow\;
\frac{1}{R}=\frac{1}{75}-\frac{1}{100}=\frac{4-3}{300}
\;\Longrightarrow\;
R=\SI{300}{\ohm}.
\]
\textbf{Viabilidade:} $\SI{300}{\ohm}$ existe na série E24, e o alvo
($\SI{75}{\ohm}$) está a $25\%$ do resistor de partida --- folga confortável.
A precisão do resultado depende principalmente do resistor de
$\SI{100}{\ohm}$, que domina o paralelo.\\
\textbf{Limite do método:} só se pode \emph{reduzir} o valor de partida; alvos
acima de $\SI{100}{\ohm}$ exigem série. E, como se verá no desafio, alvos muito
próximos de $\SI{100}{\ohm}$ tornam a solução impraticável.}
\stepcounter{questao}
\resposta{Série: $\SI{90}{\ohm}$, $\SI{1}{\ampere}$, $P=30$ e
$\SI{60}{\watt}$. Paralelo: $\SI{20}{\ohm}$, $\SI{4.5}{\ampere}$,
$P=270$ e $\SI{135}{\watt}$}
\resolucao{\textbf{Série:} $R=\SI{90}{\ohm}$, $I=\SI{1}{\ampere}$ (comum),
$P_1=30$, $P_2=\SI{60}{\watt}$; total $\SI{90}{\watt}$. O \emph{maior}
resistor dissipa mais.\\
\textbf{Paralelo:} $R=\SI{20}{\ohm}$, $I_1=3$, $I_2=\SI{1.5}{\ampere}$,
$I_t=\SI{4.5}{\ampere}$, $P_1=\SI{270}{\watt}$, $P_2=\SI{135}{\watt}$; total
$\SI{405}{\watt}$. O \emph{menor} resistor dissipa mais.\\
\textbf{Razões:} $R_{s}/R_{p}=90/20=4{,}5$, e a potência total é $4{,}5$ vezes
maior no paralelo --- o mesmo fator, pois $P=V^{2}/R$ com $V$ fixo.\\
\textbf{Síntese:} de dois resistores obtêm-se quatro valores
($30$, $60$, $20$ e $\SI{90}{\ohm}$) e duas distribuições opostas de potência.
A ligação é uma variável de projeto tão importante quanto os componentes.}
\stepcounter{questao}
\resposta{(a) $\SI{11}{\ohm}$, $\SI{1}{\ohm}$ e $\SI{0.0991}{\ohm}$;
(c) $\SI{9.9}{\ohm}$, $\SI{0.99}{\ohm}$ e $\SI{0.099}{\ohm}$}
\resolucao{(a) Com $n=I_{\max}/I_g$, vale
$R_{sh}=\dfrac{R_g}{n-1}$:
\[
n=10\Rightarrow\frac{99}{9}=\SI{11}{\ohm};\quad
n=100\Rightarrow\frac{99}{99}=\SI{1}{\ohm};\quad
n=1000\Rightarrow\frac{99}{999}=\SI{0.0991}{\ohm}.
\]
(b) \textbf{Risco:} durante a comutação, a chave passa por um instante em que
\emph{nenhum} shunt está conectado. Toda a corrente de linha --- até
$\SI{1}{\ampere}$ --- atravessaria o galvanômetro, projetado para
$\SI{1}{\milli\ampere}$: destruição imediata.\\
\textbf{Shunt de Ayrton:} os resistores ficam permanentemente ligados em série
entre si e o conjunto em paralelo com o galvanômetro; a chave apenas escolhe
\emph{em que ponto} da cadeia a corrente entra. Assim o galvanômetro nunca
fica desprotegido, mesmo durante a comutação.\\
(c) $R_A=R_g\parallel R_{sh}$: $99\parallel11=\SI{9.9}{\ohm}$;
$99\parallel1=\SI{0.99}{\ohm}$; $99\parallel0{,}0991=\SI{0.099}{\ohm}$.\\
Note que $R_A=R_g/n$ exatamente --- escalas maiores exigem instrumentos de
resistência menor, o que é justamente o desejável.}
\stepcounter{questao}
\resposta{(c) $m$ colunas em paralelo de $n$ unidades em série, com $m=n$}
\resolucao{(a) $G_{eq}=NG\Rightarrow R_{eq}=R/N$. Sob tensão $V$ comum, cada
unidade dissipa $V^{2}/R$; como todas são iguais e submetidas à mesma tensão, o
total é $N$ vezes o individual: $P_{\max}=NP_1$.\\
(b) $R_{eq}=NR$. Com corrente $I$ comum, cada unidade dissipa $RI^{2}$ e o
total é novamente $NP_1$.\\
\textbf{Fato central:} \emph{qualquer} associação de $N$ resistores iguais
suporta $NP_1$, desde que a dissipação se reparta igualmente. O que a topologia
escolhe é o \emph{valor}, não a capacidade térmica.\\
(c) Uma matriz de $n$ unidades em série, repetida em $m$ colunas paralelas:
\[
R_{eq}=\frac{nR}{m},\qquad P_{\max}=mnP_1.
\]
Para preservar o valor, $n=m$; então $R_{eq}=R$ e
$P_{\max}=n^{2}P_1$.\\
\textbf{Exemplo:} para quadruplicar a potência sem alterar o valor, use
$2\times2=4$ unidades. Para $\times9$, use $3\times3=9$. A potência cresce com
o \emph{quadrado} do lado da matriz --- e o custo, também.}
\stepcounter{questao}
\resposta{(a) $0{,}64$; (b) $S_1=R_2/(R_1+R_2)$, $S_2=R_1/(R_1+R_2)$;
(c) $S_1=0{,}8$ e $S_2=0{,}2$}
\resolucao{(a) Derivando o quociente em relação a $R_1$:
\[
\frac{\partial R_{eq}}{\partial R_1}
=\frac{R_2(R_1+R_2)-R_1R_2}{(R_1+R_2)^{2}}
=\frac{R_2^{2}}{(R_1+R_2)^{2}}
=\left(\frac{R_2}{R_1+R_2}\right)^{2}.
\]
Com $R_1=100$ e $R_2=400$: $(400/500)^{2}=0{,}64$ --- adimensional, pois é
$\si{\ohm}$ por $\si{\ohm}$.\\
(b)
\[
S_1=\frac{\partial R_{eq}}{\partial R_1}\cdot\frac{R_1}{R_{eq}}
=\frac{R_2^{2}}{(R_1+R_2)^{2}}\cdot\frac{R_1(R_1+R_2)}{R_1R_2}
=\frac{R_2}{R_1+R_2}.
\]
Por simetria, $S_2=\dfrac{R_1}{R_1+R_2}$, e a soma é $1$ --- consequência de
$R_{eq}$ ser \textbf{homogênea de grau 1}.\\
(c) $S_1=0{,}8$ e $S_2=0{,}2$: um erro de $1\%$ em $R_1$ produz $0{,}8\%$ em
$R_{eq}$, e o mesmo erro em $R_2$ produz apenas $0{,}2\%$.\\
\textbf{Regra de projeto:} o resistor \emph{menor} domina o paralelo e deve
receber a tolerância apertada --- exatamente o inverso da série, onde manda o
maior. E se $R_2\gg R_1$, então $S_1\to1$ e $S_2\to0$: o ramo grande vira
ajuste fino, o que confirma quantitativamente o observado na questão do
$\SI{1000}{\ohm}$ com $\SI{10}{\ohm}$.}
\stepcounter{questao}
\resposta{(a) $\SI{10}{\ampere}$ e $\SI{0.1}{\volt}$;
(b) $\SI{10.13}{\ampere}$ contra $\SI{9.93}{\ampere}$ --- desvio de $1{,}3\%$;
(c) o mais "fraco" aquece mais}
\resolucao{(a) $I=30/3=\SI{10}{\ampere}$;
$V=0{,}01\cdot10=\SI{0.1}{\volt}$.\\
(b) Com $R_a=\SI{0.0098}{\ohm}$ e $R_b=R_c=\SI{0.0100}{\ohm}$:
\[
G=\frac{1}{0{,}0098}+\frac{2}{0{,}0100}=102{,}04+200{,}00=\SI{302.04}{\ensuremath{\mathrm{S}}},
\]
\[
I_a=30\cdot\frac{102{,}04}{302{,}04}=\SI{10.13}{\ampere},
\qquad
I_b=I_c=30\cdot\frac{100{,}00}{302{,}04}=\SI{9.93}{\ampere}.
\]
\textbf{Verificação:} $10{,}13+9{,}93+9{,}93=\SI{29.99}{\ampere}$ \checkmark.\\
\textbf{Fato tranquilizador:} um desvio de $2\%$ na resistência produziu apenas
$1{,}3\%$ de desequilíbrio --- a divisão é \emph{mais estável} que os
componentes, porque os outros dois ramos "puxam" a média.\\
(c) $P_a=0{,}0098(10{,}13)^{2}=\SI{1.006}{\watt}$ contra
$P_b=0{,}0100(9{,}93)^{2}=\SI{0.986}{\watt}$: o shunt de menor resistência
dissipa $2\%$ a mais. Se tiver coeficiente térmico positivo, aquecerá,
aumentará sua resistência e \emph{devolverá} corrente aos vizinhos --- uma
realimentação \textbf{estabilizadora}. Com coeficiente negativo ocorreria o
oposto: fuga térmica. É por isso que resistores para divisão de corrente são
feitos em ligas de $\alpha$ positivo e pequeno, como a manganina.}
\stepcounter{questao}
\resposta{(a) $G_1\alpha_1+G_2\alpha_2=0$; (b) $R_2=\SI{25}{\ohm}$,
$R_{eq}=\SI{20}{\ohm}$; (c) exata só em primeira ordem}
\resolucao{(a) Partindo de $G_{eq}=G_1+G_2$ e de
$R_k(T)=R_k(1+\alpha_k\Delta T)$, tem-se
$G_k\approx G_k(1-\alpha_k\Delta T)$ para $\alpha_k\Delta T\ll1$. Então
\[
G_{eq}(T)\approx(G_1+G_2)-(G_1\alpha_1+G_2\alpha_2)\Delta T,
\]
e a condição de estabilidade é
\[
G_1\alpha_1+G_2\alpha_2=0
\quad\Longleftrightarrow\quad
\frac{\alpha_1}{R_1}+\frac{\alpha_2}{R_2}=0.
\]
(b) $R_2=-\dfrac{\alpha_2R_1}{\alpha_1}=\dfrac{1000\cdot100}{4000}
=\SI{25}{\ohm}$, e $R_{eq}=100\parallel25=\SI{20}{\ohm}$.\\
(c) \emph{Para $\Delta T=\SI{1}{\celsius}$:}
$R_1=\SI{100.4}{\ohm}$, $R_2=\SI{24.975}{\ohm}$, $R_{eq}=\SI{20.000}{\ohm}$ ---
compensação praticamente perfeita.\\
\emph{Para $\Delta T=\SI{50}{\celsius}$:} $R_1=\SI{120}{\ohm}$,
$R_2=\SI{23.75}{\ohm}$, $R_{eq}=\SI{19.83}{\ohm}$ --- desvio de $0{,}9\%$.\\
A compensação é de \textbf{primeira ordem}: excelente perto da temperatura de
projeto, degradando-se quadraticamente conforme se afasta.\\
\textbf{Compare com a série}, cuja condição era $R_1\alpha_1+R_2\alpha_2=0$
(ponderada por resistências). Aqui a ponderação é por \emph{condutâncias} ---
mais uma manifestação da dualidade série/paralelo.}
\stepcounter{questao}
\resposta{(a) $\SI{12}{\volt}$, $\SI{6}{\ampere}$, $\SI{2}{\ampere}$ por ramo;
(b) $\SI{14.4}{\volt}$, $\SI{4.8}{\ampere}$, $\SI{2.4}{\ampere}$ por ramo}
\resolucao{(a) $R_{eq}=6/3=\SI{2}{\ohm}$;
\[
I_t=\frac{24}{2+2}=\SI{6}{\ampere},\qquad
V_{barra}=24-2(6)=\SI{12}{\volt},\qquad
I_{ramo}=\frac{12}{6}=\SI{2}{\ampere}.
\]
(b) Com dois ramos, $R_{eq}=\SI{3}{\ohm}$:
\[
I_t=\frac{24}{2+3}=\SI{4.8}{\ampere},\qquad
V_{barra}=24-2(4{,}8)=\SI{14.4}{\volt},\qquad
I_{ramo}=\frac{14{,}4}{6}=\SI{2.4}{\ampere}.
\]
(c) Os ramos remanescentes \emph{não mudaram de valor}, mas a tensão que os
alimenta subiu de $12$ para $\SI{14.4}{\volt}$ ($+20\%$), e com ela a corrente
de cada um ($+20\%$).\\
\textbf{Mecanismo:} $r$ está em série com o conjunto. Menos ramos $\Rightarrow$
menor corrente total $\Rightarrow$ menor queda interna $\Rightarrow$ maior
tensão disponível.\\
\textbf{Implicação de projeto:} o isolamento entre ramos paralelos é uma
\emph{idealização}. Com fonte real, ramos interagem via $r$ --- e a perda de
uma carga pode sobrecarregar as demais em $20\%$ ou mais. Fontes reguladas
existem justamente para levar $r$ efetiva a valores próximos de zero e
restaurar o desacoplamento.}
\stepcounter{questao}
\resposta{(a) $\SI{900}{\ohm}$, $\SI{1900}{\ohm}$ e $\SI{9900}{\ohm}$;
(c) inviável para alvos próximos do resistor fixo}
\resolucao{(a) De $\dfrac{1}{R_{alvo}}=\dfrac{1}{R_f}+\dfrac{1}{R}$:
\[
R=\frac{R_fR_{alvo}}{R_f-R_{alvo}}.
\]
\[
90\to\frac{100\cdot90}{10}=\SI{900}{\ohm};\quad
95\to\frac{100\cdot95}{5}=\SI{1900}{\ohm};\quad
99\to\frac{100\cdot99}{1}=\SI{9900}{\ohm}.
\]
O denominador $R_f-R_{alvo}$ tende a zero e $R$ \textbf{diverge}.\\
(b) Pela sensibilidade relativa do paralelo,
$S=\dfrac{R}{R_f+R}$. Para $R=\SI{9900}{\ohm}$:
$S=9900/10000=0{,}99$.\\
Ou seja: $1\%$ de erro em $R$ produz $0{,}99\%$ de erro no alvo. Como o ajuste
pretendido é de apenas $1\%$ (de $100$ para $\SI{99}{\ohm}$), \emph{o erro tem
a mesma ordem de grandeza do ajuste} --- ele o anula.\\
(c) \textbf{Inviável.} Com resistor de $5\%$ de tolerância, o "ajuste" de
$\SI{1}{\ohm}$ vem com incerteza de $\SI{0.99}{\ohm}$: o alvo pode ficar em
qualquer ponto entre $98$ e $\SI{100}{\ohm}$. Além disso,
$\SI{9900}{\ohm}$ não é valor comercial, e a resistência de contato já é
comparável ao efeito buscado.\\
\textbf{Regra prática:} ajuste por paralelo só é confiável para correções de
\emph{dezenas} de por cento. Para correções finas, use potenciômetro de ajuste
ou selecione o componente por medição direta.}
\stepcounter{questao}
\resposta{$R_{eq}$ é homogênea de grau 1: $R_{eq}(\lambda R_k)=\lambda
R_{eq}(R_k)$; logo a tolerância relativa se conserva}
\resolucao{(a) Qualquer $R_{eq}$ de rede resistiva é construída por somas
($R_i+R_j$) e por combinações do tipo $\dfrac{R_iR_j}{R_i+R_j}$ --- ambas
\textbf{homogêneas de grau 1}. Multiplicando \emph{todos} os resistores por
$\lambda$:
\[
\lambda R_i+\lambda R_j=\lambda(R_i+R_j),
\qquad
\frac{(\lambda R_i)(\lambda R_j)}{\lambda R_i+\lambda R_j}
=\lambda\frac{R_iR_j}{R_i+R_j}.
\]
Por indução sobre a construção da rede, $R_{eq}(\lambda R_k)=\lambda R_{eq}$.
Fazendo $\lambda=1\pm t$, obtém-se $R_{eq}(1\pm t)$: tolerância relativa
exatamente $t$.\\
(b) \emph{Série:} $100\pm5\%$ e $200\pm5\%$ dão $300\pm\SI{15}{\ohm}=\pm5\%$
\checkmark. \emph{Paralelo:} $100\pm5\%$ e $400\pm5\%$ dão
$80\pm\SI{4}{\ohm}=\pm5\%$ \checkmark. A demonstração explica por que os dois
casos, aparentemente distintos, deram o mesmo resultado.\\
(c) Com tolerâncias \emph{diferentes}, a homogeneidade não se aplica --- os
resistores não escalam pelo mesmo $\lambda$. Vale então a combinação ponderada
pelas sensibilidades:
\[
t_{eq}=\sum_k S_k\,t_k,
\qquad \text{com } \sum_k S_k=1.
\]
Em série, $S_k=R_k/\sum R_j$ (manda o maior); em paralelo,
$S_k=G_k/\sum G_j$ (manda o menor). A restrição $\sum S_k=1$ é justamente a
identidade de Euler para funções homogêneas de grau 1 --- e garante que
$t_{eq}$ jamais ultrapasse a maior das tolerâncias individuais.}
\stepcounter{questao}
\resposta{(a) ramos de $25$, $10$ e $\SI{10}{\ampere}$; geral de
$\SI{32}{\ampere}$; (b) o ramo deve atuar antes do geral}
\resolucao{(a) Cada disjuntor de ramo fica acima da corrente nominal do próprio
circuito: $\SI{25}{\ampere}$, $\SI{10}{\ampere}$ e $\SI{10}{\ampere}$
(valores comerciais). O geral cobre a soma das cargas,
$\SI{27.5}{\ampere}$: disjuntor de $\SI{32}{\ampere}$.\\
(b) \textbf{Seletividade} significa que, diante de uma falta em um ramo,
\emph{apenas} o disjuntor daquele ramo abre --- os demais circuitos continuam
alimentados. Isso exige que a curva tempo--corrente do dispositivo de ramo
esteja inteiramente \emph{abaixo} da do geral.\\
Se a seletividade falhar, um curto em uma tomada derruba a instalação inteira:
além do transtorno, perde-se a informação de \emph{onde} está o defeito, e a
localização passa a exigir religamentos sucessivos.\\
(c) Porque as cargas dificilmente operam \emph{todas} no máximo ao mesmo tempo
--- é o \textbf{fator de demanda}. Aqui $25+10+10=\SI{45}{\ampere}$ contra um
geral de $\SI{32}{\ampere}$: perfeitamente admissível, desde que a demanda
simultânea real permaneça abaixo de $\SI{32}{\ampere}$.\\
\textbf{Consequência do paralelo:} como as correntes de ramo se somam no
alimentador, é o \emph{comportamento estatístico} do conjunto --- e não a soma
aritmética dos máximos --- que dimensiona a entrada. Superdimensionar o geral
para $\SI{45}{\ampere}$ exigiria condutores muito mais caros sem ganho real de
disponibilidade.}

\gabarito[3.2 Circuitos em paralelo]

\subsection{Circuitos mistos}
\stepcounter{questao}
\resposta{(c)}
\resolucao{Reduz-se o circuito por etapas, sempre pelo trecho onde a ligação
(série ou paralelo) é inequívoca --- em geral o mais interno ---, substituindo-o
por um resistor equivalente e redesenhando. Repete-se até restar um único
resistor. A transformação estrela-triângulo (alternativa d) só é necessária
quando \emph{nenhum} trecho é puramente série ou paralelo, como nas pontes.}
\stepcounter{questao}
\resposta{(a) \SI{8}{\ohm}, \SI{3}{\ampere}; (b) \SI{10}{\ohm}, \SI{2}{\ampere}}
\resolucao{\textbf{(a)} Os resistores de \SI{6}{\ohm} e \SI{3}{\ohm} da direita
estão em paralelo: $6\parallel 3 = \SI{2}{\ohm}$. Esse resultado fica em série com
o \SI{6}{\ohm} de entrada: $\Req = 6+2 = \SI{8}{\ohm}$;
$I = 24/8 = \SI{3}{\ampere}$.\\
\textbf{(b)} Os dois \SI{4}{\ohm} em paralelo dão \SI{2}{\ohm}, em série com o
\SI{8}{\ohm}: $\Req = 2+8 = \SI{10}{\ohm}$; $I = 20/10 = \SI{2}{\ampere}$.}
\stepcounter{questao}
\resposta{$\Req=\SI{20}{\ohm}$; $I_t=\SI{2}{\ampere}$; \SI{1}{\ampere} em cada \SI{20}{\ohm}}
\resolucao{Os dois \SI{20}{\ohm} em paralelo equivalem a $20/2 = \SI{10}{\ohm}$,
em série com o \SI{10}{\ohm} de entrada:
$\Req = 10+10 = \SI{20}{\ohm}$; $I_t = 40/20 = \SI{2}{\ampere}$.
Como os ramos são iguais, a corrente total se divide igualmente:
\SI{1}{\ampere} em cada resistor de \SI{20}{\ohm}.\\
\textbf{Confira as tensões:} sobre o \SI{10}{\ohm} caem $10\cdot 2 = \SI{20}{\volt}$;
sobram \SI{20}{\volt} no bloco paralelo, e $20/20 = \SI{1}{\ampere}$ em cada ramo.}
\stepcounter{questao}
\resposta{$R_{AB} = \SI{10}{\ohm}$}
\resolucao{Há dois blocos em paralelo, separados por um resistor em série:
\[
6\parallel 3=\frac{6\cdot3}{9}=\SI{2}{\ohm},
\qquad
10\parallel 15=\frac{10\cdot15}{25}=\SI{6}{\ohm}.
\]
Os três trechos estão em série (o \SI{2}{\ohm} de entrada, o primeiro bloco e o
segundo bloco):
\[
R_{AB}=2+2+6=\SI{10}{\ohm}.
\]}
\stepcounter{questao}
\resposta{$R_{AB} = \SI{4}{\ohm}$}
\resolucao{Resolve-se do trecho mais distante para o mais próximo da entrada:
\begin{align*}
6 \parallel 3 &= \SI{2}{\ohm} && \text{(último ramo)}\\
4 + 2 &= \SI{6}{\ohm} && \text{(em série com o } \SI{4}{\ohm}\text{)}\\
12 \parallel 6 &= \SI{4}{\ohm} && \text{(em paralelo com o } \SI{12}{\ohm}\text{ de entrada)}.
\end{align*}
Logo $R_{AB} = \SI{4}{\ohm}$. O nome "escada" vem justamente desse padrão de
reduzir degrau a degrau.}
\stepcounter{questao}
\resposta{$R_{AB} = \SI{10}{\ohm}$}
\resolucao{O fio ideal liga diretamente o nó A ao nó central: eles passam a ser
\emph{o mesmo nó}. Com isso, os resistores de \SI{10}{\ohm} e \SI{20}{\ohm} do topo
ficam ambos ligados entre A e o nó central --- que agora coincidem ---, ou seja,
têm os dois terminais no mesmo ponto e ficam curto-circuitados. Sobra apenas o
resistor de \SI{10}{\ohm} até B:
\[
R_{AB} = \SI{10}{\ohm}.
\]
\textbf{Lição:} identificar nós em curto \emph{antes} de somar resistências evita
contas longas e erradas.}
\stepcounter{questao}
\resposta{$R_{AB} = \SI{5}{\ohm}$}
\resolucao{Explorando a \emph{simetria} do cubo, injeta-se corrente $I$ em A e
retira-se em B. Por simetria, a corrente se distribui de forma previsível:
o resultado clássico para a resistência entre vértices de uma \emph{aresta} é
\[
R_{AB} = \frac{7}{12}\,R.
\]
Com $R = \SI{6}{\ohm}$: $R_{AB} = \dfrac{7}{12}\cdot6 = \SI{3.5}{\ohm}$.\\
(Observação: os três casos clássicos do cubo são $\tfrac{7}{12}R$ pela aresta,
$\tfrac{3}{4}R$ pela diagonal de face e $\tfrac{5}{6}R$ pela diagonal principal.
Com $R=6$: \SI{3.5}{}, \SI{4.5}{} e \SI{5}{\ohm}, respectivamente.) A técnica-chave
é usar a simetria para agrupar nós de mesmo potencial antes de calcular.}
\stepcounter{questao}
\resposta{(a) \SI{7}{\ohm}; (b) \SI{9}{\ohm}; (c) \SI{10}{\ohm}}
\resolucao{A chave é explorar a \textbf{simetria} para identificar nós de mesmo
potencial, que podem ser unidos sem alterar o circuito.

\textbf{(c) Diagonal principal --- o caso mais simétrico.} Injetando corrente $I$
em A, ela se divide igualmente pelas 3 arestas que saem de A ($I/3$ em cada). Em
seguida, cada uma se divide em duas ($I/6$ nas 6 arestas intermediárias) e,
finalmente, recombina-se em 3 arestas ($I/3$) chegando em B. Somando as quedas ao
longo de um caminho:
\[
U = \frac{I}{3}R+\frac{I}{6}R+\frac{I}{3}R
= IR\left(\frac13+\frac16+\frac13\right)=\frac{5}{6}IR,
\]
logo $R_{AB}=\dfrac{5R}{6}=\dfrac{5\cdot12}{6}=\SI{10}{\ohm}$.

\textbf{(b) Diagonal de face:} por raciocínio análogo,
$R = \dfrac{3R}{4}=\dfrac{3\cdot12}{4}=\SI{9}{\ohm}$.

\textbf{(a) Aresta:} $R = \dfrac{7R}{12}=\dfrac{7\cdot12}{12}=\SI{7}{\ohm}$.

\medskip
\begin{center}\small\sffamily
\begin{tabular}{@{}l c c@{}}
\toprule
\textbf{Par de vértices} & \textbf{Fórmula} & \textbf{Valor ($R=\SI{12}{\ohm}$)}\\
\midrule
Aresta            & $\tfrac{7}{12}R$ & \SI{7}{\ohm}\\
Diagonal de face  & $\tfrac{3}{4}R$  & \SI{9}{\ohm}\\
Diagonal principal& $\tfrac{5}{6}R$  & \SI{10}{\ohm}\\
\bottomrule
\end{tabular}
\end{center}
\textbf{Note a ordenação:} quanto mais "distantes" os vértices, maior a
resistência --- mas a diferença é modesta (7 a 10 $\Omega$), porque há sempre
muitos caminhos paralelos. \emph{A simetria é a ferramenta central: sem ela,
seria preciso resolver um sistema de 7 equações nodais.}}
\stepcounter{questao}
\resposta{$R_{AB}=\SI{6}{\ohm}$; a ponte está equilibrada}
\resolucao{\textbf{Teste de equilíbrio:} os braços opostos satisfazem
\[
\frac{R_{AC}}{R_{CB}}=\frac{4}{8}=\frac12,
\qquad
\frac{R_{AD}}{R_{DB}}=\frac{4}{8}=\frac12.
\]
As razões são \emph{iguais}, logo a ponte está \textbf{equilibrada}: os nós C e D
estão no \emph{mesmo potencial}. Sem diferença de potencial, não circula corrente
pelo resistor central --- ele pode ser \emph{removido} (ou substituído por
qualquer valor, inclusive um curto ou um circuito aberto) sem alterar nada.

Restam dois ramos em série, associados em paralelo:
\[
R_{AB}=(4+8)\parallel(4+8)=\frac{12}{2}=\SI{6}{\ohm}.
\]
\textbf{Consequência prática --- a ponte como instrumento:} é exatamente essa
insensibilidade que torna a ponte de Wheatstone tão precisa. No equilíbrio, o
detector central indica zero \emph{independentemente} de sua própria resistência
e da tensão da fonte. A medição depende apenas das \emph{razões} entre
resistores conhecidos.}
\stepcounter{questao}
\resposta{(a) 3 resistores; (b) 4 resistores; (c) 4 resistores}
\resolucao{A estratégia é decompor o valor-alvo em somas de múltiplos
($nR$: série) e submúltiplos ($R/n$: paralelo) de $R = \SI{10}{\ohm}$.

\medskip
\textbf{(a) \SI{15}{\ohm} --- 3 resistores.}
\[
15 = 10 + 5 = R + \frac{R}{2}.
\]
Um resistor em série com dois em paralelo:
$10 + \dfrac{10\cdot10}{20} = 10+5 = \SI{15}{\ohm}$.

\medskip
\textbf{(b) \SI{25}{\ohm} --- 4 resistores.}
\[
25 = 20 + 5 = 2R + \frac{R}{2}.
\]
Dois em série ($\SI{20}{\ohm}$) somados a dois em paralelo ($\SI{5}{\ohm}$).

\medskip
\textbf{(c) \SI{4}{\ohm} --- 4 resistores.} Aqui a soma direta não funciona; é
preciso um \emph{paralelo de dois grupos}:
\[
\SI{4}{\ohm} = 20 \parallel 5 = \frac{20\cdot5}{25},
\]
ou seja, \textbf{dois em série} ($\SI{20}{\ohm}$) em paralelo com \textbf{dois em
paralelo} ($\SI{5}{\ohm}$). Verificando:
$\dfrac{20\cdot5}{20+5} = \dfrac{100}{25} = \SI{4}{\ohm}$.

\medskip
\textbf{Método geral.} Uma busca exaustiva sobre todas as combinações série e
paralelo confirma que 3, 4 e 4 são de fato os mínimos. O raciocínio de projeto é:
\begin{enumerate}[leftmargin=1.7em,itemsep=1pt,topsep=2pt]
\item se o alvo for \emph{maior} que $R$, comece somando (série);
\item se for \emph{menor} que $R$, comece dividindo (paralelo);
\item se não fechar, combine \emph{grupos} em paralelo --- como no item (c), onde
      $4 < 10$ exigiu paralelo de blocos, e não de resistores individuais.
\end{enumerate}
Essa habilidade é usada na prática quando um valor comercial não existe e precisa
ser sintetizado com os valores disponíveis na série E12/E24.}
\stepcounter{questao}
\resposta{$\Req = \dfrac{R(1+\sqrt5)}{2} \approx 1{,}618\,R$}
\resolucao{\textbf{Passo 1 --- explorar a autossimilaridade.} Se a rede é
infinita, retirar a primeira célula (um $R$ em série e um $R$ em paralelo) deixa
\emph{exatamente a mesma rede}. Logo, o bloco à direita do primeiro par também
vale $\Req$:
\[
\Req = R + \big(R \parallel \Req\big) = R + \frac{R\,\Req}{R+\Req}.
\]
\textbf{Passo 2 --- resolver a equação.} Multiplicando por $(R+\Req)$:
\[
\Req(R+\Req) = R(R+\Req) + R\,\Req
\;\Rightarrow\;
\Req^{\,2} - R\,\Req - R^{2} = 0.
\]
\textbf{Passo 3 --- raiz positiva.} Pela fórmula quadrática,
\[
\Req = \frac{R \pm \sqrt{R^2+4R^2}}{2} = \frac{R(1\pm\sqrt5)}{2}.
\]
Descarta-se a raiz negativa (resistência é positiva):
\[
\boxed{\ \Req = \frac{R(1+\sqrt5)}{2} \approx 1{,}618\,R\ }
\]
\textbf{Observação notável:} o fator $\dfrac{1+\sqrt5}{2}$ é a \textbf{razão
áurea} $\varphi$ --- a mesma constante da sequência de Fibonacci e da geometria
clássica. Ela aparece aqui porque a recorrência do circuito tem a mesma forma
algébrica $x = 1 + 1/x$.}
\stepcounter{questao}
\resposta{$n = 5$ seções já bastam}
\resolucao{Aplicando a recorrência iterativamente com $R=\SI{1}{\ohm}$:
\[
\Req^{(1)}=1,\quad
\Req^{(2)}=1+\frac{1\cdot1}{2}=1{,}5,\quad
\Req^{(3)}=1+\frac{1{,}5}{2{,}5}=1{,}6,\ \dots
\]
Comparando com o limite $\varphi \approx 1{,}618034$, o erro cai
\emph{geometricamente}: já em $n=5$ o desvio é de \SI{0.024}{\percent}, abaixo do
critério de \SI{0.1}{\percent}.\\
\textbf{Lição de engenharia:} redes "infinitas" são uma \emph{idealização útil}.
Na prática, meia dúzia de seções já reproduz o comportamento assintótico dentro
da tolerância dos componentes reais (resistores comerciais têm
\SI{1}{\percent} a \SI{5}{\percent} de tolerância). É por isso que atenuadores e
conversores R-2R usam poucos estágios.}
\stepcounter{questao}
\resposta{(a) $\Req = 2R$ (exato); (b) ver comentário}
\resolucao{(a) Pela autossimilaridade,
\[
\Req = R + \big(2R \parallel \Req\big) = R + \frac{2R\,\Req}{2R+\Req}.
\]
Multiplicando por $(2R+\Req)$:
\[
\Req(2R+\Req) = R(2R+\Req) + 2R\,\Req
\;\Rightarrow\;
\Req^{\,2} - R\,\Req - 2R^{2}=0.
\]
Fatorando: $(\Req - 2R)(\Req + R) = 0$, cuja raiz positiva é
\[
\boxed{\ \Req = 2R\ }
\]
--- um valor \textbf{exato e inteiro}, sem irracionais.\\
(b) Essa é a propriedade notável da rede R-2R: a impedância vista de qualquer
ponto da escada é sempre $2R$, o que a torna \emph{perfeitamente casada} em todos
os estágios. Consequência: a corrente se divide exatamente ao meio a cada nó,
gerando os pesos binários $\tfrac12, \tfrac14, \tfrac18,\dots$ --- exatamente o
que um \textbf{conversor D/A} precisa. Além disso, o circuito usa apenas dois
valores de resistor, o que é decisivo para a precisão em circuitos integrados
(é muito mais fácil casar razões de resistores iguais do que valores absolutos).}
\stepcounter{questao}
\resposta{$\Req = \dfrac{R_s + \sqrt{R_s^2+4R_sR_p}}{2}$; racional se
$R_s^2+4R_sR_p$ for quadrado perfeito}
\resolucao{Da recorrência $\Req = R_s + \dfrac{R_p\Req}{R_p+\Req}$, multiplicando
e reorganizando:
\[
\Req^{\,2} - R_s\,\Req - R_s R_p = 0
\quad\Rightarrow\quad
\Req = \frac{R_s + \sqrt{R_s^{2}+4R_sR_p}}{2}.
\]
O resultado é racional quando o discriminante $\Delta = R_s^2+4R_sR_p$ é um
\emph{quadrado perfeito}. Casos particulares:
\begin{itemize}[leftmargin=1.6em,itemsep=1pt,topsep=2pt]
\item $R_p = R_s$: $\Delta = 5R_s^2 \Rightarrow \Req = \varphi R_s$ (irracional);
\item $R_p = 2R_s$: $\Delta = 9R_s^2 \Rightarrow \Req = 2R_s$ (racional);
\item $R_p = 6R_s$: $\Delta = 25R_s^2 \Rightarrow \Req = 3R_s$ (racional).
\end{itemize}
De modo geral, $R_p = \dfrac{k(k-1)}{1}\,R_s$ com $k$ inteiro fornece
$\Req = k\,R_s$. Verifique para $k=2$: $R_p = 2R_s \Rightarrow \Req = 2R_s$.
\textbf{Este é o tipo de raciocínio que separa o técnico do projetista}: em vez de
calcular um caso, obtém-se a \emph{família} de soluções e escolhe-se a que atende
ao requisito.}
\stepcounter{questao}
\resposta{$R_{A\infty} = \SI{1}{\ohm}$ (exatamente)}
\resolucao{\textbf{Passo 1 --- reconhecer a estrutura.} Cada estágio é um
conjunto de resistores \emph{em paralelo} entre dois barramentos. Estágios
consecutivos ficam \emph{em série}. Logo:
\[
R_{A\infty}=\sum_{n=1}^{\infty} R_n,
\qquad\text{onde } R_n \text{ é o paralelo do estágio } n.
\]

\textbf{Passo 2 --- calcular cada estágio pela condutância.} Em paralelo, somam-se
as \emph{condutâncias}. Como o resistor $k$ do estágio $n$ vale
$\dfrac{1}{\binom{n}{k}}$, sua condutância é exatamente $\binom{n}{k}$:
\[
G_n=\sum_{k=0}^{n}\binom{n}{k}.
\]

\textbf{Passo 3 --- a identidade decisiva.} A soma de uma linha inteira do
triângulo de Pascal é uma potência de dois:
\[
\sum_{k=0}^{n}\binom{n}{k}=2^{\,n}
\]
(consequência do binômio de Newton com $x=y=1$: $(1+1)^n = 2^n$). Portanto
\[
G_n = 2^{\,n}
\qquad\Longrightarrow\qquad
R_n=\frac{1}{2^{\,n}}\ \si{\ohm}.
\]
Conferindo os primeiros estágios:
\begin{center}\small\sffamily
\begin{tabular}{@{}c l c c@{}}
\toprule
$n$ & \textbf{Condutâncias} & $G_n$ & $R_n$ ($\Omega$)\\
\midrule
1 & $1,1$              & $2$  & $1/2$\\
2 & $1,2,1$            & $4$  & $1/4$\\
3 & $1,3,3,1$          & $8$  & $1/8$\\
4 & $1,4,6,4,1$        & $16$ & $1/16$\\
5 & $1,5,10,10,5,1$    & $32$ & $1/32$\\
\bottomrule
\end{tabular}
\end{center}

\textbf{Passo 4 --- somar a série.} Os estágios estão em série, então basta somar:
\[
R_{A\infty}=\sum_{n=1}^{\infty}\frac{1}{2^{\,n}}
=\frac12+\frac14+\frac18+\frac{1}{16}+\cdots
\]
É uma \textbf{série geométrica} de razão $q=\tfrac12$ e primeiro termo
$a_1=\tfrac12$:
\[
R_{A\infty}=\frac{a_1}{1-q}=\frac{1/2}{1-1/2}=1.
\]
\[
\boxed{\;R_{A\infty}=\SI{1}{\ohm}\;}
\]

\textbf{Convergência.} Somando apenas os primeiros estágios:
\[
0{,}5 \to 0{,}75 \to 0{,}875 \to 0{,}9375 \to 0{,}96875 \to \cdots \to 1.
\]
Com 8 estágios já se atinge $\SI{99.6}{\percent}$ do valor final --- o erro cai
pela metade a cada estágio acrescentado.

\medskip
\textbf{Por que esta questão é notável.} Ela reúne, em um único problema:
\begin{itemize}[leftmargin=1.7em,itemsep=1pt,topsep=2pt]
\item associação \emph{paralelo} (soma de condutâncias) e \emph{série}
      (soma de resistências), nas duas direções do circuito;
\item uma identidade combinatória (soma da linha do triângulo de Pascal $=2^n$);
\item a convergência de uma série geométrica infinita.
\end{itemize}
Um circuito de complexidade infinita colapsa em uma resposta \textbf{exata e
unitária} --- resultado que só se enxerga reconhecendo o padrão, jamais
calculando estágio a estágio.}
\stepcounter{questao}
\resposta{$R_{eq}=\SI{20}{\ohm}$}
\resolucao{Reduza primeiro o bloco \emph{interno} (o paralelo):
\[
R_{23}=\frac{20\cdot20}{40}=\SI{10}{\ohm}.
\]
Depois o bloco externo (a série):
\[
R_{eq}=10+10=\SI{20}{\ohm}.
\]
A ordem importa: somar os três de uma vez daria $\SI{50}{\ohm}$, resultado sem
sentido físico.}
\stepcounter{questao}
\resposta{$R_{eq}=\SI{25}{\ohm}$}
\resolucao{Dois blocos paralelos independentes, depois somados:
\[
R_{a}=\frac{10}{2}=\SI{5}{\ohm},
\qquad
R_{b}=\frac{30\cdot60}{90}=\SI{20}{\ohm},
\]
\[
R_{eq}=5+20=\SI{25}{\ohm}.
\]
Blocos que não compartilham nós podem ser reduzidos em qualquer ordem entre si
--- só a hierarquia interna de cada um é obrigatória.}
\stepcounter{questao}
\resposta{(b)}
\resolucao{O critério é \textbf{topológico}, não gráfico: dois elementos ligados
ao mesmo par de nós estão submetidos à mesma tensão --- definição de paralelo.
Já a série exige que compartilhem \emph{um} nó \emph{e} que nenhum outro
elemento se conecte a ele, garantindo a mesma corrente.\\
Não se guie pelo desenho: resistores que "parecem" em série no esquema podem
estar em paralelo se um fio os ligar pelos dois lados.}
\stepcounter{questao}
\resposta{(b)}
\resolucao{Reduz-se sempre \emph{de dentro para fora}, isto é, partindo do
bloco mais distante da fonte (ou mais aninhado) em direção aos terminais. Não
existe regra fixa de "paralelo antes de série": a ordem é ditada pela
hierarquia da própria rede.\\
A alternativa (d) falha, por exemplo, em $R_1+(R_2+R_3)\parallel R_4$, onde a
\emph{série} $R_2+R_3$ precisa ser feita antes do paralelo.}
\stepcounter{questao}
\resposta{$R_{eq}=\SI{15}{\ohm}$}
\resolucao{Cada ramo é reduzido primeiro:
$R_a=10+20=\SI{30}{\ohm}$ e $R_b=15+15=\SI{30}{\ohm}$. Depois o paralelo:
\[
R_{eq}=\frac{30}{2}=\SI{15}{\ohm}.
\]
Aqui a série \emph{obrigatoriamente} vem antes: só depois de reduzidos os dois
ramos é que eles ficam ligados ao mesmo par de nós.}
\stepcounter{questao}
\resposta{$R_{eq}=\SI{16}{\ohm}$}
\resolucao{Três níveis, do mais interno ao mais externo:
\[
R_{34}=4+8=\SI{12}{\ohm};
\qquad
R_{234}=\frac{6\cdot12}{18}=\SI{4}{\ohm};
\qquad
R_{eq}=12+4=\SI{16}{\ohm}.
\]}
\stepcounter{questao}
\resposta{$I=\SI{1.2}{\ampere}$; $U_{par}=\SI{12}{\volt}$}
\resolucao{Já se sabe que $R_{eq}=\SI{20}{\ohm}$, logo
\[
I=\frac{24}{20}=\SI{1.2}{\ampere}.
\]
Toda a corrente atravessa $R_1$, e depois o bloco paralelo:
\[
U_{par}=R_{23}\,I=10\cdot1{,}2=\SI{12}{\volt}.
\]
\textbf{Conferência (LKT):} $U_1=10\cdot1{,}2=\SI{12}{\volt}$, e
$12+12=\SI{24}{\volt}$ \checkmark.}
\stepcounter{questao}
\resposta{(a)}
\resolucao{Se esse resistor está \emph{em série} com a fonte --- isto é, se não
há bifurcação antes dele --- toda a corrente que sai da fonte passa por ele.
Por isso ele é o primeiro a ser recuperado na etapa de expansão, e costuma ser
também o mais solicitado termicamente.\\
Note a ressalva: se houver um nó entre a fonte e o resistor, a alternativa (c)
passaria a valer. O critério é sempre topológico.}
\stepcounter{questao}
\resposta{$I_t=\SI{1.2}{\ampere}$; $U_1=\SI{12}{\volt}$;
$U_{23}=\SI{12}{\volt}$; $I_2=I_3=\SI{0.6}{\ampere}$}
\resolucao{\textbf{Redução:} $R_{23}=\SI{10}{\ohm}$, $R_{eq}=\SI{20}{\ohm}$,
$I_t=24/20=\SI{1.2}{\ampere}$.\\
\textbf{Expansão:} $U_1=10(1{,}2)=\SI{12}{\volt}$;
$U_{23}=24-12=\SI{12}{\volt}$;
\[
I_2=I_3=\frac{12}{20}=\SI{0.6}{\ampere}.
\]
\textbf{Verificação (LKC):} $0{,}6+0{,}6=\SI{1.2}{\ampere}$ \checkmark.
Como $R_2=R_3$, a corrente se divide igualmente --- resultado que a simetria
antecipa sem nenhum cálculo.}
\stepcounter{questao}
\resposta{$R_{eq}=\SI{25}{\ohm}$; $I_t=\SI{0.96}{\ampere}$;
$I_2=\SI{0.72}{\ampere}$, $I_3=\SI{0.24}{\ampere}$}
\resolucao{$R_{23}=\dfrac{20\cdot60}{80}=\SI{15}{\ohm}$ e
$R_{eq}=10+15=\SI{25}{\ohm}$, logo $I_t=24/25=\SI{0.96}{\ampere}$.\\
$U_{23}=15(0{,}96)=\SI{14.4}{\volt}$, e daí
\[
I_2=\frac{14{,}4}{20}=\SI{0.72}{\ampere},
\qquad
I_3=\frac{14{,}4}{60}=\SI{0.24}{\ampere}.
\]
\textbf{Divisor de corrente (alternativa):}
$I_2=0{,}96\cdot\frac{60}{80}=\SI{0.72}{\ampere}$ \checkmark. Note que a razão
$I_2/I_3=3$ é exatamente $R_3/R_2$ --- inversa das resistências.}
\stepcounter{questao}
\resposta{$\SI{14.4}{\watt}$; $\SI{7.2}{\watt}$; $\SI{7.2}{\watt}$; total
$\SI{28.8}{\watt}$}
\resolucao{
\[
P_1=R_1I_t^{2}=10(1{,}2)^{2}=\SI{14.4}{\watt};
\qquad
P_2=P_3=\frac{U_{23}^{2}}{20}=\frac{144}{20}=\SI{7.2}{\watt}.
\]
Fonte: $P=24(1{,}2)=\SI{28.8}{\watt}=14{,}4+7{,}2+7{,}2$ \checkmark.\\
\textbf{Observação:} $R_1$ dissipa exatamente metade de toda a potência, embora
seja o menor resistor da rede --- porque é o único que conduz a corrente
inteira. Em redes mistas, "quem aquece mais" depende da posição, não apenas do
valor.}
\stepcounter{questao}
\resposta{$R_{eq}=\SI{10}{\ohm}$; $I_t=\SI{3}{\ampere}$; $\SI{1.5}{\ampere}$
em cada ramo}
\resolucao{Ramos: $R_a=\SI{20}{\ohm}$ e $R_b=\SI{20}{\ohm}$, logo
$R_{eq}=\SI{10}{\ohm}$ e $I_t=30/10=\SI{3}{\ampere}$.\\
Ambos os ramos estão sob os mesmos $\SI{30}{\volt}$:
$I_a=I_b=30/20=\SI{1.5}{\ampere}$.\\
\textbf{Cuidado:} apesar de $R_a=R_b$, as \emph{quedas internas} diferem. No
ramo $a$: $\SI{7.5}{\volt}$ e $\SI{22.5}{\volt}$; no ramo $b$:
$\SI{15}{\volt}$ e $\SI{15}{\volt}$. Ramos equivalentes nos terminais não são
equivalentes internamente.}
\stepcounter{questao}
\resposta{$R=\SI{48}{\ohm}$}
\resolucao{O bloco paralelo deve valer $24-8=\SI{16}{\ohm}$. Em condutâncias:
\[
\frac{1}{R}=\frac{1}{16}-\frac{1}{24}=\frac{3-2}{48}
\;\Longrightarrow\;
R=\SI{48}{\ohm}.
\]
\textbf{Conferência:} $24\parallel48=\dfrac{24\cdot48}{72}=\SI{16}{\ohm}$ e
$8+16=\SI{24}{\ohm}$ \checkmark.\\
Note que a subtração foi feita em \emph{condutâncias}: tentar subtrair
resistências ($16-24$) daria valor negativo, sinal claro de método errado.}
\stepcounter{questao}
\resposta{$R_{eq}=\SI{18}{\ohm}$}
\resolucao{
\[
R_{23}=20+20=\SI{40}{\ohm};
\qquad
R_{234}=\frac{40\cdot10}{50}=\SI{8}{\ohm};
\qquad
R_{eq}=10+8=\SI{18}{\ohm}.
\]
O bloco paralelo resultou em $\SI{8}{\ohm}$ --- menor que o menor dos dois
($\SI{10}{\ohm}$), como sempre deve ocorrer. Esse teste de sanidade a cada
etapa evita que um erro se propague por toda a redução.}
\stepcounter{questao}
\resposta{$V_A=\SI{12}{\volt}$}
\resolucao{O potencial de $A$ em relação ao terra é a tensão sobre o bloco
paralelo, pois é ele que separa $A$ do terra:
\[
V_A=U_{23}=\SI{12}{\volt}.
\]
\textbf{Alternativamente,} descendo da fonte: $V_A=24-U_1=24-12=\SI{12}{\volt}$
\checkmark.\\
Os dois caminhos precisam concordar --- é a LKT em ação. Se discordassem,
haveria erro na redução.}
\stepcounter{questao}
\resposta{$R_{eq}=\SI{10}{\ohm}$; $I_t=\SI{2.4}{\ampere}$;
$P_1=\SI{57.6}{\watt}$}
\resolucao{Um curto no bloco paralelo anula o bloco inteiro: o outro resistor
de $\SI{20}{\ohm}$ fica em paralelo com $\SI{0}{\ohm}$ e deixa de conduzir.
\[
R_{eq}=10+0=\SI{10}{\ohm},
\qquad
I_t=\frac{24}{10}=\SI{2.4}{\ampere},
\qquad
P_1=10(2{,}4)^{2}=\SI{57.6}{\watt}.
\]
$P_1$ saltou de $14{,}4$ para $\SI{57.6}{\watt}$ --- fator \textbf{quatro},
pois a corrente dobrou. Se $R_1$ fosse especificado para
$\SI{25}{\watt}$, queimaria em seguida.\\
\textbf{Padrão de falha:} o curto em um ramo interno transfere toda a
solicitação para o elemento em série --- e o defeito "migra" da causa para a
vítima, dificultando o diagnóstico.}
\stepcounter{questao}
\resposta{$R_{eq}=\SI{30}{\ohm}$; $I_t=\SI{0.8}{\ampere}$;
$P=\SI{12.8}{\watt}$}
\resolucao{Sem um dos ramos, o bloco paralelo passa a valer $\SI{20}{\ohm}$:
\[
R_{eq}=10+20=\SI{30}{\ohm},
\qquad
I_t=\frac{24}{30}=\SI{0.8}{\ampere},
\qquad
U_{par}=20(0{,}8)=\SI{16}{\volt},
\]
\[
P=\frac{16^{2}}{20}=\SI{12.8}{\watt}.
\]
\textbf{Contraste com o curto:} aqui a corrente total \emph{caiu} (de $1{,}2$
para $\SI{0.8}{\ampere}$), mas o resistor sobrevivente passou a dissipar
$12{,}8$ em vez de $\SI{7.2}{\watt}$ --- $78\%$ mais, porque agora assume
sozinho a corrente antes dividida. Falhas de abertura em blocos paralelos
\emph{sobrecarregam} os ramos irmãos.}
\stepcounter{questao}
\resposta{Total $\SI{180}{\watt}$ em ambas; pior componente:
$\SI{90}{\watt}$ em (A) e $\SI{120}{\watt}$ em (B)}
\resolucao{$I_t=60/20=\SI{3}{\ampere}$ e $P_{tot}=60(3)=\SI{180}{\watt}$ nas
duas.\\
\textbf{(A)} $P_{10}=10(3)^{2}=\SI{90}{\watt}$; bloco sob
$\SI{30}{\volt}$, cada $\SI{20}{\ohm}$ com $\SI{45}{\watt}$. Soma:
$90+45+45=\SI{180}{\watt}$ \checkmark.\\
\textbf{(B)} $P_{30}=\dfrac{3600}{30}=\SI{120}{\watt}$;
$P_{60}=\dfrac{3600}{60}=\SI{60}{\watt}$. Soma $\SI{180}{\watt}$ \checkmark.\\
\textbf{Conclusão:} resistência equivalente idêntica, potência total idêntica,
mas o ponto quente difere em $33\%$. A escolha da topologia é decisão térmica,
não elétrica.}
\stepcounter{questao}
\resposta{$I=\SI{2}{\ampere}$; $V=\SI{20}{\volt}$; $P=\SI{40}{\watt}$}
\resolucao{
\[
I=\frac{24}{2+10}=\SI{2}{\ampere};
\qquad
V=24-2(2)=\SI{20}{\volt};
\qquad
P_{rede}=20(2)=\SI{40}{\watt}.
\]
A fonte gera $24(2)=\SI{48}{\watt}$, dos quais $2(2)^{2}=\SI{8}{\watt}$ se
perdem internamente. Rendimento: $40/48=83{,}3\%$.\\
\textbf{Regra útil:} o rendimento é $R_{rede}/(R_{rede}+r)$ --- aqui
$10/12$. Ele depende apenas dessa razão, não da topologia interna da rede.}
\stepcounter{questao}
\resposta{Sim: $P=V^{2}/R_{eq}=VI=\SI{64}{\watt}$}
\resolucao{\textbf{Primeiro caminho:}
$P=\dfrac{V^{2}}{R_{eq}}=\dfrac{32^{2}}{16}=\SI{64}{\watt}$ \checkmark.\\
\textbf{Segundo caminho:} $I=32/16=\SI{2}{\ampere}$ e
$P=VI=32(2)=\SI{64}{\watt}$ \checkmark.\\
\textbf{Terceiro (se os elementos forem conhecidos):} somar $R_kI_k^{2}$ de
cada resistor deve devolver os mesmos $\SI{64}{\watt}$. Essa terceira
verificação é a mais forte, porque testa a \emph{expansão} da rede --- as duas
primeiras só testam $R_{eq}$ e passariam mesmo se a distribuição interna
estivesse errada.}
\stepcounter{questao}
\resposta{$I_3=\SI{0.667}{\ampere}$}
\resolucao{\textbf{Redução:} $R_{34}=\SI{12}{\ohm}$;
$R_{234}=6\parallel12=\SI{4}{\ohm}$; $R_{eq}=12+4=\SI{16}{\ohm}$.\\
\textbf{Expansão:} $I_t=32/16=\SI{2}{\ampere}$;
$U_{234}=4(2)=\SI{8}{\volt}$.\\
Esse é o mesmo $\SI{8}{\volt}$ aplicado ao ramo $R_3+R_4$:
\[
I_3=\frac{8}{12}=\SI{0.667}{\ampere}.
\]
\textbf{Conferência pelo divisor:}
$I_3=I_t\cdot\dfrac{6}{6+12}=2\cdot\dfrac13=\SI{0.667}{\ampere}$ \checkmark. E
$I_2=8/6=\SI{1.333}{\ampere}$, com $0{,}667+1{,}333=\SI{2}{\ampere}$
\checkmark.}
\stepcounter{questao}
\resposta{(a) $R_{eq}=\dfrac{(R_1+R_2)(R_3+R_4)}{R_1+R_2+R_3+R_4}$;
(b) $\SI{12}{\ohm}$; (c) $R_{eq}\to R_1+R_2$}
\resolucao{(a) Reduzindo cada ramo e aplicando produto pela soma:
\[
R_{eq}=\frac{(R_1+R_2)(R_3+R_4)}{(R_1+R_2)+(R_3+R_4)}.
\]
(b) $R_a=\SI{24}{\ohm}$, $R_b=\SI{24}{\ohm}$, logo
$R_{eq}=\dfrac{24\cdot24}{48}=\SI{12}{\ohm}$.\\
(c) Se $R_b\gg R_a$, divida numerador e denominador por $R_b$:
\[
R_{eq}=\frac{R_a}{\dfrac{R_a}{R_b}+1}\longrightarrow R_a=R_1+R_2.
\]
O ramo de alta resistência comporta-se como circuito aberto e não participa.
\textbf{Uso prático:} é essa aproximação que justifica desprezar a resistência
de um voltímetro ligado em paralelo a um trecho de circuito --- desde que ela
seja, digamos, cem vezes maior que o trecho medido.}
\stepcounter{questao}
\resposta{(a) $\SI{5.00}{\volt}$; (b) $\SI{4.72}{\volt}$; (c) $-5{,}5\%$}
\resolucao{(a) $U_{2}=12\cdot\dfrac{100}{240}=\SI{5.00}{\volt}$.\\
(b) A carga entra em \textbf{paralelo} com $R_2$:
\[
R_2'=\frac{100\cdot1000}{1100}=\SI{90.91}{\ohm},
\qquad
U_2=12\cdot\frac{90{,}91}{140+90{,}91}=\SI{4.72}{\volt}.
\]
(c) Erro $=(4{,}72-5{,}00)/5{,}00=-5{,}5\%$.\\
\textbf{Diagnóstico:} a carga de $\SI{1}{\kilo\ohm}$ é apenas $10\times$ maior
que $R_2$ --- insuficiente. A regra de bolso exige
$R_L\ge10R_2$ para erro em torno de $5\%$, e $R_L\ge100R_2$ para erro em torno
de $0{,}5\%$. Um divisor resistivo não é fonte de tensão: ele tem resistência
interna $R_1\parallel R_2=\SI{58.3}{\ohm}$, e é ela que produz a queda ao
ser carregada. O projeto para uma carga \emph{conhecida} é feito no bloco de
desafio.}
\stepcounter{questao}
\resposta{$\SI{24}{\watt}$ em $R_1$; $\SI{12}{\watt}$ em cada
$\SI{12}{\ohm}$; $\SI{8}{\watt}$ em cada $\SI{8}{\ohm}$}
\resolucao{\textbf{Redução:} $12\parallel12=\SI{6}{\ohm}$;
$8\parallel8=\SI{4}{\ohm}$; $R_{eq}=6+6+4=\SI{16}{\ohm}$ e
$I_t=32/16=\SI{2}{\ampere}$.\\
\textbf{Expansão:}
\begin{itemize}
\item $R_1$: conduz $\SI{2}{\ampere}$ $\Rightarrow$
      $P=6(4)=\SI{24}{\watt}$.
\item Bloco de $\SI{12}{\ohm}$: $U=6(2)=\SI{12}{\volt}$; cada ramo com
      $\SI{1}{\ampere}$ $\Rightarrow$ $P=12(1)=\SI{12}{\watt}$.
\item Bloco de $\SI{8}{\ohm}$: $U=4(2)=\SI{8}{\volt}$; cada ramo com
      $\SI{1}{\ampere}$ $\Rightarrow$ $P=8(1)=\SI{8}{\watt}$.
\end{itemize}
\textbf{Balanço:} $24+12+12+8+8=\SI{64}{\watt}=32(2)$ \checkmark.\\
Observe que $R_1$, sendo o menor resistor da rede, é o que mais dissipa ---
porque conduz o dobro da corrente dos demais e a potência vai com $I^{2}$.}
\stepcounter{questao}
\resposta{(a) de $\SI{20}{\ohm}$ a $\SI{40}{\ohm}$; (b) de $\SI{1.5}{\ampere}$
a $\SI{3}{\ampere}$}
\resolucao{(a) Com o reostato em $\SI{0}{\ohm}$, o bloco paralelo é
curto-circuitado e vale $\SI{0}{\ohm}$: $R_{eq}=\SI{20}{\ohm}$.\\
Com o reostato em $\SI{60}{\ohm}$: $30\parallel60=\SI{20}{\ohm}$ e
$R_{eq}=\SI{40}{\ohm}$.\\
(b) $I_{\max}=60/20=\SI{3}{\ampere}$ e
$I_{\min}=60/40=\SI{1.5}{\ampere}$: faixa de $2{:}1$.\\
(c) Dois limites atuam. O resistor \emph{série} de $\SI{20}{\ohm}$ impõe um
teto à corrente ($\SI{3}{\ampere}$), por mais que o bloco paralelo seja
anulado. E o resistor de $\SI{30}{\ohm}$ \emph{em paralelo} com o reostato
impõe um piso a $R_{eq}$: mesmo com o reostato em $\infty$, o bloco não
passaria de $\SI{30}{\ohm}$, e $R_{eq}$ de $\SI{50}{\ohm}$.\\
\textbf{Consequência de projeto:} para ampliar a faixa de controle, é preciso
reduzir o resistor em série \emph{e} aumentar o resistor em paralelo --- alterar
apenas o reostato tem efeito saturante.}
\stepcounter{questao}
\resposta{$V_{\max}=\SI{20}{\volt}$; o crítico é $R_1$}
\resolucao{Expresse cada potência em função da corrente total $I$:
\[
P_1=10I^{2};
\qquad
P_2=P_3=20\left(\frac{I}{2}\right)^{2}=5I^{2}.
\]
$R_1$ dissipa \textbf{o dobro} de cada ramo paralelo --- é ele que satura
primeiro:
\[
10I^{2}=10 \;\Longrightarrow\; I=\SI{1}{\ampere}
\;\Longrightarrow\;
V_{\max}=R_{eq}I=20(1)=\SI{20}{\volt}.
\]
Nessa condição os ramos paralelos operam a $\SI{5}{\watt}$ --- metade do
nominal, ou seja, subutilizados.\\
\textbf{Melhoria possível:} trocar $R_1$ por dois de $\SI{20}{\ohm}$ em
paralelo (mesmos $\SI{10}{\ohm}$, mas $\SI{20}{\watt}$ de capacidade) elevaria
$V_{\max}$ para o limite imposto pelos ramos:
$5I^{2}=10\Rightarrow I=\SI{1.41}{\ampere}$ e
$V_{\max}=\SI{28.3}{\volt}$. Identificar o componente crítico é o primeiro
passo para melhorar qualquer rede.}
\stepcounter{questao}
\resposta{Três unidades: $(\SI{36}{\ohm}+\SI{36}{\ohm})\parallel\SI{36}{\ohm}$}
\resolucao{\textbf{Com duas unidades} só existem dois valores possíveis:
$36+36=\SI{72}{\ohm}$ e $36\parallel36=\SI{18}{\ohm}$. Nenhum é
$\SI{24}{\ohm}$ --- logo duas são insuficientes.\\
\textbf{Com três:}
\[
(36+36)\parallel36=\frac{72\cdot36}{108}=\SI{24}{\ohm}.\ \checkmark
\]
\textbf{Verificação por outro caminho:} três em paralelo dariam
$\SI{12}{\ohm}$; três em série, $\SI{108}{\ohm}$; e
$36+(36\parallel36)=\SI{54}{\ohm}$. Das quatro topologias possíveis com três
unidades, apenas a indicada resolve.\\
\textbf{Observação:} $24=\frac23\cdot36$, e a razão $2/3$ é característica da
combinação "dois em série contra um em paralelo" --- padrão útil de memorizar,
pois $\dfrac{2R\cdot R}{3R}=\dfrac{2R}{3}$ para qualquer $R$.}
\stepcounter{questao}
\resposta{Comparar as tensões medidas com as previstas; o bloco defeituoso
apresenta tensão \emph{maior} que a esperada}
\resolucao{\textbf{Procedimento.} \emph{(i)} Calcule as tensões esperadas em
operação normal. \emph{(ii)} Meça, bloco a bloco, a partir da fonte.
\emph{(iii)} Compare: um aumento de tensão indica bloco com resistência
\emph{maior} que a prevista (abertura); uma queda indica resistência
\emph{menor} (curto).\\
\textbf{Aplicação.} Previsto: $U_1=\SI{12}{\volt}$ e
$U_{par}=\SI{12}{\volt}$. Com um ramo aberto,
$R_{eq}=\SI{30}{\ohm}$, $I=\SI{0.8}{\ampere}$ e:
\[
U_1=10(0{,}8)=\SI{8}{\volt}\ (\downarrow),
\qquad
U_{par}=20(0{,}8)=\SI{16}{\volt}\ (\uparrow).
\]
O voltímetro acusa $\SI{16}{\volt}$ onde se esperavam $\SI{12}{\volt}$ ---
o bloco paralelo é o suspeito.\\
\textbf{Confirmação:} isolando o bloco, um ohmímetro leria
$\SI{20}{\ohm}$ em vez de $\SI{10}{\ohm}$. E \emph{qual} dos dois ramos
abriu se identifica medindo cada um separadamente --- a tensão sozinha não
distingue ramos idênticos em paralelo.\\
\textbf{Por que $U_1$ cai:} a corrente total diminuiu. Note a assinatura
característica: em uma abertura, as tensões se \emph{redistribuem} com soma
constante ($8+16=24$), enquanto em um curto elas colapsam localmente.}
\stepcounter{questao}
\resposta{(a) $\SI{2.4}{\ampere}$, $\SI{48}{\volt}$, $80\%$;
(b) $R_{eq}=\SI{5}{\ohm}$; (c) $80\%$ contra $50\%$}
\resolucao{(a) $I=\dfrac{60}{25}=\SI{2.4}{\ampere}$;
$V=20(2{,}4)=\SI{48}{\volt}$;
$P_{rede}=48(2{,}4)=\SI{115.2}{\watt}$; rendimento
$=20/25=80\%$.\\
(b) A potência na rede é máxima quando $R_{eq}=r=\SI{5}{\ohm}$. Então
$I=60/10=\SI{6}{\ampere}$, $V=\SI{30}{\volt}$ e
$P=\SI{180}{\watt}$.\\
(c) Rendimento na condição de casamento: $5/10=50\%$. A fonte passa a
gerar $\SI{360}{\watt}$, dos quais $\SI{180}{\watt}$ se perdem internamente.\\
\textbf{Conflito de objetivos:} entregar \emph{mais} potência
($180$ contra $\SI{115.2}{\watt}$) custa despachar $\SI{360}{\watt}$ em vez de
$\SI{144}{\watt}$ --- e aquecer a fonte com $\SI{180}{\watt}$. Em eletrônica de
sinal, o casamento vale a pena; em eletrotécnica de potência, jamais.}
\stepcounter{questao}
\resposta{(a) $\SI{10}{\ohm}$ e $\SI{30}{\volt}$;
(b) $\SI{1.5}{\ampere}$ em $R_1$ e em $R_2$, $\SI{0.75}{\ampere}$ em cada
$\SI{20}{\ohm}$ do bloco; (c) $\SI{90}{\watt}$}
\resolucao{(a) O ramo vale $10+(20\parallel20)=10+10=\SI{20}{\ohm}$, e
\[
R_{eq}=20\parallel20=\SI{10}{\ohm}
\;\Longrightarrow\;
U=I_tR_{eq}=3(10)=\SI{30}{\volt}.
\]
\textbf{Diferença essencial em relação à fonte de tensão:} aqui $U$ é
\emph{resultado}, não dado. Se a rede mudar, a tensão muda --- a corrente é que
permanece fixa.\\
(b) Os dois ramos estão sob $\SI{30}{\volt}$ e têm $\SI{20}{\ohm}$ cada:
$I_{R_1}=I_{ramo}=\SI{1.5}{\ampere}$. Dentro do ramo,
$U_{R_2}=10(1{,}5)=\SI{15}{\volt}$ e restam $\SI{15}{\volt}$ para o bloco
paralelo:
\[
I=\frac{15}{20}=\SI{0.75}{\ampere}\ \text{em cada}.
\]
\textbf{LKC:} $0{,}75+0{,}75=\SI{1.5}{\ampere}$ \checkmark; e
$1{,}5+1{,}5=\SI{3}{\ampere}$ \checkmark.\\
(c) $P_1=\dfrac{30^{2}}{20}=45$; $P_2=10(1{,}5)^{2}=22{,}5$; e
$20(0{,}75)^{2}=\SI{11.25}{\watt}$ em cada resistor do bloco:
\[
45+22{,}5+11{,}25+11{,}25=\SI{90}{\watt}=30(3)\ \checkmark
\]
\textbf{Observação:} a potência da fonte de corrente só pôde ser calculada
\emph{depois} de obtida a tensão. Com fonte de tensão ocorre o inverso --- é a
corrente que precisa ser encontrada primeiro.}
\stepcounter{questao}
\resposta{(a) $R_2\approx\SI{25}{\ohm}$, $R_1=\SI{35}{\ohm}$ (erro $1{,}4\%$);
(b) $R_1=\SI{140}{\ohm}$, $R_2=\SI{111}{\ohm}$;
(c) a compensada dissipa muito menos, mas só é exata para aquela carga}
\resolucao{(a) A \textbf{corrente de sangria} é a que circula por $R_2$ sem
passar pela carga. Se ela for muito maior que a da carga, conectar a carga
perturba pouco o divisor. Com $I_L=5/1000=\SI{5}{\milli\ampere}$, exigir
sangria $\ge20I_L=\SI{100}{\milli\ampere}$ dá
$R_2\le5/0{,}1=\SI{50}{\ohm}$. Testando valores (com
$R_1=1{,}4R_2$ para manter a razão $5/12$ em vazio):
\begin{center}
\begin{tabular}{cccc}
\hline
$R_2$ & $R_1$ & $U_{\text{saída}}$ & erro\\
\hline
$\SI{100}{\ohm}$ & $\SI{140}{\ohm}$ & $\SI{4.72}{\volt}$ & $-5{,}5\%$\\
$\SI{50}{\ohm}$  & $\SI{70}{\ohm}$  & $\SI{4.86}{\volt}$ & $-2{,}8\%$\\
$\SI{25}{\ohm}$  & $\SI{35}{\ohm}$  & $\SI{4.93}{\volt}$ & $-1{,}4\%$\\
\hline
\end{tabular}
\end{center}
Escolha $R_2=\SI{25}{\ohm}$, $R_1=\SI{35}{\ohm}$.\\
(b) \textbf{Compensação:} projete com a carga \emph{já incluída}. Exigindo
$12\cdot\dfrac{x}{R_1+x}=5$ com $x=R_2\parallel1000$, vem $x=\dfrac{5R_1}{7}$.
Tomando $R_1=\SI{140}{\ohm}$: $x=\SI{100}{\ohm}$ e
\[
\frac{1}{R_2}=\frac{1}{100}-\frac{1}{1000}
\;\Longrightarrow\;
R_2=\SI{111}{\ohm}.
\]
Verificação: $111\parallel1000=\SI{100}{\ohm}$ e
$12\cdot100/240=\SI{5.00}{\volt}$ \checkmark.\\
(c) \emph{Potência:} na solução (a), a corrente do divisor é
$\approx\SI{200}{\milli\ampere}$ e a dissipação total
$\approx\SI{2.4}{\watt}$; na (b), $\approx\SI{50}{\milli\ampere}$ e
$\approx\SI{0.6}{\watt}$ --- quatro vezes menos.\\
\emph{Robustez:} a solução (b) é exata \textbf{apenas} para
$R_L=\SI{1}{\kilo\ohm}$; se a carga for retirada, a saída sobe para
$12\cdot111/251=\SI{5.31}{\volt}$ ($+6\%$). A solução (a) é menos precisa,
porém quase indiferente à carga. \textbf{A escolha depende de a carga ser fixa
e conhecida ou variável} --- e, se for variável, nenhuma das duas serve:
o caso pede um regulador.}
\stepcounter{questao}
\resposta{(a) $25$, $40$, $60$, $75$, $100$, $133{,}3$, $166{,}7$, $250$ e
$\SI{400}{\ohm}$; (b) $\SI{25}{\watt}$ em ambos}
\resolucao{(a) Enumerando as topologias possíveis:
\begin{center}
\begin{tabular}{ll}
\hline
Arranjo & $R_{eq}$\\
\hline
4 em paralelo & $\SI{25}{\ohm}$\\
$(2\text{ série})\parallel(2\text{ paralelo})$ & $\SI{40}{\ohm}$\\
$[(2\parallel)+1]\parallel 1$ & $\SI{60}{\ohm}$\\
$(3\text{ série})\parallel 1$ & $\SI{75}{\ohm}$\\
$(2\text{ série})\parallel(2\text{ série})$ & $\SI{100}{\ohm}$\\
$(2\parallel)+(2\parallel)$ & $\SI{100}{\ohm}$\\
$(3\parallel)+1$ & $\SI{133.3}{\ohm}$\\
$1+[(2\text{ série})\parallel 1]$ & $\SI{166.7}{\ohm}$\\
$(2\parallel)+2\text{ série}$ & $\SI{250}{\ohm}$\\
4 em série & $\SI{400}{\ohm}$\\
\hline
\end{tabular}
\end{center}
Nove valores distintos.\\
(b) Sob $\SI{100}{\volt}$, ambos os arranjos de $\SI{100}{\ohm}$ dão
$I_t=\SI{1}{\ampere}$ e $P_{tot}=\SI{100}{\watt}$.\\
\emph{Dois ramos de dois em série:} cada ramo conduz $\SI{0.5}{\ampere}$
$\Rightarrow$ $P=100(0{,}25)=\SI{25}{\watt}$ por resistor.\\
\emph{Dois blocos paralelos em série:} cada bloco sob $\SI{50}{\volt}$, cada
resistor com $\SI{0.5}{\ampere}$ $\Rightarrow$ também
$\SI{25}{\watt}$.\\
Os dois arranjos são \textbf{topologicamente distintos, mas termicamente
equivalentes} --- ambos distribuem a dissipação em quatro partes iguais. É o
melhor aproveitamento possível dos componentes.\\
(c) O mínimo é o paralelo total, pois cada resistor acrescentado em paralelo só
pode \emph{reduzir} $R_{eq}$; o máximo é a série total, pela razão simétrica.
Qualquer topologia mista fica necessariamente entre $R/n$ e $nR$.}
\stepcounter{questao}
\resposta{(a) $\SI{30}{\volt}$ e $\SI{1.5}{\ampere}$; (b) pior componente:
$\SI{30}{\watt}$, $\SI{22.5}{\watt}$ e $\SI{11.25}{\watt}$;
(c) a matriz $2\times2$ de $\SI{20}{\ohm}$}
\resolucao{(a) De $P=V^{2}/R_{eq}$ e $P=RI^{2}$:
\[
V=\sqrt{PR_{eq}}=\sqrt{45\cdot20}=\SI{30}{\volt},
\qquad
I=\sqrt{P/R_{eq}}=\sqrt{2{,}25}=\SI{1.5}{\ampere}.
\]
(b) Três topologias que dão $\SI{20}{\ohm}$:
\begin{center}
\begin{tabular}{lll}
\hline
Topologia & Potências & Pior\\
\hline
$\SI{30}{\ohm}\parallel\SI{60}{\ohm}$ & $30$; $\SI{15}{\watt}$
  & $\SI{30}{\watt}$\\
$\SI{10}{\ohm}+(\SI{20}{\ohm}\parallel\SI{20}{\ohm})$
  & $22{,}5$; $11{,}25$; $\SI{11.25}{\watt}$ & $\SI{22.5}{\watt}$\\
$(\SI{20}{\ohm}+\SI{20}{\ohm})\parallel(\SI{20}{\ohm}+\SI{20}{\ohm})$
  & $11{,}25\times4$ & $\SI{11.25}{\watt}$\\
\hline
\end{tabular}
\end{center}
Em todas, o total é $\SI{45}{\watt}$ --- como tem de ser, pois $R_{eq}$ e $V$
são os mesmos.\\
(c) \textbf{A matriz $2\times2$} de resistores de $\SI{20}{\ohm}$. Ela reparte
a dissipação em quatro parcelas \emph{iguais} de $\SI{11.25}{\watt}$,
permitindo componentes de $\SI{25}{\watt}$ (fator $2$ de folga) e um único
valor de estoque. As alternativas exigiriam um componente de $\SI{50}{\watt}$
ou de $\SI{50}{\watt}$/$\SI{25}{\watt}$ mistos, com ponto quente concentrado.\\
\textbf{Princípio geral:} para uma dada $R_{eq}$ e potência total, a topologia
ótima é a que \textbf{equaliza} a dissipação --- matrizes $n\times n$ de
resistores idênticos são o caso limite dessa ideia.}
\stepcounter{questao}
\resposta{$R_{eq}$ medida $=\SI{10}{\ohm}$ contra $\SI{20}{\ohm}$ prevista;
o bloco paralelo está em curto}
\resolucao{(a) $R_{eq}^{med}=24/2{,}4=\SI{10}{\ohm}$, contra
$\SI{20}{\ohm}$ previstos. Resistência \emph{menor} que a nominal
$\Rightarrow$ falha do tipo \textbf{curto}. (Se fosse maior, seria abertura.)\\
(b) A diferença é de exatamente $\SI{10}{\ohm}$, que é o valor do bloco
paralelo --- ele está anulado. \emph{Unicidade:} $R_1$ não pode estar em curto,
pois isso daria $R_{eq}=\SI{10}{\ohm}$ também... \textbf{e aqui está o ponto:}
as duas hipóteses são numericamente idênticas! Um curto em $R_1$ deixaria
$R_{eq}=0+10=\SI{10}{\ohm}$, exatamente o medido.\\
Logo a medição de corrente \textbf{não} distingue as duas falhas --- e afirmar
localização única com esse único dado seria erro de método. O item (c) é
obrigatório, não opcional.\\
(c) \textbf{Medição decisiva:} um voltímetro sobre $R_1$.
\begin{itemize}
\item Se o bloco paralelo está em curto: $U_1=10(2{,}4)=\SI{24}{\volt}$
      (toda a fonte).
\item Se $R_1$ está em curto: $U_1=\SI{0}{\volt}$, e os
      $\SI{24}{\volt}$ aparecem sobre o bloco paralelo.
\end{itemize}
\textbf{Lição de diagnóstico:} uma única medição escalar raramente identifica
uma falha em rede mista, porque topologias diferentes podem produzir a mesma
$R_{eq}$. É preciso pelo menos uma medida de \emph{distribuição} --- uma
tensão interna --- para quebrar a ambiguidade.}
\stepcounter{questao}
\resposta{(a) $R_{AB}=\SI{3}{\ohm}$; (c) $\SI{6}{\ohm}$}
\resolucao{(a) Sejam $A$ e $B$ os terminais e $C$, $D$ os outros dois vértices.
Como $C$ e $D$ ocupam posições \textbf{simétricas} em relação ao par $A$--$B$
(cada um ligado a $A$, a $B$ e ao outro, com resistências iguais), tem-se
$V_C=V_D$. A aresta $CD$ não conduz e pode ser \textbf{removida}.\\
Restam três caminhos em paralelo: a aresta direta $AB$ ($\SI{6}{\ohm}$), o
percurso $A$--$C$--$B$ ($\SI{12}{\ohm}$) e o percurso $A$--$D$--$B$
($\SI{12}{\ohm}$):
\[
R_{AB}=6\parallel12\parallel12=6\parallel6=\SI{3}{\ohm}=\frac{R}{2}.
\]
(b) Com $I_t=\SI{4}{\ampere}$, $U_{AB}=3(4)=\SI{12}{\volt}$:
\begin{itemize}
\item aresta $AB$: $12/6=\SI{2}{\ampere}$;
\item cada percurso lateral: $12/12=\SI{1}{\ampere}$;
\item aresta $CD$: $\SI{0}{\ampere}$.
\end{itemize}
Soma: $2+1+1=\SI{4}{\ampere}$ \checkmark. A aresta direta leva metade da
corrente total.\\
(c) Sem a aresta $AB$, a simetria $C\leftrightarrow D$ \emph{persiste} --- ela
não dependia daquela aresta. Logo $V_C=V_D$ ainda, e $CD$ continua sem
conduzir:
\[
R_{AB}=12\parallel12=\SI{6}{\ohm}=R.
\]
\textbf{Resultado elegante:} remover uma aresta dobra a resistência
($3\to\SI{6}{\ohm}$), porque justamente essa aresta carregava metade da
corrente.\\
\textbf{Por que a simetria funciona:} a solução de uma rede resistiva é
\emph{única}; se permutar $C$ com $D$ devolve o mesmo circuito, a solução
permutada é a mesma solução --- e isso exige $V_C=V_D$. O argumento é o mesmo
que resolve a ponte equilibrada e o cubo de resistores, e dispensa qualquer
sistema de equações.}

\gabarito[3.3 Circuitos mistos]

\subsection{Transformação estrela-triângulo}
\stepcounter{questao}
\resposta{$R_{AB}=\SI{11}{\ohm}$, $R_{BC}=\SI{33}{\ohm}$, $R_{CA}=\SI{16.5}{\ohm}$}
\resolucao{Com $P = R_1R_2+R_2R_3+R_3R_1 = 18+54+27 = \SI{99}{\ohm\squared}$, cada
resistor do triângulo é $P$ dividido pelo resistor \emph{oposto} da estrela:
\[
R_{AB}=\frac{P}{R_3}=\frac{99}{9}=\SI{11}{\ohm},\quad
R_{BC}=\frac{P}{R_1}=\frac{99}{3}=\SI{33}{\ohm},\quad
R_{CA}=\frac{P}{R_2}=\frac{99}{6}=\SI{16.5}{\ohm}.
\]}
\stepcounter{questao}
\resposta{$R_A=\SI{5}{\ohm}$, $R_B=\SI{1.67}{\ohm}$, $R_C=\SI{2.5}{\ohm}$}
\resolucao{Com $S = R_{AB}+R_{BC}+R_{CA} = 10+5+15 = \SI{30}{\ohm}$, cada resistor
da estrela é o produto dos dois lados \emph{adjacentes} ao seu vértice, dividido
por $S$:
\[
R_A=\frac{R_{AB}R_{CA}}{S}=\frac{10\cdot15}{30}=\SI{5}{\ohm},\quad
R_B=\frac{R_{AB}R_{BC}}{S}=\frac{10\cdot5}{30}\approx\SI{1.67}{\ohm},\quad
R_C=\frac{R_{BC}R_{CA}}{S}=\frac{5\cdot15}{30}=\SI{2.5}{\ohm}.
\]}
\stepcounter{questao}
\resposta{Resposta dissertativa --- ver comentário}
\resolucao{Ela permite substituir um agrupamento de três resistores em triângulo
por três em estrela (ou vice-versa) \emph{sem alterar} o comportamento elétrico
visto pelos três terminais externos. Isso "desmancha" configurações que não são
nem série nem paralelo, revelando novas associações simples. O exemplo clássico é
a \textbf{ponte de Wheatstone desequilibrada}: convertendo um dos triângulos em
estrela, os ramos passam a ficar em série/paralelo e o circuito se resolve sem
precisar montar um sistema de equações. Outra aplicação é a análise de redes
trifásicas, onde cargas ligadas em $\Delta$ e em $Y$ são convertidas entre si.}
\stepcounter{questao}
\resposta{$R_{AB} = \SI{4.5}{\ohm}$}
\resolucao{O caminho superior A--C--B está em \emph{série}:
$6+3 = \SI{9}{\ohm}$. Esse ramo está em \emph{paralelo} com o lado direto
$R_{AB} = \SI{9}{\ohm}$:
\[
R_{AB}=9\parallel 9=\frac{9}{2}=\SI{4.5}{\ohm}.
\]
Aqui a redução série/paralelo bastou. Guarde a transformação $\Delta$-$Y$ para
quando há um resistor \emph{atravessando} o triângulo (a ponte), como nas
questões seguintes.}
\stepcounter{questao}
\resposta{$I = \SI{3}{\ampere}$}
\resolucao{A estrela de três resistores iguais $R$ vira um triângulo de
resistores $3R = \SI{6}{\ohm}$ cada. Entre A e B temos o lado direto de
\SI{6}{\ohm} em paralelo com o caminho A--C--B ($6+6 = \SI{12}{\ohm}$):
\[
R_{AB}=6\parallel 12=\frac{6\cdot12}{18}=\SI{4}{\ohm}
\quad\Rightarrow\quad
I=\frac{12}{4}=\SI{3}{\ampere}.
\]
\textbf{Confira:} entre A e B, a estrela original também dá o mesmo resultado ---
$R_A + (R_B \parallel$ com o caminho por C$)$ ---, como esperado, já que as duas
representações são equivalentes vistas dos terminais.}
\stepcounter{questao}
\resposta{(a) $R_A=\SI{3}{\ohm}$, $R_B=\SI{3}{\ohm}$, $R_C=\SI{1.5}{\ohm}$;
(b) $R_{AB}=\SI{6}{\ohm}$}
\resolucao{(a) $S = 12+6+6 = \SI{24}{\ohm}$:
\[
R_A=\frac{R_{AB}R_{CA}}{S}=\frac{12\cdot6}{24}=\SI{3}{\ohm},\quad
R_B=\frac{R_{AB}R_{BC}}{S}=\frac{12\cdot6}{24}=\SI{3}{\ohm},\quad
R_C=\frac{R_{BC}R_{CA}}{S}=\frac{6\cdot6}{24}=\SI{1.5}{\ohm}.
\]
(b) Na estrela, entre A e B a corrente percorre $R_A$ e $R_B$ em série (o braço
$R_C$ fica em aberto, pois C não é terminal):
$R_{AB} = R_A + R_B = 3+3 = \SI{6}{\ohm}$.}
\stepcounter{questao}
\resposta{Equilibrada; $R_{AB} = \SI{18}{\ohm}$}
\resolucao{A ponte está \textbf{equilibrada} quando os produtos dos braços opostos
são iguais:
\[
\frac{R_{AC}}{R_{CB}}=\frac{10}{20}=\frac12,
\qquad
\frac{R_{AD}}{R_{DB}}=\frac{15}{30}=\frac12.
\]
As razões são iguais, logo C e D estão no \emph{mesmo potencial}: o resistor
central de \SI{8}{\ohm} não conduz e pode ser removido. Restam dois ramos em série,
associados em paralelo:
\[
R_{AB}=(10+20)\parallel(15+30)=30\parallel45=\frac{30\cdot45}{75}=\SI{18}{\ohm}.
\]}
\stepcounter{questao}
\resposta{$R_{AB} = \SI{10}{\ohm}$}
\resolucao{Como $\dfrac{6}{18}\neq\dfrac{18}{6}$, a ponte está
\textbf{desequilibrada} e o resistor central conduz: não dá para reduzir só por
série/paralelo. Convertemos o triângulo superior A--C--D
($R_{AC}=6$, $R_{CD}=6$, $R_{AD}=18$) em estrela, com $S = 6+6+18 = \SI{30}{\ohm}$:
\[
R_{A}=\frac{6\cdot18}{30}=\SI{3.6}{\ohm},\quad
R_{C}=\frac{6\cdot6}{30}=\SI{1.2}{\ohm},\quad
R_{D}=\frac{6\cdot18}{30}=\SI{3.6}{\ohm}.
\]
Chamando de N o centro da estrela, os caminhos até B ficam:
$N$--$C$--$B = 1{,}2+18 = \SI{19.2}{\ohm}$ e
$N$--$D$--$B = 3{,}6+6 = \SI{9.6}{\ohm}$, em paralelo:
\[
19{,}2\parallel 9{,}6 = \frac{19{,}2\cdot9{,}6}{28{,}8}=\SI{6.4}{\ohm}.
\]
Somando o braço $R_A$ que liga A ao centro:
\[
R_{AB}=3{,}6+6{,}4=\SI{10}{\ohm}.
\]}
\stepcounter{questao}
\resposta{Ver roteiro comentado}
\resolucao{\textbf{Passo 1.} Converta o triângulo XYZ em estrela, com centro N e
$S = 9+12+15 = \SI{36}{\ohm}$:
\[
R_X=\frac{R_{XY}R_{ZX}}{S}=\frac{9\cdot15}{36}=\SI{3.75}{\ohm},\quad
R_Y=\frac{R_{XY}R_{YZ}}{S}=\frac{9\cdot12}{36}=\SI{3}{\ohm},\quad
R_Z=\frac{R_{YZ}R_{ZX}}{S}=\frac{12\cdot15}{36}=\SI{5}{\ohm}.
\]
\textbf{Passo 2.} Agora cada terminal externo Y e Z leva ao terra por
$R_Y+6$ e $R_Z+6$, ambos em paralelo (pois se encontram no terra):
$(3+6)\parallel(5+6) = 9\parallel 11$. \textbf{Passo 3.} Esse resultado fica em
série com $R_X$, e todo o conjunto em paralelo com o resistor de \SI{6}{\ohm} que
liga X diretamente ao terra:
\[
R_{X\to\text{terra}}
=\Big[\,R_X+\big(9\parallel11\big)\Big]\ \big\|\ 6
=\Big[\,3{,}75+\tfrac{99}{20}\Big]\ \big\|\ 6 .
\]
O que se pretende com esta questão é o \emph{método}: transformar o triângulo em
estrela desmonta o emaranhado e transforma tudo em série/paralelo.}
\stepcounter{questao}
\resposta{(a) $R_x = \SI{694}{\ohm}$; (b) $\approx\SI{0.25}{\percent}$;
(c) ver comentário}
\resolucao{(a) No equilíbrio, os produtos cruzados se igualam:
\[
\frac{R_1}{R_2}=\frac{R_3}{R_x}
\;\Rightarrow\;
R_x = R_3\,\frac{R_2}{R_1} = 347\cdot\frac{200}{100} = \SI{694}{\ohm}.
\]
(b) Para um produto/quociente, as incertezas \emph{relativas} se somam (no pior
caso):
\[
\frac{\delta R_x}{R_x} = \frac{\delta R_3}{R_3}+\frac{\delta R_2}{R_2}
+\frac{\delta R_1}{R_1} = 0{,}05\% + 0{,}1\% + 0{,}1\% = \SI{0.25}{\percent},
\]
ou seja, $R_x = \SI{694}{\ohm} \pm \SI{1.7}{\ohm}$.
(Somando em quadratura, para erros independentes, obtém-se
$\sqrt{0{,}05^2+0{,}1^2+0{,}1^2}\approx\SI{0.15}{\percent}$ --- estimativa mais
realista.)\\
(c) Porque a ponte é um método de \textbf{zero} (ou de detecção nula): no
equilíbrio, o galvanômetro indica \emph{zero}, e a medida não depende da precisão
do próprio galvanômetro nem da tensão da fonte --- depende apenas das
\emph{razões} entre resistores conhecidos. Um ohmímetro comum, ao contrário,
depende da estabilidade da sua fonte interna e da calibração do seu indicador.
\textbf{Este é um princípio geral da metrologia:} comparar contra um padrão é
sempre mais preciso do que medir em escala absoluta.}
\stepcounter{questao}
\resposta{Três braços de $\SI{10}{\ohm}$}
\resolucao{Para o caso equilibrado, a fórmula geral colapsa:
\[
R_Y=\frac{R_\Delta\cdot R_\Delta}{3R_\Delta}=\frac{R_\Delta}{3}
=\frac{30}{3}=\SI{10}{\ohm}.
\]
\textbf{Regra rápida:} $R_Y=R_\Delta/3$. Guarde o sentido --- a estrela é
sempre \emph{menor}, porque seus dois braços em série ($2R_Y=\SI{20}{\ohm}$)
devem igualar o paralelo do triângulo
($30\parallel60=\SI{20}{\ohm}$) \checkmark.}
\stepcounter{questao}
\resposta{Três lados de $\SI{15}{\ohm}$}
\resolucao{
\[
R_\Delta=\frac{3R_Y^{2}}{R_Y}=3R_Y=3(5)=\SI{15}{\ohm}.
\]
\textbf{Regra rápida:} $R_\Delta=3R_Y$ --- inversa da anterior. Se
$R_Y=R_\Delta/3$ e $R_\Delta=3R_Y$ não fossem consistentes, uma das fórmulas
estaria trocada. É a conferência mais barata que existe para o caso
equilibrado.}
\stepcounter{questao}
\resposta{$R_B=\dfrac{R_{AB}\,R_{BC}}{R_{AB}+R_{BC}+R_{CA}}$}
\resolucao{A regra é posicional: o braço ligado ao terminal $B$ é o
\textbf{produto dos dois lados que tocam $B$} --- que são $R_{AB}$ e
$R_{BC}$ --- dividido pela \textbf{soma dos três lados}.\\
\textbf{Como não errar:} identifique o terminal, olhe quais dois lados chegam
nele e multiplique. O lado \emph{oposto} ($R_{CA}$) aparece apenas no
denominador, junto com os outros dois.}
\stepcounter{questao}
\resposta{(a)}
\resolucao{Ambas as configurações usam \textbf{três} resistores e possuem os
mesmos três terminais externos ($A$, $B$, $C$). A diferença é o \textbf{nó
central} da estrela, que não tem correspondente no triângulo.\\
\textbf{Consequência importante:} a equivalência vale apenas para o que se
observa nos \emph{três terminais externos}. O potencial do nó central da estrela
não corresponde a ponto algum do triângulo --- e nenhuma corrente interna é
preservada pela transformação.}
\stepcounter{questao}
\resposta{(a)}
\resolucao{\textbf{Triângulo $\to$ estrela:} produto de dois sobre a soma dos
três --- alternativa (a). O resultado tem dimensão de $\si{\ohm}$
($\si{\ohm^2}/\si{\ohm}$) \checkmark.\\
\textbf{Estrela $\to$ triângulo:} soma dos produtos sobre o braço oposto ---
alternativa (b). Também $\si{\ohm}$ \checkmark.\\
\textbf{Teste de sanidade:} a alternativa (d) daria $\si{\ohm^{2}}$ --- errada
por análise dimensional. E, no caso equilibrado, (a) devolve $R/3$ e (b)
devolve $3R$, o que confirma qual é qual.}
\stepcounter{questao}
\resposta{$R_A=\SI{10}{\ohm}$, $R_B=\SI{6}{\ohm}$, $R_C=\SI{15}{\ohm}$}
\resolucao{A soma dos lados é $20+30+50=\SI{100}{\ohm}$.
\[
R_A=\frac{R_{AB}R_{CA}}{100}=\frac{20(50)}{100}=\SI{10}{\ohm};
\]
\[
R_B=\frac{R_{AB}R_{BC}}{100}=\frac{20(30)}{100}=\SI{6}{\ohm};
\qquad
R_C=\frac{R_{BC}R_{CA}}{100}=\frac{30(50)}{100}=\SI{15}{\ohm}.
\]
\textbf{Conferência de plausibilidade:} cada braço é menor que os lados que o
originaram, e o maior braço ($R_C$) está no terminal tocado pelos dois maiores
lados. Coerente.}
\stepcounter{questao}
\resposta{$20$, $30$ e $\SI{50}{\ohm}$ --- os valores originais}
\resolucao{Soma dos produtos:
\[
\Sigma=R_AR_B+R_BR_C+R_CR_A=60+90+150=\SI{300}{\ohm\squared}.
\]
Cada lado é $\Sigma$ dividido pelo braço \emph{oposto}:
\[
R_{AB}=\frac{300}{R_C}=\frac{300}{15}=\SI{20}{\ohm};
\quad
R_{BC}=\frac{300}{R_A}=\SI{30}{\ohm};
\quad
R_{CA}=\frac{300}{R_B}=\SI{50}{\ohm}.
\]
\textbf{Reversibilidade confirmada.} As duas transformações são inversas exatas
--- fato garantido pela construção, já que ambas derivam do mesmo sistema de
três equações. Fazer a ida e a volta é a verificação mais completa possível de
uma conversão.}
\stepcounter{questao}
\resposta{$\SI{16}{\ohm}$ em ambas}
\resolucao{\textbf{No triângulo:} o lado direto $R_{AB}$ em paralelo com o
caminho $A$--$C$--$B$:
\[
R_{AB}^{eq}=20\parallel(50+30)=\frac{20(80)}{100}=\SI{16}{\ohm}.
\]
\textbf{Na estrela:} com $C$ aberto, o braço $R_C$ não conduz; sobram
$R_A$ e $R_B$ em série:
\[
R_{AB}^{eq}=10+6=\SI{16}{\ohm}.\ \checkmark
\]
\textbf{É exatamente esta a definição de equivalência:} a resistência vista
entre cada par de terminais, com o terceiro aberto, deve coincidir. As três
igualdades desse tipo constituem o sistema de onde as fórmulas são deduzidas ---
como se mostra no bloco de desafio.}
\stepcounter{questao}
\resposta{$R_{AB}=\SI{14.67}{\ohm}$, $R_{BC}=\SI{44}{\ohm}$,
$R_{CA}=\SI{22}{\ohm}$; o maior é $R_{BC}$}
\resolucao{
\[
\Sigma=4(8)+8(12)+12(4)=32+96+48=\SI{176}{\ohm\squared}.
\]
\[
R_{AB}=\frac{176}{12}=\SI{14.67}{\ohm};
\qquad
R_{BC}=\frac{176}{4}=\SI{44}{\ohm};
\qquad
R_{CA}=\frac{176}{8}=\SI{22}{\ohm}.
\]
\textbf{Padrão:} o maior lado do triângulo é o \emph{oposto} ao menor braço da
estrela. Como $R_A=\SI{4}{\ohm}$ é o menor braço, $R_{BC}$ é o maior lado ---
a divisão por um número pequeno produz um resultado grande.\\
\textbf{Interpretação:} um braço pequeno "aproxima" seu terminal do centro,
o que equivale, no triângulo, a afastar os outros dois entre si.}
\stepcounter{questao}
\resposta{$R_A=\dfrac{R_{AB}R_{CA}}{R_{AB}+R_{CA}}$, $R_B=R_C=0$ --- ou seja,
apenas dois resistores em série}
\resolucao{Com $R_{BC}\to\infty$, a soma dos lados também diverge. Nos numeradores:
\[
R_A=\frac{R_{AB}R_{CA}}{R_{AB}+R_{BC}+R_{CA}}\to0
\]
--- mas os braços $R_B$ e $R_C$ contêm $R_{BC}$ no numerador e o limite exige
cuidado:
\[
R_B=\frac{R_{AB}R_{BC}}{R_{AB}+R_{BC}+R_{CA}}\to R_{AB};
\qquad
R_C=\frac{R_{BC}R_{CA}}{R_{AB}+R_{BC}+R_{CA}}\to R_{CA}.
\]
\textbf{Resultado:} $R_A=0$, $R_B=R_{AB}$ e $R_C=R_{CA}$. A estrela degenera em
$R_{AB}$ e $R_{CA}$ ligados diretamente em série pelo terminal $A$ --- que é
exatamente o circuito restante. \checkmark\\
\textbf{Lição:} os casos-limite não quebram as fórmulas; eles as confirmam. Um
lado ausente do triângulo apenas devolve a associação série trivial.}
\stepcounter{questao}
\resposta{$R_A=R_B=0$ e $R_C=\dfrac{R_{BC}R_{CA}}{R_{BC}+R_{CA}}$}
\resolucao{Com $R_{AB}=0$, a soma vale $R_{BC}+R_{CA}$ e:
\[
R_A=\frac{0\cdot R_{CA}}{\Sigma}=0;
\qquad
R_B=\frac{0\cdot R_{BC}}{\Sigma}=0;
\qquad
R_C=\frac{R_{BC}R_{CA}}{R_{BC}+R_{CA}}.
\]
\textbf{Interpretação:} os terminais $A$ e $B$ estão em curto, logo ambos
coincidem com o nó central --- daí $R_A=R_B=0$. E $R_C$ é o paralelo de
$R_{BC}$ com $R_{CA}$, o que é correto: com $A$ e $B$ unidos, os dois resistores
ficam mesmo em paralelo entre $C$ e esse nó comum. \checkmark\\
\textbf{Observação:} a conversão \emph{inversa} não existe aqui --- uma estrela
com dois braços nulos daria $\Sigma=0$ e divisão por zero. A transformação
$Y\to\Delta$ exige todos os braços estritamente positivos.}
\stepcounter{questao}
\resposta{$3100=100\times31$ \checkmark}
\resolucao{\textbf{Lado esquerdo:}
\[
20(30)+30(50)+50(20)=600+1500+1000=3100.
\]
\textbf{Lado direito:} $\sum R_\Delta=100$ e
$\sum R_Y=10+6+15=31$, logo $100(31)=3100$ \checkmark.\\
\textbf{Utilidade:} a identidade permite obter a soma dos braços da estrela
\emph{sem calcular cada um}:
\[
\sum R_Y=\frac{\sum_{\text{pares}}R_\Delta R_\Delta}{\sum R_\Delta}
=\frac{3100}{100}=\SI{31}{\ohm}.
\]
É uma conferência rápida do conjunto das três conversões --- se a soma não
bater, ao menos um braço está errado.}
\stepcounter{questao}
\resposta{A estrela cria um nó interno isolado, liberando associações
série/paralelo}
\resolucao{\textbf{O problema do triângulo:} cada um de seus lados liga dois nós
que \emph{também} recebem outros elementos. Nenhum lado fica em série ou em
paralelo com nada --- a redução trava.\\
\textbf{O que a estrela resolve:} ela substitui os três lados por três braços
que se encontram em um nó \textbf{novo}, ao qual nada mais está ligado. Cada
braço passa a estar em série com o que já existia em seu terminal, e essas
séries podem então ser combinadas em paralelo.\\
\textbf{Quando fazer o contrário:} se a rede contiver uma \emph{estrela} cujo
nó central recebe outros elementos, converta-a em triângulo --- isso elimina o
nó problemático. \textbf{A regra é sempre a mesma:} transforme na direção que
\emph{elimina} o nó que está impedindo a redução.}
\stepcounter{questao}
\resposta{$R_{AB}=\SI{24}{\ohm}$, $R_{BC}=R_{CA}=\SI{12}{\ohm}$}
\resolucao{
\[
\Sigma=6(6)+6(3)+3(6)=36+18+18=\SI{72}{\ohm\squared};
\]
\[
R_{AB}=\frac{72}{3}=\SI{24}{\ohm};
\qquad
R_{BC}=R_{CA}=\frac{72}{6}=\SI{12}{\ohm}.
\]
\textbf{Verificação (entre $A$ e $B$, com $C$ aberto):}\\
\emph{Estrela:} $R_A+R_B=\SI{12}{\ohm}$.\\
\emph{Triângulo:} $24\parallel(12+12)=\dfrac{24(24)}{48}=\SI{12}{\ohm}$
\checkmark.\\
\textbf{Simetria preservada:} como $R_A=R_B$, o triângulo também sai simétrico
em relação ao par $A$--$B$. Qualquer simetria da estrela reaparece no
triângulo --- propriedade útil para conferir resultados rapidamente.}
\stepcounter{questao}
\resposta{(b) $R_{eq}=\SI{8.4}{\ohm}$; (c) $I=\SI{10}{\ampere}$}
\resolucao{(a) Equilíbrio exigiria $\dfrac{6}{12}=\dfrac{12}{6}$, ou seja
$0{,}5=2$ --- falso. \textbf{Desequilibrada}, e o resistor central conduz.
Nenhuma associação série/paralelo é possível.\\
(b) O triângulo formado por $S$, $A$ e $B$ tem lados
$R_{SA}=6$, $R_{SB}=12$ e $R_{AB}=\SI{6}{\ohm}$, com soma
$\SI{24}{\ohm}$. Convertendo em estrela de centro $N$:
\[
R_S=\frac{6(12)}{24}=\SI{3}{\ohm};
\quad
R_A=\frac{6(6)}{24}=\SI{1.5}{\ohm};
\quad
R_B=\frac{12(6)}{24}=\SI{3}{\ohm}.
\]
Agora a rede é redutível: de $N$ partem dois caminhos ao terra,
\[
(1{,}5+12)=\SI{13.5}{\ohm}
\quad\text{e}\quad
(3+6)=\SI{9}{\ohm},
\]
em paralelo: $\dfrac{13{,}5(9)}{22{,}5}=\SI{5.4}{\ohm}$. Somando $R_S$:
\[
R_{eq}=3+5{,}4=\SI{8.4}{\ohm}.
\]
(c) $I=\dfrac{84}{8{,}4}=\SI{10}{\ampere}$.}
\stepcounter{questao}
\resposta{$V_A=\SI{48}{\volt}$, $V_B=\SI{36}{\volt}$; $R_{eq}=\SI{8.4}{\ohm}$
--- confere}
\resolucao{Com $V_S=\SI{84}{\volt}$ conhecido:
\[
\begin{cases}
V_A\left(\tfrac16+\tfrac1{12}+\tfrac16\right)-\tfrac16V_B=\tfrac{84}{6}=14\\[4pt]
-\tfrac16V_A+V_B\left(\tfrac1{12}+\tfrac16+\tfrac16\right)=\tfrac{84}{12}=7
\end{cases}
\]
Resolvendo: $V_A=\SI{48}{\volt}$ e $V_B=\SI{36}{\volt}$.\\
\textbf{Corrente total} (a que sai de $S$):
\[
I=\frac{84-48}{6}+\frac{84-36}{12}=6+4=\SI{10}{\ampere}
\;\Longrightarrow\;
R_{eq}=\frac{84}{10}=\SI{8.4}{\ohm}.\ \checkmark
\]
\textbf{Corrente no ramo central:} $(48-36)/6=\SI{2}{\ampere}$, de $A$ para
$B$ --- confirmando o desequilíbrio.\\
\textbf{Vantagem da conversão:} ela entrega $R_{eq}$ sem resolver sistema
algum. \textbf{Vantagem da nodal:} ela entrega \emph{também} os potenciais
internos e a corrente do ramo central --- que a conversão destrói. Escolha
conforme o que se precisa saber.}
\stepcounter{questao}
\resposta{$R_A=R_B=\SI{49.98}{\milli\ohm}$, $R_C=\SI{49.98}{\ohm}$;
$R_A+R_B\approx R_{AB}$}
\resolucao{(a) $\Sigma=0{,}1+100+100=\SI{200.1}{\ohm}$:
\[
R_A=R_B=\frac{0{,}1(100)}{200{,}1}=\SI{49.98}{\milli\ohm};
\qquad
R_C=\frac{100(100)}{200{,}1}=\SI{49.98}{\ohm}.
\]
(b) $R_A+R_B=\SI{99.95}{\milli\ohm}\approx R_{AB}=\SI{0.1}{\ohm}$.\\
\textbf{Por quê:} a resistência entre $A$ e $B$ com $C$ aberto vale
$R_{AB}\parallel(R_{BC}+R_{CA})=0{,}1\parallel200\approx\SI{0.1}{\ohm}$, e na
estrela ela é exatamente $R_A+R_B$. O lado pequeno domina o paralelo, e a soma
dos dois braços apenas reproduz esse domínio.\\
\textbf{Consequência prática:} quando um lado do triângulo é muito menor que os
outros, a estrela resultante tem \emph{dois braços minúsculos e um braço
médio}. Os dois braços pequenos costumam ser desprezíveis diante do restante da
rede --- mas verifique antes de descartá-los, porque em redes de baixa
resistência eles podem ser exatamente o que importa.}
\stepcounter{questao}
\resposta{$R_{AB}=\SI{1020}{\ohm}$, $R_{BC}=R_{CA}=\SI{10.2}{\ohm}$;
$R_{AB}\to\infty$}
\resolucao{(a) $\Sigma=10(10)+10(0{,}1)+0{,}1(10)=100+1+1
=\SI{102}{\ohm\squared}$:
\[
R_{AB}=\frac{102}{0{,}1}=\SI{1020}{\ohm};
\qquad
R_{BC}=R_{CA}=\frac{102}{10}=\SI{10.2}{\ohm}.
\]
(b) Com $R_C\to0$, o numerador tende a $R_AR_B=100$ e
\[
R_{AB}=\frac{100}{R_C}\to\infty,
\]
enquanto $R_{BC}\to R_B$ e $R_{CA}\to R_A$.\\
\textbf{Interpretação:} $R_C=0$ significa que o terminal $C$ coincide com o nó
central. No triângulo, isso equivale a $A$ e $B$ se ligarem a $C$ diretamente
por $R_A$ e $R_B$, sem lado direto entre eles --- e $R_{AB}=\infty$ é
exatamente a ausência desse lado. \checkmark\\
\textbf{Nota numérica:} $\SI{1020}{\ohm}$ em paralelo com $\SI{10.2}{\ohm}$
contribui com menos de $1\%$. Nesses casos, o lado enorme pode ser descartado
sem prejuízo --- e vale conferir se descartá-lo simplifica a rede antes de
seguir com o valor exato.}
\stepcounter{questao}
\resposta{As fórmulas são bem condicionadas; a perda de precisão ocorre na
\emph{redução posterior} da rede}
\resolucao{(a) \textbf{Não há risco.} Ambas as fórmulas envolvem apenas
\emph{produtos} e \emph{somas} de quantidades positivas, seguidos de uma
divisão por um número positivo. Não ocorre subtração de valores próximos ---
que é a única fonte de cancelamento catastrófico. Erros relativos de entrada se
propagam com fator próximo de $1$.\\
\textbf{Verificação:} arredondando a estrela do exemplo anterior
($49{,}98\,\mathrm{m}$; $49{,}98\,\mathrm{m}$; $\SI{49.98}{\ohm}$) para três
algarismos e convertendo de volta, recupera-se
$R_{AB}=\SI{0.1000}{\ohm}$ --- erro inferior a $0{,}1\%$.\\
(b) \textbf{A precisão se perde depois.} Situações típicas:
\begin{itemize}
\item Um braço convertido, muito grande, fica em paralelo com um pequeno: sua
contribuição é desprezível, e carregar seus algarismos é inútil --- mas
\emph{descartá-lo cedo demais} pode mascarar um efeito real.
\item O resultado final é uma \emph{diferença} entre valores próximos --- por
exemplo, ao calcular a corrente de uma ponte quase equilibrada, em que
$V_A-V_B$ é pequeno diante de $V_A$.
\end{itemize}
\textbf{Regra prática:} conduza a conversão com todos os algarismos e arredonde
\emph{apenas no resultado final}. E desconfie sempre que a resposta for uma
pequena diferença de números grandes --- ali, e não na transformação, é onde a
precisão evapora.}
\stepcounter{questao}
\resposta{$36$, $110$ e $\SI{56}{\ohm}$; erros de $-0{,}2\%$, $-2{,}0\%$ e
$+1{,}7\%$ nos braços}
\resolucao{(a) Valores E24 mais próximos: $\SI{36}{\ohm}$, $\SI{110}{\ohm}$ e
$\SI{56}{\ohm}$ (desvios de $-1{,}8\%$, $0\%$ e $+1{,}8\%$ nos lados).\\
(b) Convertendo o triângulo comercial de volta a estrela
($\Sigma=36+110+56=\SI{202}{\ohm}$):
\[
R_A=\frac{36(56)}{202}=\SI{9.98}{\ohm};
\quad
R_B=\frac{36(110)}{202}=\SI{19.60}{\ohm};
\quad
R_C=\frac{110(56)}{202}=\SI{30.50}{\ohm}.
\]
Erros: $-0{,}2\%$, $-2{,}0\%$ e $+1{,}7\%$ em relação a $10$, $20$ e
$\SI{30}{\ohm}$.\\
\textbf{Observação importante:} os erros dos braços \emph{não} reproduzem os
erros dos lados --- eles se recombinam. Um lado com erro nulo
($\SI{110}{\ohm}$) contribuiu para dois braços que erraram $2\%$.\\
\textbf{Método correto de avaliação:} nunca julgue o erro pelos componentes
isolados. Converta de volta e compare com o alvo, ou --- melhor ainda --- calcule
a grandeza que realmente interessa (a resistência entre os terminais usados) nos
dois casos.}
\stepcounter{questao}
\resposta{Conversão: 3 fórmulas e reduções; nodal: sistema $2\times2$. A
conversão dá só $R_{eq}$; a nodal dá tudo}
\resolucao{(a) \textbf{Conversão:} identificar o triângulo (1); aplicar três
fórmulas (2); somar duas séries (3); um paralelo (4); uma soma final (5). Sem
sistemas de equações.\\
\textbf{Nodal:} montar a matriz $2\times2$ (1); resolver (2); calcular a
corrente total a partir dos potenciais (3).\\
(b) \textbf{Critério.}
\begin{itemize}
\item Se a pergunta é \textbf{só} $R_{eq}$ --- ou a corrente total --- a
conversão é mais rápida e não exige álgebra linear.
\item Se são necessários potenciais internos, correntes de ramo ou potências
individuais, use \textbf{nodal}. A conversão destrói a estrutura interna: o nó
central da estrela não existe no circuito real, e nenhuma corrente calculada
nele tem significado físico.
\item Para pontes com \emph{mais} de um triângulo irredutível, a conversão
precisa ser aplicada repetidamente e perde a vantagem; a nodal escala melhor.
\end{itemize}
\textbf{Estratégia mista frequente:} converter para obter $R_{eq}$ e a corrente
total, e depois voltar ao circuito \emph{original} para distribuir essa corrente
--- aproveitando o melhor dos dois.}
\stepcounter{questao}
\resposta{A condição é a desigualdade triangular sobre as resistências entre
pares}
\resolucao{(a) Sejam $R_{ab}$, $R_{bc}$ e $R_{ca}$ as resistências medidas
entre cada par de terminais (com o terceiro aberto). Da estrela, valem
$R_{ab}=R_A+R_B$, $R_{bc}=R_B+R_C$ e $R_{ca}=R_C+R_A$, de onde
\[
R_A=\frac{R_{ab}+R_{ca}-R_{bc}}{2},
\]
e analogamente para $R_B$ e $R_C$.\\
\textbf{Condição:} os três braços são não negativos se e somente se cada soma
de duas resistências medidas for maior ou igual à terceira --- exatamente a
\textbf{desigualdade triangular}.\\
(b) \textbf{Contraexemplo.} Suponha
$R_{ab}=\SI{1}{\ohm}$, $R_{bc}=\SI{10}{\ohm}$ e
$R_{ca}=\SI{1}{\ohm}$. Então
\[
R_B=\frac{1+10-1}{2}=\SI{5}{\ohm};
\qquad
R_A=\frac{1+1-10}{2}=\SI{-4}{\ohm}.
\]
Braço negativo --- fisicamente irrealizável com resistores passivos.\\
\textbf{Mas atenção:} essa combinação de medidas também não é realizável por
uma rede passiva de três terminais. A desigualdade triangular é uma
\emph{propriedade} de tais redes, não uma restrição adicional. Braços negativos
só surgem como artifício de cálculo --- em transformações intermediárias ou em
redes com elementos ativos.}
\stepcounter{questao}
\resposta{$R_A=\dfrac{R_{AB}R_{CA}}{R_{AB}+R_{BC}+R_{CA}}$ e análogas}
\resolucao{(a) A equivalência exige que a resistência entre cada par de
terminais, com o terceiro \emph{aberto}, seja a mesma nas duas configurações:
\[
\begin{cases}
R_A+R_B=R_{AB}\parallel(R_{BC}+R_{CA})\\[3pt]
R_B+R_C=R_{BC}\parallel(R_{CA}+R_{AB})\\[3pt]
R_C+R_A=R_{CA}\parallel(R_{AB}+R_{BC})
\end{cases}
\]
(b) Somando as três e dividindo por $2$, obtém-se $R_A+R_B+R_C$. Subtraindo
dessa soma a segunda equação, resta
\[
R_A=\frac{(R_A+R_B)+(R_C+R_A)-(R_B+R_C)}{2}.
\]
(c) Substituindo os paralelos, com $\Sigma=R_{AB}+R_{BC}+R_{CA}$:
\[
R_A+R_B=\frac{R_{AB}(R_{BC}+R_{CA})}{\Sigma};
\qquad
R_C+R_A=\frac{R_{CA}(R_{AB}+R_{BC})}{\Sigma};
\]
\[
R_B+R_C=\frac{R_{BC}(R_{CA}+R_{AB})}{\Sigma}.
\]
Levando à expressão de $R_A$ e simplificando o numerador:
\[
R_{AB}R_{BC}+R_{AB}R_{CA}+R_{CA}R_{AB}+R_{CA}R_{BC}
-R_{BC}R_{CA}-R_{BC}R_{AB}=2R_{AB}R_{CA},
\]
donde
\[
R_A=\frac{2R_{AB}R_{CA}}{2\Sigma}=\frac{R_{AB}R_{CA}}{\Sigma}.
\]
\textbf{A regra mnemônica agora tem origem:} o produto dos dois lados que tocam
o terminal aparece porque os demais termos se cancelam aos pares na combinação
$(R_{ab}+R_{ca}-R_{bc})/2$.}
\stepcounter{questao}
\resposta{$R_{AB}=\dfrac{R_AR_B+R_BR_C+R_CR_A}{R_C}$}
\resolucao{(a) Tomando as três fórmulas diretas e formando a soma dos produtos
dois a dois:
\[
R_AR_B+R_BR_C+R_CR_A
=\frac{R_{AB}R_{BC}R_{CA}\left(R_{AB}+R_{BC}+R_{CA}\right)}{\Sigma^{2}}
=\frac{R_{AB}R_{BC}R_{CA}}{\Sigma}.
\]
Chamando esse valor de $P$ e observando que
$R_C=\dfrac{R_{BC}R_{CA}}{\Sigma}$, tem-se
\[
\frac{P}{R_C}
=\frac{R_{AB}R_{BC}R_{CA}/\Sigma}{R_{BC}R_{CA}/\Sigma}=R_{AB}.
\]
Logo
\[
R_{AB}=\frac{R_AR_B+R_BR_C+R_CR_A}{R_C}.
\]
(b) A dedução \emph{é} a verificação: partiu-se das fórmulas diretas e
chegou-se de volta aos lados originais. Numericamente, o exemplo
$20/30/\SI{50}{\ohm}\to10/6/\SI{15}{\ohm}\to20/30/\SI{50}{\ohm}$ confirma.\\
(c) \textbf{Por que o oposto.} O braço $R_C$ é o único que \emph{não} contém
$R_{AB}$ em seu numerador --- ele é construído a partir de $R_{BC}$ e
$R_{CA}$. Ao dividir $P$ (que contém o produto dos três lados) por $R_C$,
cancelam-se exatamente $R_{BC}$ e $R_{CA}$, isolando $R_{AB}$. É por isso que
o denominador tem de ser o braço oposto, e não outro.}
\stepcounter{questao}
\resposta{Ambas as fórmulas são quocientes de quantidades positivas ---
positividade garantida}
\resolucao{(a) $R_A=\dfrac{R_{AB}R_{CA}}{R_{AB}+R_{BC}+R_{CA}}$. Com todos os
lados positivos, numerador e denominador são positivos, logo $R_A>0$. O mesmo
para $R_B$ e $R_C$.\\
\textbf{Além disso, há um limite superior:} como
$\Sigma>R_{CA}$, segue $R_A<R_{AB}$; e como $\Sigma>R_{AB}$, segue
$R_A<R_{CA}$. Ou seja, \textbf{cada braço é menor que os dois lados que o
originaram} --- teste de sanidade imediato para qualquer conversão.\\
(b) $R_{AB}=\dfrac{R_AR_B+R_BR_C+R_CR_A}{R_C}$: soma de três produtos positivos
dividida por um positivo, logo $R_{AB}>0$.\\
\textbf{E há um limite inferior:} $R_{AB}>\dfrac{R_AR_C+R_BR_C}{R_C}=R_A+R_B$
--- \textbf{cada lado é maior que a soma dos dois braços correspondentes}.
Segundo teste de sanidade, dual do anterior.\\
(c) \textbf{Casos-limite.}
\begin{itemize}
\item $R_{AB}\to0$ (lado em curto) $\Rightarrow$ $R_A,R_B\to0$: os terminais
$A$ e $B$ colapsam sobre o nó central.
\item $R_{AB}\to\infty$ (lado aberto) $\Rightarrow$ a estrela degenera em uma
série simples.
\item $R_C\to0$ na estrela $\Rightarrow$ $R_{AB}\to\infty$: o lado direto
desaparece.
\item $R_C\to\infty$ $\Rightarrow$ $R_{AB}\to R_A+R_B$, e os outros dois lados
divergem.
\end{itemize}
Em todos os casos, o limite reproduz a interpretação física esperada --- as
fórmulas nunca produzem resultados sem sentido para entradas positivas.}
\stepcounter{questao}
\resposta{Preserva o comportamento nos três terminais; \emph{não} preserva nada
do interior}
\resolucao{(a) Por construção, as três equações de equivalência impõem
igualdade das resistências entre pares. Como a rede de três terminais é linear e
tem apenas três acessos, seu comportamento externo é completamente descrito por
esses três valores --- qualquer excitação aplicada aos terminais produz a mesma
resposta nas duas configurações.\\
\emph{(Formalmente: a matriz de condutâncias de três terminais tem três
parâmetros independentes, e as três equações os fixam todos.)}\\
(b) \textbf{Não.} A estrela possui um nó central inexistente no triângulo, e o
triângulo possui um lado direto inexistente na estrela. Consequências:
\begin{itemize}
\item As correntes dos elementos são \emph{completamente diferentes}.
\item As potências individuais também --- embora a potência \emph{total} seja a
mesma, pois o comportamento externo é idêntico.
\item O potencial do nó central não corresponde a ponto algum do triângulo.
\end{itemize}
\textbf{Exemplo:} na ponte resolvida anteriormente, o resistor central de
$\SI{6}{\ohm}$ conduz $\SI{2}{\ampere}$. Após a conversão, esse resistor
\emph{não existe} --- e nenhum braço da estrela conduz $\SI{2}{\ampere}$.\\
(c) \textbf{Procedimento correto:} use a conversão para obter as grandezas
\emph{externas} (resistência equivalente, corrente total, tensões nos
terminais). Depois \textbf{volte ao circuito original} e, com essas grandezas já
conhecidas, distribua correntes e tensões pelos elementos reais.\\
Se muitas grandezas internas forem necessárias, a análise nodal costuma ser
mais direta desde o início.}
\stepcounter{questao}
\resposta{Exato: $36{,}67$, $110$ e $\SI{55}{\ohm}$; com $36/110/56$, o erro em
$R_{AB}$ é de cerca de $1\%$}
\resolucao{(a) $\Sigma=10(20)+20(30)+30(10)=200+600+300=\SI{1100}{\ohm\squared}$:
\[
R_{AB}=\frac{1100}{30}=\SI{36.67}{\ohm};
\quad
R_{BC}=\frac{1100}{10}=\SI{110}{\ohm};
\quad
R_{CA}=\frac{1100}{20}=\SI{55}{\ohm}.
\]
(b) E24: $\SI{36}{\ohm}$, $\SI{110}{\ohm}$, $\SI{56}{\ohm}$.\\
\textbf{Grandeza que interessa} --- resistência entre $A$ e $B$ com $C$ aberto:
\[
\text{alvo (estrela)}: R_A+R_B=\SI{30}{\ohm};
\]
\[
\text{obtido}: 36\parallel(110+56)=\frac{36(166)}{202}=\SI{29.58}{\ohm}.
\]
Erro de $-1{,}4\%$ --- aceitável para a maioria das aplicações, e menor que a
tolerância de $5\%$ dos próprios resistores.\\
(c) \textbf{Estratégias, em ordem de custo.}
\begin{itemize}
\item \textbf{Testar combinações:} $39/110/56$ daria
$39\parallel166=\SI{31.6}{\ohm}$ ($+5{,}3\%$) --- pior. O par
$36/110/56$ já é bom.
\item \textbf{Série E96} ($1\%$): contém $36{,}5$ e $54{,}9$, reduzindo o erro
a frações de por cento.
\item \textbf{Associação:} $\SI{36.67}{\ohm}$ sai de $\SI{33}{\ohm}$ em série
com $\SI{3.6}{\ohm}$, ou de $\SI{110}{\ohm}$ em paralelo com
$\SI{55}{\ohm}$ --- e note que esses dois valores já estão na lista!
\end{itemize}
\textbf{Critério de avaliação:} sempre meça o erro na \emph{grandeza de
interesse}, não nos componentes. Erros de lados individuais podem se compensar
ou se somar, e só o cálculo da resistência terminal revela o efeito real.}
\stepcounter{questao}
\resposta{$I_g\approx\SI{-124}{\micro\ampere}$ (de $B$ para $A$); a conversão é
inadequada para esta grandeza}
\resolucao{(a) Com $V_S=\SI{10}{\volt}$ e o terra na base, o sistema nodal
\[
\begin{cases}
V_A\left(\tfrac{1}{100}+\tfrac{1}{100}+\tfrac{1}{100}\right)
-\tfrac{V_B}{100}=\tfrac{10}{100}\\[4pt]
-\tfrac{V_A}{100}
+V_B\left(\tfrac{1}{100}+\tfrac{1}{101}+\tfrac{1}{100}\right)=\tfrac{10}{100}
\end{cases}
\]
dá $V_A=\SI{5.00621}{\volt}$ e $V_B=\SI{5.01863}{\volt}$, de onde
\[
I_g=\frac{V_A-V_B}{100}=\frac{-0{,}01242}{100}=\SI{-124}{\micro\ampere}.
\]
O sinal indica corrente de $B$ para $A$ --- coerente, pois $R_4>R_3$ eleva
$V_B$.\\
(b) \textbf{A conversão é inadequada, por duas razões independentes.}\\
\emph{(i) Ela não fornece a grandeza pedida.} Após converter o triângulo, o
resistor do galvanômetro deixa de existir como elemento --- e é justamente sua
corrente que se quer. Seria necessário converter, obter $R_{eq}$ e a corrente
total, voltar ao circuito original e só então distribuir --- caminho mais longo
que a nodal direta.\\
\emph{(ii) O problema é mal condicionado.} A resposta é a diferença
$V_A-V_B=\SI{-12.4}{\milli\volt}$ entre dois valores próximos de
$\SI{5}{\volt}$ --- uma diferença relativa de apenas $0{,}25\%$.
\begin{center}
\begin{tabular}{ccc}
\hline
Algarismos usados & $V_A-V_B$ obtido & erro em $I_g$\\
\hline
$6$ ($5{,}00621$; $5{,}01863$) & $\SI{-12.42}{\milli\volt}$ & $<0{,}1\%$\\
$4$ ($5{,}006$; $5{,}019$) & $\SI{-13}{\milli\volt}$ & $\approx5\%$\\
$3$ ($5{,}01$; $5{,}02$) & $\SI{-10}{\milli\volt}$ & $\approx20\%$\\
\hline
\end{tabular}
\end{center}
Arredondar cedo destrói o resultado.\\
\textbf{Método adequado:} tratar o desequilíbrio como \emph{perturbação}. Para
$\delta=\Delta R/R$ pequeno, $I_g$ é aproximadamente proporcional a $\delta$, e
a expressão linearizada evita completamente a subtração de valores próximos.\\
\textbf{Moral:} escolher o método é também escolher \emph{como} o erro numérico
se propaga. A conversão estrela-triângulo é excelente para $R_{eq}$ e péssima
para desequilíbrios pequenos.}
\stepcounter{questao}
\resposta{$R_A=\SI{15}{\ohm}$, $R_B=\SI{5}{\ohm}$, $R_C=\SI{25}{\ohm}$;
triângulo $23$, $38{,}33$ e $\SI{115}{\ohm}$; a estrela é preferível}
\resolucao{(a) \textbf{Realizabilidade.} A desigualdade triangular exige que
cada resistência medida seja menor que a soma das outras duas:
$20<70$, $30<60$, $40<50$ \checkmark. A rede é realizável com resistores
positivos.\\
Invertendo o sistema $R_A+R_B=R_{AB}$ e análogos:
\[
R_A=\frac{20+40-30}{2}=\SI{15}{\ohm};
\quad
R_B=\frac{20+30-40}{2}=\SI{5}{\ohm};
\quad
R_C=\frac{30+40-20}{2}=\SI{25}{\ohm}.
\]
\textbf{Conferência:} $15+5=20$; $5+25=30$; $25+15=40$ \checkmark.\\
(b) $\Sigma=15(5)+5(25)+25(15)=75+125+375=\SI{575}{\ohm\squared}$:
\[
R_{AB}=\frac{575}{25}=\SI{23}{\ohm};
\quad
R_{BC}=\frac{575}{15}=\SI{38.33}{\ohm};
\quad
R_{CA}=\frac{575}{5}=\SI{115}{\ohm}.
\]
(c) \textbf{Estrela com E24:} $\SI{15}{\ohm}$ e $\SI{24}{\ohm}$ existem, e o
mais próximo de $\SI{5}{\ohm}$ é $\SI{5.1}{\ohm}$. Com
$(15;\,5{,}1;\,24)$:
\[
R_{AB}=\SI{20.1}{\ohm}\ (+0{,}5\%);
\quad
R_{BC}=\SI{29.1}{\ohm}\ (-3{,}0\%);
\quad
R_{CA}=\SI{39.0}{\ohm}\ (-2{,}5\%).
\]
\textbf{Triângulo com E24:} $(24;\,39;\,110)$:
\[
R_{AB}=\SI{20.67}{\ohm}\ (+3{,}4\%);
\quad
R_{BC}=\SI{30.21}{\ohm}\ (+0{,}7\%);
\quad
R_{CA}=\SI{40.06}{\ohm}\ (+0{,}1\%).
\]
\textbf{Conclusão.} Os dois erram no mesmo patamar (pior caso de $3{,}0\%$
contra $3{,}4\%$), mas a \textbf{estrela é preferível} por três razões: seus
valores são menores (componentes mais baratos), o erro é dominado por um único
braço difícil ($\SI{5}{\ohm}$), e cada resistência medida depende de
\emph{apenas dois} braços --- o que torna o ajuste previsível.\\
No triângulo, cada medida envolve os \emph{três} lados através de um paralelo, e
os erros se recombinam de forma difícil de controlar.\\
\textbf{Se a precisão for crítica:} o braço de $\SI{5}{\ohm}$ pode ser obtido
associando $\SI{10}{\ohm}$ em paralelo com $\SI{10}{\ohm}$, eliminando o
único desvio significativo.}
\stepcounter{questao}
\resposta{Estrela: $n$ braços; polígono: $n(n-1)/2$ lados. A inversa só existe
para $n=3$}
\resolucao{(a) Uma estrela de $n$ braços $R_1,\dots,R_n$ ligados a um nó comum
equivale a um polígono completo em que o lado entre os terminais $i$ e $j$ vale
\[
R_{ij}=R_iR_j\sum_{k=1}^{n}\frac{1}{R_k}.
\]
Para $n=3$ isso reproduz a fórmula conhecida:
$R_{12}=R_1R_2\left(\frac1{R_1}+\frac1{R_2}+\frac1{R_3}\right)
=\frac{R_1R_2+R_2R_3+R_3R_1}{R_3}$ \checkmark.\\
(b) \textbf{Estrela:} $n$ elementos. \textbf{Polígono completo:}
$\dbinom{n}{2}=\dfrac{n(n-1)}{2}$ elementos.
\begin{center}
\begin{tabular}{ccc}
\hline
$n$ & estrela & polígono\\
\hline
$3$ & $3$ & $3$\\
$4$ & $4$ & $6$\\
$5$ & $5$ & $10$\\
\hline
\end{tabular}
\end{center}
(c) \textbf{Contagem de graus de liberdade.} A transformação
$Y\to$ polígono sempre existe. A inversa exigiria determinar $n$ braços a
partir de $n(n-1)/2$ lados --- um sistema \textbf{sobredeterminado} para
$n>3$.\\
Só quando $n=3$ os dois números coincidem ($3=3$), e a correspondência é
biunívoca. Para $n=4$, seria preciso obter $4$ incógnitas de $6$ equações: um
polígono \emph{genérico} de quatro terminais simplesmente não admite estrela
equivalente.\\
\textbf{Consequência prática:} a conversão $\Delta\to Y$ é uma ferramenta
peculiar ao caso de três terminais. Redes de quatro ou mais acessos exigem
descrição matricial completa --- e é por isso que a análise nodal, que não
depende de topologias especiais, é o método geral.}

\gabarito[3.4 Transformação estrela-triângulo]

\section{Divisor de Tensão e de Corrente}

\subsection{Divisor de tensão}
\stepcounter{questao}
\resposta{(b)}
\resolucao{Como a corrente é a mesma nos dois ($U = R\,I$), a tensão é
proporcional à resistência: o maior resistor fica com a maior tensão. A posição
em relação ao terminal positivo (alternativa d) é irrelevante.}
\stepcounter{questao}
\resposta{$U_1=\SI{30}{\volt}$, $U_2=\SI{40}{\volt}$, $U_3=\SI{50}{\volt}$}
\resolucao{A corrente é $I = \dfrac{120}{3+4+5} = \dfrac{120}{12} = \SI{10}{\ampere}$.
Aplicando o divisor (ou simplesmente $U_k = R_k I$):
\[
U_1 = 3\cdot10 = \SI{30}{\volt},\quad
U_2 = 4\cdot10 = \SI{40}{\volt},\quad
U_3 = 5\cdot10 = \SI{50}{\volt}.
\]
\textbf{Verificação:} $30+40+50 = \SI{120}{\volt}$, igual à fonte.}
\stepcounter{questao}
\resposta{(b)}
\resolucao{Aplicando diretamente a fórmula do divisor:
\[
U_2 = E\,\frac{R_2}{R_1+R_2} = 100\cdot\frac{40}{100} = \SI{40}{\volt}.
\]}
\stepcounter{questao}
\resposta{(c)}
\resolucao{Resistores iguais dividem a tensão igualmente:
$U = \dfrac{60}{3} = \SI{20}{\volt}$ em cada um. É o divisor de tensão mais
simples --- a base dos potenciômetros de precisão.}
\stepcounter{questao}
\resposta{$R_2 = \SI{1}{\kilo\ohm}$}
\resolucao{Da relação do divisor,
\[
\frac{U_{\text{out}}}{E}=\frac{R_2}{R_1+R_2}
\ \Rightarrow\
\frac{5}{15}=\frac{R_2}{2000+R_2}
\ \Rightarrow\
2000+R_2 = 3R_2
\ \Rightarrow\
R_2 = \SI{1}{\kilo\ohm}.
\]
Repare que só importa a \emph{razão} $R_1{:}R_2 = 2{:}1$; qualquer par nessa
proporção daria \SI{5}{\volt} em vazio.}
\stepcounter{questao}
\resposta{(a) \SI{6}{\volt}; (b) \SI{4}{\volt}}
\resolucao{(a) \textbf{Em vazio}, nenhuma corrente sai pelos terminais, e os dois
resistores iguais dividem a tensão ao meio:
$U_{\text{out}} = 12\cdot\dfrac{1}{2} = \SI{6}{\volt}$.\\
(b) \textbf{Com carga}, $R_L$ fica em paralelo com $R_2$:
\[
R_2 \parallel R_L = 1\text{k}\parallel 1\text{k} = \SI{0.5}{\kilo\ohm}
\ \Rightarrow\
U_{\text{out}} = 12\cdot\frac{0{,}5}{1+0{,}5} = \SI{4}{\volt}.
\]
\textbf{Conclusão:} a conexão da carga reduz a tensão de saída. Divisores resistivos só
mantêm a saída se $R_L \gg R_2$; por isso não servem como fonte de alimentação
para cargas variáveis --- para isso usa-se um regulador de tensão.}
\stepcounter{questao}
\resposta{\SI{2.7}{\volt}}
\resolucao{O cursor divide a resistência total na proporção do curso. A \SI{30}{\percent},
a saída corresponde a $\SI{30}{\percent}$ da tensão total:
\[
U_{\text{out}} = 9\cdot 0{,}30 = \SI{2.7}{\volt}.
\]
O valor da resistência total (\SI{10}{\kilo\ohm}) não entra na conta em vazio ---
apenas a \emph{fração} do curso importa. É exatamente assim que um controle de
volume ajusta o sinal.}
\stepcounter{questao}
\resposta{$I = \SI{2.38}{\ampere}$ (aprox.); $U_{8} \approx \SI{19}{\volt}$}
\resolucao{Todos os resistores estão em série:
$\Req = 10+6+8+12+6 = \SI{42}{\ohm}$.
\[
I=\frac{100}{42}\approx\SI{2.38}{\ampere},
\qquad
U_8 = 8\cdot I = \frac{800}{42}\approx\SI{19.0}{\volt}.
\]
Pelo divisor: $U_8 = 100\cdot\dfrac{8}{42}\approx\SI{19.0}{\volt}$, o mesmo
resultado.}
\stepcounter{questao}
\resposta{$R_2 = \SI{500}{\ohm}$; $R_1 = \SI{500}{\ohm}$}
\resolucao{\textbf{Ramo de $R_2$ (sangria):} $R_2$ está em paralelo com a carga e
sob \SI{5}{\volt}. Com corrente de sangria $I_2 = \SI{10}{\milli\ampere}$:
\[
R_2 = \frac{U_2}{I_2}=\frac{5}{\pot{10}{-3}}=\SI{500}{\ohm}.
\]
\textbf{Ramo de $R_1$:} por ele passa a soma da corrente de sangria com a da carga,
$I_1 = 10+10 = \SI{20}{\milli\ampere}$, sob a tensão restante
$15-5 = \SI{10}{\volt}$:
\[
R_1=\frac{U_1}{I_1}=\frac{10}{\pot{20}{-3}}=\SI{500}{\ohm}.
\]
\textbf{Por que a sangria?} Se $R_2$ conduzisse muito pouco, a tensão de saída
variaria fortemente com a carga. Uma sangria comparável à corrente da carga
estabiliza a saída, ao custo de mais potência dissipada --- o compromisso típico
de projeto.}
\stepcounter{questao}
\resposta{$R_L \gtrsim \SI{12}{\kilo\ohm}$}
\resolucao{Em vazio, $U_0 = 12\cdot\dfrac{2}{6} = \SI{4}{\volt}$. Queremos
$U \geq 0{,}9\cdot U_0 = \SI{3.6}{\volt}$. Com carga, $R_2$ é substituído por
$R_2' = R_2\parallel R_L$:
\[
U = 12\cdot\frac{R_2'}{R_1+R_2'}\geq 3{,}6.
\]
Isolando, $\dfrac{R_2'}{4000+R_2'}\geq 0{,}3 \Rightarrow R_2'\geq \SI{1714}{\ohm}$
(aprox.). De $R_2' = \dfrac{2000\,R_L}{2000+R_L}\geq 1714$ resulta
$R_L \geq \SI{12}{\kilo\ohm}$ (aproximadamente). Ou seja, a carga precisa ser cerca
de \textbf{6 vezes maior} que $R_2$ para "não pesar" no divisor. Confirma a regra
prática: divisores resistivos só servem para cargas de alta impedância, como a
entrada de um amplificador.}
\stepcounter{questao}
\resposta{Ver análise; $R_2 \ll R_L$, com alto consumo}
\resolucao{(a) A tensão de saída de um divisor \emph{depende da carga}, pois
$R_L$ entra em paralelo com $R_2$. Quando $R_L$ varia, o paralelo
$R_2 \parallel R_L$ varia, e a saída varia junto. A regulação só é boa quando
$R_2 \parallel R_L \approx R_2$, isto é, quando $R_2 \ll R_L$.\\
(b) Quantificando: para que a variação de $R_L$ (de \SI{100}{} para
\SI{200}{\ohm}) altere o paralelo em menos de \SI{2}{\percent}, precisamos de
\[
R_2 \parallel 100 \approx R_2 \parallel 200
\quad\Longrightarrow\quad
R_2 \lesssim \SI{4}{\ohm} \ (\text{ordem de } R_L/25).
\]
Com $R_2 = \SI{4}{\ohm}$ e a razão do divisor exigindo
$R_1 = 2R_2 = \SI{8}{\ohm}$, a corrente do divisor seria
$I \approx \dfrac{15}{12} = \SI{1.25}{\ampere}$, dissipando cerca de
$\SI{19}{\watt}$ --- para entregar apenas $\dfrac{5^2}{100} = \SI{0.25}{\watt}$
à carga!\\
\textbf{Conclusão de projeto:} o rendimento seria de aproximadamente
\SI{1}{\percent}. Isso demonstra por que divisores resistivos servem apenas como
\emph{referência de tensão} para entradas de alta impedância, e por que fontes
reais usam \textbf{reguladores} (série ou chaveados), que mantêm a saída
independentemente da carga sem desperdiçar potência.}
\stepcounter{questao}
\resposta{$U=\SI{4}{\volt}$}
\resolucao{A regra do divisor vale para qualquer número de elementos em série:
cada um fica com a fração de sua própria resistência sobre o total.
\[
U_2=12\cdot\frac{2}{1+2+3}=12\cdot\frac13=\SI{4}{\volt}.
\]
\textbf{Conferência:} as três tensões são $2$, $4$ e $\SI{6}{\volt}$, somando
$\SI{12}{\volt}$ \checkmark, e guardam a proporção $1{:}2{:}3$ dos
resistores.}
\stepcounter{questao}
\resposta{$R_1/R_2=3$}
\resolucao{
\[
\frac{R_2}{R_1+R_2}=0{,}25
\;\Longrightarrow\;
R_1+R_2=4R_2
\;\Longrightarrow\;
R_1=3R_2.
\]
\textbf{Observação importante:} apenas a \emph{razão} é fixada --- $3\,$k$/1\,$k
e $30\,$k$/10\,$k produzem a mesma saída em vazio. O que os distingue é a
corrente de repouso e a resistência interna do divisor, e é isso que decide o
projeto quando há carga.}
\stepcounter{questao}
\resposta{(c)}
\resolucao{A razão $\dfrac{2R_2}{2R_1+2R_2}=\dfrac{R_2}{R_1+R_2}$ é
inalterada --- a saída em vazio depende só da proporção.\\
\textbf{O que muda:} a corrente de repouso cai à metade (menos consumo) e a
resistência interna $R_1\parallel R_2$ \emph{dobra} (mais sensível à carga).
Escalar um divisor troca consumo por robustez, e essa é sempre a decisão
central do projeto.}
\stepcounter{questao}
\resposta{(a) $U_o\to0$; (b) $U_o\to U_{in}$}
\resolucao{(a) Com $R_2=0$, a saída está em \textbf{curto} com o terra:
$U_o=0$. A corrente vale $U_{in}/R_1$ --- toda a potência vai para $R_1$.\\
(b) Com $R_2\to\infty$ (circuito aberto), não circula corrente, não há queda em
$R_1$ e $U_o\to U_{in}$.\\
\textbf{Conclusão:} a saída de um divisor está sempre entre $0$ e $U_{in}$ --- ele
só \emph{reduz} tensão, nunca a eleva. Obter tensão maior que a entrada exige
elemento armazenador (indutor ou capacitor) e comutação.}
\stepcounter{questao}
\resposta{(a) $\SI{9}{\volt}$ e $\SI{6}{\volt}$; (b) $\SI{3}{\milli\ampere}$}
\resolucao{Total: $R=\SI{4}{\kilo\ohm}$ e
$I=12/4000=\SI{3}{\milli\ampere}$.\\
A tomada mais alta tem, abaixo de si, $1+2=\SI{3}{\kilo\ohm}$:
\[
U_A=12\cdot\frac{3}{4}=\SI{9}{\volt}.
\]
A tomada inferior tem $\SI{2}{\kilo\ohm}$ abaixo:
\[
U_B=12\cdot\frac{2}{4}=\SI{6}{\volt}.
\]
\textbf{Cuidado:} cada saída conta apenas a resistência \emph{entre ela e o
terra}, não a do trecho acima. E as duas tomadas se \emph{influenciam}: carregar
uma altera a outra, pois compartilham a mesma cadeia.}
\stepcounter{questao}
\resposta{$\SI{+15}{\volt}$ e $\SI{-15}{\volt}$}
\resolucao{Cada resistor cai $\SI{15}{\volt}$. Tomando o ponto médio como
$\SI{0}{\volt}$, o terminal superior fica em $\SI{+15}{\volt}$ e o inferior em
$\SI{-15}{\volt}$.\\
\textbf{Aplicação:} é assim que se obtém alimentação \emph{simétrica} para
amplificadores operacionais a partir de uma fonte simples. \\
\textbf{Limitação séria:} o "terra" assim criado é um ponto de divisor, com
resistência interna $R/2$. Qualquer corrente drenada de forma assimétrica pelos
dois ramos desloca a referência. Na prática, o divisor precisa ser reforçado por
um seguidor de tensão ou por capacitores de grande valor.}
\stepcounter{questao}
\resposta{$\SI{8}{\volt}$; $\SI{2}{\milli\ampere}$; $\SI{48}{\milli\watt}$}
\resolucao{
\[
U_o=24\cdot\frac{4}{12}=\SI{8}{\volt};
\qquad
I=\frac{24}{12000}=\SI{2}{\milli\ampere};
\]
\[
P=U_{in}I=24(\pot{2}{-3})=\SI{48}{\milli\watt}.
\]
\textbf{Distribuição:} $P_1=8000(\pot{2}{-3})^{2}=\SI{32}{\milli\watt}$ e
$P_2=\SI{16}{\milli\watt}$ --- proporcionais às resistências, como em qualquer
série.\\
Note que essa potência é consumida \emph{permanentemente}, mesmo sem carga
alguma conectada. Em equipamentos alimentados por bateria, divisores de repouso
são uma fonte silenciosa de descarga.}
\stepcounter{questao}
\resposta{de $\SI{0}{\volt}$ a $\SI{8}{\volt}$}
\resolucao{
\[
R_2=0 \Rightarrow U_o=0;
\qquad
R_2=\SI{8}{\kilo\ohm} \Rightarrow U_o=10\cdot\frac{8}{10}=\SI{8}{\volt}.
\]
\textbf{Não linearidade do ajuste:} no meio do curso
($R_2=\SI{4}{\kilo\ohm}$), a saída é $10\cdot4/6=\SI{6.67}{\volt}$ --- e não
$\SI{4}{\volt}$, como uma escala linear sugeriria. A relação
$U_o(R_2)$ é uma hipérbole, não uma reta.\\
Para obter ajuste \emph{linear}, é preciso usar um potenciômetro (em que a soma
$R_1+R_2$ permanece constante), e não um reostato em série com resistor fixo.}
\stepcounter{questao}
\resposta{$U_{in}=\SI{12}{\volt}$}
\resolucao{
\[
4{,}5=U_{in}\cdot\frac{3}{5+3}
\;\Longrightarrow\;
U_{in}=4{,}5\cdot\frac{8}{3}=\SI{12}{\volt}.
\]
\textbf{Verificação:} $12\cdot3/8=\SI{4.5}{\volt}$ \checkmark, e a corrente é
$12/8000=\SI{1.5}{\milli\ampere}$.\\
A relação do divisor pode ser resolvida em qualquer das quatro variáveis ---
$U_{in}$, $U_o$, $R_1$ ou $R_2$ --- desde que as outras três sejam conhecidas.}
\stepcounter{questao}
\resposta{$\SI{2.5}{\volt}$ e $\SI{1.43}{\volt}$}
\resolucao{
\[
\SI{25}{\celsius}:\ U_o=5\cdot\frac{10}{20}=\SI{2.5}{\volt};
\qquad
\SI{50}{\celsius}:\ U_o=5\cdot\frac{4}{14}=\SI{1.43}{\volt}.
\]
Variação de $\SI{1.07}{\volt}$ em $\SI{25}{\celsius}$, ou cerca de
$\SI{43}{\milli\volt\per\celsius}$ na média --- sinal amplo e fácil de digitalizar.\\
\textbf{Ressalva:} a relação \emph{não} é linear, pois nem o termistor é linear
em $T$, nem o divisor é linear em $R_2$. Para converter tensão em temperatura é
preciso tabela, polinômio ou a equação de Steinhart--Hart. A escolha de
$R_1=\SI{10}{\kilo\ohm}$ não é casual --- ela maximiza a sensibilidade em torno
de $\SI{25}{\celsius}$, como se demonstra no bloco de desafio.}
\stepcounter{questao}
\resposta{de $\SI{4.55}{\volt}$ (claro) a $\SI{0.45}{\volt}$ (escuro)}
\resolucao{
\[
\text{claro}:\ U_o=5\cdot\frac{10}{11}=\SI{4.55}{\volt};
\qquad
\text{escuro}:\ U_o=5\cdot\frac{10}{110}=\SI{0.45}{\volt}.
\]
Com o LDR na posição \emph{superior}, mais luz significa \emph{maior} tensão de
saída.\\
\textbf{Inversão do comportamento:} trocando as posições (LDR embaixo), a saída
passaria de $\SI{0.45}{\volt}$ no claro a $\SI{4.55}{\volt}$ no escuro. A
escolha da posição define a polaridade do sensor --- decisão de projeto que
custa nada e evita inverter o sinal depois em software.}
\stepcounter{questao}
\resposta{$U_o=\SI{3.19}{\volt}$; erro de $-3{,}2\%$}
\resolucao{
\[
U_o=12\cdot\frac{3{,}3}{9{,}1+3{,}3}=12\cdot\frac{3{,}3}{12{,}4}=\SI{3.19}{\volt};
\qquad
\text{erro}=\frac{3{,}19-3{,}30}{3{,}30}=-3{,}2\%.
\]
\textbf{Alternativa melhor:} o par $\SI{10}{\kilo\ohm}/\SI{3.9}{\kilo\ohm}$ dá
$U_o=\SI{3.37}{\volt}$, erro de apenas $+2{,}0\%$. Vale sempre testar
\emph{vários} pares E24 antes de escolher --- a razão importa mais que os
valores individuais.\\
\textbf{Se $3\%$ for demais:} use a série E96 (tolerância $1\%$), que contém
$\SI{8.66}{\kilo\ohm}$ e resolve o problema diretamente.}
\stepcounter{questao}
\resposta{(a) $\SI{3.6}{\volt}$; (b) previsão de $\SI{4}{\volt}$, erro de
$10\%$; (c) $R_{Th}=\SI{2.1}{\kilo\ohm}$}
\resolucao{(a) A resistência interna soma-se a $R_1$:
\[
U_o=12\cdot\frac{3}{1+6+3}=12\cdot\frac{3}{10}=\SI{3.6}{\volt}.
\]
(b) Ignorando $r$: $12\cdot3/9=\SI{4}{\volt}$. O erro é de $10\%$ --- nada
desprezível.\\
(c) Zerando a f.e.m. (mas mantendo $r$), a saída enxerga $R_2$ em paralelo com
$(r+R_1)$:
\[
R_{Th}=\frac{3\cdot7}{10}=\SI{2.1}{\kilo\ohm}.
\]
\textbf{Leitura:} a resistência interna da fonte se comporta como parte de
$R_1$ --- ela piora a razão \emph{e} aumenta a impedância de saída. Em divisores
de baixa impedância (resistores de dezenas de ohms), $r$ pode dominar
completamente o projeto.}
\stepcounter{questao}
\resposta{(a) $\SI{2}{\volt}$ e $\SI{2.146}{\volt}$; diferença de
$\SI{146}{\milli\volt}$}
\resolucao{(a)
\[
U_A=8\cdot\frac{1}{4}=\SI{2}{\volt};
\qquad
U_B=8\cdot\frac{1{,}1}{4{,}1}=\SI{2.146}{\volt};
\qquad
\Delta U=\SI{146}{\milli\volt}.
\]
(b) Derivando $U=U_{in}R_2/(R_1+R_2)$:
\[
\frac{\mathrm{d}U}{\mathrm{d}R_2}
=\frac{U_{in}R_1}{(R_1+R_2)^{2}}
=\frac{8(3000)}{(4000)^{2}}=\SI{1.5}{\milli\volt\per\ohm}.
\]
Para $\Delta R_2=\SI{100}{\ohm}$, prevê-se $\SI{150}{\milli\volt}$ --- muito
próximo dos $\SI{146}{\milli\volt}$ exatos, pois a variação é pequena.\\
\textbf{Esta é a estrutura da ponte:} dois divisores comparados, em que o
\emph{desequilíbrio} carrega a informação. A grande vantagem é que variações
comuns aos dois ramos --- deriva da fonte, temperatura ambiente --- se cancelam
na diferença.}
\stepcounter{questao}
\resposta{(a) $\SI{4.76}{\volt}$; (b) $-4{,}8\%$; (c)
$\ge\SI{5}{\mega\ohm}$}
\resolucao{(a) O voltímetro entra em paralelo com $R_2$:
\[
R_2'=\frac{100\cdot1000}{1100}=\SI{90.9}{\kilo\ohm};
\qquad
U=10\cdot\frac{90{,}9}{100+90{,}9}=\SI{4.76}{\volt}.
\]
(b) Erro: $(4{,}76-5{,}00)/5{,}00=-4{,}8\%$.\\
(c) A resistência interna do divisor é
$R_{Th}=100\parallel100=\SI{50}{\kilo\ohm}$. O erro relativo é
aproximadamente $R_{Th}/R_V$, então
\[
\frac{50\,\mathrm{k}}{R_V}\le0{,}01
\;\Longrightarrow\;
R_V\ge\SI{5}{\mega\ohm}.
\]
\textbf{Lição de instrumentação:} o instrumento \emph{perturba} o circuito que
mede. Multímetros digitais comuns têm $\SI{10}{\mega\ohm}$ --- suficientes aqui,
mas insuficientes para divisores de megaohms. E note que a especificação
depende de $R_{Th}$ do circuito, não da tensão medida.}
\stepcounter{questao}
\resposta{(a) $\SI{2}{\volt}$, $80\%$ do fundo; (b) $\SI{2.44}{\milli\volt}$;
(c) resistores dominam ($\approx\SI{16}{\milli\volt}$)}
\resolucao{(a) $U_o=10\cdot1/5=\SI{2}{\volt}$, ou $80\%$ do fundo de escala ---
boa utilização da faixa, sem risco de saturação.\\
(b) Com $10$ bits, $2^{10}=1024$ degraus:
\[
\Delta=\frac{2{,}5}{1024}=\SI{2.44}{\milli\volt}.
\]
(c) Com resistores de $1\%$, o pior erro de razão vale
$\dfrac{R_1}{R_1+R_2}(t_1+t_2)=0{,}8(2\%)=1{,}6\%$, ou
$0{,}016\times2=\SI{32}{\milli\volt}$ --- \textbf{treze vezes} o degrau de
quantização.\\
\textbf{Conclusão de projeto:} não adianta usar um conversor de $12$ ou
$16$ bits com esse divisor. O gargalo é o divisor, e a solução é resistores de
$0{,}1\%$, calibração por software, ou uma referência de tensão dedicada.
Aumentar a resolução do conversor sem corrigir o divisor é gastar dinheiro em
dígitos sem significado.}
\stepcounter{questao}
\resposta{(a) $R_1=\SI{300}{\ohm}$, $R_2=\SI{100}{\ohm}$;
(b) $\SI{6.86}{\volt}$; (c) inviável}
\resolucao{(a) $R_2=12/0{,}120=\SI{100}{\ohm}$;
$R_1=(48-12)/0{,}120=\SI{300}{\ohm}$.\\
(b) A carga de $\SI{100}{\ohm}$ fica em paralelo com $R_2$:
\[
R_2'=\frac{100\cdot100}{200}=\SI{50}{\ohm};
\qquad
U_o=48\cdot\frac{50}{350}=\SI{6.86}{\volt}.
\]
Queda de $43\%$ em relação ao projetado.\\
(c) \textbf{Inviável.} A carga consome $\SI{120}{\milli\ampere}$ a
$\SI{12}{\volt}$ --- exatamente a corrente de repouso, e não uma fração dela. A
regra de bolso exige repouso de $10$ a $20$ vezes a corrente da carga, o que
levaria $R_2$ a $\SI{10}{\ohm}$ e a dissipação total a
$48\times1{,}2=\SI{57.6}{\watt}$ para entregar $\SI{1.44}{\watt}$ --- rendimento
de $2{,}5\%$.\\
\textbf{Diagnóstico geral:} divisores resistivos servem para gerar
\emph{referências} de alta impedância, não para \emph{alimentar} cargas. Aqui o
caso pede regulador.}
\stepcounter{questao}
\resposta{$\SI{1}{\watt}$ e $\SI{100}{\volt}$ em cada; exige resistor de
potência com tensão de trabalho adequada}
\resolucao{(a) $I=200/20000=\SI{10}{\milli\ampere}$;
\[
U_1=U_2=\SI{100}{\volt};
\qquad
P_1=P_2=10^{4}(10^{-2})^{2}=\SI{1}{\watt}.
\]
(b) \textbf{Dois critérios independentes:}
\begin{itemize}
\item \textbf{Potência:} com fator $2$ de folga, resistores de
$\SI{2}{\watt}$.
\item \textbf{Tensão máxima de trabalho:} um resistor de filme comum de
$\frac14\,\si{\watt}$ suporta cerca de $\SI{250}{\volt}$; os de
$\SI{2}{\watt}$ chegam a $\SI{500}{\volt}$. Os $\SI{100}{\volt}$ exigidos são
atendidos com folga.
\end{itemize}
\textbf{Por que o segundo critério existe:} a tensão máxima é limitada pela
ruptura do revestimento e pelo arco entre as espiras do filme --- fenômeno que
independe da potência. Em divisores de \emph{alta} tensão (quilovolts), é comum
que a tensão, e não a potência, force o uso de vários resistores em série,
mesmo com dissipação irrisória.}
\stepcounter{questao}
\resposta{(a) $R_1=\SI{1}{\mega\ohm}$, $R_2=\SI{500}{\kilo\ohm}$;
(b) erro de $\approx3\%$}
\resolucao{(a) $R_2=5/\pot{10}{-6}=\SI{500}{\kilo\ohm}$;
$R_1=10/\pot{10}{-6}=\SI{1}{\mega\ohm}$.\\
(b) A corrente de entrada de $\SI{1}{\micro\ampere}$ é $10\%$ da corrente de
repouso --- longe de desprezível. Modelando a entrada como resistência
equivalente $R_{in}=5/\pot{1}{-6}=\SI{5}{\mega\ohm}$:
\[
R_2'=\frac{500\cdot5000}{5500}=\SI{454.5}{\kilo\ohm};
\qquad
U_o=15\cdot\frac{454{,}5}{1454{,}5}=\SI{4.69}{\volt},
\]
erro de $-6{,}2\%$.\\
(c) \textbf{Compromisso.} Reduzir o erro exige aumentar a corrente de repouso ---
com $\SI{100}{\micro\ampere}$ ($R_1=\SI{100}{\kilo\ohm}$,
$R_2=\SI{50}{\kilo\ohm}$), o erro cai para cerca de $0{,}7\%$, mas o consumo em
repouso sobe dez vezes.\\
Em um dispositivo alimentado por bateria de $\SI{1000}{\milli\ampere\hour}$,
$\SI{100}{\micro\ampere}$ contínuos consomem a bateria em pouco mais de um ano
\emph{só com o divisor}. É o tipo de detalhe que decide a autonomia de um
produto.}
\stepcounter{questao}
\resposta{$R_{Th}=R_1\parallel R_2=\SI{2}{\kilo\ohm}$}
\resolucao{(a) Zerando a fonte de entrada (curto), a saída enxerga $R_1$ e
$R_2$ em paralelo:
\[
R_{Th}=\frac{R_1R_2}{R_1+R_2}=\frac{6\cdot3}{9}=\SI{2}{\kilo\ohm}.
\]
(b) O paralelo de dois resistores positivos é sempre menor que o menor deles ---
resultado geral já estabelecido para associações. Aqui,
$\SI{2}{\kilo\ohm}<\SI{3}{\kilo\ohm}$ \checkmark.\\
(c) \textbf{Papel no carregamento.} Conectada uma carga $R_L$, o divisor se
comporta como fonte de Thévenin:
\[
U_L=U_{Th}\cdot\frac{R_L}{R_L+R_{Th}}.
\]
O erro relativo é aproximadamente $R_{Th}/R_L$ para $R_L\gg R_{Th}$. Logo:
\begin{itemize}
\item $R_L=10R_{Th}$ $\Rightarrow$ erro $\approx9\%$;
\item $R_L=100R_{Th}$ $\Rightarrow$ erro $\approx1\%$.
\end{itemize}
\textbf{Consequência de projeto:} o único parâmetro que importa para o
carregamento é $R_{Th}$ --- não $R_1$, não $R_2$, não a razão. Reduzir os dois
resistores proporcionalmente reduz $R_{Th}$ e melhora a robustez, ao custo de
consumo.}
\stepcounter{questao}
\resposta{(a) de $\SI{4.17}{\volt}$ a $\SI{4.90}{\volt}$;
(b) $\approx15\%$}
\resolucao{(a)
\[
R_L=\SI{10}{\kilo\ohm}:\ U=5\cdot\frac{10}{12}=\SI{4.17}{\volt};
\qquad
R_L=\SI{100}{\kilo\ohm}:\ U=5\cdot\frac{100}{102}=\SI{4.90}{\volt}.
\]
(b) Regulação:
\[
\frac{4{,}90-4{,}17}{4{,}90}=15{,}0\%.
\]
(c) \textbf{Como melhorar.}
\begin{itemize}
\item \textbf{Reduzir $R_{Th}$.} Com $R_{Th}=\SI{200}{\ohm}$, a faixa passa a
$4{,}90$--$\SI{4.99}{\volt}$ e a regulação cai para $1{,}8\%$ --- ao custo de
dez vezes mais consumo de repouso.
\item \textbf{Inserir um seguidor de tensão} (amp-op em configuração de ganho
unitário). A impedância de saída cai para frações de ohm e a regulação
praticamente desaparece, sem aumentar o consumo do divisor.
\item \textbf{Usar regulador}, se a carga também exigir corrente apreciável.
\end{itemize}
\textbf{Observação:} a segunda opção é a mais elegante --- ela desacopla
\emph{razão} de \emph{impedância}, que no divisor puro estão inevitavelmente
amarradas.}
\stepcounter{questao}
\resposta{$\varepsilon=\dfrac{R_{Th}}{R_L+R_{Th}}$; aproximadamente
$R_L\ge R_{Th}/\varepsilon$}
\resolucao{(a) Com carga,
$U_L=U_{Th}\dfrac{R_L}{R_L+R_{Th}}$, logo
\[
\varepsilon=\frac{U_{Th}-U_L}{U_{Th}}=1-\frac{R_L}{R_L+R_{Th}}
=\frac{R_{Th}}{R_L+R_{Th}}.
\]
(b) Para $R_L\gg R_{Th}$, o denominador tende a $R_L$:
\[
\varepsilon\approx\frac{R_{Th}}{R_L}
\;\Longrightarrow\;
R_L\ge\frac{R_{Th}}{\varepsilon}.
\]
(c)
\begin{center}
\begin{tabular}{ccc}
\hline
$\varepsilon$ & $R_L/R_{Th}$ (exato) & $R_L/R_{Th}$ (aprox.)\\
\hline
$10\%$ & $9$ & $10$\\
$1\%$ & $99$ & $100$\\
$0{,}1\%$ & $999$ & $1000$\\
\hline
\end{tabular}
\end{center}
A aproximação erra por exatamente $1$ unidade e é conservadora --- pode ser
usada sem receio.\\
\textbf{Regra de bolso resultante:} para $1\%$ de erro,
$R_L\ge100\,R_{Th}$. Note que $R_{Th}=R_1\parallel R_2$ é sempre \emph{menor}
que qualquer um dos resistores, o que torna a condição menos severa do que
comparar $R_L$ com $R_2$ diretamente --- erro comum que leva a
superdimensionar o divisor.}
\stepcounter{questao}
\resposta{(a) $R_1=\SI{19}{\kilo\ohm}$, $R_2=\SI{1}{\kilo\ohm}$;
(b) $0{,}475$ e $\SI{0.025}{\watt}$; (c) $R_L\ge\SI{95}{\kilo\ohm}$}
\resolucao{(a) A restrição de potência fixa a \emph{soma}:
\[
P=\frac{U_{in}^{2}}{R_1+R_2}\le0{,}5
\;\Longrightarrow\;
R_1+R_2\ge\frac{100^{2}}{0{,}5}=\SI{20}{\kilo\ohm}.
\]
A razão fixa a \emph{proporção}: $R_2/(R_1+R_2)=0{,}05$. Adotando o mínimo,
\[
R_2=\SI{1}{\kilo\ohm};\qquad R_1=\SI{19}{\kilo\ohm}.
\]
(b) $I=100/20000=\SI{5}{\milli\ampere}$;
\[
P_1=19000(\pot{5}{-3})^{2}=\SI{0.475}{\watt};
\qquad
P_2=\SI{0.025}{\watt};
\]
com $U_1=\SI{95}{\volt}$ e $U_2=\SI{5}{\volt}$.\\
\textbf{Especificação:} $R_1$ concentra $95\%$ da dissipação --- use
$\SI{1}{\watt}$ nele e $\frac18\,\si{\watt}$ em $R_2$. Alternativamente,
dividir $R_1$ em dois de $\SI{9.5}{\kilo\ohm}$ reparte o calor e a tensão.\\
(c) $R_{Th}=19\parallel1=\SI{950}{\ohm}$; pela condição anterior,
\[
R_L\ge100\,R_{Th}=\SI{95}{\kilo\ohm}.
\]
\textbf{Tensão de projeto:} o divisor entrega $\SI{5}{\volt}$ a uma carga que
consome no máximo $\SI{53}{\micro\ampere}$ --- serve para um conversor A/D, não
para alimentar nada.}
\stepcounter{questao}
\resposta{Em $x=0{,}5$: $\SI{4.00}{\volt}$ contra $\SI{5}{\volt}$ ideal ---
desvio de $20\%$}
\resolucao{(a) A porção inferior vale $xR_p$ e fica em paralelo com $R_L$; a
superior vale $(1-x)R_p$:
\[
U_o=U_{in}\,\frac{xR_p\parallel R_L}{(1-x)R_p+\left(xR_p\parallel R_L\right)}.
\]
Com $R_p=R_L=\SI{10}{\kilo\ohm}$, tem-se
$xR_p\parallel R_L=\dfrac{10x}{1+x}\,\si{\kilo\ohm}$ e
\[
U_o=10\cdot\frac{\dfrac{10x}{1+x}}{10(1-x)+\dfrac{10x}{1+x}}
=10\cdot\frac{x}{(1-x)(1+x)+x}
=\frac{10x}{1+x-x^{2}}.
\]
(b)
\begin{center}
\begin{tabular}{cccc}
\hline
$x$ & ideal & com carga & desvio\\
\hline
$0{,}25$ & $\SI{2.50}{\volt}$ & $\SI{2.105}{\volt}$ & $-15{,}8\%$\\
$0{,}50$ & $\SI{5.00}{\volt}$ & $\SI{4.000}{\volt}$ & $-20{,}0\%$\\
$0{,}75$ & $\SI{7.50}{\volt}$ & $\SI{6.316}{\volt}$ & $-15{,}8\%$\\
\hline
\end{tabular}
\end{center}
(c) O desvio \emph{absoluto} é máximo perto do meio do curso, onde a
resistência de Thévenin do potenciômetro ($x(1-x)R_p$) atinge seu máximo,
$R_p/4=\SI{2.5}{\kilo\ohm}$. Nos extremos ($x\to0$ ou $x\to1$) a impedância
tende a zero e a carga deixa de importar.\\
\textbf{A curva deixa de ser reta:} o potenciômetro linear, carregado, comporta-se
de forma \textbf{não linear}, e o erro depende da posição --- o que é
inaceitável em um controle graduado.\\
\textbf{Correções:} \emph{(i)} usar $R_L\ge100R_p$, o que reduz o desvio máximo
a menos de $0{,}25\%$; \emph{(ii)} inserir um seguidor de tensão entre o cursor
e a carga, solução padrão em instrumentação; \emph{(iii)} usar potenciômetro de
lei "logarítmica reversa" pré-distorcida --- gambiarra que só funciona para uma
carga específica.}
\stepcounter{questao}
\resposta{$\varepsilon=\dfrac{R_1}{R_1+R_2}(t_1+t_2)$; $1{,}00\%$,
$1{,}17\%$ e $1{,}90\%$}
\resolucao{(a) De $U_o=U_{in}\dfrac{R_2}{R_1+R_2}$, tomando logaritmos e
diferenciando:
\[
\frac{\Delta U_o}{U_o}
=\frac{R_1}{R_1+R_2}\left(\frac{\Delta R_2}{R_2}-\frac{\Delta R_1}{R_1}\right),
\]
cujo módulo máximo (desvios opostos) é
$\varepsilon=\dfrac{R_1}{R_1+R_2}(t_1+t_2)$.\\
(b) O fator $\dfrac{R_1}{R_1+R_2}$ vale $1-U_o/U_{in}$:
\begin{center}
\begin{tabular}{ccc}
\hline
Razão & $R_1/(R_1+R_2)$ & Pior erro\\
\hline
$12\to6$ & $0{,}500$ & $1{,}00\%$\\
$12\to5$ & $0{,}583$ & $1{,}17\%$\\
$100\to5$ & $0{,}950$ & $1{,}90\%$\\
\hline
\end{tabular}
\end{center}
\textbf{Padrão:} quanto \emph{maior} a razão de divisão, pior o erro --- ele
tende a $t_1+t_2=2\%$ quando $U_o\ll U_{in}$, e nunca o ultrapassa.\\
(c) \textbf{Resistores casados.} Em uma rede resistiva integrada, os dois
resistores são fabricados lado a lado, no mesmo processo e no mesmo substrato.
Seus desvios são \textbf{correlacionados}: se um sai $0{,}8\%$ acima do
nominal, o outro também sai. O termo entre parênteses,
$\dfrac{\Delta R_2}{R_2}-\dfrac{\Delta R_1}{R_1}$, tende a \emph{zero}.\\
Na prática, redes casadas oferecem razão com precisão de $0{,}05\%$ mesmo com
valores absolutos de $\pm20\%$. \textbf{O que importa no divisor é a razão, não
os valores} --- e o casamento é a única forma barata de obtê-la com precisão.
Esse é o princípio que viabiliza conversores A/D e amplificadores de
instrumentação.}
\stepcounter{questao}
\resposta{(a) $1/10$ e $\SI{10}{\mega\ohm}$; (b) $1{,}0\%$ com sonda,
$9{,}1\%$ sem}
\resolucao{(a) A entrada do osciloscópio \emph{é} o $R_2$ do divisor:
\[
\frac{U_o}{U_{in}}=\frac{1}{9+1}=\frac{1}{10}\ \checkmark;
\qquad
R_{entrada}=9+1=\SI{10}{\mega\ohm}.
\]
(b) \textbf{Com sonda:} o instrumento apresenta $\SI{10}{\mega\ohm}$ à fonte:
\[
\varepsilon=\frac{100\,\mathrm{k}}{10\,\mathrm{M}+100\,\mathrm{k}}=1{,}0\%.
\]
\textbf{Sem sonda} (cabo direto, $\SI{1}{\mega\ohm}$):
\[
\varepsilon=\frac{100\,\mathrm{k}}{1\,\mathrm{M}+100\,\mathrm{k}}=9{,}1\%.
\]
A sonda reduz o erro de carregamento por um fator $9$.\\
(c) \textbf{O que se perde: amplitude.} O sinal chega ao instrumento dez vezes
menor, o que reduz a relação sinal-ruído em $\SI{20}{\decibel}$ e torna
inviável medir sinais de poucos milivolts.\\
\textbf{O compromisso central da sonda} é exatamente esse: troca-se sensibilidade
por não perturbação. Por isso os osciloscópios trazem seletor $1\times/10\times$
--- sinais pequenos em fontes de baixa impedância pedem $1\times$; sinais grandes
em circuitos de alta impedância pedem $10\times$.\\
\emph{(Em sinais variáveis há ainda a capacitância parasita, que exige o
capacitor de compensação ajustável presente em toda sonda real --- assunto de
corrente alternada.)}}
\stepcounter{questao}
\resposta{(b) $R=R_t$; (c) $\SI{-50}{\milli\volt\per\celsius}$}
\resolucao{(a) De $U_o=U\dfrac{R_t}{R+R_t}$:
\[
\frac{\mathrm{d}U_o}{\mathrm{d}R_t}
=U\,\frac{(R+R_t)-R_t}{(R+R_t)^{2}}
=\frac{U\,R}{(R+R_t)^{2}}.
\]
(b) Maximizando em relação a $R$, com $R_t$ fixo:
\[
\frac{\mathrm{d}}{\mathrm{d}R}\left[\frac{R}{(R+R_t)^{2}}\right]
=\frac{(R+R_t)-2R}{(R+R_t)^{3}}
=\frac{R_t-R}{(R+R_t)^{3}}=0
\;\Longrightarrow\;
\boxed{R=R_t}
\]
--- e a saída no ponto ótimo é exatamente $U/2$.\\
\textbf{Mesma estrutura} do resultado de máxima potência: a derivada de
$R/(R+R_t)^{2}$ é a mesma que a de $P_R$ em uma série, e por isso o ótimo cai
em igualdade.\\
(c) Com $R=R_t=\SI{10}{\kilo\ohm}$:
\[
\frac{\mathrm{d}U_o}{\mathrm{d}R_t}=\frac{5(10^{4})}{(\pot{2}{4})^{2}}
=\SI{1.25e-4}{\volt\per\ohm},
\]
\[
\frac{\mathrm{d}U_o}{\mathrm{d}T}
=(\pot{1.25}{-4})(-400)=\SI{-50}{\milli\volt\per\celsius}.
\]
\textbf{Excelente sensibilidade:} um conversor A/D de $10$ bits com fundo de
$\SI{5}{\volt}$ tem degrau de $\SI{4.9}{\milli\volt}$, resolvendo cerca de
$\SI{0.1}{\celsius}$.\\
\textbf{Ressalva:} o ótimo vale para \emph{um} ponto de operação. Em faixa
ampla de temperatura, $R_t$ varia por ordens de grandeza e a sensibilidade cai
nas extremidades --- escolha $R$ igual a $R_t$ na temperatura de maior interesse,
não na média da faixa.}
\stepcounter{questao}
\resposta{(a) $R_1=\SI{70}{\ohm}$, $R_2=\SI{50}{\ohm}$; regulação
$\approx5{,}3\%$; (b) $3{,}8\%$ contra $42\%$}
\resolucao{(a) Repouso de $\SI{100}{\milli\ampere}$:
$R_2=5/0{,}1=\SI{50}{\ohm}$ e
$R_1=7/0{,}1=\SI{70}{\ohm}$.
\[
R_{Th}=50\parallel70=\SI{29.2}{\ohm};
\qquad
\Delta U=R_{Th}\,\Delta I_L=29{,}2(\pot{9}{-3})=\SI{0.26}{\volt},
\]
ou seja, cerca de $5{,}3\%$ de variação.\\
(b) \textbf{Divisor:} consumo total
$\approx12\times0{,}11=\SI{1.32}{\watt}$ para entregar no máximo
$5\times0{,}01=\SI{50}{\milli\watt}$ --- rendimento de $3{,}8\%$.\\
\textbf{Regulador linear:} consome a mesma corrente da carga,
$12\times0{,}01=\SI{120}{\milli\watt}$ para entregar
$\SI{50}{\milli\watt}$ --- rendimento de $41{,}7\%$ (igual à razão
$U_o/U_{in}$), com regulação típica melhor que $0{,}1\%$.\\
(c) \textbf{Critério.}
\begin{itemize}
\item \textbf{Divisor} apenas quando a carga é de \emph{altíssima impedância} e
praticamente constante: entradas de conversores A/D, polarização de bases,
referências internas. Vantagens: dois componentes, custo desprezível, sem ruído
de comutação.
\item \textbf{Regulador} sempre que a carga consumir corrente apreciável ou
variar. O regulador linear resolve com rendimento $U_o/U_{in}$; o chaveado
ultrapassa $90\%$, ao custo de complexidade e ruído.
\end{itemize}
\textbf{Fronteira prática:} se a corrente da carga for maior que cerca de
$1\%$ da corrente de repouso que o divisor exigiria, o divisor já não é a
escolha certa.}
\stepcounter{questao}
\resposta{(a) razão $\approx1/10$: $R_1=\SI{91}{\kilo\ohm}$,
$R_2=\SI{10}{\kilo\ohm}$; (c) $\SI{297}{\micro\ampere}$ no fundo de escala}
\resolucao{(a) É preciso $30\to\SI{3.3}{\volt}$, razão $0{,}11$. Com
$R_2=\SI{10}{\kilo\ohm}$ e $R_1=\SI{91}{\kilo\ohm}$ (ambos E24):
\[
U_o=30\cdot\frac{10}{101}=\SI{2.97}{\volt},
\]
$10\%$ abaixo do fundo de escala --- margem prudente, que acomoda tolerância dos
resistores sem saturar o conversor.\\
(b) Com $\SI{50}{\volt}$ na entrada, a saída iria a
$50\times10/101=\SI{4.95}{\volt}$, acima do limite absoluto do conversor.\\
\textbf{Proteção:} um diodo Zener de $\SI{3.3}{\volt}$ (ou um diodo de
grampeamento para a alimentação) em paralelo com $R_2$. Em sobretensão, ele
conduz e limita a saída; a corrente é confinada por $R_1$ a
\[
I=\frac{50-3{,}3}{91\,\mathrm{k}}=\SI{0.51}{\milli\ampere},
\]
valor inofensivo tanto para o Zener quanto para $R_1$
($P=\SI{24}{\milli\watt}$). \textbf{É $R_1$ que faz a proteção funcionar} ---
sem ele, o diodo teria de absorver corrente ilimitada.\\
(c) No fundo de escala, $I=30/101\,\mathrm{k}=\SI{297}{\micro\ampere}$ e
$P\approx\SI{8.9}{\milli\watt}$.\\
\textbf{Compromisso.} $R_{Th}=91\parallel10=\SI{9}{\kilo\ohm}$. Conversores A/D
de aproximações sucessivas exigem impedância de fonte baixa (tipicamente
$<\SI{10}{\kilo\ohm}$) para carregar o capacitor de amostragem no tempo
disponível --- este projeto fica no limite.\\
\textbf{Soluções:} reduzir os resistores por um fator $10$ (custa dez vezes mais
consumo), acrescentar um capacitor de $\SI{100}{\nano\farad}$ na entrada do
conversor (fornece a carga instantânea de amostragem), ou inserir um seguidor
de tensão. A segunda opção é a mais comum, e a mais barata.}

\gabarito[4.1 Divisor de tensão]

\subsection{Divisor de corrente}
\stepcounter{questao}
\resposta{(b)}
\resolucao{Ambos os ramos têm a mesma tensão; por $I = U/R$, o menor resistor
conduz a maior corrente. A corrente se reparte de forma \emph{inversamente}
proporcional às resistências.}
\stepcounter{questao}
\resposta{$i_1 = \SI{6}{\ampere}$; $i_2 = \SI{4}{\ampere}$}
\resolucao{Pelo divisor de corrente (índices \emph{trocados}):
\[
i_1 = I_t\,\frac{R_2}{R_1+R_2}=10\cdot\frac{3}{5}=\SI{6}{\ampere},
\qquad
i_2 = I_t\,\frac{R_1}{R_1+R_2}=10\cdot\frac{2}{5}=\SI{4}{\ampere}.
\]
\textbf{Verificação:} $6+4 = \SI{10}{\ampere}$; e o menor resistor
(\SI{2}{\ohm}) ficou com a maior corrente.}
\stepcounter{questao}
\resposta{(b)}
\resolucao{$i_{3\Omega} = 15\cdot\dfrac{6}{9} = \SI{10}{\ampere}$ e
$i_{6\Omega} = 15\cdot\dfrac{3}{9} = \SI{5}{\ampere}$. O resistor de \SI{3}{\ohm},
por ser menor, conduz o dobro do de \SI{6}{\ohm}.}
\stepcounter{questao}
\resposta{$i_1=\SI{15}{\ampere}$, $i_2=\SI{10}{\ampere}$, $i_3=\SI{5}{\ampere}$}
\resolucao{Com três ou mais ramos, o mais prático é achar a tensão comum.
$\Req = 2\parallel3\parallel6$: como
$\dfrac{1}{2}+\dfrac{1}{3}+\dfrac{1}{6}=\dfrac{3+2+1}{6}=1$, temos
$\Req=\SI{1}{\ohm}$ e $U = 30\cdot1 = \SI{30}{\volt}$. Então:
\[
i_1=\frac{30}{2}=\SI{15}{\ampere},\quad
i_2=\frac{30}{3}=\SI{10}{\ampere},\quad
i_3=\frac{30}{6}=\SI{5}{\ampere}.
\]
\textbf{Verificação:} $15+10+5 = \SI{30}{\ampere}$.}
\stepcounter{questao}
\resposta{$i_1=i_3\approx\SI{4.17}{\ampere}$; $i_2\approx\SI{16.7}{\ampere}$}
\resolucao{$\dfrac{1}{\Req}=\dfrac{1}{12}+\dfrac{1}{3}+\dfrac{1}{12}
=\dfrac{1+4+1}{12}=\dfrac{6}{12}$, logo $\Req=\SI{2}{\ohm}$ e
$U = 25\cdot2 = \SI{50}{\volt}$.
\[
i_1=i_3=\frac{50}{12}\approx\SI{4.17}{\ampere},
\qquad
i_2=\frac{50}{3}\approx\SI{16.7}{\ampere}.
\]
Os dois ramos de \SI{12}{\ohm} conduzem o mesmo (por simetria); o ramo de
\SI{3}{\ohm}, quatro vezes menor, conduz quatro vezes mais.}
\stepcounter{questao}
\resposta{(a) \SI{20}{\ampere}, \SI{10}{\ampere}, \SI{5}{\ampere};
(b) \SI{3500}{\watt} nos dois lados}
\resolucao{(a) $\dfrac{1}{\Req}=\dfrac{1}{5}+\dfrac{1}{10}+\dfrac{1}{20}
=\dfrac{4+2+1}{20}=\dfrac{7}{20}$, então $\Req = \dfrac{20}{7}\,\si{\ohm}$ e
$U = 35\cdot\dfrac{20}{7} = \SI{100}{\volt}$. As correntes:
\[
i_5=\frac{100}{5}=\SI{20}{\ampere},\quad
i_{10}=\frac{100}{10}=\SI{10}{\ampere},\quad
i_{20}=\frac{100}{20}=\SI{5}{\ampere}.
\]
(b) Potência em cada ramo ($P = U\,i$, com $U=\SI{100}{\volt}$):
\[
100\cdot20 + 100\cdot10 + 100\cdot5 = 2000+1000+500 = \SI{3500}{\watt}.
\]
Potência da fonte: $P = U\,I_t = 100\cdot35 = \SI{3500}{\watt}$ --- confere com a
conservação de energia.}
\stepcounter{questao}
\resposta{$R_s \approx \SI{0.05}{\ohm}$}
\resolucao{O galvanômetro e o shunt estão em paralelo, logo têm a \emph{mesma
tensão}. No fundo de escala, o medidor conduz $\SI{1}{\milli\ampere}$ e o shunt
deve desviar o restante:
\[
i_s = 1 - 0{,}001 = \SI{0.999}{\ampere}.
\]
Igualando as tensões nos dois ramos ($R_g\,i_g = R_s\,i_s$):
\[
R_s = \frac{R_g\,i_g}{i_s} = \frac{50\cdot0{,}001}{0{,}999}\approx\SI{0.05}{\ohm}.
\]
\textbf{Ideia central:} o shunt é um divisor de corrente que amplia a faixa do
instrumento por um fator $\dfrac{I}{i_g} = \dfrac{1}{0{,}001} = 1000$. É assim que
um mesmo galvanômetro vira amperímetro de várias escalas.}
\stepcounter{questao}
\resposta{(a) $R_s = \dfrac{R_g\,i_g}{I_{\max}-i_g}$;
(b) \SI{0.505}{\ohm}, \SI{50.05}{\milli\ohm}, \SI{5.0}{\milli\ohm}; (c) ver comentário}
\resolucao{\textbf{(a)} O shunt está em paralelo com o galvanômetro, logo ambos
têm a mesma tensão. No fundo de escala, o galvanômetro conduz $i_g$ e o shunt
desvia o restante, $I_{\max}-i_g$:
\[
R_g\,i_g = R_s\,(I_{\max}-i_g)
\;\Longrightarrow\;
\boxed{\,R_s=\frac{R_g\,i_g}{I_{\max}-i_g}\,}
\]
\textbf{(b)} Aplicando com $R_g i_g = 50\cdot\pot{1}{-3}=\SI{0.05}{\volt}$:
\begin{center}\small\sffamily
\begin{tabular}{@{}c c c@{}}
\toprule
\textbf{Escala} & \textbf{$R_s$} & \textbf{Fator de multiplicação}\\
\midrule
\SI{100}{\milli\ampere} & $\dfrac{0{,}05}{0{,}099}\approx\SI{0.505}{\ohm}$ & 100$\times$\\
\SI{1}{\ampere}   & $\dfrac{0{,}05}{0{,}999}\approx\SI{50.05}{\milli\ohm}$ & 1000$\times$\\
\SI{10}{\ampere}  & $\dfrac{0{,}05}{9{,}999}\approx\SI{5.0}{\milli\ohm}$ & 10\,000$\times$\\
\bottomrule
\end{tabular}
\end{center}
\textbf{(c)} A resistência total do instrumento é $R_g \parallel R_s$. Como $R_s$
diminui a cada escala maior, o paralelo também diminui --- na escala de
\SI{10}{\ampere}, o amperímetro apresenta apenas $\approx\SI{5}{\milli\ohm}$.\\
\textbf{Por que isso é desejável:} o amperímetro é ligado \emph{em série} com o
circuito, então sua resistência se soma à do circuito e \emph{altera} a corrente
que se pretende medir --- o chamado \textbf{erro de inserção}. Quanto menor a
resistência do instrumento, menor a perturbação. O amperímetro ideal teria
resistência nula (e o voltímetro ideal, resistência infinita).\\
\textbf{Estimativa do erro:} medindo \SI{10}{\ampere} em um circuito de
$\SI{1}{\ohm}$, o erro relativo seria $\dfrac{0{,}005}{1}=\SI{0.5}{\percent}$ ---
aceitável. Mas em um circuito de $\SI{0.05}{\ohm}$, o erro subiria para
\SI{10}{\percent}, tornando a medida inútil.}
\stepcounter{questao}
\resposta{$\SI{5}{\ampere}$, $\SI{3}{\ampere}$ e $\SI{2}{\ampere}$}
\resolucao{Com $\sum G=0{,}5+0{,}3+0{,}2=\SI{1}{\ensuremath{\mathrm{S}}}$:
\[
I_k=I_T\frac{G_k}{\sum G}
\;\Longrightarrow\;
5;\ 3;\ \SI{2}{\ampere}.
\]
Em condutâncias a conta é imediata: a corrente se reparte na \emph{mesma
proporção} das condutâncias. É por isso que essa é a forma natural do divisor de
corrente, e não a expressão com resistências.}
\stepcounter{questao}
\resposta{$I_1/I_2=3$}
\resolucao{Como a tensão é comum, $I=U/R$ e
\[
\frac{I_1}{I_2}=\frac{R_2}{R_1}=\frac{12}{4}=3.
\]
A razão das correntes é a \textbf{inversa} da razão das resistências --- o ramo
de menor resistência leva a maior corrente. Compare com a série, em que a razão
das \emph{tensões} é \emph{direta}.}
\stepcounter{questao}
\resposta{$I_T=\SI{2.5}{\ampere}$}
\resolucao{A tensão comum sai do ramo conhecido:
\[
U=20(1{,}5)=\SI{30}{\volt}
\;\Longrightarrow\;
I_2=\frac{30}{30}=\SI{1}{\ampere},
\]
\[
I_T=1{,}5+1=\SI{2.5}{\ampere}.
\]
Um único ramo conhecido determina a tensão comum e, com ela, todo o resto ---
propriedade que torna qualquer resistor de um paralelo um sensor de corrente em
potencial.}
\stepcounter{questao}
\resposta{Toda a corrente passa pelo curto; os demais ramos ficam sem corrente}
\resolucao{O ramo em curto tem $R=0$, logo $G\to\infty$. Pelo divisor,
\[
I_k=I_T\frac{G_k}{\infty+\sum_{j\ne k}G_j}\to0
\]
para todos os ramos finitos, e $I_{curto}\to I_T$.\\
\textbf{Fisicamente:} o curto impõe $U=0$ à associação inteira, e sem tensão
nenhum resistor conduz. Os demais ramos continuam íntegros --- eles apenas
deixam de funcionar, sendo vítimas e não causa do defeito.}
\stepcounter{questao}
\resposta{(b)}
\resolucao{A fonte de corrente ideal mantém $I_T$ fixo. Com um caminho a menos,
a mesma corrente se distribui entre menos ramos --- cada um recebe \emph{mais}.\\
\textbf{Contraste essencial:} sob fonte de \emph{tensão} ideal, os ramos
restantes ficariam \emph{inalterados} (a tensão comum não muda) e apenas a
corrente total cairia. A mesma falha produz efeitos opostos conforme o tipo de
alimentação --- ponto explorado quantitativamente no bloco de
aprofundamento.}
\stepcounter{questao}
\resposta{$U=\SI{30}{\volt}$; $6$, $3$, $1{,}5$ e $\SI{1.5}{\ampere}$}
\resolucao{
\[
\sum G=0{,}20+0{,}10+0{,}05+0{,}05=\SI{0.40}{\ensuremath{\mathrm{S}}}
\;\Longrightarrow\;
U=\frac{I_T}{\sum G}=\frac{12}{0{,}40}=\SI{30}{\volt}.
\]
\[
I_k=\frac{U}{R_k}\Rightarrow 6;\ 3;\ 1{,}5;\ \SI{1.5}{\ampere}.
\]
\textbf{Verificação:} $6+3+1{,}5+1{,}5=\SI{12}{\ampere}$ \checkmark.\\
\textbf{Método:} calcular a tensão comum primeiro é quase sempre mais rápido do
que aplicar a fórmula do divisor ramo a ramo --- e produz um valor intermediário
que serve de conferência.}
\stepcounter{questao}
\resposta{$50\%$, $25\%$, $12{,}5\%$ e $12{,}5\%$}
\resolucao{
\[
\frac{I_k}{I_T}=\frac{G_k}{\sum G}
\Rightarrow \frac{0{,}20}{0{,}40}=50\%;\quad25\%;\quad12{,}5\%;\quad12{,}5\%.
\]
A razão depende \emph{apenas} das condutâncias --- $I_T$ cancela. Dobrar a
corrente total dobra todas as parcelas, preservando as proporções.\\
\textbf{Consequência prática:} o perfil de divisão é uma propriedade fixa da
rede. Ele só muda se algum ramo mudar de valor, for aberto ou entrar em curto
--- o que faz da monitoração dessas frações uma técnica de diagnóstico.}
\stepcounter{questao}
\resposta{$R=\SI{45}{\ohm}$}
\resolucao{As correntes são $I_R=\SI{2}{\ampere}$ e
$I_{15}=\SI{6}{\ampere}$. A tensão é comum:
\[
U=15(6)=\SI{90}{\volt}
\;\Longrightarrow\;
R=\frac{90}{2}=\SI{45}{\ohm}.
\]
\textbf{Conferência pela razão:} $I_R/I_{15}=R_{15}/R=15/45=1/3$, e de fato
$2/6=1/3$ \checkmark.\\
Note que $R$ é o \emph{triplo} do outro ramo para conduzir um \emph{terço} da
corrente dele --- a proporcionalidade inversa em ação.}
\stepcounter{questao}
\resposta{$I_T=\SI{6}{\ampere}$; $\SI{3.6}{\ampere}$ e $\SI{2.4}{\ampere}$}
\resolucao{(a) $R_{par}=\dfrac{10\cdot15}{25}=\SI{6}{\ohm}$ e
\[
I_T=\frac{60}{4+6}=\SI{6}{\ampere}.
\]
(b) Pelo divisor de dois ramos:
\[
I_1=6\cdot\frac{15}{25}=\SI{3.6}{\ampere};
\qquad
I_2=6\cdot\frac{10}{25}=\SI{2.4}{\ampere}.
\]
\textbf{Ponto de atenção:} $R_s$ afeta a corrente \emph{total}, mas não a
\emph{proporção} da divisão --- ela é determinada só pelos dois ramos. Essa
separação entre "quanto chega" e "como se reparte" simplifica muito a análise de
circuitos reais.}
\stepcounter{questao}
\resposta{$U=\SI{36}{\volt}$; $\SI{129.6}{\watt}$ e $\SI{86.4}{\watt}$}
\resolucao{
\[
U=R_{par}I_T=6(6)=\SI{36}{\volt};
\qquad
P_1=\frac{36^{2}}{10}=\SI{129.6}{\watt};
\qquad
P_2=\frac{36^{2}}{15}=\SI{86.4}{\watt}.
\]
Total nos ramos: $\SI{216}{\watt}=36(6)$ \checkmark.\\
\textbf{Observe:} o ramo de \emph{menor} resistência dissipa mais --- regra do
paralelo ($P\propto1/R$ com $U$ comum). O resistor em série dissipa ainda
$4(36)=\SI{144}{\watt}$, e a fonte entrega $60(6)=\SI{360}{\watt}$
\checkmark.}
\stepcounter{questao}
\resposta{$\SI{9.99}{\ampere}$ e $\SI{9.99}{\milli\ampere}$; erro de
$0{,}1\%$}
\resolucao{(a)
\[
I_1=10\cdot\frac{1000}{1001}=\SI{9.990}{\ampere};
\qquad
I_2=10\cdot\frac{1}{1001}=\SI{9.99}{\milli\ampere}.
\]
(b) Desprezar o segundo ramo daria $I_1=\SI{10}{\ampere}$ --- erro de
$0{,}1\%$, muito abaixo da tolerância de qualquer resistor real.\\
\textbf{Regra prática:} um ramo mil vezes maior desvia um milésimo da corrente.
Se a razão for de $100{:}1$, o erro sobe para $1\%$; se for $10{:}1$, para
$9\%$ --- aí já não se pode desprezar.\\
\textbf{Mas atenção:} desprezar o ramo para calcular $I_1$ é legítimo;
\emph{esquecê-lo} pode não ser, pois aqueles $\SI{10}{\milli\ampere}$ podem ser
exatamente o sinal que interessa medir.}
\stepcounter{questao}
\resposta{Antes: $6$, $3$, $\SI{3}{\ampere}$. Depois: $\SI{6}{\ampere}$ em cada}
\resolucao{\textbf{Antes:} $\sum G=\frac16+\frac1{12}+\frac1{12}
=\SI{0.3333}{\ensuremath{\mathrm{S}}}$, $U=12/0{,}3333=\SI{36}{\volt}$ e
$I=6;\ 3;\ \SI{3}{\ampere}$.\\
\textbf{Depois:} $\sum G=\SI{0.1667}{\ensuremath{\mathrm{S}}}$,
$U=12/0{,}1667=\SI{72}{\volt}$ e $I=\SI{6}{\ampere}$ em cada.\\
\textbf{A tensão dobrou} e a corrente de cada ramo remanescente também. Note que
os resistores não mudaram --- foi a \emph{fonte de corrente} que, obrigada a
manter $\SI{12}{\ampere}$, elevou a tensão para forçar a mesma corrente por
menos caminhos.\\
\textbf{Alerta de projeto:} em circuitos alimentados por fonte de corrente,
abrir um ramo \emph{sobrecarrega} os demais. É o oposto do que ocorre com fonte
de tensão.}
\stepcounter{questao}
\resposta{$\SI{0}{\ampere}$ no resistor; $\SI{4}{\ampere}$ no indutor}
\resolucao{Em regime permanente CC, a corrente do indutor é constante e
$v=L\,\mathrm{d}i/\mathrm{d}t=0$: ele se comporta como um
\textbf{curto-circuito}.\\
Pelo divisor, $G_L\to\infty$ e toda a corrente vai para o indutor:
\[
I_L=\SI{4}{\ampere};
\qquad
I_R=0.
\]
\textbf{Consequência:} um indutor em paralelo com um resistor "apaga" o resistor
em CC. Isso é usado deliberadamente para curto-circuitar componentes em regime e
deixá-los ativos apenas durante transitórios.\\
\emph{(Um capacitor em paralelo faria o oposto: em CC ele é um aberto e toda a
corrente iria ao resistor.)}}
\stepcounter{questao}
\resposta{Consistente; $\SI{5}{\ohm}$, $\SI{10}{\ohm}$ e $\SI{15}{\ohm}$}
\resolucao{\textbf{Resistências:} $R_k=U/I_k$:
\[
\frac{20}{4}=5;\qquad \frac{20}{2}=10;\qquad \frac{20}{1{,}33}=\SI{15}{\ohm}.
\]
\textbf{Consistência:} $\sum G=0{,}2+0{,}1+0{,}0667=\SI{0.3667}{\ensuremath{\mathrm{S}}}$, logo
$I_T=20(0{,}3667)=\SI{7.33}{\ampere}$, que confere com
$4+2+1{,}33$ \checkmark.\\
\textbf{Dois testes independentes:} a tensão comum deve ser a mesma em todos os
ramos (senão as leituras são incompatíveis) e a LKC deve fechar. Se apenas o
segundo fechasse, o erro estaria na leitura de tensão.}
\stepcounter{questao}
\resposta{A forma correta usa condutâncias; a versão com resistências vale só
para dois ramos}
\resolucao{\textbf{Contraexemplo.} Ramos de $\SI{6}{\ohm}$, $\SI{12}{\ohm}$ e
$\SI{12}{\ohm}$ com $I_T=\SI{12}{\ampere}$.\\
\emph{Forma correta:} $I_1=12\cdot\dfrac{1/6}{1/3}=\SI{6}{\ampere}$.\\
\emph{"Generalização" ingênua} $I_1=I_T\dfrac{R_2+R_3}{R_1+R_2+R_3}$:
$12\cdot\dfrac{24}{30}=\SI{9.6}{\ampere}$ --- errado, e nem sequer a soma das
três parcelas assim calculadas daria $\SI{12}{\ampere}$.\\
\textbf{Origem da confusão:} com \emph{dois} ramos,
$\dfrac{G_1}{G_1+G_2}=\dfrac{1/R_1}{1/R_1+1/R_2}=\dfrac{R_2}{R_1+R_2}$ --- a
simplificação funciona por acaso algébrico. Com três, a manipulação equivalente
produziria $\dfrac{R_2R_3}{R_2R_3+R_1R_3+R_1R_2}$, expressão que ninguém
memoriza.\\
\textbf{Recomendação:} use sempre condutâncias. A fórmula com resistências é
uma conveniência de dois ramos, não a regra geral.}
\stepcounter{questao}
\resposta{$109{,}4$; $99{,}5$ e $\SI{91.2}{\ampere}$; desvios de
$+9{,}4\%$, $-0{,}5\%$ e $-8{,}8\%$}
\resolucao{(a) $\sum G=10+9{,}091+8{,}333=\SI{27.42}{\ensuremath{\mathrm{S}}}$ e
$U=300/27{,}42=\SI{10.94}{\volt}$. As correntes valem
\[
109{,}4;\qquad 99{,}5;\qquad \SI{91.2}{\ampere}.
\]
(b) Em relação a $\SI{100}{\ampere}$ ideais: $+9{,}4\%$, $-0{,}5\%$ e
$-8{,}8\%$.\\
(c) \textbf{Implicação térmica.} A dissipação vai com $I^{2}$: o condutor mais
carregado dissipa $(109{,}4/100)^{2}=1{,}20$ vezes o previsto --- $20\%$ a
mais. Ele aquece mais que os outros dois.\\
\textbf{E aí surge a pergunta decisiva:} o cobre tem coeficiente térmico
\emph{positivo}. Aquecendo, sua resistência sobe, ele devolve corrente aos
vizinhos e o sistema se \textbf{estabiliza}. Se o material tivesse coeficiente
negativo, ocorreria o contrário --- fuga térmica.\\
\textbf{Prática de instalação:} condutores em paralelo devem ter mesmo
comprimento, seção, material e trajeto. Uma diferença de $\SI{20}{\percent}$ no
comprimento produz o desequilíbrio calculado aqui, e é por isso que a norma
exige lances idênticos.}
\stepcounter{questao}
\resposta{$2$, $4$ e $\SI{6}{\ampere}$; $\SI{6}{\ohm}$, $\SI{3}{\ohm}$ e
$\SI{2}{\ohm}$; $24$, $48$ e $\SI{72}{\watt}$}
\resolucao{(a) As frações são $\frac16$, $\frac26$ e $\frac36$ de
$\SI{12}{\ampere}$: $2$, $4$ e $\SI{6}{\ampere}$.\\
(b) Como $U$ é comum, $R_k=U/I_k$:
\[
\frac{12}{2}=6;\qquad \frac{12}{4}=3;\qquad \frac{12}{6}=\SI{2}{\ohm}.
\]
As resistências ficam na razão $6{:}3{:}2$ --- \emph{inversa} de
$1{:}2{:}3$, como esperado.\\
(c) $P_k=UI_k$: $24$, $48$ e $\SI{72}{\watt}$; total $\SI{144}{\watt}=12(12)$
\checkmark.\\
\textbf{Observação de projeto:} a tensão comum foi \emph{escolhida}. Adotando
$\SI{1.2}{\volt}$ em vez de $\SI{12}{\volt}$, as resistências cairiam por dez
($0{,}6$; $0{,}3$; $\SI{0.2}{\ohm}$) e a dissipação total também --- de
$144$ para $\SI{14.4}{\watt}$. A razão de divisão é a mesma; o que muda é a
perda. Esse grau de liberdade é explorado no bloco de desafio.}
\stepcounter{questao}
\resposta{$5{,}25$ e $\SI{4.75}{\ampere}$; desvio de $\pm5\%$}
\resolucao{(a) O pior caso ocorre com desvios opostos: $R_1=\SI{9.5}{\ohm}$ e
$R_2=\SI{10.5}{\ohm}$.
\[
I_1=10\cdot\frac{10{,}5}{20}=\SI{5.25}{\ampere};
\qquad
I_2=\SI{4.75}{\ampere}.
\]
(b) Em relação aos $\SI{5}{\ampere}$ ideais, o desvio é de
$\pm5\%$ --- \emph{igual} à tolerância dos resistores.\\
\textbf{Por que dá exatamente isso:} a sensibilidade relativa da divisão é
$\dfrac{R_1}{R_1+R_2}=0{,}5$, e o pior caso combina as duas tolerâncias:
$0{,}5\times(5\%+5\%)=5\%$.\\
\textbf{Generalização:} com ramos \emph{desiguais}, o fator muda. Em um divisor
$1{:}9$, a sensibilidade do ramo pequeno é $0{,}9$ e o desequilíbrio chega a
$9\%$ --- o ramo minoritário é sempre o mais afetado por tolerâncias.}
\stepcounter{questao}
\resposta{$\SI{47.6}{\ampere}$ e $\SI{52.4}{\ampere}$; desvio de
$\mp4{,}8\%$; $\SI{113}{\watt}$ no contato}
\resolucao{(a) O ramo afetado passa a $\SI{0.55}{\ohm}$:
\[
I_1=100\cdot\frac{0{,}50}{1{,}05}=\SI{47.6}{\ampere};
\qquad
I_2=100\cdot\frac{0{,}55}{1{,}05}=\SI{52.4}{\ampere}.
\]
(b) Desvio de $-4{,}8\%$ e $+4{,}8\%$ em relação aos
$\SI{50}{\ampere}$ ideais.\\
No contato: $P=0{,}05(47{,}6)^{2}=\SI{113}{\watt}$ --- concentrados em uma área
minúscula.\\
\textbf{Por que é grave.} A resistência de contato \emph{cresce} com a corrosão
e com a temperatura. Mais resistência $\to$ mais aquecimento $\to$ mais
oxidação $\to$ mais resistência: realimentação positiva que termina em falha
térmica. Curiosamente, o desvio de corrente é \emph{estabilizador} (o ramo ruim
conduz menos), mas a concentração de potência no ponto de contato domina o
processo.\\
\textbf{Conclusão:} em divisores de alta corrente, a qualidade das conexões
importa mais que a tolerância dos resistores.}
\stepcounter{questao}
\resposta{$8$ e $\SI{2}{\ampere}$; $320$ e $\SI{80}{\watt}$ --- o menor
resistor dissipa $4\times$ mais}
\resolucao{(a)
\[
I_1=10\cdot\frac{20}{25}=\SI{8}{\ampere};
\qquad
I_2=\SI{2}{\ampere};
\]
\[
P_1=5(8)^{2}=\SI{320}{\watt};
\qquad
P_2=20(2)^{2}=\SI{80}{\watt}.
\]
(b) O ramo de $\SI{5}{\ohm}$ dissipa quatro vezes mais.\\
\textbf{Razão algébrica:} com $U$ comum, $P=U^{2}/R$, logo
$P_1/P_2=R_2/R_1=4$. A razão das potências é a \emph{inversa} da razão das
resistências --- e igual à razão das correntes ($8/2=4$).\\
\textbf{Erro comum:} supor que o resistor maior aquece mais. Isso vale em
\emph{série} ($P\propto R$), não em paralelo. Em um divisor de corrente, o
componente crítico para dimensionamento térmico é sempre o de \emph{menor}
resistência.}
\stepcounter{questao}
\resposta{(a) tensão dobra ($60\to\SI{120}{\volt}$), ramos vão de $3$ a
$\SI{6}{\ampere}$; (b) tensão e ramos inalterados, total cai de $12$ a
$\SI{6}{\ampere}$}
\resolucao{\textbf{(a) Fonte de corrente.} Antes:
$\sum G=\SI{0.2}{\ensuremath{\mathrm{S}}}$, $U=\SI{60}{\volt}$,
$I=6;3;\SI{3}{\ampere}$.\\
Depois: $\sum G=\SI{0.1}{\ensuremath{\mathrm{S}}}$, e a fonte \emph{eleva a tensão} para
manter $\SI{12}{\ampere}$:
\[
U=\frac{12}{0{,}1}=\SI{120}{\volt};
\qquad
I=6;\ \SI{6}{\ampere}.
\]
Cada ramo remanescente \textbf{dobra} de corrente, e a potência quadruplica
(de $180$ para $\SI{720}{\watt}$ em cada).\\
\textbf{(b) Fonte de tensão.} $U=\SI{60}{\volt}$ permanece.
Antes: $6;3;\SI{3}{\ampere}$, total $\SI{12}{\ampere}$. Depois:
$3;\SI{3}{\ampere}$, total $\SI{6}{\ampere}$. \textbf{Os ramos remanescentes
não sentem nada.}\\
\textbf{Síntese.} A mesma falha, no mesmo circuito, produz sobrecarga severa em
um caso e nenhum efeito no outro. O que decide é a natureza da fonte:
\begin{itemize}
\item fonte de \emph{tensão} $\Rightarrow$ ramos desacoplados, falha localizada;
\item fonte de \emph{corrente} $\Rightarrow$ ramos acoplados, falha propagada.
\end{itemize}
É por isso que instalações elétricas são alimentadas por tensão, e por que
circuitos em série de LEDs acionados por corrente exigem proteção contra
abertura.}
\stepcounter{questao}
\resposta{(a) de $0$ a $\SI{4}{\ampere}$; (b) $R=\SI{10}{\ohm}$}
\resolucao{(a) Com o reostato em $\SI{0}{\ohm}$, ele curto-circuita o ramo
fixo: $I_{fixo}=0$. Com $\SI{40}{\ohm}$:
\[
I_{fixo}=5\cdot\frac{40}{50}=\SI{4}{\ampere}.
\]
A faixa é $0$ a $\SI{4}{\ampere}$ --- \emph{nunca} os $\SI{5}{\ampere}$
completos, pois o reostato sempre desvia alguma corrente.\\
(b) Divisão igual exige resistências iguais: $R=\SI{10}{\ohm}$, um quarto do
curso.\\
(c) \textbf{Fortemente não linear.} A relação
$I_{fixo}=5R/(R+10)$ é hiperbólica:
\begin{center}
\begin{tabular}{ccc}
\hline
$R$ & posição do curso & $I_{fixo}$\\
\hline
$\SI{10}{\ohm}$ & $25\%$ & $\SI{2.5}{\ampere}$\\
$\SI{20}{\ohm}$ & $50\%$ & $\SI{3.33}{\ampere}$\\
$\SI{40}{\ohm}$ & $100\%$ & $\SI{4.0}{\ampere}$\\
\hline
\end{tabular}
\end{center}
Metade da faixa útil de corrente é percorrida no primeiro quarto do curso --- o
ajuste é grosseiro no início e quase inerte no fim.}
\stepcounter{questao}
\resposta{$I_k=I_T\dfrac{R_p}{R_k+R_p}$; para $R_k=\SI{6}{\ohm}$ e
$R_p=\SI{6}{\ohm}$, $I_k=\SI{6}{\ampere}$}
\resolucao{\textbf{Dedução.} Agrupando todos os ramos exceto o $k$-ésimo em
$R_p$, sobram \emph{dois} ramos em paralelo --- e aí a fórmula de dois ramos é
válida:
\[
I_k=I_T\frac{R_p}{R_k+R_p},
\qquad\text{com }
\frac{1}{R_p}=\sum_{j\ne k}\frac{1}{R_j}.
\]
\textbf{Verificação.} Para $R_k=\SI{6}{\ohm}$, os demais dão
$R_p=12\parallel12=\SI{6}{\ohm}$:
\[
I_k=12\cdot\frac{6}{6+6}=\SI{6}{\ampere}.\ \checkmark
\]
\textbf{Utilidade:} é a maneira correta de \emph{recuperar} a fórmula de dois
ramos em uma rede com muitos ramos --- e mostra exatamente onde a
"generalização" ingênua erra: $R_p$ é o \emph{paralelo} dos demais, não a
\emph{soma}. No exemplo, a soma daria $\SI{24}{\ohm}$ e a corrente
$\SI{9.6}{\ampere}$, valor incorreto.}
\stepcounter{questao}
\resposta{$I_1=\SI{1}{\ampere}$ e $I_2=\SI{3}{\ampere}$}
\resolucao{(a) O divisor de corrente pressupõe ramos \textbf{puramente
resistivos}, submetidos à mesma tensão e obedecendo $I=UG$. Um ramo com fonte
segue $I=(U-E)G$ --- relação \emph{afim}, não proporcional. A repartição deixa
de depender só das condutâncias.\\
(b) Chamando $U$ a tensão comum:
\[
\underbrace{\frac{U-10}{5}}_{I_1}+\underbrace{\frac{U}{5}}_{I_2}=4
\;\Longrightarrow\;
2U-10=20
\;\Longrightarrow\;
U=\SI{15}{\volt},
\]
\[
I_1=\frac{15-10}{5}=\SI{1}{\ampere};
\qquad
I_2=\frac{15}{5}=\SI{3}{\ampere}.
\]
\textbf{Compare:} o divisor simples, com dois ramos iguais, preveria
$\SI{2}{\ampere}$ em cada. A fonte de $\SI{10}{\volt}$ \emph{se opõe} à corrente
em seu ramo e desvia $\SI{1}{\ampere}$ para o outro.\\
\textbf{Método correto:} volte à LKC com a tensão comum como incógnita --- ou
seja, análise nodal. O divisor de corrente é um atalho, e atalhos têm
hipóteses.}
\stepcounter{questao}
\resposta{$I_k=I_T\dfrac{G_k}{\sum_j G_j}$}
\resolucao{(a) Todos os ramos compartilham a tensão $U$. Pela LKC,
\[
I_T=\sum_j I_j=\sum_j UG_j=U\sum_j G_j
\;\Longrightarrow\;
U=\frac{I_T}{\sum_j G_j}.
\]
Logo
\[
I_k=UG_k=I_T\frac{G_k}{\sum_j G_j}.
\]
(b) Para $n=2$, substituindo $G=1/R$:
\[
\frac{G_1}{G_1+G_2}=\frac{1/R_1}{\dfrac{R_1+R_2}{R_1R_2}}
=\frac{R_2}{R_1+R_2}.
\]
A simplificação depende de o denominador ter \emph{exatamente} dois termos.
Para $n=3$, o resultado análogo seria
\[
I_1=I_T\frac{R_2R_3}{R_2R_3+R_1R_3+R_1R_2},
\]
que não tem a forma "resistência do outro sobre a soma" --- e para $n$ geral
envolve todos os produtos de $n-1$ resistências.\\
(c)
\[
\sum_k I_k=I_T\frac{\sum_k G_k}{\sum_j G_j}=I_T.
\]
A LKC é satisfeita \textbf{identicamente}, para quaisquer condutâncias --- ela
não é um resultado do cálculo, mas uma consequência da forma da expressão.
Verificar a soma, portanto, não detecta erro nas condutâncias; ela só detecta
erro de aritmética.}
\stepcounter{questao}
\resposta{$R\le\SI{8.33}{\milli\ohm}$; queda de $\SI{0.167}{\volt}$;
$\SI{3.33}{\watt}$ por ramo no limite}
\resolucao{(a) Divisão igual exige três resistências iguais, com
$\SI{20}{\ampere}$ cada. A perda total é
\[
P=3RI_k^{2}=3R(400)=1200R\le10
\;\Longrightarrow\;
R\le\SI{8.33}{\milli\ohm}.
\]
Adotando $R=\SI{8.33}{\milli\ohm}$, a queda é
$U=8{,}33\,\mathrm{m}\times20=\SI{0.167}{\volt}$.\\
(b) No limite, cada ramo dissipa $\SI{3.33}{\watt}$ --- exigindo resistores de
pelo menos $\SI{10}{\watt}$ para operar com folga, e provavelmente ventilados.\\
(c) \textbf{O compromisso central.} O desequilíbrio relativo entre os ramos é
governado pela \emph{razão} entre a dispersão dos resistores de balanceamento e
as resistências \emph{parasitas} (contatos, trechos de barramento, cabos):
\begin{itemize}
\item $R$ \textbf{grande}: as parasitas tornam-se desprezíveis em comparação, e
a divisão fica precisa --- mas a perda cresce com $R$.
\item $R$ \textbf{pequeno}: pouca perda, mas as parasitas passam a dominar e o
balanceamento se degrada.
\end{itemize}
Com $R=\SI{8.33}{\milli\ohm}$ e contatos de $\SI{1}{\milli\ohm}$, a incerteza
de divisão já é da ordem de $12\%$ --- o resistor de balanceamento praticamente
não cumpre sua função.\\
\textbf{Conclusão:} não se pode ter simultaneamente divisão precisa e perda
desprezível com resistores passivos. Aplicações críticas usam balanceamento
\emph{ativo} (realimentação de corrente por ramo), que rompe o compromisso.}
\stepcounter{questao}
\resposta{Desvio relativo $\approx\dfrac{n-1}{n}\,\delta$; para $n=3$:
$3{,}4\%$, $7{,}1\%$ e $15{,}4\%$}
\resolucao{(a) As condutâncias são $\dfrac{1}{R(1-\delta)}$ e
$(n-1)$ vezes $\dfrac1R$:
\[
\frac{I_1}{I_T}=\frac{\dfrac{1}{1-\delta}}{\dfrac{1}{1-\delta}+(n-1)}
=\frac{1}{1+(n-1)(1-\delta)}.
\]
(b) Para $\delta\ll1$, expandindo em torno de $\delta=0$ (onde a fração vale
$1/n$):
\[
\frac{I_1}{I_T}\approx\frac1n\left(1+\frac{n-1}{n}\,\delta\right)
\;\Longrightarrow\;
\frac{\Delta I_1}{I_1}\approx\frac{n-1}{n}\,\delta.
\]
O \textbf{fator de amplificação} é $\dfrac{n-1}{n}$: sempre menor que $1$, mas
crescente com $n$ --- tende a $1$ para muitos condutores.\\
(c) Com $n=3$, o fator é $2/3$:
\begin{center}
\begin{tabular}{ccc}
\hline
$\delta$ & exato & aproximado\\
\hline
$5\%$ & $+3{,}45\%$ & $+3{,}33\%$\\
$10\%$ & $+7{,}14\%$ & $+6{,}67\%$\\
$20\%$ & $+15{,}4\%$ & $+13{,}3\%$\\
\hline
\end{tabular}
\end{center}
\textbf{Duas leituras.} A aproximação é boa até cerca de $10\%$ e subestima o
desvio real --- portanto \emph{não} é conservadora, e deve ser usada com
cuidado em verificação de projeto.\\
E o desequilíbrio de corrente é sempre \emph{menor} que o de resistência (fator
$<1$), o que é uma boa notícia. A má notícia é que a \textbf{potência} vai com
$I^{2}$: um desvio de $7\%$ na corrente significa $15\%$ a mais de
aquecimento no condutor mais carregado.}
\stepcounter{questao}
\resposta{$S_1=-\dfrac{R_1}{R_1+R_2}$, $S_2=+\dfrac{R_2}{R_1+R_2}$}
\resolucao{(a) De $I_1=I_T\dfrac{R_2}{R_1+R_2}$:
\[
\frac{\partial I_1}{\partial R_1}=-\frac{I_TR_2}{(R_1+R_2)^{2}}.
\]
(b) Sensibilidade relativa:
\[
S_1=\frac{\partial I_1}{\partial R_1}\cdot\frac{R_1}{I_1}
=-\frac{I_TR_2}{(R_1+R_2)^{2}}\cdot\frac{R_1(R_1+R_2)}{I_TR_2}
=-\frac{R_1}{R_1+R_2}.
\]
Analogamente, $S_2=+\dfrac{R_2}{R_1+R_2}$, e
\[
|S_1|+|S_2|=\frac{R_1+R_2}{R_1+R_2}=1.
\]
\textbf{Mesmo padrão} já encontrado em série, paralelo e divisor de tensão ---
consequência de $I_1/I_T$ ser função homogênea de grau zero nas resistências.\\
(c) \textbf{$R_1\ll R_2$:} $S_1\to0$. O ramo pequeno leva quase toda a corrente
e é \emph{insensível} ao próprio valor --- mas $S_2\to1$, ou seja, sua corrente
depende integralmente do ramo grande. Contraintuitivo e importante.\\
\textbf{$R_1\gg R_2$:} $S_1\to-1$. O ramo grande conduz pouco, e essa pouca
corrente é inteiramente determinada pelo seu próprio valor.\\
\textbf{Regra de projeto:} para dividir corrente com precisão, aperte a
tolerância do ramo que conduz \emph{menos} --- é ele que carrega a maior
sensibilidade relativa.}
\stepcounter{questao}
\resposta{Ramo de medição de $\SI{100}{\milli\ohm}$; $\SI{99.9}{\milli\ampere}$;
perda de $\SI{1}{\watt}$ contra $\SI{10}{\watt}$}
\resolucao{(a) A razão de correntes é a razão inversa das resistências:
\[
\frac{I_{med}}{I_{princ}}=\frac{R_{princ}}{R_{med}}=\frac{1}{1000}
\;\Longrightarrow\;
R_{med}=1000\times0{,}1\,\mathrm{m}=\SI{100}{\milli\ohm}.
\]
(b)
\[
G_{tot}=10^{4}+10=\SI{10010}{\ensuremath{\mathrm{S}}};
\qquad
I_{med}=100\cdot\frac{10}{10010}=\SI{99.9}{\milli\ampere}.
\]
\[
R_{eq}=\SI{99.9}{\micro\ohm};
\qquad
U=100(\pot{99.9}{-6})=\SI{9.99}{\milli\volt};
\qquad
P=100(\pot{9.99}{-3})=\SI{0.999}{\watt}.
\]
(c) \textbf{Shunt direto de $\SI{1}{\milli\ohm}$:}
$U=\SI{100}{\milli\volt}$ e $P=\SI{10}{\watt}$ --- \textbf{dez vezes} mais
perda.\\
\textbf{O compromisso, explicitado.} A perda vale $P=I_T^{2}R_{eq}$, e
$R_{eq}$ é essencialmente o ramo principal. Reduzi-lo reduz a perda \emph{e} a
tensão de medição --- que já é de apenas $\SI{10}{\milli\volt}$, exigindo
amplificação de baixo ruído e cuidado com termopares parasitas nas junções.\\
\textbf{Alternativa moderna:} sensores de efeito Hall ou transformadores de
corrente medem sem inserir resistência alguma no caminho principal, eliminando a
perda por completo --- ao custo de exatidão e de resposta em CC.}
\stepcounter{questao}
\resposta{$\alpha>0$: estável (realimentação negativa); $\alpha<0$: instável
(fuga térmica)}
\resolucao{(a) \textbf{Coeficiente positivo} (cobre, manganina, a maioria dos
metais). O ramo mais quente tem sua resistência \emph{aumentada}; sua
condutância cai e ele passa a conduzir \emph{menos}. Menos corrente significa
menos aquecimento --- o desvio se corrige sozinho.\\
Trata-se de \textbf{realimentação negativa}: o sistema converge para um novo
equilíbrio, com desequilíbrio limitado.\\
(b) \textbf{Coeficiente negativo} (semicondutores, termistores NTC, carbono). O
ramo mais quente reduz sua resistência, atrai \emph{mais} corrente, aquece
ainda mais --- e o ciclo se acelera. \textbf{Realimentação positiva}, ou fuga
térmica: um ramo acaba conduzindo praticamente toda a corrente e se destrói.\\
(c) \textbf{Condição de estabilidade:} $\alpha>0$ nos ramos, ou --- se
$\alpha<0$ for inevitável --- resistências de balanceamento em série, de valor
suficiente para dominar a variação térmica.\\
\textbf{Exemplos práticos.}
\begin{itemize}
\item \textbf{Condutores em paralelo:} cobre tem $\alpha\approx+0{,}004/\si{\celsius}$
--- naturalmente estáveis.
\item \textbf{Transistores bipolares em paralelo:} o coeficiente da junção é
\emph{negativo}, e por isso cada emissor recebe um resistor de balanceamento.
Sem ele, um transistor assume toda a corrente e queima.
\item \textbf{Diodos e LEDs em paralelo:} mesmo problema --- daí a exigência de
resistor individual por ramo, e não um único resistor comum.
\item \textbf{Resistores de precisão para divisão:} usa-se manganina, que
combina $\alpha$ pequeno \emph{e} positivo.
\end{itemize}}
\stepcounter{questao}
\resposta{(b) $U=\SI{120}{\volt}$; $\SI{12}{\ampere}$ no resistor e
$\SI{6}{\ampere}$ injetados pelo ramo ativo; (c) $g_c=\SI{0.1}{\ensuremath{\mathrm{S}}}$}
\resolucao{(a) O divisor pressupõe que \emph{cada} ramo obedeça $I=UG$ com
$G>0$ fixo. A fonte dependente também é proporcional a $U$, mas com sinal que
pode \emph{retirar} corrente do nó --- ela funciona como condutância
\textbf{negativa} de valor $-g$. A soma $\sum G$ deixa de ser a soma de
condutâncias físicas.\\
(b) LKC no nó, com a fonte dependente \textbf{injetando} $gU$:
\[
\frac{U}{10}=6+gU
\;\Longrightarrow\;
U(0{,}1-0{,}05)=6
\;\Longrightarrow\;
U=\frac{6}{0{,}05}=\SI{120}{\volt},
\]
\[
I_R=\frac{120}{10}=\SI{12}{\ampere};
\qquad
I_{dep}=0{,}05(120)=\SI{6}{\ampere}.
\]
\textbf{Verificação da LKC:} entram $6$ (fonte externa) mais $6$ (fonte
dependente), totalizando $\SI{12}{\ampere}$ --- exatamente o que sai pelo
resistor \checkmark.\\
\textbf{Observe o absurdo aparente:} o resistor conduz $\SI{12}{\ampere}$, o
\textbf{dobro} da corrente injetada externamente. Nenhum divisor passivo
produziria isso, pois nele vale sempre $I_k\le I_T$.\\
(c) A condutância líquida é $\dfrac1{10}-g$, que se anula em
\[
g_c=\SI{0.1}{\ensuremath{\mathrm{S}}}.
\]
Nesse ponto o nó "flutua" e não há solução finita; acima dele, a resistência
equivalente torna-se negativa e o circuito é instável.\\
\textbf{Lição:} elementos ativos podem produzir correntes de ramo
\emph{maiores} que a corrente total injetada, ou de sinal invertido --- ambos
impossíveis em rede puramente passiva. O divisor de corrente é uma ferramenta
de redes passivas.}
\stepcounter{questao}
\resposta{$P=I_T^{2}R_{eq}$, e $R_{eq}$ escala com o fator adotado; escalas de
$0{,}6/0{,}3/\SI{0.2}{\ohm}$ e $6/3/\SI{2}{\ohm}$}
\resolucao{(a) Multiplicando todas as resistências por $\lambda$, as
\emph{frações} $G_k/\sum G$ não mudam --- a divisão é preservada. Mas
\[
R_{eq}=\lambda R_{eq,0}
\;\Longrightarrow\;
P=I_T^{2}R_{eq}=\lambda\,I_T^{2}R_{eq,0}.
\]
A perda é \textbf{diretamente proporcional} à escala. (É a mesma propriedade de
homogeneidade de grau 1 já usada para tolerâncias em associações.)\\
(b) Com $R=6/3/\SI{2}{\ohm}$: $R_{eq}=\SI{1}{\ohm}$ e
$P=144(1)=\SI{144}{\watt}$.\\
Com $R=0{,}6/0{,}3/\SI{0.2}{\ohm}$: $R_{eq}=\SI{0.1}{\ohm}$ e
$P=\SI{14.4}{\watt}$ --- dez vezes menos, com \emph{exatamente} a mesma divisão
$2{:}4{:}6$ ampères.\\
(c) \textbf{Critério.} A escala é limitada por baixo e por cima:
\begin{itemize}
\item \textbf{Limite inferior:} resistências pequenas demais tornam-se
comparáveis às parasitas (contatos, trilhas, cabos), que passam a governar a
divisão. Regra prática: mantenha $R\ge20$ vezes a maior parasita estimada.
\item \textbf{Limite superior:} a perda $\lambda I_T^{2}R_{eq,0}$ e a queda de
tensão crescem linearmente, exigindo componentes maiores e comprometendo o
rendimento.
\end{itemize}
\textbf{Escolha:} a menor escala que ainda domine as parasitas com folga. Com
contatos de $\SI{1}{\milli\ohm}$, a escala de $0{,}6/0{,}3/\SI{0.2}{\ohm}$ é
segura (o menor ramo é $200\times$ a parasita); reduzir mais dez vezes já
deixaria a divisão à mercê da qualidade das conexões.}

\gabarito[4.2 Divisor de corrente]

\section{Potência Elétrica e Máxima Transferência}
\stepcounter{questao}
\resposta{(c)}
\resolucao{$P = R\,I^2 = 7\cdot 3^2 = \SI{63}{\watt}$.}
\stepcounter{questao}
\resposta{(c)}
\resolucao{A potência na carga é máxima quando $R_L = R_{\text{Th}}$ (casamento de
impedância). Cargas muito pequenas recebem alta corrente mas baixa tensão; cargas
muito grandes, o contrário. O equilíbrio ótimo está na igualdade.}
\stepcounter{questao}
\resposta{(c)}
\resolucao{$P = U\,I = 12\cdot3 = \SI{36}{\watt}$.}
\stepcounter{questao}
\resposta{(a) \SI{4}{\ohm}; (b) \SI{9}{\watt}}
\resolucao{(a) $R_L = R_{\text{Th}} = \SI{4}{\ohm}$.\\
(b) $P_{\max} = \dfrac{V_{\text{Th}}^{\,2}}{4R_{\text{Th}}}
= \dfrac{12^2}{4\cdot4} = \dfrac{144}{16} = \SI{9}{\watt}$.\\
\textbf{Confira:} com $R_L = R_{\text{Th}}$, a corrente é
$I = \dfrac{12}{4+4} = \SI{1.5}{\ampere}$ e
$P_L = R_L I^2 = 4\cdot1{,}5^2 = \SI{9}{\watt}$.}
\stepcounter{questao}
\resposta{$R_L = \SI{6}{\ohm}$; $P_{\max} = \SI{13.5}{\watt}$}
\resolucao{$R_L = R_{\text{int}} = \SI{6}{\ohm}$, e
\[
P_{\max}=\frac{V^2}{4R_{\text{int}}}=\frac{18^2}{4\cdot6}=\frac{324}{24}
=\SI{13.5}{\watt}.
\]}
\stepcounter{questao}
\resposta{(c)}
\resolucao{Para uma fonte \emph{absorver} potência, a corrente deve entrar pelo
seu terminal positivo (convenção do receptor). Aqui a fonte de \SI{90}{\volt}
impõe a corrente no sentido horário, que de fato entra pelo $+$ da fonte de
\SI{60}{\volt} --- por isso esta última é carregada.
A corrente vale $I = \dfrac{90-60}{10} = \SI{3}{\ampere}$, e a fonte de
\SI{60}{\volt} absorve $P = 60\cdot3 = \SI{180}{\watt}$. É exatamente o que ocorre
ao carregar uma bateria: uma fonte de maior tensão força corrente "para dentro"
dela.}
\stepcounter{questao}
\resposta{(a) \SI{16}{\watt}; (b) máximo é \SI{18}{\watt} em $R_L=\SI{2}{\ohm}$}
\resolucao{(a) $I = \dfrac{12}{2+4} = \SI{2}{\ampere}$;
$P_L = R_L I^2 = 4\cdot4 = \SI{16}{\watt}$.\\
(b) O máximo ocorreria em $R_L = R_{\text{int}} = \SI{2}{\ohm}$:
$P_{\max} = \dfrac{12^2}{4\cdot2} = \SI{18}{\watt}$. A carga de \SI{4}{\ohm}
dissipa \SI{16}{\watt}, um pouco abaixo do máximo --- coerente com a curva suave
de $P_L$ em torno do ponto ótimo.}
\stepcounter{questao}
\resposta{$P(5)\approx\SI{8.9}{\watt}$, $P(10)=\SI{10}{\watt}$,
$P(20)\approx\SI{8.9}{\watt}$; máximo em \SI{10}{\ohm}}
\resolucao{A corrente é $I = \dfrac{V_{\text{Th}}}{R_{\text{Th}}+R_L}$ e a potência
na carga, $P_L = R_L I^2 = \dfrac{V_{\text{Th}}^{\,2}\,R_L}{(R_{\text{Th}}+R_L)^2}$:
\[
P(5)=\frac{400\cdot5}{15^2}=\frac{2000}{225}\approx\SI{8.9}{\watt},\quad
P(10)=\frac{400\cdot10}{20^2}=\SI{10}{\watt},\quad
P(20)=\frac{400\cdot20}{30^2}=\frac{8000}{900}\approx\SI{8.9}{\watt}.
\]
O valor máximo é $P(10) = \SI{10}{\watt}$, confirmando
$R_L = R_{\text{Th}} = \SI{10}{\ohm}$. Note a \emph{simetria}: cargas de
\SI{5}{} e \SI{20}{\ohm} (uma metade, outra o dobro de $R_{\text{Th}}$) dão a mesma
potência. A curva $P_L \times R_L$ tem um máximo suave em $R_L = R_{\text{Th}}$.}
\stepcounter{questao}
\resposta{Ver comentário}
\resolucao{Com $R_L = R_{\text{Th}}$, a mesma corrente percorre as duas resistências e,
como elas têm o mesmo valor, a potência dissipada na carga é igual à potência
dissipada \emph{internamente} em $R_{\text{Th}}$.
O rendimento é
\[
\eta=\frac{P_{\text{carga}}}{P_{\text{total}}}
=\frac{R_L}{R_L+R_{\text{Th}}}=\frac{1}{2}=\SI{50}{\percent}.
\]
Metade da energia vira calor na fonte --- aceitável em sinais fracos (áudio,
antenas, onde importa extrair o \emph{máximo} sinal), mas inaceitável na rede
elétrica, onde se quer \emph{eficiência}. Por isso, em sistemas de potência,
opera-se com $R_L \gg R_{\text{Th}}$ (rendimento próximo de \SI{100}{\percent}),
sacrificando potência máxima em favor de baixo desperdício.}
\stepcounter{questao}
\resposta{(a) $I=\SI{2}{\ampere}$, $P_L=\SI{32}{\watt}$, $\eta=\SI{67}{\percent}$;
(b) e (c) ver comentário}
\resolucao{\textbf{(a)} Para $R_L=\SI{8}{\ohm}$:
\[
I=\frac{24}{4+8}=\SI{2}{\ampere},\quad
P_L=R_LI^2=8\cdot4=\SI{32}{\watt},\quad
\eta=\frac{R_L}{R_{\text{Th}}+R_L}=\frac{8}{12}\approx\SI{67}{\percent}.\ \checkmark
\]
\textbf{(b) A simetria.} Observe que $R_L=\SI{2}{\ohm}$ e $R_L=\SI{8}{\ohm}$ dão a
\emph{mesma} potência (\SI{32}{\watt}); o mesmo ocorre com \SI{1}{} e
\SI{16}{\ohm} (\SI{23}{\watt}). Isso porque cargas que são \emph{recíprocas} em
relação a $R_{\text{Th}}$ --- isto é, $R_L$ e $R_{\text{Th}}^2/R_L$ --- produzem
potências idênticas:
\[
P_L=\frac{V^2R_L}{(R_{\text{Th}}+R_L)^2}
\quad\text{é invariante sob}\quad
R_L \leftrightarrow \frac{R_{\text{Th}}^{2}}{R_L}.
\]
Verifique: $2$ e $16/2=8$ \checkmark; $1$ e $16/1=16$ \checkmark. A curva é
simétrica em escala \emph{logarítmica}, com máximo em $R_L=R_{\text{Th}}$.

\textbf{(c) O conflito central de projeto.}
\begin{itemize}[leftmargin=1.6em,itemsep=1pt,topsep=2pt]
\item \textbf{Máxima potência} $\Rightarrow R_L = R_{\text{Th}} = \SI{4}{\ohm}$,
      mas com apenas \SI{50}{\percent} de rendimento --- metade da energia vira
      calor na fonte.
\item \textbf{Máxima eficiência} $\Rightarrow R_L$ o \emph{maior} possível
      ($\eta \to \SI{100}{\percent}$), mas a potência entregue tende a zero.
\end{itemize}
\textbf{Não existe ponto que otimize ambos} --- são objetivos conflitantes.
A escolha depende do que é escasso:
\begin{itemize}[leftmargin=1.6em,itemsep=1pt,topsep=2pt]
\item em \emph{telecomunicações} e áudio, o sinal é fraco e a energia é barata:
      casa-se a impedância para extrair o máximo sinal;
\item em \emph{sistemas de potência}, a energia é o produto: opera-se com
      $R_L \gg R_{\text{Th}}$, com rendimento acima de \SI{95}{\percent}.
\end{itemize}
\emph{Saber qual critério aplicar é mais importante que saber calcular ambos.}}
\stepcounter{questao}
\resposta{$R_L = \SI{2}{\ohm}$; $P = \SI{72}{\watt}$; $I = \SI{6}{\ampere}$
(no limite)}
\resolucao{\textbf{Condição de máxima transferência:}
$R_L = R_{\text{int}} = \SI{2}{\ohm}$.\\
\textbf{Corrente resultante:}
\[
I = \frac{24}{R_{\text{int}}+R_L}=\frac{24}{4}=\SI{6}{\ampere}.
\]
A restrição é atendida \emph{exatamente no limite} --- não há margem.\\
\textbf{Potência na carga:}
\[
P_L = R_L I^2 = 2\cdot6^2 = \SI{72}{\watt}
\quad\text{(confere com } P_{\max}=\tfrac{24^2}{4\cdot2}=\SI{72}{\watt}).
\]
\textbf{Análise crítica do projeto:} a fonte dissipa internamente os mesmos
\SI{72}{\watt} (rendimento de \SI{50}{\percent}), totalizando \SI{144}{\watt}
consumidos. Como a corrente está no limite especificado, qualquer tolerância
para baixo em $R_L$ (por aquecimento, por exemplo) violaria a restrição.
\textbf{Recomendação de engenharia:} adotar $R_L$ ligeiramente \emph{acima} de
\SI{2}{\ohm} --- por exemplo \SI{2.2}{\ohm} --- sacrifica menos de
\SI{1}{\percent} da potência ($\SI{71.5}{\watt}$) e traz a corrente para
\SI{5.7}{\ampere}, criando margem de segurança. \emph{Projetar no limite exato é
um erro clássico de quem só sabe calcular; o profissional experiente sempre
reserva margem.}}
\stepcounter{questao}
\resposta{$R_L=\SI{6}{\ohm}$; $P_{\max}=\SI{24}{\watt}$}
\resolucao{A condição é $R_L=R_{Th}=\SI{6}{\ohm}$, e então
\[
P_{\max}=\frac{V_{Th}^{2}}{4R_{Th}}=\frac{576}{24}=\SI{24}{\watt}.
\]
\textbf{Guarde a fórmula do máximo:} ela dispensa calcular corrente e tensão, e
mostra que a potência máxima disponível de uma fonte depende \emph{apenas} de
$V_{Th}$ e $R_{Th}$ --- é uma característica da fonte, não da carga.}
\stepcounter{questao}
\resposta{$I=\SI{2}{\ampere}$; $U_L=\SI{12}{\volt}$}
\resolucao{Com $R_L=R_{Th}$, a resistência total é $2R_{Th}$:
\[
I=\frac{V_{Th}}{2R_{Th}}=\frac{24}{12}=\SI{2}{\ampere};
\qquad
U_L=R_LI=6(2)=\SI{12}{\volt}=\frac{V_{Th}}{2}.
\]
\textbf{Regra do ponto ótimo:} a tensão na carga é sempre \emph{metade} da
tensão em vazio, e a corrente é metade da corrente de curto-circuito. Esses dois
fatos identificam o ponto ótimo em qualquer medição de bancada.}
\stepcounter{questao}
\resposta{$V_{Th}=\SI{40}{\volt}$}
\resolucao{De $P_{\max}=\dfrac{V_{Th}^{2}}{4R_{Th}}$:
\[
V_{Th}=\sqrt{4R_{Th}P_{\max}}=\sqrt{4(8)(50)}=\sqrt{1600}=\SI{40}{\volt}.
\]
\textbf{Conferência:} $\dfrac{40^{2}}{4(8)}=\dfrac{1600}{32}=\SI{50}{\watt}$
\checkmark.}
\stepcounter{questao}
\resposta{$R_{Th}=\SI{6}{\ohm}$}
\resolucao{
\[
R_{Th}=\frac{V_{Th}^{2}}{4P_{\max}}=\frac{576}{96}=\SI{6}{\ohm}.
\]
\textbf{Método experimental:} medindo $V_{oc}$ e a potência máxima entregue a
uma carga ajustável, obtém-se $R_{Th}$ sem nunca curto-circuitar a fonte ---
alternativa útil quando o curto seria destrutivo.}
\stepcounter{questao}
\resposta{(c)}
\resolucao{Com $x=R_L/R_{Th}=2$:
\[
\frac{P_L}{P_{\max}}=\frac{4x}{(1+x)^{2}}=\frac{8}{9}=88{,}9\%.
\]
\textbf{Fato importante:} o máximo é \textbf{achatado}. Errar a carga por um
fator $2$ custa apenas $11\%$ da potência --- o que torna o casamento muito
mais tolerante do que se costuma supor. A quantificação completa dessa
tolerância está no bloco de aprofundamento.}
\stepcounter{questao}
\resposta{$\SI{21.33}{\watt}$ nos dois casos}
\resolucao{
\[
R_L=\SI{3}{\ohm}:\quad I=\frac{24}{9}=\SI{2.667}{\ampere},\quad
P=3(2{,}667)^{2}=\SI{21.33}{\watt};
\]
\[
R_L=\SI{12}{\ohm}:\quad I=\frac{24}{18}=\SI{1.333}{\ampere},\quad
P=12(1{,}333)^{2}=\SI{21.33}{\watt}.
\]
\textbf{Potências idênticas} --- e note que $3\times12=36=R_{Th}^{2}$. Cargas
que são \emph{recíprocas} em relação a $R_{Th}$ recebem exatamente a mesma
potência. A demonstração está no bloco de desafio.\\
\textbf{Mas os dois casos não são equivalentes:} com $\SI{3}{\ohm}$ a corrente
é o dobro e o rendimento é $33\%$; com $\SI{12}{\ohm}$, o rendimento é
$67\%$. Mesma potência entregue, consumo da fonte muito diferente.}
\stepcounter{questao}
\resposta{$\SI{48}{\watt}$, $\SI{24}{\watt}$ e $\SI{24}{\watt}$}
\resolucao{Com $I=\SI{2}{\ampere}$:
\[
P_{tot}=V_{Th}I=24(2)=\SI{48}{\watt};
\qquad
P_L=R_LI^{2}=\SI{24}{\watt};
\qquad
P_{int}=R_{Th}I^{2}=\SI{24}{\watt}.
\]
\textbf{Metade da energia é perdida dentro da fonte} --- consequência direta de
$R_L=R_{Th}$ e da corrente comum. É o preço do casamento, e a razão pela qual
ele nunca é usado em eletrotécnica de potência.}
\stepcounter{questao}
\resposta{$\SI{2}{\ohm}$ e $\SI{8}{\ohm}$}
\resolucao{
\[
R_{Th}=\frac{V_{Th}^{2}}{4P_{\max}}
\Rightarrow
\frac{144}{72}=\SI{2}{\ohm};
\qquad
\frac{576}{72}=\SI{8}{\ohm}.
\]
\textbf{Leitura:} para a mesma potência disponível, dobrar a tensão permite
\emph{quadruplicar} a resistência interna. Fontes de tensão mais alta toleram
resistência interna maior --- princípio que fundamenta a transmissão de energia
em alta tensão.\\
\textbf{Diferença prática:} a primeira entrega $\SI{18}{\watt}$ a
$\SI{3}{\ampere}$; a segunda, a $\SI{1.5}{\ampere}$. Correntes menores
significam condutores mais finos.}
\stepcounter{questao}
\resposta{$\SI{27.8}{\watt}$ e $\SI{10}{\watt}$}
\resolucao{
\[
R_{Th}=\SI{2}{\ohm}:\quad I=\frac{20}{12}=\SI{1.667}{\ampere},\quad
P_L=10(1{,}667)^{2}=\SI{27.8}{\watt};
\]
\[
R_{Th}=\SI{10}{\ohm}:\quad I=\SI{1}{\ampere},\quad P_L=\SI{10}{\watt}.
\]
\textbf{Aparente paradoxo:} a segunda fonte está \emph{casada} com a carga
($R_L=R_{Th}$) e mesmo assim entrega bem menos que a primeira, descasada.\\
\textbf{Resolução:} o casamento maximiza a potência quando $R_L$ é o parâmetro
livre e a fonte é dada. Aqui a carga é fixa e a \emph{fonte} varia --- e nesse
caso quanto \emph{menor} $R_{Th}$, melhor. São problemas de otimização
diferentes, com respostas diferentes. O ponto é retomado no bloco de
aprofundamento.}
\stepcounter{questao}
\resposta{$\SI{120}{\watt\hour}$ e $\SI{240}{\watt\hour}$}
\resolucao{
\[
E_L=P_Lt=24(5)=\SI{120}{\watt\hour}=\SI{0.12}{\kilo\watt\hour};
\]
\[
E_{tot}=48(5)=\SI{240}{\watt\hour}.
\]
Metade da energia --- $\SI{120}{\watt\hour}$ --- foi dissipada internamente e
perdida.\\
\textbf{Consequência para sistemas alimentados por bateria:} operar casado
\emph{reduz pela metade} a autonomia. Em um rádio portátil, isso significa
trocar tempo de operação por volume de som --- decisão que às vezes vale a pena,
mas que precisa ser consciente.}
\stepcounter{questao}
\resposta{$P_L=\SI{18}{\watt}$; $\eta=75\%$; $75\%$ do máximo}
\resolucao{
\[
I=\frac{24}{24}=\SI{1}{\ampere};
\qquad
P_L=18(1)^{2}=\SI{18}{\watt};
\qquad
\eta=\frac{18}{24}=75\%.
\]
Comparando com $P_{\max}=\SI{24}{\watt}$: entrega-se $75\%$ da potência
disponível, com rendimento de $75\%$ em vez de $50\%$.\\
\textbf{Compromisso quantificado:} abrir mão de $25\%$ da potência comprou
$25$ pontos percentuais de rendimento. Em quase toda aplicação prática esse é
um bom negócio --- e é por isso que sistemas reais operam com
$R_L\gg R_{Th}$.}
\stepcounter{questao}
\resposta{$R_{Th}=\SI{6}{\ohm}$; $P_{\max}=\SI{37.5}{\watt}$}
\resolucao{
\[
R_{Th}=\frac{V_{oc}}{I_{sc}}=\frac{30}{5}=\SI{6}{\ohm};
\qquad
P_{\max}=\frac{30^{2}}{4(6)}=\SI{37.5}{\watt}.
\]
\textbf{Forma alternativa e elegante:}
\[
P_{\max}=\frac{V_{oc}I_{sc}}{4}=\frac{30(5)}{4}=\SI{37.5}{\watt}.
\]
A potência máxima é um quarto do produto dos dois valores extremos --- e faz
sentido: no ponto ótimo, tensão e corrente valem, cada uma, metade do seu
extremo.}
\stepcounter{questao}
\resposta{Fonte: $\SI{72}{\watt}$; carga: $\SI{18}{\watt}$; interna:
$\SI{54}{\watt}$}
\resolucao{
\[
I=\frac{24}{8}=\SI{3}{\ampere};
\qquad
P_{fonte}=24(3)=\SI{72}{\watt};
\]
\[
P_L=2(9)=\SI{18}{\watt};
\qquad
P_{int}=6(9)=\SI{54}{\watt};
\qquad 18+54=72\ \checkmark
\]
\textbf{Rendimento de apenas $25\%$} --- três quartos da energia aquecem a
fonte. Compare com a carga de $\SI{18}{\ohm}$ da questão anterior, que entrega
a \emph{mesma} $\SI{18}{\watt}$ com rendimento de $75\%$. Duas cargas
recíprocas, mesma potência, rendimentos opostos.}
\stepcounter{questao}
\resposta{(a) $\SI{18}{\watt}$ nos dois; (b) $\SI{15.36}{\watt}$ nos dois;
(c) $R_{L1}R_{L2}=R_{Th}^{2}$}
\resolucao{(a) $R_L=2$: $I=\SI{3}{\ampere}$, $P=\SI{18}{\watt}$.
$R_L=18$: $I=\SI{1}{\ampere}$, $P=\SI{18}{\watt}$.\\
(b) $R_L=1{,}5$: $I=24/7{,}5=\SI{3.2}{\ampere}$,
$P=1{,}5(10{,}24)=\SI{15.36}{\watt}$.
$R_L=24$: $I=\SI{0.8}{\ampere}$, $P=24(0{,}64)=\SI{15.36}{\watt}$.\\
(c) Em ambos os casos, $R_{L1}R_{L2}=R_{Th}^{2}=36$:
\[
2\times18=36;\qquad 1{,}5\times24=36.
\]
Em termos de $x=R_L/R_{Th}$, os pares são $x$ e $1/x$ --- \textbf{recíprocos}.\\
\textbf{Consequência para medição:} se uma varredura de carga encontra a mesma
potência em dois valores $R_{L1}$ e $R_{L2}$, então
$R_{Th}=\sqrt{R_{L1}R_{L2}}$ --- a \emph{média geométrica}. É um método de
determinar $R_{Th}$ sem localizar o máximo com precisão, o que é vantajoso
justamente porque o máximo é achatado.}
\stepcounter{questao}
\resposta{$0{,}52\le x\le1{,}93$ --- razão de quase $4{:}1$}
\resolucao{(a) Impondo $\dfrac{4x}{(1+x)^{2}}=0{,}9$:
\[
4x=0{,}9(1+x)^{2}
\;\Longrightarrow\;
0{,}9x^{2}-2{,}2x+0{,}9=0,
\]
\[
x=\frac{2{,}2\pm\sqrt{4{,}84-3{,}24}}{1{,}8}
=\frac{2{,}2\pm1{,}265}{1{,}8}
\;\Longrightarrow\;
x=0{,}520\ \text{ou}\ 1{,}925.
\]
Note que $0{,}520\times1{,}925=1{,}000$ --- os limites são recíprocos, como a
simetria exige.\\
(b) \textbf{A faixa cobre uma razão de $3{,}7{:}1$.} Com
$R_{Th}=\SI{6}{\ohm}$, qualquer carga entre $\SI{3.1}{\ohm}$ e
$\SI{11.6}{\ohm}$ entrega ao menos $90\%$ do máximo.\\
\textbf{Implicação prática:} não vale a pena esforço para casar cargas com
precisão melhor que uns $20\%$. O máximo é achatado porque a função tem
derivada nula ali --- perto do topo, a curva é essencialmente uma parábola
suave. Esse é o motivo pelo qual "casamento aproximado" costuma bastar.}
\stepcounter{questao}
\resposta{$\eta$: $33$, $50$, $67$, $80$ e $90\%$; $P/P_{\max}$: $89$, $100$,
$89$, $64$ e $36\%$}
\resolucao{Com $\eta=\dfrac{x}{1+x}$ e
$\dfrac{P_L}{P_{\max}}=\dfrac{4x}{(1+x)^{2}}$:
\begin{center}
\begin{tabular}{ccc}
\hline
$x=R_L/R_{Th}$ & $\eta$ & $P_L/P_{\max}$\\
\hline
$0{,}5$ & $33{,}3\%$ & $88{,}9\%$\\
$1$ & $50{,}0\%$ & $100\%$\\
$2$ & $66{,}7\%$ & $88{,}9\%$\\
$4$ & $80{,}0\%$ & $64{,}0\%$\\
$9$ & $90{,}0\%$ & $36{,}0\%$\\
\hline
\end{tabular}
\end{center}
\textbf{Três leituras.}\\
\emph{(i)} $x=0{,}5$ e $x=2$ dão a mesma potência ($88{,}9\%$) mas rendimentos
opostos ($33\%$ e $67\%$). \textbf{Entre dois pares recíprocos, escolha sempre
o maior} --- mesma entrega, metade da perda.\\
\emph{(ii)} As duas curvas puxam em sentidos contrários: $\eta$ cresce
monotonicamente com $x$, enquanto $P_L/P_{\max}$ tem máximo em $x=1$ e decresce.
Não existe ponto que otimize ambos.\\
\emph{(iii)} A região $2\le x\le4$ é o compromisso usual: ainda se entrega de
$64\%$ a $89\%$ da potência disponível, com rendimento entre $67\%$ e
$80\%$.}
\stepcounter{questao}
\resposta{(b) $R_{Th}\to0$; a potência cresce monotonicamente conforme
$R_{Th}$ diminui}
\resolucao{(a)
\begin{center}
\begin{tabular}{cc}
\hline
$R_{Th}$ & $P_L$\\
\hline
$\SI{16}{\ohm}$ & $\SI{8}{\watt}$\\
$\SI{8}{\ohm}$ & $\SI{18}{\watt}$\\
$\SI{4}{\ohm}$ & $\SI{32}{\watt}$\\
$\SI{2}{\ohm}$ & $\SI{46.1}{\watt}$\\
$\SI{0}{\ohm}$ & $\SI{72}{\watt}$\\
\hline
\end{tabular}
\end{center}
(b) A potência cresce sem parar conforme $R_{Th}$ diminui: o ótimo é
$R_{Th}\to0$, no \emph{contorno} do domínio, não em um ponto interior.\\
(c) \textbf{A condição $R_L=R_{Th}$ não é simétrica.} Ela resolve o problema
"dada a fonte, escolha a carga". Aqui o problema é o inverso --- "dada a carga,
escolha a fonte" --- e a função
\[
P_L=\frac{V_{Th}^{2}R_L}{(R_L+R_{Th})^{2}}
\]
é \textbf{monotonicamente decrescente} em $R_{Th}$: sua derivada,
$-\dfrac{2V_{Th}^{2}R_L}{(R_L+R_{Th})^{3}}$, nunca se anula para
$R_{Th}\ge0$.\\
\textbf{Regra:} o casamento de impedâncias só faz sentido quando a fonte é
\emph{dada e imutável} --- uma antena, um transdutor, um sensor. Se a fonte pode
ser projetada, sempre convém $R_{Th}$ o menor possível.}
\stepcounter{questao}
\resposta{(a) $\SI{72}{\watt}$ com $R_L=\SI{0.5}{\ohm}$; (b)
$\SI{12}{\ampere}$; (c) $\SI{46.1}{\watt}$ com $\eta=80\%$}
\resolucao{(a) $P_{\max}=\dfrac{144}{4(0{,}5)}=\SI{72}{\watt}$ com
$R_L=\SI{0.5}{\ohm}$.\\
(b) $I=\dfrac{12}{1}=\SI{12}{\ampere}$, com $\SI{72}{\watt}$ dissipados
\emph{dentro} da bateria. Uma bateria pequena não suporta isso por muito tempo:
o eletrólito aquece, $r$ cresce, e a potência disponível cai durante a própria
operação.\\
(c) $I=\dfrac{12}{2{,}5}=\SI{4.8}{\ampere}$;
$P_L=2(4{,}8)^{2}=\SI{46.1}{\watt}$; $\eta=2/2{,}5=80\%$.\\
\textbf{Comparação decisiva:} entrega-se $64\%$ da potência máxima, mas a perda
interna cai de $\SI{72}{\watt}$ para $\SI{11.5}{\watt}$ --- \emph{seis vezes}
menos aquecimento.\\
\textbf{Além disso:} o modelo de $r$ constante já é otimista a
$\SI{12}{\ampere}$. Baterias reais apresentam $r$ crescente com a corrente e
com o estado de descarga, o que torna a operação casada ainda menos atraente do
que o cálculo sugere.}
\stepcounter{questao}
\resposta{(a) $V_{Th}=\SI{9.82}{\volt}$, $R_{Th}=\SI{4.36}{\ohm}$;
(b) $R_L=\SI{4.36}{\ohm}$ e $P_{\max}=\SI{5.53}{\watt}$; (c) contra
$\SI{81}{\watt}$ disponíveis na fonte}
\resolucao{(a) Com a carga removida, a rede é um divisor de
$4+12=\SI{16}{\ohm}$ e $\SI{6}{\ohm}$:
\[
V_{Th}=36\cdot\frac{6}{22}=\SI{9.82}{\volt};
\qquad
R_{Th}=\frac{16\cdot6}{22}=\SI{4.36}{\ohm}.
\]
(b) $R_L=\SI{4.36}{\ohm}$ e
\[
P_{\max}=\frac{(9{,}82)^{2}}{4(4{,}36)}=\SI{5.53}{\watt}.
\]
(c) A fonte original disponibilizaria
$\dfrac{36^{2}}{4(4)}=\SI{81}{\watt}$.\\
\textbf{A rede intermediária destruiu $93\%$ da potência disponível.} O
resistor em série de $\SI{12}{\ohm}$ derruba a tensão e o de $\SI{6}{\ohm}$ desvia
corrente --- ambos dissipando.\\
\textbf{Conclusão:} redes resistivas \emph{não} casam impedâncias sem perda.
Elas alteram $R_{Th}$, sim, mas sempre ao custo de reduzir $V_{Th}$ e,
com ele, a potência disponível. Casamento sem perda exige elementos reativos
(transformador, rede LC) --- possíveis apenas em corrente alternada.}
\stepcounter{questao}
\resposta{(a) $\SI{12}{\ohm}$ ($\SI{21.3}{\watt}$); (b) $\SI{40}{\ohm}$
($\eta=87\%$)}
\resolucao{(a)
\[
R_L=12:\ I=\frac{24}{18}=1{,}333;\ P_L=12(1{,}778)=\SI{21.3}{\watt};
\]
\[
R_L=40:\ I=\frac{24}{46}=0{,}522;\ P_L=40(0{,}272)=\SI{10.9}{\watt}.
\]
A de $\SI{12}{\ohm}$ é a mais próxima do ótimo ($x=2$) e entrega quase o dobro.\\
(b) $\eta=\dfrac{R_L}{R_L+6}$: $67\%$ contra $87\%$. A de $\SI{40}{\ohm}$
vence.\\
(c) \textbf{O critério depende do que é escasso.}
\begin{itemize}
\item Se a fonte é abundante (rede elétrica) e o custo está na energia, escolha
o \emph{rendimento}.
\item Se a fonte é limitada (bateria, sensor, antena) e o que importa é extrair
o máximo de um sinal fraco, escolha a \emph{potência}.
\item Se a limitação é térmica na própria fonte, escolha rendimento --- a de
$\SI{40}{\ohm}$ dissipa apenas $\SI{1.6}{\watt}$ internamente, contra
$\SI{10.7}{\watt}$ da outra.
\end{itemize}
Sem saber o que é escasso, a pergunta "qual é melhor" não tem resposta.}
\stepcounter{questao}
\resposta{(a) não --- exigiria $\SI{2}{\ampere}$; (b)
$R_L=\SI{10}{\ohm}$; (c) $\SI{22.5}{\watt}$ contra $\SI{24}{\watt}$}
\resolucao{(a) O ótimo teórico exige $I=V_{Th}/(2R_{Th})=\SI{2}{\ampere}$,
acima do limite de $\SI{1.5}{\ampere}$. \textbf{O ponto ótimo é inatingível.}\\
(b) Como $P_L=R_LI^{2}$ cresce quando $R_L$ se aproxima de $R_{Th}$ por cima, o
melhor admissível é a \emph{menor} carga que ainda respeite o limite:
\[
I=\frac{24}{6+R_L}\le1{,}5
\;\Longrightarrow\;
6+R_L\ge16
\;\Longrightarrow\;
R_L\ge\SI{10}{\ohm}.
\]
Adotando $R_L=\SI{10}{\ohm}$: $I=\SI{1.5}{\ampere}$ e
$P_L=10(2{,}25)=\SI{22.5}{\watt}$.\\
(c) Isso é $93{,}8\%$ do máximo irrestrito --- perda de apenas $6\%$.\\
\textbf{Duas observações.} O ótimo restrito fica na \emph{fronteira} do domínio
admissível, não onde a derivada se anula: é a marca de um problema com
restrição ativa. E o custo da restrição é pequeno justamente porque o máximo é
achatado --- a mesma propriedade que torna o casamento tolerante torna as
restrições baratas.}
\stepcounter{questao}
\resposta{$\eta=80\%$ (com $x=4$) ou $\eta=20\%$ (com $x=0{,}25$)}
\resolucao{Impondo $\dfrac{4x}{(1+x)^{2}}=0{,}64$:
\[
0{,}64x^{2}-2{,}72x+0{,}64=0
\;\Longrightarrow\;
x=4\ \text{ou}\ x=0{,}25,
\]
recíprocos, como esperado.\\
\textbf{Rendimentos:}
\[
x=4:\ \eta=\frac{4}{5}=80\%;
\qquad
x=0{,}25:\ \eta=\frac{0{,}25}{1{,}25}=20\%.
\]
\textbf{A mesma potência, dois rendimentos opostos} --- e a soma é
$100\%$, o que não é coincidência: se $x$ e $1/x$ são recíprocos, então
$\dfrac{x}{1+x}+\dfrac{1/x}{1+1/x}=\dfrac{x}{1+x}+\dfrac{1}{1+x}=1$.\\
\textbf{Decisão óbvia:} escolha sempre o ramo $x>1$. A escolha $x=0{,}25$
entregaria a mesma potência à carga dissipando \emph{quatro vezes} mais na
fonte.}
\stepcounter{questao}
\resposta{$R_L=R_{Th}$; $P_{\max}=\dfrac{V_{Th}^{2}}{4R_{Th}}$;
$\eta=50\%$}
\resolucao{(a) Com $I=\dfrac{V_{Th}}{R_{Th}+R_L}$:
\[
P_L=R_LI^{2}=\frac{V_{Th}^{2}R_L}{(R_{Th}+R_L)^{2}}.
\]
Derivando pela regra do quociente:
\[
\frac{\mathrm{d}P_L}{\mathrm{d}R_L}
=V_{Th}^{2}\,\frac{(R_{Th}+R_L)^{2}-R_L\cdot2(R_{Th}+R_L)}{(R_{Th}+R_L)^{4}}
=V_{Th}^{2}\,\frac{R_{Th}-R_L}{(R_{Th}+R_L)^{3}}.
\]
Anulando: $R_L=R_{Th}$.\\
(b) \textbf{Pelo sinal da derivada}, que é mais simples que a segunda derivada:
o denominador é sempre positivo, logo o sinal é o de $(R_{Th}-R_L)$ ---
\emph{positivo} para $R_L<R_{Th}$ e \emph{negativo} para $R_L>R_{Th}$. A função
cresce, atinge o topo e decresce: máximo confirmado.\\
\textbf{Confirmação pelos limites:} $P_L\to0$ quando $R_L\to0$ (tensão nula na
carga) e quando $R_L\to\infty$ (corrente nula). Entre dois zeros, uma função
positiva e suave tem máximo interior.\\
(c)
\[
P_{\max}=\frac{V_{Th}^{2}R_{Th}}{(2R_{Th})^{2}}=\frac{V_{Th}^{2}}{4R_{Th}};
\qquad
\eta=\frac{R_L}{R_L+R_{Th}}=\frac{1}{2}.
\]
\textbf{Observação estrutural:} a mesma derivada aparece na potência de um
resistor em série (máxima quando ele iguala o restante) e na sensibilidade do
divisor de tensão. É a mesma função $R/(R+a)^{2}$ em três contextos --- vale
reconhecê-la.}
\stepcounter{questao}
\resposta{$P(x)=P(1/x)$; $\eta(x)+\eta(1/x)=1$;
$R_{Th}=\sqrt{R_{L1}R_{L2}}$}
\resolucao{(a) Em variável adimensional $x=R_L/R_{Th}$:
\[
\frac{P_L}{P_{\max}}=\frac{4x}{(1+x)^{2}}.
\]
Substituindo $x\to1/x$:
\[
\frac{4/x}{(1+1/x)^{2}}
=\frac{4/x}{\dfrac{(x+1)^{2}}{x^{2}}}
=\frac{4x}{(1+x)^{2}}.
\]
Idêntico --- a função é \textbf{invariante} sob $x\to1/x$, e o ponto fixo dessa
inversão é $x=1$, que é justamente o máximo.\\
(b)
\[
\eta(x)+\eta(1/x)=\frac{x}{1+x}+\frac{1/x}{1+1/x}
=\frac{x}{1+x}+\frac{1}{1+x}=1.
\]
Os rendimentos são complementares: $33\%$ e $67\%$; $20\%$ e $80\%$; e assim
por diante.\\
(c) \textbf{Método de medição.} Varie a carga e encontre \emph{dois} valores
$R_{L1}\ne R_{L2}$ que produzam a mesma potência. Como
$\dfrac{R_{L1}}{R_{Th}}=\dfrac{R_{Th}}{R_{L2}}$, segue
\[
R_{Th}=\sqrt{R_{L1}R_{L2}}\quad\text{(média geométrica)}.
\]
\textbf{Vantagem sobre buscar o máximo:} perto do topo a curva é chata, e
localizar o pico com precisão é difícil. Já os dois pontos de igual potência
podem ser tomados bem afastados do máximo, onde a curva é \emph{íngreme} --- e a
determinação fica muito mais precisa. É um caso em que fugir do ponto de
interesse melhora a medição.}
\stepcounter{questao}
\resposta{$x=\dfrac{2-f\pm2\sqrt{1-f}}{f}$; faixas de $3{,}7{:}1$,
$2{,}5{:}1$ e $1{,}5{:}1$}
\resolucao{(a) De $\dfrac{4x}{(1+x)^{2}}=f$:
\[
fx^{2}+(2f-4)x+f=0
\;\Longrightarrow\;
x=\frac{(4-2f)\pm\sqrt{(2f-4)^{2}-4f^{2}}}{2f}
=\frac{2-f\pm2\sqrt{1-f}}{f}.
\]
Como o produto das raízes vale $f/f=1$, elas são sempre \textbf{recíprocas} ---
confirmação independente da simetria.\\
(b)
\begin{center}
\begin{tabular}{cccc}
\hline
$f$ & $x_{\min}$ & $x_{\max}$ & razão\\
\hline
$0{,}90$ & $0{,}520$ & $1{,}925$ & $3{,}70$\\
$0{,}95$ & $0{,}635$ & $1{,}576$ & $2{,}48$\\
$0{,}99$ & $0{,}818$ & $1{,}222$ & $1{,}49$\\
\hline
\end{tabular}
\end{center}
(c) \textbf{A largura encolhe como $\sqrt{1-f}$.} Passar de $90\%$ para
$99\%$ da potência exige apertar a faixa por um fator $2{,}5$; chegar a
$99{,}9\%$ exigiria outro fator $3$.\\
\textbf{Origem:} perto do máximo a derivada é nula, e a curva se comporta como
uma parábola --- a perda cresce com o \emph{quadrado} do desvio relativo.
Consequentemente, um desvio de $10\%$ na carga custa cerca de $0{,}25\%$ da
potência.\\
\textbf{Conclusão de engenharia:} casar impedâncias com precisão melhor que
$20\%$ raramente compensa. O esforço deve ir para reduzir $R_{Th}$, não para
refinar o casamento.}
\stepcounter{questao}
\resposta{$V_{Th}=\SI{20}{\volt}$, $R_{Th}=\SI{4}{\ohm}$;
$I=\SI{2.5}{\ampere}$; $P_{tot}=\SI{50}{\watt}$}
\resolucao{(a) Para que $\SI{4}{\ohm}$ seja a carga ótima, é preciso
$R_{Th}=\SI{4}{\ohm}$. Então
\[
P_{\max}=\frac{V_{Th}^{2}}{4(4)}=25
\;\Longrightarrow\;
V_{Th}=\sqrt{400}=\SI{20}{\volt}.
\]
(b) $I=\dfrac{20}{8}=\SI{2.5}{\ampere}$;
$U_L=\SI{10}{\volt}$; $P_{tot}=20(2{,}5)=\SI{50}{\watt}$.\\
(c) \textbf{A fonte dissipa $\SI{25}{\watt}$ internamente} --- tanto quanto
entrega. Isso significa:
\begin{itemize}
\item O elemento que realiza $R_{Th}=\SI{4}{\ohm}$ precisa ser dimensionado para
$\SI{25}{\watt}$ contínuos (especifique $\SI{50}{\watt}$, com dissipador).
\item A fonte primária deve fornecer $\SI{50}{\watt}$ e
$\SI{2.5}{\ampere}$.
\item O rendimento global é $50\%$ por construção.
\end{itemize}
\textbf{Crítica ao enunciado do projeto.} Deliberadamente inserir
$R_{Th}=\SI{4}{\ohm}$ para "casar" é quase sempre um erro. Se o objetivo é
entregar $\SI{25}{\watt}$ a $\SI{4}{\ohm}$, basta uma fonte de
$\SI{10}{\volt}$ com $R_{Th}$ pequeno: ela entrega os mesmos
$\SI{25}{\watt}$ consumindo pouco mais que isso.\\
O casamento só se justifica quando $R_{Th}$ é \emph{imposto} pela física do
dispositivo --- resistência de saída de um transdutor, impedância característica
de uma linha --- e não quando é escolhido por projeto.}
\stepcounter{questao}
\resposta{$x\ge4$; $R_L=\SI{24}{\ohm}$ com $P_L=\SI{15.36}{\watt}$
($64\%$ do máximo)}
\resolucao{(a) $\eta=\dfrac{x}{1+x}\ge0{,}8
\Rightarrow x\ge0{,}8+0{,}8x
\Rightarrow 0{,}2x\ge0{,}8
\Rightarrow \boxed{x\ge4}$.\\
(b) No domínio admissível $x\ge4$, a função
$\dfrac{4x}{(1+x)^{2}}$ é \emph{decrescente} (pois já se passou do máximo em
$x=1$). O melhor valor está, portanto, na \textbf{fronteira}: $x=4$, isto é
$R_L=\SI{24}{\ohm}$.
\[
I=\frac{24}{30}=\SI{0.8}{\ampere};
\qquad
P_L=24(0{,}64)=\SI{15.36}{\watt};
\qquad
\eta=80\%\ \checkmark
\]
(c) O máximo irrestrito era $\SI{24}{\watt}$ com $\eta=50\%$. A restrição
custou $36\%$ da potência e comprou $30$ pontos de rendimento.\\
\textbf{Estrutura do problema.} Como as duas curvas são monótonas em sentidos
opostos na região $x>1$, \textbf{a restrição de rendimento é sempre ativa} --- o
ótimo cai exatamente sobre ela, nunca no interior. Isso simplifica muito o
projeto: basta resolver a igualdade $\eta=\eta_{\min}$.\\
\textbf{Generalização:} para rendimento mínimo $\eta_0$, o ótimo é
$x=\dfrac{\eta_0}{1-\eta_0}$, e a fração de potência obtida é
$4\eta_0(1-\eta_0)$. Para $\eta_0=0{,}5$ isso dá $100\%$ (o casamento); para
$\eta_0=0{,}9$, apenas $36\%$.}
\stepcounter{questao}
\resposta{Não há máximo interior; $P_L\to\infty$ com
$R_{Th}\to0$ e $R_L\to0$}
\resolucao{(a) Com
$P_L=\dfrac{V_{Th}^{2}R_L}{(R_L+R_{Th})^{2}}$, as duas derivadas parciais são
\[
\frac{\partial P_L}{\partial R_L}=V_{Th}^{2}\frac{R_{Th}-R_L}{(R_L+R_{Th})^{3}};
\qquad
\frac{\partial P_L}{\partial R_{Th}}=-\frac{2V_{Th}^{2}R_L}{(R_L+R_{Th})^{3}}.
\]
A segunda é \textbf{estritamente negativa} para todo $R_L>0$ --- ela nunca se
anula. Não existe ponto em que ambas as derivadas sejam nulas, logo não há
máximo interior.\\
(b) O ótimo migra para o contorno $R_{Th}\to0$. Nesse limite,
\[
P_L\to\frac{V_{Th}^{2}}{R_L},
\]
que por sua vez cresce sem limite conforme $R_L\to0$. \textbf{A potência
diverge} --- o modelo perde sentido físico, pois pressupõe uma fonte de tensão
ideal capaz de fornecer corrente infinita.\\
(c) \textbf{Relevância.} A análise revela que "casamento de impedâncias" não é
um princípio universal de otimização, mas a solução de um problema \emph{muito
específico}: fonte fixa, carga livre.
\begin{itemize}
\item \emph{Fonte fixa, carga livre} $\Rightarrow$ $R_L=R_{Th}$ (casamento).
\item \emph{Carga fixa, fonte livre} $\Rightarrow$ $R_{Th}\to0$ (fonte rígida).
\item \emph{Ambos livres} $\Rightarrow$ problema mal posto, sem solução finita.
\end{itemize}
\textbf{Onde cada caso ocorre:} o primeiro em antenas, transdutores e sensores,
cuja impedância é imposta pela física. O segundo em fontes de alimentação, em
que o projeto persegue $R_{Th}$ o menor possível. Confundir os dois é o erro
conceitual mais comum do tópico --- e leva à ideia equivocada de que fontes de
alimentação deveriam ser "casadas" com suas cargas.}
\stepcounter{questao}
\resposta{$\SI{81}{\watt}$ reais contra $\SI{27.5}{\watt}$ previstos --- o
painel não é uma fonte linear}
\resolucao{(a) \textbf{Modelo de Thévenin:}
$R_{Th}=22/5=\SI{4.4}{\ohm}$ e
\[
P_{\max}^{Th}=\frac{V_{oc}I_{sc}}{4}=\frac{22(5)}{4}=\SI{27.5}{\watt}.
\]
\textbf{Real:} $P=18(4{,}5)=\SI{81}{\watt}$ --- quase \textbf{três vezes}
mais.\\
(b) \textbf{O painel não é linear.} Sua curva $I\times V$ não é uma reta: ela é
quase horizontal em $I\approx I_{sc}$ até perto de $V_{oc}$, onde despenca
abruptamente. O modelo de Thévenin, que supõe a reta
$I=I_{sc}(1-V/V_{oc})$, subestima grosseiramente a área sob a curva.\\
A razão $\dfrac{P_{\max}}{V_{oc}I_{sc}}$ --- chamada \textbf{fator de forma} ---
vale $0{,}25$ para uma fonte linear e $0{,}74$ para este painel. Painéis bons
chegam a $0{,}80$.\\
(c) \textbf{Implicações para o rastreamento.}
\begin{itemize}
\item A carga ótima é $R_L=18/4{,}5=\SI{4}{\ohm}$ --- próxima de
$R_{Th}=\SI{4.4}{\ohm}$ por coincidência numérica, não por casamento.
\item O máximo é \textbf{muito mais agudo} que o de uma fonte linear, pois a
curva tem um "joelho" pronunciado. A tolerância generosa deduzida para fontes
lineares \emph{não} se aplica: errar a carga custa caro.
\item $V_{oc}$ e $I_{sc}$ variam com irradiância e temperatura, deslocando o
ponto ótimo continuamente ao longo do dia.
\end{itemize}
\textbf{Consequência:} é necessário um rastreador ativo (MPPT), que ajusta a
carga aparente em tempo real. E toda a teoria deste capítulo serve aqui apenas
como \emph{referência conceitual} --- o dimensionamento exige a curva real do
dispositivo.}

\gabarito[Capítulo 5 --- Potência e máxima transferência]

\section{Análise Nodal}
\stepcounter{questao}
\resposta{(b)}
\resolucao{O terra é uma referência arbitrária de potencial nulo; os potenciais
dos outros nós são medidos em relação a ele. A escolha não altera as
\emph{diferenças} de potencial nem as correntes --- apenas simplifica as
equações.}
\stepcounter{questao}
\resposta{$V = \SI{18}{\volt}$}
\resolucao{A fonte injeta \SI{9}{\ampere} no nó, que escoam pelos dois resistores
em paralelo:
\[
9 = \frac{V}{6}+\frac{V}{3}=\frac{V+2V}{6}=\frac{3V}{6}=\frac{V}{2}
\ \Rightarrow\ V = \SI{18}{\volt}.
\]
Equivale a $V = I\cdot(6\parallel3) = 9\cdot2 = \SI{18}{\volt}$.}
\stepcounter{questao}
\resposta{$V_1 \approx \SI{9.82}{\volt}$; $V_2 \approx \SI{6.55}{\volt}$}
\resolucao{Com o terra na base, escrevemos a lei dos nós em cada nó.\\
\textbf{Nó $V_1$} (a fonte injeta \SI{6}{\ampere}):
\[
6=\frac{V_1}{2}+\frac{V_1-V_2}{3}.
\]
\textbf{Nó $V_2$:}
\[
\frac{V_1-V_2}{3}=\frac{V_2}{6}.
\]
Da segunda equação, $2(V_1-V_2)=V_2 \Rightarrow V_1 = \dfrac{3V_2}{2}$.
Substituindo na primeira e resolvendo:
\[
V_1=\frac{108}{11}\approx\SI{9.82}{\volt},
\qquad
V_2=\frac{72}{11}\approx\SI{6.55}{\volt}.
\]
\textbf{Verificação:} corrente da fonte $= \dfrac{V_1}{2}+\dfrac{V_1-V_2}{3}
= \dfrac{108}{22}+\dfrac{36}{33}\approx 4{,}91+1{,}09 = \SI{6}{\ampere}$. Resultado consistente.}
\stepcounter{questao}
\resposta{$V = \SI{4}{\volt}$}
\resolucao{A corrente líquida injetada no nó é $3 - 1 = \SI{2}{\ampere}$, que
escoa pelos dois resistores em paralelo:
\[
V = I_{\text{líq}}\cdot(6\parallel3) = 2\cdot2 = \SI{4}{\volt}.
\]
Conferindo pela lei dos nós: $\dfrac{V}{6}+\dfrac{V}{3} = \dfrac{4}{6}+\dfrac{4}{3}
= 0{,}67 + 1{,}33 = \SI{2}{\ampere}$, igual à corrente líquida injetada.}
\stepcounter{questao}
\resposta{$V_A = \SI{12}{\volt}$}
\resolucao{Chamando de $V_A$ o potencial do nó central e notando que o terminal
superior da fonte de tensão está fixo em \SI{12}{\volt}, a lei dos nós em $A$
(correntes saindo $=0$; a fonte de corrente injeta \SI{2}{\ampere} \emph{no} nó):
\[
\frac{V_A-12}{4}+\frac{V_A}{6}-2=0.
\]
Multiplicando por 12: $3(V_A-12)+2V_A-24=0 \Rightarrow 5V_A = 60
\Rightarrow V_A = \SI{12}{\volt}$.\\
\textbf{Interpretação:} nesse ponto, a corrente que a fonte de tensão forneceria
pelo resistor de \SI{4}{\ohm} é nula ($V_A = \SI{12}{\volt}$ iguala os dois lados),
e toda a corrente de \SI{2}{\ampere} da fonte de corrente escoa pelo resistor de
\SI{6}{\ohm}: $2 = 12/6$. Consistente.}
\stepcounter{questao}
\resposta{$V_1 = \dfrac{76}{9}\approx\SI{8.44}{\volt}$,
$V_2=\dfrac{16}{3}\approx\SI{5.33}{\volt}$,
$V_3=\dfrac{40}{9}\approx\SI{4.44}{\volt}$}
\resolucao{\textbf{Montagem.} Com o terra na base e aplicando a lei dos nós
(soma das correntes que \emph{saem} igual a zero) em cada nó:
\[
\text{Nó 1: }\ \frac{V_1-12}{2}+\frac{V_1-V_2}{4}+\frac{V_1-V_3}{4}=0,
\]
\[
\text{Nó 2: }\ \frac{V_2-V_1}{4}+\frac{V_2}{8}+\frac{V_2-V_3}{8}=0,
\]
\[
\text{Nó 3: }\ \frac{V_3-V_1}{4}+\frac{V_3-V_2}{8}+\frac{V_3}{4}=0.
\]
\textbf{Forma matricial} (multiplicando para eliminar denominadores e agrupando
as condutâncias):
\[
\begin{bmatrix}
1{,}0 & -0{,}25 & -0{,}25\\
-0{,}25 & 0{,}5 & -0{,}125\\
-0{,}25 & -0{,}125 & 0{,}625
\end{bmatrix}
\begin{bmatrix}V_1\\V_2\\V_3\end{bmatrix}
=
\begin{bmatrix}6\\0\\0\end{bmatrix}
\]
Note a \textbf{simetria} da matriz de condutâncias --- verificação de que a
montagem está correta.

\textbf{Solução:}
\[
V_1=\frac{76}{9}\approx\SI{8.44}{\volt},\quad
V_2=\frac{16}{3}\approx\SI{5.33}{\volt},\quad
V_3=\frac{40}{9}\approx\SI{4.44}{\volt}.
\]
\textbf{Verificação no nó 2:}
$\dfrac{5{,}33-8{,}44}{4}+\dfrac{5{,}33}{8}+\dfrac{5{,}33-4{,}44}{8}
= -0{,}778+0{,}667+0{,}111 = 0$. \checkmark\\
\textbf{Observação estrutural:} a matriz nodal é o \emph{dual} da matriz de
malhas --- diagonal com a soma das condutâncias do nó, fora da diagonal o negativo
da condutância compartilhada.}
\stepcounter{questao}
\resposta{$V_1 = \SI{8}{\volt}$; $V_2 = \SI{2}{\volt}$}
\resolucao{\textbf{Por que o método normal falha:} a corrente que atravessa a
fonte de \SI{6}{\volt} é desconhecida, então não se pode escrever a lei dos nós
isoladamente em 1 ou em 2.

\medskip
\textbf{Equação 1 --- lei dos nós no supernó.} Traçamos uma superfície fechada
englobando os nós 1 e 2 (tracejado laranja). Tudo o que \emph{entra} deve
\emph{sair}. Entra \SI{2}{\ampere} da fonte de corrente; saem as correntes por
$R_1$ e $R_2$:
\[
\frac{V_1}{R_1}+\frac{V_2}{R_2} = 2
\quad\Longrightarrow\quad
\frac{V_1}{6}+\frac{V_2}{3} = 2.
\]
Note que a corrente \emph{interna} à fonte não aparece --- ela entra e sai da
mesma superfície, cancelando-se. \textbf{Este é o truque do método.}

\medskip
\textbf{Equação 2 --- restrição da fonte.} A fonte impõe diretamente:
\[
V_1 - V_2 = 6.
\]

\medskip
\textbf{Resolvendo.} Da segunda, $V_2 = V_1 - 6$. Substituindo na primeira e
multiplicando por 6:
\[
V_1 + 2(V_1-6) = 12
\;\Rightarrow\;
3V_1 = 24
\;\Rightarrow\;
V_1 = \SI{8}{\volt},\quad V_2 = \SI{2}{\volt}.
\]
\textbf{Verificação:} $\dfrac{8}{6}+\dfrac{2}{3} = 1{,}33+0{,}67 =
\SI{2}{\ampere}$. Confere com a fonte de corrente.}
\stepcounter{questao}
\resposta{4}
\resolucao{Com um nó fixado em zero, restam $5-1=4$ potenciais desconhecidos.}
\stepcounter{questao}
\resposta{$V_A=\SI{18}{\volt}$}
\resolucao{$V_A(1/10+1/15)=3\Rightarrow V_A=\SI{18}{\volt}$.}
\stepcounter{questao}
\resposta{$i_{AB}=(V_A-V_B)/R$}
\resolucao{É a lei de Ohm escrita diretamente em função das tensões nodais.}
\stepcounter{questao}
\resposta{Negativo}
\resolucao{O vetor $\mathbf I$ representa corrente líquida injetada. Uma fonte que retira corrente entra como $-I$.}
\stepcounter{questao}
\resposta{$G_{eq}=\SI{0.5}{\ensuremath{\mathrm{S}}}$}
\resolucao{$G_{eq}=1/4+1/5+1/20=0{,}5\,\ensuremath{\mathrm{S}}$.}
\stepcounter{questao}
\resposta{$\SI{12}{\volt}$}
\resolucao{O terminal positivo é um nó de potencial conhecido; não é necessário escrever uma equação nodal para determiná-lo.}
\stepcounter{questao}
\resposta{Supernó}
\resolucao{A corrente interna da fonte não é conhecida por uma relação resistiva; escreve-se a LKC no contorno do supernó e acrescenta-se a restrição de tensão.}
\stepcounter{questao}
\resposta{$\SI{1.2}{\ampere}$, de $A$ para o terra}
\resolucao{$i=18/15=\SI{1.2}{\ampere}$.}
\stepcounter{questao}
\resposta{$V=\SI{33.33}{\volt}$; $i_{10}=\SI{3.33}{\ampere}$; $i_{20}=\SI{1.67}{\ampere}$}
\resolucao{$5=V/10+V/20\Rightarrow V=100/3\,\si{\volt}$. As correntes são $V/R$.}
\stepcounter{questao}
\resposta{$V=\SI{16}{\volt}$}
\resolucao{A corrente líquida injetada é $6-2=4$ A. Logo $V(1/8+1/8)=4\Rightarrow V=16$ V.}
\stepcounter{questao}
\resposta{$V_A=\SI{8}{\volt}$}
\resolucao{$\frac{V_A-12}{4}+\frac{V_A-6}{6}+\frac{V_A}{12}=0\Rightarrow V_A=8$ V.}
\stepcounter{questao}
\resposta{$V_1=\SI{20}{\volt}$; $V_2=\SI{10}{\volt}$}
\resolucao{$V_1/10+(V_1-V_2)/10=3$ e $V_2/10+(V_2-V_1)/10=0$. Da segunda, $V_1=2V_2$; então $V_2=10$ V.}
\stepcounter{questao}
\resposta{$V_A=\SI{18}{\volt}$}
\resolucao{$\frac{V_A-30}{10}+\frac{V_A}{15}=0\Rightarrow V_A=18$ V.}
\stepcounter{questao}
\resposta{$V_A=\SI{24}{\volt}$; $i=\SI{0.8}{\ampere}$}
\resolucao{$\frac{V_A-20}{5}+\frac{V_A}{20}=2\Rightarrow V_A=24$ V; $i=(24-20)/5=0{,}8$ A.}
\stepcounter{questao}
\resposta{$V_1=\SI{16}{\volt}$; $V_2=\SI{8}{\volt}$}
\resolucao{$V_1/5+(V_1-V_2)/10=4$ e $V_2/10+(V_2-V_1)/10=0$.}
\stepcounter{questao}
\resposta{$V_1=\SI{15.6}{\volt}$; $V_2=\SI{7.2}{\volt}$}
\resolucao{Nó 1: $V_1/6+(V_1-V_2)/6+1=5$. Nó 2: $V_2/3+(V_2-V_1)/6=1$. A solução é $V_1=78/5$ V e $V_2=36/5$ V.}
\stepcounter{questao}
\resposta{$V_1=\SI{25}{\volt}$; $V_2=\SI{10}{\volt}$; $V_3=\SI{5}{\volt}$}
\resolucao{As três LKC formam uma matriz tridiagonal. Da última, $V_2=2V_3$; da segunda, $V_1=2{,}5V_2$; na primeira obtém-se $V_2=10$ V.}
\stepcounter{questao}
\resposta{$V_1=\SI{14}{\volt}$; $V_2=\SI{10}{\volt}$}
\resolucao{$V_1/4+(V_1-V_2)/8=4$ e $V_2/4+(V_2-V_1)/8=2$.}
\stepcounter{questao}
\resposta{$V_1=\SI{10.5}{\volt}$; $V_2=\SI{8.25}{\volt}$}
\resolucao{$\frac{V_1-18}{6}+\frac{V_1}{12}+\frac{V_1-V_2}{6}=0$ e $\frac{V_2}{6}+\frac{V_2-V_1}{6}=1$.}
\stepcounter{questao}
\resposta{$V_1=\SI{8}{\volt}$; $V_2=\SI{12}{\volt}$}
\resolucao{Nó 1: $(V_1-24)/6+V_1/12+2=0\Rightarrow V_1=8$ V. Nó 2: $V_2/6=2\Rightarrow V_2=12$ V.}
\stepcounter{questao}
\resposta{$V_1=\SI{6}{\volt}$; $V_2=\SI{4}{\volt}$; $i_{12}=\SI{0.5}{\ampere}$}
\resolucao{$\frac{V_1-10}{2}+\frac{V_1}{4}+\frac{V_1-V_2}{4}=0$ e $\frac{V_2}{8}+\frac{V_2-V_1}{4}=0$. Logo $V_1=6$ V, $V_2=4$ V e $i_{12}=(6-4)/4=0{,}5$ A.}
\stepcounter{questao}
\resposta{$i_1+i_2=\SI{15}{\ampere}$}
\resolucao{Os quatro ramos estão em paralelo, portanto $i_1=V/1$ e $i_2=V/5$, logo $i_1=5i_2$. Pela LKC no nó superior, $5+4i_2=i_1+i_2$. Assim $5+4i_2=6i_2$, de onde $i_2=\SI{2.5}{\ampere}$, $i_1=\SI{12.5}{\ampere}$ e $i_1+i_2=\SI{15}{\ampere}$.}
\stepcounter{questao}
\resposta{$V_1=\SI{16.8}{\volt}$; $V_2=\SI{13.2}{\volt}$; $i_{12}=\SI{1.2}{\ampere}$ de 1 para 2}
\resolucao{As equações são $V_1/6+(V_1-V_2)/3=4$ e $V_2/6+(V_2-V_1)/3=1$. A solução é $V_1=84/5$ V e $V_2=66/5$ V. Assim $i_{12}=(V_1-V_2)/3=\SI{1.2}{\ampere}$.}
\stepcounter{questao}
\resposta{$V_1=\dfrac{56}{3}\,\si{\volt}$; $V_2=\dfrac{20}{3}\,\si{\volt}$; $i_s=\dfrac{5}{3}\,\si{\ampere}$ de 1 para 2}
\resolucao{No supernó: $V_1/8+V_2/4=4$. A fonte impõe $V_1-V_2=12$. Resulta $V_1=56/3$ V e $V_2=20/3$ V. No nó 1, $i_s=4-V_1/8=5/3$ A.}
\stepcounter{questao}
\resposta{$V_1=\dfrac{40}{3}\,\si{\volt}$; $V_2=\dfrac{20}{3}\,\si{\volt}$; $V_3=\dfrac{10}{3}\,\si{\volt}$}
\resolucao{As LKC são $V_1/4+(V_1-V_2)/4=5$, $V_2/8+(V_2-V_1)/4+(V_2-V_3)/4=0$ e $V_3/4+(V_3-V_2)/4=0$. A solução é $(40/3,20/3,10/3)$ V.}
\stepcounter{questao}
\resposta{$V=\SI{9}{\volt}$}
\resolucao{Com a fonte de corrente retirando $\SI{1}{\ampere}$ do nó: $(V-18)/6+(V-12)/3+V/6+1=0$. Daí $V=\SI{9}{\volt}$.}
\stepcounter{questao}
\resposta{$V_1=\SI{14}{\volt}$; $V_2=\SI{6}{\volt}$; $V_3=\SI{3}{\volt}$}
\resolucao{Supernó 1--2: $V_1/10+(V_2-V_3)/5=2$; restrição $V_1-V_2=8$; nó 3: $V_3/5+(V_3-V_2)/5=0$. A solução é $(14,6,3)$ V.}
\stepcounter{questao}
\resposta{$V_1=\SI{16}{\volt}$; $V_2=\SI{16}{\volt}$}
\resolucao{$V_1/4+(V_1-V_2)/8=4$ e $V_2/8+(V_2-V_1)/8=0{,}5(V_1/4)$. A solução é $V_1=V_2=\SI{16}{\volt}$; portanto o resistor entre os nós conduz corrente nula.}
\stepcounter{questao}
\resposta{$V_1=\SI{24}{\volt}$; $V_2=\SI{48}{\volt}$; $i_{12}=\SI{-3}{\ampere}$}
\resolucao{A LKC no nó 1 é $V_1/4+(V_1-V_2)/8=3$. Com $V_2=2V_1$, resulta $V_1=24$ V e $V_2=48$ V. Logo $i_{12}=(24-48)/8=-3$ A, ou seja, a corrente real flui do nó 2 para o nó 1.}
\stepcounter{questao}
\resposta{$V_1=\SI{43.2}{\volt}$; $V_2=V_3=\SI{28.8}{\volt}$; $i_{23}=0$}
\resolucao{As duas divisões laterais têm a mesma razão: $4/8=6/12$. A solução nodal confirma $V_2=V_3=28{,}8$ V e $V_1=43{,}2$ V. Consequentemente, $i_{23}=(V_2-V_3)/10=0$.}
\stepcounter{questao}
\resposta{$V_1=\SI{12}{\volt}$; $V_2=\SI{8}{\volt}$}
\resolucao{Diretamente: $(V_1-24)/6+V_1/12+(V_1-V_2)/4=0$ e $V_2/8+(V_2-V_1)/4=0$, resultando $(12,8)$ V. A transformação Norton produz uma fonte de $4$ A em paralelo com $6\,\ohm$ e conduz ao mesmo sistema.}
\stepcounter{questao}
\resposta{$V_1=\SI{10}{\volt}$; $V_2=\SI{0}{\volt}$; $i_s=0$}
\resolucao{Supernó: $V_1/5+V_2/10=2$ e $V_1-V_2=10$. A solução é $V_1=10$ V e $V_2=0$. Como o resistor de $5\,\ohm$ já conduz os $2$ A injetados, a corrente na fonte de tensão é nula.}
\stepcounter{questao}
\resposta{$V_1=\SI{18.75}{\volt}$; $V_2=\SI{12.5}{\volt}$; $P=\SI{93.75}{\watt}$ fornecidos}
\resolucao{As LKC são $V_1/5+(V_1-V_2)/10+0{,}05V_2=5$ e $V_2/10+(V_2-V_1)/10-0{,}05V_2=0$. Resulta $V_1=75/4$ V e $V_2=25/2$ V. A fonte independente fornece $P=V_1I=18{,}75\cdot5=\SI{93.75}{\watt}$.}
\stepcounter{questao}
\resposta{$V_1=\dfrac{70}{3}\,\si{\volt}$; $V_2=\dfrac{40}{3}\,\si{\volt}$; $i_s=\dfrac{2}{3}\,\si{\ampere}$ de 1 para 2}
\resolucao{Com $i_s$ positivo de 1 para 2:
\[
\frac{V_1}{10}+i_s=3,\qquad \frac{V_2}{20}-i_s=0,\qquad V_1-V_2=10.
\]
Na forma matricial,
\[
\begin{bmatrix}0{,}1&0&1\\0&0{,}05&-1\\1&-1&0\end{bmatrix}
\begin{bmatrix}V_1\\V_2\\i_s\end{bmatrix}=
\begin{bmatrix}3\\0\\10\end{bmatrix}.
\]
A solução é $V_1=70/3$ V, $V_2=40/3$ V e $i_s=2/3$ A.}
\stepcounter{questao}
\resposta{$V_{Th}=\SI{40}{\volt}$; $R_{Th}=\SI{30}{\ohm}$}
\resolucao{Com a porta aberta não circula corrente no resistor de $20\,\ohm$, logo $V_a=V_1$. Toda a corrente de $4$ A passa pelo resistor de $10\,\ohm$, então $V_{Th}=40$ V. Para $R_{Th}$, zere a fonte de corrente (abra-a) e aplique uma fonte de teste na porta: a porta vê $20\,\ohm$ em série com $10\,\ohm$, portanto $R_{Th}=30\,\ohm$.}
\stepcounter{questao}
\resposta{$R_L\le\SI{20}{\ohm}$; no limite, $P_L=\SI{28.8}{\watt}$}
\resolucao{A tensão é $V=2(30\parallel R_L)$. Impor $V\le24$ fornece
\[
2\frac{30R_L}{30+R_L}\le24\Rightarrow R_L\le20\,\ohm.
\]
No limite $V=24$ V e $P_L=V^2/R_L=576/20=\SI{28.8}{\watt}$.}
\stepcounter{questao}
\resposta{$R_L=\SI{8}{\ohm}$; $\left.dV/dR_L\right|_{8\Omega}=\SI{0.333}{\volt\per\ohm}$}
\resolucao{$V=12R_L/(4+R_L)$. De $8=12R_L/(4+R_L)$ resulta $R_L=8\,\ohm$. A derivada é
\[
\frac{dV}{dR_L}=\frac{48}{(4+R_L)^2},
\]
logo, em $R_L=8\,\ohm$, $dV/dR_L=48/144=1/3\,\si{\volt\per\ohm}$.}
\stepcounter{questao}
\resposta{$V_1=\dfrac{520}{21}\,\si{\volt}$; $V_2=\dfrac{200}{21}\,\si{\volt}$; $V_3=\dfrac{80}{21}\,\si{\volt}$; $V_4=\dfrac{40}{21}\,\si{\volt}$; $R_{in}=\dfrac{130}{21}\,\ohm\approx\SI{6.19}{\ohm}$}
\resolucao{Multiplicando as LKC por $10$, obtém-se
\[
\begin{bmatrix}2&-1&0&0\\-1&3&-1&0\\0&-1&3&-1\\0&0&-1&2\end{bmatrix}
\begin{bmatrix}V_1\\V_2\\V_3\\V_4\end{bmatrix}
=\begin{bmatrix}40\\0\\0\\0\end{bmatrix}.
\]
A solução é $(520,200,80,40)/21$ V. Assim $R_{in}=V_1/4=130/21\,\ohm\approx6{,}19\,\ohm$.}
\stepcounter{questao}
\resposta{$V_1=\SI{25}{\volt}$; $V_2=\SI{12.5}{\volt}$; $i_s=\SI{2.5}{\ampere}$ de 1 para 2; a fonte dependente absorve $\SI{31.25}{\watt}$}
\resolucao{Supernó: $V_1/10+V_2/5=5$. A restrição é $V_1-V_2=5(V_2/5)=V_2$, portanto $V_1=2V_2$. Da LKC, $V_2=12{,}5$ V e $V_1=25$ V. No nó 1, $i_s=5-25/10=2{,}5$ A. A tensão da fonte é $12{,}5$ V e a corrente entra em seu terminal positivo, logo ela absorve $12{,}5\cdot2{,}5=\SI{31.25}{\watt}$.}
\stepcounter{questao}
\resposta{$V(g)=1/(0{,}05-g)$; $g_c=\SI{0.05}{\ensuremath{\mathrm{S}}}$; para $g=\SI{0.04}{\ensuremath{\mathrm{S}}}$, $V=\SI{100}{\volt}$}
\resolucao{A LKC fornece $V/20=1+gV$, então
\[
V(g)=\frac{1}{0{,}05-g}.
\]
Em $g_c=0{,}05$ S o coeficiente nodal se anula e a solução linear deixa de ser finita. Para $g=0{,}04$ S, $V=100$ V.}
\stepcounter{questao}
\resposta{$V_{Th}=\SI{50}{\volt}$; $R_{Th}=\SI{30}{\ohm}$}
\resolucao{Com a porta aberta, $(V_2-V_1)/5=0{,}05V_1$, logo $V_2=1{,}25V_1$. No nó 1, $V_1/10+(V_1-V_2)/5=2$, resultando $V_1=40$ V e $V_{Th}=50$ V. Para $R_{Th}$, abra a fonte independente e aplique uma corrente de teste $I_t$ na porta. A LKC do nó 1 dá $V_1=(2/3)V_2$ e, no nó 2, obtém-se $I_t=V_2/30$. Portanto $R_{Th}=V_2/I_t=30\,\ohm$.}
\stepcounter{questao}
\resposta{Equilíbrio quando $R_1/R_2=R_3/R_4$; para a figura, $V_1=\SI{43.2}{\volt}$ e $V_2=V_3=\SI{28.8}{\volt}$}
\resolucao{Corrente nula em $R_b$ exige $V_2=V_3$. Pelos divisores laterais, isso ocorre quando
\[
\frac{R_2}{R_1+R_2}=\frac{R_4}{R_3+R_4}
\quad\Longleftrightarrow\quad
\frac{R_1}{R_2}=\frac{R_3}{R_4}.
\]
Na figura, $4/8=6/12$, portanto a ponte está equilibrada. As duas pernas equivalem a $12\,\ohm$ e $18\,\ohm$ em paralelo, logo $R_{eq}=7{,}2\,\ohm$ e $V_1=6\cdot7{,}2=43{,}2$ V. Os pontos médios ficam em $V_2=V_3=28{,}8$ V.}
\stepcounter{questao}
\resposta{$I_1=\SI{5}{\ampere}$ e $I_2=\SI{1}{\ampere}$, ambas injetando corrente; potência total fornecida $=\SI{102}{\watt}$}
\resolucao{Pela LKC,
\[
I_1=\frac{18}{6}+\frac{18-12}{3}=3+2=5\,\text{A},
\]
\[
I_2=\frac{12}{4}+\frac{12-18}{3}=3-2=1\,\text{A}.
\]
As fontes fornecem $18\cdot5+12\cdot1=102$ W. Os resistores dissipam $18^2/6+12^2/4+(18-12)^2/3=54+36+12=102$ W.}

\gabarito[Capítulo 7 --- Análise Nodal]

\section{Método de Maxwell (Análise de Malhas)}
\stepcounter{questao}
\resposta{$I_1=\SI{3.2}{\ampere}$; $I_2=\SI{2.4}{\ampere}$; $I_c=\SI{0.8}{\ampere}$}
\resolucao{Com correntes de malha horárias,
\[
10I_1-5I_2=20,
\qquad
-5I_1+10I_2=2{,}5I_1.
\]
A segunda equação dá $I_2=0{,}75I_1$. Substituindo na primeira,
$I_1=\SI{3.2}{\ampere}$ e $I_2=\SI{2.4}{\ampere}$. No ramo comum,
$I_c=I_1-I_2=\SI{0.8}{\ampere}$.}
\stepcounter{questao}
\resposta{$I_1=\SI{4}{\ampere}$; $I_2=\SI{2}{\ampere}$; $|v_s|=\SI{10}{\volt}$}
\resolucao{Por estar na periferia, a fonte fixa diretamente
$I_2=0{,}5I_1$. Na malha 1,
\[
10I_1-5I_2=30.
\]
Logo $10I_1-2{,}5I_1=30$, de onde $I_1=\SI{4}{\ampere}$ e
$I_2=\SI{2}{\ampere}$. Na malha 2, as quedas resistivas totalizam
$10(2)+5(2-4)=\SI{10}{\volt}$ em módulo; a fonte de corrente sustenta essa tensão.}
\stepcounter{questao}
\resposta{$I_1=\SI{4}{\ampere}$; $I_2=\SI{2}{\ampere}$}
\resolucao{A malha 1 fornece $8I_1-4I_2=24$. Na malha 2,
\[
-4I_1+12I_2=4(I_1-I_2),
\]
ou $-8I_1+16I_2=0$. Assim $I_1=2I_2$ e, pela primeira equação,
$I_1=\SI{4}{\ampere}$ e $I_2=\SI{2}{\ampere}$.}
\stepcounter{questao}
\resposta{$I_1=\SI{4.8}{\ampere}$; $I_2=\SI{3.6}{\ampere}$}
\resolucao{Como $v_x=4I_1$, a fonte dependente vale $2I_1$ volts. Logo,
\[
8I_1-4I_2=24,
\qquad
-4I_1+8I_2=2I_1.
\]
Da segunda, $I_2=0{,}75I_1$. Substituindo na primeira, obtêm-se
$I_1=\SI{4.8}{\ampere}$ e $I_2=\SI{3.6}{\ampere}$.}
\stepcounter{questao}
\resposta{$I_1=\SI{3.2}{\ampere}$; $I_2=\SI{4.8}{\ampere}$}
\resolucao{Como $v_x=5I_1$, a restrição imposta pela fonte é
\[
I_2-I_1=0{,}1(5I_1)=0{,}5I_1,
\]
portanto $I_2=1{,}5I_1$. A LKT na supermalha externa é
$40=5I_1+5I_2$. Substituindo, $I_1=\SI{3.2}{\ampere}$ e
$I_2=\SI{4.8}{\ampere}$.}
\stepcounter{questao}
\resposta{$I_1=\SI{2}{\ampere}$; $I_2=\SI{6}{\ampere}$}
\resolucao{A restrição é $I_2-I_1=2I_1$, portanto $I_2=3I_1$. A LKT na
supermalha é $40=5I_1+5I_2$. Assim $40=20I_1$, resultando
$I_1=\SI{2}{\ampere}$ e $I_2=\SI{6}{\ampere}$.}
\stepcounter{questao}
\resposta{$I_1=\dfrac{8}{3}\,\si{\ampere}$; $I_2=I_3=\dfrac{4}{3}\,\si{\ampere}$}
\resolucao{As três equações são
\[
10I_1-5I_2=20,
\]
\[
-5I_1+15I_2-5I_3=0,
\]
\[
-5I_2+10I_3=2{,}5I_1.
\]
A solução do sistema fornece $I_1=8/3$ A e $I_2=I_3=4/3$ A.}
\stepcounter{questao}
\resposta{$I_1=\SI{3}{\ampere}$; $I_2=0$; $I_3=\SI{1.5}{\ampere}$}
\resolucao{A malha 1 fornece $10I_1-5I_2=30$. A supermalha 2--3 fornece
\[
-5I_1+10I_2+10I_3=0,
\]
e a fonte impõe $I_3-I_2=0{,}5I_1$. Resolvendo o sistema,
$I_1=\SI{3}{\ampere}$, $I_2=0$ e $I_3=\SI{1.5}{\ampere}$.}
\stepcounter{questao}
\resposta{$I_1=\SI{3}{\ampere}$; $k=\dfrac{5}{3}\,\ohm\approx\SI{1.67}{\ohm}$}
\resolucao{A malha 1 dá $10I_1-5I_2=20$. Impondo $I_2=2$ A,
$I_1=3$ A. Na malha 2,
\[
-5I_1+10I_2=kI_1.
\]
Logo $-15+20=3k$, e portanto $k=5/3\,\ohm$.}
\stepcounter{questao}
\resposta{$I_1=\dfrac{10}{3}\,\si{\ampere}$; $I_2=\dfrac{5}{3}\,\si{\ampere}$;
$i_d=\dfrac{5}{3}\,\si{\ampere}$ para baixo; $P_d=\dfrac{100}{9}\,\si{\watt}$ absorvidos}
\resolucao{Na malha 1, $4I_1+v_d=20$ e $v_d=2I_1$, portanto
$I_1=10/3$ A. Na malha 2, $4I_2-v_d=0$, então $I_2=5/3$ A. A corrente real no
ramo comum é $i_d=I_1-I_2=5/3$ A para baixo. Como $v_d=20/3$ V e a corrente
entra no terminal positivo superior da fonte dependente,
\[
P_d=v_di_d=\frac{20}{3}\frac{5}{3}=\frac{100}{9}\,\si{\watt},
\]
logo a fonte dependente \textbf{absorve} potência.}
\stepcounter{questao}
\resposta{Horário}
\resolucao{A seta percorre o laço no sentido horário. Essa é apenas uma
convenção de referência; se o resultado algébrico de $I_1$ for negativo, a
corrente real circula no sentido anti-horário.}
\stepcounter{questao}
\resposta{$i_c=I_1-I_2$}
\resolucao{As correntes horárias atravessam o ramo compartilhado em sentidos
opostos; por isso a corrente de ramo é a diferença algébrica.}
\stepcounter{questao}
\resposta{$m=5$}
\resolucao{$m=b-n+1=11-7+1=5$.}
\stepcounter{questao}
\resposta{$I_2=I_s$}
\resolucao{A fonte pertence somente àquela malha; portanto ela fixa diretamente
a corrente de malha e elimina uma incógnita do sistema.}
\stepcounter{questao}
\resposta{Supermalha, mais a equação de restrição da fonte}
\resolucao{A tensão sobre uma fonte ideal de corrente é desconhecida; por isso
não se escreve LKT separadamente nas duas malhas. Contorna-se o ramo da fonte e
acrescenta-se a relação entre $I_1$ e $I_2$.}
\stepcounter{questao}
\resposta{$R_{kk}$ é a soma das resistências da malha $k$; $R_{jk}$ é o negativo
da resistência compartilhada}
\resolucao{A diagonal reúne todos os resistores percorridos por $I_k$. Fora da
diagonal aparece $-R_c$ porque as correntes horárias percorrem o ramo comum em
sentidos opostos.}
\stepcounter{questao}
\resposta{A corrente real tem módulo $\SI{0.8}{\ampere}$ e circula no sentido anti-horário}
\resolucao{O sinal negativo não indica erro: ele informa que a escolha inicial
do sentido foi oposta à corrente física.}
\stepcounter{questao}
\resposta{$I_1=\dfrac{18}{7}\,\si{\ampere}$; $I_2=\dfrac{9}{7}\,\si{\ampere}$;
$i_c=\dfrac{9}{7}\,\si{\ampere}$}
\resolucao{As equações são
\[
10I_1-6I_2=18,\qquad -6I_1+12I_2=0.
\]
Da segunda, $I_1=2I_2$. Logo $14I_2=18$, resultando
$I_2=9/7$ A e $I_1=18/7$ A. No ramo comum,
$i_c=I_1-I_2=9/7$ A.}
\stepcounter{questao}
\resposta{$I_1=\SI{3.75}{\ampere}$; $I_2=\SI{1.5}{\ampere}$;
$i_c=\SI{2.25}{\ampere}$}
\resolucao{A fonte periférica fixa $I_2=1{,}5$ A. Para a malha 1,
\[
5I_1+5(I_1-I_2)=30,
\]
portanto $10I_1-7{,}5=30$ e $I_1=3{,}75$ A. Assim
$i_c=I_1-I_2=2{,}25$ A.}
\stepcounter{questao}
\resposta{$I_1=\SI{3}{\ampere}$; $I_2=\SI{1.5}{\ampere}$}
\resolucao{Percorrendo as duas malhas no sentido horário:
\[
8I_1-4I_2=18,
\qquad
-4I_1+12I_2=6.
\]
A solução é $I_1=3$ A e $I_2=1{,}5$ A. A fonte compartilhada entra com sinais
opostos nas duas equações porque é atravessada em sentidos opostos.}
\stepcounter{questao}
\resposta{$I_1=\SI{3}{\ampere}$; $I_2=\SI{2}{\ampere}$; $I_3=\SI{1}{\ampere}$}
\resolucao{Por inspeção,
\[
\begin{bmatrix}
7&-2&-1\\
-2&10&-3\\
-1&-3&10
\end{bmatrix}
\begin{bmatrix}I_1\\I_2\\I_3\end{bmatrix}
=
\begin{bmatrix}16\\11\\1\end{bmatrix}.
\]
Substituindo $(3,2,1)$ verifica-se cada linha; logo essa é a solução. O desenho
mostra que a matriz depende de quais raios são compartilhados, e não de uma
sequência linear de malhas.}
\stepcounter{questao}
\resposta{$I_1=\dfrac{38}{21}\,\si{\ampere}$;
$I_2=-\dfrac{10}{21}\,\si{\ampere}$}
\resolucao{O sistema é
\[
\begin{bmatrix}10&-4\\-4&10\end{bmatrix}
\begin{bmatrix}I_1\\I_2\end{bmatrix}
=
\begin{bmatrix}20\\-12\end{bmatrix}.
\]
Daí $I_1=38/21$ A e $I_2=-10/21$ A. Portanto a corrente física da segunda
malha tem módulo $10/21$ A e circula no sentido anti-horário.}
\stepcounter{questao}
\resposta{$E=\SI{48}{\volt}$, impulsionando $I_2$ no sentido horário}
\resolucao{Corrente nula no ramo comum exige $I_1=I_2=I$. A primeira malha dá
$(8-4)I=24$, ou $4I=24$, logo $I=6$ A. Na segunda,
$(12-4)I=E$, portanto $E=8\cdot6=48$ V, com polaridade que impulsione $I_2$.}
\stepcounter{questao}
\resposta{$I_1=\SI{5}{\ampere}$; $I_2=\SI{3}{\ampere}$;
$P_{fontes}=P_R=\SI{216}{\watt}$}
\resolucao{O sistema
\[
\begin{bmatrix}9&-3\\-3&9\end{bmatrix}
\begin{bmatrix}I_1\\I_2\end{bmatrix}
=
\begin{bmatrix}36\\12\end{bmatrix}
\]
fornece $(I_1,I_2)=(5,3)$ A. A corrente no resistor comum é $2$ A. Assim
\[
P_R=6(5)^2+6(3)^2+3(2)^2=150+54+12=216\,\text{W}.
\]
As fontes fornecem $36(5)+12(3)=216$ W.}
\stepcounter{questao}
\resposta{$I_1=\SI{3}{\ampere}$; $v_s=\SI{-15}{\volt}$}
\resolucao{Na malha 1,
$5I_1+5(I_1-2)=20$, então $I_1=3$ A. Na malha 2, percorrida no sentido horário,
\[
10(2)+v_s-5(3-2)=0,
\]
portanto $v_s=-15$ V. Isso significa que o terminal inferior da fonte está
$15$ V acima do terminal superior.}
\stepcounter{questao}
\resposta{$I_1=\SI{1.5}{\ampere}$; $I_2=\SI{3.5}{\ampere}$;
$I_3=\SI{0.75}{\ampere}$}
\resolucao{As equações são
\[
4I_1+6I_2-4I_3=24,
\qquad I_2-I_1=2,
\qquad -4I_2+8I_3=-8.
\]
Resolvendo, obtém-se $I_1=3/2$ A, $I_2=7/2$ A e $I_3=3/4$ A.}
\stepcounter{questao}
\resposta{$I_1=\SI{2}{\ampere}$; $I_2=\SI{-1.5}{\ampere}$;
$P_{6V}=\SI{21}{\watt}$ absorvidos}
\resolucao{Na malha 1, $6I_1+6=18$, logo $I_1=2$ A. Na malha 2,
$4I_2+12-6=0$, então $I_2=-1{,}5$ A. A corrente real na fonte comum, para
baixo, é $I_1-I_2=3{,}5$ A. Como ela entra no terminal positivo superior,
\[
P=6(3{,}5)=21\,\text{W},
\]
portanto a fonte de $6$ V absorve potência.}
\stepcounter{questao}
\resposta{$I_2=\SI{2}{\ampere}$; o resistor de $\SI{2}{\ohm}$ conduz corrente nula}
\resolucao{A equação da malha central é
\[
(2+4+3)I_2-2I_1-3I_3=11.
\]
Com $I_1=2$ A e $I_3=1$ A:
$9I_2-4-3=11$, portanto $I_2=2$ A. Assim, no resistor compartilhado entre as
malhas 1 e 2, $I_1-I_2=0$.}
\stepcounter{questao}
\resposta{$I_1=\SI{5}{\ampere}$; $I_2=I_3=\SI{2}{\ampere}$;
$i_{ponte}=0$}
\resolucao{A matriz é
\[
\begin{bmatrix}
6&-2&-4\\
-2&10&-5\\
-4&-5&15
\end{bmatrix}
\begin{bmatrix}I_1\\I_2\\I_3\end{bmatrix}
=
\begin{bmatrix}18\\0\\0\end{bmatrix}.
\]
A solução é $(5,2,2)$ A. O resistor de ponte é comum às malhas 2 e 3, logo
$i_{ponte}=I_2-I_3=0$. O resultado expressa o equilíbrio da ponte.}
\stepcounter{questao}
\resposta{$I_1=\SI{2}{\ampere}$; $I_2=\SI{1}{\ampere}$;
$R_{eq}=\SI{6}{\ohm}$}
\resolucao{As equações são
\[
9I_1-6I_2=12,\qquad -6I_1+12I_2=0.
\]
Da segunda, $I_1=2I_2$; substituindo na primeira, $I_2=1$ A e $I_1=2$ A. A
fonte de teste fornece $2$ A, portanto $R_{eq}=12/2=6\,\ohm$.}
\stepcounter{questao}
\resposta{$I_1=\SI{4}{\ampere}$; $I_2=\SI{3}{\ampere}$;
$I_3=\SI{2}{\ampere}$; $I_4=\SI{1}{\ampere}$}
\resolucao{Cada malha contém $4\,\ohm$ exclusivo e dois raios de
$2\,\ohm$. A matriz é cíclica:
\[
\begin{bmatrix}
8&-2&0&-2\\
-2&8&-2&0\\
0&-2&8&-2\\
-2&0&-2&8
\end{bmatrix}
\begin{bmatrix}I_1\\I_2\\I_3\\I_4\end{bmatrix}
=
\begin{bmatrix}24\\12\\8\\-4\end{bmatrix}.
\]
A solução é $(4,3,2,1)$ A. O termo $-4$ representa a fonte da quarta malha
orientada contra a corrente horária de referência.}
\stepcounter{questao}
\resposta{$R=\SI{6}{\ohm}$; $I_1=\dfrac{20}{7}\,\si{\ampere}$;
$I_2=\dfrac{10}{7}\,\si{\ampere}$}
\resolucao{Na malha 2,
\[
-RI_1+(6+R)I_2=0.
\]
Impondo $I_2=I_1/2$:
$-R+(6+R)/2=0$, logo $R=6\,\ohm$. Na malha 1,
$10I_1-6I_2=20$; com $I_2=I_1/2$, resulta $7I_1=20$.}
\stepcounter{questao}
\resposta{$R_L=\dfrac{8}{3}\,\ohm\approx\SI{2.67}{\ohm}$;
$P_{L,\max}=\SI{13.5}{\watt}$}
\resolucao{As equações são
\[
(4+R_L)I_1-R_LI_2=24,
\]
\[
-R_LI_1+(8+R_L)I_2=12.
\]
A corrente na carga é
\[
i_L=I_1-I_2=\frac{36}{3R_L+8}.
\]
Logo
\[
P_L=R_Li_L^2=\frac{1296R_L}{(3R_L+8)^2}.
\]
Derivando e igualando a zero, obtém-se $R_L=8/3\,\ohm$. Nesse ponto,
$i_L=36/16=2{,}25$ A e $P_{L,\max}=13{,}5$ W.}
\stepcounter{questao}
\resposta{$I_1=\SI{3}{\ampere}$; $I_2=0$; $I_3=\SI{-6}{\ampere}$;
$E=\SI{30}{\volt}$ orientada contra $I_3$ horário}
\resolucao{O sistema é
\[
8I_1-2I_2=24,
\qquad
-2I_1+7I_2-I_3=0,
\qquad
-I_2+5I_3=E_s,
\]
onde $E_s$ é positivo se impulsiona $I_3$ horário. Impondo $I_2=0$:
$I_1=3$ A e, pela segunda equação, $I_3=-6$ A. Assim
$E_s=5(-6)=-30$ V. Portanto a fonte deve ter magnitude $30$ V e polaridade
oposta à referência horária.}
\stepcounter{questao}
\resposta{$R=\SI{8}{\ohm}$; $P_R=\SI{8}{\watt}$}
\resolucao{Pela malha 1,
\[
(4+R)3-R(2)=20 \Rightarrow 12+R=20 \Rightarrow R=8\,\ohm.
\]
A malha 2 confirma o mesmo valor:
\[
-3R+(6+R)2=4 \Rightarrow 12-R=4 \Rightarrow R=8\,\ohm.
\]
A corrente no ramo comum é $I_1-I_2=1$ A; portanto
$P_R=8(1)^2=8$ W.}

\gabarito[Capítulo 9 --- Método de Maxwell]

\section{Capacitores e Indutores}

\subsection{Capacitores}
\stepcounter{questao}
\resposta{(b)}
\resolucao{De $C = \dfrac{\varepsilon_0 A}{d}$: a capacitância cresce com a área
$A$ e \emph{diminui} com a distância $d$ --- o que elimina (a). Ela não depende da
tensão (elimina e), e inserir um dielétrico \emph{aumenta} $C$ (elimina d).}
\stepcounter{questao}
\resposta{$C \approx \SI{88.5}{\pico\farad}$}
\resolucao{Convertendo: $A = \SI{100}{\centi\meter\squared} = \pot{100}{-4}\,\si{\meter\squared}
= \pot{1}{-2}\,\si{\meter\squared}$ e $d = \pot{1}{-3}\,\si{\meter}$:
\[
C=\frac{\varepsilon_0 A}{d}
=\frac{\pot{8.85}{-12}\cdot\pot{1}{-2}}{\pot{1}{-3}}
=\pot{8.85}{-11}\,\si{\farad}\approx\SI{88.5}{\pico\farad}.
\]}
\stepcounter{questao}
\resposta{$Q = \SI{120}{\micro\coulomb}$; $E = \SI{720}{\micro\joule}$}
\resolucao{$Q = C\,U = \pot{10}{-6}\cdot12 = \SI{120}{\micro\coulomb}$.\\
$E = \tfrac12 C\,U^2 = \tfrac12\cdot\pot{10}{-6}\cdot12^2
= \SI{720}{\micro\joule}$.}
\stepcounter{questao}
\resposta{(a)}
\resolucao{Em série, $\Ceq = \dfrac{C_1 C_2}{C_1 + C_2} = \dfrac{6\cdot3}{9}
= \SI{2}{\micro\farad}$. Note que, no capacitor, a série se comporta como o
paralelo do resistor: o equivalente é \emph{menor} que o menor deles.}
\stepcounter{questao}
\resposta{(c)}
\resolucao{No circuito RLC em paralelo, as capacitâncias \emph{somam}:
$\Ceq = 4 + 12 = \SI{16}{\micro\farad}$.}
\stepcounter{questao}
\resposta{(c)}
\resolucao{$E = \tfrac12 C U^2 = \tfrac12\cdot\pot{100}{-6}\cdot50^2
= \tfrac12\cdot\pot{100}{-6}\cdot2500 = \SI{0.125}{\joule} = \SI{125}{\milli\joule}$.
Essa energia fica armazenada no campo elétrico entre as placas e pode ser liberada
rapidamente --- princípio do flash fotográfico e do desfibrilador.}
\stepcounter{questao}
\resposta{$\Ceq \approx \SI{1.64}{\micro\farad}$}
\resolucao{\textbf{Paralelo} $C_1$ e $C_2$: somam-se,
$C_{12} = 6 + 3 = \SI{9}{\micro\farad}$.\\
\textbf{Série} com $C_3$:
\[
\Ceq=\frac{C_{12}\,C_3}{C_{12}+C_3}=\frac{9\cdot2}{11}=\frac{18}{11}
\approx\SI{1.64}{\micro\farad}.
\]}
\stepcounter{questao}
\resposta{$\Ceq = \SI{8}{\micro\farad}$}
\resolucao{\textbf{Série} dos dois de \SI{8}{\micro\farad}:
$\dfrac{8}{2} = \SI{4}{\micro\farad}$.\\
\textbf{Paralelo} com o de \SI{4}{\micro\farad}:
$4 + 4 = \SI{8}{\micro\farad}$.}
\stepcounter{questao}
\resposta{(a) $\tau = \SI{1}{\milli\second}$; (b) $\approx\SI{6.32}{\volt}$;
(c) \SI{10}{\volt}}
\resolucao{(a) $\tau = RC = \pot{1}{3}\cdot\pot{1}{-6} = \SI{1}{\milli\second}$.\\
(b) Na carga, $v_C(\tau) = U_0(1 - e^{-1}) = 10\cdot0{,}632 \approx \SI{6.32}{\volt}$
--- o capacitor atinge cerca de $\SI{63}{\percent}$ da tensão final em um $\tau$.\\
(c) Em regime permanente ($t\to\infty$), o capacitor está totalmente carregado,
não circula corrente, e $v_C = U_0 = \SI{10}{\volt}$: o capacitor se comporta como
um circuito \emph{aberto}.}
\stepcounter{questao}
\resposta{$v_C = \SI{4}{\volt}$}
\resolucao{Em regime permanente, o capacitor está carregado e \emph{não conduz}
(ramo aberto). Toda a corrente passa por $R_1$ e $R_2$ em série:
\[
I=\frac{12}{4+2}=\SI{2}{\ampere}.
\]
A tensão no capacitor iguala a tensão sobre $R_2$ (estão em paralelo):
\[
v_C = R_2\,I = 2\cdot2 = \SI{4}{\volt}.
\]
\textbf{Princípio-chave:} em regime permanente CC, capacitor = circuito aberto
(e indutor = curto-circuito). Isso transforma o problema em uma simples análise
resistiva.}
\stepcounter{questao}
\resposta{(a) \SI{2}{\micro\farad}; (b) \SI{18}{\micro\farad}; (c) \SI{4}{\micro\farad}}
\resolucao{(a) Série de três iguais: $\Ceq = \dfrac{C}{3} = \SI{2}{\micro\farad}$.\\
(b) Paralelo de três iguais: $\Ceq = 3C = \SI{18}{\micro\farad}$.\\
(c) Dois em paralelo: $6+6 = \SI{12}{\micro\farad}$; em série com o terceiro:
$\dfrac{12\cdot6}{18} = \SI{4}{\micro\farad}$.\\
Repare como a topologia muda drasticamente o resultado, de \SI{2}{} a
\SI{18}{\micro\farad} com os mesmos três componentes.}
\stepcounter{questao}
\resposta{(a) $\tau = \SI{1}{\milli\second}$; (b) $I_0 = \SI{12}{\milli\ampere}$;
(c) $\approx\SI{4.4}{\micro\coulomb}$}
\resolucao{(a) $\tau = RC = 500\cdot\pot{2}{-6} = \SI{1}{\milli\second}$.\\
(b) A corrente inicial é limitada apenas pelo resistor:
$I_0 = \dfrac{U_0}{R} = \dfrac{6}{500} = \SI{12}{\milli\ampere}$.\\
(c) A carga inicial é $Q_0 = C\,U_0 = \pot{2}{-6}\cdot6 = \SI{12}{\micro\coulomb}$.
Após um $\tau$, resta $Q(\tau) = Q_0\,e^{-1} = 12\cdot0{,}368
\approx\SI{4.4}{\micro\coulomb}$ --- cerca de $\SI{37}{\percent}$ do inicial.}
\stepcounter{questao}
\resposta{$\tau = \SI{1}{\milli\second}$; $U_0 = \SI{10}{\volt}$}
\resolucao{O valor $\SI{6.32}{\volt}$ corresponde a $\SI{63}{\percent}$ da tensão
final --- por definição, esse é o instante $t = \tau$. Logo
$\tau = \SI{1}{\milli\second}$. A curva satisfaz $v_C = U_0(1 - e^{-t/\tau})$, e
sua assíntota horizontal é a tensão da fonte, $U_0 = \SI{10}{\volt}$. Lendo o
gráfico: a $\SI{63}{\percent}$ em $\SI{1}{\milli\second}$ e saturando em
\SI{10}{\volt}, confirma-se $\tau = \SI{1}{\milli\second}$.}
\stepcounter{questao}
\resposta{(a) $V_f \approx \SI{66.7}{\volt}$; (b) \SI{50}{\milli\joule} $\to$
\SI{33.3}{\milli\joule}; (c) ver comentário}
\resolucao{\textbf{(a)} A carga total se \emph{conserva} e se redistribui até que
as tensões se igualem:
\[
Q_0 = C_1V_0=(\pot{10}{-6})(100)=\SI{1}{\milli\coulomb},
\]
\[
V_f=\frac{Q_0}{C_1+C_2}=\frac{\pot{1}{-3}}{\pot{15}{-6}}
=\frac{200}{3}\approx\SI{66.7}{\volt}.
\]
\textbf{(b) Energias:}
\[
E_i=\tfrac12 C_1V_0^{\,2}=\tfrac12(\pot{10}{-6})(100)^2=\SI{50}{\milli\joule},
\]
\[
E_f=\tfrac12 (C_1+C_2)V_f^{\,2}
=\tfrac12(\pot{15}{-6})\left(\tfrac{200}{3}\right)^{2}
=\frac{100}{3}\approx\SI{33.3}{\milli\joule}.
\]
Perderam-se $\SI{16.7}{\milli\joule}$ --- exatamente $\dfrac{1}{3}$ da energia
inicial.

\textbf{(c) Resolução do paradoxo.} A carga se conserva, mas a energia
\textbf{não} precisa se conservar \emph{dentro do subsistema dos capacitores} ---
ela foi \emph{dissipada}. Os mecanismos:
\begin{itemize}[leftmargin=1.6em,itemsep=1pt,topsep=2pt]
\item \textbf{Resistência dos fios e da chave:} por menor que seja, a corrente
      transitória a atravessa e gera calor $\int Ri^2\,\mathrm{d}t$. O
      resultado notável é que \emph{a energia perdida independe do valor de $R$}
      --- só o \emph{tempo} do transitório muda.
\item \textbf{Radiação eletromagnética:} no limite $R\to0$, o circuito LC
      oscilaria e irradiaria a energia como ondas.
\end{itemize}
\textbf{Fórmula geral da perda:}
\[
\Delta E=\frac{1}{2}\,\frac{C_1C_2}{C_1+C_2}\,V_0^{\,2}
=\frac12\cdot\frac{10\cdot5}{15}\cdot10^{-6}\cdot 100^2
\approx\SI{16.7}{\milli\joule}.\ \checkmark
\]
\textbf{Relevância prática:} esse é o motivo pelo qual bancos de capacitores nunca
são conectados em paralelo sem resistores de pré-carga (\emph{soft-start}) ---
a corrente de equalização pode ser destrutiva.}
\stepcounter{questao}
\resposta{(b) $Q=\SI{1}{\milli\coulomb}$ (constante), $C=\SI{40}{\micro\farad}$,
$U=\SI{25}{\volt}$, $E$ cai de \SI{50}{\milli\joule} para \SI{12.5}{\milli\joule};
(c) e (d) ver comentário}
\resolucao{\textbf{(a) e (b) --- capacitor isolado ($Q$ constante).}
Desconectado, a carga não tem para onde ir:
\[
Q = C_0U_0 = \pot{10}{-6}\cdot100 = \SI{1}{\milli\coulomb}\ \text{(constante)}.
\]
O dielétrico multiplica a capacitância por $\kappa$:
\[
C = \kappa C_0 = \SI{40}{\micro\farad},
\qquad
U = \frac{Q}{C}=\frac{\pot{1}{-3}}{\pot{40}{-6}}=\SI{25}{\volt}
\quad\left(=\frac{U_0}{\kappa}\right).
\]
Energias:
\[
E_i=\tfrac12 C_0U_0^{\,2}=\SI{50}{\milli\joule},
\qquad
E_f=\tfrac12 CU^{2}=\tfrac12(\pot{40}{-6})(25)^2=\SI{12.5}{\milli\joule}
=\frac{E_i}{\kappa}.
\]
\textbf{(c) Para onde foi a energia?} \emph{Não houve violação alguma.} O
dielétrico é \textbf{atraído} para dentro do capacitor pelo campo de borda ---
existe uma força elétrica real puxando-o. Se ele for solto, essa força realiza
trabalho e o acelera; a energia "perdida" pelo campo converte-se em energia
cinética do dielétrico (e, ao final, em calor por atrito e vibração).
Para inseri-lo \emph{quase-estaticamente}, um agente externo precisa segurá-lo,
e é ele quem \emph{recebe} os \SI{37.5}{\milli\joule}.

\textbf{(d) Caso alternativo --- capacitor ligado à fonte ($U$ constante).}
Agora a tensão é fixada em \SI{100}{\volt}:
\[
C = \SI{40}{\micro\farad},\qquad
Q = CU = \SI{4}{\milli\coulomb}\ (\text{quadruplicou}),\qquad
E_f=\tfrac12 CU^{2}=\SI{200}{\milli\joule}\ (\text{quadruplicou}).
\]
A energia armazenada \emph{aumentou} \SI{150}{\milli\joule}. A fonte, porém,
forneceu $\Delta Q\cdot U = (\pot{3}{-3})(100)=\SI{300}{\milli\joule}$ --- o
dobro. A metade restante (\SI{150}{\milli\joule}) foi entregue como trabalho
mecânico ao dielétrico (que também é sugado para dentro).

\smallskip
\begin{center}\small\sffamily
\begin{tabular}{@{}l c c@{}}
\toprule
 & \textbf{Isolado ($Q$ const.)} & \textbf{Ligado ($U$ const.)}\\
\midrule
Carga & constante & $\times\kappa$\\
Tensão & $\div\kappa$ & constante\\
Energia & $\div\kappa$ (diminui) & $\times\kappa$ (aumenta)\\
\bottomrule
\end{tabular}
\end{center}
\textbf{Moral:} a mesma ação física (inserir o dielétrico) produz resultados
\emph{opostos} para a energia conforme a condição de contorno. Identificar
\emph{qual grandeza permanece constante} é o passo decisivo do problema.}
\stepcounter{questao}
\resposta{$C=\SI{15}{\micro\farad}$}
\resolucao{
\[
C=\frac{Q}{U}=\frac{\pot{300}{-6}}{20}=\pot{15}{-6}\,\si{\farad}
=\SI{15}{\micro\farad}.
\]
A capacitância é uma propriedade \emph{geométrica} do componente: ela não
depende de $Q$ nem de $U$, apenas da razão entre eles. Aplicar o dobro da
tensão dobra a carga e deixa $C$ inalterada.}
\stepcounter{questao}
\resposta{$U\approx\SI{12.6}{\volt}$}
\resolucao{De $W=\frac12CU^{2}$,
\[
U=\sqrt{\frac{2W}{C}}=\sqrt{\frac{2(\pot{8}{-3})}{\pot{100}{-6}}}
=\sqrt{160}=\SI{12.65}{\volt}.
\]
A raiz quadrada indica que a energia é pouco sensível à tensão em valores
baixos e muito sensível em valores altos: dobrar a tensão \emph{quadruplica}
a energia armazenada.}
\stepcounter{questao}
\resposta{(b)}
\resolucao{Um salto de tensão implicaria
$\mathrm{d}u/\mathrm{d}t\to\infty$ e, por
$i=C\,\mathrm{d}u/\mathrm{d}t$, corrente infinita --- impossível. Logo
$u(0^{+})=u(0^{-})$.\\
\textbf{Dualidade com o indutor:} no capacitor é a \emph{tensão} que é
contínua; no indutor, a \emph{corrente}. Em regime CC o capacitor tem corrente
nula (comporta-se como circuito aberto) e tensão em geral não nula --- o que
descarta a alternativa (d).}
\stepcounter{questao}
\resposta{$\tau=\SI{47}{\milli\second}$}
\resolucao{
\[
\tau=RC=4700\cdot\pot{10}{-6}=\pot{4.7}{-2}\,\si{\second}
=\SI{47}{\milli\second}.
\]
\textbf{Atenção à estrutura:} no RC, resistência \emph{maior} torna o
transitório \emph{mais lento} ($\tau=RC$); no RL ocorre o contrário
($\tau=L/R$). Verificar a dimensão ajuda:
$\si{\ohm}\times\si{\farad}=\si{\second}$.}
\stepcounter{questao}
\resposta{$C_{eq}=\SI{20}{\micro\farad}$}
\resolucao{\textbf{Série} (regra inversa, ao contrário dos resistores):
\[
C_s=\frac{10\cdot10}{20}=\SI{5}{\micro\farad}.
\]
\textbf{Paralelo} (soma direta):
\[
C_{eq}=5+15=\SI{20}{\micro\farad}.
\]
Memorize a inversão: capacitores em \emph{série} combinam como resistores em
paralelo, e vice-versa. A razão física é que a série reduz a capacitância
(equivale a afastar as placas), enquanto o paralelo a aumenta (equivale a somar
áreas).}
\stepcounter{questao}
\resposta{$C_{eq}=\SI{0.667}{\micro\farad}$; $Q=\SI{20}{\micro\coulomb}$;
$U_1=\SI{20}{\volt}$, $U_2=\SI{10}{\volt}$}
\resolucao{
\[
C_{eq}=\frac{1\cdot2}{3}=\SI{0.667}{\micro\farad};
\qquad
Q=C_{eq}U=0{,}667\cdot30=\SI{20}{\micro\coulomb}.
\]
Em série, a \textbf{carga é a mesma} em todos (é a mesma corrente que os
carregou):
\[
U_1=\frac{20}{1}=\SI{20}{\volt},
\qquad
U_2=\frac{20}{2}=\SI{10}{\volt}.
\]
\textbf{Verificação:} $20+10=\SI{30}{\volt}$ \checkmark.\\
A tensão se reparte \emph{inversamente} à capacitância --- o capacitor
\textbf{menor} fica com a maior tensão. É o oposto da série de resistores, e é
a causa de um problema de projeto tratado adiante.}
\stepcounter{questao}
\resposta{$i=\SI{5}{\milli\ampere}$}
\resolucao{
\[
i=C\frac{\mathrm{d}u}{\mathrm{d}t}=\pot{10}{-6}\cdot500
=\pot{5}{-3}\,\si{\ampere}=\SI{5}{\milli\ampere}.
\]
Enquanto a rampa durar, a corrente é \emph{constante}. Se a tensão parasse de
variar, a corrente cairia a zero instantaneamente --- o capacitor responde à
\textbf{variação} da tensão, não ao seu valor.}
\stepcounter{questao}
\resposta{Onda \textbf{quadrada} de $\pm\SI{22}{\milli\ampere}$}
\resolucao{Na subida,
$\mathrm{d}u/\mathrm{d}t=10/10^{-3}=\SI{10000}{\volt\per\second}$:
\[
i=2{,}2\times10^{-6}\cdot10^{4}=\SI{22}{\milli\ampere}.
\]
Na descida a taxa inverte e $i=\SI{-22}{\milli\ampere}$.\\
\textbf{Resultado geral:} o capacitor \emph{deriva} a tensão --- triangular
entra, quadrada sai. O indutor faz o oposto (integra): tensão quadrada nele
produz corrente triangular. É a dualidade dos dois elementos em sua forma mais
visível no osciloscópio.}
\stepcounter{questao}
\resposta{$\SI{12.64}{\volt}$, $\SI{17.29}{\volt}$ e $\SI{19.00}{\volt}$}
\resolucao{Com $u_C=U\left(1-e^{-t/\tau}\right)$ e $U=\SI{20}{\volt}$:
\[
u(\tau)=20(0{,}632)=\SI{12.64}{\volt};
\quad
u(2\tau)=20(0{,}865)=\SI{17.29}{\volt};
\quad
u(3\tau)=20(0{,}950)=\SI{19.00}{\volt}.
\]
As frações $63{,}2\%$, $86{,}5\%$ e $95{,}0\%$ são as mesmas do circuito RL ---
elas dependem só da forma exponencial, não do tipo de elemento.}
\stepcounter{questao}
\resposta{$t\approx\SI{2.77}{\milli\second}$}
\resolucao{
\[
15=20\left(1-e^{-t/\tau}\right)
\;\Longrightarrow\;
e^{-t/\tau}=0{,}25
\;\Longrightarrow\;
t=\tau\ln4=2(1{,}386)=\SI{2.77}{\milli\second}.
\]
\textbf{Atalho útil:} atingir $75\%$ leva exatamente $\ln4\approx1{,}39\tau$;
atingir $50\%$ leva $\ln2\approx0{,}69\tau$. Esse último valor é a "meia-vida"
do transitório e vale a pena memorizar.}
\stepcounter{questao}
\resposta{$U_C=\SI{8}{\volt}$; $i_C=0$}
\resolucao{Em regime CC o capacitor é um \textbf{circuito aberto}: não conduz.
Toda a corrente passa pelos dois resistores, e a tensão do capacitor é a de
$R_2$:
\[
U_C=24\cdot\frac{4}{8+4}=\SI{8}{\volt};
\qquad
i_C=0.
\]
\textbf{Método:} para determinar o estado final de qualquer circuito CC com
capacitores, basta \emph{removê-los} do esquema e resolver a rede resistiva
restante. Para indutores, o procedimento dual é substituí-los por curtos.}
\stepcounter{questao}
\resposta{$U_2=\SI{20}{\volt}$}
\resolucao{
\[
\frac12C_1U_1^{2}=\frac12C_2U_2^{2}
\;\Longrightarrow\;
U_2=U_1\sqrt{\frac{C_1}{C_2}}=10\sqrt{4}=\SI{20}{\volt}.
\]
Ambos armazenam $\SI{5}{\milli\joule}$. Reduzir $C$ por um fator $4$ exige
\emph{dobrar} a tensão --- e a tensão é justamente o parâmetro caro e limitado
em capacitores reais, pois exige dielétrico mais espesso, o que reduz $C$.}
\stepcounter{questao}
\resposta{(a) $2C$; (b) $2C$; (c) $4C$}
\resolucao{De $C=\dfrac{\varepsilon A}{d}$: $C\propto A$ e
$C\propto1/d$.\\
(a) $A\to2A$ dá $2C$. (b) $d\to d/2$ dá $2C$. (c) os efeitos se multiplicam:
$4C$.\\
\textbf{Limite prático de (b):} reduzir $d$ aumenta o campo
$E=U/d$ para a mesma tensão. Chega-se rapidamente à rigidez dielétrica do
material, e o capacitor perfura. Há, portanto, um teto para essa estratégia ---
motivo pelo qual capacitores de alta capacitância usam grandes áreas
(folhas enroladas) em vez de dielétricos ultrafinos.}
\stepcounter{questao}
\resposta{$i_0=\SI{22.7}{\milli\ampere}$, decaindo exponencialmente}
\resolucao{No instante inicial a tensão ainda é $\SI{50}{\volt}$
(continuidade), e ela se aplica integralmente ao resistor:
\[
i_0=\frac{50}{2200}=\SI{22.7}{\milli\ampere}.
\]
À medida que o capacitor se descarrega, sua tensão cai e a corrente a acompanha:
$i(t)=i_0e^{-t/\tau}$.\\
\textbf{Observação:} a corrente é máxima no início e a tensão é máxima no
início --- diferentemente do circuito RL, em que a corrente parte de zero. Essa
diferença decorre justamente de qual grandeza é contínua em cada elemento.}
\stepcounter{questao}
\resposta{(a) $\SI{3}{\micro\farad}$; $\SI{75}{\volt}$ e $\SI{25}{\volt}$;
(b) $\SI{66.7}{\volt}$}
\resolucao{(a) $C_{eq}=\dfrac{4\cdot12}{16}=\SI{3}{\micro\farad}$. Como a carga
é comum, $U\propto1/C$:
\[
U_1=100\cdot\frac{12}{16}=\SI{75}{\volt},
\qquad
U_2=100\cdot\frac{4}{16}=\SI{25}{\volt}.
\]
(b) Sob $\SI{100}{\volt}$, $C_1$ já está a $\SI{75}{\volt}$ --- $50\%$ acima do
limite. Impondo $U_1\le\SI{50}{\volt}$:
\[
U_{\max}=50\cdot\frac{16}{12}=\SI{66.7}{\volt}.
\]
(c) \textbf{É contraintuitivo} porque associar dois componentes de
$\SI{50}{\volt}$ em série sugere $\SI{100}{\volt}$ de capacidade --- e isso
\emph{só} é verdade se eles forem \textbf{iguais}. Com valores diferentes, o
menor capacitor absorve a maior parte da tensão e limita todo o conjunto.\\
\textbf{Regra prática:} em série de capacitores, use sempre unidades idênticas
--- e ainda assim adicione resistores de equalização, como se detalha no bloco
de desafio.}
\stepcounter{questao}
\resposta{$40$, $60$ e $\SI{100}{\micro\coulomb}$; total
$\SI{200}{\micro\coulomb}$; $C_{eq}=\SI{10}{\micro\farad}$}
\resolucao{Em paralelo a \textbf{tensão} é comum:
\[
Q_k=C_kU \Rightarrow 40,\ 60,\ \SI{100}{\micro\coulomb};
\qquad
Q_{tot}=\SI{200}{\micro\coulomb}.
\]
\[
C_{eq}=2+3+5=\SI{10}{\micro\farad}
\quad\text{(confere: }200/20=10\text{)}.
\]
\textbf{Comparação com a série:} lá, a \emph{carga} é comum e a tensão se
divide; aqui, a \emph{tensão} é comum e a carga se divide. O paralelo soma
capacitâncias e mantém a tensão de trabalho da menor unidade; a série reduz a
capacitância e (idealmente) soma as tensões. São estratégias opostas para
problemas opostos.}
\stepcounter{questao}
\resposta{(a) $\SI{4}{\volt}$; (b) $\tau=\SI{0.2}{\milli\second}$;
(c) $u_C=4\left(1-e^{-t/0{,}2\,\mathrm{ms}}\right)$}
\resolucao{(a) Com o capacitor aberto em regime,
$U_C=12\cdot\dfrac{3}{9}=\SI{4}{\volt}$.\\
(b) Zerando a fonte, o capacitor vê $R_1\parallel R_2$:
\[
R_{Th}=\frac{6\cdot3}{9}=\SI{2}{\kilo\ohm}
\;\Longrightarrow\;
\tau=R_{Th}C=2000\cdot\pot{100}{-9}=\SI{0.2}{\milli\second}.
\]
(c) $u_C(t)=4\left(1-e^{-t/0{,}0002}\right)$ volts.\\
\textbf{Erro comum:} usar $\tau=R_1C$ ou $\tau=R_2C$. A constante de tempo é
governada pela resistência \emph{vista pelo capacitor} com as fontes zeradas
--- aqui, o paralelo das duas, que é menor que qualquer uma delas. O transitório
é, portanto, mais rápido do que qualquer estimativa ingênua sugeriria.}
\stepcounter{questao}
\resposta{(a) $u_C=10-8e^{-t/\tau}$; (b) $\SI{7.06}{\volt}$;
(c) $\SI{1.6}{\milli\ampere}$}
\resolucao{(a) A forma geral de qualquer transitório de primeira ordem é
\[
u(t)=U_f+\left(U_0-U_f\right)e^{-t/\tau}=10-8e^{-t/\tau}.
\]
(b) $u(5\,\mathrm{ms})=10-8e^{-1}=10-2{,}94=\SI{7.06}{\volt}$.\\
(c) Em $t=0^{+}$ a tensão do capacitor é $\SI{2}{\volt}$, e sobre o resistor
sobram $10-2=\SI{8}{\volt}$:
\[
i_0=\frac{8}{5000}=\SI{1.6}{\milli\ampere}.
\]
\textbf{Vantagem da forma geral:} ela cobre carga ($U_0=0$), descarga
($U_f=0$) e qualquer condição intermediária. Basta identificar os três
parâmetros $U_0$, $U_f$ e $\tau$ --- não há três fórmulas a memorizar.}
\stepcounter{questao}
\resposta{(a) $\tau=\SI{940}{\second}\approx\SI{15.7}{\minute}$;
(b) $\approx\SI{38}{\volt}$}
\resolucao{(a) A fuga é modelada por um resistor em paralelo:
\[
\tau=R_{iso}C=\pot{2}{6}\cdot\pot{470}{-6}=\SI{940}{\second}.
\]
(b) Com $t=\SI{900}{\second}$:
\[
u=100\,e^{-900/940}=100\,e^{-0{,}957}=\SI{38.4}{\volt}.
\]
(c) \textbf{Depois de quinze minutos desligado, o capacitor ainda guarda
$\SI{38}{\volt}$} --- e, em bancos de fonte chaveada com centenas de volts, o
resíduo é letal por horas. É por isso que equipamentos com barramento CC trazem
\emph{resistores de descarga} (bleeders) permanentes e o aviso de aguardar
antes de abrir o gabinete.\\
\textbf{Nota:} capacitores eletrolíticos ainda apresentam \emph{absorção
dielétrica} --- após descarga completa, a tensão "ressurge" parcialmente em
minutos. Curto-circuitar antes de manusear é procedimento obrigatório, não
precaução exagerada.}
\stepcounter{questao}
\resposta{(a) $P=\SI{0.225}{\watt}$}
\resolucao{(a)
\[
P=\mathrm{ESR}\cdot I_{rms}^{2}=0{,}1(1{,}5)^{2}=\SI{0.225}{\watt}.
\]
(b) Essa potência é dissipada \emph{dentro} do capacitor, elevando a
temperatura do eletrólito. A vida útil de eletrolíticos segue uma regra
empírica severa: ela \textbf{cai pela metade a cada $\SI{10}{\celsius}$} de
elevação. Um componente de $\SI{5000}{\hour}$ a $\SI{105}{\celsius}$ dura
$\SI{20000}{\hour}$ a $\SI{85}{\celsius}$ --- e o inverso também vale.
Pior: o eletrólito evapora, a ESR \emph{aumenta}, a dissipação cresce e o
processo se realimenta.\\
(c) \textbf{(i) Dividir a corrente:} usar $n$ capacitores em paralelo reduz a
corrente por unidade a $I/n$ e a dissipação individual a $P/n^{2}$.
Com dois, cada um dissipa $\SI{56}{\milli\watt}$ --- um quarto do original.\\
\textbf{(ii) Escolher tecnologia de baixa ESR:} capacitores de polímero sólido
têm ESR uma ordem de grandeza menor, ao custo de capacitância e tensão
menores.\\
\textbf{Consequência de diagnóstico:} em fontes chaveadas, o capacitor de saída
é quase sempre o primeiro componente a falhar --- e a ESR crescente, não a
perda de capacitância, é o sintoma que primeiro aparece.}
\stepcounter{questao}
\resposta{(a) $\SI{40}{\milli\joule}$; (b) $\SI{20}{\milli\joule}$;
(c) $\SI{20}{\milli\joule}$ --- exatamente metade}
\resolucao{(a) A fonte mantém $\SI{20}{\volt}$ enquanto entrega toda a carga
final $Q=CU=\SI{2}{\milli\coulomb}$:
\[
W_{fonte}=QU=CU^{2}=\pot{100}{-6}(400)=\SI{40}{\milli\joule}.
\]
(b) $W_C=\frac12CU^{2}=\SI{20}{\milli\joule}$.\\
(c) Por conservação, $W_R=40-20=\SI{20}{\milli\joule}$.\\
\textbf{Resultado notável:} o rendimento do carregamento de um capacitor por
fonte de tensão através de resistor é \textbf{sempre $50\%$} --- não importa o
valor de $R$, nem de $C$, nem da tensão. Diminuir $R$ apenas acelera o
processo; a energia perdida é a mesma. A demonstração formal por integração
está no bloco de desafio.\\
\textbf{Consequência prática:} carregadores eficientes não usam resistor em série
--- usam conversores chaveados, que transferem energia em pacotes através de um
indutor e escapam desse limite de $50\%$.}
\stepcounter{questao}
\resposta{(a) $\SI{2}{\kilo\volt}$; (b) $\SI{1}{\kilo\volt}$}
\resolucao{(a) $U_{rup}=E_{rig}\cdot d=20\cdot0{,}1=\SI{2}{\kilo\volt}$.\\
(b) Com margem de $50\%$, especifica-se $\SI{1}{\kilo\volt}$ --- e ainda assim
é prudente derrateá-la em alta temperatura e sob ondulação.\\
(c) Combinando $C=\dfrac{\varepsilon A}{d}$ com
$U_{\max}=E_{rig}d$:
\[
C\,U_{\max}=\varepsilon A E_{rig}
\quad\Longrightarrow\quad
W_{\max}=\frac12CU_{\max}^{2}=\frac12\varepsilon E_{rig}^{2}\,(A\,d).
\]
O termo $Ad$ é o \textbf{volume} de dielétrico. Ou seja, a energia máxima
armazenável depende apenas do volume e das propriedades do material
($\varepsilon$ e $E_{rig}$) --- \emph{não} da geometria escolhida.\\
\textbf{Consequência:} não existe truque geométrico para ganhar energia. Só
avança quem melhora o material. É por isso que o desenvolvimento de capacitores
é, essencialmente, ciência de materiais.}
\stepcounter{questao}
\resposta{$R=\SI{5}{\kilo\ohm}$; $C=\SI{0.6}{\micro\farad}$}
\resolucao{\textbf{Resistência:} em $t=0^{+}$ o capacitor está descarregado e
toda a tensão cai em $R$:
\[
R=\frac{U_f}{i_0}=\frac{10}{\pot{2}{-3}}=\SI{5}{\kilo\ohm}.
\]
\textbf{Constante de tempo:} $6{,}32/10=0{,}632$ --- exatamente a fração de
$1\tau$. Logo $\tau=\SI{3}{\milli\second}$ e
\[
C=\frac{\tau}{R}=\frac{\pot{3}{-3}}{5000}=\pot{6}{-7}\,\si{\farad}
=\SI{0.6}{\micro\farad}.
\]
\textbf{Estrutura da medição:} a corrente inicial isola $R$; o transitório
isola $\tau$; e só a combinação entrega $C$. Nenhuma medida sozinha determina
o capacitor --- é o mesmo padrão do circuito RL, em que o regime permanente dá
$R$ e o transitório dá $L$.}
\stepcounter{questao}
\resposta{Capacitor: curto em $t=0^{+}$ (se descarregado), aberto em regime;
indutor: aberto em $t=0^{+}$ (se sem corrente), curto em regime}
\resolucao{\textbf{Capacitor descarregado em $t=0^{+}$:} $u_C=0$, e um elemento
com tensão nula equivale a um \textbf{curto-circuito}. Em regime CC,
$\mathrm{d}u/\mathrm{d}t=0$ implica $i=0$: equivale a \textbf{circuito
aberto}.\\
\textbf{Indutor sem corrente em $t=0^{+}$:} $i_L=0$, o que equivale a
\textbf{circuito aberto}. Em regime, $\mathrm{d}i/\mathrm{d}t=0$ implica
$v=0$: equivale a \textbf{curto}.\\
\textbf{Atenção à condição inicial:} as equivalências em $t=0^{+}$ pressupõem
estado inicial nulo. Se o capacitor já estiver a $U_0$, ele equivale a uma
\emph{fonte de tensão} $U_0$; se o indutor já conduzir $I_0$, equivale a uma
\emph{fonte de corrente} $I_0$.\\
\textbf{Por que isso é útil:} esses quatro substitutos permitem obter os
valores inicial e final de qualquer transitório de primeira ordem resolvendo
apenas circuitos \emph{resistivos}. Junto com $\tau$, eles determinam a solução
completa sem integrar nada.}
\stepcounter{questao}
\resposta{(a) $f_c\approx\SI{339}{\hertz}$; (b) $f_c=1/(2\pi\tau)$}
\resolucao{(a)
\[
f_c=\frac{1}{2\pi RC}=\frac{1}{2\pi(10^{4})(\pot{47}{-9})}
=\SI{339}{\hertz}.
\]
(b) Como $\tau=RC=\SI{0.47}{\milli\second}$, vale
$f_c=\dfrac{1}{2\pi\tau}$: um circuito \emph{lento} (τ grande) filtra a partir
de frequências \emph{baixas}. As duas descrições --- temporal e frequencial ---
são a mesma informação.\\
(c) \textbf{Muito abaixo de $f_c$:} o capacitor tem reatância alta, quase não
carrega, e o sinal passa praticamente intacto (ganho $\approx1$).\\
\textbf{Muito acima:} o capacitor não tem tempo de carregar entre as variações;
a saída é fortemente atenuada, a $-\SI{20}{\decibel}$ por década.\\
\textbf{Exatamente em $f_c$:} o ganho vale $1/\sqrt2\approx0{,}707$
($-\SI{3}{\decibel}$), e a defasagem é de $\SI{45}{\degree}$. É a definição
convencional de frequência de corte.}
\stepcounter{questao}
\resposta{$u_C=U\left(1-e^{-t/RC}\right)$;
$i=\dfrac{U}{R}e^{-t/RC}$}
\resolucao{(a) Pela LKT, com $i=C\,\mathrm{d}u_C/\mathrm{d}t$:
\[
U=Ri+u_C=RC\frac{\mathrm{d}u_C}{\mathrm{d}t}+u_C
\;\Longrightarrow\;
\frac{\mathrm{d}u_C}{U-u_C}=\frac{\mathrm{d}t}{RC}.
\]
Integrando de $0$ a $t$:
\[
-\ln\!\left(\frac{U-u_C}{U}\right)=\frac{t}{RC}
\;\Longrightarrow\;
u_C(t)=U\left(1-e^{-t/RC}\right).
\]
(b)
\[
i=C\frac{\mathrm{d}u_C}{\mathrm{d}t}
=CU\cdot\frac{1}{RC}e^{-t/RC}=\frac{U}{R}e^{-t/RC}.
\]
\emph{Em $t=0^{+}$:} $u_C=0$ e $Ri=U$; soma $=U$ \checkmark.\\
\emph{Em $t\to\infty$:} $u_C=U$ e $Ri=0$; soma $=U$ \checkmark.\\
(c) \textbf{Dualidade estrutural.} No RL, a \emph{corrente} sobe como
$1-e^{-t/\tau}$ e a tensão do indutor decai como $e^{-t/\tau}$. No RC, é a
\emph{tensão} que sobe e a corrente que decai. Trocando
$L\leftrightarrow C$, $i\leftrightarrow u$ e $R\leftrightarrow1/R$, uma solução
vira a outra --- razão pela qual basta dominar um dos dois casos.}
\stepcounter{questao}
\resposta{$W_R=\frac12CU^{2}$; rendimento sempre $50\%$}
\resolucao{(a) Com $i=\dfrac{U}{R}e^{-t/\tau}$ e $\tau=RC$:
\[
W_R=\int_0^{\infty}R\left(\frac{U}{R}\right)^{2}e^{-2t/\tau}\,\mathrm{d}t
=\frac{U^{2}}{R}\left[-\frac{\tau}{2}e^{-2t/\tau}\right]_0^{\infty}
=\frac{U^{2}\tau}{2R}.
\]
(b) Substituindo $\tau=RC$:
\[
W_R=\frac{U^{2}}{2R}\cdot RC=\frac12CU^{2}.
\]
O $R$ \textbf{cancela}: ele aparece na potência ($\propto1/R$) e desaparece na
duração ($\tau\propto R$). Como a fonte forneceu $CU^{2}$ e o capacitor reteve
$\frac12CU^{2}$, o rendimento é exatamente $50\%$ --- sempre.\\
(c) \textbf{O limite decorre de forçar a tensão do capacitor a saltar de $0$
para $U$ contra uma fonte rígida.} Conversores chaveados evitam isso: um
indutor recebe energia da fonte em regime de corrente, e depois a entrega ao
capacitor. Como o indutor é um elemento reativo (não dissipa), a transferência
pode, em princípio, ser feita sem perda --- e na prática rendimentos de $90\%$
a $95\%$ são rotineiros.\\
\textbf{Observação:} o mesmo resultado de $50\%$ aparece no compartilhamento de
carga entre dois capacitores e na descarga de um indutor. Em todos os casos, a
independência em relação a $R$ tem a mesma origem: potência e duração variam
inversamente.}
\stepcounter{questao}
\resposta{(a) duas séries de dois, em paralelo --- 4 unidades;
(c) $R\le\SI{50}{\kilo\ohm}$, dissipando $\SI{50}{\milli\watt}$ cada}
\resolucao{(a) \textbf{Série de dois:} suporta $\SI{100}{\volt}$, mas a
capacitância cai para $\SI{50}{\micro\farad}$.\\
\textbf{Duas dessas séries em paralelo:} $50+50=\SI{100}{\micro\farad}$,
mantendo $\SI{100}{\volt}$. Total: \textbf{4 capacitores}.\\
Note o custo: quadruplicou-se o número de componentes para dobrar a tensão. É o
preço inevitável --- a energia armazenada, essa sim, quadruplicou.\\
(b) \textbf{Sem equalização, a divisão de tensão em série é instável.} A tensão
de repouso é fixada pelas \emph{correntes de fuga}, não pelas capacitâncias: o
capacitor de menor fuga assume mais tensão. Como a fuga cresce com a tensão e
com a temperatura, o desequilíbrio se realimenta e pode levar um dos capacitores
a ultrapassar $\SI{50}{\volt}$ e perfurar --- o que joga toda a tensão no outro
e o destrói em seguida.\\
(c) Ligue um resistor em paralelo com cada capacitor, dimensionado para que sua
corrente domine a fuga. Adotando um fator $100$:
\[
I_R=100\cdot\pot{10}{-6}=\SI{1}{\milli\ampere}
\;\Longrightarrow\;
R\le\frac{50}{\pot{1}{-3}}=\SI{50}{\kilo\ohm}.
\]
Dissipação: $P=\dfrac{50^{2}}{\pot{50}{3}}=\SI{50}{\milli\watt}$ por resistor
--- especifique $\frac14\,\si{\watt}$.\\
\textbf{Bônus:} esses resistores também funcionam como \emph{bleeders},
descarregando o banco quando o equipamento é desligado. Com
$\tau=RC=50\,\mathrm{k}\times100\,\mu=\SI{5}{\second}$, o banco fica seguro em
menos de meio minuto.}
\stepcounter{questao}
\resposta{(a) $R=\SI{15.9}{\kilo\ohm}\to\SI{16}{\kilo\ohm}$;
(b) $f_c=\SI{99.5}{\hertz}$, erro $0{,}5\%$;
(c) $\approx-\SI{20}{\decibel}$}
\resolucao{(a)
\[
R=\frac{1}{2\pi f_cC}=\frac{1}{2\pi(100)(\pot{100}{-9})}
=\SI{15915}{\ohm}.
\]
O valor E24 mais próximo é $\SI{16}{\kilo\ohm}$.\\
(b) $f_c=\dfrac{1}{2\pi(16000)(\pot{100}{-9})}=\SI{99.5}{\hertz}$ --- erro de
$0{,}5\%$, muito abaixo da tolerância do próprio capacitor (tipicamente
$\pm10\%$ ou pior). \textbf{O capacitor, e não o resistor, domina a incerteza
do projeto.}\\
(c) Uma década acima de $f_c$, o ganho é
\[
|H|=\frac{1}{\sqrt{1+(f/f_c)^{2}}}\approx\frac{1}{10}
\;\Longrightarrow\;
-\SI{20}{\decibel}.
\]
\textbf{Limitação:} $\SI{20}{\decibel}$ por década é pouco. Para atenuar
$\SI{60}{\decibel}$ seria preciso $f=\SI{100}{\kilo\hertz}$ --- três décadas
acima. Se a especificação exigir corte mais abrupto, é necessário aumentar a
\emph{ordem} do filtro (cascatear seções ou usar topologia ativa), não
apertar $R$ e $C$.\\
\textbf{Cuidado ao cascatear:} duas seções RC ligadas diretamente
\emph{carregam} uma à outra e o $f_c$ resultante não é o previsto. É preciso
desacoplá-las com um buffer, ou projetar a rede como um todo.}
\stepcounter{questao}
\resposta{(a) $\dfrac{\Delta\tau}{\tau}=\dfrac{\Delta R}{R}+\dfrac{\Delta C}{C}$;
(b) $21\%$ (pior caso) e $20{,}0\%$ (estatístico)}
\resolucao{(a) Como $\tau=RC$ é um \emph{produto}, tome o logaritmo e
diferencie:
\[
\ln\tau=\ln R+\ln C
\;\Longrightarrow\;
\frac{\Delta\tau}{\tau}=\frac{\Delta R}{R}+\frac{\Delta C}{C}.
\]
As incertezas \textbf{relativas} se somam --- diferentemente da soma, em que se
somam as absolutas.\\
(b) \emph{Pior caso:} $1\%+20\%=21\%$.\\
\emph{Estatístico:} $\sqrt{1^{2}+20^{2}}=\sqrt{401}=20{,}02\%$.\\
(c) \textbf{Os dois critérios praticamente coincidem} --- e isso é a
informação relevante. Quando uma parcela domina de forma tão desproporcional,
combinar em quadratura não ajuda: o resultado é essencialmente a maior
tolerância isolada.\\
\textbf{Conclusão de projeto:} trocar o resistor de $1\%$ por um de $0{,}1\%$
não muda nada ($20{,}00\%$ contra $20{,}02\%$). Todo o ganho está no capacitor.
Se a aplicação exigir temporização precisa, é preciso abandonar o eletrolítico
--- filme de poliéster oferece $\pm5\%$, e polipropileno chega a $\pm1\%$ com
excelente estabilidade térmica.\\
\textbf{Regra geral:} identifique a maior contribuição \emph{antes} de
especificar componentes. Precisão gasta em parcelas menores é dinheiro
jogado fora.}
\stepcounter{questao}
\resposta{$F=\dfrac{\varepsilon_0AU^{2}}{2d^{2}}\approx\SI{44}{\milli\newton}$,
atrativa}
\resolucao{(a) Com o capacitor \textbf{isolado}, $Q$ é constante e convém usar
\[
W=\frac{Q^{2}}{2C}=\frac{Q^{2}d}{2\varepsilon_0A}.
\]
A força que o sistema exerce é
\[
F=-\frac{\mathrm{d}W}{\mathrm{d}d}=-\frac{Q^{2}}{2\varepsilon_0A},
\]
negativa, isto é, \textbf{atrativa} --- as placas tendem a se aproximar,
reduzindo a energia.\\
(b) Substituindo $Q=CU=\dfrac{\varepsilon_0AU}{d}$:
\[
F=\frac{\varepsilon_0AU^{2}}{2d^{2}}\quad\text{(em módulo)}.
\]
\textbf{Sutileza importante:} a tensão constante, a energia \emph{aumenta} ao
aproximar as placas ($W=\frac12CU^{2}$ com $C$ crescendo), o que sugeriria
força repulsiva. O paradoxo se resolve incluindo a bateria: ela fornece
$U\,\mathrm{d}Q$, o dobro da variação de energia do campo. Feito o balanço
completo, a força é atrativa e de \emph{mesmo módulo} --- como tem de ser, pois
a força é uma grandeza física e não depende do que se escolhe manter
constante.\\
(c)
\[
F=\frac{(\pot{8.85}{-12})(\pot{1}{-2})(10^{6})}{2(10^{-6})}
=\SI{44.3}{\milli\newton}.
\]
\textbf{Ordem de grandeza:} cerca de $\SI{4.5}{\gram}$-força para
$\SI{1}{\kilo\volt}$ --- pequena, mas suficiente para fazer as placas vibrarem
audivelmente sob tensão alternada. É o "zumbido" de capacitores cerâmicos
classe 2, e o princípio de funcionamento de microfones e atuadores
eletrostáticos de MEMS.}
\stepcounter{questao}
\resposta{(a) $\tau=\SI{0.83}{\second}$; (b) $C=\SI{100}{\micro\farad}$ e
$R=\SI{8.2}{\kilo\ohm}$; (c) erro dominado pelo capacitor}
\resolucao{(a) De $0{,}70=1-e^{-t/\tau}$:
\[
e^{-t/\tau}=0{,}30
\;\Longrightarrow\;
\tau=\frac{t}{\ln(1/0{,}30)}=\frac{1}{1{,}204}=\SI{0.83}{\second}.
\]
(b) Com $C=\SI{100}{\micro\farad}$:
$R=\tau/C=0{,}83/\pot{100}{-6}=\SI{8306}{\ohm}$. O valor E24 mais próximo é
$\SI{8.2}{\kilo\ohm}$, que dá $\tau=\SI{0.82}{\second}$ e
\[
t_{70\%}=0{,}82(1{,}204)=\SI{0.987}{\second}\quad(\text{erro de }1{,}3\%).
\]
(c) \textbf{O erro de $1{,}3\%$ do arredondamento é irrelevante} diante da
tolerância do eletrolítico ($\pm20\%$), que se transfere integralmente para
$\tau$: o tempo real ficará entre $0{,}79$ e $\SI{1.19}{\second}$. Some-se a
deriva térmica (vários por cento entre $0$ e $\SI{60}{\celsius}$) e o
envelhecimento.\\
\textbf{Conclusão:} um temporizador RC com eletrolítico serve para "cerca de um
segundo", não para "um segundo". Para precisão, as alternativas são capacitor
de filme (levando a $R$ muito maior e $C$ menor) ou --- o caminho moderno ---
um contador digital com oscilador a cristal, cuja estabilidade é de partes por
milhão.\\
\textbf{Observação de método:} note que o requisito foi dado em $70\%$, e não
em $\tau$. Confundir os dois levaria a um erro de $20\%$ no projeto, pois
$t_{70\%}=1{,}204\,\tau$.}
\stepcounter{questao}
\resposta{$u=U_0e^{-t/R_fC}$; $\tau=\SI{100}{\second}$; após
$\SI{5}{\minute}$ restam $\approx\SI{15}{\volt}$}
\resolucao{(a) Em circuito aberto, a única corrente é a de fuga:
\[
C\frac{\mathrm{d}u}{\mathrm{d}t}=-\frac{u}{R_f}
\;\Longrightarrow\;
u(t)=U_0e^{-t/R_fC}.
\]
(b) $\tau_f=R_fC$. Catálogos raramente informam $R_f$; informam a
\emph{corrente de fuga} $I_f$ a uma dada tensão. A conversão é imediata:
$R_f=U/I_f$, e portanto
\[
\tau_f=\frac{U\,C}{I_f}.
\]
Note que $\tau_f$ depende do \emph{produto} $R_fC$, e em eletrolíticos $R_f$
cai quando $C$ cresce --- razão pela qual é a autodescarga de capacitores grandes
não é proporcionalmente mais lenta.\\
(c) $\tau_f=\pot{100}{3}\cdot\pot{1000}{-6}=\SI{100}{\second}$. Após
$\SI{300}{\second}$:
\[
u=300\,e^{-3}=300(0{,}0498)=\SI{14.9}{\volt}.
\]
\textbf{Leitura de segurança:} cinco minutos depois de desligado, o barramento
ainda guarda $\SI{15}{\volt}$ --- e, o que é pior, $\SI{0.11}{\joule}$
concentrados. Aos $\SI{60}{\second}$ ainda havia $\SI{110}{\volt}$, plenamente
letais.\\
\textbf{Por isso o resistor de descarga é obrigatório}, e não opcional: com um
bleeder de $\SI{10}{\kilo\ohm}$, $\tau$ cai para $\SI{10}{\second}$ e o
equipamento fica seguro em menos de um minuto, ao custo de
$\SI{9}{\watt}$ apenas enquanto energizado.}
\stepcounter{questao}
\resposta{(a) $\SI{0.2}{\watt}$ e $\approx\SI{100}{\milli\volt}$ de pico;
(b) $\approx\SI{5}{\milli\volt}$ --- a ESR domina}
\resolucao{(a) $P=\mathrm{ESR}\cdot I_{rms}^{2}=0{,}05(4)=\SI{0.2}{\watt}$.\\
A ondulação de tensão que a ESR produz acompanha a corrente instantaneamente:
\[
\Delta u_{ESR}\approx\mathrm{ESR}\cdot I_{pico}\approx0{,}05(2)=\SI{100}{\milli\volt}.
\]
(b) A parcela capacitiva depende da carga transferida por semiciclo. Tomando
$\Delta t\approx1/(2f)=\SI{5}{\micro\second}$:
\[
\Delta u_C\approx\frac{I\,\Delta t}{C}
=\frac{2(\pot{5}{-6})}{\pot{470}{-6}}\approx\SI{21}{\milli\volt}.
\]
(c) \textbf{A ESR domina} --- ela responde por cerca de $80\%$ da ondulação
total. Aumentar $C$ para $\SI{1000}{\micro\farad}$ reduziria $\Delta u_C$ pela
metade e deixaria a ondulação praticamente inalterada.\\
\textbf{Critério correto de seleção:} em fontes chaveadas de alta frequência,
escolhe-se o capacitor pela \textbf{ESR e pela corrente de ondulação
admissível}, não pela capacitância. É comum usar vários capacitores menores em
paralelo --- não para somar $C$, mas para dividir a ESR por $n$ e a dissipação
individual por $n^{2}$.\\
\textbf{Nota:} em frequências baixas ($\SI{120}{\hertz}$ de retificação de
rede), a relação se inverte e a capacitância volta a ser o critério dominante.
O parâmetro decisivo depende da faixa de operação.}
\stepcounter{questao}
\resposta{(a) $\SI{10.9}{\kilo\joule}=\SI{3.04}{\watt\hour}$, ou
$\SI{6.1}{\watt\hour\per\kilogram}$; (b) cerca de $40$ vezes menos}
\resolucao{(a)
\[
W=\frac12CU^{2}=\frac12(3000)(2{,}7)^{2}=\SI{10935}{\joule}
=\frac{10935}{3600}=\SI{3.04}{\watt\hour}.
\]
Densidade: $3{,}04/0{,}50=\SI{6.1}{\watt\hour\per\kilogram}$.\\
(b) A bateria armazena cerca de \textbf{40 vezes mais} energia por quilograma.
Para igualar uma célula de íon-lítio de $\SI{20}{\gram}$, seriam necessários
$\SI{0.8}{\kilogram}$ de supercapacitores.\\
(c) \textbf{Potência específica e vida útil.} O supercapacitor pode ser
carregado e descarregado em \emph{segundos}, com correntes de centenas de
ampères, porque armazena energia \textbf{eletrostaticamente} --- sem reações
químicas, sem difusão iônica lenta, sem degradação de eletrodos. Suporta
centenas de milhares de ciclos, contra alguns milhares da bateria, e funciona
bem a baixas temperaturas.\\
\textbf{Consequência de aplicação:} baterias são \emph{reservatórios} de
energia; supercapacitores são \emph{amortecedores} de potência. Por isso
convivem em veículos elétricos --- o supercapacitor absorve o pico da frenagem
regenerativa e entrega o pico da aceleração, poupando a bateria justamente dos
regimes que mais a degradam.\\
\textbf{Ressalva:} a tensão do supercapacitor cai linearmente com a carga
(pois $Q=CU$), ao contrário do platô da bateria. Aproveitar mais de $75\%$ da
energia exige um conversor CC-CC de entrada ampla.}

\gabarito[10.1 Capacitores]

\subsection{Indutores}
\stepcounter{questao}
\resposta{(b)}
\resolucao{Como $v_L = L\dfrac{\mathrm{d}i}{\mathrm{d}t}$ e, em regime permanente,
a corrente é constante ($\mathrm{d}i/\mathrm{d}t = 0$), a tensão no indutor é
\emph{nula}: ele se comporta como um \textbf{curto-circuito}. (É o oposto do
capacitor, que vira circuito aberto.)}
\stepcounter{questao}
\resposta{$L \approx \SI{0.52}{\milli\henry}$}
\resolucao{Com $A = \SI{5}{\centi\meter\squared} = \pot{5}{-4}\,\si{\meter\squared}$:
\[
L=\frac{\mu_0 N^2 A}{\ell}
=\frac{(\pot{4\pi}{-7})(500)^2(\pot{5}{-4})}{0{,}3}
\approx\pot{5.24}{-4}\,\si{\henry}=\SI{0.52}{\milli\henry}.
\]}
\stepcounter{questao}
\resposta{$v_L = \SI{1}{\volt}$}
\resolucao{$v_L = L\dfrac{\Delta i}{\Delta t}
= 0{,}1\cdot\dfrac{2}{0{,}2} = \SI{1}{\volt}$.}
\stepcounter{questao}
\resposta{(c)}
\resolucao{Indutores seguem a mesma regra dos resistores. Em paralelo:
\[
\Leq=\frac{L_1 L_2}{L_1+L_2}=\frac{3\cdot9}{12}=\frac{27}{12}=\SI{2.25}{\henry}.
\]}
\stepcounter{questao}
\resposta{(a) $\tau = \SI{10}{\milli\second}$; (b) \SI{0.12}{\ampere};
(c) $\approx\SI{75.8}{\milli\ampere}$}
\resolucao{(a) $\tau = \dfrac{L}{R} = \dfrac{1}{100} = \SI{10}{\milli\second}$.\\
(b) Em regime permanente, o indutor vira curto e a corrente é limitada só por $R$:
$I_\infty = \dfrac{12}{100} = \SI{0.12}{\ampere}$.\\
(c) $i(\tau) = I_\infty(1 - e^{-1}) = 0{,}12\cdot0{,}632
\approx\SI{75.8}{\milli\ampere}$ --- cerca de $\SI{63}{\percent}$ do valor final,
o mesmo comportamento do capacitor na carga.}
\stepcounter{questao}
\resposta{(a) \SI{5}{\milli\henry}; (b) \SI{1.2}{\milli\henry}}
\resolucao{(a) Série: $\Leq = L_1 + L_2 = 2 + 3 = \SI{5}{\milli\henry}$.\\
(b) Paralelo: $\Leq = \dfrac{2\cdot3}{5} = \SI{1.2}{\milli\henry}$.}
\stepcounter{questao}
\resposta{$L_{\mathrm{eq}}=\SI{4.5}{\henry}$}
\resolucao{Os dois indutores iguais em paralelo resultam em
$L_p=5/2=\SI{2.5}{\henry}$. Em série com $\SI{2}{\henry}$:
$L_{\mathrm{eq}}=2{,}5+2=\SI{4.5}{\henry}$.}
\stepcounter{questao}
\resposta{Ver comentário}
\resolucao{A tensão do capacitor está ligada à sua carga por $v = Q/C$; mudar $v$
instantaneamente exigiria corrente infinita ($i = C\,\mathrm{d}v/\mathrm{d}t$),
o que é fisicamente impossível. De modo \emph{dual}, a corrente do indutor não
pode saltar, pois isso exigiria tensão infinita
($v = L\,\mathrm{d}i/\mathrm{d}t$). Por isso o capacitor "segura" a tensão e o
indutor "segura" a corrente --- é essa continuidade que dá origem às respostas
exponenciais com constante de tempo $\tau$. A tabela evidencia a bela
\textbf{dualidade} entre os dois componentes: trocando tensão$\leftrightarrow$corrente
e $C\leftrightarrow L$, uma descrição vira a outra.}
\stepcounter{questao}
\resposta{(a) $E=\SI{4}{\joule}$; (c) $\approx\SI{2000}{\volt}$; (b) e (d) ver comentário}
\resolucao{\textbf{(a)} $E = \tfrac12 L I^2 = \tfrac12(0{,}5)(4)^2
= \SI{4}{\joule}$ --- armazenados no \emph{campo magnético}.

\textbf{(b) Origem do arco.} A corrente em um indutor \textbf{não pode variar
instantaneamente}, pois isso exigiria tensão infinita:
$v_L = L\,\dfrac{\mathrm{d}i}{\mathrm{d}t}$. Ao abrir a chave, tenta-se impor
$\mathrm{d}i/\mathrm{d}t \to \infty$; o indutor reage gerando uma tensão
altíssima, de polaridade \emph{invertida}, que tenta manter a corrente
circulando. Essa tensão ioniza o ar entre os contatos e a corrente continua
fluindo pelo \textbf{arco} --- que é o caminho que o circuito "encontra" para
dissipar os \SI{4}{\joule} armazenados.

\textbf{(c) Estimativa da sobretensão.} Com $\Delta i = \SI{4}{\ampere}$ em
$\Delta t = \SI{1}{\milli\second}$:
\[
|v_L| = L\,\frac{\Delta i}{\Delta t}
= 0{,}5\cdot\frac{4}{\pot{1}{-3}} = \SI{2000}{\volt}.
\]
Duzentas vezes maior que uma alimentação típica de \SI{12}{\volt}! Se a abertura
fosse dez vezes mais rápida, a tensão seria dez vezes maior.

\textbf{(d) Proteção --- o diodo de roda livre (\emph{freewheeling}).}
Liga-se um diodo em \textbf{antiparalelo} com o indutor (catodo no positivo).
Em operação normal ele fica reversamente polarizado e não conduz. No instante da
abertura, a tensão do indutor inverte, o diodo entra em condução e oferece um
caminho fechado para a corrente, que decai suavemente com constante de tempo
$\tau = L/R$, dissipando a energia no próprio circuito em vez de no arco.\\
Alternativas: \emph{snubber} RC, varistor (MOV) ou diodo Zener --- este último
quando se deseja desmagnetização \emph{rápida}, aceitando uma sobretensão
controlada.\\
\textbf{Onde isso aparece na prática:} todo relé, contator, solenoide ou motor
CC acionado por transistor precisa desse diodo. Sem ele, o transistor de
acionamento é destruído na primeira comutação --- um dos erros mais comuns de
quem começa em eletrônica de potência.}
\stepcounter{questao}
\resposta{$i_L(0^-)=i_L(0^+)=\SI{2}{\milli\ampere}$; $\tau=\SI{1}{\micro\second}$;
$i_L(t)=\SI{2}{\milli\ampere}e^{-t/(\SI{1}{\micro\second})}$; $|v_R(0^+)|=\SI{4}{\volt}$}
\resolucao{Em regime permanente, o indutor é curto-circuito: $i=12/6000=\SI{2}{\milli\ampere}$.
A corrente do indutor é contínua na comutação. Na descarga, $\tau=L/R=2\,\mathrm{mH}/2\,\mathrm{k\Omega}=\SI{1}{\micro\second}$.
A tensão inicial no resistor vale $Ri=2000\times0{,}002=\SI{4}{\volt}$; sua polaridade é tal que mantém a corrente.}
\stepcounter{questao}
\resposta{$L=\SI{0.2}{\henry}$}
\resolucao{Isolando $L$ na relação fundamental:
\[
v=L\frac{\mathrm{d}i}{\mathrm{d}t}
\;\Longrightarrow\;
L=\frac{v}{\mathrm{d}i/\mathrm{d}t}=\frac{50}{250}=\SI{0.2}{\henry}.
\]
O henry é, portanto, volt-segundo por ampère: $\SI{1}{\henry}$ é a indutância
que produz $\SI{1}{\volt}$ para uma variação de $\SI{1}{\ampere\per\second}$.}
\stepcounter{questao}
\resposta{$W=\SI{5}{\joule}$}
\resolucao{
\[
W=\frac12LI^{2}=\frac12(0{,}4)(5)^{2}=\SI{5}{\joule}.
\]
A dependência é \emph{quadrática} na corrente: dobrar $I$ quadruplica a
energia. Note que, embora a tensão sobre o indutor ideal seja nula em regime CC
(corrente constante), a energia armazenada não é --- ela ficou guardada no
campo durante a fase de crescimento da corrente.}
\stepcounter{questao}
\resposta{(b)}
\resolucao{Um salto de corrente significaria
$\mathrm{d}i/\mathrm{d}t\to\infty$ e, por $v=L\,\mathrm{d}i/\mathrm{d}t$, uma
tensão infinita --- fisicamente impossível. Logo a corrente no indutor é uma
função \textbf{contínua} do tempo: $i(0^{+})=i(0^{-})$.\\
Cuidado com (c) e (d): a corrente pode variar (e varia, nos transitórios), e em
CC ela é constante e em geral \emph{não nula}. O que é nula em CC é a
\emph{tensão}.}
\stepcounter{questao}
\resposta{$\tau=\SI{4}{\milli\second}$}
\resolucao{
\[
\tau=\frac{L}{R}=\frac{0{,}200}{50}=\SI{4e-3}{\second}=\SI{4}{\milli\second}.
\]
\textbf{Atenção à estrutura:} no circuito RL, $\tau=L/R$ --- resistência
\emph{maior} torna o transitório \emph{mais rápido}. No RC é o oposto
($\tau=RC$). Trocar as duas fórmulas é o erro mais comum entre os dois
elementos armazenadores.}
\stepcounter{questao}
\resposta{$\lambda=\SI{2}{\weber}$}
\resolucao{
\[
\lambda=LI=0{,}25\cdot8=\SI{2}{\weber}.
\]
O fluxo concatenado $\lambda=N\Phi$ é o produto do fluxo por espira pelo número
de espiras. A indutância é justamente a constante de proporcionalidade entre
$\lambda$ e $I$ --- é essa a definição primária de $L$, da qual
$v=L\,\mathrm{d}i/\mathrm{d}t$ decorre pela lei de Faraday.}
\stepcounter{questao}
\resposta{$\mathrm{d}i/\mathrm{d}t=\SI{200}{\ampere\per\second}$}
\resolucao{
\[
\frac{\mathrm{d}i}{\mathrm{d}t}=\frac{v}{L}=\frac{30}{0{,}15}
=\SI{200}{\ampere\per\second}.
\]
No instante inicial a corrente é nula, mas sua \emph{taxa} é máxima. Se a
tensão se mantivesse constante e não houvesse resistência, a corrente cresceria
indefinidamente em rampa --- é a resistência do circuito que estabelece o
patamar final.}
\stepcounter{questao}
\resposta{(a) $\SI{2.25}{\joule}$; (b) $\SI{30}{\volt}$}
\resolucao{(a) $W=\frac12(0{,}5)(9)=\SI{2.25}{\joule}$.\\
(b) A taxa é $\mathrm{d}i/\mathrm{d}t=(0-3)/0{,}050
=\SI{-60}{\ampere\per\second}$, e
\[
|v|=L\left|\frac{\mathrm{d}i}{\mathrm{d}t}\right|=0{,}5\cdot60=\SI{30}{\volt}.
\]
O sinal negativo indica \textbf{inversão de polaridade}: o indutor passa a
funcionar como fonte, devolvendo ao circuito os $\SI{2.25}{\joule}$
armazenados. É esse comportamento que produz as sobretensões de abertura.}
\stepcounter{questao}
\resposta{$\SI{20}{\volt}$, com polaridade \emph{invertida} em relação à fase
de crescimento}
\resolucao{
\[
|v|=L\left|\frac{\Delta i}{\Delta t}\right|
=0{,}100\cdot\frac{4}{0{,}020}=\SI{20}{\volt}.
\]
Durante o \emph{crescimento} da corrente, o indutor absorve energia e sua
tensão se opõe ao aumento. Durante o \emph{decrescimento}, ele devolve energia
e a tensão inverte, tentando sustentar a corrente.\\
Em ambos os casos vale a lei de Lenz: o indutor se opõe à \emph{variação}, não
ao sentido da corrente.}
\stepcounter{questao}
\resposta{Série: $\SI{30}{\milli\henry}$; paralelo: $\SI{3}{\milli\henry}$}
\resolucao{\textbf{Série:} $L_{eq}=6+12+12=\SI{30}{\milli\henry}$.\\
\textbf{Paralelo:}
\[
\frac{1}{L_{eq}}=\frac16+\frac1{12}+\frac1{12}=\frac{2+1+1}{12}=\frac13
\;\Longrightarrow\;
L_{eq}=\SI{3}{\milli\henry}.
\]
As regras são \emph{idênticas} às dos resistores --- e opostas às dos
capacitores. A ressalva é a hipótese de \textbf{não acoplamento}: se houver
fluxo mútuo, as fórmulas mudam (ver bloco de aprofundamento).}
\stepcounter{questao}
\resposta{(a) $\tau=\SI{20}{\milli\second}$, $I_f=\SI{0.8}{\ampere}$;
(b) $\SI{24}{\volt}$}
\resolucao{(a) $\tau=L/R=0{,}6/30=\SI{20}{\milli\second}$;
$I_f=V/R=24/30=\SI{0.8}{\ampere}$.\\
(b) Em $t=0^{+}$ a corrente ainda é nula (continuidade), logo o resistor não
tem queda e \textbf{toda} a tensão da fonte aparece sobre o indutor:
$v_L(0^{+})=\SI{24}{\volt}$.\\
\textbf{Os dois extremos:} em $t=0^{+}$ o indutor se comporta como circuito
\emph{aberto} (corrente nula); em $t\to\infty$, como \emph{curto}
(tensão nula). Reconhecer esses dois limites resolve a maioria dos problemas
sem integrar nada.}
\stepcounter{questao}
\resposta{$\SI{0.506}{\ampere}$, $\SI{0.692}{\ampere}$ e
$\SI{0.760}{\ampere}$}
\resolucao{Com $i(t)=I_f\left(1-e^{-t/\tau}\right)$ e
$I_f=\SI{0.8}{\ampere}$:
\[
i(\tau)=0{,}8(1-e^{-1})=0{,}8(0{,}632)=\SI{0.506}{\ampere},
\]
\[
i(2\tau)=0{,}8(0{,}865)=\SI{0.692}{\ampere},
\qquad
i(3\tau)=0{,}8(0{,}950)=\SI{0.760}{\ampere}.
\]
As frações $63{,}2\%$, $86{,}5\%$ e $95{,}0\%$ são universais --- valem para
qualquer RL ou RC. Após $5\tau$ atinge-se $99{,}3\%$, e por convenção o
transitório é dado por encerrado.}
\stepcounter{questao}
\resposta{$\Delta i=\SI{2}{\ampere}$}
\resolucao{Durante o semiciclo positivo, a tensão constante produz rampa de
corrente:
\[
\Delta i=\frac{v\,\Delta t}{L}=\frac{10\cdot0{,}001}{0{,}005}=\SI{2}{\ampere}.
\]
No semiciclo negativo a corrente desce pela mesma quantidade, e o regime é
simétrico.\\
\textbf{Resultado geral:} tensão quadrada no indutor produz corrente
\textbf{triangular} --- o indutor \emph{integra} a tensão. O capacitor faz o
oposto: corrente quadrada nele produz tensão triangular.}
\stepcounter{questao}
\resposta{(a) $\SI{3}{\ampere}$; (b) $W=\SI{2.25}{\joule}$ e
$P=\SI{72}{\watt}$}
\resolucao{(a) Em regime, $\mathrm{d}i/\mathrm{d}t=0$ e o indutor ideal não cai
tensão: só $r$ limita a corrente.
\[
I=\frac{24}{8}=\SI{3}{\ampere}.
\]
(b) $W=\frac12(0{,}5)(9)=\SI{2.25}{\joule}$;
$P=rI^{2}=8(9)=\SI{72}{\watt}$.\\
\textbf{Contraste essencial:} a energia é \emph{armazenada} e recuperável; a
potência é \emph{dissipada} continuamente e perdida. Manter esses
$\SI{2.25}{\joule}$ no campo custa $\SI{72}{\watt}$ permanentes --- razão pela
qual eletroímãs de grande porte usam supercondutores, em que $r=0$ e a
manutenção do campo é gratuita.}
\stepcounter{questao}
\resposta{$I_2=\SI{2}{\ampere}$}
\resolucao{
\[
\frac12L_1I_1^{2}=\frac12L_2I_2^{2}
\;\Longrightarrow\;
I_2=I_1\sqrt{\frac{L_1}{L_2}}=1\cdot\sqrt{4}=\SI{2}{\ampere}.
\]
Ambos armazenam $\SI{1}{\joule}$. Reduzir $L$ por um fator $4$ exige
\emph{dobrar} a corrente --- e, como as perdas vão com $I^{2}$, o indutor menor
custa quatro vezes mais em dissipação para armazenar o mesmo tanto. É o
compromisso central no projeto de indutores de potência.}
\stepcounter{questao}
\resposta{$4L$}
\resolucao{Para um solenoide longo,
\[
L=\frac{\mu N^{2}A}{\ell},
\]
de modo que $L\propto N^{2}$. Dobrar $N$ quadruplica $L$.\\
\textbf{Por que o quadrado:} dobrar as espiras dobra o campo produzido por
uma dada corrente \emph{e} dobra o número de espiras que concatenam esse campo.
Os dois efeitos se multiplicam.\\
\textbf{Ressalva prática:} com o mesmo espaço de bobina, dobrar $N$ exige fio
de seção menor, o que aumenta $r$ e as perdas. Na prática, o ganho de $L$ vem
acompanhado de penalidade em corrente admissível.}
\stepcounter{questao}
\resposta{$L=\SI{1}{\henry}$}
\resolucao{Como $L\propto\mu=\mu_r\mu_0$:
\[
L=\mu_rL_0=500\cdot2=\SI{1000}{\milli\henry}=\SI{1}{\henry}.
\]
\textbf{Três limitações do núcleo.} \emph{(i)} \textbf{Saturação:} acima de
certa corrente, $\mu_r$ despenca e $L$ colapsa. \emph{(ii)} \textbf{Perdas:}
histerese e correntes de Foucault dissipam energia proporcionalmente à
frequência. \emph{(iii)} \textbf{Não linearidade:} $L$ deixa de ser constante,
e as fórmulas lineares passam a ser aproximações válidas apenas em torno de um
ponto de operação.\\
Núcleos de ar têm $L$ minúsculo, mas são perfeitamente lineares e sem perdas
magnéticas --- por isso dominam aplicações de radiofrequência.}
\stepcounter{questao}
\resposta{$t\approx\SI{19.6}{\milli\second}$}
\resolucao{De $i=I_f(1-e^{-t/\tau})$, isole a exponencial:
\[
e^{-t/\tau}=1-\frac{0{,}5}{0{,}8}=0{,}375
\;\Longrightarrow\;
t=\tau\ln\!\left(\frac{1}{0{,}375}\right)=0{,}020\cdot0{,}981
=\SI{19.6}{\milli\second}.
\]
Ou seja, praticamente $1\tau$ --- coerente com o fato de que
$\SI{0.5}{\ampere}$ corresponde a $62{,}5\%$ de $I_f$, muito próximo dos
$63{,}2\%$ atingidos em exatamente $\tau$.}
\stepcounter{questao}
\resposta{$\tau=\SI{20}{\milli\second}$; $i=\SI{0.27}{\ampere}$}
\resolucao{$\tau=L/R=0{,}4/20=\SI{20}{\milli\second}$.\\
Na descarga, a corrente decai exponencialmente a partir do valor inicial:
\[
i(t)=I_0e^{-t/\tau}=2\,e^{-40/20}=2e^{-2}=2(0{,}135)=\SI{0.27}{\ampere}.
\]
\textbf{Continuidade:} no instante da comutação a corrente permanece
$\SI{2}{\ampere}$ --- ela não pode saltar. O que muda instantaneamente é a
\emph{tensão}, que assume $-RI_0=\SI{-40}{\volt}$ para manter essa corrente
circulando pelo resistor.}
\stepcounter{questao}
\resposta{(a) $\SI{2.5}{\joule}$; (b) conservação de energia; (c) $R$ altera a
\emph{duração} e a \emph{potência}, não a energia total}
\resolucao{(a) Toda a energia armazenada acaba no resistor:
\[
W=\frac12LI_0^{2}=\frac12(0{,}2)(25)=\SI{2.5}{\joule}.
\]
(b) O indutor ideal não dissipa e o circuito é fechado: a energia magnética não
tem outro destino. A independência em relação a $R$ é consequência direta da
conservação --- e é confirmada por integração no bloco de desafio.\\
(c) $R$ governa $\tau=L/R$ e, portanto, o \emph{ritmo}:
\begin{itemize}
\item $R$ pequeno $\Rightarrow$ $\tau$ grande: descarga lenta, potência baixa e
prolongada.
\item $R$ grande $\Rightarrow$ $\tau$ pequeno: descarga rápida, com pico de
potência $P_0=RI_0^{2}$ muito elevado.
\end{itemize}
\textbf{Exemplo:} com $R=\SI{10}{\ohm}$, $P_0=\SI{250}{\watt}$ por
$\SI{20}{\milli\second}$; com $R=\SI{1000}{\ohm}$, $P_0=\SI{25}{\kilo\watt}$
por $\SI{0.2}{\milli\second}$. Mesma energia, potências que diferem em cem
vezes.}
\stepcounter{questao}
\resposta{(a) $\SI{200}{\kilo\volt}$ (estimativa idealizada);
(b) o ar rompe muito antes; (c) formação de arco elétrico nos contatos}
\resolucao{(a)
\[
|v|=L\frac{\Delta i}{\Delta t}=0{,}1\cdot\frac{2}{10^{-6}}
=\pot{2}{5}\,\si{\volt}=\SI{200}{\kilo\volt}.
\]
(b) A rigidez dielétrica do ar é da ordem de
$\SI{3}{\kilo\volt\per\milli\meter}$. Com contatos separados por décimos de
milímetro, o ar ioniza a alguns \emph{quilovolts} --- muito antes dos
$\SI{200}{\kilo\volt}$. O cálculo é uma \textbf{estimativa por absurdo}: ele
mostra que a hipótese "a corrente cessa em $\SI{1}{\micro\second}$" é
insustentável.\\
(c) Forma-se um \textbf{arco elétrico}, que mantém a corrente circulando
enquanto o campo se esvazia. O arco limita a tensão a algumas centenas de
volts, mas corrói os contatos, gera interferência eletromagnética e reduz
drasticamente a vida da chave.\\
\textbf{Leitura correta:} não é o indutor que "gera" $\SI{200}{\kilo\volt}$ ---
é a impossibilidade de interromper a corrente que \emph{força} o circuito a
encontrar um caminho, e ele o encontra no ar.}
\stepcounter{questao}
\resposta{(b) $\approx\SI{0.7}{\volt}$ acima da alimentação;
(c) $\tau=\SI{20}{\milli\second}$}
\resolucao{(a) Com a chave aberta, a corrente da bobina precisa de um caminho.
O diodo, que estava reversamente polarizado, entra em condução e fecha uma
malha bobina--diodo. A corrente continua circulando e decai exponencialmente,
dissipando-se em $r$ e no diodo --- sem arco.\\
(b) A tensão sobre a bobina fica grampeada em $-V_D\approx\SI{-0.7}{\volt}$, de
modo que a chave suporta apenas a tensão de alimentação acrescida de
$\SI{0.7}{\volt}$. Compare com os quilovolts do caso anterior.\\
(c) $\tau=L/r=0{,}1/5=\SI{20}{\milli\second}$.\\
\textbf{Custo do método:} a extinção fica \emph{lenta}. Em um relé, isso
prolonga o tempo de desarme; em acionamentos rápidos, prefere-se um diodo em
série com um resistor (ou um Zener), que eleva a tensão de grampeamento e
acelera a queda --- trocando proteção máxima por velocidade.}
\stepcounter{questao}
\resposta{(a) $\SI{19}{\milli\henry}$ e $\SI{7}{\milli\henry}$; (b) $k=0{,}5$}
\resolucao{(a) Em série,
\[
L_{eq}=L_1+L_2\pm2M
\;\Longrightarrow\;
4+9+6=\SI{19}{\milli\henry}
\ \text{ou}\ 4+9-6=\SI{7}{\milli\henry}.
\]
O sinal $+$ vale para \textbf{acoplamento aditivo} (os fluxos se reforçam; a
corrente entra pelos dois pontos marcados); o sinal $-$, para
\textbf{subtrativo}.\\
(b)
\[
k=\frac{M}{\sqrt{L_1L_2}}=\frac{3}{\sqrt{36}}=0{,}5.
\]
Como $0\le k\le1$, este é um acoplamento médio --- típico de bobinas
próximas com núcleo de ar. Transformadores de potência atingem $k>0{,}99$.\\
(c) \textbf{Método experimental:} meça $L_{eq}$ nas duas ligações possíveis. A
\emph{maior} corresponde ao acoplamento aditivo e identifica os pontos.
Além disso, $M$ sai de graça:
\[
M=\frac{L_{aditivo}-L_{subtrativo}}{4}=\frac{19-7}{4}=\SI{3}{\milli\henry}.
\]
É assim que se determina $M$ sem acesso à geometria interna das bobinas.}
\stepcounter{questao}
\resposta{(a) $R_{Th}=\SI{45}{\ohm}$; (b) $\tau=\SI{20}{\milli\second}$ e
$I_f=\SI{0.2}{\ampere}$}
\resolucao{(a) Zerando a fonte (curto), o indutor vê $R_2$ em série com
$R_1\parallel R_3$:
\[
R_{Th}=30+\frac{20\cdot60}{80}=30+15=\SI{45}{\ohm}.
\]
(b) $\tau=L/R_{Th}=0{,}9/45=\SI{20}{\milli\second}$.\\
A tensão de Thévenin sai com o indutor removido (sem corrente em $R_2$, logo
sem queda nele):
\[
V_{Th}=12\cdot\frac{60}{20+60}=\SI{9}{\volt}
\;\Longrightarrow\;
I_f=\frac{9}{45}=\SI{0.2}{\ampere}.
\]
\textbf{Método geral:} em qualquer circuito com \emph{um} elemento armazenador,
substitua todo o restante pelo equivalente de Thévenin visto por ele. O
problema se reduz a um RL série, com $\tau=L/R_{Th}$ e
$I_f=V_{Th}/R_{Th}$ --- e não é preciso resolver a rede inteira em regime
transitório.}
\stepcounter{questao}
\resposta{(a) $i(t)=3-2e^{-t/\tau}$ com $\tau=\SI{20}{\milli\second}$;
(b) $\SI{2.26}{\ampere}$; (c) $\SI{50}{\volt}$}
\resolucao{(a) A forma geral do transitório de primeira ordem é
\[
i(t)=I_f+\left(I_0-I_f\right)e^{-t/\tau},
\]
com $I_0=\SI{1}{\ampere}$, $I_f=\SI{3}{\ampere}$ e
$\tau=0{,}5/25=\SI{20}{\milli\second}$:
\[
i(t)=3-2e^{-t/0{,}020}.
\]
(b) $i(0{,}020)=3-2e^{-1}=3-0{,}736=\SI{2.26}{\ampere}$.\\
(c) Em $t=0^{+}$, $i=\SI{1}{\ampere}$ e o resistor cai
$25(1)=\SI{25}{\volt}$; a fonte vale $RI_f=25(3)=\SI{75}{\volt}$. Logo
\[
v_L(0^{+})=75-25=\SI{50}{\volt}.
\]
\textbf{Generalidade:} a fórmula de (a) cobre carga, descarga e condições
iniciais quaisquer --- basta identificar $I_0$, $I_f$ e $\tau$. Não é preciso
memorizar três expressões diferentes.}
\stepcounter{questao}
\resposta{$R=\SI{25}{\ohm}$; $L=\SI{0.2}{\henry}$}
\resolucao{\textbf{Resistência:} do regime permanente,
$R=V/I_f=50/2=\SI{25}{\ohm}$.\\
\textbf{Constante de tempo:} note que
$1{,}264/2=0{,}632$ --- exatamente a fração característica de $1\tau$. Logo
$\tau=\SI{8}{\milli\second}$, e
\[
L=R\tau=25(0{,}008)=\SI{0.2}{\henry}.
\]
\textbf{Se a fração não fosse reconhecível}, o caminho geral seria
$\tau=-t/\ln(1-i/I_f)$.\\
\textbf{Método:} o regime permanente entrega $R$; o transitório entrega $\tau$
e, por consequência, $L$. Duas medidas de naturezas diferentes são necessárias,
porque nenhuma delas isola $L$ sozinha.}
\stepcounter{questao}
\resposta{$p>0$ na carga (absorve), $p<0$ na descarga (devolve); a integral em
um ciclo completo é nula}
\resolucao{\textbf{Carga:} $i>0$ e crescente $\Rightarrow$
$v=L\,\mathrm{d}i/\mathrm{d}t>0$ $\Rightarrow$ $p>0$: o indutor \emph{absorve}
energia da fonte e a armazena no campo.\\
\textbf{Descarga:} $i>0$ mas decrescente $\Rightarrow$ $v<0$ $\Rightarrow$
$p<0$: o indutor \emph{fornece} energia ao circuito.\\
\textbf{Integral:} como
\[
\int p\,\mathrm{d}t=\int Li\,\frac{\mathrm{d}i}{\mathrm{d}t}\,\mathrm{d}t
=\int Li\,\mathrm{d}i=\frac12Li^{2},
\]
a energia depende \emph{apenas} da corrente instantânea, não do caminho
percorrido. Se o ciclo retorna à mesma corrente, a integral é nula: o indutor
ideal não consome energia líquida.\\
\textbf{Consequência:} indutor e capacitor são elementos \textbf{reativos} ---
trocam energia com o circuito sem consumi-la. Só a resistência dissipa. Essa
distinção é a base de todo o tratamento de potência ativa e reativa em corrente
alternada.}
\stepcounter{questao}
\resposta{(a) $u\approx\SI{573}{\kilo\joule\per\meter\cubed}$;
(b) $\approx\SI{40}{\joule\per\meter\cubed}$}
\resolucao{(a)
\[
u_B=\frac{B^{2}}{2\mu_0}=\frac{(1{,}2)^{2}}{2(4\pi\times10^{-7})}
=\frac{1{,}44}{\pot{2.513}{-6}}
\approx\pot{5.73}{5}\,\si{\joule\per\meter\cubed}.
\]
(b)
\[
u_E=\frac12\varepsilon_0E^{2}
=\frac12(\pot{8.85}{-12})(3\times10^{6})^{2}
\approx\SI{40}{\joule\per\meter\cubed}.
\]
\textbf{Razão: cerca de $14\,000$ vezes.} O campo elétrico usado é o limite de
ruptura do ar, e o magnético é típico de aço-silício --- ou seja, comparam-se
dois valores \emph{praticáveis}.\\
\textbf{Consequência tecnológica:} armazenar energia magneticamente é muito
mais denso do que eletrostaticamente. É por isso que transformadores e motores
operam com campos magnéticos, e por que capacitores só competem em energia
quando se usa dielétrico de altíssima constante e espessura nanométrica ---
como nos supercapacitores.}
\stepcounter{questao}
\resposta{(a) $M_{\max}=\SI{100}{\milli\henry}$; (b) $k=0{,}8$}
\resolucao{(a) Como $k\le1$,
\[
M_{\max}=\sqrt{L_1L_2}=\sqrt{50\cdot200}=\sqrt{10\,000}
=\SI{100}{\milli\henry}.
\]
(b) $k=80/100=0{,}8$ --- acoplamento forte, típico de bobinas sobre um mesmo
núcleo ferromagnético com pequeno entreferro.\\
(c) $k=1$ significa \textbf{acoplamento perfeito}: \emph{todo} o fluxo produzido
por uma bobina concatena a outra, sem dispersão. É o transformador ideal ---
inatingível, mas aproximado em núcleos toroidais fechados de alta
permeabilidade.\\
$k=0$ significa \textbf{desacoplamento total}: nenhuma influência mútua. Obtém-se
afastando as bobinas ou dispondo seus eixos perpendicularmente, de modo que o
fluxo de uma atravesse a outra tangencialmente. É a razão de indutores vizinhos
em placas serem montados em orientações ortogonais.}
\stepcounter{questao}
\resposta{(a) $\SI{3.2}{\ampere}$; (b) $\SI{5}{\ampere}$; (c) risco de
destruição do interruptor}
\resolucao{(a) Com $L$ constante,
\[
\Delta i=\frac{v\,\Delta t}{L}=\frac{12\cdot10^{-5}}{10^{-4}}
=\SI{1.2}{\ampere}
\;\Longrightarrow\;
i_{final}=2+1{,}2=\SI{3.2}{\ampere}.
\]
(b) A corrente sobe de $2$ a $\SI{3}{\ampere}$ com
$L=\SI{100}{\micro\henry}$, o que consome
\[
\Delta t_1=\frac{L\,\Delta i}{v}=\frac{10^{-4}\cdot1}{12}
=\SI{8.33}{\micro\second}.
\]
Restam $\SI{1.67}{\micro\second}$, agora com $L=\SI{20}{\micro\henry}$:
\[
\Delta i_2=\frac{12\cdot\pot{1.67}{-6}}{\pot{2}{-5}}=\SI{1}{\ampere}
\;\Longrightarrow\;
i_{final}=\SI{5}{\ampere}.
\]
(c) A corrente real é $56\%$ maior que a prevista pelo modelo linear, e a
subida final é \emph{cinco vezes} mais íngreme. Em um conversor chaveado, isso
significa que o transistor pode ver o dobro da corrente esperada em uma fração
do tempo --- destruição por sobrecorrente antes que qualquer proteção lenta
atue.\\
\textbf{Regra de projeto:} dimensione o indutor para que a corrente de pico
fique confortavelmente abaixo do joelho de saturação, e verifique
$L$ \emph{na corrente de operação}, nunca no valor de catálogo medido a
corrente nula.}
\stepcounter{questao}
\resposta{(a) $\SI{2}{\ampere}$ no indutor, $\SI{0}{\ampere}$ no resistor;
(b) $\SI{-20}{\volt}$; (c) $\tau=\SI{50}{\milli\second}$}
\resolucao{(a) Em regime CC o indutor ideal é um \textbf{curto}: ele desvia
toda a corrente, e o resistor em paralelo fica com tensão nula --- logo
corrente nula. Todos os $\SI{2}{\ampere}$ passam por $L$.\\
(b) Ao abrir a fonte, a corrente do indutor \emph{não pode saltar}: continua
$\SI{2}{\ampere}$, e o único caminho disponível é o resistor. A corrente agora
o percorre em sentido oposto ao anterior:
\[
v_R=-RI_0=-10(2)=\SI{-20}{\volt}.
\]
(c) $\tau=L/R=0{,}5/10=\SI{50}{\milli\second}$; a extinção é praticamente
completa em $5\tau=\SI{250}{\milli\second}$.\\
\textbf{Observação:} o resistor, que parecia inútil em regime (corrente nula),
é justamente o que evita a sobretensão destrutiva na abertura. Ele é um
\emph{caminho de escape} projetado --- exatamente a função do diodo de roda
livre, aqui desempenhada por um resistor.}
\stepcounter{questao}
\resposta{$i(t)=\dfrac{V}{R}\left(1-e^{-Rt/L}\right)$;
$\tau=L/R$; $v_L=Ve^{-Rt/L}$}
\resolucao{(a) Pela LKT,
\[
V=Ri+L\frac{\mathrm{d}i}{\mathrm{d}t}
\;\Longrightarrow\;
\frac{\mathrm{d}i}{V/R-i}=\frac{R}{L}\,\mathrm{d}t.
\]
Integrando de $0$ a $t$ (com $i(0)=0$):
\[
-\ln\!\left(\frac{V/R-i}{V/R}\right)=\frac{R}{L}t
\;\Longrightarrow\;
i(t)=\frac{V}{R}\left(1-e^{-Rt/L}\right).
\]
(b) O expoente deve ser adimensional, o que identifica
$\tau=L/R$; e $i(\infty)=V/R$, coerente com o indutor tornando-se um curto.\\
(c)
\[
v_L=L\frac{\mathrm{d}i}{\mathrm{d}t}
=L\cdot\frac{V}{R}\cdot\frac{R}{L}e^{-Rt/L}=V\,e^{-Rt/L}.
\]
\emph{Em $t=0^{+}$:} $v_L=V$ e $Ri=0$; soma $=V$ \checkmark.\\
\emph{Em $t\to\infty$:} $v_L=0$ e $Ri=V$; soma $=V$ \checkmark.\\
\textbf{Estrutura geral:} a solução é (regime permanente) $+$ (transitório
exponencial). Essa decomposição vale para \emph{qualquer} circuito de primeira
ordem e dispensa refazer a integração a cada problema.}
\stepcounter{questao}
\resposta{$W=\frac12LI_0^{2}$, independente de $R$}
\resolucao{(a) Sem fonte, $Ri+L\,\mathrm{d}i/\mathrm{d}t=0$, cuja solução é
$i(t)=I_0e^{-t/\tau}$ com $\tau=L/R$. Então
\[
p_R(t)=Ri^{2}=RI_0^{2}e^{-2t/\tau}.
\]
(b)
\[
W=\int_0^{\infty}RI_0^{2}e^{-2t/\tau}\,\mathrm{d}t
=RI_0^{2}\left[-\frac{\tau}{2}e^{-2t/\tau}\right]_0^{\infty}
=\frac{RI_0^{2}\tau}{2}.
\]
Substituindo $\tau=L/R$:
\[
W=\frac{RI_0^{2}}{2}\cdot\frac{L}{R}=\frac12LI_0^{2}.
\]
(c) O $R$ \textbf{cancela} exatamente: ele aparece na potência
($\propto R$) e desaparece na duração ($\tau\propto1/R$). Aumentar $R$ dez
vezes decuplica a potência inicial e divide o tempo por dez --- a área sob a
curva não muda.\\
\textbf{Significado:} a energia dissipada é determinada pelo \emph{estado
inicial} do indutor, não pelo caminho de descarga. É a mesma lógica que faz um
capacitor descarregado entregar sempre $\frac12CV_0^{2}$, e confirma por
integração o argumento de conservação usado no bloco anterior.}
\stepcounter{questao}
\resposta{(a) $\SI{0.24}{\ampere}$, $\SI{5.76}{\milli\joule}$;
(b) $R_g\le\SI{150}{\ohm}$; (c) $\SI{0.8}{\milli\second}$ contra
$\SI{2}{\milli\second}$}
\resolucao{(a) $I=24/100=\SI{0.24}{\ampere}$;
$W=\frac12(0{,}2)(0{,}24)^{2}=\SI{5.76}{\milli\joule}$.\\
(b) Na abertura, a corrente $\SI{0.24}{\ampere}$ é forçada pelo diodo e por
$R_g$. A tensão no coletor sobe para $V+R_gI$ (desprezando o diodo). Exigindo
$\le\SI{60}{\volt}$:
\[
24+R_g(0{,}24)\le60
\;\Longrightarrow\;
R_g\le\frac{36}{0{,}24}=\SI{150}{\ohm}.
\]
Adote $R_g=\SI{150}{\ohm}$ e verifique a dissipação: o pulso entrega
$\SI{5.76}{\milli\joule}$ e, a $\SI{10}{\hertz}$ de comutação, isso são
$\SI{57.6}{\milli\watt}$ médios --- um resistor de
$\frac14\,\si{\watt}$ basta, desde que suporte o pico de
$0{,}24^{2}(150)=\SI{8.6}{\watt}$ instantâneos.\\
(c) \emph{Com grampeador:} $\tau=L/(r+R_g)=0{,}2/250=\SI{0.8}{\milli\second}$.\\
\emph{Com diodo simples:} $\tau=L/r=0{,}2/100=\SI{2}{\milli\second}$.\\
\textbf{Compromisso quantificado:} o grampeador reduz o tempo de desarme em
$60\%$, ao custo de submeter o transistor a $\SI{60}{\volt}$ em vez de
$\SI{24.7}{\volt}$. Quanto maior $R_g$, mais rápido o relé desarma e maior a
tensão no interruptor --- o projeto consiste em usar toda a margem do
transistor, e nada além dela.}
\stepcounter{questao}
\resposta{(a) $L=\SI{0.1}{\henry}$; (b) $N\approx35$ espiras;
(c) $B\approx\SI{3.5}{\tesla}$ --- \textbf{inviável}, o núcleo satura}
\resolucao{(a) $W=\frac12LI^{2}\Rightarrow
L=\dfrac{2W}{I^{2}}=\dfrac{10}{100}=\SI{0.1}{\henry}$.\\
(b) De $L=\dfrac{\mu_r\mu_0N^{2}A}{\ell}$:
\[
N=\sqrt{\frac{L\ell}{\mu_r\mu_0A}}
=\sqrt{\frac{0{,}1\cdot0{,}25}{2000(4\pi\times10^{-7})(10^{-3})}}
=\sqrt{9{,}95}\approx35\ \text{espiras}.
\]
(c) O fluxo é
\[
B=\frac{\mu_r\mu_0NI}{\ell}
=\frac{2000(4\pi\times10^{-7})(35)(10)}{0{,}25}\approx\SI{3.5}{\tesla}.
\]
Materiais ferromagnéticos comuns saturam entre $1{,}2$ e
$\SI{1.8}{\tesla}$. O projeto é \textbf{inviável}: muito antes de atingir
$\SI{10}{\ampere}$, o núcleo satura, $\mu_r$ despenca e $L$ some.\\
\textbf{Diagnóstico e correção.} A energia armazenada em material
ferromagnético é limitada justamente porque $\mu_r$ alto significa $B$ alto para
pouca corrente. A solução clássica é introduzir um \textbf{entreferro}: ele
reduz a permeabilidade efetiva (exigindo mais espiras), mas permite corrente
muito maior antes da saturação --- e a maior parte da energia passa a residir
no entreferro, onde $u=B^{2}/2\mu_0$ é enorme.\\
\textbf{Lição:} calcular $L$ e $N$ não basta. Sem verificar $B$, obtém-se um
projeto aritmeticamente correto e fisicamente impossível.}
\stepcounter{questao}
\resposta{(a) $L=\SI{1}{\henry}$; (c) aumentar $R$, ou usar RC equivalente}
\resolucao{(a) $L=\tau R=0{,}020(50)=\SI{1}{\henry}$.\\
(b) \textbf{Um henry é um componente grande.} Exige núcleo ferromagnético e
centenas de espiras; a resistência do próprio enrolamento facilmente atinge
dezenas de ohms --- comparável aos $\SI{50}{\ohm}$ do projeto, o que
\emph{altera} o $\tau$ pretendido. Some-se a isso massa, custo, sensibilidade a
campos externos e saturação.\\
(c) \textbf{Alternativa 1 --- elevar $R$.} Com $R=\SI{500}{\ohm}$,
$L=\SI{10}{\henry}$... pior. O caminho correto é o inverso: fixar um $L$
razoável e ajustar $R$. Com $L=\SI{10}{\milli\henry}$, obter
$\tau=\SI{20}{\milli\second}$ exigiria $R=\SI{0.5}{\ohm}$ --- baixo demais para
ser prático em circuitos de sinal.\\
\textbf{Alternativa 2 --- usar RC.} Um capacitor de
$\SI{20}{\micro\farad}$ com $\SI{1}{\kilo\ohm}$ dá $\tau=\SI{20}{\milli\second}$
em um componente barato, pequeno e estável.\\
\textbf{Conclusão de engenharia:} para constantes de tempo da ordem de
milissegundos ou mais, \textbf{RC vence RL} quase sempre. Indutores são
preferíveis quando se precisa de armazenamento de energia com corrente elevada
ou de comportamento em frequência --- não para gerar atrasos.}
\stepcounter{questao}
\resposta{(b) da condição $W\ge0$ para toda corrente segue
$M^{2}\le L_1L_2$}
\resolucao{(a) Cada bobina sofre autoindução \emph{e} indução mútua:
\[
v_1=L_1\frac{\mathrm{d}i}{\mathrm{d}t}\pm M\frac{\mathrm{d}i}{\mathrm{d}t},
\qquad
v_2=L_2\frac{\mathrm{d}i}{\mathrm{d}t}\pm M\frac{\mathrm{d}i}{\mathrm{d}t}.
\]
Somando e comparando com
$v=L_{eq}\,\mathrm{d}i/\mathrm{d}t$:
\[
L_{eq}=L_1+L_2\pm2M.
\]
(b) Com correntes independentes $i_1$ e $i_2$, a energia do par é
\[
W=\frac12L_1i_1^{2}+\frac12L_2i_2^{2}+Mi_1i_2.
\]
Como um sistema passivo não pode armazenar energia negativa, $W\ge0$ para
\emph{quaisquer} $i_1,i_2$. Tomando $i_1=x$ e $i_2=-1$:
\[
\frac12L_1x^{2}-Mx+\frac12L_2\ge0\quad\forall x,
\]
o que exige discriminante não positivo:
\[
M^{2}-L_1L_2\le0
\;\Longrightarrow\;
M\le\sqrt{L_1L_2}
\;\Longrightarrow\;
k=\frac{M}{\sqrt{L_1L_2}}\le1.
\]
(c) Os pontos marcam terminais \textbf{magneticamente correspondentes}:
correntes que entram pelos terminais pontuados produzem fluxos que se
\emph{reforçam}. Se ambas as correntes entram (ou ambas saem) pelos pontos, o
termo mútuo é positivo; caso contrário, negativo.\\
\textbf{Por que a convenção é necessária:} o esquema elétrico não informa o
sentido de enrolamento das bobinas, que é uma propriedade \emph{física} da
construção. Sem os pontos, $L_{eq}$ seria ambíguo entre
$\SI{19}{\milli\henry}$ e $\SI{7}{\milli\henry}$ --- e, em um transformador, a
polaridade da saída ficaria indeterminada.}
\stepcounter{questao}
\resposta{(a) $\SI{5}{\ampere}$, $\SI{6.25}{\joule}$, $\SI{50}{\watt}$;
(b) $\approx\SI{1.15}{\second}$}
\resolucao{(a) $I=10/2=\SI{5}{\ampere}$;
$W=\frac12(0{,}5)(25)=\SI{6.25}{\joule}$;
$P=rI^{2}=2(25)=\SI{50}{\watt}$.\\
(b) $\tau=L/r=0{,}5/2=\SI{0.25}{\second}$, e $99\%$ correspondem a
$t=\tau\ln100=0{,}25(4{,}605)=\SI{1.15}{\second}$.\\
(c) \textbf{O contraste é gritante.} Carregar o campo custou
$\SI{6.25}{\joule}$; mantê-lo custa $\SI{50}{\watt}$ \emph{continuamente}.
Em apenas $\SI{0.125}{\second}$ de operação, a dissipação já iguala toda a
energia armazenada; em um minuto, dissipam-se $\SI{3000}{\joule}$ --- $480$
vezes o conteúdo do campo.\\
\textbf{Figura de mérito:} a razão $W/P=L/(2r)=\tau/2$ tem dimensão de tempo e
mede a "qualidade" do indutor em CC. Aqui vale $\SI{0.125}{\second}$.\\
\textbf{Consequência:} indutores são péssimos \emph{reservatórios} de energia
para armazenamento prolongado --- ao contrário de baterias ou capacitores, que
não exigem corrente para manter o estado. Eles são excelentes para
\emph{transferir} energia em ciclos curtos, que é exatamente o que fazem em
conversores chaveados, onde o ciclo dura microssegundos.}
\stepcounter{questao}
\resposta{(a) $\Delta i=\SI{1.2}{\ampere}$; (b) $I_{pico}=\SI{3.6}{\ampere}$;
(c) satura --- elevar $L$ ou a frequência}
\resolucao{(a) Durante o tempo de condução
$t_{on}=D/f=0{,}5/\pot{5}{4}=\SI{10}{\micro\second}$:
\[
\Delta i=\frac{V\,t_{on}}{L}=\frac{12(10^{-5})}{10^{-4}}=\SI{1.2}{\ampere}.
\]
(b) A ondulação se distribui simetricamente em torno da média:
\[
I_{pico}=3+\frac{1{,}2}{2}=\SI{3.6}{\ampere}\ >\ \SI{3.5}{\ampere}.
\]
\textbf{O indutor satura} --- e note que a corrente \emph{média}
($\SI{3}{\ampere}$) estava confortavelmente abaixo do limite. Dimensionar pela
média é o erro clássico.\\
(c) Como $\Delta i=\dfrac{VD}{fL}$, há dois caminhos:\\
\textbf{(i) Aumentar $L$.} Para $I_{pico}\le3{,}5$ é preciso
$\Delta i\le\SI{1}{\ampere}$, ou seja
\[
L\ge\frac{12(10^{-5})}{1}=\SI{120}{\micro\henry}.
\]
Custo: indutor maior e mais caro.\\
\textbf{(ii) Aumentar a frequência.} Mantendo
$L=\SI{100}{\micro\henry}$, exige-se
\[
f\ge\frac{12(0{,}5)}{10^{-4}(1)}=\SI{60}{\kilo\hertz}.
\]
Custo: mais perdas de comutação no transistor e maior interferência
eletromagnética.\\
\textbf{Compromisso:} $\Delta i\propto1/(fL)$ --- o produto $fL$ é o que
importa. Elevar a frequência permite indutores menores, e é exatamente essa a
tendência que tornou as fontes chaveadas modernas tão compactas.}
\stepcounter{questao}
\resposta{Um salto exigiria tensão e potência infinitas; a energia é finita,
mas a taxa de transferência não pode ser}
\resolucao{(a) Um salto implica
\[
v=L\lim_{\Delta t\to0}\frac{I_2-I_1}{\Delta t}\to\infty.
\]
Nenhuma fonte real fornece tensão infinita.\\
(b) A potência $p=vi$ também diverge. A \emph{energia} necessária é finita
---
\[
\Delta W=\frac12L\left(I_2^{2}-I_1^{2}\right),
\]
--- mas ela teria de ser transferida em tempo nulo, exigindo potência
infinita. \textbf{É a taxa, não a quantidade, que torna o salto impossível.}\\
(c) \textbf{Problema dual:} dois capacitores com tensões diferentes ligados em
paralelo exigiriam salto de \emph{tensão}, o que demandaria corrente infinita
($i=C\,\mathrm{d}v/\mathrm{d}t$). Na prática, a resistência inevitável dos
condutores limita a corrente e dissipa a diferença de energia --- e demonstra-se
que a perda independe do valor dessa resistência, resultado análogo ao da
descarga do indutor.\\
\textbf{Regra prática:} nunca ligue em série indutores com correntes distintas,
nem em paralelo capacitores com tensões distintas. Em ambos os casos, o modelo
ideal falha e o circuito real encontra uma saída --- arco elétrico no primeiro
caso, pico de corrente no segundo.}
\stepcounter{questao}
\resposta{(a) $\SI{20}{\milli\henry}$ e $\SI{12}{\milli\henry}$;
(b) $-\SI{10}{\milli\henry\per\milli\meter}$ e
$-\SI{0.4}{\milli\henry\per\milli\meter}$}
\resolucao{(a)
\[
L(1)=10(1+1)=\SI{20}{\milli\henry};
\qquad
L(5)=10(1+0{,}2)=\SI{12}{\milli\henry}.
\]
(b)
\[
\frac{\mathrm{d}L}{\mathrm{d}x}=-\frac{L_0a}{x^{2}}.
\]
Em $x=\SI{1}{\milli\meter}$:
$-10(1)/1=-\SI{10}{\milli\henry\per\milli\meter}$.\\
Em $x=\SI{5}{\milli\meter}$:
$-10(1)/25=-\SI{0.4}{\milli\henry\per\milli\meter}$.\\
A sensibilidade cai com $1/x^{2}$: é \textbf{25 vezes menor} a
$\SI{5}{\milli\meter}$.\\
(c) \textbf{Três consequências.} \emph{(i)} O sensor é fortemente
\textbf{não linear} --- a leitura bruta de $L$ não é proporcional a $x$, e a
conversão exige linearização ou tabela. \emph{(ii)} A \textbf{faixa útil} é
curta: além de poucos milímetros, a variação se confunde com a deriva térmica
da própria bobina. \emph{(iii)} A \textbf{resolução} depende da posição: perto
do alvo, décimos de micrômetro são detectáveis; longe, nem décimos de
milímetro.\\
\textbf{Solução usual:} operar sempre na região próxima e usar o sensor como
detector de \emph{limiar} (presença/ausência), em vez de medidor de distância
--- que é exatamente como os sensores indutivos industriais são
especificados.}

\gabarito[10.2 Indutores]

\section{Circuitos de Primeira e Segunda Ordem}
\stepcounter{questao}
\resposta{$v_C(0^+)=\SI{7}{\volt}$}
\resolucao{A tensão em um capacitor ideal não pode variar instantaneamente, pois
$i_C=C\,\mathrm{d}v_C/\mathrm{d}t$ exigiria corrente impulsiva infinita. Logo,
$v_C(0^+)=v_C(0^-)=\SI{7}{\volt}$.}
\stepcounter{questao}
\resposta{$i_L(0^+)=\SI{350}{\milli\ampere}$}
\resolucao{A corrente de um indutor ideal é contínua. Como
$v_L=L\,\mathrm{d}i_L/\mathrm{d}t$, uma mudança instantânea de corrente exigiria
tensão impulsiva infinita. Portanto $i_L(0^+)=i_L(0^-)$.}
\stepcounter{questao}
\resposta{$\tau=\SI{100}{\milli\second}$}
\resolucao{$\tau=RC=(\pot{2}{3})(\pot{50}{-6})=\SI{0.1}{\second}
=\SI{100}{\milli\second}$.}
\stepcounter{questao}
\resposta{$\tau=\SI{4}{\milli\second}$}
\resolucao{$\tau=L/R_{\mathrm{eq}}=0{,}08/20=\SI{0.004}{\second}
=\SI{4}{\milli\second}$.}
\stepcounter{questao}
\resposta{(c)}
\resolucao{Para uma carga partindo de zero,
$v_C(\tau)=V_f(1-e^{-1})\approx0{,}632V_f$. Portanto cerca de
\SI{63.2}{\percent} da variação total já ocorreu.}
\stepcounter{questao}
\resposta{$e^{-3}\times100\%\approx\SI{4.98}{\percent}$}
\resolucao{$e^{-3}\approx0{,}0498$. Assim resta cerca de \SI{4.98}{\percent} do
valor inicial.}
\stepcounter{questao}
\resposta{(c)}
\resolucao{Para calcular $R_{\mathrm{eq}}$ com fontes independentes suprimidas,
uma fonte ideal de tensão é substituída por curto-circuito e uma fonte ideal de
corrente por circuito aberto. Fontes dependentes, quando existirem, permanecem
ativas.}
\stepcounter{questao}
\resposta{Segunda ordem}
\resolucao{Cada variável independente de armazenamento de energia pode introduzir
um estado. Um capacitor e um indutor independentes produzem duas variáveis de
estado, caracterizando um circuito de segunda ordem.}
\stepcounter{questao}
\resposta{$\omega_0\approx\SI{316.2}{\radian\per\second}$}
\resolucao{$LC=0{,}2\cdot\pot{50}{-6}=\pot{1}{-5}$ e
$\omega_0=1/\sqrt{LC}=1/\sqrt{\pot{1}{-5}}\approx\SI{316.2}{\radian\per\second}$.}
\stepcounter{questao}
\resposta{Subamortecido}
\resolucao{$\alpha=R/(2L)=40/0{,}2=\SI{200}{\per\second}$ e
$\omega_0=1/\sqrt{0{,}1\cdot\pot{100}{-6}}\approx\SI{316.2}{\radian\per\second}$.
Como $\alpha<\omega_0$, a resposta é subamortecida.}
\stepcounter{questao}
\resposta{$R_{\mathrm{crit}}\approx\SI{31.6}{\ohm}$}
\resolucao{$R_{\mathrm{crit}}=2\sqrt{L/C}=2\sqrt{0{,}05/(\pot{200}{-6})}
=2\sqrt{250}\approx\SI{31.6}{\ohm}$.}

\gabarito[Nível 1]
\stepcounter{questao}
\resposta{$v_C(\SI{0.2}{\second})\approx\SI{2.81}{\volt}$}
\resolucao{$\tau=RC=\SI{0.1}{\second}$. Logo
\[
v_C(t)=2+(8-2)e^{-t/0{,}1}.
\]
Em $t=\SI{0.2}{\second}$, $v_C=2+6e^{-2}\approx\SI{2.81}{\volt}$.}
\stepcounter{questao}
\resposta{$i_L\approx\SI{3.59}{\ampere}$}
\resolucao{$\tau=L/R=0{,}6/30=\SI{20}{\milli\second}$. Assim
$i_L(t)=4+(1-4)e^{-t/\tau}$. Para $t=2\tau$,
$i_L=4-3e^{-2}\approx\SI{3.59}{\ampere}$.}
\stepcounter{questao}
\resposta{$t\approx\SI{230}{\milli\second}$}
\resolucao{$\tau=\SI{0.1}{\second}$. De
$0{,}9=1-e^{-t/\tau}$ resulta $e^{-t/\tau}=0{,}1$ e
$t=\tau\ln 10\approx\SI{0.230}{\second}$.}
\stepcounter{questao}
\resposta{$t\approx\SI{69.3}{\milli\second}$}
\resolucao{$3=12e^{-t/\tau}$, então $e^{-t/\tau}=1/4$ e
$t=\tau\ln4=\SI{50}{\milli\second}\cdot1{,}386\approx\SI{69.3}{\milli\second}$.}
\stepcounter{questao}
\resposta{$t\approx\SI{57.6}{\milli\second}$}
\resolucao{$0{,}2=2e^{-t/\tau}$ implica $e^{-t/\tau}=0{,}1$. Logo
$t=\tau\ln10\approx25\cdot2{,}303=\SI{57.6}{\milli\second}$.}
\stepcounter{questao}
\resposta{$\tau=\SI{100}{\milli\second}$; $v_C(\tau)\approx\SI{4.21}{\volt}$}
\resolucao{Em regime permanente, a tensão no capacitor corresponde à tensão do divisor:
$V_f=20\cdot3/(6+3)=\SI{6.67}{\volt}$. Suprimindo a fonte,
$R_{\mathrm{eq}}=R_1\parallel R_2=\SI{2}{\kilo\ohm}$, então
$\tau=R_{\mathrm{eq}}C=\SI{0.1}{\second}$. Como $v_C(0)=0$,
$v_C(\tau)=V_f(1-e^{-1})\approx\SI{4.21}{\volt}$.}
\stepcounter{questao}
\resposta{$v_C\approx-\SI{0.90}{\volt}$}
\resolucao{$v_C(t)=-10+[5-(-10)]e^{-t/\tau}=-10+15e^{-t/\tau}$.
Para $t/\tau=0{,}5$, obtém-se $v_C\approx-10+15e^{-0{,}5}
\approx-\SI{0.90}{\volt}$.}
\stepcounter{questao}
\resposta{$i_L\approx\SI{1.80}{\ampere}$}
\resolucao{$\tau=L/R=\SI{10}{\milli\second}$ e
$i_L(t)=2+(0{,}5-2)e^{-t/\tau}$. Em $t=2\tau$,
$i_L=2-1{,}5e^{-2}\approx\SI{1.80}{\ampere}$.}
\stepcounter{questao}
\resposta{$\alpha=\SI{100}{\per\second}$; $\omega_0\approx\SI{316.2}{\radian\per\second}$;
$\omega_d=\SI{300}{\radian\per\second}$; subamortecida}
\resolucao{$\alpha=R/(2L)=10/0{,}1=100$. Além disso,
$\omega_0=1/\sqrt{0{,}05\cdot\pot{200}{-6}}\approx316{,}2$. Como
$\alpha<\omega_0$, a resposta é subamortecida e
$\omega_d=\sqrt{316{,}2^2-100^2}=\SI{300}{\radian\per\second}$.}
\stepcounter{questao}
\resposta{Superamortecida}
\resolucao{$\alpha=40/(2\cdot0{,}05)=\SI{400}{\per\second}$, enquanto
$\omega_0\approx\SI{316.2}{\radian\per\second}$. Como $\alpha>\omega_0$, a
resposta é superamortecida.}
\stepcounter{questao}
\resposta{$R_{\mathrm{crit}}\approx\SI{63.25}{\ohm}$}
\resolucao{$R_{\mathrm{crit}}=2\sqrt{L/C}=2\sqrt{0{,}1/\pot{100}{-6}}
=2\sqrt{1000}\approx\SI{63.25}{\ohm}$.}
\stepcounter{questao}
\resposta{$B\approx\SI{3.33}{\volt}$}
\resolucao{Derivando e avaliando em zero,
$\dot v(0)=-100\cdot10+300B=0$. Portanto
$B=1000/300\approx\SI{3.33}{\volt}$.}
\stepcounter{questao}
\resposta{$\mathrm{d}v_C/\mathrm{d}t|_{0^+}=\SI{200}{\volt\per\second}$}
\resolucao{Pela convenção passiva, $i_C=C\dot v_C$. Logo
$\dot v_C=0{,}02/(\pot{100}{-6})=\SI{200}{\volt\per\second}$.}
\stepcounter{questao}
\resposta{Subamortecida}
\resolucao{No circuito RLC em paralelo, $\alpha=1/(2RC)=1/[2\cdot200\cdot\pot{100}{-6}]
=\SI{25}{\per\second}$. Já $\omega_0\approx\SI{316.2}{\radian\per\second}$.
Como $\alpha<\omega_0$, a resposta é subamortecida.}
\stepcounter{questao}
\resposta{$\approx\SI{0.674}{\percent}$}
\resolucao{O erro normalizado é $e^{-t/\tau}$. Em $5\tau$,
$e^{-5}\approx0{,}00674$, ou \SI{0.674}{\percent}.}
\stepcounter{questao}
\resposta{$E=\SI{5}{\milli\joule}$; $I_{\max}\approx\SI{0.447}{\ampere}$}
\resolucao{$E=\tfrac12CV^2=\tfrac12\pot{100}{-6}\cdot100=\SI{5}{\milli\joule}$.
Quando toda a energia está no indutor,
$\tfrac12LI_{\max}^2=E$, então
$I_{\max}=V\sqrt{C/L}=10\sqrt{\pot{100}{-6}/0{,}05}\approx\SI{0.447}{\ampere}$.}
\stepcounter{questao}
\resposta{$T_0\approx\SI{6.28}{\milli\second}$}
\resolucao{$\omega_0=1/\sqrt{LC}=1/\sqrt{\pot{1}{-6}}=\SI{1000}{\radian\per\second}$.
Logo $T_0=2\pi/\omega_0\approx\SI{6.28}{\milli\second}$.}

\gabarito[Nível 2]
\stepcounter{questao}
\resposta{$v_C\approx\SI{4.95}{\volt}$}
\resolucao{Na primeira etapa, $\tau_1=\SI{20}{\milli\second}$ e
$v_C(50\,\mathrm{ms})=12(1-e^{-2{,}5})\approx\SI{11.015}{\volt}$. Esse valor é a
condição inicial da segunda etapa. Agora $\tau_2=\SI{50}{\milli\second}$ e
$v_C=11{,}015e^{-40/50}\approx\SI{4.95}{\volt}$.}
\stepcounter{questao}
\resposta{$i_L\approx\SI{0.210}{\ampere}$}
\resolucao{Na energização, $I_f=20/10=\SI{2}{\ampere}$ e
$\tau_1=0{,}1/10=\SI{10}{\milli\second}$. Assim
$i_L(15\,\mathrm{ms})=2(1-e^{-1{,}5})\approx\SI{1.554}{\ampere}$. Na descarga,
$\tau_2=0{,}1/25=\SI{4}{\milli\second}$, portanto após $2\tau_2$:
$i_L=1{,}554e^{-2}\approx\SI{0.210}{\ampere}$.}
\stepcounter{questao}
\resposta{$R_{\mathrm{eq}}=\SI{3}{\kilo\ohm}$; $\tau=\SI{60}{\milli\second}$}
\resolucao{$R_1\parallel R_2=(6\cdot3)/(6+3)=\SI{2}{\kilo\ohm}$. Em série com
$R_3$, $R_{\mathrm{eq}}=\SI{3}{\kilo\ohm}$. Logo
$\tau=R_{\mathrm{eq}}C=\pot{3}{3}\pot{20}{-6}=\SI{60}{\milli\second}$.}
\stepcounter{questao}
\resposta{$t\approx\SI{40.5}{\milli\second}$}
\resolucao{$\tau=\SI{0.1}{\second}$ e
$v_C=8+(-4-8)e^{-t/0{,}1}=8-12e^{-t/0{,}1}$. No cruzamento,
$e^{-t/0{,}1}=2/3$. Portanto $t=0{,}1\ln(3/2)\approx\SI{40.5}{\milli\second}$.}
\stepcounter{questao}
\resposta{$t\approx\SI{7.19}{\milli\second}$}
\resolucao{$i_L=3+(-1-3)e^{-t/\tau}=3-4e^{-t/\tau}$. Para $i_L=0$,
$e^{-t/\tau}=3/4$. Logo $t=\tau\ln(4/3)\approx\SI{7.19}{\milli\second}$.}
\stepcounter{questao}
\resposta{$i(0^+)=\SI{4}{\milli\ampere}$; $i(150\,\mathrm{ms})\approx\SI{0.893}{\milli\ampere}$}
\resolucao{$\tau=RC=\SI{0.1}{\second}$. A corrente é
$i(t)=(10-2)R^{-1}e^{-t/\tau}=\SI{4}{\milli\ampere}e^{-t/\tau}$. Em
$150\,\mathrm{ms}=1{,}5\tau$, $i\approx4e^{-1{,}5}=\SI{0.893}{\milli\ampere}$.}
\stepcounter{questao}
\resposta{$E=\SI{20}{\milli\joule}$}
\resolucao{$v_C=V_0e^{-t/RC}$ e $i=(V_0/R)e^{-t/RC}$. Logo
$p_R=Ri^2=(V_0^2/R)e^{-2t/RC}$. Integrando de zero a infinito:
\[
E_R=\frac{V_0^2}{R}\int_0^\infty e^{-2t/RC}\,\mathrm{d}t
=\frac{V_0^2}{R}\frac{RC}{2}=\frac12CV_0^2.
\]
Numericamente, $E_R=\tfrac12\pot{100}{-6}\cdot20^2=\SI{20}{\milli\joule}$.}
\stepcounter{questao}
\resposta{$C\approx\SI{4.63}{\micro\farad}$; valor comercial: \SI{4.7}{\micro\farad}}
\resolucao{$3{,}3=5(1-e^{-t/RC})$, então
$C=-t/[R\ln(1-3{,}3/5)]$. Substituindo,
$C\approx\SI{4.63}{\micro\farad}$. Um valor comercial próximo é
\SI{4.7}{\micro\farad}.}
\stepcounter{questao}
\resposta{$R\approx\SI{7.49}{\ohm}$; $V\approx\SI{7.49}{\volt}$}
\resolucao{$0{,}95=1-e^{-t/\tau}$ implica
$\tau=-t/\ln0{,}05\approx\SI{6.68}{\milli\second}$. Como $\tau=L/R$,
$R=0{,}05/0{,}00668\approx\SI{7.49}{\ohm}$. Para corrente final de \SI{1}{\ampere},
$V=RI_f\approx\SI{7.49}{\volt}$.}
\stepcounter{questao}
\resposta{$s_{1,2}=-100\pm j300\ \si{\per\second}$}
\resolucao{$\alpha=R/(2L)=100$ e $\omega_0\approx316{,}2$. Assim
$\omega_d=\sqrt{316{,}2^2-100^2}=300$. Para o caso subamortecido,
$s_{1,2}=-\alpha\pm j\omega_d=-100\pm j300\ \si{\per\second}$.}
\stepcounter{questao}
\resposta{$A=\SI{10}{\volt}$; $B=\SI{1500}{\volt\per\second}$}
\resolucao{De $v(0)=A$, segue $A=10$. Derivando,
$\dot v(0)=B-200A=-500$. Logo $B=-500+2000=\SI{1500}{\volt\per\second}$.}
\stepcounter{questao}
\resposta{$s_1\approx-\SI{112.7}{\per\second}$; $s_2\approx-\SI{887.3}{\per\second}$}
\resolucao{A equação característica é
$s^2+(R/L)s+1/(LC)=s^2+1000s+100000=0$. Portanto
\[
s=\frac{-1000\pm\sqrt{1000^2-4\cdot100000}}{2}
=\frac{-1000\pm\sqrt{600000}}{2},
\]
resultando aproximadamente em $-112{,}7$ e $-887{,}3\ \si{\per\second}$.}
\stepcounter{questao}
\resposta{$R_{\mathrm{crit}}\approx\SI{7.91}{\ohm}$}
\resolucao{No circuito RLC em paralelo, a condição crítica é
$1/(2RC)=1/\sqrt{LC}$. Logo
$R_{\mathrm{crit}}=\tfrac12\sqrt{L/C}
=\tfrac12\sqrt{0{,}05/\pot{200}{-6}}\approx\SI{7.91}{\ohm}$.}
\stepcounter{questao}
\resposta{$M_p\approx\SI{16.3}{\percent}$}
\resolucao{Substituindo $\zeta=0{,}5$:
$M_p=e^{-\pi(0{,}5)/\sqrt{0{,}75}}\approx0{,}163$. Portanto o sobressinal é
aproximadamente \SI{16.3}{\percent}.}
\stepcounter{questao}
\resposta{$\zeta\approx0{,}316$}
\resolucao{$\alpha=R/(2L)=100$ e $\omega_0\approx316{,}2$. Assim
$\zeta=\alpha/\omega_0\approx100/316{,}2\approx0{,}316$.}
\stepcounter{questao}
\resposta{A subamortecido; B aproximadamente crítico; C superamortecido; D subamortecido}
\resolucao{A, B e C possuem $\omega_0\approx316{,}2$. No circuito RLC em série,
$\alpha=R/(0{,}2)$: A dá $100<316{,}2$; B dá aproximadamente $316{,}25$; C dá
$500>316{,}2$. No caso D, $\alpha=1/(2RC)=125$ e
$\omega_0=1/\sqrt{0{,}05\cdot\pot{200}{-6}}\approx316{,}2$, logo também é
subamortecido.}
\stepcounter{questao}
\resposta{$E=\SI{6}{\milli\joule}$}
\resolucao{$E_C=\tfrac12CV^2=\SI{5}{\milli\joule}$ e
$E_L=\tfrac12LI^2=\tfrac12(0{,}05)(0{,}2)^2=\SI{1}{\milli\joule}$. A soma é
\SI{6}{\milli\joule}.}
\stepcounter{questao}
\resposta{$v_C(0^+)=\SI{6}{\volt}$ e $i_L(0^+)=\SI{0.3}{\ampere}$}
\resolucao{As variáveis de estado associadas aos elementos ideais armazenadores
são contínuas: a tensão do capacitor e a corrente do indutor mantêm seus valores
no instante da comutação. A nova topologia altera suas derivadas e seus valores
finais, não os valores instantâneos em $0^+$.}
\stepcounter{questao}
\resposta{$\ddot v_C+200\dot v_C+100000v_C=1000000$}
\resolucao{Pela LKT,
$LC\ddot v_C+RC\dot v_C+v_C=V_s$. Como $LC=\pot{1}{-5}$ e
$RC=\pot{2}{-3}$, dividindo por $LC$ resulta
\[
\ddot v_C+200\dot v_C+100000v_C=10\cdot100000=1000000.
\]}
\stepcounter{questao}
\resposta{$v_C(0^+)=0$; $\dot v_C(0^+)=0$; $\ddot v_C(0^+)=\SI{1e6}{\volt\per\second\squared}$}
\resolucao{Energia inicial nula implica $v_C(0^+)=0$ e $i_L(0^+)=0$. Como
$i_C=C\dot v_C$, segue $\dot v_C(0^+)=0$. Substituindo esses valores na equação
$\ddot v_C+200\dot v_C+100000v_C=1000000$, obtém-se
$\ddot v_C(0^+)=\SI{1e6}{\volt\per\second\squared}$.}

\gabarito[Nível 3]
\stepcounter{questao}
\resposta{$R_{\mathrm{eq}}\approx\SI{8}{\kilo\ohm}$; há cerca de \SI{3}{\kilo\ohm}
adicionais na trajetória equivalente}
\resolucao{$R_{\mathrm{eq}}=\tau/C=0{,}08/(\pot{10}{-6})=\SI{8}{\kilo\ohm}$.
Como o resistor nominal é de \SI{5}{\kilo\ohm}, há uma resistência
adicional de aproximadamente \SI{3}{\kilo\ohm}. Essa diferença pode decorrer da resistência de saída da fonte, de uma
rede de polarização não considerada, de um resistor em série, de uma chave ou
sensor, ou de uma topologia
que não foi corretamente reduzida ao equivalente de Thévenin.}
\stepcounter{questao}
\resposta{$R=\SI{240}{\ohm}$; $t_{95}\approx\SI{0.338}{\second}$; $E_R\approx\SI{0.135}{\joule}$}
\resolucao{Pela corrente inicial, $R\ge24/0{,}1=\SI{240}{\ohm}$. Com esse valor,
$\tau=RC=240\cdot\pot{470}{-6}=\SI{112.8}{\milli\second}$. Para \SI{95}{\percent},
$t=-\tau\ln0{,}05\approx\SI{0.338}{\second}$, atendendo à especificação. A
energia dissipada no resistor durante a carga ideal por fonte de tensão é igual à
energia final do capacitor:
$E_R=\tfrac12CV^2\approx\SI{0.135}{\joule}$. A potência inicial é
$V^2/R=\SI{2.4}{\watt}$, importante para a escolha da tecnologia do resistor.}
\stepcounter{questao}
\resposta{$I_0=\SI{0.4}{\ampere}$; $E_L=\SI{9.6}{\milli\joule}$; $\tau_A\approx\SI{2}{\milli\second}$;
$\tau_B\approx\SI{0.5}{\milli\second}$; $v_{R,\mathrm{clamp}}(0^+)\approx\SI{72}{\volt}$}
\resolucao{Em regime, $I_0=24/60=\SI{0.4}{\ampere}$ e
$E_L=\tfrac12LI_0^2=\SI{9.6}{\milli\joule}$. Com o diodo ideal, a resistência no caminho de descarga é aproximadamente
a resistência do enrolamento da bobina: $\tau_A=L/60=\SI{2}{\milli\second}$. Com o
resistor adicional, $R_{\mathrm{tot}}=60+180=\SI{240}{\ohm}$ e
$\tau_B=L/240=\SI{0.5}{\milli\second}$. A tensão inicial no resistor externo é
$0{,}4\cdot180=\SI{72}{\volt}$. O clamp resistivo desliga o relé mais rápido,
mas produz tensão maior; o diodo reduz a sobretensão, porém prolonga a corrente.}
\stepcounter{questao}
\resposta{$\alpha\approx\SI{69.3}{\per\second}$; $\omega_d\approx\SI{1571}{\radian\per\second}$;
$\omega_0\approx\SI{1572}{\radian\per\second}$; $\zeta\approx0{,}044$;
$L\approx\SI{40.45}{\milli\henry}$; $R\approx\SI{5.61}{\ohm}$}
\resolucao{Da envoltória, $e^{-\alpha(0{,}01)}=1/2$, então
$\alpha=\ln2/0{,}01\approx\SI{69.3}{\per\second}$. Do período,
$\omega_d=2\pi/0{,}004\approx\SI{1571}{\radian\per\second}$. Logo
$\omega_0=\sqrt{\omega_d^2+\alpha^2}\approx\SI{1572}{\radian\per\second}$ e
$\zeta\approx0{,}044$. Como $\omega_0^2=1/(LC)$,
$L\approx1/(C\omega_0^2)\approx\SI{40.45}{\milli\henry}$. Por fim,
$R=2\alpha L\approx\SI{5.61}{\ohm}$.}
\stepcounter{questao}
\resposta{$R_{\mathrm{crit}}\approx\SI{89.4}{\ohm}$; \SI{45}{\ohm} é subamortecido;
\SI{180}{\ohm} é superamortecido}
\resolucao{$R_{\mathrm{crit}}=2\sqrt{L/C}=2\sqrt{0{,}02/\pot{10}{-6}}
\approx\SI{89.4}{\ohm}$. Com \SI{45}{\ohm}, $R<R_{\mathrm{crit}}$ e haverá
oscilações amortecidas (ringing) e sobressinal. Com \SI{180}{\ohm}, $R>R_{\mathrm{crit}}$ e a resposta será
monotônica, porém mais lenta. O amortecimento crítico é o limite de resposta
monotônica mais rápida para o modelo ideal.}
\stepcounter{questao}
\resposta{$\tau\approx\SI{20}{\milli\second}$; $R\approx\SI{4.26}{\kilo\ohm}$}
\resolucao{Em $t=\tau$, uma carga RC atinge \SI{63.2}{\percent}, exatamente o
primeiro ponto. Em $2\tau$, atinge $1-e^{-2}=0{,}8647$, coerente com
\SI{8.65}{\volt}. Logo $\tau\approx\SI{20}{\milli\second}$. Assim
$R=\tau/C=0{,}02/(\pot{4.7}{-6})\approx\SI{4.26}{\kilo\ohm}$.}
\stepcounter{questao}
\resposta{$v_C(t)=e^{-200t}[5\cos(400t)+5\sin(400t)]\ \si{\volt}$}
\resolucao{$\alpha=R/(2L)=200$, $\omega_0=1/\sqrt{LC}\approx447{,}2$ e
$\omega_d=\sqrt{\omega_0^2-\alpha^2}=400$. Portanto
$v_C=e^{-200t}(A\cos400t+B\sin400t)$. De $v_C(0)=5$, $A=5$. Como
$i=C\dot v_C$, $\dot v_C(0)=0{,}1/(\pot{100}{-6})=\SI{1000}{\volt\per\second}$.
Então $\dot v_C(0)=-200A+400B=1000$, resultando em $B=5$.}
\stepcounter{questao}
\resposta{$R_{\mathrm{carga}}\approx\SI{8.59}{\kilo\ohm}$; $R_{\mathrm{descarga}}\lesssim\SI{2.65}{\kilo\ohm}$;
valores práticos: \SI{8.2}{\kilo\ohm} e \SI{2.4}{\kilo\ohm}}
\resolucao{Na carga,
$2=3{,}3(1-e^{-8\mathrm{ms}/RC})$, fornecendo
$R\approx\SI{8.59}{\kilo\ohm}$. Se esse mesmo resistor fosse usado na descarga,
o tempo até \SI{0.5}{\volt} seria $RC\ln(3{,}3/0{,}5)\approx\SI{16.2}{\milli\second}$,
acima do limite. Para a descarga em no máximo \SI{5}{\milli\second},
$R\le t/[C\ln(3{,}3/0{,}5)]\approx\SI{2.65}{\kilo\ohm}$. Um diodo permite
separar os caminhos; \SI{8.2}{\kilo\ohm} e \SI{2.4}{\kilo\ohm} são escolhas
comerciais coerentes.}
\stepcounter{questao}
\resposta{A: $\zeta\approx0{,}447$, subamortecido; B: crítico; C: $\zeta\approx2{,}01$, superamortecido; escolher B}
\resolucao{Como $R_{\mathrm{crit}}\approx\SI{89.4}{\ohm}$, no RLC em série vale
$\zeta=R/R_{\mathrm{crit}}$. Assim A fornece $\zeta\approx0{,}447$ e oscila; B
fica praticamente em $\zeta=1$; C fornece $\zeta\approx2{,}01$ e é lento por
superamortecimento. Para rapidez sem sobressinal, B é a escolha adequada.}
\stepcounter{questao}
\resposta{$R_{\mathrm{efetivo}}\approx\SI{100}{\ohm}$; há cerca de \SI{80}{\ohm}
de resistência em série não contabilizada}
\resolucao{Para $s^2+(R/L)s+1/(LC)=0$, a soma dos polos é $-R/L$. A soma medida é
$-1000\ \si{\per\second}$; com $L=\SI{0.1}{\henry}$,
$R=1000L=\SI{100}{\ohm}$. O produto dos polos é aproximadamente
$100000\ \si{\per\second\squared}$, coerente com $1/(LC)$, confirmando que $L$ e
$C$ são plausíveis. Portanto existe cerca de \SI{80}{\ohm} extra, por ESR,
resistência do indutor, resistor de proteção, fonte ou fiação.}
\stepcounter{questao}
\resposta{$R_{\max}\approx\SI{61.4}{\kilo\ohm}$; escolha \SI{56}{\kilo\ohm};
$t\approx\SI{54.7}{\second}$; $P_0\approx\SI{2.86}{\watt}$; $E\approx\SI{37.6}{\joule}$}
\resolucao{Na descarga, $50=400e^{-t/RC}$. Assim
$R=t/[C\ln(400/50)]\approx\SI{61.4}{\kilo\ohm}$. Um valor de
\SI{56}{\kilo\ohm} fornece $t=RC\ln8\approx\SI{54.7}{\second}$. A potência
inicial é $P_0=V^2/R=400^2/56000\approx\SI{2.86}{\watt}$, portanto é prudente
usar resistor com margem de potência e tensão. A energia inicialmente armazenada,
e posteriormente dissipada, é $E=\tfrac12CV^2\approx\SI{37.6}{\joule}$.}

\gabarito[Nível 4]

\section{Eletromagnetismo}

\subsection{Conceitos iniciais e materiais magnéticos}
\stepcounter{questao}
\resposta{(b)}
\resolucao{Diferentemente das cargas elétricas, não existem monopolos magnéticos:
ao partir um ímã, cada pedaço volta a ter polos norte e sul. Por isso as linhas de
campo magnético são sempre \emph{fechadas} (sem começo nem fim), o que elimina
(d).}
\stepcounter{questao}
\resposta{(c)}
\resolucao{Ferromagnéticos (ferro, níquel, cobalto) têm $\mu_r$ da ordem de
centenas ou milhares --- por isso concentram intensamente as linhas de campo e são
usados em núcleos de transformadores e motores. Paramagnéticos têm
$\mu_r\gtrsim1$; diamagnéticos, $\mu_r\lesssim1$.}
\stepcounter{questao}
\resposta{O de $B_2 = \SI{1.8}{\tesla}$}
\resolucao{Como $B = \mu\,H$, para o mesmo $H$ o material com maior $B$ tem maior
permeabilidade. O valor $B_2 = \SI{1.8}{\tesla}$ é muito superior, indicando
$\mu$ elevado: material \textbf{ferromagnético}. O de $B_1 = \SI{0.002}{\tesla}$ é
paramagnético ou diamagnético (resposta muito fraca ao campo).}
\stepcounter{questao}
\resposta{$A$: paramagnético; $B$: diamagnético}
\resolucao{Materiais com $\mu_r$ ligeiramente maior que 1 são paramagnéticos e são
fracamente atraídos pelo campo. Materiais com $\mu_r$ ligeiramente menor que 1 são
diamagnéticos e são fracamente repelidos.}
\stepcounter{questao}
\resposta{Ver comentário}
\resolucao{Por terem $\mu_r$ muito elevado, os ferromagnéticos \emph{concentram} e
intensificam o campo magnético para uma mesma corrente de excitação, o que aumenta
o fluxo e, portanto, a eficiência do acoplamento magnético (mais tensão induzida,
mais torque). Além disso, oferecem um "caminho" de baixa relutância para o fluxo,
canalizando-o pelo núcleo em vez de dispersá-lo pelo ar. A desvantagem são as
perdas por histerese e correntes parasitas, minimizadas usando núcleos laminados
e ligas específicas (aço-silício).}
\stepcounter{questao}
\resposta{(a) $B_{\text{Fe}}/B_{\text{ar}} = 4000$;
(b) a razão \emph{cai}; (c) ver comentário}
\resolucao{(a) Como $B = \mu H = \mu_r\mu_0 H$, com $H$ fixo (mesma corrente e
geometria):
\[
\frac{B_{\text{Fe}}}{B_{\text{ar}}}=\frac{\mu_{r,\text{Fe}}}{\mu_{r,\text{ar}}}
=4000.
\]
O núcleo ferromagnético multiplica o campo por quatro mil --- é o que torna
viáveis transformadores e motores compactos.\\
(b) Na \textbf{saturação}, praticamente todos os domínios magnéticos já estão
alinhados; o material não consegue mais contribuir. A partir daí, aumentos de
corrente produzem acréscimos de $B$ apenas na proporção do ar
($\mu_r \to 1$ incremental). A razão $B_{\text{Fe}}/B_{\text{ar}}$, portanto,
\emph{diminui} progressivamente --- a curva $B\times H$ "achata".\\
(c) Porque o torque de um motor e a tensão de um transformador dependem do fluxo,
e o fluxo não pode crescer indefinidamente. Ultrapassado o joelho da curva de
saturação:
\begin{itemize}[leftmargin=1.6em,itemsep=1pt,topsep=2pt]
\item a corrente de magnetização aumenta acentuadamente (aquecendo o enrolamento) sem ganho
      proporcional de fluxo;
\item a forma de onda distorce, gerando harmônicos;
\item as perdas aumentam desproporcionalmente.
\end{itemize}
Por isso o projetista dimensiona a \emph{seção do núcleo} para operar
\textbf{abaixo} do joelho --- tipicamente em torno de \SI{1.5}{\tesla} para
aço-silício. \emph{Saturação é o limite físico que define o tamanho mínimo de
uma máquina elétrica.}}
\stepcounter{questao}
\resposta{Ver análise comentada}
\resolucao{\textbf{(a)} $B_r$ é a \textbf{remanência}: a indução que \emph{permanece}
no material quando o campo externo é removido ($H = 0$). É o que faz um prego
continuar magnetizado depois de encostar num ímã.\\
$H_c$ é a \textbf{coercividade}: o campo externo, de sentido oposto, necessário
para \emph{anular} a magnetização ($B = 0$). Quantifica a resistência do material à desmagnetização.

\textbf{(b)} A área interna do ciclo representa a \textbf{energia dissipada por
unidade de volume e por ciclo} --- as chamadas \emph{perdas por histerese}. Cada
inversão do campo exige reorientar os domínios magnéticos, e esse processo é
irreversível, gerando calor. A potência perdida é
\[
P_h = k\,f\,V\,(\text{área do ciclo}),
\]
proporcional à \emph{frequência} $f$: por isso máquinas de alta frequência exigem
materiais especiais.

\textbf{(c) Requisitos opostos.}
\begin{center}\small\sffamily
\begin{tabular}{@{}l l l@{}}
\toprule
 & \textbf{Núcleo de transformador} & \textbf{Ímã permanente}\\
\midrule
Ciclo desejado & \textbf{estreito} (fino) & \textbf{largo} (amplo)\\
$H_c$ & baixa & alta\\
$B_r$ & baixa & alta\\
Objetivo & minimizar perdas & reter magnetização\\
Material típico & ferro-silício, ferrite & alnico, neodímio (NdFeB)\\
Denominação & magneticamente \emph{mole} & magneticamente \emph{duro}\\
\bottomrule
\end{tabular}
\end{center}
\textbf{Raciocínio central:} o transformador inverte o campo 50 ou 60 vezes por
segundo --- cada inversão custa a área do ciclo em calor, logo o ciclo deve ser o
mais estreito possível. O ímã permanente, ao contrário, precisa \emph{resistir} à
desmagnetização: quanto maior a área do ciclo, maior é sua resistência à desmagnetização. \emph{O mesmo
gráfico, lido com objetivos opostos, leva a escolhas opostas de material.}}
\stepcounter{questao}
\resposta{$B=\SI{62.8}{\milli\tesla}$}
\resolucao{
\[
B=\mu_r\mu_0H=(500)(4\pi\times10^{-7})(100)=\pot{6.28}{-2}\,\si{\tesla}.
\]
\textbf{Sem o material} ($\mu_r=1$), o mesmo $H$ produziria apenas
$\SI{126}{\micro\tesla}$. O material multiplicou $B$ por $500$ --- e é
exatamente esse o papel de um núcleo magnético.\\
\textbf{Atenção:} $H$ é fixado pelas correntes; $B$ é o que o material
\emph{responde}. Trocar o núcleo muda $B$, mas não muda $H$.}
\stepcounter{questao}
\resposta{$M=\SI{49.9}{\kilo\ampere\per\meter}$}
\resolucao{De $B=\mu_0(H+M)$:
\[
M=\frac{B}{\mu_0}-H=\mu_rH-H=(\mu_r-1)H
=(499)(100)=\pot{4.99}{4}\,\si{\ampere\per\meter}.
\]
\textbf{Interpretação:} a magnetização do material equivale a uma corrente
superficial de $\SI{49.9}{\kilo\ampere\per\meter}$ --- quinhentas vezes maior
que a corrente de condução que a provocou.\\
\textbf{É por isso que o núcleo "amplifica":} ele não cria energia, apenas
alinha os momentos magnéticos já existentes no material, cujo efeito somado
supera em muito o do enrolamento.}
\stepcounter{questao}
\resposta{Paramagnético e diamagnético, respectivamente}
\resolucao{\textbf{Paramagnético} ($\chi>0$, pequeno): os momentos atômicos
alinham-se \emph{fracamente} com o campo, reforçando-o de forma quase
imperceptível. Exemplos: alumínio, platina, oxigênio líquido.\\
\textbf{Diamagnético} ($\chi<0$): o campo induz momentos \emph{opostos}, e o
material é levemente repelido. Exemplos: cobre, água, bismuto, grafite. Todo
material tem componente diamagnética; ela só é observável quando não há
efeito maior a mascará-la.\\
\textbf{Ordens de grandeza:} $|\chi|\sim10^{-5}$ a $10^{-4}$ nesses dois casos,
contra $\chi\sim10^{3}$ em ferromagnéticos --- diferença de sete ordens de
grandeza.}
\stepcounter{questao}
\resposta{$H=\SI{800}{\ampere\per\meter}$}
\resolucao{Da lei de Ampère aplicada ao caminho médio,
\[
H=\frac{NI}{\ell}=\frac{500(0{,}4)}{0{,}25}=\SI{800}{\ampere\per\meter}.
\]
\textbf{Observação central:} $H$ \textbf{não depende do material}. Ele é
determinado exclusivamente pelas correntes e pela geometria.\\
Trocar o núcleo de ar por ferro não altera $H$ --- altera $B$, e portanto o
fluxo. É essa separação que torna útil distinguir as duas grandezas.}
\stepcounter{questao}
\resposta{$\mathcal{R}=\pot{7.96}{5}\,\si{\per\henry}$;
$\Phi=\SI{0.251}{\milli\weber}$; $B=\SI{2.51}{\tesla}$ --- \textbf{saturado}}
\resolucao{(a)
\[
\mathcal{R}=\frac{\ell}{\mu_r\mu_0A}
=\frac{0{,}10}{(1000)(4\pi\times10^{-7})(\pot{1}{-4})}
=\pot{7.96}{5}\,\si{\per\henry}.
\]
(b)
\[
\Phi=\frac{NI}{\mathcal{R}}=\frac{200}{\pot{7.96}{5}}
=\pot{2.51}{-4}\,\si{\weber};
\qquad
B=\frac{\Phi}{A}=\SI{2.51}{\tesla}.
\]
(c) \textbf{Resultado impossível.} Nenhum aço comum sustenta
$\SI{2.51}{\tesla}$ --- todos saturam entre $1{,}5$ e $\SI{2}{\tesla}$.\\
\textbf{O que realmente ocorreria:} conforme $B$ se aproxima da saturação,
$\mu_r$ desaba de $1000$ para valores próximos de $1$. A relutância dispara e o
fluxo estaciona em torno de $B_{sat}A=\SI{0.15}{\milli\weber}$.\\
\textbf{Regra de método:} calcular com $\mu_r$ constante e \emph{depois}
verificar $B$ contra $B_{sat}$. Se o valor exceder, o modelo linear não se
aplica e é preciso recorrer à curva $B\times H$ real.}
\stepcounter{questao}
\resposta{$\SI{628}{\micro\tesla}$; $\SI{628}{\micro\tesla}$;
$\SI{2.51}{\tesla}$}
\resolucao{
\[
\text{ar}:\ B=\mu_0(500)=\SI{628}{\micro\tesla};
\]
\[
\text{alumínio}:\ B=(1{,}00002)\mu_0(500)=\SI{628}{\micro\tesla};
\]
\[
\text{ferro-silício}:\ B=(4000)\mu_0(500)=\SI{2.51}{\tesla}
\ \text{(saturaria antes)}.
\]
\textbf{Duas leituras.}\\
\emph{(i)} O alumínio é \emph{indistinguível} do ar para fins práticos --- sua
suscetibilidade de $\pot{2}{-5}$ altera $B$ na quinta casa decimal. É por isso
que se pode enrolar bobinas em carretéis de alumínio sem consequência magnética.\\
\emph{(ii)} O ferro-silício produziria $\SI{2.51}{\tesla}$ \emph{se fosse
linear} --- mas ele satura antes. Com $H=\SI{500}{\ampere\per\meter}$, um
ferro-silício real já está bem dentro do joelho da curva, com $B$ real em torno
de $1{,}5$ a $\SI{1.7}{\tesla}$.}
\stepcounter{questao}
\resposta{$\mu_r=2390$ e $995$}
\resolucao{(a)
\[
\mu_r=\frac{B}{\mu_0H}:\quad
\frac{1{,}2}{(4\pi\times10^{-7})(400)}=2390;
\qquad
\frac{1{,}5}{(4\pi\times10^{-7})(1200)}=995.
\]
(b) \textbf{A permeabilidade caiu para menos da metade} enquanto $H$
triplicou.\\
\textbf{Interpretação:} o material está entrando na região de \emph{joelho} da
curva de magnetização. Os domínios já estão em boa parte alinhados, e cada
acréscimo de $H$ produz cada vez menos acréscimo de $B$.\\
\textbf{Consequência prática:} triplicar a corrente aumentou o fluxo em apenas
$25\%$. Operar nessa região é ineficiente --- o cobre gasta muito para
render pouco.\\
\textbf{Lição de fundo:} $\mu_r$ de catálogo é sempre um valor \emph{para uma
condição específica}. Usá-lo fora dessa condição é um erro comum e silencioso.}
\stepcounter{questao}
\resposta{Normal para o ponto de operação; diferencial para pequenos sinais
sobrepostos}
\resolucao{\textbf{Permeabilidade normal:} razão entre $B$ e $H$ no ponto de
operação. Ela governa o \emph{fluxo total} e, portanto, a relutância e o
dimensionamento do núcleo.\\
\textbf{Permeabilidade diferencial:} inclinação da curva naquele ponto. Ela
governa a resposta a \emph{pequenas variações} sobrepostas --- e portanto a
indutância incremental.\\
\textbf{Onde diferem drasticamente:} perto da saturação, $B/H$ ainda é grande
(centenas), mas $\mathrm{d}B/\mathrm{d}H$ já caiu para quase $\mu_0$. Um
indutor com corrente contínua elevada pode ter indutância incremental dezenas
de vezes menor que a nominal.\\
\textbf{Consequência de projeto:} em filtros com componente contínua --- muito
comuns em fontes chaveadas --- é a permeabilidade \emph{diferencial} que
determina o desempenho, e é ela que precisa ser consultada no catálogo.}
\stepcounter{questao}
\resposta{Acima de $T_C$ o material torna-se paramagnético e perde
praticamente toda a permeabilidade}
\resolucao{\textbf{O que ocorre acima de $T_C$.} A agitação térmica supera a
interação de troca que mantém os momentos atômicos alinhados dentro dos
domínios. A estrutura de domínios se desfaz, e o material passa de
\textbf{ferromagnético} a \textbf{paramagnético}:
\[
\mu_r:\ 10^{3}\text{--}10^{4}\ \longrightarrow\ \approx1.
\]
\textbf{A transição é abrupta} --- trata-se de uma transição de fase de segunda
ordem, e $\mu_r$ despenca em poucos graus em torno de $T_C$.\\
\textbf{Consequências práticas.}
\begin{itemize}
\item \textbf{Ferrites são as mais vulneráveis.} Com $T_C$ próxima de
$\SI{250}{\celsius}$, e temperatura de operação já em torno de
$\SI{100}{\celsius}$, a margem é pequena. Um transformador de fonte chaveada
que superaqueça perde indutância, a corrente dispara e ele se destrói --- fuga
térmica.
\item \textbf{Aplicação deliberada:} panelas de indução com fundo cuja $T_C$ é
próxima da temperatura desejada param de aquecer sozinhas ao atingi-la. É um
termostato sem sensor.
\item \textbf{Desmagnetização por aquecimento:} aquecer um ímã acima de
$T_C$ e resfriá-lo sem campo o deixa completamente desmagnetizado.
\end{itemize}}
\stepcounter{questao}
\resposta{Redução de $20\%$}
\resolucao{
\[
\frac{\Delta B_r}{B_r}=(-0{,}002)(100)=-0{,}20=-20\%.
\]
\textbf{Reversível ou não?} Até certo limite, a perda é \textbf{reversível} --- o
ímã recupera a magnetização ao esfriar. Acima de uma temperatura crítica
(bem abaixo de $T_C$), parte da perda torna-se \textbf{irreversível}, porque
domínios se reorientam permanentemente.\\
\textbf{Comparação entre materiais:}
\begin{center}
\begin{tabular}{lcc}
\hline
Material & Coef. de $B_r$ & $T_{\max}$ típica\\
\hline
Ferrite & $\SI{-0.2}{\percent\per\celsius}$ & $\SI{250}{\celsius}$\\
Neodímio (NdFeB) & $\SI{-0.12}{\percent\per\celsius}$ & $\SI{80}{}$ a
 $\SI{200}{\celsius}$\\
Samário-cobalto & $\SI{-0.03}{\percent\per\celsius}$ & $\SI{300}{\celsius}$\\
\hline
\end{tabular}
\end{center}
\textbf{Compromisso revelador:} o neodímio tem a maior densidade de energia mas
péssima estabilidade térmica; o samário-cobalto é bem mais estável e mais caro.
Em motores que aquecem, a escolha raramente é o material mais forte.}
\stepcounter{questao}
\resposta{$\SI{12}{\watt}$ e $\SI{80}{\watt}$}
\resolucao{
\[
P_h=f\times(\text{área})\times V:
\quad
60(200)(\pot{1}{-3})=\SI{12}{\watt};
\qquad
400(200)(\pot{1}{-3})=\SI{80}{\watt}.
\]
\textbf{Proporcionalidade linear com $f$} --- cada ciclo custa a mesma energia,
e mais ciclos por segundo significam mais potência.\\
\textbf{Contraste com as correntes parasitas}, cujas perdas vão com
$f^{2}$. Em baixa frequência a histerese domina; em alta, as parasitas.\\
\textbf{Consequência:} o mesmo núcleo que funciona bem em $\SI{60}{\hertz}$
pode ser inviável em $\SI{400}{\hertz}$ --- e é por isso que equipamentos
aeronáuticos, que operam nessa frequência, usam ligas especiais de lâminas
finíssimas.}
\stepcounter{questao}
\resposta{Razão de $10^{4}$ na coercividade; moles para núcleos, duros para ímãs
permanentes}
\resolucao{\textbf{Coercividade $H_c$} é o campo necessário para \emph{anular} a
magnetização remanente. Ela mede o quanto o material "resiste" a mudar de
estado.
\begin{center}
\begin{tabular}{lccl}
\hline
Tipo & $H_c$ & Área do ciclo & Aplicação\\
\hline
Mole & baixa & pequena & núcleos de transformador e motor\\
Duro & alta & grande & ímãs permanentes\\
\hline
\end{tabular}
\end{center}
\textbf{Materiais moles} (ferro-silício, permalloy, ferrites moles):
magnetizam e desmagnetizam facilmente, com pouca perda por ciclo. É o que se
quer em algo que inverte a magnetização $120$ vezes por segundo.\\
\textbf{Materiais duros} (alnico, ferrite dura, NdFeB): resistem à
desmagnetização, mantendo o campo sem corrente alguma.\\
\textbf{O critério é a área do ciclo:} pequena para quem vai percorrê-lo muitas
vezes, grande para quem precisa permanecer magnetizado.}
\stepcounter{questao}
\resposta{A curva inicial parte da origem; o ciclo nunca mais passa por ela}
\resolucao{\textbf{Material virgem:} os domínios estão orientados
aleatoriamente, e a magnetização líquida é nula. Aplicando $H$ crescente, a
curva parte da \textbf{origem} e sobe até a saturação --- é a \emph{curva de
magnetização inicial}.\\
\textbf{Ao reduzir $H$ a zero}, a magnetização \emph{não} retorna a zero: resta
a \textbf{remanência} $B_r$. Os domínios que se reorientaram não voltam
espontaneamente.\\
\textbf{A partir daí}, ciclando $H$ entre $\pm H_{\max}$, o material percorre
um laço fechado que \textbf{nunca mais passa pela origem}.\\
\textbf{Como voltar ao estado virgem:} aplicar campo alternado de amplitude
decrescente até zero (desmagnetização por ciclos decrescentes), ou aquecer acima
de $T_C$ e resfriar sem campo.\\
\textbf{Importância prática:} ao levantar experimentalmente a curva
$B\times H$, é preciso desmagnetizar o material antes --- caso contrário, os
primeiros pontos refletem a história anterior, não a propriedade do material.}
\stepcounter{questao}
\resposta{O fluxo remanente soma-se ao fluxo imposto, podendo levar o núcleo a
saturação profunda}
\resolucao{\textbf{Ao desligar}, o núcleo retém fluxo remanente
$\Phi_r$ --- consequência direta da histerese.\\
\textbf{Ao religar}, o fluxo evolui a partir desse valor. No pior caso, a tensão
é aplicada no instante em que exigiria fluxo de mesma polaridade, e o fluxo
total pode chegar a
\[
\Phi\approx\Phi_r+2\Phi_{\max},
\]
muito além da saturação.\\
\textbf{Consequência:} saturado, o núcleo perde permeabilidade, a indutância
desaba e a corrente é limitada apenas pela resistência do enrolamento. A
corrente de \emph{inrush} pode atingir dez a vinte vezes a nominal, por alguns
ciclos.\\
\textbf{Soluções empregadas:}
\begin{itemize}
\item \textbf{Chaveamento sincronizado:} energizar no instante correto da onda
de tensão.
\item \textbf{Resistor ou NTC de pré-carga}, curto-circuitado após alguns
ciclos.
\item \textbf{Proteção com curva temporizada}, que tolera o transitório sem
atuar.
\end{itemize}
\textbf{Note a conexão:} este fenômeno só existe \emph{porque} há histerese ---
um núcleo hipotético sem remanência não teria \emph{inrush} por essa causa.}
\stepcounter{questao}
\resposta{$\mathcal{R}_{fe}=\pot{2.39}{5}$, $\mathcal{R}_{ar}=\pot{3.18}{6}$;
o ar responde por $93\%$}
\resolucao{
\[
\mathcal{R}_{fe}=\frac{0{,}30}{(2000)\mu_0(\pot{5}{-4})}
=\pot{2.39}{5}\,\si{\per\henry};
\]
\[
\mathcal{R}_{ar}=\frac{\pot{2}{-3}}{\mu_0(\pot{5}{-4})}
=\pot{3.18}{6}\,\si{\per\henry}.
\]
\[
\mathcal{R}_{tot}=\pot{3.42}{6};
\qquad
\frac{\mathcal{R}_{ar}}{\mathcal{R}_{tot}}=93\%.
\]
\textbf{Dois milímetros de ar dominam trinta centímetros de ferro} --- porque a
razão de permeabilidades ($2000$) supera de longe a razão de comprimentos
($150$).\\
\textbf{As relutâncias somam-se em série}, exatamente como resistências. E o
fluxo é o mesmo nos dois trechos, o que garante $B$ igual --- já que a seção
também é a mesma.\\
\textbf{Consequência:} o comportamento do circuito é governado quase
inteiramente pelo entreferro, e a não linearidade do ferro fica mascarada. É
esse efeito que \emph{lineariza} indutores com entreferro.}
\stepcounter{questao}
\resposta{$\vec B=\mu_0(\vec H+\vec M)$; ao inserir o núcleo, $H$ não muda e
$B$ aumenta}
\resolucao{(a)
\[
\vec B=\mu_0\left(\vec H+\vec M\right),
\]
com $[H]=[M]=\si{\ampere\per\meter}$ e $[B]=\si{\tesla}$.\\
\emph{($H$ é às vezes chamado "intensidade de campo magnético" e $B$
"densidade de fluxo" --- nomenclatura infeliz, porque sugere que $H$ é o campo
principal.)}\\
(b) \textbf{Ao inserir o núcleo com corrente constante:}
\begin{itemize}
\item $H=NI/\ell$ \textbf{não muda} --- ele depende só das correntes de
condução.
\item $M$ \textbf{passa de zero a $(\mu_r-1)H$} --- o material se magnetiza.
\item $B$ \textbf{aumenta por} $\mu_r$ --- consequência das duas anteriores.
\end{itemize}
(c) \textbf{Qual é mais fundamental: $\vec B$.} Argumentos:
\begin{itemize}
\item É $\vec B$ que aparece na força de Lorentz
($\vec F=q\vec v\times\vec B$) --- é ele que age sobre cargas.
\item É $\vec B$ que aparece na lei de Faraday --- é seu fluxo que induz f.e.m.
\item $\nabla\cdot\vec B=0$ é lei universal; $\nabla\cdot\vec H=-\nabla\cdot
\vec M$ não é nula em geral.
\end{itemize}
\textbf{Para que serve $\vec H$, então:} ele é conveniente porque sua circulação
depende \emph{apenas} das correntes de condução ($\oint\vec H\cdot
\mathrm{d}\vec\ell=NI$), sem precisar contabilizar as correntes de
magnetização do material. É uma grandeza \emph{auxiliar} de enorme utilidade
prática --- mas auxiliar.}
\stepcounter{questao}
\resposta{Deslocamento reversível, deslocamento irreversível e rotação;
$\mu_r$ máxima no trecho intermediário}
\resolucao{(a) \textbf{Três regiões.}
\begin{enumerate}
\item \textbf{Campos baixos --- deslocamento reversível de paredes.} As paredes
entre domínios deslocam-se ligeiramente, favorecendo os domínios alinhados com
$H$. Removido o campo, elas retornam. A curva é pouco inclinada.
\item \textbf{Campos médios --- deslocamento irreversível (saltos de
Barkhausen).} As paredes superam obstáculos (impurezas, contornos de grão) em
saltos abruptos. É aqui que a curva é mais \textbf{íngreme} e onde nasce a
histerese, pois os saltos não se desfazem ao reduzir $H$.
\item \textbf{Campos altos --- rotação dos domínios.} Já não há paredes a
mover; os momentos giram progressivamente até se alinharem com $H$. A curva
achata-se: é o \emph{joelho} e depois a saturação.
\end{enumerate}
(b) \textbf{$\mu_r$ é máxima na região intermediária}, onde
$\mathrm{d}B/\mathrm{d}H$ é maior. Nem no início (paredes presas) nem no fim
(domínios já alinhados).\\
\textbf{Consequência contraintuitiva:} a permeabilidade em campos muito baixos é
\emph{menor} que a máxima --- daí catálogos especificarem $\mu_i$ (inicial) e
$\mu_{\max}$ separadamente.\\
(c) \textbf{Região recomendada:} abaixo do joelho, tipicamente com
$B_{\max}$ entre $60$ e $80\%$ de $B_{sat}$. Isso aproveita a alta
permeabilidade sem risco de saturação em transitórios.}
\stepcounter{questao}
\resposta{Interseção da curva $B\times H$ do material com a reta imposta pelo
entreferro}
\resolucao{(a) \textbf{Formulação.} A lei de Ampère no circuito magnético dá
\[
NI=H_{fe}\,\ell_{fe}+H_g\,g,
\]
e no entreferro $H_g=B/\mu_0$ (com $B$ igual ao do ferro, por conservação do
fluxo e seção constante). Então
\[
NI=H_{fe}\ell_{fe}+\frac{B}{\mu_0}g.
\]
Duas incógnitas ($B$ e $H_{fe}$) e uma equação --- a segunda relação é a
\emph{curva do material}, que não tem forma algébrica simples.\\
(b) \textbf{Método gráfico.} Reescrevendo,
\[
B=\frac{\mu_0\,NI}{g}-\frac{\mu_0\ell_{fe}}{g}\,H_{fe},
\]
que é uma \textbf{reta} no plano $B\times H$ --- a \emph{reta de entreferro} ou
\emph{reta de carga magnética}.\\
Traçando-a sobre a curva $B\times H$ do material, o \textbf{ponto de operação} é
a interseção.\\
\emph{(É exatamente o mesmo procedimento da reta de carga usada com diodos: uma
característica não linear do dispositivo e uma reta imposta pelo circuito.)}\\
(c) \textbf{Por que o entreferro estabiliza.} A inclinação da reta é
$-\mu_0\ell_{fe}/g$: quanto \emph{maior} o entreferro, mais \textbf{deitada} a
reta.\\
Uma reta deitada cruza a curva do material em uma região de inclinação
comparável --- e ali a interseção é \emph{pouco sensível} a variações de
$\mu_r$ do material. Trocar o lote de ferro, ou aquecê-lo, quase não desloca o
ponto.\\
\textbf{Sem entreferro}, a reta é quase vertical e a interseção depende
inteiramente da curva do material --- que varia com temperatura, lote e
história magnética. \textbf{O entreferro compra previsibilidade ao custo de mais
ampère-espiras.}}
\stepcounter{questao}
\resposta{Aço-silício: alto $B_{sat}$, baixa resistividade, até
$\si{\kilo\hertz}$; ferrite: baixo $B_{sat}$, altíssima resistividade,
até $\si{\mega\hertz}$}
\resolucao{(a)
\begin{center}
\begin{tabular}{lcc}
\hline
Propriedade & Aço-silício & Ferrite (MnZn)\\
\hline
$B_{sat}$ & $1{,}5$--$\SI{2}{\tesla}$ & $0{,}3$--$\SI{0.5}{\tesla}$\\
Resistividade & $\approx\pot{5}{-7}\,\si{\ohm\meter}$ &
 $10^{-1}$ a $\SI{10}{\ohm\meter}$\\
$\mu_r$ típica & $2000$--$8000$ & $1000$--$5000$\\
Faixa útil & até $\approx\SI{1}{\kilo\hertz}$ & $\SI{20}{\kilo\hertz}$ a
 $\si{\mega\hertz}$\\
\hline
\end{tabular}
\end{center}
\textbf{Note a resistividade:} a ferrite é um \emph{cerâmico}, com
resistividade até $10^{7}$ vezes maior que a do aço.\\
(b) \textbf{Baixa frequência} ($\SI{50}{}$--$\SI{400}{\hertz}$): aço-silício.
Seu $B_{sat}$ elevado permite núcleos menores para a mesma potência, e as perdas
parasitas são controladas por laminação.\\
\textbf{Alta frequência} ($>\SI{20}{\kilo\hertz}$): ferrite. Nenhuma laminação
seria fina o bastante para conter as correntes parasitas no aço.\\
(c) \textbf{Justificativa física.} As perdas por correntes parasitas vão com
\[
P_F\propto\frac{B^{2}f^{2}t^{2}}{\rho}.
\]
Em alta frequência, o termo $f^{2}$ explode. Há duas formas de compensar:
reduzir $t$ (laminação) ou aumentar $\rho$.\\
A ferrite ataca $\rho$ --- e ganha por muitas ordens de grandeza, o que nenhuma
laminação alcança.\\
\textbf{O preço:} $B_{sat}$ baixo obriga a usar seção maior ou aceitar menos
fluxo. \textbf{Mas em alta frequência isso não importa}, porque
$V=4{,}44Nf\Phi_{\max}$ --- com $f$ mil vezes maior, o fluxo necessário é mil
vezes menor. \textbf{É a alta frequência que torna a ferrite viável, e a ferrite
que torna a alta frequência viável.}}
\stepcounter{questao}
\resposta{Ao ar, o ponto de operação desce na curva de desmagnetização; o ferro
o eleva}
\resolucao{(a) \textbf{Ao ar livre.} O ímã precisa "empurrar" o fluxo por um
caminho de alta relutância --- o próprio ar em torno dele. Isso equivale a um
entreferro enorme, e a reta de carga fica muito deitada.\\
O ponto de operação desce na \emph{curva de desmagnetização} (segundo quadrante
do laço), afastando-se de $B_r$. O ímã fica submetido a campo desmagnetizante
produzido por ele mesmo.\\
(b) \textbf{Aproximando ferro.} O ferro oferece caminho de baixa relutância, e a
reta de carga se levanta. O ponto de operação sobe, aproximando-se de
$B_r$ --- o ímã fica "mais forte".\\
É por isso que um ímã em ferradura fechado por uma armadura de ferro sustenta
muito mais peso que aberto.\\
(c) \textbf{Quebra-fluxo (\emph{keeper}).} É uma peça de ferro doce colocada
entre os polos durante o armazenamento. Ela:
\begin{itemize}
\item mantém o ponto de operação alto, evitando desmagnetização progressiva;
\item protege contra campos externos adversos;
\item reduz a atração de limalha e detritos.
\end{itemize}
\textbf{Ressalva histórica:} a prática era essencial com alnico, cuja curva de
desmagnetização é frágil. Com neodímio moderno --- cuja curva é praticamente
reta no segundo quadrante --- o efeito é bem menor, embora o quebra-fluxo
continue útil por outras razões.}
\stepcounter{questao}
\resposta{$(BH)_{\max}\approx\SI{270}{\kilo\joule\per\meter\cubed}$ contra
$\SI{25}{\kilo\joule\per\meter\cubed}$}
\resolucao{(a) Para curva de desmagnetização aproximadamente \emph{reta} (caso
do neodímio), o produto máximo ocorre no ponto médio:
\[
(BH)_{\max}\approx\frac{B_rH_c}{4}
=\frac{(1{,}2)(\pot{9}{5})}{4}
=\SI{270}{\kilo\joule\per\meter\cubed}.
\]
(b) Para a ferrite:
\[
(BH)_{\max}\approx\frac{(0{,}4)(\pot{2.5}{5})}{4}
=\SI{25}{\kilo\joule\per\meter\cubed},
\]
cerca de \textbf{onze vezes menor}.\\
(c) \textbf{Significado.} O produto de energia mede a energia magnética
disponível \emph{por unidade de volume} do ímã. Ele é a figura de mérito que
determina quão pequeno o ímã pode ser para produzir dado campo em dado
entreferro.\\
\textbf{Consequência prática:} para o mesmo desempenho, um ímã de neodímio ocupa
cerca de um décimo do volume de um de ferrite. Isso é o que viabilizou
servomotores compactos, fones de ouvido pequenos e discos rígidos.\\
\textbf{Contrapartida:} o neodímio é caro, oxida (precisa de revestimento),
desmagnetiza acima de $\SI{80}{}$--$\SI{200}{\celsius}$ e depende de
terras-raras. A ferrite é barata, quimicamente estável e térmicamente mais
tolerante --- e por isso continua dominando em volume de produção.}
\stepcounter{questao}
\resposta{Aquecimento acima de $T_C$ desordena os domínios; campo alternado os
distribui simetricamente}
\resolucao{(a) \textbf{Aquecimento acima de $T_C$.} A agitação térmica destrói o
alinhamento dentro dos domínios. Resfriando \emph{sem campo aplicado}, os
domínios se reformam em orientações aleatórias, e a magnetização líquida é
nula.\\
\textbf{Campo alternado decrescente.} Aplica-se campo alternado de amplitude
inicialmente alta (o bastante para saturar) e reduz-se lentamente a zero. O
material percorre laços de histerese cada vez menores, espiralando até a origem.
Ao final, tantos domínios apontam em um sentido quanto no outro.\\
(b) \textbf{Comparação.}
\begin{center}
\begin{tabular}{lcc}
\hline
& Aquecimento & Campo alternado\\
\hline
Eficácia & completa & muito boa\\
Danifica a peça? & possivelmente & não\\
Aplicável a peça montada? & raramente & sim\\
Rapidez & lento & segundos\\
\hline
\end{tabular}
\end{center}
\textbf{Na prática industrial, usa-se campo alternado} --- em desmagnetizadores
de ferramentas, de peças usinadas e de fitas magnéticas. O aquecimento é
inviável para peças que não toleram $\SI{770}{\celsius}$.\\
\textbf{Cuidado essencial:} a redução do campo deve ser \emph{lenta} em relação
ao período. Desligar bruscamente o desmagnetizador pode deixar a peça
magnetizada no valor correspondente à última meia-onda --- pior que antes.\\
\textbf{Alternativa equivalente:} afastar lentamente a peça de um campo
alternado fixo, o que produz o mesmo decaimento gradual.}
\stepcounter{questao}
\resposta{$N$ cai $20\%$; perdas por histerese sobem cerca de $60\%$ e
parasitas $56\%$}
\resolucao{(a) De $V=4{,}44Nf\Phi_{\max}=4{,}44NfB_{\max}A$, com $V$, $f$ e
$A$ fixos:
\[
N\propto\frac{1}{B_{\max}}
\;\Longrightarrow\;
\frac{N'}{N}=\frac{1{,}2}{1{,}5}=0{,}80.
\]
\textbf{Vinte por cento menos espiras} --- menos cobre, menos resistência, menos
perdas no enrolamento.\\
(b) \textbf{Perdas no núcleo aumentam.}
\begin{itemize}
\item \textbf{Parasitas:} $P_F\propto B_{\max}^{2}$, logo
$(1{,}5/1{,}2)^{2}=1{,}56$ --- aumento de $56\%$.
\item \textbf{Histerese:} $P_h\propto B_{\max}^{n}$ com $n\approx1{,}6$ a
$2$ (expoente de Steinmetz), dando aumento de $60$ a $56\%$.
\end{itemize}
(c) \textbf{O compromisso.}
\begin{center}
\begin{tabular}{lc}
\hline
Elevar $B_{\max}$ & Efeito\\
\hline
Espiras & $-20\%$\\
Perdas no cobre & diminuem\\
Perdas no núcleo & $+\approx58\%$\\
Volume e custo & diminuem\\
Margem para saturação & \textbf{diminui}\\
\hline
\end{tabular}
\end{center}
\textbf{Existe um ótimo econômico}, em que a soma das perdas é mínima --- e ele
depende do custo relativo do cobre e do aço e do regime de carga.\\
\textbf{Mas há um limite duro:} $\SI{1.5}{\tesla}$ já está próximo do joelho. A
margem para sobretensões e para o transitório de energização encolhe, e o risco
de \emph{inrush} severo cresce. \textbf{Projetar perto da saturação economiza
material e compra problemas operacionais.}}
\stepcounter{questao}
\resposta{$\mu$ depende de $H$, da temperatura, da frequência e da história
magnética}
\resolucao{\textbf{Quatro dependências, todas relevantes.}
\begin{enumerate}
\item \textbf{Da intensidade $H$.} É a mais óbvia: $\mu_r$ parte de um valor
inicial, cresce até um máximo e despenca na saturação. Variação de uma ordem de
grandeza ou mais.
\item \textbf{Da temperatura.} $\mu_r$ tende a \emph{crescer} ao se aproximar de
$T_C$ (efeito Hopkinson) e depois colapsa abruptamente. Isso torna o
comportamento térmico não monotônico e difícil de prever.
\item \textbf{Da frequência.} Em alta frequência, as paredes de domínio não
conseguem acompanhar o campo --- há \emph{relaxação}. A permeabilidade efetiva
cai, e aparece uma componente de perda. Catálogos de ferrite trazem
$\mu$ complexo, com parte real e imaginária.
\item \textbf{Da história magnética.} Por causa da histerese, o mesmo $H$ pode
corresponder a diferentes $B$, conforme o material esteja subindo ou descendo o
laço.
\end{enumerate}
\textbf{Consequência para o projetista.} Um valor único de $\mu_r$ em catálogo é
sempre uma \emph{simplificação} referida a condições específicas --- geralmente
campo baixo, temperatura ambiente e baixa frequência.\\
\textbf{Prática correta:} usar $\mu_r$ para estimativas iniciais, e recorrer às
\emph{curvas} do fabricante (de $B\times H$, de perdas por frequência, de
$\mu$ por temperatura) para o dimensionamento final. Projetos de fontes
chaveadas que ignoram isso costumam funcionar na bancada e falhar em campo.}
\stepcounter{questao}
\resposta{Ferromagnéticos aumentam $L$ por permeabilidade; não ferromagnéticos
reduzem $L$ e o fator de qualidade por correntes parasitas}
\resolucao{(a) \textbf{Ferromagnéticos.} A aproximação do material oferece
caminho de menor relutância, aumentando a indutância da bobina sensora. A
variação de $L$ é detectada pelo circuito.\\
Simultaneamente, correntes parasitas no aço amortecem o oscilador.\\
(b) \textbf{Não ferromagnéticos.} Alumínio e cobre têm $\mu_r\approx1$ --- não
alteram a relutância. Mas são \emph{bons condutores}, e o campo alternado induz
neles correntes parasitas intensas.\\
Pela lei de Lenz, essas correntes criam campo oposto, o que \textbf{reduz} a
indutância efetiva e, sobretudo, \textbf{drena energia} do oscilador, reduzindo
seu fator de qualidade.\\
(c) \textbf{Por que o alcance difere.} Sensores indutivos comuns são
especificados para aço, e apresentam \emph{fatores de correção}:
\begin{center}
\begin{tabular}{lc}
\hline
Material & Fator de alcance\\
\hline
Aço & $1{,}00$\\
Aço inoxidável & $\approx0{,}70$\\
Latão & $\approx0{,}50$\\
Alumínio & $\approx0{,}40$\\
Cobre & $\approx0{,}30$\\
\hline
\end{tabular}
\end{center}
\textbf{Razão:} no aço, os dois efeitos (permeabilidade e parasitas) somam-se.
Nos não ferromagnéticos, resta apenas o segundo --- e ele é tanto mais fraco
quanto \emph{melhor} o condutor, porque as correntes se concentram numa
película superficial cada vez mais fina.\\
\textbf{Consequência de aplicação:} um sensor ajustado para detectar peças de
aço a $\SI{10}{\milli\meter}$ só detectará cobre a $\SI{3}{\milli\meter}$ ---
erro de projeto frequente em linhas de montagem com peças mistas.}
\stepcounter{questao}
\resposta{$w=\oint H\,\mathrm{d}B$ por ciclo; $P=fVw$}
\resolucao{(a) \textbf{Dedução.} A potência instantânea fornecida ao núcleo pela
fonte é $p=\varepsilon i$. Com $\varepsilon=NA\,\mathrm{d}B/\mathrm{d}t$ e
$i=H\ell/N$:
\[
p=NA\frac{\mathrm{d}B}{\mathrm{d}t}\cdot\frac{H\ell}{N}
=(A\ell)\,H\,\frac{\mathrm{d}B}{\mathrm{d}t}
=V\,H\,\frac{\mathrm{d}B}{\mathrm{d}t}.
\]
A energia por ciclo é
\[
W=\int p\,\mathrm{d}t=V\oint H\,\mathrm{d}B,
\]
e $\oint H\,\mathrm{d}B$ é precisamente a \textbf{área do laço} no plano
$B\times H$, com unidade $\si{\joule\per\meter\cubed}$.\\
\textbf{Se o material fosse sem histerese}, o laço colapsaria em uma curva
única, a integral de ida cancelaria a de volta e a energia por ciclo seria
\textbf{zero} --- toda a energia magnética seria devolvida.\\
(b)
\[
P_h=f\,V\oint H\,\mathrm{d}B,
\]
\textbf{linear em $f$}. (As perdas por correntes parasitas, distintas destas,
vão com $f^{2}$.)\\
(c) \textbf{Origem microscópica.} A perda vem dos \textbf{saltos de Barkhausen}:
as paredes de domínio ficam presas em impurezas, defeitos e contornos de grão, e
se libertam abruptamente quando o campo cresce o suficiente.\\
Cada salto é um processo \textbf{irreversível} --- a energia elástica acumulada
converte-se em vibração da rede, ou seja, em \textbf{calor}.\\
\textbf{Consequência de engenharia:} materiais mais puros e com grãos maiores e
melhor orientados têm menos obstáculos, laços mais estreitos e menores perdas.
É exatamente isso que o processo de \emph{grão orientado} do aço-silício
persegue.}
\stepcounter{questao}
\resposta{$B$ é contínuo; $H$ é descontínuo, na razão inversa das
permeabilidades}
\resolucao{(a) \textbf{Duas equações.}\\
\emph{Conservação do fluxo} (equivalente à LKC):
\[
\Phi_1=\Phi_2
\;\Longrightarrow\;
B_1A=B_2A
\;\Longrightarrow\;
B_1=B_2=B.
\]
\emph{Lei de Ampère} (equivalente à LKT):
\[
NI=H_1\ell_1+H_2\ell_2.
\]
(b) Com $H_i=B/\mu_i$:
\[
NI=B\left(\frac{\ell_1}{\mu_1}+\frac{\ell_2}{\mu_2}\right)
\;\Longrightarrow\;
B=\frac{NI}{\dfrac{\ell_1}{\mu_1}+\dfrac{\ell_2}{\mu_2}}.
\]
\[
H_1=\frac{B}{\mu_1};
\qquad
H_2=\frac{B}{\mu_2};
\qquad
\frac{H_1}{H_2}=\frac{\mu_2}{\mu_1}.
\]
(c) \textbf{Qual grandeza é contínua --- e por quê.}
\begin{itemize}
\item \textbf{$B$ é contínuo} através da interface (na componente normal),
porque $\nabla\cdot\vec B=0$: o fluxo que chega precisa continuar. Não há
"acúmulo de fluxo" possível.
\item \textbf{$H$ é descontínuo}, na razão inversa das permeabilidades. O
material de menor $\mu$ exige \emph{muito mais} $H$ para sustentar o mesmo
$B$.
\end{itemize}
\textbf{Exemplo numérico revelador.} Com $\mu_{r1}=2000$ (ferro) e
$\mu_{r2}=2$ (um material quase não magnético), $H_2$ seria \emph{mil vezes}
$H_1$ --- e praticamente toda a força magnetomotriz seria consumida naquele
trecho, ainda que ele fosse curtíssimo.\\
\textbf{É exatamente o que acontece no entreferro} ($\mu_r=1$), e por isso ele
domina o circuito. Um entreferro não é um caso especial: é apenas o caso extremo
deste problema de dois materiais.}
\stepcounter{questao}
\resposta{Canal horizontal $\propto H$ (corrente por shunt); vertical
$\propto B$ (f.e.m. integrada)}
\resolucao{(a) \textbf{Circuito.} Um núcleo toroidal com dois enrolamentos:
primário $N_1$ excitado por fonte alternada e secundário $N_2$ para medição.
\begin{itemize}
\item \textbf{Canal horizontal ($X$):} tensão sobre um resistor \emph{shunt}
$R_s$ em série com o primário. Como $H=N_1i_1/\ell$, essa tensão é proporcional
a $H$.
\item \textbf{Canal vertical ($Y$):} a f.e.m. do secundário passa por um
\textbf{integrador} $RC$. Como
$\varepsilon_2=N_2A\,\mathrm{d}B/\mathrm{d}t$, sua integral é proporcional a
$B$.
\end{itemize}
Em modo $XY$, o osciloscópio traça diretamente o laço.\\
(b) \textbf{Constantes de conversão.}
\[
H=\frac{N_1}{\ell R_s}\,v_X;
\qquad
B=\frac{R_iC_i}{N_2A}\,v_Y,
\]
com $R_iC_i$ os componentes do integrador.\\
\textbf{Condição do integrador:} $R_iC_i\gg1/f$, para que ele integre de fato em
vez de apenas atenuar. Com $f=\SI{60}{\hertz}$, adote
$R_iC_i\ge\SI{0.1}{\second}$.\\
(c) \textbf{Cuidados experimentais.}
\begin{itemize}
\item \textbf{Desmagnetizar antes} --- caso contrário os primeiros ciclos
refletem a história anterior. Comece com amplitude alta e reduza.
\item \textbf{$R_s$ pequeno} (fração de ohm), para não distorcer a corrente nem
dissipar demais. Use quatro fios se possível.
\item \textbf{Núcleo toroidal}, que garante $H$ uniforme e caminho médio bem
definido --- em núcleos em E o caminho é ambíguo.
\item \textbf{Deriva do integrador:} qualquer \emph{offset} faz o laço "subir"
lentamente na tela. Use acoplamento CA ou integrador com realimentação CC.
\item \textbf{Varrer a amplitude} para obter a família de laços e localizar a
curva de magnetização normal, que é o lugar geométrico dos vértices.
\end{itemize}
\textbf{O que se extrai:} $B_r$ (interseção com o eixo vertical), $H_c$
(interseção com o horizontal), $B_{sat}$ e a área --- e daí as perdas por
ciclo.}
\stepcounter{questao}
\resposta{Domínios minimizam a energia magnetostática; a mobilidade das paredes
determina $\mu_r$ e $H_c$}
\resolucao{(a) \textbf{Por que existem.} Um cristal ferromagnético
uniformemente magnetizado criaria campo externo intenso, com alta energia
magnetostática armazenada no espaço.\\
Dividir-se em \textbf{domínios} de orientações opostas faz esse campo externo
quase desaparecer --- as linhas se fecham internamente. O custo é a energia das
\emph{paredes de domínio} (paredes de Bloch), onde os momentos giram
gradualmente.\\
\textbf{O equilíbrio} entre esses dois termos fixa o tamanho típico dos
domínios, tipicamente de micrômetros a dezenas de micrômetros.\\
(b) \textbf{Relação com as propriedades macroscópicas.}
\begin{itemize}
\item \textbf{$\mu_r$ alta} exige paredes \emph{móveis} --- material puro, sem
tensões internas, com grãos grandes e bem orientados.
\item \textbf{$H_c$ baixa} tem a mesma origem: paredes que deslizam facilmente
não resistem à inversão.
\item \textbf{$B_r$ alta} exige que os domínios permaneçam alinhados após a
retirada do campo --- o que requer \emph{anisotropia} cristalina forte,
"ancorando" a magnetização em direções preferenciais.
\end{itemize}
(c) \textbf{Moles contra duros, no nível microscópico.}
\begin{center}
\begin{tabular}{lll}
\hline
& Moles & Duros\\
\hline
Paredes & muito móveis & fixadas ou inexistentes\\
Impurezas & minimizadas & introduzidas de propósito\\
Anisotropia & baixa & alta\\
Grãos & grandes, orientados & finos, monodomínio\\
\hline
\end{tabular}
\end{center}
\textbf{O contraste é elegante:} ímãs modernos de neodímio são feitos de
partículas tão pequenas que cada uma é um \textbf{monodomínio} --- não há
paredes para mover, e inverter a magnetização exige girar todos os momentos
contra a anisotropia. Daí a coercividade enorme.\\
\textbf{Já o aço-silício de grão orientado} persegue o oposto: grãos grandes,
alinhados por laminação a frio, com paredes livres para deslizar.\\
\textbf{Mesma física, objetivos opostos --- e processos de fabricação
inteiramente diferentes.}}
\stepcounter{questao}
\resposta{$\SI{1.33}{\kilogram}$ a $\SI{1.5}{\tesla}$;
$\SI{2.22}{\kilogram}$ a $\SI{1.2}{\tesla}$}
\resolucao{(a)
\[
m\le\frac{2}{1{,}5}=\SI{1.33}{\kilogram}\ \text{a } \SI{1.5}{\tesla};
\qquad
m\le\frac{2}{0{,}9}=\SI{2.22}{\kilogram}\ \text{a } \SI{1.2}{\tesla}.
\]
(b) De $V=4{,}44NfB_{\max}A$, com tensão e frequência fixas,
\[
N\,B_{\max}A=\text{const}.
\]
Reduzir $B_{\max}$ de $1{,}5$ para $\SI{1.2}{\tesla}$ exige aumentar
$NA$ em $25\%$ --- mais espiras, ou seção maior, ou ambos.\\
(c) \textbf{Critério de escolha.} As duas opções levam a projetos diferentes:
\begin{center}
\begin{tabular}{lcc}
\hline
& $B=\SI{1.5}{\tesla}$ & $B=\SI{1.2}{\tesla}$\\
\hline
Massa de núcleo & menor & maior ($+67\%$)\\
Espiras & menos & mais ($+25\%$)\\
Perdas no cobre & menores & maiores\\
Margem para saturação & pequena & confortável\\
Corrente de \emph{inrush} & severa & moderada\\
\hline
\end{tabular}
\end{center}
\textbf{Note o conflito:} operar em $B$ alto economiza \emph{ferro} mas gasta
mais \emph{cobre} por unidade de perda; operar em $B$ baixo faz o inverso.\\
\textbf{Decisão:}
\begin{itemize}
\item \textbf{Se o custo de material dominar} (transformadores de baixo custo,
uso intermitente), escolha $B$ alto.
\item \textbf{Se o custo de energia dominar} (operação contínua por décadas),
escolha $B$ baixo --- as perdas em vazio existem $\SI{8760}{\hour}$ por ano.
\item \textbf{Se houver risco de sobretensão}, escolha $B$ baixo por segurança:
um transformador operando a $\SI{1.5}{\tesla}$ satura com apenas $10\%$ de
sobretensão.
\end{itemize}}
\stepcounter{questao}
\resposta{Grãos alinhados numa direção preferencial; excelente ao longo dela,
ruim transversalmente}
\resolucao{(a) \textbf{Processo e efeito.} O aço-silício é laminado a frio e
recozido de modo que os grãos cristalinos se alinhem com o eixo de fácil
magnetização do ferro (a direção $\langle100\rangle$ da célula cúbica) paralelo
à direção de laminação.\\
\textbf{Resultado:} nessa direção, magnetizar é muito mais fácil --- as paredes
de domínio se movem livremente, $\mu_r$ é altíssima e as perdas caem
drasticamente.\\
(b) \textbf{Comparação por direção.}
\begin{center}
\begin{tabular}{lcc}
\hline
& Direção de laminação & Transversal\\
\hline
Perdas específicas & $\approx\SI{1}{\watt\per\kilogram}$ &
 $2$ a $3\times$ maiores\\
$\mu_r$ & muito alta & bem menor\\
$B$ para dado $H$ & alto & reduzido\\
\hline
\end{tabular}
\end{center}
\textbf{A anisotropia é grande} --- o material é excelente em uma direção e
medíocre nas outras.\\
(c) \textbf{Por que serve a transformadores e não a motores.}
\begin{itemize}
\item \textbf{Transformador:} o fluxo percorre um caminho \emph{fixo e
conhecido} --- as colunas e as culatras do núcleo. É possível cortar as lâminas
de modo que o fluxo siga sempre a direção favorável. \emph{(Nos cantos, onde o
fluxo precisa dobrar, usam-se juntas em $\ang{45}$ para minimizar a
transição.)}
\item \textbf{Motor:} o fluxo no estator \textbf{gira}, percorrendo todas as
direções ao longo de cada revolução. Não existe direção preferencial a
privilegiar --- metade do tempo o fluxo estaria na direção ruim.
\end{itemize}
\textbf{Por isso motores usam aço de grão \emph{não} orientado}, deliberadamente
\textbf{isotrópico}: pior que o grão orientado em sua melhor direção, mas
uniforme em todas --- e portanto melhor \emph{na média} sobre uma rotação
completa.\\
\textbf{Princípio geral:} otimizar para o caso médio quando a solicitação é
variável, e para o caso específico quando ela é fixa.}
\stepcounter{questao}
\resposta{Fator $\approx\mu_r t/D$; o mecanismo é desvio de fluxo, não
cancelamento}
\resolucao{(a) \textbf{Mecanismo.} Ao contrário da blindagem elétrica --- que
\emph{cancela} o campo por redistribuição de cargas --- a blindagem magnética
\textbf{desvia} o fluxo.\\
O invólucro de alta $\mu_r$ oferece caminho de relutância muito menor que o ar
interno. O fluxo "prefere" circular pela parede e contorna a região protegida.\\
\textbf{Analogia:} é um divisor de fluxo, em que o material é o ramo de baixa
relutância e a cavidade, o de alta.\\
(b) \textbf{Fator de blindagem.} Para um cilindro de diâmetro $D$ e parede
$t$, com $t\ll D$, a estimativa clássica é
\[
S\approx\frac{\mu_r\,t}{D}.
\]
Com $\mu_r=20000$ (mu-metal recozido), $t=\SI{1}{\milli\meter}$ e
$D=\SI{100}{\milli\meter}$:
\[
S\approx\frac{20000(1)}{100}=200,
\]
ou seja, atenuação de cerca de $\SI{46}{\decibel}$.\\
\textbf{Compare com a blindagem elétrica}, que atinge muitas ordens de grandeza
com uma folha fina de qualquer metal. A magnética é \emph{muito} mais modesta.\\
(c) \textbf{Limitações práticas.}
\begin{itemize}
\item \textbf{Saturação.} Se o campo externo for intenso, a parede satura,
$\mu_r$ desaba e a blindagem falha justamente quando mais se precisa dela.
\emph{Solução:} camadas múltiplas, com a externa de material de alto
$B_{sat}$ (aço comum) e a interna de mu-metal.
\item \textbf{Sensibilidade mecânica.} O mu-metal perde permeabilidade se for
dobrado, cortado ou submetido a choque. Exige \textbf{recozimento após a
conformação}, em atmosfera de hidrogênio --- processo caro.
\item \textbf{Aberturas.} Furos e fendas quebram o caminho de baixa relutância e
degradam a blindagem muito mais que no caso elétrico.
\item \textbf{Custo e peso}, ambos elevados.
\end{itemize}
\textbf{Conclusão:} blindar campo magnético estático é caro, frágil e de eficácia
limitada. Sempre que possível, é preferível \emph{afastar} a fonte, reduzir a
área das espiras ou compensar ativamente o campo com bobinas.}
\stepcounter{questao}
\resposta{A reta de carga tem inclinação
$-\dfrac{\ell_m A_g}{g A_m}$; o aproveitamento é máximo em
$(BH)_{\max}$}
\resolucao{(a) \textbf{Reta de carga.} Duas condições:
\emph{conservação do fluxo} ($B_mA_m=B_gA_g$) e \emph{lei de Ampère}
($H_m\ell_m+H_gg=0$, pois não há corrente).\\
Da segunda, $H_m\ell_m=-H_gg=-\dfrac{B_g}{\mu_0}g$. Combinando com a primeira:
\[
B_m=-\mu_0\,\frac{\ell_m A_g}{g\,A_m}\,H_m.
\]
É uma reta pela \textbf{origem} no segundo quadrante ($H_m<0$, $B_m>0$) ---
a chamada \emph{linha de carga} do ímã.\\
\textbf{Note:} o ímã opera com $H$ \emph{negativo} --- ele está sujeito ao
próprio campo desmagnetizante. Sem entreferro ($g\to0$), a reta seria vertical e
o ponto de operação estaria em $B_r$.\\
(b) \textbf{Máximo aproveitamento.} A energia disponível no entreferro é
proporcional ao produto $B_mH_m$. O ponto ótimo é aquele em que
$|B_mH_m|$ é \textbf{máximo} --- o já mencionado $(BH)_{\max}$.\\
Dimensiona-se o circuito \emph{para que a reta de carga passe por esse ponto},
escolhendo a razão $\dfrac{\ell_m A_g}{g A_m}$ adequada. Isso minimiza o volume
de ímã necessário.\\
(c) \textbf{Efeito de aumentar o entreferro.} A inclinação
$-\mu_0\ell_mA_g/(gA_m)$ diminui em módulo --- a reta \textbf{deita}. O ponto de
operação desce na curva de desmagnetização:
\begin{itemize}
\item $B_m$ diminui (menos fluxo);
\item $|H_m|$ aumenta (mais campo desmagnetizante sobre o próprio ímã);
\item se o entreferro for grande demais, o ponto pode ultrapassar o
\textbf{joelho} da curva de desmagnetização --- e aí ocorre perda
\textbf{irreversível} de magnetização.
\end{itemize}
\textbf{Consequência prática decisiva:} desmontar um circuito magnético com ímã
de alnico pode danificá-lo permanentemente, pois o ponto de operação despenca.
Ímãs de neodímio, cuja curva de desmagnetização é praticamente reta até o
segundo quadrante inteiro, toleram isso muito melhor --- e é uma das razões de
sua popularidade, além do produto de energia.}

\gabarito[12.1 Conceitos iniciais e materiais magnéticos]

\subsection{Campo magnético}
\stepcounter{questao}
\resposta{(b)}
\resolucao{Para um fio reto longo, $B = \dfrac{\mu_0 I}{2\pi r}$: o campo é
\emph{inversamente} proporcional à distância $r$ (não ao quadrado, como no campo
elétrico de uma carga puntiforme).}
\stepcounter{questao}
\resposta{$B = \pot{4}{-5}\,\si{\tesla} = \SI{40}{\micro\tesla}$}
\resolucao{$B = \dfrac{\mu_0 I}{2\pi r}
= \dfrac{(\pot{4\pi}{-7})(10)}{2\pi(0{,}05)}
= \dfrac{\pot{2}{-6}}{0{,}05} = \pot{4}{-5}\,\si{\tesla}$.\\
Um atalho útil: $\dfrac{\mu_0}{2\pi} = \pot{2}{-7}$, então
$B = \pot{2}{-7}\cdot\dfrac{I}{r}$.}
\stepcounter{questao}
\resposta{(c)}
\resolucao{O campo de um fio reto forma \emph{círculos concêntricos} em planos
perpendiculares ao fio. Pela regra da mão direita (polegar no sentido da corrente,
dedos envolvendo o fio), determina-se o sentido dessas linhas. O campo nunca é
radial (isso é característico do campo \emph{elétrico} de uma carga).}
\stepcounter{questao}
\resposta{$B \approx \SI{2.51}{\milli\tesla}$}
\resolucao{$B = \mu_0\,n\,I = (\pot{4\pi}{-7})(1000)(2)
= \pot{8\pi}{-4} \approx\SI{2.51}{\milli\tesla}$. O campo no interior de um
solenoide é \emph{uniforme} e independe do raio.}
\stepcounter{questao}
\resposta{$B \approx \SI{31.4}{\micro\tesla}$}
\resolucao{$B = \dfrac{\mu_0 I}{2R}
= \dfrac{(\pot{4\pi}{-7})(5)}{2(0{,}1)}
= \dfrac{\pot{2\pi}{-6}}{0{,}2}=\pot{\pi}{-5}
\approx\SI{31.4}{\micro\tesla}$.}
\stepcounter{questao}
\resposta{$B = \SI{40}{\micro\tesla}$}
\resolucao{$B = \dfrac{\mu_0 I}{2\pi r} = \pot{2}{-7}\cdot\dfrac{4}{0{,}02}
= \pot{4}{-5}\,\si{\tesla} = \SI{40}{\micro\tesla}$
(usando o atalho $\mu_0/2\pi = \pot{2}{-7}$).}
\stepcounter{questao}
\resposta{$B \approx \SI{12}{\micro\tesla}$; os fios se atraem}
\resolucao{No ponto médio ($r = \SI{5}{\centi\meter}$ de cada fio), os campos dos
dois fios têm \emph{sentidos opostos} (correntes no mesmo sentido), então os
módulos se subtraem:
\[
B = \frac{\mu_0}{2\pi r}(I_2 - I_1)
= \pot{2}{-7}\cdot\frac{6-3}{0{,}05}
= \pot{2}{-7}\cdot60 = \pot{1.2}{-5}\,\si{\tesla}\approx\SI{12}{\micro\tesla}.
\]
Quanto à força \emph{entre} os fios: correntes de mesmo sentido se
\textbf{atraem} (regra fundamental do eletromagnetismo, base da definição do
ampère). Note a distinção: o \emph{campo} no ponto médio se subtrai, mas a
\emph{força} entre os fios é atrativa --- são perguntas diferentes.}
\stepcounter{questao}
\resposta{$F/\ell = \SI{4e-4}{\newton\per\meter}$, repulsão}
\resolucao{A força por comprimento entre dois fios paralelos é
\[
\frac{F}{\ell}=\frac{\mu_0 I_1 I_2}{2\pi d}
=\pot{2}{-7}\cdot\frac{10\cdot10}{0{,}05}
=\pot{4}{-4}\,\si{\newton\per\meter}.
\]
Correntes em sentidos \emph{opostos} se \textbf{repelem} (e as de mesmo sentido se
atraem). Esta força é a base da definição do ampère no SI.}
\stepcounter{questao}
\resposta{(a) $B = \dfrac{\mu_0 N I}{2\pi r} = \SI{1.2}{\milli\tesla}$;
(b) e (c) ver comentário}
\resolucao{\textbf{(a)} A lei de Ampère estabelece
$\oint \vec B\cdot\mathrm{d}\vec\ell = \mu_0 I_{\text{env}}$. Escolhendo uma
circunferência de raio $r$ concêntrica ao toroide, por simetria $B$ é constante e
tangente ao longo dela:
\[
B\,(2\pi r) = \mu_0\,N I
\;\Longrightarrow\;
B=\frac{\mu_0 N I}{2\pi r}
=\frac{(\pot{4\pi}{-7})(200)(3)}{2\pi(0{,}1)}
=\pot{1.2}{-3}\,\si{\tesla}=\SI{1.2}{\milli\tesla}.
\]
\textbf{(b) Campo externo nulo.} Para um percurso \emph{fora} do toroide há dois
casos: (i) circunferência de raio maior que o externo --- a corrente entra e sai
do enrolamento, de modo que a corrente líquida envolvida é \emph{zero};
(ii) circunferência de raio menor que o interno --- não envolve corrente alguma.
Em ambos, $\oint \vec B\cdot\mathrm{d}\vec\ell = 0$ e, pela simetria,
$B = 0$.\\
\textbf{Consequência prática:} o toroide \textbf{confina} o campo em seu interior,
praticamente sem dispersão. Por isso transformadores toroidais produzem muito
menos interferência eletromagnética que os de núcleo E-I, e são preferidos em
áudio e instrumentação de precisão.

\textbf{(c) Comparação com o solenoide.} O perímetro médio é
$\ell = 2\pi r = \SI{0.628}{\meter}$. Um solenoide reto com o mesmo $N$ e esse
comprimento teria
\[
B = \mu_0\,\frac{N}{\ell}\,I = \frac{\mu_0 N I}{2\pi r},
\]
\emph{exatamente a mesma expressão}. Ou seja: o toroide é um solenoide
"fechado sobre si mesmo". A diferença essencial é que o solenoide reto tem
\textbf{extremidades}, onde as linhas de campo escapam (efeito de borda) e o
campo é menos uniforme; o toroide não tem bordas, sendo o dispositivo ideal para
estudar o comportamento magnético puro de um material --- razão pela qual as
curvas de histerese são levantadas em amostras toroidais.}
\stepcounter{questao}
\resposta{$I_2=\SI{0.5}{\ampere}$, em sentido oposto a $I_1$}
\resolucao{No centro de uma espira, $B=\mu_0 I/(2R)$. Para cancelamento,
$I_1/R_1=I_2/R_2$. Assim,
$I_2=I_1(R_2/R_1)=2(2{,}5/10)=\SI{0.5}{\ampere}$. Os sentidos devem ser opostos.}
\stepcounter{questao}
\resposta{$B=\SI{40}{\micro\tesla}$}
\resolucao{
\[
B=\frac{\mu_0 I}{2\pi r}=\frac{(4\pi\times10^{-7})(2)}{2\pi(0{,}01)}
=\frac{\pot{2}{-7}\times2}{0{,}01}=\pot{4}{-5}\,\si{\tesla}.
\]
\textbf{Atalho útil:} $\dfrac{\mu_0}{2\pi}=\pot{2}{-7}$ exatamente, de modo que
$B=\pot{2}{-7}\dfrac{I}{r}$ --- evita carregar $\pi$ pela conta toda.\\
\textbf{Ordem de grandeza:} $\SI{40}{\micro\tesla}$ é comparável ao campo
magnético da Terra ($\approx\SI{50}{\micro\tesla}$). Uma bússola perto desse fio
já seria perturbada.}
\stepcounter{questao}
\resposta{$B=\SI{7.54}{\milli\tesla}$}
\resolucao{
\[
B=\mu_0 nI=(4\pi\times10^{-7})(2000)(3)=\pot{7.54}{-3}\,\si{\tesla}.
\]
\textbf{Note o que \emph{não} aparece na fórmula:} o raio do solenoide e o
número total de espiras. Só importa a \textbf{densidade linear}
$n=N/L$.\\
Dois solenoides, um fino e um grosso, com a mesma densidade de espiras e a mesma
corrente, produzem o \emph{mesmo} campo interno.}
\stepcounter{questao}
\resposta{$I=\SI{10}{\ampere}$}
\resolucao{
\[
I=\frac{2\pi rB}{\mu_0}=\frac{rB}{\pot{2}{-7}}
=\frac{(0{,}02)(\pot{1}{-4})}{\pot{2}{-7}}=\SI{10}{\ampere}.
\]
\textbf{Leitura prática:} são necessários $\SI{10}{\ampere}$ para produzir, a
dois centímetros, apenas o dobro do campo terrestre. Campos magnéticos
intensos exigem correntes altas --- ou muitas espiras.}
\stepcounter{questao}
\resposta{(a)}
\resolucao{Aplicando a lei de Ampère a um retângulo com um lado dentro e outro
fora, obtém-se $B_{int}=\mu_0nI$ e $B_{ext}\approx0$ para um solenoide
\emph{ideal} (infinitamente longo).\\
\textbf{Por que o campo externo quase se anula:} as linhas que saem de uma
extremidade retornam pela outra, espalhando-se por um volume enorme --- a
densidade de linhas fora é desprezível diante da interna.\\
\textbf{Em solenoides reais}, o campo externo não é exatamente nulo, e nas
extremidades o campo interno cai para cerca de \emph{metade} do valor central.}
\stepcounter{questao}
\resposta{$\SI{0.5}{\gauss}$}
\resolucao{
\[
B=\pot{5}{-5}\,\si{\tesla}\times10^{4}=\SI{0.5}{\gauss}.
\]
\textbf{Escala de referência para o tópico:}
\begin{center}
\begin{tabular}{lcc}
\hline
Situação & Tesla & Gauss\\
\hline
Campo terrestre & $\pot{5}{-5}$ & $0{,}5$\\
Ímã de geladeira & $\pot{5}{-3}$ & $50$\\
Ímã de neodímio (superfície) & $0{,}5$ & $5000$\\
Ressonância magnética & $1{,}5$ a $3$ & $\pot{1.5}{4}$ a $\pot{3}{4}$\\
\hline
\end{tabular}
\end{center}
O tesla é uma unidade \emph{grande}: campos de laboratório raramente passam de
alguns tesla, e o gauss continua em uso justamente por ser mais próximo dos
valores cotidianos.}
\stepcounter{questao}
\resposta{$n=\SI{1000}{\per\meter}$; $B=\SI{1.26}{\milli\tesla}$}
\resolucao{(a) $n=\dfrac{100}{0{,}10}=\SI{1000}{\per\meter}$.\\
(b) $B=\mu_0nI=(4\pi\times10^{-7})(1000)(1)
=\pot{1.257}{-3}\,\si{\tesla}$.\\
\textbf{Comparação instrutiva:} isso é apenas $25$ vezes o campo terrestre ---
um solenoide modesto produz campo modesto. Para chegar a
$\SI{0.1}{\tesla}$ seriam necessários $\SI{80}{\ampere}$ com a mesma
geometria, o que exigiria refrigeração.\\
\textbf{Caminho alternativo:} inserir núcleo ferromagnético, que multiplica
$B$ por $\mu_r$ --- centenas ou milhares de vezes.}
\stepcounter{questao}
\resposta{$B_1/B_2=3$}
\resolucao{Como $B\propto1/r$:
\[
\frac{B_1}{B_2}=\frac{r_2}{r_1}=\frac{12}{4}=3.
\]
\textbf{Contraste importante com o campo elétrico:} o campo de uma carga
puntiforme cai com $1/r^{2}$, mas o de um fio \emph{infinito} cai com
$1/r$.\\
A razão é a mesma da eletrostática: o fio é uma fonte \emph{estendida} em uma
dimensão, e a superfície amperiana (um cilindro) tem área proporcional a $r$,
não a $r^{2}$.}
\stepcounter{questao}
\resposta{$I\approx\SI{7.96}{\ampere}$}
\resolucao{Do campo no centro de uma espira, $B=\dfrac{\mu_0I}{2R}$:
\[
I=\frac{2RB}{\mu_0}=\frac{2(0{,}05)(\pot{1}{-4})}{4\pi\times10^{-7}}
=\SI{7.96}{\ampere}.
\]
\textbf{Comparação com o fio reto:} um fio retilíneo a
$\SI{5}{\centi\meter}$ precisaria de $\SI{25}{\ampere}$ para o mesmo campo. A
espira é mais eficiente porque \emph{todos} os seus elementos contribuem no
mesmo sentido no centro --- a razão entre as duas fórmulas é exatamente
$\pi$.}
\stepcounter{questao}
\resposta{O campo não se altera}
\resolucao{Como $B=\mu_0\dfrac{N}{L}I$, dobrar $N$ e $L$ mantém a razão
$N/L$ inalterada:
\[
B'=\mu_0\frac{2N}{2L}I=B.
\]
\textbf{Interpretação:} o solenoide dobrado equivale a dois solenoides
originais colocados em sequência. Cada trecho "enxerga" a mesma vizinhança de
espiras, e o campo local não muda.\\
\textbf{O que muda:} a região de campo uniforme fica mais longa, e a resistência
do enrolamento dobra --- exigindo o dobro da tensão para a mesma corrente.}
\stepcounter{questao}
\resposta{$B=0$}
\resolucao{Cada fio produz, no ponto médio,
\[
B=\pot{2}{-7}\frac{4}{0{,}04}=\pot{2}{-5}\,\si{\tesla}.
\]
Pela regra da mão direita, com as correntes no \emph{mesmo} sentido, os dois
campos têm sentidos \textbf{opostos} no ponto médio --- e módulos iguais.
Cancelam-se:
\[
B_{total}=0.
\]
\textbf{Cuidado com a intuição:} correntes de mesmo sentido \emph{atraem-se},
mas seus campos se \emph{cancelam} entre elas. Força e campo são coisas
diferentes.\\
\textbf{Fora do par}, ao contrário, os campos se somam --- e é por isso que dois
fios juntos conduzindo no mesmo sentido produzem, de longe, o campo de um único
fio de corrente $2I$.}
\stepcounter{questao}
\resposta{O campo fica $200$ vezes maior}
\resolucao{Em meio linear,
\[
B=\frac{\mu_r\mu_0I}{2\pi r},
\]
de modo que $B$ é multiplicado por $\mu_r=200$.\\
\textbf{Mecanismo:} o material se \emph{magnetiza}, e seus dipolos alinhados
somam sua própria contribuição ao campo aplicado. Diferentemente do dielétrico
--- que \emph{reduz} o campo elétrico --- o material ferromagnético o
\textbf{amplifica}.\\
\textbf{Ressalva séria:} $\mu_r$ não é constante. Ele cai drasticamente quando o
material satura, e a relação $B\times H$ deixa de ser linear. A fórmula acima
vale apenas em campos baixos.}
\stepcounter{questao}
\resposta{$r=\SI{20}{\centi\meter}$}
\resolucao{
\[
r=\frac{\pot{2}{-7}I}{B}=\frac{(\pot{2}{-7})(50)}{\pot{5}{-5}}
=\SI{0.20}{\meter}.
\]
\textbf{Consequência prática:} a menos de $\SI{20}{\centi\meter}$ de um
condutor de $\SI{50}{\ampere}$, uma bússola aponta mais para o fio do que para
o norte.\\
É por isso que instrumentos magnéticos sensíveis --- e monitores de tubo, na
época em que existiam --- precisam ficar afastados de cabos de potência.}
\stepcounter{questao}
\resposta{Campo interno praticamente igual; o toroide tem campo externo
\emph{rigorosamente} nulo}
\resolucao{\textbf{Campo interno:} ambos valem $\approx\mu_0nI$. No toroide,
$n=N/(2\pi r)$ varia ligeiramente com o raio, de modo que o campo não é
perfeitamente uniforme na seção --- mas para toroides finos a diferença é
pequena.\\
\textbf{Campo externo:} aqui está a diferença essencial.
\begin{itemize}
\item \textbf{Solenoide:} as linhas precisam retornar por fora, e o campo
externo, embora fraco, \emph{existe}.
\item \textbf{Toroide:} as linhas fecham-se \emph{inteiramente} dentro do
núcleo. Aplicando a lei de Ampère a um círculo externo, a corrente líquida
envolvida é \textbf{zero} --- e o campo também.
\end{itemize}
\textbf{Consequência de projeto:} indutores toroidais praticamente não irradiam
nem captam interferência, e podem ser montados lado a lado sem acoplamento
mútuo. É por isso que dominam aplicações de filtragem em fontes chaveadas.}
\stepcounter{questao}
\resposta{$B=\SI{56.6}{\micro\tesla}$}
\resolucao{Cada fio produz campo perpendicular ao plano definido por ele e pelo
ponto. Como os fios são perpendiculares entre si e coplanares, ambos os campos
saem (ou entram) \emph{perpendicularmente à folha} --- portanto são
\textbf{paralelos} e se somam algebricamente:
\[
B_1=\pot{2}{-7}\frac{6}{0{,}03}=\SI{40}{\micro\tesla};
\qquad
B_2=\pot{2}{-7}\frac{8}{0{,}04}=\SI{40}{\micro\tesla}.
\]
Se os sentidos coincidirem: $B=\SI{80}{\micro\tesla}$; se forem opostos,
$B=0$.\\
\textbf{Atenção ao caso geral:} se os fios \emph{não} fossem coplanares, os
campos teriam direções diferentes e a soma seria vetorial --- com
$\SI{40}{\micro\tesla}$ perpendiculares, o resultado seria
$40\sqrt2=\SI{56.6}{\micro\tesla}$. \textbf{A geometria decide se a soma é
algébrica ou vetorial}, e é ela que precisa ser estabelecida antes de qualquer
conta.}
\stepcounter{questao}
\resposta{$5$, $10$ e $\SI{4}{\milli\tesla}$; máximo na superfície}
\resolucao{\textbf{Dentro ($r<a$):} a superfície amperiana envolve apenas a
fração $(r/a)^{2}$ da corrente:
\[
B=\frac{\mu_0 I r}{2\pi a^{2}}
=\pot{2}{-7}\frac{100\,r}{(0{,}002)^{2}}.
\]
Em $r=\SI{1}{\milli\meter}$: $B=\SI{5}{\milli\tesla}$.\\
\textbf{Fora ($r>a$):} toda a corrente é envolvida, e vale a fórmula do fio
fino:
\[
B=\pot{2}{-7}\frac{100}{r}.
\]
Em $r=\SI{5}{\milli\meter}$: $B=\SI{4}{\milli\tesla}$.\\
\textbf{Na superfície ($r=a$):} as duas expressões coincidem em
$B=\SI{10}{\milli\tesla}$ --- o campo é contínuo.\\
(b) \textbf{O máximo está na superfície.} Dentro, $B$ cresce linearmente com
$r$; fora, decai com $1/r$. O pico ocorre exatamente em $r=a$.\\
\textbf{Analogia com a esfera isolante carregada:} lá o campo elétrico também
cresce linearmente dentro e cai com $1/r^{2}$ fora, com máximo na superfície. A
estrutura matemática é a mesma --- muda apenas a lei de Gauss ou de Ampère
aplicada.}
\stepcounter{questao}
\resposta{$B=\SI{0.5}{\milli\tesla}$ entre os condutores; $B=0$ fora}
\resolucao{(a) Entre os condutores, a superfície amperiana envolve apenas a
corrente central:
\[
B=\pot{2}{-7}\frac{5}{0{,}002}=\pot{5}{-4}\,\si{\tesla}=\SI{0.5}{\milli\tesla}.
\]
(b) \textbf{Fora do cabo}, a superfície amperiana envolve $+I$ e $-I$ --- a
corrente líquida é \textbf{zero}:
\[
\oint\vec B\cdot\mathrm{d}\vec\ell=\mu_0(0)=0
\;\Longrightarrow\;
B=0.
\]
(c) \textbf{Consequência prática: o cabo coaxial não irradia nem capta campo
magnético.} Todo o campo fica confinado entre os condutores.\\
Isso o torna a escolha padrão para transmitir sinais em ambientes com
interferência --- e explica por que um cabo paralelo comum, cujos condutores
estão afastados, é muito pior: ali as correntes opostas não se cancelam
perfeitamente e o par forma uma espira que irradia.\\
\textbf{Mesma ideia, outra forma:} o par trançado obtém efeito semelhante ao
inverter a orientação da espira a cada torção, de modo que as contribuições
sucessivas se anulem.}
\stepcounter{questao}
\resposta{$31{,}4$; $11{,}1$ e $\SI{2.81}{\micro\tesla}$}
\resolucao{(a) $B_0=\dfrac{\mu_0I}{2R}=\dfrac{(4\pi\times10^{-7})(5)}{0{,}20}
=\pot{3.14}{-5}\,\si{\tesla}$.\\
(b) No eixo,
\[
B=\frac{\mu_0IR^{2}}{2\left(R^{2}+x^{2}\right)^{3/2}}.
\]
Em $x=R=\SI{0.10}{\meter}$:
$B=\pot{1.11}{-5}\,\si{\tesla}$, ou seja
$B_0/2^{3/2}=0{,}354\,B_0$.\\
Em $x=2R=\SI{0.20}{\meter}$: $B=\pot{2.81}{-6}\,\si{\tesla}$, cerca de
$0{,}089\,B_0$.\\
(c) \textbf{A queda é rápida.} A apenas um raio de distância, o campo já caiu
para $35\%$; a dois raios, para $9\%$.\\
Para $x\gg R$, a expressão tende a
$B\to\dfrac{\mu_0IR^{2}}{2x^{3}}$ --- decaimento com $1/x^{3}$, típico de
\textbf{dipolo}. É o mesmo expoente do dipolo elétrico, e pela mesma razão: a
espira não tem "monopolo magnético" para produzir termo $1/x^{2}$.}
\stepcounter{questao}
\resposta{Na extremidade, $B\approx\mu_0nI/2$ --- metade do valor central}
\resolucao{(a) \textbf{Condição:} $L\gg D$, isto é, comprimento muito maior que
o diâmetro. Na prática, $L/D>10$ garante erro inferior a $1\%$ no centro.\\
(b) \textbf{Na extremidade}, apenas "metade" do solenoide contribui --- as
espiras estão todas de um lado. O resultado exato para solenoide longo é
\[
B_{extremidade}=\frac{\mu_0nI}{2},
\]
exatamente \textbf{metade} do valor central.\\
(c) \textbf{Consequências de projeto.}
\begin{itemize}
\item A região de campo verdadeiramente uniforme é \emph{menor} que o
solenoide: ela se restringe à parte central, tipicamente à metade do
comprimento.
\item Para uniformidade melhor que $1\%$ em uma região de comprimento
$\ell$, o solenoide precisa ter comprimento bem maior que $\ell$ --- ou usar
enrolamento com densidade \emph{corrigida} nas extremidades.
\item A alternativa clássica é o par de \textbf{bobinas de Helmholtz}: duas
bobinas iguais separadas por uma distância igual ao raio, arranjo que anula a
segunda derivada do campo no centro e produz região uniforme com pouco
material.
\end{itemize}}
\stepcounter{questao}
\resposta{$B=0$ no fio central}
\resolucao{O fio central não produz campo sobre si mesmo. Restam as
contribuições dos dois vizinhos, ambos a $\SI{5}{\centi\meter}$:
\[
B_1=B_2=\pot{2}{-7}\frac{10}{0{,}05}=\SI{40}{\micro\tesla}.
\]
Pela regra da mão direita, com os dois vizinhos conduzindo no mesmo sentido e
estando em \emph{lados opostos} do fio central, seus campos ali têm sentidos
\textbf{contrários}:
\[
B=40-40=0.
\]
\textbf{Mas a força não é nula --- ela também é.} Como $\vec F=I\vec\ell\times
\vec B$ e $\vec B=\vec 0$ no fio central, a força sobre ele é zero: as atrações
dos dois vizinhos se equilibram.\\
\textbf{Nos fios externos, porém}, a situação é diferente: cada um sofre atração
resultante para o centro. O conjunto tende a se \emph{comprimir} --- efeito real
em barramentos de alta corrente, que precisam de espaçadores mecânicos.}
\stepcounter{questao}
\resposta{$NI=\SI{15.9}{\kilo\ampere}$-espira; exigiria $\SI{15.9}{\ampere}$
com $1000$ espiras}
\resolucao{(a) De $B=\mu_0NI/L$:
\[
NI=\frac{BL}{\mu_0}=\frac{(0{,}1)(0{,}20)}{4\pi\times10^{-7}}
=\pot{1.59}{4}\ \text{A-espira}.
\]
(b) Com $N=1000$: $I=\SI{15.9}{\ampere}$.\\
\textbf{Viável, mas problemático.} Mil espiras conduzindo
$\SI{16}{\ampere}$ exigem fio grosso (a densidade de corrente admissível limita
a seção), o que aumenta o volume do enrolamento. E a dissipação é considerável:
com resistência de apenas $\SI{1}{\ohm}$, seriam
$\SI{256}{\watt}$ --- exigindo refrigeração.\\
(c) \textbf{Alternativa: núcleo ferromagnético.} Com $\mu_r=1000$, a mesma
corrente produziria $\SI{100}{\tesla}$ --- valor impossível, porque o material
\textbf{satura} em torno de $1{,}5$ a $\SI{2}{\tesla}$.\\
\textbf{O que o núcleo realmente entrega:} os $\SI{0.1}{\tesla}$ pedidos com
corrente cerca de mil vezes menor ($\SI{16}{\milli\ampere}$), enquanto o
material permanecer na região linear.\\
\textbf{Limite físico:} acima da saturação, o núcleo deixa de ajudar e o
eletroímã volta a se comportar como se fosse de ar. É por isso que campos acima
de $\SI{2}{\tesla}$ exigem bobinas supercondutoras, sem núcleo.}
\stepcounter{questao}
\resposta{$\SI{400}{\micro\tesla}$ no meio; $\approx\SI{0.04}{\micro\tesla}$ a
$\SI{10}{\centi\meter}$}
\resolucao{(a) No ponto médio, cada fio dista $\SI{1}{\milli\meter}$, e com
correntes \emph{opostas} os campos se \textbf{somam}:
\[
B=2\times\pot{2}{-7}\frac{1}{0{,}001}=\SI{400}{\micro\tesla}.
\]
(b) A $\SI{10}{\centi\meter}$, os dois fios estão praticamente à mesma
distância e seus campos quase se cancelam. O resultado é aproximadamente o de um
\textbf{dipolo}:
\[
B\approx\frac{\mu_0 I d}{2\pi r^{2}}
=\pot{2}{-7}\frac{(1)(0{,}002)}{(0{,}10)^{2}}
=\pot{4}{-8}\,\si{\tesla}.
\]
(c) \textbf{Um fio único} de $\SI{1}{\ampere}$ produziria, a
$\SI{10}{\centi\meter}$:
$B=\pot{2}{-6}\,\si{\tesla}$ --- \textbf{cinquenta vezes mais}.\\
\textbf{Conclusão:} manter os condutores de ida e volta \emph{juntos} reduz
drasticamente o campo irradiado, porque o par se comporta como dipolo
($1/r^{2}$) em vez de fio isolado ($1/r$). Quanto menor a separação $d$, melhor
--- e é exatamente por isso que se recomenda não separar os condutores de um
mesmo circuito.}
\stepcounter{questao}
\resposta{$\SI{25}{\micro\volt\per\milli\tesla}$; $B=\SI{24}{\milli\tesla}$}
\resolucao{(a)
\[
S=\frac{2{,}5\,\si{\milli\volt}}{100\,\si{\milli\tesla}}
=\SI{25}{\micro\volt\per\milli\tesla}.
\]
(b) $B=\dfrac{0{,}6}{0{,}025}=\SI{24}{\milli\tesla}$.\\
(c) \textbf{Fontes de erro, em ordem de importância.}
\begin{itemize}
\item \textbf{Tensão de \emph{offset}:} sensores reais indicam algo mesmo com
$B=0$, por assimetria construtiva. Valores de dezenas de microvolts são comuns
--- comparáveis ao sinal de alguns militesla. \emph{Correção:} medir com o
sensor blindado e subtrair.
\item \textbf{Deriva térmica:} tanto o \emph{offset} quanto a sensibilidade
variam com a temperatura.
\item \textbf{Orientação:} o sensor responde apenas à componente
\emph{perpendicular} à sua face. Um erro angular $\theta$ introduz fator
$\cos\theta$ --- desprezível para ângulos pequenos, mas $13\%$ a
$\ang{30}$.
\item \textbf{Campo terrestre} somado à medida, se não for compensado.
\end{itemize}
\textbf{Procedimento robusto:} medir com o sensor em uma orientação, depois
girá-lo $\ang{180}$ e medir de novo. A média das duas leituras cancela o
\emph{offset}, e a semidiferença dá o campo.}
\stepcounter{questao}
\resposta{$B=\dfrac{\mu_0 I}{2R}$}
\resolucao{(a) A lei de Biot--Savart estabelece que um elemento
$I\,\mathrm{d}\vec\ell$ produz, a uma distância $\vec r$,
\[
\mathrm{d}\vec B=\frac{\mu_0}{4\pi}\,
\frac{I\,\mathrm{d}\vec\ell\times\hat r}{r^{2}}.
\]
Para o \textbf{centro} da espira, todo elemento está à mesma distância $R$, e
$\mathrm{d}\vec\ell$ é sempre \emph{perpendicular} a $\hat r$ --- logo
$|\mathrm{d}\vec\ell\times\hat r|=\mathrm{d}\ell$.\\
(b) Além disso, todos os $\mathrm{d}\vec B$ apontam na \emph{mesma} direção
(perpendicular ao plano da espira), de modo que a integral vetorial vira
escalar:
\[
B=\frac{\mu_0I}{4\pi R^{2}}\oint\mathrm{d}\ell
=\frac{\mu_0I}{4\pi R^{2}}(2\pi R)
=\frac{\mu_0 I}{2R}.
\]
(c) \textbf{Três simplificações coincidiram:} distância constante, produto
vetorial de módulo máximo (ângulo reto) e contribuições todas paralelas. É a
combinação mais favorável possível.\\
\textbf{Fora do centro isso se perde:} no eixo, a distância deixa de ser $R$ e
as contribuições ganham componentes radiais que se cancelam --- exigindo
projeção. Fora do eixo, a integral não tem forma fechada elementar e recorre-se a
métodos numéricos ou a funções elípticas.}
\stepcounter{questao}
\resposta{Ampère exige simetria suficiente para tirar $B$ da integral;
Biot--Savart é geral mas trabalhoso}
\resolucao{(a) \textbf{Biot--Savart:}
$\mathrm{d}\vec B=\dfrac{\mu_0}{4\pi}
\dfrac{I\,\mathrm{d}\vec\ell\times\hat r}{r^{2}}$ --- vale \emph{sempre}, para
qualquer distribuição de corrente. É a lei fundamental.\\
\textbf{Ampère:} $\oint\vec B\cdot\mathrm{d}\vec\ell=\mu_0I_{env}$ --- também
sempre verdadeira, mas só \emph{útil} para calcular $B$ em casos especiais.\\
(b) \textbf{Ampère é praticável quando} existe uma curva fechada ao longo da
qual $|\vec B|$ é constante e $\vec B$ é paralelo (ou perpendicular) ao
caminho. Só assim $B$ pode ser retirado da integral:
\[
\oint\vec B\cdot\mathrm{d}\vec\ell=B\oint\mathrm{d}\ell=B\,\ell.
\]
Isso exige \textbf{simetria} --- cilíndrica, plana ou toroidal --- que precisa
ser \emph{conhecida de antemão}, por argumento físico.\\
(c)
\begin{center}
\begin{tabular}{lll}
\hline
Configuração & Método & Motivo\\
\hline
Fio retilíneo infinito & Ampère & simetria cilíndrica\\
Solenoide/toroide ideal & Ampère & simetria evidente\\
Espira circular (centro) & Biot--Savart & sem curva de $B$ constante\\
Segmento finito de fio & Biot--Savart & sem simetria alguma\\
\hline
\end{tabular}
\end{center}
\textbf{Erro comum:} tentar aplicar Ampère a uma espira circular usando a
própria espira como caminho. O campo \emph{não} é constante nem tangente ao
longo dela --- a integral existe, mas não permite isolar $B$.\\
\textbf{Analogia com a eletrostática:} Gauss está para Coulomb assim como
Ampère está para Biot--Savart. Em ambos os casos, a lei integral é elegante e
poderosa \emph{apenas} quando a simetria coopera.}
\stepcounter{questao}
\resposta{$I=\SI{0.398}{\ampere}$; $P=\SI{0.32}{\watt}$}
\resolucao{(a) $n=\dfrac{500}{0{,}25}=\SI{2000}{\per\meter}$, e
\[
I=\frac{B}{\mu_0 n}=\frac{\pot{1}{-3}}{(4\pi\times10^{-7})(2000)}
=\SI{0.398}{\ampere}.
\]
(b) $P=RI^{2}=2(0{,}398)^{2}=\SI{0.32}{\watt}$ --- desprezível, sem necessidade
de refrigeração.\\
(c) \textbf{Uniformidade.} Supondo diâmetro de $\SI{4}{\centi\meter}$, a razão
$L/D=6{,}25$ --- razoável, mas não excelente. Consequências:
\begin{itemize}
\item No \textbf{centro}, o campo fica cerca de $1$ a $2\%$ abaixo do valor
ideal $\mu_0nI$.
\item Nas \textbf{extremidades}, cai para aproximadamente metade.
\item A região com uniformidade melhor que $1\%$ limita-se ao terço central,
cerca de $\SI{8}{\centi\meter}$.
\end{itemize}
\textbf{Se for necessária melhor uniformidade:} alongar o solenoide (mantendo
$n$, o que exige mais espiras e mais tensão), ou adotar bobinas de Helmholtz.\\
\textbf{Verificação de coerência:} $\SI{1}{\milli\tesla}$ é vinte vezes o campo
terrestre --- o experimento precisa ser orientado ou compensado, sob pena de o
campo ambiente introduzir erro de $5\%$.}
\stepcounter{questao}
\resposta{$B\to\dfrac{\mu_0 m}{2\pi x^{3}}$, com $m=I\pi R^{2}$}
\resolucao{(a) Partindo de
$B=\dfrac{\mu_0IR^{2}}{2(R^{2}+x^{2})^{3/2}}$ e fazendo $x\gg R$:
\[
B\to\frac{\mu_0IR^{2}}{2x^{3}}.
\]
(b) Definindo o \textbf{momento de dipolo magnético}
$m=I\,A=I\pi R^{2}$:
\[
B=\frac{\mu_0 m}{2\pi x^{3}}.
\]
\textbf{Verificação numérica:} para $R=\SI{10}{\centi\meter}$ e
$I=\SI{5}{\ampere}$, tem-se $m=\SI{0.157}{\ampere\meter\squared}$ e, a
$\SI{1}{\meter}$, $B=\pot{3.14}{-8}\,\si{\tesla}$. O valor exato é
$\pot{3.10}{-8}$ --- a aproximação erra $1{,}5\%$ a dez raios de distância.\\
(c) \textbf{Comparação com o dipolo elétrico.} Ambos decaem com $1/r^{3}$, e as
expressões são formalmente análogas ($p=2qa$ lá, $m=IA$ aqui).\\
\textbf{Mas há uma diferença fundamental.} O dipolo elétrico é feito de duas
cargas que \emph{podem} ser separadas --- e então cada uma produz campo
$1/r^{2}$. Já o dipolo magnético \textbf{não pode} ser separado: não existem
monopolos magnéticos.\\
\textbf{Consequência:} o dipolo é o termo \emph{dominante} de qualquer fonte
magnética estacionária. Não há termo $1/r^{2}$ em magnetostática --- e é por isso
que $\nabla\cdot\vec B=0$ figura entre as equações de Maxwell, sem análogo
elétrico.}
\stepcounter{questao}
\resposta{Não há "cargas magnéticas" a redistribuir; usa-se desvio por material
de alto $\mu_r$ ou correntes induzidas}
\resolucao{(a) \textbf{Por que a gaiola falha.} A blindagem elétrica funciona
porque o condutor possui \emph{cargas livres} que se redistribuem até anular o
campo interno. Não existe equivalente magnético: \textbf{não há monopolos
magnéticos livres} para migrar.\\
Um campo magnético \emph{estático} atravessa o cobre e o alumínio praticamente
sem atenuação.\\
(b) \textbf{Dois mecanismos disponíveis.}
\begin{itemize}
\item \textbf{Desvio de fluxo} (campos estáticos ou lentos). Um invólucro de
material de alto $\mu_r$ --- \emph{mu-metal}, com $\mu_r$ de dezenas de
milhares --- oferece caminho de baixa relutância. O fluxo prefere circular pelo
material e "contorna" a região protegida. Não há cancelamento: há \emph{desvio}.
\item \textbf{Correntes induzidas} (campos variáveis). Pela lei de Lenz, um
campo alternado induz correntes em um invólucro condutor, e essas correntes
geram campo oposto. A atenuação cresce com a frequência e com a condutividade.
\end{itemize}
(c) \textbf{Dependência com a frequência.}
\begin{center}
\begin{tabular}{lll}
\hline
Faixa & Mecanismo eficaz & Material\\
\hline
Estático (CC) & desvio de fluxo & mu-metal\\
Baixa (dezenas de Hz) & desvio, pouca indução & mu-metal\\
Alta (kHz e acima) & correntes induzidas & cobre, alumínio\\
\hline
\end{tabular}
\end{center}
\textbf{Consequência prática:} blindar campo elétrico é fácil e barato --- basta
qualquer metal aterrado. Blindar campo magnético estático é caro, exige material
especial, e a atenuação típica fica em uma ou duas ordens de grandeza, contra
muitas ordens no caso elétrico.\\
É por isso que salas de ressonância magnética têm blindagem magnética de
toneladas, enquanto a blindagem elétrica de um cabo é uma malha fina.}
\stepcounter{questao}
\resposta{Medir com inversão de orientação para cancelar \emph{offset};
três orientações ortogonais para as componentes}
\resolucao{(a) \textbf{Procedimento.}
\begin{enumerate}
\item \textbf{Caracterizar o zero:} com a bobina \emph{desligada}, registrar a
leitura em cada ponto. Isso captura tanto o \emph{offset} do sensor quanto o
campo ambiente (terrestre e de equipamentos próximos).
\item \textbf{Medir com a bobina ligada} nos mesmos pontos.
\item \textbf{Subtrair} ponto a ponto. O resultado é o campo da bobina,
livre de \emph{offset} e de campo ambiente.
\end{enumerate}
\textbf{Alternativa mais robusta:} inverter o sentido da corrente na bobina e
tomar a \emph{semidiferença} das duas leituras. Isso cancela tudo o que não
depende da corrente.\\
(b) \textbf{Três componentes.} O sensor Hall responde apenas à componente
perpendicular à sua face. Para obter $\vec B$ completo, meça em três orientações
ortogonais no mesmo ponto:
\[
|\vec B|=\sqrt{B_x^{2}+B_y^{2}+B_z^{2}}.
\]
Um sensor triaxial faz isso simultaneamente e evita erros de reposicionamento.\\
(c) \textbf{Critérios de validação.}
\begin{itemize}
\item \textbf{Linearidade com a corrente:} $B$ deve ser proporcional a $I$.
Meça em três correntes e verifique --- desvio indica saturação de núcleo ou
não linearidade do sensor.
\item \textbf{Simetria:} pontos simétricos em relação ao eixo devem dar valores
iguais. Discrepância revela desalinhamento do sensor ou da bobina.
\item \textbf{Valor central:} deve concordar com $\mu_0nI$ dentro da incerteza,
descontada a correção de comprimento finito.
\item \textbf{Reprodutibilidade:} repita o primeiro ponto ao final. Diferença
indica deriva térmica do sensor.
\end{itemize}
\textbf{Registre sempre a temperatura} --- a sensibilidade Hall varia
tipicamente algumas centenas de ppm por grau.}
\stepcounter{questao}
\resposta{Núcleos saturam entre $1{,}5$ e $\SI{2}{\tesla}$; acima disso é
preciso bobina sem núcleo, com corrente muito alta}
\resolucao{(a) \textbf{Saturação.} Materiais ferromagnéticos alinham seus
domínios completamente em torno de $1{,}5$ a $\SI{2}{\tesla}$ (ferro-silício,
$\approx\SI{2}{\tesla}$; ferrites, bem menos). Acima disso, $\mu_r$ cai para
próximo de $1$ e o núcleo deixa de contribuir --- \textbf{acrescentar corrente
passa a render tão pouco quanto no ar}.\\
Este é um limite \emph{de material}, não de projeto: nenhuma geometria o
contorna.\\
(b) \textbf{Limite térmico.} Sem núcleo, $B=\mu_0nI$ exige corrente enorme. A
potência dissipada vai com $I^{2}$, e a densidade de corrente admissível em
cobre refrigerado a água fica em torno de
$\SI{20}{\ampere\per\milli\meter\squared}$.\\
Eletroímãs resistivos de laboratório atingem $1$ a $\SI{2}{\tesla}$ consumindo
\emph{megawatts} --- e boa parte do custo operacional é a refrigeração.\\
(c) \textbf{Soluções empregadas.}
\begin{itemize}
\item \textbf{Supercondutores:} resistência nula elimina o limite térmico.
Bobinas de NbTi e $\mathrm{Nb_3Sn}$ chegam a $\SI{20}{\tesla}$ e são a base de
equipamentos de ressonância magnética e de aceleradores. O custo passa a ser a
criogenia.
\item \textbf{Campos pulsados:} correntes altíssimas por milissegundos, antes
que o calor se acumule. Alcançam $\SI{100}{\tesla}$, com a bobina
frequentemente destruída no processo.
\item \textbf{Ímãs híbridos:} bobina supercondutora externa somada a uma
resistiva interna --- combinação que detém os recordes de campo contínuo, acima
de $\SI{45}{\tesla}$.
\end{itemize}
\textbf{Limite último:} acima de algumas dezenas de tesla, a \emph{pressão
magnética} $B^{2}/2\mu_0$ supera a resistência mecânica de qualquer material. A
$\SI{100}{\tesla}$ ela vale cerca de $\SI{4}{\giga\pascal}$ --- a bobina
literalmente se rompe.}
\stepcounter{questao}
\resposta{O campo na cavidade é \textbf{uniforme}, de módulo
$\mu_0Jd/2$}
\resolucao{(a) \textbf{Estratégia.} O condutor furado equivale à
\emph{superposição} de:
\begin{itemize}
\item um cilindro \textbf{maciço} de raio $a$ com densidade de corrente
$+J$; \emph{mais}
\item um cilindro de raio $b$, na posição da cavidade, com densidade
$-J$.
\end{itemize}
A soma reproduz corrente $J$ na região sólida e zero na cavidade --- exatamente
a situação real.\\
(b) Dentro de um cilindro maciço de densidade $J$, a lei de Ampère dá
$B=\dfrac{\mu_0 J r}{2}$, ou vetorialmente
$\vec B=\dfrac{\mu_0}{2}\,\vec J\times\vec r$, com $\vec r$ medido a partir do
\emph{eixo daquele} cilindro.\\
Somando as duas contribuições em um ponto da cavidade, com $\vec r_1$ e
$\vec r_2$ medidos a partir de cada eixo:
\[
\vec B=\frac{\mu_0}{2}\vec J\times\vec r_1
-\frac{\mu_0}{2}\vec J\times\vec r_2
=\frac{\mu_0}{2}\vec J\times(\vec r_1-\vec r_2)
=\frac{\mu_0}{2}\vec J\times\vec d.
\]
\textbf{O vetor $\vec d$ é constante} --- ele liga os dois eixos e não depende do
ponto.\\
(c) \textbf{Interpretação: o campo na cavidade é rigorosamente uniforme}, de
módulo $\mu_0Jd/2$ e direção perpendicular a $\vec d$.\\
Isso é notável: uma geometria sem simetria evidente produz, na região vazia, o
campo mais simples possível.\\
\textbf{Aplicação:} o mesmo resultado, com $\vec J$ substituído por densidade de
carga, dá campo elétrico uniforme na cavidade de uma esfera carregada --- e é
usado no projeto de aceleradores e de ímãs de deflexão, em que se busca campo
uniforme em uma região de acesso.}

\gabarito[12.2 Campo magnético]

\subsection{Fluxo magnético}
\stepcounter{questao}
\resposta{(b)}
\resolucao{Como $\Phi = B\,A\cos\theta$, sendo $\theta$ o ângulo entre o campo e a
\emph{normal} à espira, o fluxo é máximo quando $\theta = 0$ --- ou seja, quando o
campo atravessa a espira perpendicularmente ao seu plano. Se o plano fica paralelo
ao campo, $\theta = \SI{90}{\degree}$ e o fluxo é nulo.}
\stepcounter{questao}
\resposta{$\Phi = \SI{1}{\milli\weber}$}
\resolucao{Com $\theta = 0$ ($\cos\theta = 1$) e
$A = \SI{20}{\centi\meter\squared} = \pot{20}{-4}\,\si{\meter\squared}$:
\[
\Phi = B\,A = 0{,}5\cdot\pot{20}{-4} = \pot{1}{-3}\,\si{\weber}=\SI{1}{\milli\weber}.
\]}
\stepcounter{questao}
\resposta{(a) $\Phi \approx \SI{6.03}{\milli\weber}$; (b) $\varepsilon \approx \SI{8.04}{\milli\volt}$}
\resolucao{(a) A área é $A = \pi R^2 = \pi(0{,}08)^2 \approx \SI{0.0201}{\meter\squared}$:
\[
\Phi = B\,A = 0{,}3\cdot0{,}0201 \approx \SI{6.03}{\milli\weber}.
\]
(b) A fem é o módulo da taxa de variação do fluxo. Como só $B$ muda:
\[
\varepsilon = \left|\frac{\Delta\Phi}{\Delta t}\right|
= A\,\frac{|\Delta B|}{\Delta t}
= 0{,}0201\cdot\frac{0{,}2}{0{,}5}\approx\SI{8.04}{\milli\volt}.
\]}
\stepcounter{questao}
\resposta{$\Phi=\SI{7.07e-10}{\weber}$}
\resolucao{$\Phi=BA\cos\theta=(10^{-4})(10^{-5})\cos45^\circ
=7{,}07\times10^{-10}\,\si{\weber}$.}
\stepcounter{questao}
\resposta{$0$--$2\,$s: $\varepsilon=\SI{2}{\milli\volt}$; $2$--$4\,$s:
$\varepsilon=0$; $4$--$6\,$s: $\varepsilon=\SI{-2}{\milli\volt}$}
\resolucao{A fem é a \emph{inclinação} do gráfico $\Phi \times t$ (com sinal
trocado, por Lenz):
\begin{itemize}[leftmargin=1.6em,itemsep=1pt,topsep=2pt]
\item \textbf{$0$ a $\SI{2}{\second}$:} $\Phi$ sobe de $0$ a
      \SI{4}{\milli\weber}, inclinação $\dfrac{4}{2} = \SI{2}{\milli\weber\per\second}$,
      logo $|\varepsilon| = \SI{2}{\milli\volt}$.
\item \textbf{$2$ a $\SI{4}{\second}$:} $\Phi$ constante, inclinação nula,
      $\varepsilon = 0$.
\item \textbf{$4$ a $\SI{6}{\second}$:} $\Phi$ desce de \SI{4}{} a $0$, inclinação
      $-\SI{2}{\milli\weber\per\second}$, $\varepsilon = \SI{-2}{\milli\volt}$
      (sentido oposto).
\end{itemize}
A fem só existe quando há \emph{variação} de fluxo, e seu sinal se inverte entre a
subida e a descida --- exatamente o princípio de um gerador.}
\stepcounter{questao}
\resposta{(a) $\Phi_{\max}=\SI{8}{\milli\weber}$;
(b) $\varepsilon_{\text{pico}}\approx\SI{1.51}{\volt}$; (c) \SI{30}{\hertz};
(d) tudo $\times100$}
\resolucao{\textbf{(a)} Com a espira girando, o ângulo entre a normal e o campo
varia: $\theta = \omega t$. Logo
\[
\Phi(t)=B\,A\cos(\omega t),
\qquad
\Phi_{\max}=B\,A=0{,}4\cdot\pot{200}{-4}=\SI{8}{\milli\weber}.
\]
\textbf{(b)} A velocidade angular é
\[
\omega = \frac{2\pi\cdot1800}{60}=60\pi\approx\SI{188.5}{\radian\per\second}.
\]
Pela lei de Faraday,
\[
\varepsilon(t)=-\frac{\mathrm{d}\Phi}{\mathrm{d}t}
= B A\,\omega\sin(\omega t),
\qquad
\varepsilon_{\text{pico}}=BA\omega
=(\pot{8}{-3})(188{,}5)\approx\SI{1.51}{\volt}.
\]
\textbf{(c)} A frequência é
$f = \dfrac{1800}{60} = \SI{30}{\hertz}$ (uma volta $=$ um ciclo completo).\\
\textbf{Por que senoidal?} Porque o \emph{fluxo} varia cossenoidalmente com o
ângulo (projeção da área sobre a direção do campo), e a fem é sua
\emph{derivada}. A derivada do cosseno é o seno --- daí a defasagem de
$\SI{90}{\degree}$ entre fluxo e tensão. Este é literalmente o princípio de
funcionamento do \textbf{alternador}: a forma de onda senoidal da rede elétrica
nasce da geometria circular do movimento.

\textbf{(d)} Com $N$ espiras, todas concatenam o mesmo fluxo e suas fem se somam
em série:
\[
\varepsilon_{\text{pico}} = N B A \omega \approx 100\cdot1{,}51 = \SI{151}{\volt}.
\]
O fluxo \emph{por espira} não muda; o que muda é o \emph{fluxo concatenado}
$N\Phi$. É assim que geradores reais atingem centenas de volts com campos
modestos: aumentando $N$, $B$, $A$ ou $\omega$ --- os quatro parâmetros de
projeto de uma máquina elétrica.}
\stepcounter{questao}
\resposta{$\Phi=\SI{0.5}{\milli\weber}$}
\resolucao{
\[
\Phi=BA\cos\theta=(0{,}1)(0{,}01)\cos\ang{60}
=(0{,}1)(0{,}01)(0{,}5)=\pot{5}{-4}\,\si{\weber}.
\]
\textbf{Atenção ao ângulo:} $\theta$ é medido entre $\vec B$ e a \textbf{normal}
à superfície. Se o enunciado der o ângulo com o \emph{plano}, use o complemento
--- é a troca de $\cos$ por $\sin$ que mais gera erro neste tópico.}
\stepcounter{questao}
\resposta{$\Phi=0$}
\resolucao{Campo paralelo ao plano significa campo \emph{perpendicular} à
normal, logo $\theta=\ang{90}$ e $\cos\theta=0$:
\[
\Phi=0.
\]
\textbf{Interpretação geométrica:} nenhuma linha de campo \emph{atravessa} a
superfície --- elas passam rentes a ela.\\
\textbf{Casos-limite úteis:} $\Phi$ é máximo com $\vec B$ perpendicular ao plano
(paralelo à normal) e nulo com $\vec B$ contido no plano.}
\stepcounter{questao}
\resposta{$A=\SI{50}{\centi\meter\squared}$}
\resolucao{
\[
A=\frac{\Phi}{B}=\frac{\pot{2}{-3}}{0{,}4}=\pot{5}{-3}\,\si{\meter\squared}
=\SI{50}{\centi\meter\squared}.
\]
\textbf{Lembre da conversão:}
$\SI{1}{\meter\squared}=\SI{e4}{\centi\meter\squared}$, pois a área é o
\emph{quadrado} de um comprimento.}
\stepcounter{questao}
\resposta{$B=\SI{0.5}{\tesla}$}
\resolucao{
\[
B=\frac{\Phi}{A}=\frac{\pot{1.5}{-3}}{\pot{3}{-3}}=\SI{0.5}{\tesla}.
\]
\textbf{Nomenclatura:} $B$ é chamado \emph{densidade de fluxo} justamente por
ser fluxo por unidade de área --- e $\SI{1}{\tesla}=\SI{1}{\weber\per\meter\squared}$.\\
As duas grandezas descrevem a mesma realidade em escalas diferentes: $B$ é
local, $\Phi$ é integral.}
\stepcounter{questao}
\resposta{(b)}
\resolucao{As linhas de campo magnético são \textbf{fechadas} --- não têm início
nem fim. Toda linha que entra em uma superfície fechada necessariamente sai dela,
e as contribuições se cancelam:
\[
\oint\vec B\cdot\mathrm{d}\vec A=0.
\]
\textbf{Contraste com o campo elétrico:} lá, a lei de Gauss dá
$\oint\vec E\cdot\mathrm{d}\vec A=Q_{env}/\varepsilon_0$ --- diferente de zero
sempre que houver carga envolvida.\\
\textbf{A diferença tem uma causa única:} existem cargas elétricas isoladas
(monopolos elétricos), mas \textbf{não existem monopolos magnéticos}. Esta é uma
das equações de Maxwell, e sua forma diferencial é $\nabla\cdot\vec B=0$.}
\stepcounter{questao}
\resposta{$\pot{5}{5}\,\si{\maxwell}$}
\resolucao{
\[
\Phi=\pot{5}{-3}\times10^{8}=\pot{5}{5}\,\si{\maxwell}.
\]
\textbf{Origem da unidade:} o maxwell pertence ao sistema CGS, e vale
$\SI{1}{\gauss\centi\meter\squared}$. A conversão decorre diretamente de
$\SI{1}{\tesla}=\SI{e4}{\gauss}$ e
$\SI{1}{\meter\squared}=\SI{e4}{\centi\meter\squared}$:
\[
10^{4}\times10^{4}=10^{8}.
\]
Embora obsoleto no SI, o maxwell ainda aparece em catálogos de ímãs e em
literatura mais antiga.}
\stepcounter{questao}
\resposta{$\lambda=\SI{50}{\milli\weber}$-espira; $L=\SI{25}{\milli\henry}$}
\resolucao{(a)
\[
\lambda=N\Phi=50(\pot{1}{-3})=\pot{5}{-2}\,\si{\weber}\text{-espira}.
\]
(b)
\[
L=\frac{\lambda}{I}=\frac{0{,}05}{2}=\SI{25}{\milli\henry}.
\]
\textbf{Distinção essencial.} $\Phi$ é o fluxo através de \emph{uma} espira;
$\lambda=N\Phi$ é o \emph{fluxo concatenado}, que conta cada espira
separadamente. É $\lambda$, e não $\Phi$, que entra na definição de indutância e
na lei de Faraday para bobinas ($\varepsilon=-\mathrm{d}\lambda/\mathrm{d}t$).}
\stepcounter{questao}
\resposta{$\Phi=\SI{2.5}{\milli\weber}$; máximo de $\SI{5}{\milli\weber}$}
\resolucao{
\[
\Phi=BA\cos\ang{60}=(0{,}25)(\pot{2}{-2})(0{,}5)=\pot{2.5}{-3}\,\si{\weber};
\]
\[
\Phi_{\max}=BA=\pot{5}{-3}\,\si{\weber}\quad(\theta=0).
\]
\textbf{Aplicação:} girar a espira entre essas duas posições varia o fluxo de
$\SI{2.5}{\milli\weber}$. Feito em $\SI{10}{\milli\second}$, isso induziria
$\varepsilon=\SI{0.25}{\volt}$ por espira --- princípio do gerador.}
\stepcounter{questao}
\resposta{$\Phi=\SI{2.51}{\micro\weber}$}
\resolucao{
\[
B=\mu_0nI=(4\pi\times10^{-7})(1000)(2)=\pot{2.51}{-3}\,\si{\tesla};
\]
\[
\Phi=BA=(\pot{2.51}{-3})(\pot{1}{-3})=\pot{2.51}{-6}\,\si{\weber}.
\]
\textbf{Ordem de grandeza:} apenas $\SI{2.5}{\micro\weber}$ --- fluxos em
núcleo de ar são pequenos. Com núcleo ferromagnético de
$\mu_r=1000$, o mesmo solenoide produziria $\SI{2.5}{\milli\weber}$.\\
\textbf{É por isso que máquinas elétricas usam núcleo:} sem ele, o fluxo (e
portanto a f.e.m. induzida e o torque) seria mil vezes menor.}
\stepcounter{questao}
\resposta{$\varepsilon=\SI{20}{\volt}$}
\resolucao{
\[
|\varepsilon|=N\frac{\Delta\Phi}{\Delta t}
=100\,\frac{\pot{4}{-3}}{\pot{2}{-2}}=\SI{20}{\volt}.
\]
\textbf{Como a variação é linear}, a f.e.m. é \emph{constante} durante todo o
intervalo. Se o fluxo variasse de outra forma, seria preciso derivar ponto a
ponto.\\
\textbf{Note o papel de $N$:} cem espiras multiplicam a f.e.m. por cem, embora o
fluxo \emph{por espira} seja o mesmo. É o fluxo concatenado que importa.}
\stepcounter{questao}
\resposta{São iguais}
\resolucao{O fluxo através de qualquer superfície depende apenas de sua
\textbf{borda}.\\
\textbf{Demonstração:} as duas superfícies, juntas, formam uma superfície
\emph{fechada}. Como
$\oint\vec B\cdot\mathrm{d}\vec A=0$, o fluxo que entra por uma sai pela outra
--- e portanto os fluxos através delas são iguais (atentando para a orientação
das normais).\\
\textbf{Consequência prática:} na lei de Faraday, pode-se escolher
\emph{qualquer} superfície apoiada no circuito. Escolha a que tornar a
integral mais simples --- geralmente a plana, mas às vezes uma superfície
curva evita regiões de campo complicado.\\
\textbf{Contraste:} isso \emph{não} vale para o fluxo elétrico, que depende da
carga envolvida e portanto da superfície escolhida.}
\stepcounter{questao}
\resposta{$\Delta\Phi=\SI{-24}{\milli\weber}$}
\resolucao{
\[
\Phi_i=BA=(0{,}6)(0{,}04)=\SI{24}{\milli\weber};
\qquad
\Phi_f=BA\cos\ang{90}=0.
\]
\[
\Delta\Phi=0-0{,}024=\SI{-24}{\milli\weber}.
\]
\textbf{Duas formas de variar o fluxo aparecem no tópico:} mudar $B$, mudar
$A$, ou mudar a \emph{orientação}. Esta última é a base do gerador rotativo, em
que $\theta=\omega t$ e o fluxo varia senoidalmente.\\
\textbf{Sinal:} o fluxo diminuiu. Pela lei de Lenz, a corrente induzida
circularia no sentido de \emph{manter} o fluxo original.}
\stepcounter{questao}
\resposta{$B=\SI{1.5}{\tesla}$ --- próximo da saturação típica}
\resolucao{
\[
B=\frac{\Phi}{A}=\frac{\pot{6}{-4}}{\pot{4}{-4}}=\SI{1.5}{\tesla}.
\]
\textbf{Avaliação:} aços-silício saturam entre $1{,}5$ e $\SI{2}{\tesla}$;
ferrites, bem antes (tipicamente $0{,}3$ a $\SI{0.5}{\tesla}$).\\
Portanto, para aço-silício este núcleo opera \textbf{no limite}. Qualquer
aumento de corrente levaria à saturação, com queda abrupta de $\mu_r$ e
distorção severa da forma de onda.\\
\textbf{Prática de projeto:} dimensione para $B_{\max}$ entre $60$ e
$80\%$ da saturação, deixando margem para transitórios e para a variação com a
temperatura.}
\stepcounter{questao}
\resposta{$\lambda=\SI{120}{\milli\weber}$-espira}
\resolucao{
\[
\Phi=BA=(0{,}4)(\pot{1.5}{-3})=\pot{6}{-4}\,\si{\weber};
\]
\[
\lambda=N\Phi=200(\pot{6}{-4})=\pot{0.12}{}\,\si{\weber}\text{-espira}.
\]
\textbf{Se essa bobina fosse retirada do campo em $\SI{50}{\milli\second}$:}
\[
|\varepsilon|=\frac{\Delta\lambda}{\Delta t}=\frac{0{,}12}{0{,}05}
=\SI{2.4}{\volt}.
\]
É esse o princípio do \emph{fluxímetro}: integrando a f.e.m. medida, recupera-se
$\Delta\lambda$ e, com $N$ e $A$ conhecidos, o campo.}
\stepcounter{questao}
\resposta{Cresce de zero a $BA$, com taxa constante enquanto a espira atravessa
a fronteira}
\resolucao{\textbf{Três fases:}
\begin{itemize}
\item \textbf{Totalmente fora:} $\Phi=0$ e $\mathrm{d}\Phi/\mathrm{d}t=0$ ---
nenhuma f.e.m.
\item \textbf{Atravessando a fronteira:} a área \emph{dentro} do campo cresce
linearmente com o deslocamento. Se a velocidade for constante,
$\Phi$ cresce linearmente e a f.e.m. é \textbf{constante}.
\item \textbf{Totalmente dentro:} $\Phi=BA$ e volta a ser constante ---
nenhuma f.e.m., mesmo que a espira continue se movendo.
\end{itemize}
\textbf{Ponto conceitual importante:} não é o \emph{movimento} que induz f.e.m.,
e sim a \textbf{variação de fluxo}. Uma espira movendo-se rapidamente dentro de
um campo uniforme não gera nada.\\
Com largura $\ell$ e velocidade $v$, a f.e.m. na fase intermediária vale
$\varepsilon=B\ell v$.}
\stepcounter{questao}
\resposta{$\Phi$ é integral (weber); $B$ é local (tesla). $\Phi$ governa a
indução; $B$ governa forças e saturação}
\resolucao{\textbf{Densidade de fluxo $\vec B$:} grandeza \emph{local} e
vetorial, medida em tesla. Ela é o que determina:
\begin{itemize}
\item a força sobre cargas e condutores ($\vec F=q\vec v\times\vec B$);
\item a saturação do material ($B$ próximo de $B_{sat}$);
\item o que um sensor Hall mede.
\end{itemize}
\textbf{Fluxo $\Phi$:} grandeza \emph{integral} e escalar, medida em weber.
Ela é o que determina:
\begin{itemize}
\item a f.e.m. induzida ($\varepsilon=-\mathrm{d}\Phi/\mathrm{d}t$);
\item a indutância ($L=\lambda/I$);
\item o acoplamento entre bobinas.
\end{itemize}
\textbf{Relação:} $\Phi=\int\vec B\cdot\mathrm{d}\vec A$, e no caso uniforme
$B=\Phi/A$.\\
\textbf{Erro comum:} falar em "fluxo" quando se quer dizer "campo". Um ímã
pequeno e um grande podem ter o mesmo $B$ na superfície e fluxos totais muito
diferentes --- e é o fluxo que determina o desempenho em um circuito magnético.}
\stepcounter{questao}
\resposta{$\Phi=\SI{5}{\milli\weber}$ --- igual à estimativa pelo valor médio}
\resolucao{(a) Dividindo em faixas verticais de largura $\mathrm{d}x$ e área
$h\,\mathrm{d}x$:
\[
\Phi=\int_0^{L}B_0\frac{x}{L}\,h\,\mathrm{d}x
=\frac{B_0h}{L}\left[\frac{x^{2}}{2}\right]_0^{L}
=\frac{B_0hL}{2}.
\]
\[
\Phi=\frac{(0{,}5)(0{,}10)(0{,}20)}{2}=\pot{5}{-3}\,\si{\weber}.
\]
(b)
\begin{center}
\begin{tabular}{lc}
\hline
Estimativa & $\Phi$\\
\hline
Usando $B_0$ (valor máximo) & $\SI{10}{\milli\weber}$\\
\textbf{Integração exata} & $\SI{5}{\milli\weber}$\\
Usando $\bar B=B_0/2$ & $\SI{5}{\milli\weber}$\\
\hline
\end{tabular}
\end{center}
\textbf{A média coincide com o exato} --- mas apenas porque $B$ varia
\emph{linearmente}. Para $B\propto x^{2}$, a integração daria $B_0hL/3$,
enquanto a média aritmética dos extremos daria $B_0hL/2$ --- erro de
$50\%$.\\
\textbf{Regra:} substituir o campo pela média só é legítimo em variação linear.
Fora disso, integre.}
\stepcounter{questao}
\resposta{$\Phi=\dfrac{\mu_0Ia}{2\pi}\ln\dfrac{d+b}{d}=\SI{0.55}{\micro\weber}$}
\resolucao{(a) O campo do fio \textbf{não é uniforme} na região da espira:
$B(r)=\dfrac{\mu_0I}{2\pi r}$. Dividindo a espira em faixas de largura
$\mathrm{d}r$ e comprimento $a$:
\[
\Phi=\int_d^{d+b}\frac{\mu_0I}{2\pi r}\,a\,\mathrm{d}r
=\frac{\mu_0Ia}{2\pi}\ln\!\left(\frac{d+b}{d}\right).
\]
(b) Com $a=\SI{0.20}{\meter}$, $d=\SI{0.05}{\meter}$ e
$b=\SI{0.15}{\meter}$ (borda externa em $\SI{0.20}{\meter}$):
\[
\Phi=\pot{2}{-7}(10)(0{,}20)\ln 4
=(\pot{4}{-7})(1{,}386)=\pot{5.55}{-7}\,\si{\weber}.
\]
\textbf{Observe o logaritmo:} afastar a espira reduz o fluxo muito devagar.
Dobrar $d$ e $b$ juntos (mantendo a razão) \emph{não muda nada} --- o fluxo
depende apenas da razão $(d+b)/d$ e do comprimento $a$.\\
\textbf{Consequência:} a indutância mútua entre um fio e uma espira próxima é
notavelmente insensível à distância absoluta. É por isso que interferência
magnética de cabos de potência é difícil de eliminar apenas afastando-os.}
\stepcounter{questao}
\resposta{$\mathcal{R}=\pot{4.97}{5}\,\si{\per\henry}$;
$\Phi=\SI{2}{\milli\weber}$; $B=\SI{5}{\tesla}$ --- \textbf{impossível}}
\resolucao{(a) Em analogia ao circuito elétrico, a \textbf{relutância} vale
\[
\mathcal{R}=\frac{\ell}{\mu_r\mu_0A}
=\frac{0{,}50}{(2000)(4\pi\times10^{-7})(\pot{4}{-4})}
=\pot{4.97}{5}\,\si{\per\henry}.
\]
(b) Com a força magnetomotriz $\mathcal{F}=NI=\SI{1000}{}$ ampère-espiras:
\[
\Phi=\frac{\mathcal{F}}{\mathcal{R}}=\frac{1000}{\pot{4.97}{5}}
=\pot{2.01}{-3}\,\si{\weber};
\qquad
B=\frac{\Phi}{A}=\SI{5.03}{\tesla}.
\]
(c) \textbf{O resultado é impossível.} Nenhum material ferromagnético comum
sustenta $\SI{5}{\tesla}$ --- todos saturam entre $1{,}5$ e
$\SI{2}{\tesla}$.\\
\textbf{O que realmente aconteceria:} conforme $B$ se aproxima de
$\SI{2}{\tesla}$, $\mu_r$ despenca de $2000$ para valores próximos de $1$. A
relutância dispara, e o fluxo estaciona em torno de
$B_{sat}A=\SI{0.8}{\milli\weber}$.\\
\textbf{Lição de método:} o modelo linear com $\mu_r$ constante é válido
\emph{apenas} abaixo da saturação. Calcular e depois \textbf{verificar $B$}
contra $B_{sat}$ é etapa obrigatória --- exatamente como se verifica a rigidez
dielétrica em problemas de campo elétrico.}
\stepcounter{questao}
\resposta{$\mathcal{R}_g=\pot{1.99}{6}$, quatro vezes a do ferro;
$B=\SI{1.0}{\tesla}$}
\resolucao{(a) No entreferro, $\mu_r=1$:
\[
\mathcal{R}_g=\frac{g}{\mu_0A}
=\frac{\pot{1}{-3}}{(4\pi\times10^{-7})(\pot{4}{-4})}
=\pot{1.99}{6}\,\si{\per\henry}.
\]
Em série com o ferro:
$\mathcal{R}_{tot}=\pot{4.97}{5}+\pot{1.99}{6}=\pot{2.49}{6}$.\\
(b)
\[
\Phi=\frac{1000}{\pot{2.49}{6}}=\pot{4.02}{-4}\,\si{\weber};
\qquad
B=\SI{1.005}{\tesla}.
\]
\textbf{Agora sim, valor fisicamente sustentável} --- e com margem para
saturação.\\
(c) \textbf{Por que o entreferro domina.} Ele tem apenas
$\SI{1}{\milli\meter}$ contra $\SI{500}{\milli\meter}$ de ferro --- é
$500$ vezes mais curto. Mas sua permeabilidade é $2000$ vezes menor. O saldo:
\[
\frac{\mathcal{R}_g}{\mathcal{R}_{fe}}=\frac{2000}{500}=4.
\]
\textbf{Um milímetro de ar vale quatro vezes meio metro de ferro.}\\
\textbf{Consequências de projeto.}
\begin{itemize}
\item O entreferro \emph{lineariza} o circuito magnético: como ele domina, a
relação $\Phi\times I$ torna-se quase reta, mesmo com o ferro não linear.
\item Ele permite armazenar muito mais energia (que se concentra no ar) --- daí
seu uso em indutores de conversores chaveados.
\item Em contrapartida, reduz a indutância e exige mais ampère-espiras.
\item \textbf{Entreferros parasitas} (juntas mal ajustadas de núcleos em E ou U)
degradam severamente o desempenho, mesmo com dezenas de micrômetros.
\end{itemize}}
\stepcounter{questao}
\resposta{$\SI{1}{\milli\weber}$ de dispersão; $k\approx0{,}8$}
\resolucao{(a) \textbf{Dispersão:} a parcela que não atravessa a outra bobina:
\[
\Phi_{disp}=5-4=\SI{1}{\milli\weber}\quad(20\%).
\]
(b) O coeficiente de acoplamento é a fração de fluxo compartilhada:
\[
k\approx\frac{\Phi_{12}}{\Phi_1}=\frac{4}{5}=0{,}8.
\]
\emph{(A definição rigorosa é $k=M/\sqrt{L_1L_2}$, que coincide com esta razão
quando as bobinas são simétricas.)}\\
(c) \textbf{Como aumentar $k$.}
\begin{itemize}
\item \textbf{Núcleo fechado} de alta permeabilidade: o fluxo prefere o caminho
de baixa relutância e "não escapa". Toroides atingem $k>0{,}99$.
\item \textbf{Enrolamentos concêntricos ou intercalados}, em vez de lado a lado
--- reduz drasticamente o caminho de fuga.
\item \textbf{Eliminar entreferros} desnecessários, que forçam o fluxo a
divergir.
\end{itemize}
\textbf{Compromisso:} $k$ alto é desejável em transformadores (transferência
eficiente), mas dispersão \emph{controlada} é útil em alguns casos --- ela
funciona como indutância série natural, limitando corrente de curto em
transformadores de solda e em reatores.}
\stepcounter{questao}
\resposta{Mesmo fluxo; $B$ três vezes maior na seção estreita, que satura
primeiro}
\resolucao{(a) \textbf{O fluxo é o mesmo.} Sem dispersão, todas as linhas que
passam por um trecho passam pelo outro --- é a consequência direta de
$\oint\vec B\cdot\mathrm{d}\vec A=0$ aplicada a uma superfície que envolve a
junção.\\
\textbf{Analogia elétrica:} é o equivalente magnético da LKC. O fluxo se
conserva ao longo de um circuito magnético sem ramificação, como a corrente em
um circuito série.\\
(b) Como $B=\Phi/A$:
\[
\frac{B_2}{B_1}=\frac{A_1}{A_2}=\frac{6}{2}=3.
\]
(c) \textbf{A seção estreita satura primeiro}, por ter $B$ três vezes maior.\\
\textbf{Consequência de projeto:} o dimensionamento de um núcleo é governado
pela sua \textbf{menor} seção. Aumentar a seção larga não adianta nada ---
é preciso alargar o gargalo.\\
\textbf{Onde isso aparece na prática:} cantos de núcleos em E, junções mal
projetadas e regiões com furos de fixação. São esses pontos que saturam
primeiro e produzem aquecimento e distorção localizados.}
\stepcounter{questao}
\resposta{$\Phi=\SI{1.2}{\milli\weber}$; $\varepsilon=\SI{30}{\milli\volt}$}
\resolucao{(a) Apenas a área \emph{dentro} do campo contribui:
\[
\Phi=B\,A_{efetiva}=(0{,}8)(0{,}60)(\pot{2.5}{-3})
=\pot{1.2}{-3}\,\si{\weber}.
\]
(b)
\[
|\varepsilon|=\frac{\Delta\Phi}{\Delta t}
=\frac{\pot{1.2}{-3}}{\pot{4}{-2}}=\pot{3}{-2}\,\si{\volt}.
\]
\textbf{Observação importante:} o fluxo depende da área \emph{efetivamente}
imersa no campo, não da área geométrica total da espira. Fora da região,
$B=0$ e a contribuição é nula.\\
\textbf{Se a espira tivesse $N$ espiras}, a f.e.m. seria $N$ vezes maior --- e é
assim que se aumenta a sensibilidade de sensores de fluxo sem aumentar a área.}
\stepcounter{questao}
\resposta{$\Phi_{\max}=\SI{1.2}{\milli\weber}$; $N\approx688$ espiras}
\resolucao{(a) $\Phi_{\max}=B_{\max}A=(1{,}2)(\pot{1}{-3})
=\pot{1.2}{-3}\,\si{\weber}$.\\
(b) Para fluxo senoidal, a f.e.m. eficaz vale
\[
V_{ef}=4{,}44\,N\,f\,\Phi_{\max},
\]
\emph{(o fator $4{,}44=2\pi/\sqrt2$ decorre de derivar uma senoide e converter
para valor eficaz)}. Logo
\[
N=\frac{220}{4{,}44(60)(\pot{1.2}{-3})}=\frac{220}{0{,}3197}\approx688.
\]
\textbf{Compromisso central do projeto:} menos espiras significam $B$ maior (e
risco de saturação); mais espiras significam mais cobre, mais resistência e mais
perdas.\\
\textbf{Por que a frequência importa:} $N\propto1/f$. Em $\SI{400}{\hertz}$
(aviação) ou em dezenas de quilohertz (fontes chaveadas), o mesmo transformador
precisa de muito menos espiras --- e o núcleo pode ser muito menor. É essa a
razão física de fontes chaveadas serem tão mais compactas que as lineares.}
\stepcounter{questao}
\resposta{Nulo, independentemente da posição e da intensidade do ímã}
\resolucao{\textbf{Resultado:} $\oint\vec B\cdot\mathrm{d}\vec A=0$, sempre.\\
\textbf{Justificativa:} as linhas de campo de um ímã \emph{não começam nem
terminam} nele --- elas saem do polo norte, percorrem o espaço externo, entram
pelo sul e \textbf{continuam dentro do material}, fechando-se sobre si mesmas.\\
Toda linha que atravessa a superfície para fora atravessa também para dentro,
em algum outro ponto. O balanço é rigorosamente nulo.\\
\textbf{Contraste com a carga elétrica.} Se a superfície envolvesse uma carga
$+Q$, o fluxo elétrico seria $Q/\varepsilon_0\ne0$ --- porque as linhas
elétricas \emph{nascem} na carga.\\
\textbf{O que isso significa fisicamente:} não existe "carga magnética". Cortar
um ímã ao meio não separa os polos --- produz dois ímãs completos, cada um com
seus dois polos. Não há como isolar um norte de um sul.\\
\textbf{Consequência para a teoria:} $\nabla\cdot\vec B=0$ é uma equação de
Maxwell, e é ela que garante a existência do potencial vetor $\vec A$ com
$\vec B=\nabla\times\vec A$.}
\stepcounter{questao}
\resposta{$\lambda=\SI{6}{\milli\weber}$-espira; $B=\SI{15}{\milli\tesla}$}
\resolucao{(a) Integrando a lei de Faraday:
\[
\int\varepsilon\,\mathrm{d}t=\Delta\lambda
\;\Longrightarrow\;
\lambda_i=\pot{6}{-3}\,\si{\weber}\text{-espira}
\]
(pois o fluxo final é zero).\\
(b)
\[
\Phi=\frac{\lambda}{N}=\frac{\pot{6}{-3}}{500}=\pot{1.2}{-5}\,\si{\weber};
\qquad
B=\frac{\Phi}{A}=\frac{\pot{1.2}{-5}}{\pot{8}{-4}}=\pot{1.5}{-2}\,\si{\tesla}.
\]
(c) \textbf{Vantagem do método integral.} O resultado depende \emph{apenas} da
variação total de fluxo --- \textbf{não} de \emph{como} a bobina foi retirada.
Rápido ou devagar, com solavancos ou suavemente, a integral é a mesma.\\
Isso torna a medição notavelmente robusta: não é preciso controlar a velocidade
nem a trajetória.\\
\textbf{Limitações:} o integrador acumula qualquer \emph{offset} do
amplificador, o que limita o tempo de medição (problema já discutido na
integração de sinais ruidosos). E o método mede apenas \emph{variações} --- para
campos estáticos permanentes, é preciso mover a bobina ou usar sensor Hall.}
\stepcounter{questao}
\resposta{Erro de cerca de $+2{,}5\%$ na estimativa pelo centro}
\resolucao{\textbf{Exato} (com $I$ genérico e $a=\SI{0.10}{\meter}$):
\[
\Phi_{ex}=\frac{\mu_0Ia}{2\pi}\ln\frac{0{,}15}{0{,}05}
=\pot{2}{-8}I\ln3=\pot{2.197}{-8}\,I.
\]
\textbf{Estimativa pelo centro} ($r=\SI{0.10}{\meter}$):
\[
\Phi_{est}=\frac{\mu_0I}{2\pi(0{,}10)}\,a^{2}
=\pot{2}{-7}\frac{I}{0{,}10}(0{,}01)=\pot{2}{-8}\,I.
\]
\[
\text{erro}=\frac{2{,}00-2{,}197}{2{,}197}=-9{,}0\%.
\]
\textbf{A estimativa pelo centro \emph{subestima}} o fluxo --- porque
$B\propto1/r$ é uma função \textbf{convexa}, e a média de uma função convexa é
maior que a função da média (desigualdade de Jensen).\\
\textbf{Quando a aproximação é aceitável:} quando a variação de $B$ ao longo da
espira for pequena. Aqui, $B$ varia por um fator $3$ --- muito. Se a espira
estivesse entre $95$ e $\SI{105}{\centi\meter}$, a variação seria de apenas
$10\%$ e o erro cairia para frações de por cento.\\
\textbf{Critério prático:} use campo uniforme se
$B_{\max}/B_{\min}<1{,}1$; caso contrário, integre.}
\stepcounter{questao}
\resposta{Linhas fechadas $\Rightarrow$ fluxo líquido nulo $\Rightarrow$ o fluxo
se conserva ao longo de um circuito magnético}
\resolucao{(a) \textbf{Demonstração.} As linhas de $\vec B$ são curvas
\emph{fechadas} --- fato experimental, equivalente à inexistência de monopolos
magnéticos. Uma curva fechada que penetre uma superfície fechada precisa
atravessá-la um número \emph{par} de vezes, alternando entrada e saída.\\
Cada entrada contribui com fluxo negativo e cada saída com positivo, de módulos
iguais. A soma se cancela:
\[
\oint\vec B\cdot\mathrm{d}\vec A=0.
\]
Pelo teorema da divergência, isso equivale a $\nabla\cdot\vec B=0$ em todo
ponto.\\
(b) \textbf{Conservação do fluxo.} Considere uma junção em um núcleo magnético,
onde um trecho de seção $A_1$ encontra outro de seção $A_2$. Envolvendo a junção
com uma superfície fechada:
\[
-\Phi_1+\Phi_2=0
\;\Longrightarrow\;
\Phi_1=\Phi_2.
\]
\textbf{O fluxo se conserva} --- exatamente como a corrente em um nó. Se o
núcleo se ramificar em dois caminhos:
\[
\Phi_{total}=\Phi_a+\Phi_b,
\]
que é a \textbf{lei dos nós magnética}, análoga à LKC.\\
(c) \textbf{Comparação com Gauss elétrica.}
\begin{center}
\begin{tabular}{ll}
\hline
Elétrico & Magnético\\
\hline
$\oint\vec E\cdot\mathrm{d}\vec A=Q/\varepsilon_0$
 & $\oint\vec B\cdot\mathrm{d}\vec A=0$\\
Linhas nascem e morrem em cargas & Linhas são fechadas\\
Existem monopolos (cargas) & Não existem monopolos\\
$\nabla\cdot\vec E=\rho/\varepsilon_0$ & $\nabla\cdot\vec B=0$\\
\hline
\end{tabular}
\end{center}
\textbf{Consequência profunda:} a ausência de monopolos permite escrever
$\vec B=\nabla\times\vec A$, base de toda a formulação moderna do
eletromagnetismo. Se monopolos existissem, essa construção seria impossível.}
\stepcounter{questao}
\resposta{Duas superfícies de mesma borda formam superfície fechada, cujo fluxo
é nulo}
\resolucao{(a) \textbf{Demonstração.} Sejam $S_1$ e $S_2$ duas superfícies
apoiadas na mesma curva $C$. Juntas, elas delimitam um \textbf{volume fechado}.
Aplicando o resultado anterior:
\[
\oint\vec B\cdot\mathrm{d}\vec A
=\Phi_{S_2}^{(saindo)}-\Phi_{S_1}^{(saindo)}=0
\;\Longrightarrow\;
\Phi_{S_1}=\Phi_{S_2}.
\]
(b) \textbf{Importância para Faraday.} A lei
$\varepsilon=-\dfrac{\mathrm{d}\Phi}{\mathrm{d}t}$ só faz sentido se
$\Phi$ estiver \emph{bem definido} --- ou seja, se não depender da superfície
escolhida.\\
Se dependesse, a f.e.m. de um mesmo circuito teria valores diferentes conforme a
superfície imaginada --- absurdo, pois a f.e.m. é medível.\\
\textbf{Liberdade prática:} escolha a superfície que simplificar a integral. Ao
calcular o fluxo através de uma espira que envolve um solenoide, por exemplo,
convém uma superfície que corte o solenoide perpendicularmente.\\
(c) \textbf{Orientação das normais.} A igualdade exige que as normais sejam
orientadas \emph{coerentemente} com o sentido de percurso da borda, pela regra
da mão direita.\\
Se uma delas for invertida, aparece um sinal negativo --- e a conclusão errônea
de que os fluxos são opostos. \textbf{Fixe o sentido de percurso da curva
primeiro}, e derive as normais dele.}
\stepcounter{questao}
\resposta{$\mathcal{F}=\Phi\mathcal{R}$, com
$\mathcal{R}=\ell/\mu A$; a analogia falha por não linearidade, dispersão e
ausência de dissipação}
\resolucao{(a)
\begin{center}
\begin{tabular}{ll}
\hline
Circuito elétrico & Circuito magnético\\
\hline
F.e.m. $\varepsilon$ & Força magnetomotriz $\mathcal{F}=NI$\\
Corrente $I$ & Fluxo $\Phi$\\
Resistência $R=\ell/\sigma A$ & Relutância
 $\mathcal{R}=\ell/\mu A$\\
$\varepsilon=RI$ & $\mathcal{F}=\mathcal{R}\Phi$\\
Condutividade $\sigma$ & Permeabilidade $\mu$\\
LKC nos nós & Conservação do fluxo\\
Resistências em série somam & Relutâncias em série somam\\
\hline
\end{tabular}
\end{center}
(b) \textbf{Dedução da relutância.} Da lei de Ampère aplicada ao caminho médio
do núcleo, $NI=H\ell$. Com $B=\mu H$ e $\Phi=BA$:
\[
NI=\frac{B\ell}{\mu}=\frac{\Phi\,\ell}{\mu A}
\;\Longrightarrow\;
\mathcal{R}=\frac{\ell}{\mu A}.
\]
(c) \textbf{Três diferenças que quebram a analogia.}
\begin{enumerate}
\item \textbf{Não linearidade.} $\sigma$ é constante em metais, mas
$\mu$ varia enormemente com $B$ --- e desaba na saturação. A "resistência
magnética" depende da própria corrente que a atravessa.
\item \textbf{Dispersão.} A corrente elétrica fica confinada no condutor porque
a razão de condutividades metal/ar é de $10^{20}$. Já a razão de
permeabilidades ferro/ar é de apenas $10^{3}$ --- e por isso parte do fluxo
\emph{escapa} do núcleo. Não existe "isolante magnético".
\item \textbf{Ausência de dissipação.} A corrente dissipa $RI^{2}$ em calor.
O fluxo \emph{não dissipa} nada --- manter $\Phi$ constante não custa energia
(as perdas por histerese e Foucault ocorrem apenas quando o fluxo \emph{varia},
e por outros mecanismos).
\end{enumerate}
\textbf{Conclusão:} a analogia é excelente como \emph{ferramenta de cálculo}
abaixo da saturação, e enganosa como \emph{modelo físico}. Use-a para
dimensionar; não a use para raciocinar sobre energia.}
\stepcounter{questao}
\resposta{$W=\SI{1.25}{\milli\joule}$; $N=17$ espiras;
$g\approx\SI{0.33}{\milli\meter}$; $92\%$ da energia no entreferro}
\resolucao{(a)
\[
W=\tfrac12LI^{2}=\tfrac12(\pot{1}{-4})(25)=\SI{1.25}{\milli\joule}.
\]
O fluxo concatenado é $\lambda=LI=\pot{5}{-4}\,\si{\weber}$-espira, e o fluxo
máximo por espira é limitado pela saturação:
\[
\Phi_{\max}=B_{sat}A=(0{,}3)(\pot{1}{-4})=\SI{30}{\micro\weber}.
\]
\[
N\ge\frac{\lambda}{\Phi_{\max}}=\frac{\pot{5}{-4}}{\pot{3}{-5}}=16{,}7
\;\Longrightarrow\;
N=17\ \text{espiras}.
\]
\emph{(Verificação: com $N=17$, $B=\lambda/(NA)=\SI{0.294}{\tesla}$ ---
logo abaixo do limite.)}\\
(b) De $L=N^{2}/\mathcal{R}$:
\[
\mathcal{R}_{tot}=\frac{17^{2}}{\pot{1}{-4}}=\pot{2.89}{6}\,\si{\per\henry}.
\]
A parcela do ferrite é
$\mathcal{R}_{fe}=\dfrac{0{,}06}{(2000)\mu_0(\pot{1}{-4})}
=\pot{2.39}{5}$, restando para o entreferro
$\mathcal{R}_g=\pot{2.65}{6}$:
\[
g=\mathcal{R}_g\,\mu_0A=(\pot{2.65}{6})(4\pi\times10^{-7})(\pot{1}{-4})
=\SI{0.33}{\milli\meter}.
\]
\textbf{Note que o entreferro responde por $92\%$ da relutância total}, embora
seja $180$ vezes mais curto que o caminho de ferrite.\\
(c) Como $B$ é o mesmo nos dois meios e $u=B^{2}/2\mu$:
\[
\frac{W_g}{W_{fe}}=\frac{g/\mu_0}{\ell_{fe}/(\mu_r\mu_0)}
=\frac{g\,\mu_r}{\ell_{fe}}
=\frac{(\pot{3.33}{-4})(2000)}{0{,}06}=11{,}1,
\]
ou seja, $\SI{91.7}{\percent}$ da energia está no \textbf{entreferro}.\\
\textbf{Por que isso importa.} A densidade de energia
$u=B^{2}/2\mu$ é \emph{inversamente} proporcional a $\mu$ --- para o mesmo
$B$, o ar armazena duas mil vezes mais energia por volume que o ferrite.\\
\textbf{O entreferro é o reservatório; o núcleo apenas conduz o fluxo até
ele.} Um indutor sem entreferro satura com corrente mínima e armazena quase
nada --- ao contrário de um transformador, que \emph{não deve} armazenar
energia e por isso dispensa entreferro.}
\stepcounter{questao}
\resposta{$k=M/\sqrt{L_1L_2}$; a dispersão aparece como indutância série e
causa queda de tensão com a carga}
\resolucao{(a) Por definição,
\[
k=\frac{M}{\sqrt{L_1L_2}},
\qquad 0\le k\le1,
\]
resultado já demonstrado a partir da não negatividade da energia do par de
bobinas.\\
Decompondo cada indutância em parcela \emph{magnetizante} (fluxo compartilhado)
e parcela de \emph{dispersão}:
\[
L_1=L_{m1}+L_{d1},
\qquad
k^{2}=\frac{L_{m1}L_{m2}}{L_1L_2}.
\]
(b) \textbf{Efeito na saída.} A indutância de dispersão comporta-se como uma
\emph{impedância em série} com o enrolamento. Sob carga, ela produz queda de
tensão proporcional à corrente:
\begin{itemize}
\item \textbf{Em vazio}, a relação de tensões é praticamente $N_2/N_1$ --- a
dispersão não se manifesta.
\item \textbf{Com carga}, a tensão de saída \emph{cai}, e a regulação piora
proporcionalmente à dispersão.
\end{itemize}
Um transformador com $k=0{,}98$ tem cerca de $2\%$ de dispersão e regulação
tipicamente na faixa de $2$ a $5\%$.\\
(c) \textbf{Quando a dispersão é desejável.}
\begin{itemize}
\item \textbf{Transformadores de solda:} a dispersão elevada (às vezes
$k\approx0{,}8$) limita naturalmente a corrente de arco, dispensando reator
externo. A característica "mole" é justamente o que se quer.
\item \textbf{Transformadores de forno e de descarga:} mesma lógica --- a carga
é essencialmente um curto controlado.
\item \textbf{Limitação de curto-circuito:} em transformadores de distribuição,
uma dispersão mínima é \emph{exigida} por norma, para limitar a corrente de
falta a valores que a proteção consiga interromper.
\end{itemize}
\textbf{Conclusão:} $k\to1$ não é universalmente melhor. A dispersão é um
parâmetro de projeto, e não apenas um defeito a minimizar.}
\stepcounter{questao}
\resposta{$u=B^{2}/2\mu$; praticamente toda a energia fica no entreferro}
\resolucao{(a) Com $L=N^{2}/\mathcal{R}$ e $\lambda=N\Phi=LI$:
\[
W=\frac12LI^{2}=\frac{\lambda^{2}}{2L}
=\frac{(N\Phi)^{2}}{2N^{2}/\mathcal{R}}
=\frac{1}{2}\mathcal{R}\,\Phi^{2}.
\]
\textbf{Analogia:} é a forma magnética de $W=\frac12 CU^{2}$ escrita como
$Q^{2}/2C$ --- o fluxo faz o papel da carga e a relutância, o do inverso da
capacitância.\\
(b) Escrevendo $\mathcal{R}=\ell/\mu A$ e $\Phi=BA$:
\[
W=\frac{1}{2}\frac{\ell}{\mu A}(BA)^{2}=\frac{B^{2}}{2\mu}(A\ell)
\;\Longrightarrow\;
u=\frac{B^{2}}{2\mu}.
\]
\textbf{Ponto decisivo:} a densidade de energia é \emph{inversamente}
proporcional a $\mu$. Para o mesmo $B$, o ar armazena $\mu_r$ vezes \emph{mais}
energia por volume que o ferro.\\
(c) \textbf{Fração no entreferro.} Como $B$ é o mesmo nos dois meios (fluxo
conservado e mesma seção), a razão das energias é a razão dos produtos
$\ell/\mu$:
\[
\frac{W_g}{W_{fe}}
=\frac{g/\mu_0}{\ell_{fe}/(\mu_r\mu_0)}
=\frac{g\,\mu_r}{\ell_{fe}}
=\frac{(\pot{0.56}{-3})(2000)}{0{,}50}=2{,}24.
\]
Ou seja, cerca de $69\%$ da energia está no entreferro, contra $31\%$ no
ferro --- e a proporção cresce rapidamente com $g$.\\
\textbf{Por isso o entreferro é o elemento de projeto de um indutor de
potência:} ele é o \emph{reservatório}. O ferro apenas conduz o fluxo até ele.}
\stepcounter{questao}
\resposta{$\mathcal{R}=\displaystyle\int\frac{\mathrm{d}\ell}{\mu A(\ell)}$;
dimensiona-se pela \emph{menor} seção}
\resolucao{(a) Cada trecho infinitesimal contribui como uma relutância em
série:
\[
\mathcal{R}=\int_0^{L}\frac{\mathrm{d}\ell}{\mu\,A(\ell)}.
\]
Para trechos de seção constante, isso vira a soma
$\sum\ell_i/(\mu A_i)$.\\
(b) \textbf{Critério de dimensionamento.} Sem dispersão, o fluxo é o
\emph{mesmo} em toda a extensão. Logo
\[
B(\ell)=\frac{\Phi}{A(\ell)},
\]
e $B$ é máximo onde $A$ é \textbf{mínimo}. A condição de não saturação é
\[
\Phi\le B_{sat}\,A_{\min}.
\]
\textbf{Toda a capacidade do núcleo é ditada pela sua menor seção} --- alargar
o resto não ajuda em nada.\\
(c) \textbf{Estreitamento localizado.} Suponha um núcleo de
$\SI{6}{\centi\meter\squared}$ com um trecho curto de
$\SI{2}{\centi\meter\squared}$:
\begin{itemize}
\item \textbf{Efeito na relutância:} se o trecho estreito for curto, sua
contribuição a $\mathcal{R}$ é pequena --- o fluxo total quase não muda.
\item \textbf{Efeito na saturação:} ali $B$ é \emph{três vezes} maior, e o
trecho satura com um terço do fluxo que o resto suportaria.
\item \textbf{Efeito local:} saturado, aquele trecho aquece por perdas
concentradas e sua permeabilidade cai --- criando um entreferro
\emph{efetivo} não projetado.
\end{itemize}
\textbf{Onde isso ocorre na prática:} furos de fixação, cantos internos de
núcleos em E, e emendas de lâminas. São defeitos que passam despercebidos no
cálculo de relutância (pequeno efeito) mas dominam o comportamento na saturação
(efeito grande).\\
\textbf{Regra:} calcule $\mathcal{R}$ pelo caminho todo, mas verifique
$B$ na \emph{seção crítica}.}

\gabarito[12.3 Fluxo magnético]

\subsection{Força magnética}
\stepcounter{questao}
\resposta{(b)}
\resolucao{Como $F = qvB\sin\theta$, a força se anula quando $\sin\theta = 0$, ou
seja, quando $\vec{v}$ é paralela (ou antiparalela) a $\vec{B}$. É máxima quando
$\vec{v}\perp\vec{B}$ ($\theta = \SI{90}{\degree}$).}
\stepcounter{questao}
\resposta{$F = \SI{0.3}{\newton}$}
\resolucao{$F = B\,I\,\ell = 0{,}5\cdot2\cdot0{,}3 = \SI{0.3}{\newton}$
(com $\sin\theta = 1$, pois são perpendiculares).}
\stepcounter{questao}
\resposta{(a)}
\resolucao{Pela regra da mão direita para $\vec F = q\vec v\times\vec B$: dedos
apontando no sentido de $\vec v$ (direita), curvando-se para $\vec B$ (para dentro
da página), o polegar aponta para \emph{cima}. Como a carga é positiva, a força
tem esse mesmo sentido. (Se fosse negativa, seria para baixo.)}
\stepcounter{questao}
\resposta{$F = \SI{3.2e-14}{\newton}$; trajetória circular}
\resolucao{$F = q\,v\,B = \pot{1.6}{-19}\cdot\pot{1}{6}\cdot0{,}2
= \pot{3.2}{-14}\,\si{\newton}$.\\
Como a força é sempre \emph{perpendicular} à velocidade, ela não altera o módulo
de $v$ (não realiza trabalho), apenas curva a trajetória: o resultado é um
\textbf{movimento circular uniforme}. É o princípio dos cíclotrons e dos
espectrômetros de massa.}
\stepcounter{questao}
\resposta{Ver comentário}
\resolucao{Em um campo uniforme, os dois lados opostos da espira conduzem corrente
em sentidos \emph{contrários}, sofrendo forças magnéticas iguais e opostas
($\vec F = B I \ell$). Essas forças formam um \textbf{binário} (conjugado): a
resultante é nula (não há translação), mas o \emph{torque} não é --- a espira
gira. Esse é exatamente o princípio do \textbf{motor de corrente contínua}: o
torque $\tau = N B I A \sin\theta$ faz o rotor girar, e o comutador inverte a
corrente a cada meia-volta para manter a rotação no mesmo sentido.}
\stepcounter{questao}
\resposta{$r \approx \SI{4.2}{\centi\meter}$}
\resolucao{A força magnética atua como resultante centrípeta:
$qvB = \dfrac{mv^2}{r}$. Isolando $r$:
\[
r=\frac{mv}{qB}
=\frac{(\pot{1.67}{-27})(\pot{2}{6})}{(\pot{1.6}{-19})(0{,}5)}
\approx\pot{4.2}{-2}\,\si{\meter}=\SI{4.2}{\centi\meter}.
\]
O raio cresce com a quantidade de movimento ($mv$) e diminui com a carga e o
campo. Medindo $r$ conhece-se a razão carga/massa da partícula --- é assim que
funciona o espectrômetro de massa.}
\stepcounter{questao}
\resposta{(a) $qE = qvB \Rightarrow v = E/B$; (c) $v = \SI{2e6}{\meter\per\second}$;
(b) e (d) ver comentário}
\resolucao{\textbf{(a)} Sobre a partícula atuam duas forças perpendiculares à
trajetória e \emph{opostas} entre si:
\[
F_e = qE \quad(\text{para baixo, se } q>0),
\qquad
F_m = qvB \quad(\text{para cima, pela regra da mão direita}).
\]
Para não haver desvio, a resultante deve ser nula:
\[
qE = qvB
\;\Longrightarrow\;
\boxed{\,v=\frac{E}{B}\,}
\]
\textbf{(b)} A carga $q$ \textbf{cancela-se} nos dois lados da igualdade, e a
massa sequer aparece --- não há aceleração, logo a segunda lei de Newton não é
invocada. Consequência notável: o seletor filtra por \emph{velocidade}, e apenas
por ela. Elétrons, prótons e íons pesados com a mesma $v$ atravessam igualmente.
Partículas mais rápidas são desviadas pela força magnética (que cresce com $v$);
mais lentas, pela força elétrica.

\textbf{(c)} $v = \dfrac{E}{B} = \dfrac{\pot{3}{5}}{0{,}15}
= \pot{2}{6}\,\si{\meter\per\second}$.

\textbf{(d) O espectrômetro de massa.} Com a velocidade agora \emph{conhecida e
uniforme}, a partícula entra em uma região de campo $B'$ puro, onde a força
magnética atua como centrípeta:
\[
qvB' = \frac{mv^{2}}{r}
\;\Longrightarrow\;
\frac{q}{m}=\frac{v}{B'r}=\frac{E}{B\,B'\,r}.
\]
Medindo o raio $r$ da trajetória (pela marca deixada no detector) e conhecendo
$E$, $B$ e $B'$, obtém-se $q/m$ diretamente. Íons de massas diferentes descrevem
raios diferentes e se separam espacialmente --- daí o nome
\emph{espectrômetro}.\\
\textbf{Impacto histórico:} foi com esse arranjo que Thomson mediu $q/m$ do
elétron (1897) e que Aston descobriu os \textbf{isótopos} (1919), rendendo-lhe o
Nobel. Hoje é ferramenta padrão em química analítica, datação e antidoping.}
\stepcounter{questao}
\resposta{$F=\SI{0.2}{\milli\newton}$}
\resolucao{
\[
F=qvB\sin\theta=(\pot{2}{-6})(500)(0{,}2)(1)=\pot{2}{-4}\,\si{\newton}.
\]
\textbf{Direção:} perpendicular \emph{simultaneamente} a $\vec v$ e a
$\vec B$, dada pela regra da mão direita aplicada a
$\vec v\times\vec B$ (e invertida, se a carga for negativa).}
\stepcounter{questao}
\resposta{$F=\SI{0.3}{\newton}$}
\resolucao{
\[
F=BIL\sin\theta=(0{,}25)(3)(0{,}40)(1)=\SI{0.3}{\newton}.
\]
\textbf{Origem da fórmula:} ela é a soma das forças sobre todos os portadores do
condutor. Escrevendo $I=nqv_dA$ e $L$ o comprimento, o produto
$nqv_dA\cdot L$ reproduz $qv_d$ somado sobre $nAL$ portadores.\\
A força é transmitida ao \emph{fio} pelas colisões dos portadores com a rede
cristalina --- não são os elétrons que "empurram" o fio diretamente.}
\stepcounter{questao}
\resposta{$F=\SI{0.1}{\milli\newton}$}
\resolucao{
\[
F=qvB\sin\ang{30}=(\pot{2}{-4})(0{,}5)=\pot{1}{-4}\,\si{\newton}.
\]
\textbf{Metade} do valor perpendicular.\\
\textbf{Casos-limite:} $\theta=\ang{90}$ dá força máxima; $\theta=0$ ou
$\ang{180}$ (velocidade \emph{paralela} ao campo) dá força \textbf{nula} --- a
carga segue em linha reta, sem sentir o campo.}
\stepcounter{questao}
\resposta{(a)}
\resolucao{Pela regra da mão direita para $\vec v\times\vec B$: dedos apontando
no sentido de $\vec v$ (direita), curvando-se para $\vec B$ (cima) --- o polegar
aponta para \textbf{dentro} do plano.\\
\textbf{Verificação:} a força deve ser perpendicular a ambos. "Direita" e
"cima" definem o plano da folha, logo a força só pode ser perpendicular a ele
--- o que elimina (c) e (d) imediatamente.\\
\textbf{Se a carga fosse negativa}, a força seria para fora do plano.}
\stepcounter{questao}
\resposta{$\si{\tesla}=\dfrac{\si{\newton}}{\si{\coulomb}\cdot
\si{\meter\per\second}}$}
\resolucao{
\[
B=\frac{F}{qv}
\;\Longrightarrow\;
[\si{\tesla}]=\frac{\si{\newton}}{\si{\coulomb}\cdot\si{\meter\per\second}}
=\frac{\si{\newton}\cdot\si{\second}}{\si{\coulomb}\cdot\si{\meter}}
=\frac{\si{\newton}}{\si{\ampere}\cdot\si{\meter}}.
\]
A última forma, $\si{\newton\per\ampere\per\meter}$, decorre diretamente de
$F=BIL$ --- e é a mais prática para conferir contas com condutores.\\
\textbf{O tesla é uma unidade grande:} $\SI{1}{\tesla}$ exerce
$\SI{1}{\newton}$ sobre um metro de fio conduzindo $\SI{1}{\ampere}$.}
\stepcounter{questao}
\resposta{$F=\SI{20}{\milli\newton}$}
\resolucao{
\[
F=BIL\sin\ang{30}=(0{,}2)(2)(0{,}10)(0{,}5)=\pot{2}{-2}\,\si{\newton}.
\]
\textbf{Interpretação do seno:} apenas a componente de $\vec\ell$
\emph{perpendicular} a $\vec B$ contribui. Equivalentemente, o comprimento
\emph{efetivo} é $L\sin\theta=\SI{5}{\centi\meter}$.\\
\textbf{Direção da força:} perpendicular ao plano definido pelo fio e pelo
campo --- e \emph{não} perpendicular apenas ao fio.}
\stepcounter{questao}
\resposta{$r=\SI{4.18}{\centi\meter}$; $T=\SI{131}{\nano\second}$}
\resolucao{(a) A força magnética é centrípeta:
\[
qvB=\frac{mv^{2}}{r}
\;\Longrightarrow\;
r=\frac{mv}{qB}=\frac{(\pot{1.67}{-27})(\pot{2}{6})}{(\pot{1.6}{-19})(0{,}5)}
=\SI{0.0418}{\meter}.
\]
(b)
\[
T=\frac{2\pi r}{v}=\frac{2\pi m}{qB}
=\frac{2\pi(\pot{1.67}{-27})}{(\pot{1.6}{-19})(0{,}5)}
=\pot{1.31}{-7}\,\si{\second}.
\]
\textbf{Fato notável:} o período \textbf{não depende da velocidade} nem do raio.
Partículas rápidas descrevem círculos maiores, mas levam o mesmo tempo para
completá-los.\\
É esse fato que torna possível o \textbf{cíclotron}: a frequência de aceleração
pode ser fixa, embora a partícula ganhe energia a cada volta.}
\stepcounter{questao}
\resposta{$r_p/r_e=1836$; $T_p/T_e=1836$}
\resolucao{Como $r=mv/(qB)$ e $T=2\pi m/(qB)$, ambos são
\emph{proporcionais à massa} (as cargas têm o mesmo módulo):
\[
\frac{r_p}{r_e}=\frac{T_p}{T_e}=\frac{m_p}{m_e}
=\frac{\pot{1.67}{-27}}{\pot{9.11}{-31}}=1836.
\]
\textbf{Sentidos opostos:} tendo cargas de sinais contrários, as duas partículas
curvam-se para lados opostos --- é assim que se identifica o sinal da carga em
câmaras de detecção.\\
\textbf{Consequência experimental:} para conter elétrons e prótons no mesmo
aparelho, o raio do próton é quase dois mil vezes maior. Aceleradores de prótons
são, por isso, muito maiores que os de elétrons para a mesma energia.}
\stepcounter{questao}
\resposta{$F/L=\SI{0.4}{\milli\newton\per\meter}$}
\resolucao{Um fio cria campo $B=\dfrac{\mu_0I_1}{2\pi d}$ na posição do outro,
que sofre $F=BI_2L$:
\[
\frac{F}{L}=\frac{\mu_0I_1I_2}{2\pi d}
=\pot{2}{-7}\frac{(10)(10)}{0{,}05}=\pot{4}{-4}\,\si{\newton\per\meter}.
\]
\textbf{Sentido:} correntes de \emph{mesmo} sentido \textbf{atraem-se};
de sentidos opostos, \textbf{repelem-se}.\\
\textbf{Cuidado com a intuição elétrica:} cargas de mesmo sinal se repelem, mas
correntes de mesmo sentido se atraem. As duas situações não são análogas.}
\stepcounter{questao}
\resposta{$m=\SI{1}{\ampere\meter\squared}$; $\tau=\SI{0.3}{\newton\meter}$}
\resolucao{(a) $A=\SI{100}{\centi\meter\squared}=\pot{1}{-2}\,\si{\meter\squared}$:
\[
m=NIA=50(2)(\pot{1}{-2})=\SI{1}{\ampere\meter\squared}.
\]
(b)
\[
\tau=mB\sin\theta
\;\Longrightarrow\;
\tau_{\max}=mB=(1)(0{,}3)=\SI{0.3}{\newton\meter}.
\]
\textbf{Quando ocorre o máximo:} com o plano da espira \emph{paralelo} ao campo,
isto é, $\vec m$ perpendicular a $\vec B$. O torque é \textbf{nulo} quando
$\vec m$ está alinhado com $\vec B$ --- posição de equilíbrio.\\
\textbf{Analogia útil:} a espira comporta-se como uma agulha de bússola, e
$\vec m$ faz o papel do "norte" dela.}
\stepcounter{questao}
\resposta{Ela é sempre perpendicular a $\vec v$, logo não realiza trabalho}
\resolucao{A potência entregue por uma força é $P=\vec F\cdot\vec v$. Como
\[
\vec F=q\,\vec v\times\vec B
\;\Longrightarrow\;
\vec F\perp\vec v
\;\Longrightarrow\;
\vec F\cdot\vec v=0.
\]
Sem potência, não há variação de energia cinética --- e portanto o \emph{módulo}
da velocidade permanece constante.\\
\textbf{O que a força magnética faz:} muda a \emph{direção} do movimento, e só
isso. Ela curva a trajetória sem acelerar nem frear.\\
\textbf{Consequência prática:} um campo magnético sozinho jamais acelera
partículas. Aceleradores usam campos \emph{elétricos} para ganhar energia e
campos magnéticos apenas para \emph{guiar} o feixe.}
\stepcounter{questao}
\resposta{$F=\SI{8e-17}{\newton}$ (só elétrica);
$a=\pot{8.8}{13}\,\si{\meter\per\second\squared}$}
\resolucao{(a) A força total é a \textbf{força de Lorentz}:
\[
\vec F=q\left(\vec E+\vec v\times\vec B\right).
\]
Com $\vec v=\vec 0$, a parcela magnética \textbf{se anula} --- resta apenas a
elétrica:
\[
F=eE=(\pot{1.6}{-19})(500)=\pot{8}{-17}\,\si{\newton};
\qquad
a=\frac{F}{m_e}=\pot{8.8}{13}\,\si{\meter\per\second\squared}.
\]
(b) \textbf{Movimento subsequente.} Assim que o elétron adquire velocidade, a
parcela magnética \emph{deixa de ser nula} e passa a encurvar a trajetória. A
força elétrica continua acelerando; a magnética continua desviando.\\
O resultado é uma trajetória \textbf{cicloidal}, com deriva média
\[
v_d=\frac{E}{B}=\frac{500}{0{,}02}=\pot{2.5}{4}\,\si{\meter\per\second},
\]
perpendicular a ambos os campos.\\
\textbf{Ponto essencial:} a força magnética \emph{depende da velocidade} --- ela
não age sobre cargas paradas. Já a elétrica age sempre. É essa assimetria que
faz o movimento começar de um jeito e evoluir para outro.}
\stepcounter{questao}
\resposta{$F=2RIB$, como se o fio fosse retilíneo de comprimento $2R$}
\resolucao{Para um fio de forma qualquer em campo \textbf{uniforme},
\[
\vec F=I\left(\int\mathrm{d}\vec\ell\right)\times\vec B
=I\,\vec L_{ef}\times\vec B,
\]
onde $\vec L_{ef}$ é o vetor que liga o \textbf{início ao fim} do fio --- a
\emph{corda}, não o comprimento percorrido.\\
Para a semicircunferência, a corda vale $2R$:
\[
F=2RIB.
\]
\textbf{Poder do resultado:} a forma do fio é irrelevante em campo uniforme.
Um fio retorcido, ondulado ou em zigue-zague sofre a mesma força que o segmento
reto entre suas extremidades.\\
\textbf{Corolário importante:} em uma \emph{espira fechada}, a corda é nula ---
logo a força resultante sobre qualquer espira em campo uniforme é
\textbf{zero}. Só há torque.}
\stepcounter{questao}
\resposta{$I=\SI{0.245}{\ampere}$}
\resolucao{Equilíbrio entre peso e força magnética:
\[
BIL=mg
\;\Longrightarrow\;
I=\frac{mg}{BL}=\frac{(\pot{2}{-3})(9{,}8)}{(0{,}4)(0{,}20)}
=\SI{0.245}{\ampere}.
\]
\textbf{Uso do arranjo:} conhecendo $m$, $L$ e $I$, obtém-se $B$ com boa
exatidão --- é o método clássico de calibração absoluta de campo magnético, sem
depender de sensores.\\
\textbf{Inversamente}, conhecendo $B$, mede-se corrente. Foi assim que o ampère
foi definido durante décadas: pela força entre condutores.}
\stepcounter{questao}
\resposta{$v=\pot{2}{5}\,\si{\meter\per\second}$}
\resolucao{Para não haver desvio, as forças elétrica e magnética devem se
cancelar:
\[
qE=qvB
\;\Longrightarrow\;
v=\frac{E}{B}=\frac{\pot{2}{4}}{0{,}1}=\pot{2}{5}\,\si{\meter\per\second}.
\]
\textbf{Duas propriedades notáveis:}
\begin{itemize}
\item A condição \textbf{não depende da carga} --- $q$ cancela. Íons e elétrons
com a mesma velocidade passam igualmente.
\item Também \textbf{não depende da massa}. É pura seleção de \emph{velocidade}.
\end{itemize}
\textbf{Consequência:} partículas mais rápidas são desviadas pela força
magnética (maior), e mais lentas pela elétrica. Apenas a velocidade $E/B$
atravessa em linha reta.}
\stepcounter{questao}
\resposta{$\vec F=\vec 0$; $\tau=\SI{10}{\milli\newton\meter}$}
\resolucao{(a) \textbf{Força resultante nula.} Em campo uniforme, os quatro
lados sofrem forças que se cancelam duas a duas --- lados opostos conduzem em
sentidos contrários.\\
(b) \textbf{Mas o torque não é nulo.} As forças nos dois lados perpendiculares
ao campo formam um \emph{binário}:
\[
\tau=mB\sin\ang{90}=IAB=(5)(\pot{1}{-2})(0{,}2)
=\pot{1}{-2}\,\si{\newton\meter}.
\]
\textbf{Situação de torque máximo:} o campo está no plano da espira, ou seja,
$\vec m$ (perpendicular ao plano) está a $\ang{90}$ de $\vec B$.\\
\textbf{Este é o princípio do motor elétrico} --- e também sua limitação: o
torque cai a zero quando a espira gira até o alinhamento, o que exige comutador
para inverter a corrente e manter a rotação.}
\stepcounter{questao}
\resposta{$r=\dfrac{mv}{qB}$; $T=\dfrac{2\pi m}{qB}$;
$\omega=\dfrac{qB}{m}$}
\resolucao{(a) A força magnética é perpendicular a $\vec v$ e de módulo
constante --- exatamente as condições de um movimento \textbf{circular
uniforme}. Igualando à resultante centrípeta:
\[
qvB=\frac{mv^{2}}{r}
\;\Longrightarrow\;
r=\frac{mv}{qB}.
\]
\[
T=\frac{2\pi r}{v}=\frac{2\pi m}{qB}.
\]
(b)
\[
\omega=\frac{2\pi}{T}=\frac{qB}{m},
\]
a \textbf{frequência angular de cíclotron}.\\
(c) \textbf{Por que $T$ independe de $v$.} Dobrar a velocidade dobra a força
magnética \emph{e} dobra o raio necessário --- os dois efeitos se compensam
exatamente. A partícula percorre um círculo duas vezes maior a velocidade duas
vezes maior, no mesmo tempo.\\
Algebricamente, $v$ aparece em $r$ e é cancelado por $v$ em $T=2\pi r/v$.\\
\textbf{Consequência:} o cíclotron pode operar com frequência \emph{fixa}.
\emph{(Isso vale enquanto $v\ll c$; em regime relativístico a massa efetiva
cresce, $T$ aumenta e é preciso modular a frequência --- o síncrotron.)}}
\stepcounter{questao}
\resposta{Hélice com $r=\SI{5.2}{\centi\meter}$ e passo de
$\SI{0.57}{\meter}$}
\resolucao{(a) \textbf{Decomposição:}
\[
v_\perp=v\sin\ang{30}=\pot{5}{5}\,\si{\meter\per\second};
\qquad
v_\parallel=v\cos\ang{30}=\pot{8.66}{5}\,\si{\meter\per\second}.
\]
A componente \textbf{perpendicular} sofre força e produz movimento circular. A
componente \textbf{paralela} não sofre força alguma ($\vec v\times\vec B=0$
para ela) e produz movimento retilíneo uniforme.\\
A composição é uma \textbf{hélice} de eixo paralelo a $\vec B$.\\
(b)
\[
r=\frac{mv_\perp}{qB}
=\frac{(\pot{1.67}{-27})(\pot{5}{5})}{(\pot{1.6}{-19})(0{,}1)}
=\SI{0.052}{\meter};
\]
\[
T=\frac{2\pi m}{qB}=\pot{6.56}{-7}\,\si{\second};
\qquad
\text{passo}=v_\parallel T=(\pot{8.66}{5})(\pot{6.56}{-7})=\SI{0.568}{\meter}.
\]
\textbf{Aplicação:} é essa trajetória que confina partículas nos cinturões de
Van Allen e em espelhos magnéticos --- as partículas espiralam ao longo das
linhas de campo, e o passo diminui onde o campo se intensifica, até a reversão
do movimento.}
\stepcounter{questao}
\resposta{$r_{20}=\SI{20.8}{\centi\meter}$, $r_{22}=\SI{22.8}{\centi\meter}$;
separação de $\SI{4.2}{\centi\meter}$}
\resolucao{(a) Com $r=mv/(qB)$ e velocidade já selecionada:
\[
r_{20}=\frac{20(\pot{1.66}{-27})(\pot{1}{5})}{(\pot{1.6}{-19})(0{,}1)}
=\SI{0.2075}{\meter};
\qquad
r_{22}=\SI{0.2283}{\meter}.
\]
(b) Após meia volta, cada íon atinge o detector a uma distância $2r$ do ponto
de entrada:
\[
\Delta=2\left(r_{22}-r_{20}\right)=2(0{,}0208)=\SI{4.15}{\centi\meter}.
\]
\textbf{Por que o seletor é indispensável:} sem ele, íons de massas diferentes
mas velocidades diferentes poderiam produzir o \emph{mesmo} raio --- pois
$r\propto mv$. O seletor fixa $v$, tornando $r$ função apenas de $m/q$.\\
\textbf{Resolução do instrumento:} a separação relativa é
$\Delta r/r=\Delta m/m=10\%$ para esses isótopos --- larga. Separar isótopos
com $\Delta m/m$ de $0{,}1\%$ exigiria feixe muito bem colimado e detector de
alta resolução espacial.}
\stepcounter{questao}
\resposta{$F/L=\pot{2}{-7}\,\si{\newton\per\meter}$}
\resolucao{(a)
\[
\frac{F}{L}=\frac{\mu_0I^{2}}{2\pi d}
=\frac{(4\pi\times10^{-7})(1)}{2\pi(1)}=\pot{2}{-7}\,\si{\newton\per\meter}.
\]
(b) \textbf{Definição antiga (até 2019):} o ampère era \emph{definido} como a
corrente que, em dois condutores paralelos infinitos separados por
$\SI{1}{\meter}$ no vácuo, produzisse exatamente
$\pot{2}{-7}\,\si{\newton}$ por metro.\\
Essa definição \textbf{fixava $\mu_0$} em $4\pi\times10^{-7}$ por construção ---
não era um valor medido, mas uma consequência da escolha.\\
(c) \textbf{Desde 2019}, o ampère é definido fixando-se a \emph{carga
elementar} em $e=\pot{1.602176634}{-19}\,\si{\coulomb}$ exatamente.\\
\textbf{Consequência da mudança:} $\mu_0$ deixou de ser exato e passou a ser uma
grandeza \textbf{medida}, com incerteza experimental --- embora ainda
igual a $4\pi\times10^{-7}$ dentro de partes por bilhão.\\
A troca reflete uma mudança de filosofia: definir unidades por
\emph{constantes fundamentais} em vez de por experimentos difíceis de
reproduzir.}
\stepcounter{questao}
\resposta{$\SI{0.37}{\micro\volt}$ no cobre; $\SI{3.1}{\volt}$ no
semicondutor}
\resolucao{(a) A tensão Hall vale
\[
V_H=\frac{IB}{nqt}
=\frac{(1)(0{,}5)}{(\pot{8.5}{28})(\pot{1.6}{-19})(\pot{1}{-4})}
=\pot{3.68}{-7}\,\si{\volt}.
\]
(b) Com $n=\pot{1}{22}$:
\[
V_H=\frac{0{,}5}{(\pot{1}{22})(\pot{1.6}{-19})(\pot{1}{-4})}=\SI{3.1}{\volt}.
\]
\textbf{Diferença de sete ordens de grandeza.}\\
(c) \textbf{Explicação.} A tensão Hall é inversamente proporcional a $n$. Metais
têm um portador por átomo --- densidade altíssima --- e cada portador precisa se
mover \emph{devagar} para produzir a mesma corrente. Velocidade de deriva baixa
significa força magnética pequena e, portanto, separação de cargas mínima.\\
Semicondutores têm densidade de portadores milhões de vezes menor; eles se movem
muito mais rápido, e o efeito Hall se torna macroscópico.\\
\textbf{Consequência tecnológica:} \emph{todos} os sensores Hall comerciais são
de semicondutor. Um sensor de cobre exigiria amplificação de nanovolts,
inviável na presença de ruído térmico e de tensões termoelétricas.}
\stepcounter{questao}
\resposta{$\tau_{\max}=\SI{1.2}{\newton\meter}$;
$P\approx\SI{96}{\watt}$}
\resolucao{(a)
\[
m=NIA=200(3)(\pot{5}{-3})=\SI{3}{\ampere\meter\squared};
\qquad
\tau_{\max}=mB=3(0{,}4)=\SI{1.2}{\newton\meter}.
\]
\emph{(Atenção à conversão: $\SI{50}{\centi\meter\squared}
=\pot{5}{-3}\,\si{\meter\squared}$, e não $\pot{5}{-2}$.)}\\
(b) O torque varia como $\sin\theta$ ao longo da volta, e seu valor médio sobre
meio ciclo é $\dfrac{2}{\pi}\tau_{\max}=\SI{0.764}{\newton\meter}$.
\[
\omega=1200\ \text{rpm}=\frac{1200(2\pi)}{60}=\SI{125.7}{\radian\per\second};
\]
\[
P=\tau_{\text{méd}}\,\omega=0{,}764(125{,}7)=\SI{96}{\watt}.
\]
\textbf{Por que o comutador é necessário:} sem inverter a corrente a cada meia
volta, o torque mudaria de sinal e a espira apenas oscilaria em torno da posição
de alinhamento.\\
\textbf{Por que motores reais usam muitas bobinas:} com uma só, o torque
\emph{pulsa} entre zero e o máximo. Várias bobinas defasadas produzem torque
quase constante --- e é isso que se busca em qualquer máquina prática.}
\stepcounter{questao}
\resposta{Há força resultante, dirigida para a região de campo mais intenso (se
$\vec m$ estiver alinhado com $\vec B$)}
\resolucao{(a) \textbf{Sim.} Em campo \emph{uniforme}, as forças nos lados
opostos da espira se cancelam exatamente e resta apenas torque. Em campo
\emph{não uniforme}, os lados estão em campos de módulos diferentes --- e o
cancelamento é incompleto.\\
(b) \textbf{Sentido.} A força resultante é dada pelo gradiente:
\[
\vec F=\nabla\left(\vec m\cdot\vec B\right).
\]
Se $\vec m$ estiver \emph{alinhado} com $\vec B$, a espira é atraída para a
região de campo mais intenso. Se estiver \emph{antialinhado}, é repelida.\\
(c) \textbf{Relação com ímãs.} Um ímã permanente é, essencialmente, um conjunto
de espiras microscópicas (correntes de magnetização) com momentos alinhados.
Daí:
\begin{itemize}
\item \textbf{Dois ímãs se atraem ou repelem} conforme a orientação relativa dos
momentos --- exatamente o comportamento da espira.
\item \textbf{Um ímã atrai ferro não magnetizado} porque primeiro o magnetiza
(alinhando seus domínios com o campo aplicado) e depois o puxa para a região de
campo mais forte. A atração é sempre no sentido do gradiente.
\item \textbf{Em campo uniforme, um ímã não é atraído} --- ele apenas gira até
se alinhar. É por isso que a agulha de uma bússola aponta, mas não é puxada em
direção ao norte.
\end{itemize}}
\stepcounter{questao}
\resposta{Mesmo sentido: atração; opostos: repulsão. Módulos iguais}
\resolucao{(a) e (b) O \textbf{módulo} é o mesmo nos dois casos:
\[
\frac{F}{L}=\frac{\mu_0I_1I_2}{2\pi d}.
\]
O que muda é o \emph{sentido}: correntes paralelas \textbf{atraem-se},
antiparalelas \textbf{repelem-se}.\\
(c) \textbf{Explicação pelo campo.} Considere o fio 2 imerso no campo do fio 1.
\begin{itemize}
\item \textbf{Mesmo sentido:} na posição do fio 2, o campo de 1 aponta em certo
sentido; aplicando $\vec F=I\vec\ell\times\vec B$, a força aponta \emph{para} o
fio 1.
\item \textbf{Sentidos opostos:} $\vec\ell$ inverte, e a força inverte junto.
\end{itemize}
\textbf{Leitura pelas linhas de campo.} Com correntes paralelas, o campo se
\emph{cancela} entre os fios e se \emph{reforça} fora --- as linhas envolvem o
par, e a "tensão" ao longo delas puxa os fios um contra o outro. Com correntes
opostas, ocorre o inverso: campo intenso entre os fios, que os afasta.\\
\textbf{Consequência prática:} em curtos-circuitos, barramentos paralelos
sofrem forças enormes. Com $\SI{20}{\kilo\ampere}$ a $\SI{10}{\centi\meter}$,
a força chega a $\SI{800}{\newton\per\meter}$ --- razão de existirem
espaçadores mecânicos robustos em painéis de potência.}
\stepcounter{questao}
\resposta{Domina a elétrica em $\pot{1}{5}$, equilíbrio em $\pot{2}{5}$ e a
magnética em $\pot{4}{5}\,\si{\meter\per\second}$}
\resolucao{(a) A parcela elétrica \textbf{não depende de $v$}:
\[
F_E=qE=(\pot{1.6}{-19})(\pot{2}{4})=\pot{3.2}{-15}\,\si{\newton}.
\]
A magnética é \textbf{proporcional a $v$}: $F_B=qvB$.
\begin{center}
\begin{tabular}{cccl}
\hline
$v$ ($\si{\meter\per\second}$) & $F_E$ ($\si{\newton}$)
 & $F_B$ ($\si{\newton}$) & Resultante\\
\hline
$\pot{1}{5}$ & $\pot{3.2}{-15}$ & $\pot{1.6}{-15}$ & elétrica domina\\
$\pot{2}{5}$ & $\pot{3.2}{-15}$ & $\pot{3.2}{-15}$ & \textbf{nula}\\
$\pot{4}{5}$ & $\pot{3.2}{-15}$ & $\pot{6.4}{-15}$ & magnética domina\\
\hline
\end{tabular}
\end{center}
(b) e (c) \textbf{Interpretação.} O equilíbrio ocorre em
\[
v=\frac{E}{B}=\frac{\pot{2}{4}}{0{,}1}=\pot{2}{5}\,\si{\meter\per\second},
\]
e apenas nessa velocidade a partícula segue reta.\\
\textbf{As três situações da tabela descrevem o mesmo dispositivo}: partículas
lentas são empurradas para um lado pela força elétrica, rápidas para o outro
pela magnética, e só as de velocidade $E/B$ atravessam sem desvio. É a
\emph{seleção de velocidades}, vista agora como competição entre as duas
parcelas de $\vec F=q(\vec E+\vec v\times\vec B)$.\\
\textbf{Observe a estrutura:} uma parcela constante e outra linear em $v$. A
igualdade só pode ocorrer em \emph{um} valor de velocidade --- e é isso que dá
seletividade ao instrumento.}
\stepcounter{questao}
\resposta{Movimento cicloidal, com deriva média
$\vec v_d=\dfrac{\vec E\times\vec B}{B^{2}}$}
\resolucao{(a) \textbf{Movimento.} A partícula descreve uma trajetória
\textbf{cicloidal}: um movimento circular (devido a $\vec B$) superposto a um
deslocamento uniforme. Se ela parte do repouso, a curva é uma cicloide clássica
--- como um ponto na borda de uma roda que rola.\\
(b) \textbf{Velocidade de deriva.} O centro do círculo desloca-se com
\[
\vec v_d=\frac{\vec E\times\vec B}{B^{2}},
\qquad
|v_d|=\frac{E}{B},
\]
perpendicular a \emph{ambos} os campos.\\
\textbf{Duas propriedades notáveis:} a deriva \textbf{não depende da carga}
(íons e elétrons derivam no mesmo sentido) nem da \textbf{massa}.\\
(c) \textbf{Aplicações.}
\begin{itemize}
\item \textbf{Magnetrons} (fornos de micro-ondas): elétrons em campos cruzados
descrevem trajetórias cicloidais que excitam cavidades ressonantes.
\item \textbf{Confinamento de plasma:} a deriva $\vec E\times\vec B$ é um dos
mecanismos que governam o transporte em tokamaks --- e como ela não separa
cargas, não gera corrente própria.
\item \textbf{Magnetosfera terrestre:} explica parte do movimento de partículas
aprisionadas.
\end{itemize}}
\stepcounter{questao}
\resposta{$I=\SI{0.613}{\ampere}$; $a\approx\SI{2.9}{\meter\per\second\squared}$}
\resolucao{(a) A força magnética deve vencer o atrito estático:
\[
BIL=\mu_e mg
\;\Longrightarrow\;
I=\frac{0{,}25(0{,}100)(9{,}8)}{(0{,}8)(0{,}50)}=\SI{0.613}{\ampere}.
\]
(b) Com $I=\SI{1.226}{\ampere}$:
\[
F=BIL=(0{,}8)(1{,}226)(0{,}50)=\SI{0.490}{\newton};
\]
\[
f_c=\mu_c mg=0{,}20(0{,}100)(9{,}8)=\SI{0.196}{\newton};
\]
\[
a=\frac{0{,}490-0{,}196}{0{,}100}=\SI{2.94}{\meter\per\second\squared}.
\]
\textbf{Este é o princípio do canhão eletromagnético (\emph{railgun})} --- e
também do motor linear. \\
\textbf{Ressalva importante:} conforme o condutor acelera, ele passa a induzir
f.e.m. contrária ($\varepsilon=BLv$), que reduz a corrente. A aceleração
calculada vale apenas no \emph{instante inicial}; o movimento tende a uma
velocidade-limite em que a f.e.m. induzida equilibra a fonte.}
\stepcounter{questao}
\resposta{$\vec F\cdot\vec v=0$ sempre; o trabalho no motor vem da fonte, não do
campo}
\resolucao{(a) A potência é
\[
P=\vec F\cdot\vec v=q\left(\vec v\times\vec B\right)\cdot\vec v.
\]
O produto misto $(\vec v\times\vec B)\cdot\vec v$ é \textbf{identicamente
nulo}, pois $\vec v\times\vec B$ é perpendicular a $\vec v$ por construção. Logo
$P=0$ e $\Delta E_c=0$: o \emph{módulo} da velocidade nunca muda.\\
(b) \textbf{O aparente paradoxo do motor.} Um motor claramente realiza trabalho
--- ele levanta cargas e move veículos. Como conciliar?\\
\textbf{Resolução:} a força magnética atua sobre os \emph{portadores} do
condutor, e é perpendicular à velocidade \emph{deles}. Mas o fio se move, e a
velocidade total de um portador é a soma da deriva ao longo do fio com a
velocidade do fio.\\
A força magnética \emph{redistribui} energia: ela empurra o fio (realizando
trabalho positivo sobre ele) e, simultaneamente, freia a deriva dos portadores
(trabalho negativo de mesmo módulo). A soma continua sendo zero.\\
(c) \textbf{De onde vem a energia.} O freio à deriva manifesta-se como
\textbf{f.e.m. contraeletromotriz}, e é a \emph{fonte} que precisa vencê-la:
\[
P_{fonte}=\varepsilon_{fem}\,I=P_{\text{mecânica}}+P_{Joule}.
\]
\textbf{O campo magnético é um intermediário, não uma fonte.} Ele transfere
energia da fonte elétrica para o eixo, sem fornecer nem armazenar nada
líquido.\\
\emph{(É o mesmo papel de uma polia em mecânica: ela redireciona força sem criar
energia.)}}
\stepcounter{questao}
\resposta{$\vec\tau=\vec m\times\vec B$, com $\vec m=NI\vec A$}
\resolucao{(a) Seja a espira de lados $a$ e $b$, com $\vec B$ no plano
horizontal e o eixo de rotação vertical. Chamando $\theta$ o ângulo entre
$\vec B$ e a \emph{normal} à espira:
\begin{itemize}
\item Os lados de comprimento $a$ \emph{paralelos} ao eixo sofrem forças
$F=BIa$, perpendiculares tanto ao fio quanto ao campo, em sentidos opostos.
Elas formam um \textbf{binário}.
\item Os lados de comprimento $b$ sofrem forças que se cancelam e cuja linha de
ação passa pelo eixo --- não produzem torque.
\end{itemize}
(b) O braço de alavanca de cada força do binário é
$\dfrac{b}{2}\sin\theta$:
\[
\tau=2\left(BIa\right)\frac{b}{2}\sin\theta=BI(ab)\sin\theta=BIA\sin\theta.
\]
Com $N$ espiras, $\tau=NIAB\sin\theta$.\\
(c) \textbf{Generalização.} Definindo o momento magnético
$\vec m=NI\vec A$ (com $\vec A$ normal à espira, sentido dado pela regra da mão
direita):
\[
\boxed{\vec\tau=\vec m\times\vec B}
\]
Essa forma vale para espira de \textbf{qualquer} formato --- basta decompô-la em
retângulos infinitesimais, cujas correntes internas se cancelam aos pares.\\
\textbf{Analogia formal com o dipolo elétrico:} $\vec\tau=\vec p\times\vec E$.
As duas expressões são idênticas em estrutura, e por isso a espira é
literalmente um \emph{dipolo magnético}.}
\stepcounter{questao}
\resposta{$\Delta r/r=\Delta m/m$; a resolução exige feixe com dispersão angular
e de velocidade menores que $1\%$}
\resolucao{(a) De $r=\dfrac{mv}{qB}$ com $v$ e $B$ fixos:
\[
\frac{\Delta r}{r}=\frac{\Delta m}{m}=1\%.
\]
Com $r=\SI{20}{\centi\meter}$, a separação de raios é
$\SI{2}{\milli\meter}$, e no detector (após meia volta) a separação é
$2\Delta r=\SI{4}{\milli\meter}$.\\
(b) \textbf{Requisito de colimação.} Para \emph{resolver} duas linhas separadas
por $\SI{4}{\milli\meter}$, a largura de cada linha deve ser bem menor. Duas
fontes de alargamento:
\begin{itemize}
\item \textbf{Dispersão de velocidade:} como $r\propto v$, uma dispersão
$\Delta v/v$ produz alargamento $\Delta r/r$ de mesma magnitude. É preciso
$\Delta v/v\ll1\%$ --- daí a necessidade do seletor de velocidades.
\item \textbf{Dispersão angular:} íons entrando com pequenos ângulos diferentes
descrevem círculos deslocados. \emph{Felizmente}, a geometria de meia volta é
\textbf{autofocalizante}: íons que divergem de um ponto reconvergem após
$\ang{180}$, em primeira ordem. O alargamento residual é de segunda ordem no
ângulo.
\end{itemize}
(c) \textbf{Limitações práticas.}
\begin{itemize}
\item \textbf{Vácuo:} colisões com gás residual espalham o feixe. É preciso
pressão abaixo de $\pot{1}{-5}\,\si{\pascal}$.
\item \textbf{Uniformidade do campo:} variações de $B$ ao longo da trajetória
introduzem erro proporcional. Exige-se uniformidade melhor que a resolução
desejada.
\item \textbf{Carga espacial:} feixes intensos se repelem eletrostaticamente e
se alargam --- limitando a corrente utilizável.
\item \textbf{Íons de mesma razão $m/q$:} não são separáveis. $\mathrm{N_2^+}$ e
$\mathrm{CO^+}$ (ambos $28\,u$) exigem instrumentos de altíssima resolução para
distinguir.
\end{itemize}}
\stepcounter{questao}
\resposta{Com $B=\SI{0.5}{\tesla}$ e $L=\SI{0.2}{\meter}$: $I=\SI{50}{\ampere}$
com uma espira, ou $\SI{2.5}{\ampere}$ com vinte}
\resolucao{(a) De $F=NBIL$:
\[
NIL=\frac{F}{B}=\frac{5}{0{,}5}=\SI{10}{\ampere\meter}.
\]
Qualquer combinação que satisfaça esse produto serve --- há dois graus de
liberdade.\\
(b) \textbf{Compromisso.} Com $L=\SI{0.2}{\meter}$ fixo:
\begin{center}
\begin{tabular}{ccc}
\hline
$N$ & $I$ & Observação\\
\hline
$1$ & $\SI{50}{\ampere}$ & fio muito grosso, fonte de alta corrente\\
$20$ & $\SI{2.5}{\ampere}$ & equilíbrio razoável\\
$200$ & $\SI{0.25}{\ampere}$ & bobina volumosa, alta indutância\\
\hline
\end{tabular}
\end{center}
\textbf{Mais espiras} reduzem a corrente (fio mais fino, fonte mais simples),
mas aumentam volume, massa, resistência e \emph{indutância} --- o que torna o
atuador mais lento a responder.\\
\textbf{Menos espiras} dão resposta rápida, ao custo de corrente elevada.\\
(c) \textbf{Dissipação.} Com $N=20$, $I=\SI{2.5}{\ampere}$ e resistência total
estimada em $\SI{2}{\ohm}$:
\[
P=RI^{2}=2(6{,}25)=\SI{12.5}{\watt}.
\]
\textbf{Trabalho útil:}
\[
W=F\,d=5(0{,}02)=\SI{0.1}{\joule}.
\]
\textbf{Comparação reveladora:} se o curso levar $\SI{0.1}{\second}$, a potência
mecânica é $\SI{1}{\watt}$ contra $\SI{12.5}{\watt}$ dissipados --- rendimento
de $7\%$. Atuadores eletromagnéticos de curso curto são intrinsecamente
ineficientes, e é por isso que costumam operar em regime \emph{pulsado}, e não
contínuo.}
\stepcounter{questao}
\resposta{Equilíbrio entre força magnética e força elétrica transversal}
\resolucao{(a) \textbf{Dedução.} Os portadores, com velocidade de deriva
$v_d$, sofrem força magnética $qv_dB$ transversal e acumulam-se em uma face.
Isso cria campo elétrico transversal $E_H$, que cresce até equilibrar:
\[
qE_H=qv_dB
\;\Longrightarrow\;
E_H=v_dB.
\]
Com largura $w$, a tensão é $V_H=E_Hw=v_dBw$. Usando
$I=nqv_d(wt)$, elimina-se $v_d$:
\[
v_d=\frac{I}{nqwt}
\;\Longrightarrow\;
V_H=\frac{IB}{nqt}.
\]
\textbf{Note:} a largura $w$ cancela --- só a \emph{espessura} $t$ importa.\\
(b) \textbf{Sinal e tipo de portador.} Elétrons (negativos) e lacunas
(positivas) movem-se em sentidos opostos para produzir a mesma corrente. A força
$q\vec v\times\vec B$ os desvia para o \textbf{mesmo} lado --- mas com cargas de
sinais contrários, a polaridade de $V_H$ \textbf{se inverte}.\\
\textbf{É o único método simples de determinar experimentalmente o sinal dos
portadores} --- e foi assim que se estabeleceu que semicondutores tipo P
conduzem por lacunas.\\
(c) \textbf{Três aplicações.}
\begin{itemize}
\item \textbf{Medição de campo magnético} (gaussímetro): com $I$, $n$ e
$t$ conhecidos, $V_H\propto B$.
\item \textbf{Medição de corrente sem contato:} coloca-se o sensor no entreferro
de um núcleo que envolve o condutor. Mede $B$, e daí $I$ --- com isolação
galvânica completa.
\item \textbf{Caracterização de materiais:} com $B$ e $I$ conhecidos, obtém-se
$n$ (densidade de portadores) e o \emph{sinal} deles. Combinado com a
condutividade, fornece a mobilidade.
\end{itemize}
\textbf{Vantagem sobre alternativas:} sensores Hall respondem a campo
\emph{estático}, ao contrário de bobinas, que só detectam variação.}
\stepcounter{questao}
\resposta{$\theta=0$ é estável, $\theta=\pi$ é instável;
$T=2\pi\sqrt{J/(mB)}$}
\resolucao{(a) O torque é $\tau=-mB\sin\theta$, com $\theta$ o ângulo entre
$\vec m$ e $\vec B$. Ele se anula em
\[
\theta=0 \quad\text{e}\quad \theta=\pi.
\]
(b) \textbf{Classificação pela energia.} A energia potencial de um dipolo
magnético é
\[
U=-\vec m\cdot\vec B=-mB\cos\theta.
\]
\begin{itemize}
\item $\theta=0$ ($\vec m$ paralelo a $\vec B$): $U=-mB$, \textbf{mínimo}
$\Rightarrow$ equilíbrio \textbf{estável}. Deslocada, a espira sofre torque
restaurador.
\item $\theta=\pi$ (antiparalelo): $U=+mB$, \textbf{máximo} $\Rightarrow$
equilíbrio \textbf{instável}. Qualquer perturbação a faz girar $\ang{180}$.
\end{itemize}
(c) \textbf{Pequenas oscilações.} Para $\theta\ll1$, $\sin\theta\approx\theta$ e
\[
J\ddot\theta=-mB\theta
\;\Longrightarrow\;
\ddot\theta+\frac{mB}{J}\theta=0,
\]
equação do MHS com
\[
\omega=\sqrt{\frac{mB}{J}};
\qquad
T=2\pi\sqrt{\frac{J}{mB}},
\]
onde $J$ é o momento de inércia da espira.\\
\textbf{Aplicação prática:} é assim que funcionam o \textbf{galvanômetro de
bobina móvel} --- em que uma mola acrescenta torque restaurador proporcional a
$\theta$, tornando a deflexão proporcional à corrente --- e a \textbf{bússola},
cujo período de oscilação permite medir o campo terrestre se $J$ e $m$ forem
conhecidos.\\
\textbf{Diferença de $U$ entre os dois equilíbrios:} $2mB$ --- energia
necessária para inverter a espira, e a mesma que um ímã precisa receber para ter
sua polarização revertida.}
\stepcounter{questao}
\resposta{$F_E=\SI{8e-17}{\newton}$, $F_B=\SI{9.6e-17}{\newton}$,
perpendiculares; $|F|=\SI{1.25e-16}{\newton}$ a $\ang{50.2}$ de $\vec E$}
\resolucao{(a) \textbf{Comparação das duas parcelas.}
\begin{center}
\begin{tabular}{lll}
\hline
& Parcela elétrica & Parcela magnética\\
\hline
Expressão & $q\vec E$ & $q\,\vec v\times\vec B$\\
Depende de $v$? & \textbf{não} & \textbf{sim} (linear)\\
Age em repouso? & sim & não\\
Direção & paralela a $\vec E$ & $\perp$ a $\vec v$ \emph{e} a $\vec B$\\
Realiza trabalho? & \textbf{sim} & \textbf{nunca}\\
Altera $|v|$? & sim & não\\
\hline
\end{tabular}
\end{center}
\textbf{A linha decisiva é a do trabalho.} Como
$\left(\vec v\times\vec B\right)\cdot\vec v=0$ identicamente, a parcela
magnética tem potência nula --- ela desvia sem acelerar. A elétrica, ao
contrário, é a única capaz de variar a energia cinética.\\
(b) \textbf{Avaliação numérica.}
\[
F_E=qE=(\pot{1.6}{-19})(500)=\pot{8}{-17}\,\si{\newton};
\]
\[
F_B=qvB=(\pot{1.6}{-19})(\pot{3}{4})(0{,}02)=\pot{9.6}{-17}\,\si{\newton}.
\]
Com $\vec v$ paralelo a $\vec E$ e $\vec B$ perpendicular a ambos, a força
magnética é \textbf{perpendicular} à elétrica:
\[
|\vec F|=\sqrt{F_E^{2}+F_B^{2}}=\pot{1.25}{-16}\,\si{\newton},
\qquad
\theta=\arctan\frac{9{,}6}{8{,}0}=\ang{50.2}
\]
em relação a $\vec E$.\\
\textbf{Razão entre as parcelas:}
$\dfrac{F_B}{F_E}=\dfrac{vB}{E}=1{,}20$ --- adimensional, e igual a $1$
exatamente na velocidade $E/B$ do seletor.\\
\textbf{Trabalho em um deslocamento de $\SI{1}{\centi\meter}$ ao longo de
$\vec E$:} apenas a parcela elétrica contribui,
\[
W=F_E d=(\pot{8}{-17})(0{,}01)=\pot{8}{-19}\,\si{\joule}=\SI{5}{\electronvolt},
\]
que é simplesmente $qEd$ --- ou seja, $q$ vezes a tensão percorrida.\\
(c) \textbf{Divisão de papéis em aceleradores.}
\begin{itemize}
\item \textbf{Campo elétrico acelera.} É o único que realiza trabalho, e o ganho
de energia por passagem é $qU$ --- daí as cavidades de radiofrequência.
\item \textbf{Campo magnético guia.} Ele curva a trajetória sem custo
energético, permitindo fechar o percurso em círculo e reaproveitar as mesmas
cavidades milhares de vezes.
\item \textbf{Campo magnético também focaliza}, com quadrupolos que comprimem o
feixe transversalmente --- de novo, sem alterar a energia.
\end{itemize}
\textbf{Se fosse possível acelerar com campo magnético}, aceleradores lineares
de quilômetros seriam desnecessários. A impossibilidade não é técnica: é
consequência direta de $\left(\vec v\times\vec B\right)\perp\vec v$.\\
\textbf{Nota histórica:} a expressão completa reúne dois fenômenos que foram
descobertos e estudados separadamente. Sua unificação --- e a constatação de que
o que é campo elétrico em um referencial pode ser magnético em outro --- é um
dos pontos de partida da relatividade restrita.}
\stepcounter{questao}
\resposta{$f=\SI{22.9}{\mega\hertz}$; $r=\SI{30.5}{\centi\meter}$;
$100$ voltas; correção de $1{,}1\%$}
\resolucao{(a) A frequência de cíclotron independe da energia:
\[
f=\frac{qB}{2\pi m}=\frac{(\pot{1.6}{-19})(1{,}5)}{2\pi(\pot{1.67}{-27})}
=\SI{22.9}{\mega\hertz}.
\]
\textbf{É isso que torna o aparelho viável:} o oscilador que alimenta os "dês"
opera em frequência \emph{fixa}, embora a partícula ganhe energia
continuamente.\\
(b) Da energia final,
\[
v=\sqrt{\frac{2E}{m}}
=\sqrt{\frac{2(10)(\pot{1.6}{-13})}{\pot{1.67}{-27}}}
=\pot{4.38}{7}\,\si{\meter\per\second};
\]
\[
r=\frac{mv}{qB}=\SI{0.305}{\meter}
\;\Longrightarrow\;
\text{diâmetro}\approx\SI{61}{\centi\meter}.
\]
Com dois cruzamentos do intervalo por volta, cada volta rende
$\SI{100}{\kilo\electronvolt}$:
\[
N=\frac{10\,\si{\mega\electronvolt}}{0{,}1\,\si{\mega\electronvolt}}
=100\ \text{voltas}.
\]
\textbf{Tempo total:} $100/f=\SI{4.4}{\micro\second}$.\\
(c) \textbf{Correção relativística.} Com
$v/c=0{,}146$:
\[
\gamma=\frac{1}{\sqrt{1-(v/c)^{2}}}=1{,}0108,
\]
ou seja, a massa efetiva é $1{,}1\%$ maior no fim da aceleração.\\
\textbf{Consequência:} a frequência de cíclotron cai na mesma proporção, e a
partícula começa a chegar \emph{atrasada} ao intervalo acelerador. Após muitas
voltas o defasamento acumula e a aceleração cessa.\\
\textbf{Aqui, $1{,}1\%$ é tolerável} --- o cíclotron clássico funciona. Mas para
$\SI{100}{\mega\electronvolt}$ o desvio passaria de $10\%$ e seria proibitivo.\\
\textbf{Soluções:} o \textbf{sincrociclotron} modula a frequência ao longo do
ciclo; o \textbf{síncrotron} mantém o raio fixo e varia o campo. Ambos abandonam
a simplicidade do cíclotron justamente porque o período deixa de ser constante
--- e a raiz disso está na expressão $T=2\pi m/qB$, válida apenas enquanto
$m$ o for.}
\stepcounter{questao}
\resposta{$F/L\approx\SI{2.1}{\kilo\newton\per\meter}$ --- comparável ao peso
de uma pessoa por metro}
\resolucao{(a)
\[
\frac{F}{L}=\frac{\mu_0I^{2}}{2\pi d}
=\pot{2}{-7}\frac{(\pot{4}{4})^{2}}{0{,}15}
=\pot{2}{-7}\frac{\pot{1.6}{9}}{0{,}15}=\SI{2133}{\newton\per\meter}.
\]
Mais de duas toneladas-força por metro de barramento --- e \textbf{repulsiva},
se as correntes forem opostas (que é o caso em um curto entre fases).\\
(b) \textbf{Espaçamento entre suportes.} Tratando o barramento como viga
biapoiada com carga distribuída $w=F/L$, a flecha é
\[
\delta=\frac{5wL_s^{4}}{384EJ}.
\]
Para deflexão admissível $\delta$, o vão máximo cresce apenas com
$L_s\propto\delta^{1/4}$ --- \textbf{dependência muito fraca}. Dobrar a flecha
admissível permite aumentar o vão em apenas $19\%$.\\
Na prática, os suportes ficam a cada $\SI{0.5}{}$ a $\SI{1}{\meter}$.\\
(c) \textbf{Caráter da solicitação.}
\begin{itemize}
\item \textbf{É transitória} --- dura o tempo de atuação da proteção, tipicamente
$\SI{100}{\milli\second}$. Mas a força \emph{de pico} pode ser o dobro da
calculada, devido à componente assimétrica da corrente de curto.
\item \textbf{Vai com $I^{2}$} --- dobrar a corrente de curto quadruplica a
força. É por isso que a corrente de curto-circuito presumida é um dado de
projeto tão crítico.
\item \textbf{É dinâmica:} se a frequência natural da estrutura estiver próxima
de $\SI{120}{\hertz}$ (dobro da rede, pois $F\propto i^{2}$), pode haver
ressonância mecânica.
\end{itemize}
\textbf{Conclusão:} em regime normal ($\SI{1}{\kilo\ampere}$, digamos), a força
seria $\SI{1.3}{\newton\per\meter}$ --- desprezível. O dimensionamento mecânico
de barramentos é governado inteiramente pela condição de \emph{falta}.}

\gabarito[12.4 Força magnética]

\subsection{Lei de Lenz e indução eletromagnética}
\stepcounter{questao}
\resposta{(b)}
\resolucao{A lei de Lenz é uma manifestação da conservação de energia: a corrente
induzida cria um campo que \emph{se opõe} à variação de fluxo que a gerou. Se
reforçasse a variação (alternativa a), teríamos energia criada do nada.}
\stepcounter{questao}
\resposta{$\varepsilon = \SI{20}{\milli\volt}$}
\resolucao{A fem é a taxa de variação do fluxo:
\[
\varepsilon=\left|\frac{\mathrm{d}\Phi}{\mathrm{d}t}\right|
=\SI{0.02}{\weber\per\second}=\SI{20}{\milli\volt}.
\]
Como o fluxo cresce linearmente, a fem é \emph{constante}.}
\stepcounter{questao}
\resposta{$\varepsilon = \SI{25}{\milli\volt}$}
\resolucao{Com $N = 50$ espiras e
$A = \pot{20}{-4}\,\si{\meter\squared}$:
\[
\varepsilon = N\,A\,\frac{\Delta B}{\Delta t}
= 50\cdot\pot{20}{-4}\cdot\frac{0{,}5}{2}
= 50\cdot\pot{20}{-4}\cdot0{,}25 = \SI{25}{\milli\volt}.
\]}
\stepcounter{questao}
\resposta{$\varepsilon = \SI{30}{\milli\volt}$, constante}
\resolucao{A taxa de variação do campo é
$\dfrac{\mathrm{d}B}{\mathrm{d}t} = \SI{-0.1}{\tesla\per\second}$ (constante):
\[
|\varepsilon| = N\,A\,\left|\frac{\mathrm{d}B}{\mathrm{d}t}\right|
= 100\cdot\pot{30}{-4}\cdot0{,}1 = \SI{30}{\milli\volt}.
\]
Como a taxa de variação é constante, a fem é constante --- e, num circuito de
resistência fixa, a corrente induzida também é \emph{constante}. O fluxo está
diminuindo, então a corrente induzida atua no sentido de \emph{reforçá-lo}
(Lenz).}
\stepcounter{questao}
\resposta{(a)}
\resolucao{Ao aproximar o polo norte, o fluxo através da espira \emph{aumenta}.
Pela lei de Lenz, a corrente induzida se opõe a esse aumento, criando um campo que
\emph{repele} o ímã --- ou seja, a face da espira voltada para o ímã se comporta
como um polo \textbf{norte}. Consequentemente, surge uma força de \emph{repulsão}
que dificulta a aproximação. É a manifestação da conservação de energia: para
aproximar o ímã é preciso realizar trabalho, que se converte na energia elétrica
da corrente induzida.}
\stepcounter{questao}
\resposta{$\varepsilon = \SI{2}{\volt}$}
\resolucao{$\varepsilon = N\dfrac{\mathrm{d}\Phi}{\mathrm{d}t}
= 200\cdot0{,}01 = \SI{2}{\volt}$. É assim que o transformador transfere energia:
o fluxo variável criado pelo primário induz tensão no secundário,
proporcionalmente ao número de espiras de cada lado.}
\stepcounter{questao}
\resposta{(a) $\varepsilon(t) = \Phi_0\omega\sin(\omega t)$;
(b) $\varepsilon_{\max} = \SI{1}{\volt}$}
\resolucao{(a) A fem é a derivada do fluxo (com sinal trocado):
\[
\varepsilon = -\frac{\mathrm{d}\Phi}{\mathrm{d}t}
= -\Phi_0\cdot(-\omega\sin(\omega t)) = \Phi_0\,\omega\sin(\omega t).
\]
(b) O valor máximo ocorre quando $\sin(\omega t) = 1$:
\[
\varepsilon_{\max} = \Phi_0\,\omega = 0{,}01\cdot100 = \SI{1}{\volt}.
\]
\textbf{Observação:} um fluxo cossenoidal gera uma fem \emph{senoidal} --- há uma
defasagem de $\SI{90}{\degree}$ entre eles. É exatamente o que ocorre em um
gerador de corrente alternada: a espira girando em campo uniforme produz fluxo
senoidal e, portanto, tensão senoidal.}
\stepcounter{questao}
\resposta{$\mathcal{E}(t)=RI_0\sin(\omega t)$; inversões em $t=0$, $\pi/\omega$ e $2\pi/\omega$}
\resolucao{Como o circuito é resistivo, $\mathcal{E}=Ri$. O sinal da senoide muda a
cada múltiplo de $\pi$, portanto o sentido da corrente se inverte a cada meio período.}
\stepcounter{questao}
\resposta{(a) $\varepsilon=\SI{0.6}{\volt}$, $i=\SI{0.3}{\ampere}$;
(b) $F=\SI{0.06}{\newton}$, opondo-se ao movimento;
(c) ambas $=\SI{0.18}{\watt}$; (d) ver comentário}
\resolucao{\textbf{(a) Fem de movimento.} À medida que a barra avança, a área do
circuito aumenta a uma taxa $\dfrac{\mathrm{d}A}{\mathrm{d}t} = \ell v$. Logo
\[
\varepsilon=\left|\frac{\mathrm{d}\Phi}{\mathrm{d}t}\right|
= B\,\ell\,v = 0{,}5\cdot0{,}4\cdot3=\SI{0.6}{\volt},
\qquad
i=\frac{\varepsilon}{R}=\frac{0{,}6}{2}=\SI{0.3}{\ampere}.
\]
\textbf{(b) Força sobre a barra.} A barra, agora percorrida por corrente,
encontra-se no campo $B$ e sofre
\[
F = B\,i\,\ell = 0{,}5\cdot0{,}3\cdot0{,}4=\SI{0.06}{\newton}.
\]
Pela lei de Lenz, essa força \textbf{opõe-se ao movimento} --- ela freia a barra.
(Se favorecesse o movimento, teríamos aceleração espontânea e energia criada do
nada.)

\textbf{(c) Balanço de potência.}
\[
P_{\text{dissipada}} = R\,i^{2} = 2(0{,}3)^{2}=\SI{0.18}{\watt},
\]
\[
P_{\text{mecânica}} = F\,v = 0{,}06\cdot3 = \SI{0.18}{\watt}.
\]
\textbf{São exatamente iguais.}

\textbf{(d) Interpretação.} O resultado não é coincidência: é a
\textbf{conservação de energia} em ação. Para manter a barra em movimento
\emph{uniforme}, um agente externo deve aplicar uma força igual e oposta à força
magnética de frenagem, realizando trabalho a uma taxa $Fv$. Toda essa potência
mecânica é convertida em potência elétrica e, finalmente, dissipada como calor no
resistor. Nada se perde, nada se cria:
\[
\underbrace{P_{\text{mec}}}_{\text{trabalho externo}}
\;\longrightarrow\;
\underbrace{\varepsilon\,i}_{\text{potência elétrica}}
\;\longrightarrow\;
\underbrace{R i^{2}}_{\text{calor}}.
\]
Verificando pela via elétrica: $\varepsilon\,i = 0{,}6\cdot0{,}3
= \SI{0.18}{\watt}$ \checkmark\\
\textbf{Este é o princípio do gerador elétrico} --- e, invertendo o raciocínio, o
do \emph{freio eletromagnético} usado em trens e em esteiras de academia: quanto
maior a velocidade, maior a corrente induzida e maior a frenagem, sem contato
mecânico nem desgaste.}
\stepcounter{questao}
\resposta{$\varepsilon=\SI{80}{\volt}$}
\resolucao{
\[
|\varepsilon|=N\frac{\Delta\Phi}{\Delta t}
=200\,\frac{\pot{2}{-3}}{\pot{5}{-3}}=\SI{80}{\volt}.
\]
\textbf{Note o papel de $N$:} cada espira contribui com a mesma f.e.m., e elas
estão em série. Duzentas espiras multiplicam por duzentos --- é assim que se
obtêm tensões altas a partir de variações de fluxo modestas.}
\stepcounter{questao}
\resposta{$\varepsilon=\SI{0.6}{\volt}$}
\resolucao{
\[
\varepsilon=B\ell v=(0{,}5)(0{,}30)(4)=\SI{0.6}{\volt}.
\]
\textbf{Origem da fórmula:} os portadores da barra movem-se com ela e sofrem
força magnética $qvB$ ao longo do comprimento. Eles se acumulam nas
extremidades até que o campo elétrico resultante equilibre a força magnética ---
e a diferença de potencial daí resultante é $B\ell v$.\\
\textbf{Equivalência com Faraday:} se a barra fizer parte de um circuito, a área
varia à taxa $\ell v$ e $\mathrm{d}\Phi/\mathrm{d}t=B\ell v$ --- o mesmo
resultado.}
\stepcounter{questao}
\resposta{$N=300$}
\resolucao{
\[
N=\frac{\varepsilon}{\mathrm{d}\Phi/\mathrm{d}t}=\frac{12}{0{,}04}=300.
\]
\textbf{Compromisso de projeto:} mais espiras dão mais tensão, mas aumentam a
resistência do enrolamento (mais perdas) e o volume. Aumentar
$\mathrm{d}\Phi/\mathrm{d}t$ --- por campo mais intenso ou variação mais rápida
--- costuma ser preferível quando possível.}
\stepcounter{questao}
\resposta{$\Delta t=\SI{10}{\milli\second}$}
\resolucao{
\[
\Delta t=\frac{N\,\Delta\Phi}{\varepsilon}
=\frac{100(\pot{5}{-3})}{50}=\pot{1}{-2}\,\si{\second}.
\]
\textbf{Leitura:} a f.e.m. é inversamente proporcional ao tempo. A \emph{mesma}
variação de fluxo, feita dez vezes mais rápido, produz dez vezes mais tensão ---
princípio de funcionamento das bobinas de ignição e dos geradores de alta
tensão por interrupção de corrente.}
\stepcounter{questao}
\resposta{(a)}
\resolucao{O sinal indica que a corrente induzida circula no sentido de
\textbf{opor-se} à variação de fluxo que a produziu.\\
\textbf{Por que precisa ser assim:} se a corrente induzida \emph{reforçasse} a
variação, ela aumentaria o fluxo, que induziria mais corrente, e assim por
diante --- energia surgiria do nada. A lei de Lenz é a garantia de que isso não
ocorre.\\
\textbf{Consequência prática:} qualquer variação de fluxo encontra
\emph{resistência}. Aproximar um ímã de uma espira fechada exige força; girar
um gerador sob carga exige torque. É de onde vem a energia elétrica.}
\stepcounter{questao}
\resposta{$\varepsilon=0$}
\resolucao{O fluxo é $\Phi=BA$, com $B$ e $A$ constantes --- ele \textbf{não
varia}, e portanto
\[
\varepsilon=-\frac{\mathrm{d}\Phi}{\mathrm{d}t}=0.
\]
\textbf{Ponto conceitual central:} \emph{movimento não induz f.e.m.}; variação
de fluxo é que induz. Uma espira pode mover-se a grande velocidade sem gerar
nada, desde que o fluxo através dela permaneça o mesmo.\\
\textbf{Onde há f.e.m.:} apenas ao \emph{atravessar a fronteira} da região de
campo, quando a área imersa muda.}
\stepcounter{questao}
\resposta{$\si{\weber}=\si{\volt\second}$}
\resolucao{Da lei de Faraday, $\varepsilon=\mathrm{d}\Phi/\mathrm{d}t$, logo
\[
[\si{\volt}]=\frac{[\si{\weber}]}{[\si{\second}]}
\;\Longrightarrow\;
\si{\weber}=\si{\volt}\cdot\si{\second}.
\]
\textbf{Leitura útil:} o weber é um "volt-segundo". Isso torna imediato o
resultado de que a \emph{integral} da f.e.m. no tempo é a variação de fluxo ---
base do fluxímetro e do método de medição por integração.}
\stepcounter{questao}
\resposta{(b)}
\resolucao{A corrente induzida produz campo que se \emph{opõe ao aumento} --- ou
seja, no sentido contrário ao campo externo dentro da espira.\\
\textbf{Se o fluxo diminuísse}, a corrente inverteria, criando campo no
\emph{mesmo} sentido do externo, para tentar sustentá-lo.\\
\textbf{Regra prática:} a espira "não gosta de mudanças". Ela sempre reage no
sentido de manter o fluxo como estava --- nunca no de acelerar a mudança.}
\stepcounter{questao}
\resposta{$\varepsilon=\SI{0.6}{\volt}$; $I=\SI{0.3}{\ampere}$;
$F=\SI{60}{\milli\newton}$; $P=\SI{0.18}{\watt}$ pelos dois caminhos}
\resolucao{(a)
\[
\varepsilon=B\ell v=(0{,}4)(0{,}5)(3)=\SI{0.6}{\volt};
\qquad
I=\frac{0{,}6}{2}=\SI{0.3}{\ampere}.
\]
(b) A corrente induzida na barra, imersa no campo, sofre força magnética
\emph{oposta} ao movimento (lei de Lenz):
\[
F=BI\ell=(0{,}4)(0{,}3)(0{,}5)=\SI{0.06}{\newton}.
\]
Para manter a velocidade constante, o agente externo deve aplicar força igual e
contrária.\\
(c) \textbf{Balanço.}
\[
P_{\text{mecânica}}=Fv=(0{,}06)(3)=\SI{0.18}{\watt};
\qquad
P_{\text{elétrica}}=\frac{\varepsilon^{2}}{R}=\frac{0{,}36}{2}=\SI{0.18}{\watt}.
\]
\textbf{Coincidem exatamente} --- toda a energia mecânica fornecida vira calor no
resistor. É a lei de Lenz em forma quantitativa: a força de oposição é
precisamente a necessária para que a conta feche.}
\stepcounter{questao}
\resposta{$\varepsilon_{\max}=\SI{377}{\volt}$;
$\varepsilon_{ef}=\SI{267}{\volt}$}
\resolucao{(a) Com $\Phi=BA\cos\omega t$:
\[
\varepsilon=-N\frac{\mathrm{d}\Phi}{\mathrm{d}t}=NBA\omega\sin\omega t
\;\Longrightarrow\;
\varepsilon_{\max}=NBA\omega.
\]
\[
\varepsilon_{\max}=100(0{,}5)(\pot{2}{-2})(377)=\SI{377}{\volt}.
\]
(b) Para forma senoidal,
\[
\varepsilon_{ef}=\frac{\varepsilon_{\max}}{\sqrt2}=\SI{267}{\volt}.
\]
\textbf{Note:} $\omega=\SI{377}{\radian\per\second}$ corresponde a
$\SI{60}{\hertz}$ --- é o valor $2\pi(60)$, e aparece constantemente em
eletrotécnica.\\
\textbf{Compromisso do gerador:} $\varepsilon\propto\omega$, mas a frequência é
fixada pela rede. Resta aumentar $N$, $B$ ou $A$ --- e é por isso que geradores
de grande porte têm núcleos volumosos e campos próximos da saturação.}
\stepcounter{questao}
\resposta{$\varepsilon_2=\SI{10}{\volt}$}
\resolucao{
\[
\varepsilon_2=M\frac{\mathrm{d}i_1}{\mathrm{d}t}=(0{,}05)(200)=\SI{10}{\volt}.
\]
\textbf{Simetria da indução mútua:} $M_{12}=M_{21}$ sempre, mesmo que as bobinas
sejam muito diferentes. Uma corrente variando na bobina pequena induz na grande
exatamente a mesma f.e.m. que o inverso produziria --- resultado de
reciprocidade, análogo ao que vale em redes resistivas.\\
\textbf{Aplicação:} é este o princípio do transformador, e também da
interferência indesejada entre circuitos vizinhos.}
\stepcounter{questao}
\resposta{$V_2=\SI{22}{\volt}$; $I_1=\SI{0.5}{\ampere}$}
\resolucao{(a) Como o \emph{mesmo} fluxo atravessa os dois enrolamentos,
\[
\frac{V_2}{V_1}=\frac{N_2}{N_1}
\;\Longrightarrow\;
V_2=220\left(\frac{100}{1000}\right)=\SI{22}{\volt}.
\]
(b) No transformador ideal não há perdas, e a potência se conserva:
\[
V_1I_1=V_2I_2
\;\Longrightarrow\;
I_1=\frac{22(5)}{220}=\SI{0.5}{\ampere}.
\]
\textbf{Relação inversa:} a corrente transforma-se na razão \emph{oposta} à da
tensão. Abaixar tensão dez vezes eleva a corrente dez vezes.\\
\textbf{Consequência para os enrolamentos:} o secundário precisa de fio dez
vezes mais grosso, embora tenha dez vezes menos espiras. O volume de cobre acaba
sendo comparável nos dois lados.}
\stepcounter{questao}
\resposta{$\varepsilon(2)=\SI{0.2}{\volt}$; média de $\SI{0.1}{\volt}$}
\resolucao{\textbf{Instantânea:}
\[
\varepsilon=-\frac{\mathrm{d}\Phi}{\mathrm{d}t}=-0{,}10\,t
\;\Longrightarrow\;
|\varepsilon(2)|=\SI{0.2}{\volt}.
\]
\textbf{Média:}
\[
|\bar\varepsilon|=\frac{\Delta\Phi}{\Delta t}
=\frac{0{,}05(4)-0}{2}=\SI{0.1}{\volt}.
\]
\textbf{As duas diferem por um fator $2$} --- e ambas estão corretas, para
perguntas diferentes.\\
\textbf{Quando coincidem:} apenas se o fluxo variar \emph{linearmente}. Aqui a
variação é quadrática, e a f.e.m. cresce ao longo do intervalo, partindo de zero
e chegando a $\SI{0.2}{\volt}$ --- a média é o valor intermediário.}
\stepcounter{questao}
\resposta{$q=\SI{20}{\milli\coulomb}$}
\resolucao{Integrando a corrente induzida:
\[
q=\int i\,\mathrm{d}t=\int\frac{|\varepsilon|}{R}\,\mathrm{d}t
=\frac{N}{R}\int\left|\frac{\mathrm{d}\Phi}{\mathrm{d}t}\right|\mathrm{d}t
=\frac{N\,\Delta\Phi}{R}.
\]
\[
q=\frac{50(\pot{4}{-3})}{10}=\pot{2}{-2}\,\si{\coulomb}.
\]
\textbf{Resultado notável: a carga \emph{não depende do tempo}.} Rápida ou
lenta, a mesma variação de fluxo faz circular a mesma carga --- porque
f.e.m. e duração variam inversamente e o produto se conserva.\\
\textbf{É o princípio do fluxímetro balístico:} integrando a corrente com um
galvanômetro de bobina móvel, mede-se $\Delta\Phi$ sem controlar a velocidade do
movimento.}
\stepcounter{questao}
\resposta{$\SI{0.12}{\volt}$ durante a entrada; zero depois}
\resolucao{(a) Enquanto atravessa a fronteira, apenas o lado dianteiro está no
campo e funciona como barra móvel:
\[
\varepsilon=B\ell v=(0{,}6)(0{,}10)(2)=\SI{0.12}{\volt}.
\]
\emph{(Equivalentemente: a área imersa cresce a $\ell v$, e
$\mathrm{d}\Phi/\mathrm{d}t=B\ell v$.)}\\
(b) \textbf{Totalmente dentro}, os lados dianteiro e traseiro estão ambos no
campo e suas f.e.m. se \emph{cancelam}. O fluxo é constante e
\[
\varepsilon=0.
\]
\textbf{Perfil completo:} pulso retangular de $\SI{0.12}{\volt}$ durante a
entrada, zero no percurso interno, e pulso de \emph{polaridade invertida}
durante a saída. Cada pulso dura $\ell/v=\SI{50}{\milli\second}$.}
\stepcounter{questao}
\resposta{Ele freia --- correntes de Foucault opõem-se ao movimento}
\resolucao{O disco, ao girar, tem diferentes regiões atravessadas por fluxos
variáveis. Isso induz \textbf{correntes de Foucault} (ou parasitas) --- redemoinhos
de corrente no próprio material.\\
Pela lei de Lenz, essas correntes produzem forças que se \textbf{opõem ao
movimento}: o disco desacelera, e a energia cinética vira calor.\\
\textbf{Duas leituras.}
\begin{itemize}
\item \textbf{Como recurso:} é o \emph{freio magnético}, usado em trens, em
esteiras ergométricas e no disco do antigo medidor de energia. Não há contato,
portanto não há desgaste.
\item \textbf{Como problema:} em núcleos de transformadores e motores, as
mesmas correntes representam perda pura. Combate-se \textbf{laminando} o núcleo
--- assunto retomado no bloco de desafio.
\end{itemize}
\textbf{Propriedade importante:} a força de frenagem é proporcional à
\emph{velocidade}. O freio é forte em alta velocidade e desaparece no repouso ---
por isso ele nunca segura um veículo parado.}
\stepcounter{questao}
\resposta{$\varepsilon=\SI{75}{\volt}$}
\resolucao{
\[
|\varepsilon|=N\frac{\Delta\Phi}{\Delta t}
=500\,\frac{\pot{1.2}{-3}}{\pot{8}{-3}}=\SI{75}{\volt}.
\]
\textbf{Papel do núcleo:} sem ele, o mesmo enrolamento e a mesma corrente
produziriam fluxo centenas de vezes menor --- e f.e.m. proporcionalmente menor.
O ferro é o que torna a indução tecnologicamente útil.\\
\textbf{Ressalva:} se o fluxo variar a partir de um valor já próximo da
saturação, a relação deixa de ser linear e a f.e.m. real fica abaixo da
prevista.}
\stepcounter{questao}
\resposta{Atrativa --- ela tenta impedir o afastamento}
\resolucao{Ao afastar o ímã, o fluxo através da espira \textbf{diminui}. Pela lei
de Lenz, a corrente induzida circula no sentido de \emph{sustentar} o fluxo ---
criando campo no mesmo sentido do original.\\
A espira comporta-se então como um ímã com o polo oposto voltado ao ímã que se
afasta, e portanto o \textbf{atrai}.\\
\textbf{Pela terceira lei de Newton}, o ímã atrai a espira com força igual e
contrária.\\
\textbf{Conclusão geral:} a espira sempre \emph{resiste} ao que se está fazendo.
Aproximar exige empurrar contra repulsão; afastar exige puxar contra atração.
Em ambos os casos, o trabalho realizado é o que aparece como energia elétrica
--- e o resistor aquece.}
\stepcounter{questao}
\resposta{$\varepsilon=\SI{80}{\volt}$; $I=\SI{2}{\ampere}$;
$q=\SI{40}{\milli\coulomb}$}
\resolucao{(a)
\[
|\varepsilon|=N\frac{\Delta\Phi}{\Delta t}
=200\,\frac{(0{,}8)(\pot{1}{-2})}{\pot{2}{-2}}=\SI{80}{\volt};
\qquad
I=\frac{80}{40}=\SI{2}{\ampere}.
\]
(b) O fluxo \textbf{diminui}, logo a corrente induzida circula no sentido de
\emph{sustentá-lo} --- criando campo próprio no \textbf{mesmo} sentido do campo
original.\\
\emph{(Se o campo estivesse sendo ligado, a corrente circularia ao contrário.)}\\
(c)
\[
q=\frac{N\,\Delta\Phi}{R}=\frac{200(\pot{8}{-3})}{40}
=\pot{4}{-2}\,\si{\coulomb}.
\]
\textbf{Observação de segurança:} $\SI{80}{\volt}$ surgem em uma bobina que
antes não tinha tensão alguma. Desligar bruscamente eletroímãs de grande porte
produz tensões muito superiores --- razão pela qual eles exigem circuitos de
desmagnetização controlada.}
\stepcounter{questao}
\resposta{$v_{lim}=\SI{12.25}{\meter\per\second}$; $P=\SI{12}{\watt}$}
\resolucao{(a) A barra sofre o peso e a força magnética de oposição:
\[
m\frac{\mathrm{d}v}{\mathrm{d}t}=mg-\frac{B^{2}\ell^{2}v}{R},
\]
pois $F=BI\ell$ com $I=B\ell v/R$.\\
\textbf{Estrutura da equação:} é idêntica à de um corpo em queda com arrasto
\emph{linear} --- a força de oposição é proporcional à velocidade.\\
(b) Na velocidade-limite, a aceleração se anula:
\[
v_{lim}=\frac{mgR}{B^{2}\ell^{2}}
=\frac{(0{,}1)(9{,}8)(0{,}5)}{(0{,}25)(0{,}16)}
=\frac{0{,}49}{0{,}04}=\SI{12.25}{\meter\per\second}.
\]
(c)
\[
P=mgv_{lim}=(0{,}1)(9{,}8)(12{,}25)=\SI{12}{\watt},
\]
que é exatamente $\dfrac{(B\ell v_{lim})^{2}}{R}$ \checkmark.\\
\textbf{Interpretação:} em regime, toda a energia potencial gravitacional
perdida vira calor no resistor. A barra não ganha mais energia cinética.\\
\textbf{Efeito de $R$:} reduzir $R$ \emph{diminui} a velocidade-limite --- o
freio fica mais forte. Com $R\to0$ (trilhos supercondutores), a barra mal se
moveria.}
\stepcounter{questao}
\resposta{$q=\dfrac{N\Delta\Phi}{R}$ --- independe de $\Delta t$}
\resolucao{(a)
\[
q=\int i\,\mathrm{d}t
=\int\frac{N}{R}\frac{\mathrm{d}\Phi}{\mathrm{d}t}\,\mathrm{d}t
=\frac{N}{R}\int\mathrm{d}\Phi
=\frac{N\,\Delta\Phi}{R}.
\]
O tempo \textbf{cancela} na integração.\\
(b) \textbf{Explicação física.} Retirar a bobina depressa produz f.e.m. alta
durante pouco tempo; devagar, f.e.m. baixa durante muito tempo. O
\emph{produto} --- que é a carga --- permanece o mesmo.\\
Algebricamente: $\varepsilon\propto1/\Delta t$ e a duração é
$\Delta t$; o produto não depende dela.\\
(c) \textbf{Aplicação: fluxímetro balístico.} Um galvanômetro balístico responde
à \emph{carga} total, não à corrente instantânea. Conectado à bobina, sua
deflexão máxima é proporcional a $\Delta\Phi$ --- e o operador não precisa
controlar a velocidade com que retira a bobina.\\
\textbf{Robustez notável:} nem a velocidade, nem a suavidade, nem a trajetória
influenciam o resultado. É uma das medições mais tolerantes a erro de
procedimento em todo o eletromagnetismo.}
\stepcounter{questao}
\resposta{Redução de $400$ vezes}
\resolucao{(a) As perdas por correntes de Foucault em lâminas obedecem a
\[
P\propto\frac{B_{\max}^{2}f^{2}t^{2}}{\rho},
\]
com $t$ a espessura da lâmina. Reduzindo $t$ de $10$ para
$\SI{0.5}{\milli\meter}$:
\[
\frac{P'}{P}=\left(\frac{0{,}5}{10}\right)^{2}=\frac{1}{400}.
\]
(b) \textbf{Mecanismo.} As correntes parasitas circulam em laços
\emph{perpendiculares} ao fluxo. Laços maiores encerram mais fluxo (mais f.e.m.)
e encontram menos resistência (caminho mais largo) --- e por isso a perda cresce
com o \emph{quadrado} da dimensão.\\
Isolar as lâminas quebra os laços grandes em muitos laços pequenos, cada um com
f.e.m. menor e resistência maior.\\
\textbf{Importante:} as lâminas devem ser paralelas ao fluxo e isoladas
\emph{entre si} --- um verniz de poucos micrômetros basta. Se houver rebarba
metálica unindo-as, o efeito se perde.\\
(c) \textbf{Limitações.}
\begin{itemize}
\item \textbf{Fator de empilhamento:} o isolamento ocupa espaço, e a seção
magnética efetiva cai para $95$ a $97\%$ da geométrica.
\item \textbf{Custo:} lâminas mais finas são mais caras de produzir e montar.
\item \textbf{Frequência:} como $P\propto f^{2}t^{2}$, em alta frequência nem
lâminas finíssimas bastam --- daí o uso de \textbf{ferrites}, que são cerâmicas
de altíssima resistividade e praticamente não conduzem correntes parasitas.
\end{itemize}}
\stepcounter{questao}
\resposta{$M=\mu_0n_1N_2A$; $M=\SI{1.26}{\milli\henry}$}
\resolucao{(a) A corrente $i_1$ no solenoide cria $B=\mu_0n_1i_1$ em seu
interior, e fluxo $\Phi=BA$ por espira da bobina externa. O fluxo concatenado
nela é
\[
\lambda_2=N_2\mu_0n_1i_1A
\;\Longrightarrow\;
M=\frac{\lambda_2}{i_1}=\mu_0n_1N_2A.
\]
(b)
\[
M=(4\pi\times10^{-7})(2000)(100)(\pot{5}{-4})
=\pot{1.26}{-4}\,\si{\henry}=\SI{0.126}{\milli\henry}.
\]
(c) \textbf{Por que o raio externo não importa.} O campo do solenoide longo é
praticamente nulo do lado de fora. Toda a contribuição ao fluxo vem da região
\emph{interna} ao solenoide, cuja área é $A$ --- não a área da bobina
externa.\\
Enrolar a bobina justa ou folgada dá o \emph{mesmo} $M$.\\
\textbf{Consequência de projeto:} para aumentar o acoplamento, aumente a seção
do solenoide, não o diâmetro do enrolamento secundário. E note que
$M=M_{21}=M_{12}$ --- calcular pelo caminho fácil (do solenoide para a bobina) dá
o mesmo resultado que o caminho difícil.}
\stepcounter{questao}
\resposta{$\varepsilon=-M\omega I_0\cos\omega t$, com
$M=\dfrac{\mu_0a}{2\pi}\ln\dfrac{d+b}{d}$}
\resolucao{(a) Do resultado já obtido para o fluxo de um fio através de espira
retangular:
\[
\Phi(t)=\frac{\mu_0\,i(t)\,a}{2\pi}\ln\!\left(\frac{d+b}{d}\right)
=M\,i(t),
\]
com $M$ a indutância mútua entre fio e espira.\\
(b)
\[
\varepsilon=-\frac{\mathrm{d}\Phi}{\mathrm{d}t}
=-M\,\frac{\mathrm{d}i}{\mathrm{d}t}
=-M\omega I_0\cos\omega t.
\]
A amplitude é $M\omega I_0$ --- \textbf{proporcional à frequência}.\\
(c) \textbf{Por que é difícil eliminar.}
\begin{itemize}
\item $M$ depende do \textbf{logaritmo} da razão de distâncias. Afastar a espira
reduz a interferência muito devagar --- dobrar $d$ e $b$ juntos não muda
\emph{nada}.
\item A amplitude cresce com $\omega$: transitórios rápidos e harmônicos de alta
ordem acoplam-se muito mais que a fundamental.
\item Blindagem magnética é cara e pouco eficaz em baixa frequência, ao
contrário da blindagem elétrica.
\end{itemize}
\textbf{O que realmente funciona:} reduzir a \emph{área} da espira receptora
(cabos trançados, roteamento junto ao retorno) e reduzir a área da fonte (par
trançado no lado do fio). Atacar a geometria, não a distância.}
\stepcounter{questao}
\resposta{$\varepsilon=\SI{80}{\volt}$; $W=\SI{0.4}{\joule}$}
\resolucao{(a)
\[
|\varepsilon|=L\left|\frac{\Delta i}{\Delta t}\right|
=0{,}200\,\frac{2}{\pot{5}{-3}}=\SI{80}{\volt}.
\]
(b) $W=\frac12Li^{2}=\frac12(0{,}2)(4)=\SI{0.4}{\joule}$.\\
(c) \textbf{Risco.} A f.e.m. autoinduzida opõe-se à \emph{redução} da corrente
--- ela tenta mantê-la circulando. Se a interrupção for mais rápida, a tensão
sobe proporcionalmente: interromper em $\SI{5}{\micro\second}$ produziria
$\SI{80}{\kilo\volt}$.\\
Na prática, o ar rompe muito antes e forma-se \textbf{arco elétrico} nos
contatos, que os corrói e gera interferência.\\
\textbf{Proteção:} diodo de roda livre em antiparalelo, que oferece caminho à
corrente e limita a tensão. Os $\SI{0.4}{\joule}$ então se dissipam na
resistência do próprio circuito, sem arco.\\
\textbf{Ponto conceitual:} a autoindução é indução mútua da bobina
\emph{consigo mesma} --- a mesma lei de Faraday, aplicada ao próprio fluxo.}
\stepcounter{questao}
\resposta{$\tau=\SI{2.95}{\newton\meter}$; sem carga, o torque cai quase a
zero}
\resolucao{(a) A potência mecânica de entrada é
\[
P_{mec}=\frac{500}{0{,}90}=\SI{556}{\watt};
\qquad
\omega=\frac{1800(2\pi)}{60}=\SI{188.5}{\radian\per\second};
\]
\[
\tau=\frac{P_{mec}}{\omega}=\frac{556}{188{,}5}=\SI{2.95}{\newton\meter}.
\]
(b) \textbf{Sem carga}, não circula corrente no enrolamento. Sem corrente, não
há força magnética de oposição --- e o torque necessário cai a praticamente
zero (resta apenas o atrito e as perdas em vazio).\\
\textbf{É a lei de Lenz em ação.} O torque resistente \emph{surge} porque a
corrente induzida cria campo que se opõe ao movimento. Mais carga significa mais
corrente, mais oposição, mais torque exigido.\\
\textbf{Experiência reveladora:} girar um gerador manualmente com os terminais
abertos é fácil; curto-circuitá-los torna a rotação quase impossível. A energia
elétrica não vem "do campo" --- vem do braço de quem gira.}
\stepcounter{questao}
\resposta{Repulsão por corrente induzida; anel cortado não salta; em CC de
regime, nada acontece}
\resolucao{(a) \textbf{Explicação.} A corrente alternada no solenoide produz
fluxo variável através do anel, induzindo nele corrente. Pela lei de Lenz, essa
corrente cria campo que se \emph{opõe} ao do solenoide --- e dois campos opostos
se \textbf{repelem}.\\
Como a repulsão persiste em ambos os semiciclos (a corrente induzida inverte
junto com a indutora), a força média é sempre para cima.\\
\emph{(É o experimento clássico do "anel saltante" de Elihu Thomson.)}\\
(b) \textbf{Anel cortado:} não há caminho fechado, logo não circula corrente. Sem
corrente, não há campo induzido nem força. \textbf{O anel não salta} --- embora a
f.e.m. continue sendo induzida em seus terminais.\\
Este é o teste decisivo: mostra que a força depende da \emph{corrente}, não da
f.e.m.\\
(c) \textbf{Corrente contínua em regime:} o fluxo é constante, não há indução,
não há força. \textbf{Nada acontece.}\\
\emph{Mas no instante em que a corrente CC é \textbf{ligada}, há um transitório
de fluxo --- e o anel dá um salto único.} Isso reforça o ponto: é a
\emph{variação} que induz, não o campo.}
\stepcounter{questao}
\resposta{$\varepsilon=-\left(A\dfrac{\mathrm{d}B}{\mathrm{d}t}
+B\dfrac{\mathrm{d}A}{\mathrm{d}t}\right)$}
\resolucao{(a) Como $\Phi=BA$ com \emph{ambos} variáveis, aplica-se a regra do
produto:
\[
\varepsilon=-\frac{\mathrm{d}}{\mathrm{d}t}(BA)
=-\left(A\frac{\mathrm{d}B}{\mathrm{d}t}+B\frac{\mathrm{d}A}{\mathrm{d}t}\right).
\]
(b) \textbf{Duas contribuições de naturezas diferentes.}
\begin{itemize}
\item \textbf{Termo $A\,\mathrm{d}B/\mathrm{d}t$:} f.e.m. induzida por
\emph{campo variável} --- o mecanismo do transformador. Aqui,
$\mathrm{d}B/\mathrm{d}t=40\pi\cos(100\pi t)$, com amplitude
$\SI{125.7}{\tesla\per\second}$, dando até
$\SI{0.63}{\volt}$.
\item \textbf{Termo $B\,\mathrm{d}A/\mathrm{d}t$:} f.e.m. por \emph{movimento}
--- o mecanismo do gerador. Com
$\mathrm{d}A/\mathrm{d}t=\SI{-1e-3}{\meter\squared\per\second}$, contribui com
até $\SI{0.4}{\milli\volt}$.
\end{itemize}
\textbf{Neste caso o primeiro termo domina} por três ordens de grandeza.\\
\textbf{Ponto conceitual:} os dois mecanismos são fisicamente distintos --- um
envolve campo elétrico induzido, o outro força magnética sobre portadores --- mas
a lei de Faraday os unifica em uma única expressão. Essa unificação é uma das
observações que levaram Einstein à relatividade restrita.}
\stepcounter{questao}
\resposta{$\Delta\Phi=\SI{1.2}{\milli\weber}$; $q=\SI{24}{\milli\coulomb}$}
\resolucao{(a) Girar $\ang{180}$ inverte o sinal do fluxo:
\[
\Phi_i=BA=(0{,}3)(\pot{2}{-3})=\SI{0.6}{\milli\weber};
\qquad
\Phi_f=-\SI{0.6}{\milli\weber};
\]
\[
|\Delta\Phi|=\SI{1.2}{\milli\weber}\quad\text{--- o \textbf{dobro} de }
\Phi_i.
\]
(b)
\[
q=\frac{N|\Delta\Phi|}{R}=\frac{500(\pot{1.2}{-3})}{25}
=\pot{2.4}{-2}\,\si{\coulomb}.
\]
\textbf{Erro comum:} usar $\Delta\Phi=BA$ em vez de $2BA$. A inversão completa
faz o fluxo passar de $+\Phi$ a $-\Phi$, e a variação é o dobro do valor
inicial.\\
\textbf{Vantagem prática do giro de $\ang{180}$:} ele dobra o sinal em relação a
simplesmente retirar a bobina do campo --- técnica padrão em medições com
fluxímetro.}
\stepcounter{questao}
\resposta{$M=\SI{47.5}{\milli\henry}$; $\varepsilon_2=\SI{2.38}{\volt}$;
$5\%$ de dispersão}
\resolucao{(a)
\[
M=k\sqrt{L_1L_2}=0{,}95\sqrt{(0{,}100)(0{,}025)}
=0{,}95(0{,}05)=\SI{47.5}{\milli\henry}.
\]
(b)
\[
\varepsilon_2=M\frac{\mathrm{d}i_1}{\mathrm{d}t}=(0{,}0475)(50)
=\SI{2.38}{\volt}.
\]
(c) A fração acoplada é $k=95\%$; portanto \textbf{$5\%$ do fluxo se
dispersa}, não atravessando o secundário.\\
\textbf{Efeito da dispersão:} ela aparece como indutância em série
(\emph{indutância de dispersão}) e provoca queda de tensão sob carga. Um
transformador com $k=0{,}95$ tem regulação bem pior que um com
$k=0{,}999$.\\
\textbf{Comparação de topologias:} núcleos toroidais fechados atingem
$k>0{,}99$; núcleos em E com entreferro parasita ficam em torno de $0{,}95$;
transformadores com enrolamentos lado a lado, bem menos.}
\stepcounter{questao}
\resposta{$+0{,}3$; $0$; $\SI{-0.6}{\volt}$; área líquida nula}
\resolucao{(a)
\begin{center}
\begin{tabular}{lccc}
\hline
Trecho & $\mathrm{d}\Phi/\mathrm{d}t$ & $\varepsilon$ & Duração\\
\hline
Subida & $\SI{3}{\milli\weber\per\second}$ & $\SI{-0.3}{\volt}$
 & $\SI{2}{\second}$\\
Platô & $0$ & $0$ & $\SI{3}{\second}$\\
Descida & $\SI{-6}{\milli\weber\per\second}$ & $\SI{+0.6}{\volt}$
 & $\SI{1}{\second}$\\
\hline
\end{tabular}
\end{center}
(Os sinais decorrem de $\varepsilon=-N\,\mathrm{d}\Phi/\mathrm{d}t$.)\\
\textbf{Observe:} a descida é duas vezes mais rápida, e a f.e.m. é o
\emph{dobro} em módulo.\\
(b) \textbf{Área sob a curva:}
\[
(-0{,}3)(2)+(0)(3)+(0{,}6)(1)=-0{,}6+0{,}6=0.
\]
\textbf{Isso não é coincidência.} A área é $\int\varepsilon\,\mathrm{d}t
=-N\Delta\Phi_{total}$, e o fluxo terminou onde começou. \textbf{Sempre que o
fluxo retorna ao valor inicial, a integral da f.e.m. é nula.}\\
\textbf{Consequência prática:} em regime periódico, a tensão média sobre um
indutor é sempre zero --- princípio do \emph{balanço volt-segundo}, base do
projeto de conversores chaveados.}
\stepcounter{questao}
\resposta{$v=v_0e^{-t/\tau}$ com $\tau=\SI{0.25}{\second}$;
$d=\SI{2.5}{\meter}$}
\resolucao{(a)
\[
m\frac{\mathrm{d}v}{\mathrm{d}t}=-bv
\;\Longrightarrow\;
v(t)=v_0e^{-bt/m}.
\]
(b) $\tau=m/b=5/20=\SI{0.25}{\second}$.\\
A distância total é
\[
d=\int_0^{\infty}v_0e^{-t/\tau}\mathrm{d}t=v_0\tau=10(0{,}25)=\SI{2.5}{\meter}.
\]
\textbf{Mesma estrutura} da descarga exponencial de um capacitor --- e a mesma
propriedade: a "área" (distância) equivale à velocidade inicial mantida por
$\tau$.\\
(c) \textbf{Por que não para.} A força é proporcional a $v$: quando a velocidade
tende a zero, a força também tende a zero. Matematicamente, o corpo leva tempo
\emph{infinito} para parar --- ele apenas se aproxima assintoticamente do
repouso.\\
\textbf{Consequência prática decisiva:} freios magnéticos são excelentes para
\emph{desacelerar} e péssimos para \emph{segurar}. Trens de alta velocidade os
usam na faixa alta e completam a parada com freios de atrito --- que funcionam
justamente onde o magnético falha.}
\stepcounter{questao}
\resposta{Sinal oposto levaria a crescimento ilimitado de energia}
\resolucao{(a) \textbf{Demonstração por absurdo.} Suponha que a corrente
induzida \emph{reforçasse} a variação de fluxo. Ao aproximar um ímã de uma
espira, ela o \textbf{atrairia} --- acelerando a aproximação, o que aumentaria
$\mathrm{d}\Phi/\mathrm{d}t$, que aumentaria a corrente, que aumentaria a
atração.\\
O sistema entraria em fuga: energia cinética \emph{e} energia elétrica cresceriam
simultaneamente, sem fonte. \textbf{Contradição} com a conservação da energia.\\
Portanto a corrente deve \emph{opor-se} --- e o sinal negativo não é convenção,
é necessidade.\\
(b) \textbf{Quantificação.} Na barra em trilhos:
\[
P_{mec}=Fv=\left(\frac{B^{2}\ell^{2}v}{R}\right)v
=\frac{B^{2}\ell^{2}v^{2}}{R};
\qquad
P_{el}=\frac{(B\ell v)^{2}}{R}=\frac{B^{2}\ell^{2}v^{2}}{R}.
\]
\textbf{Idênticas, termo a termo.} A força de oposição tem exatamente o módulo
necessário para que a potência mecânica iguale a elétrica --- nem mais, nem
menos.\\
\textbf{Se a força fosse maior}, sobraria energia mecânica sem destino; se fosse
menor, apareceria energia elétrica do nada.\\
(c) \textbf{O que a lei de Lenz não determina.} Ela fixa apenas o \emph{sentido}
da corrente. O \textbf{módulo} vem da lei de Faraday combinada com a
resistência do circuito.\\
E ela não diz nada sobre \emph{quanta} energia é convertida --- isso depende da
carga. Uma espira aberta induz f.e.m. mas não corrente, e não há oposição
alguma: aproximar um ímã de uma espira cortada não exige força.}
\stepcounter{questao}
\resposta{$N=100$ espiras}
\resolucao{(a)
\[
N=\frac{\varepsilon\,\Delta t}{\Delta\Phi}
=\frac{5(\pot{4}{-3})}{\pot{2}{-4}}=100\ \text{espiras}.
\]
(b) Com carga de $\SI{10}{\kilo\ohm}$ e $\SI{5}{\volt}$, a corrente é
$\SI{0.5}{\milli\ampere}$ --- já abaixo do limite. A queda na resistência
interna deve ser pequena:
\[
\frac{R_{bobina}}{R_{bobina}+10\,\mathrm{k}}\le1\%
\;\Longrightarrow\;
R_{bobina}\le\SI{100}{\ohm}.
\]
Com $100$ espiras, isso é facilmente atingível --- fio de $\SI{0.2}{\milli\meter}$
e alguns metros de comprimento dão poucos ohms.\\
(c) \textbf{Resposta em frequência.} A bobina é um derivador: sua saída é
proporcional a $\mathrm{d}\Phi/\mathrm{d}t$, logo a sensibilidade
\textbf{cresce com a frequência}.
\begin{itemize}
\item \textbf{Vantagem:} excelente para detectar transitórios rápidos e
descargas.
\item \textbf{Limitação inferior:} ela \emph{não} detecta campo estático ---
para isso é preciso sensor Hall ou fluxgate.
\item \textbf{Limitação superior:} a indutância própria da bobina, associada à
capacitância entre espiras, forma um ressonador. Acima da frequência de
autorressonância a resposta deixa de ser proporcional a
$\mathrm{d}\Phi/\mathrm{d}t$.
\end{itemize}
\textbf{Para recuperar $\Phi$ (e não sua derivada):} acrescente um integrador na
saída --- é assim que funciona a bobina de Rogowski para medição de corrente.}
\stepcounter{questao}
\resposta{$F\propto\sigma B^{2}tv$; proporcionalidade a $v$ decorre de
$\varepsilon\propto v$ e $i\propto\varepsilon$}
\resolucao{(a) \textbf{Dependência.} A f.e.m. induzida localmente é
proporcional a $Bv$; a corrente parasita, a $\sigma Bv$ (com $\sigma$ a
condutividade); e a força, ao produto da corrente pelo campo:
\[
F\propto\sigma\,B^{2}\,t\,v\,A,
\]
com $t$ a espessura e $A$ a área sob o polo.\\
\textbf{Note a dependência quadrática em $B$} --- dobrar o campo quadruplica a
frenagem, o que torna ímãs de neodímio muito mais eficazes que ferrites.\\
(b) \textbf{Por que $F\propto v$.} Dois fatores lineares se combinam:
\[
\varepsilon\propto v
\;\Longrightarrow\;
i\propto v
\;\Longrightarrow\;
F=Bi\ell\propto v.
\]
É o mesmo mecanismo do arrasto viscoso --- e leva à mesma equação diferencial,
com decaimento exponencial da velocidade.\\
(c) \textbf{Velocidades muito altas.} A proporcionalidade \emph{deixa de valer}
por dois motivos:
\begin{itemize}
\item \textbf{Reação de armadura:} as correntes parasitas tornam-se intensas o
bastante para criar campo próprio comparável ao aplicado, distorcendo-o. A força
satura e chega a \emph{diminuir}.
\item \textbf{Efeito pelicular:} em alta frequência efetiva, as correntes se
concentram na superfície do disco, reduzindo a espessura útil $t$.
\end{itemize}
\textbf{Existe, portanto, uma velocidade de força máxima}, acima da qual é o freio
perde eficácia --- comportamento oposto ao do freio de atrito e que precisa ser
levado em conta no dimensionamento.}
\stepcounter{questao}
\resposta{Pontos marcam terminais magneticamente correspondentes; o sinal de
$M$ depende de as correntes entrarem ou não pelos pontos}
\resolucao{(a) \textbf{Convenção.} Correntes que entram pelos terminais
\emph{pontuados} produzem fluxos que se \textbf{reforçam} no núcleo comum.\\
Isso é informação \emph{física} (sentido de enrolamento), que o desenho
esquemático não transmite --- daí a necessidade da marcação.\\
(b) \textbf{Sinal nas equações.} Para duas bobinas acopladas:
\[
v_1=L_1\frac{\mathrm{d}i_1}{\mathrm{d}t}
\pm M\frac{\mathrm{d}i_2}{\mathrm{d}t};
\qquad
v_2=L_2\frac{\mathrm{d}i_2}{\mathrm{d}t}
\pm M\frac{\mathrm{d}i_1}{\mathrm{d}t}.
\]
Usa-se $+M$ se \emph{ambas} as correntes entram pelos pontos (ou ambas saem), e
$-M$ caso contrário.\\
\textbf{Em série}, isso dá $L_{eq}=L_1+L_2\pm2M$ --- resultado já obtido no banco
de indutores.\\
(c) \textbf{Procedimento experimental.} Meça $L_{eq}$ com as bobinas ligadas em
série nas \emph{duas} formas possíveis:
\[
M=\frac{L_{aditivo}-L_{subtrativo}}{4}.
\]
A ligação de \emph{maior} indutância identifica o acoplamento aditivo --- e
portanto os terminais correspondentes.\\
\textbf{Método alternativo, mais direto:} aplique um pulso de tensão a uma
bobina e observe a polaridade do pulso induzido na outra com osciloscópio. Se
ambas as pontas positivas estiverem nos terminais correspondentes, os sinais têm
a mesma polaridade.\\
\textbf{Por que isso importa na prática:} inverter a marcação de um
transformador em um conversor inverte a fase da saída --- e pode transformar
realimentação negativa em positiva, destruindo o circuito.}
\stepcounter{questao}
\resposta{Ímã caindo por tubo com bobinas de $N$ variável; verificar
linearidade em $N$ e em $v$}
\resolucao{(a) \textbf{Arranjo.} Um ímã permanente cai por dentro de um tubo
isolante vertical, atravessando bobinas montadas ao longo dele. Um osciloscópio
registra a f.e.m. de cada bobina.
\begin{enumerate}
\item \textbf{Variar $N$:} bobinas de $50$, $100$, $200$ e $400$ espiras, todas
com a mesma área e posição relativa. Compare os picos de f.e.m.
\item \textbf{Variar $v$:} use \emph{a mesma} bobina em alturas diferentes. Como
$v=\sqrt{2gh}$, a velocidade é conhecida a partir da altura de queda.
\end{enumerate}
(b) \textbf{Critério de validação.}
\begin{itemize}
\item Gráfico de $\varepsilon_{pico}$ contra $N$: deve ser uma \textbf{reta pela
origem}. Ajuste linear com $R^{2}>0{,}99$ e intercepto compatível com zero.
\item Gráfico de $\varepsilon_{pico}$ contra $\sqrt{h}$: também reta pela
origem.
\item \textbf{Teste mais forte:} a \emph{área} sob cada pulso de f.e.m. deve ser
$N\Delta\Phi$ --- \emph{independente} da velocidade. Verificar isso confirma
simultaneamente a lei e a independência da carga induzida.
\end{itemize}
(c) \textbf{Fontes de erro.}
\begin{itemize}
\item \textbf{Frenagem do próprio ímã} pelas correntes induzidas: se as bobinas
estiverem em curto ou a carga for baixa, o ímã desacelera e $v$ deixa de ser
$\sqrt{2gh}$. \emph{Solução:} medir em circuito aberto (osciloscópio de alta
impedância).
\item \textbf{Tubo condutor:} um tubo de alumínio ou cobre freia
drasticamente o ímã. Use acrílico ou vidro.
\item \textbf{Posição de disparo do osciloscópio} e largura de banda
insuficiente, que arredondam o pico.
\item \textbf{Orientação do ímã} variando durante a queda --- use guia.
\end{itemize}}
\stepcounter{questao}
\resposta{Fluxo constante $\Rightarrow$ $\mathrm{d}\Phi/\mathrm{d}t=0
\Rightarrow$ nenhuma f.e.m. no secundário}
\resolucao{(a) \textbf{Regime permanente.} Com corrente contínua constante, o
fluxo no núcleo é constante:
\[
\frac{\mathrm{d}\Phi}{\mathrm{d}t}=0
\;\Longrightarrow\;
\varepsilon_2=0.
\]
\textbf{Nenhuma tensão no secundário}, por maior que seja a corrente no
primário. A indução exige \emph{variação}.\\
(b) \textbf{Transitório de ligação.} No instante em que a fonte CC é ligada, o
fluxo sobe de zero ao valor final --- e \emph{há} f.e.m. no secundário durante
esse intervalo. Aparece um pulso, cuja integral é $N_2\Delta\Phi$.\\
O mesmo ocorre, com polaridade invertida, ao desligar.\\
\textbf{É esse pulso que se explora na bobina de ignição} de motores: uma
corrente contínua é interrompida abruptamente, e o pulso resultante atinge
dezenas de quilovolts.\\
(c) \textbf{O que ocorre com o primário --- e é o ponto crítico.} Em regime CC, a
única coisa que limita a corrente é a \textbf{resistência ôhmica} do
enrolamento, que é propositalmente pequena (poucos ohms).\\
Um primário projetado para $\SI{220}{\volt}$ em $\SI{60}{\hertz}$ tem sua
corrente limitada pela \emph{reatância} $\omega L$, tipicamente centenas de
ohms. Em CC, $\omega=0$ e essa limitação \textbf{desaparece}:
\[
I_{CC}=\frac{220}{R_{enrolamento}}\gg I_{CA}.
\]
\textbf{Resultado:} corrente dezenas de vezes maior que a nominal, saturação
imediata do núcleo e queima do enrolamento em segundos.\\
\textbf{Conclusão:} não é apenas que o transformador "não funciona" em CC --- ele
se \emph{destrói}.}
\stepcounter{questao}
\resposta{$\varepsilon=B\ell v$ pelos dois caminhos;
$F=\dfrac{B^{2}\ell^{2}v}{R}$}
\resolucao{(a) \textbf{Pela variação de área.} Com o lado dianteiro a uma
profundidade $x$ dentro do campo, $A_{imersa}=\ell x$ e
\[
\Phi=B\ell x
\;\Longrightarrow\;
\varepsilon=-\frac{\mathrm{d}\Phi}{\mathrm{d}t}=-B\ell\frac{\mathrm{d}x}{\mathrm{d}t}
=-B\ell v.
\]
\textbf{Pela força sobre portadores.} Os portadores do lado dianteiro movem-se
com velocidade $v$ e sofrem força $qvB$ ao longo do lado. O trabalho por unidade
de carga ao percorrer o comprimento $\ell$ é
\[
\varepsilon=\frac{W}{q}=vB\ell.
\]
\textbf{Os dois caminhos coincidem} --- e essa coincidência não é trivial: um
usa a lei de Faraday (variação de fluxo), o outro a força de Lorentz. Sua
equivalência é uma das simetrias mais notáveis do eletromagnetismo.\\
(b) \textbf{Força necessária.} A corrente $i=B\ell v/R$ percorre o lado
dianteiro, imerso no campo, sofrendo
\[
F=Bi\ell=\frac{B^{2}\ell^{2}v}{R},
\]
oposta ao movimento. É essa força que o agente externo deve equilibrar.\\
\textbf{Note:} os lados laterais também estão parcialmente no campo, mas suas
forças são opostas entre si e se cancelam.\\
(c) \textbf{Balanço.}
\[
P_{mec}=Fv=\frac{B^{2}\ell^{2}v^{2}}{R}
=\frac{(B\ell v)^{2}}{R}=\frac{\varepsilon^{2}}{R}=P_{el}.
\]
\textbf{Fecha exatamente.} E note o que acontece quando a espira fica
\emph{totalmente dentro}: os dois lados no campo geram f.e.m. opostas, a
corrente cessa, a força desaparece --- e a espira desliza livremente. A oposição
existe apenas enquanto o fluxo varia.}
\stepcounter{questao}
\resposta{$\bar v_L=0$; daí $V_1t_1=V_2t_2$ em regime}
\resolucao{(a) \textbf{Demonstração.} Em regime periódico de período $T$, a
corrente retorna ao mesmo valor a cada ciclo: $i(T)=i(0)$. Como
$v_L=L\,\mathrm{d}i/\mathrm{d}t$,
\[
\int_0^{T}v_L\,\mathrm{d}t=L\int_{i(0)}^{i(T)}\mathrm{d}i=L\left[i(T)-i(0)\right]=0.
\]
Logo $\bar v_L=0$. Equivalentemente, o \textbf{produto volt-segundo} aplicado
em um sentido deve igualar o aplicado no outro.\\
(b) \textbf{Aplicação a conversor.} Se o indutor recebe $+V_1$ durante
$t_1$ e $-V_2$ durante $t_2$:
\[
V_1t_1=V_2t_2.
\]
\textbf{Exemplo (conversor abaixador):} com entrada $\SI{12}{\volt}$, saída
$\SI{5}{\volt}$ e razão cíclica $D$, o indutor vê $(12-5)$ durante $DT$ e
$-5$ durante $(1-D)T$:
\[
7D=5(1-D)
\;\Longrightarrow\;
D=\frac{5}{12}=0{,}417.
\]
\textbf{O balanço volt-segundo produziu a relação de conversão} $V_o=DV_i$ sem
nenhuma análise de circuito --- é o método padrão de projeto.\\
(c) \textbf{Se o balanço não for respeitado.} A corrente do indutor não retorna
ao valor inicial: ela cresce (ou decresce) a cada ciclo, acumulando
indefinidamente.\\
Em um transformador, isso significa \textbf{saturação progressiva} do núcleo até
o colapso da indutância e a destruição do interruptor --- o fenômeno conhecido
como \emph{staircase saturation}.\\
\textbf{É por isso que conversores em ponte exigem simetria rigorosa} entre os
semiciclos, e por que se acrescenta capacitor de bloqueio em série com o
primário: ele impede qualquer componente contínua de tensão, garantindo o
balanço por construção.}
\stepcounter{questao}
\resposta{$\delta=\sqrt{\dfrac{2\rho}{\omega\mu}}\approx\SI{8.5}{\milli\meter}$
para cobre a $\SI{60}{\hertz}$}
\resolucao{(a) \textbf{Efeito pelicular.} A corrente alternada em um condutor
cria fluxo variável \emph{dentro dele}. Esse fluxo induz correntes parasitas
que, pela lei de Lenz, se opõem à corrente no centro e a reforçam na periferia.\\
Resultado: a corrente concentra-se junto à superfície, a seção efetiva diminui e
a \textbf{resistência aumenta}.\\
(b) \textbf{Profundidade de penetração.}
\[
\delta=\sqrt{\frac{2\rho}{\omega\mu}}.
\]
Para cobre ($\rho=\pot{1.7}{-8}$, $\mu\approx\mu_0$) a
$\SI{60}{\hertz}$:
\[
\delta=\sqrt{\frac{2(\pot{1.7}{-8})}{2\pi(60)(4\pi\times10^{-7})}}
\approx\SI{8.5}{\milli\meter}.
\]
\textbf{Interpretação:} condutores com raio muito menor que
$\SI{8.5}{\milli\meter}$ praticamente não sofrem o efeito em
$\SI{60}{\hertz}$. Já um barramento de $\SI{50}{\milli\meter}$ tem seu centro
essencialmente inútil.\\
\textbf{Dependência com a frequência:} $\delta\propto1/\sqrt{f}$. A
$\SI{1}{\mega\hertz}$, $\delta$ cai para cerca de
$\SI{66}{\micro\meter}$.\\
(c) \textbf{Soluções.}
\begin{itemize}
\item \textbf{Condutores tubulares} ou perfis chatos, que aproveitam melhor a
periferia --- comuns em barramentos de subestação.
\item \textbf{Fio Litz:} muitos fios finos isolados e \emph{transpostos}, de modo
que cada um ocupe alternadamente centro e periferia. Usado em enrolamentos de
alta frequência.
\item \textbf{Vários condutores em paralelo}, transpostos ao longo do percurso.
\end{itemize}
\textbf{Observação:} apenas colocar fios em paralelo \emph{sem} transposição não
resolve --- os externos conduziriam mais que os internos, pelo mesmo mecanismo.
É a transposição que equaliza.}
\stepcounter{questao}
\resposta{$v=\dfrac{E}{B}=\SI{1.4e6}{\meter\per\second}$,
\quad
$r=\dfrac{m_pv}{q_pB}\approx\SI{0.292}{\meter}$,
\quad
$2r\approx\SI{0.585}{\meter}$}
\resolucao{
(a) \textbf{Velocidade dos prótons no seletor.}
Para que o próton atravesse a região de campos cruzados sem sofrer desvio,
a força elétrica e a força magnética devem possuir o mesmo módulo e sentidos
opostos:
\[
F_E=F_B.
\]
Como
\[
F_E=qE
\qquad\text{e}\qquad
F_B=qvB,
\]
tem-se
\[
qE=qvB.
\]
Cancelando a carga $q$:
\[
v=\frac{E}{B}.
\]
Substituindo
\[
E=\SI{70}{\kilo\volt\per\meter}
\qquad\text{e}\qquad
B=\SI{50}{\milli\tesla}=\SI{0.050}{\tesla},
\]
resulta
\[
v=
\frac{70\times10^{3}}{0.050}
=
\boxed{\SI{1.4e6}{\meter\per\second}}.
\]

\textbf{Interpretação:} somente os prótons com essa velocidade atravessam o
seletor em linha reta. Prótons mais rápidos ou mais lentos sofrem uma força
resultante e são desviados.\\

(b) \textbf{Raio da trajetória na região magnética.}
Após deixar o seletor, o próton entra em uma região onde atua apenas o campo
magnético. Como a força magnética é sempre perpendicular à velocidade, ela não
altera o módulo da velocidade, mas muda continuamente sua direção, atuando como
força centrípeta:
\[
F_B=F_c.
\]
Portanto,
\[
q_pvB=\frac{m_pv^2}{r}.
\]
Isolando o raio:
\[
r=\frac{m_pv}{q_pB}.
\]
Para
\[
m_p=\SI{1.67e-27}{\kilogram},
\qquad
q_p=\SI{1.60e-19}{\coulomb},
\]
tem-se
\[
r=
\frac{
(\pot{1.67}{-27})(\pot{1.4}{6})
}{
(\pot{1.60}{-19})(0.050)
}
\approx
\boxed{\SI{0.292}{\meter}}.
\]
Assim,
\[
r\approx\SI{29.2}{\centi\meter}.
\]

(c) \textbf{Distância até a placa detectora.}
A trajetória mostrada corresponde a uma semicircunferência. Dessa forma, a
distância vertical entre o ponto de entrada na região magnética e o ponto de
impacto na placa detectora corresponde ao diâmetro da trajetória:
\[
d=2r.
\]
Logo,
\[
d=2(0.292)
\approx
\boxed{\SI{0.585}{\meter}}
\]
ou, equivalentemente,
\[
\boxed{d\approx\SI{58.5}{\centi\meter}}.
\]

\textbf{Observação:} o seletor de velocidades utiliza a compensação entre as
forças elétrica e magnética para selecionar partículas com uma velocidade
específica. Na segunda região, a presença apenas do campo magnético permite
determinar a razão entre momento linear e carga da partícula a partir da
curvatura de sua trajetória.
}
\stepcounter{questao}
\resposta{$v\approx\SI{2.52e7}{\meter\per\second}$,
\quad
$r\approx\SI{5.72}{\milli\meter}$,
\quad
$T\approx\SI{1.43}{\nano\second}$,
\quad
$t_{1/4}\approx\SI{0.357}{\nano\second}$}
\resolucao{
(a) \textbf{Velocidade após a aceleração.}
A energia potencial elétrica adquirida pelo elétron transforma-se em energia
cinética:
\[
|q_e|\Delta V=\frac{1}{2}m_ev^2.
\]
Portanto,
\[
v=\sqrt{\frac{2|q_e|\Delta V}{m_e}}.
\]
Substituindo os valores:
\[
v=
\sqrt{
\frac{
2(1.60\times10^{-19})(1.8\times10^3)
}{
9.11\times10^{-31}
}}
\approx
\boxed{\SI{2.52e7}{\meter\per\second}}.
\]

(b) \textbf{Raio da trajetória.}
Como $\vec v\perp\vec B$, a força magnética atua como força centrípeta:
\[
|q_e|vB=\frac{m_ev^2}{r}.
\]
Assim,
\[
r=\frac{m_ev}{|q_e|B}.
\]
Logo,
\[
r=
\frac{(9.11\times10^{-31})(2.52\times10^7)}
{(1.60\times10^{-19})(25\times10^{-3})}
\approx
\boxed{\SI{5.72}{\milli\meter}}.
\]

(c) \textbf{Período do movimento.}
\[
T=\frac{2\pi m_e}{|q_e|B}.
\]
Portanto,
\[
T\approx
\boxed{\SI{1.43e-9}{\second}}
=
\boxed{\SI{1.43}{\nano\second}}.
\]

(d) Para um quarto de circunferência:
\[
t_{1/4}=\frac{T}{4}.
\]
Assim,
\[
t_{1/4}\approx
\boxed{\SI{3.57e-10}{\second}}
=
\boxed{\SI{0.357}{\nano\second}}.
\]

\textbf{Interpretação:} a diferença de potencial determina a energia cinética
da partícula, enquanto o campo magnético modifica apenas a direção da
velocidade, produzindo uma trajetória circular.
}
\stepcounter{questao}
\resposta{$v=\SI{1.50e6}{\meter\per\second}$,
\quad
$r_{20}\approx\SI{0.889}{\meter}$,
\quad
$r_{22}\approx\SI{0.977}{\meter}$,
\quad
$\Delta d\approx\SI{0.178}{\meter}$}
\resolucao{
(a) \textbf{Velocidade selecionada.}
No seletor:
\[
qE=qvB_s.
\]
Logo,
\[
v=\frac{E}{B_s}.
\]
Assim,
\[
v=
\frac{30\times10^3}{20\times10^{-3}}
=
\boxed{\SI{1.50e6}{\meter\per\second}}.
\]

(b) \textbf{Raios das trajetórias.}
Para um íon de massa $m$:
\[
r=\frac{mv}{qB_a}.
\]

Para o $^{20}\mathrm{Ne}^{+}$:
\[
m_{20}=20u=20(1.66\times10^{-27}),
\]
resultando em
\[
r_{20}\approx
\boxed{\SI{0.889}{\meter}}.
\]

Para o $^{22}\mathrm{Ne}^{+}$:
\[
m_{22}=22u,
\]
e
\[
r_{22}\approx
\boxed{\SI{0.977}{\meter}}.
\]

(c) Como cada partícula descreve uma semicircunferência, a posição de impacto
é determinada pelo diâmetro:
\[
d=2r.
\]
A separação entre os impactos é
\[
\Delta d=2(r_{22}-r_{20}).
\]
Logo,
\[
\Delta d
=
2(0.977-0.889)
\approx
\boxed{\SI{0.178}{\meter}}.
\]

(d) Como
\[
r=\frac{mv}{qB},
\]
mantidos $v$, $q$ e $B$, o raio é diretamente proporcional à massa. Portanto,
o isótopo
\[
\boxed{^{22}\mathrm{Ne}^{+}}
\]
atinge a região mais distante.

\textbf{Interpretação:} o espectrômetro converte diferenças de massa em
diferenças geométricas mensuráveis na posição de impacto das partículas.
}
\stepcounter{questao}
\resposta{$v_{\parallel}\approx\SI{5.36e5}{\meter\per\second}$,
\quad
$v_{\perp}\approx\SI{4.50e5}{\meter\per\second}$,
\quad
$r\approx\SI{5.18}{\centi\meter}$,
\quad
$T\approx\SI{0.724}{\micro\second}$,
\quad
$p\approx\SI{0.388}{\meter}$}
\resolucao{
(a) \textbf{Decomposição da velocidade.}
A componente paralela ao campo é
\[
v_{\parallel}=v\cos\theta,
\]
e a componente perpendicular:
\[
v_{\perp}=v\sin\theta.
\]

Logo,
\[
v_{\parallel}
=
(7.0\times10^5)\cos40^\circ
\approx
\boxed{\SI{5.36e5}{\meter\per\second}},
\]
e
\[
v_{\perp}
=
(7.0\times10^5)\sin40^\circ
\approx
\boxed{\SI{4.50e5}{\meter\per\second}}.
\]

(b) \textbf{Raio da hélice.}
Somente $v_{\perp}$ produz movimento circular:
\[
q_\alpha v_{\perp}B
=
\frac{m_\alpha v_{\perp}^2}{r}.
\]
Portanto,
\[
r=
\frac{m_\alpha v_{\perp}}{q_\alpha B}.
\]
Substituindo:
\[
r\approx
\boxed{\SI{5.18}{\centi\meter}}.
\]

(c) \textbf{Período.}
\[
T=
\frac{2\pi m_\alpha}{q_\alpha B}.
\]
Assim,
\[
T\approx
\boxed{\SI{7.24e-7}{\second}}
=
\boxed{\SI{0.724}{\micro\second}}.
\]

(d) \textbf{Passo da hélice.}
Durante uma volta, o deslocamento paralelo ao campo é
\[
p=v_{\parallel}T.
\]
Logo,
\[
p=
(5.36\times10^5)(7.24\times10^{-7})
\approx
\boxed{\SI{0.388}{\meter}}.
\]

(e) Para $\theta=90^\circ$:
\[
v_{\parallel}=0.
\]
Consequentemente, desaparece o avanço longitudinal da hélice e a trajetória
passa a ser
\[
\boxed{\text{circular}}.
\]

\textbf{Interpretação:} o movimento helicoidal resulta da superposição de um
movimento retilíneo uniforme paralelo ao campo com um movimento circular
uniforme perpendicular ao campo.
}
\stepcounter{questao}
\resposta{$f\approx\SI{12.2}{\mega\hertz}$,
\quad
$T\approx\SI{82.0}{\nano\second}$,
\quad
$v_{\max}\approx\SI{1.92e7}{\meter\per\second}$,
\quad
$K_{\max}\approx\SI{3.07e-13}{\joule}
\approx\SI{1.92}{\mega\electronvolt}$}
\resolucao{
(a) \textbf{Frequência ciclotrônica.}
A força magnética fornece a força centrípeta:
\[
q_pvB=\frac{m_pv^2}{r}.
\]
Como
\[
v=\omega r,
\]
obtém-se
\[
\omega=\frac{q_pB}{m_p}.
\]
Portanto,
\[
f=\frac{\omega}{2\pi}
=
\frac{q_pB}{2\pi m_p}.
\]
Substituindo:
\[
f=
\frac{
(1.60\times10^{-19})(0.80)
}{
2\pi(1.67\times10^{-27})
}
\approx
\boxed{\SI{12.2}{\mega\hertz}}.
\]

(b) \textbf{Período.}
\[
T=\frac{1}{f}.
\]
Logo,
\[
T\approx
\boxed{\SI{82.0}{\nano\second}}.
\]

(c) \textbf{Velocidade máxima.}
Da condição de movimento circular:
\[
r=\frac{m_pv}{q_pB}.
\]
Para $r=R$:
\[
v_{\max}=
\frac{q_pBR}{m_p}.
\]
Assim,
\[
v_{\max}
=
\frac{
(1.60\times10^{-19})(0.80)(0.25)
}{
1.67\times10^{-27}
}
\approx
\boxed{\SI{1.92e7}{\meter\per\second}}.
\]

(d) \textbf{Energia cinética máxima.}
\[
K_{\max}
=
\frac{1}{2}m_pv_{\max}^2.
\]
Substituindo:
\[
K_{\max}
\approx
\boxed{\SI{3.07e-13}{\joule}}.
\]

Como
\[
\SI{1}{\electronvolt}
=
\SI{1.60e-19}{\joule},
\]
tem-se
\[
K_{\max}
=
\frac{3.07\times10^{-13}}
{1.60\times10^{-19}}
\approx
1.92\times10^6\,\mathrm{eV}.
\]
Portanto,
\[
\boxed{K_{\max}\approx\SI{1.92}{\mega\electronvolt}}.
\]

\textbf{Interpretação:} a frequência ciclotrônica permanece aproximadamente
constante enquanto os efeitos relativísticos forem desprezíveis. A cada
passagem pela região entre os eletrodos, o campo elétrico fornece energia à
partícula, aumentando sua velocidade e, consequentemente, o raio da órbita.
}

\gabarito[12.5 Lei de Lenz e indução eletromagnética]
